% ============================================================
%  DISPENSA DI PROGRAMMAZIONE
%  Universita' degli Studi di Napoli Federico II
%  Compilare con LuaLaTeX (necessario: fontspec + luacode)
% ============================================================
\documentclass{unina_doc_class}

\gdef\documentname{Programmazione}
\gdef\documentyear{A.A. \ 2024-2025}
\gdef\author{Fabio Rosiello}
\pagestyle{plain}

\gdef\xauthors{\hi{\iuser Fabio Rosiello}}

% Didascalie delle figure
\usepackage[font=small,labelfont={bf,color=primary!80!black},
            textfont=it,skip=6pt]{caption}

% Preambolo specifico della dispensa (font, listati, box, stili TikZ)
% ============================================================
%  Preambolo aggiuntivo per la dispensa di Programmazione
%  Estende unina_doc_class.cls (stile ripreso dalle Dispense APA e
%  dalla dispensa di Tecnologie Web, di cui riusa listati e riquadri)
% ============================================================

% Tabelle con i filetti di booktabs (\toprule, \midrule, \bottomrule)
\usepackage{booktabs}

% ─── Font aggiuntivi ─────────────────────────────────────────
% \Mononoki e' richiamato dalla classe nelle macro \dir e \file
\newfontfamily\Mononoki{mononoki}[
  Path = fonts/, Extension = .ttf,
  UprightFont = *-Regular, BoldFont = *-Bold,
  ItalicFont  = *-Italic,  BoldItalicFont = *-BoldItalic]

% FiraCode SENZA legature di programmazione.
% La classe dichiara \FiraCode con Contextuals=Alternate: in un testo
% sul C questo e' dannoso, perche' "!=" verrebbe disegnato come "≠", "->"
% come una freccia e "<=" come "≤", cioe' non come cio' che lo studente
% deve digitare. Qui si ridefinisce \FiraCode senza legature.
\newfontfamily\FiraCodeNL{FiraCode}[
  Path = fonts/, Extension = .ttf,
  UprightFont = *-Regular, BoldFont = *-Bold,
  Ligatures = NoCommon, Contextuals = NoAlternate]
\let\FiraCode\FiraCodeNL

% Nessun "$" deve aprire la modalita' matematica dentro un listato
\lstset{mathescape=false}

% ─── Etichette richieste dalla classe per l'elenco dei listati ─
\gdef\documentlsttitle{Elenco dei listati}
\gdef\documentlstlabel{Listato}

% ─── Colori per il codice su sfondo chiaro (figure e riquadri) ─
\definecolor{codestring}{HTML}{1F7A5A}
\definecolor{codecomment}{HTML}{8A8A8A}
\definecolor{codekey}{HTML}{C2185B}

% ============================================================
%  ASPETTO DEI LISTATI: UNA FINESTRA DI VS CODE (TEMA DARK+)
% ============================================================
% Ogni listato e' disegnato come una piccola finestra dell'editor Visual
% Studio Code con il tema scuro: sfondo quasi nero, numeri di riga grigi a
% sinistra e, se il listato ha un nome, la barra delle schede in alto con
% la scheda del file aperto (icona del linguaggio, nome del file e la
% crocetta di chiusura). E' cio' che lo studente vede sul proprio schermo.
\definecolor{vscbg}{HTML}{1E1E1E}        % sfondo dell'editor
\definecolor{vsctabbar}{HTML}{252526}    % barra delle schede
\definecolor{vscfg}{HTML}{D4D4D4}        % testo normale, punteggiatura
\definecolor{vsclinenum}{HTML}{858585}   % numeri di riga
\definecolor{vscblue}{HTML}{569CD6}      % tipi base, modificatori, macro
\definecolor{vscpurple}{HTML}{C586C0}    % if, for, return, #include...
\definecolor{vscstring}{HTML}{CE9178}    % stringhe e caratteri
\definecolor{vsccomment}{HTML}{6A9955}   % commenti
\definecolor{vscfunc}{HTML}{DCDCAA}      % nomi di funzione
\definecolor{vscvar}{HTML}{9CDCFE}       % variabili e parametri
\definecolor{vsctype}{HTML}{4EC9B0}      % tipi definiti (typedef, struct)
\definecolor{vscnum}{HTML}{B5CEA8}       % numeri
\definecolor{vscaccent}{HTML}{0078D4}    % filetto della scheda attiva
\definecolor{vsctabtext}{HTML}{FFFFFF}   % nome del file nella scheda
\definecolor{vsctabdim}{HTML}{9D9D9D}    % crocetta di chiusura
% colori delle icone dei file (tema di icone Seti, quello predefinito)
\definecolor{vsciconorange}{HTML}{E37933}
\definecolor{vsciconblue}{HTML}{519ABA}
\definecolor{vsciconyellow}{HTML}{CBCB41}
\definecolor{vscicongreen}{HTML}{8DC149}
\definecolor{vsciconpurple}{HTML}{A074C4}
\definecolor{vsciconred}{HTML}{CC3E44}

% ============================================================
%  LINGUAGGI PER LISTINGS
% ============================================================
% Listings non conosce la grammatica del C: colora le parole che compaiono
% in elenchi fissi. Le classi di parole chiave riproducono le categorie del
% tema Dark+ di VS Code (con l'estensione C/C++ installata):
%   [1] blu       tipi base, modificatori, sizeof e macro (NULL, EOF, PI...)
%   [2] viola     controllo del flusso (if, for, return...) e direttive (#include)
%   [3] giallo    funzioni chiamate o definite negli esempi della dispensa
%   [4] verde     tipi definiti con typedef e nomi delle struct
% Tutti gli altri identificatori sono variabili e prendono l'azzurro chiaro.
\lstset{keywords={}}

\lstdefinelanguage{C99}{
  sensitive=true,
  alsoletter={\#},
  morecomment=[l]{//},
  morecomment=[s]{/*}{*/},
  morestring=[b]",
  morestring=[b]',
  keywords=[1]{
    int,char,float,double,void,short,long,signed,unsigned,struct,union,
    enum,typedef,const,static,extern,register,auto,volatile,sizeof,inline,
    restrict,_Bool,bool,true,false,
    % macro di libreria e degli esempi
    NULL,EOF,PI,CLASS_SIZE,RAND_MAX,INT_MAX,INT_MIN,CHAR_MAX,CHAR_MIN,
    LONG_MAX,UINT_MAX,SHRT_MAX,EXIT_SUCCESS,EXIT_FAILURE,FLT_EPSILON,
    DBL_EPSILON,LIMIT,MAXWORD,N,M,DIM,MAX,SQ,SQUARE,DEBUG,BUFSIZ,
    STUDENTE_H,SORT_H,AREE_H,BIG,PAUSA,SCHAR_MIN,SCHAR_MAX,UCHAR_MAX,
    SHRT_MIN,LONG_MIN,ULONG_MAX,FLT_MIN,FLT_MAX,FLT_DIG,DBL_MIN,DBL_MAX,
    DBL_DIG,INITIAL_SEED,MULTIPLIER,INCREMENT,MODULUS,FLOATING_MODULUS},
  keywords=[2]{
    if,else,for,while,do,switch,case,default,break,continue,return,goto,
    \#include,\#define,\#ifndef,\#ifdef,\#endif,\#if,\#else,\#elif,\#undef,
    \#pragma},
  keywords=[3]{
    main,printf,scanf,getchar,putchar,puts,gets,fgets,fputs,fprintf,
    fscanf,sprintf,sscanf,fopen,fclose,malloc,calloc,realloc,free,exit,
    abs,fabs,sqrt,pow,sin,cos,exp,log,floor,ceil,rand,srand,time,
    strlen,strcpy,strncpy,strcat,strncat,strcmp,strncmp,strchr,strstr,
    strcspn,strspn,strdup,memcpy,memset,isalpha,isdigit,isspace,isupper,
    islower,isalnum,ispunct,toupper,tolower,assert,feof,qsort,
    % funzioni degli esempi
    max,min,maxmin,prn_info,area,area_cerchio,area_quadrato,perimetro,
    fattoriale,factorial,potenza,power,swap,sum,somma,sum_array,media,
    stampa,stampa_array,leggi,leggi_array,fail,update,inc,incrementa,
    wrt_address,twice,all_add,compute_sum,f,g,h,sf,count,counter,
    random_seed,my_random,ordina,bubble_sort,selection_sort,bubble,
    sort_words,wrt_words,error_exit,word_cnt,my_strlen,my_strcpy,
    my_strcat,crea_matrice,libera_matrice,get_matrix_space,
    release_matrix_space,fill_matrix,wrt_matrix,
    nuovoNodo,inserisciTesta,inserisciCoda,inserisciPosizione,
    inserisciOrdinato,cerca,lunghezza,eliminaValore,eliminaPosizione,
    eliminaTesta,eliminaCoda,unisci,stampaLista,liberaLista,
    leggiStudenti,stampaStudenti,inserisci,rimuovi,inserisciNodo,
    rimuoviNodo,stampaLista,cancellaLista,cnt_char,capitalize,
    trial,mult,prova,converti,fahrenheit,celsius,cubo,quadrato,
    leggi_intero,conta,scambia,inverti,cerca_massimo,mcd,primo,
    stampa_matrice,somma_matrici,somma_righe,trasponi,setbuf,sleep,pausa,
    errore,nothing,absolute_value,probability,random},
  keywords=[4]{
    size_t,FILE,ptrdiff_t,int8_t,int16_t,int32_t,int64_t,uint8_t,uint32_t,
    uint64_t,carta,complex,student,date,department,address,employee_data,
    fruit,vegetable,card,Node,Nodo,studente,Studente,lista,Lista,punto,
    Punto,persona,Persona,matrice,elemento,elem,frazione,data,
    stack,queue,Coda,Pila,uppercase,inches,u_char,tipoInt},
}

% ─── Pseudocodice ────────────────────────────────────────────
% Le parole chiave (in maiuscolo) sono viola come le istruzioni di
% controllo del C; i testi fra virgolette sono stringhe.
\lstdefinelanguage{Pseudo}{
  sensitive=true,
  morestring=[b]",
  morecomment=[l]{//},
  keywords=[1]{BEGIN,END,LEGGI,STAMPA,SE,ALLORA,ALTRIMENTI,MENTRE,
               FINCHE,RIPETI,PER,DA,A,ESEGUI,FINE,RESTITUISCI},
  keywords=[2]{PI,AND,OR,NOT},
}

% ─── Comandi del terminale e script di shell ─────────────────
% I comandi sono gialli, come le funzioni (VS Code fa lo stesso).
\lstdefinelanguage{Shell}{
  sensitive=true,
  morecomment=[l]{\#},
  morestring=[b]",
  morestring=[b]',
  morekeywords={gcc,cc,clang,make,cd,ls,mkdir,cat,echo,cp,mv,rm,pwd,
                sudo,apt,wsl,code,chmod,touch,man,clear,gdb,diff,
                xcode-select,bash,sh,history,less,head,tail},
}

% ─── Makefile ────────────────────────────────────────────────
% Commenti verdi, variabili $(...) azzurre, obiettivi gialli. Il TAB
% obbligatorio prima dei comandi si scrive (*@\TAB@*) ed e' disegnato
% come una freccia, come fa VS Code mostrando gli spazi bianchi.
\lstdefinelanguage{Make}{
  sensitive=true,
  alsoletter={.-},
  morecomment=[l]{\#},
  morestring=[b]",
  moredelim=**[s][\color{vscvar}]{\$(}{)},
  keywords=[1]{CC,CFLAGS,OBJ,OBJS,LDLIBS,LDFLAGS,TARGET,SRC},
  keywords=[2]{gcc,rm,echo,cc},
}

% ============================================================
%  STILI DI LISTATO
% ============================================================
% Ogni stile sceglie anche l'icona mostrata nella scheda (chiave vscicon,
% definita piu' sotto).
% I nomi delle intestazioni standard dopo #include (<stdio.h>...) sono
% stringhe per VS Code: li si colora con delle sostituzioni ("literate"),
% che con l'asterisco non agiscono dentro commenti e stringhe.
\lstdefinestyle{c}{
  language=C99,
  vscicon=c,
  keywordstyle=[1]\color{vscblue},
  keywordstyle=[2]\color{vscpurple},
  keywordstyle=[3]\color{vscfunc},
  keywordstyle=[4]\color{vsctype},
  identifierstyle=\color{vscvar},
  stringstyle=\color{vscstring},
  commentstyle=\color{vsccomment},
  literate=*%
    {<stdio.h>}{{\textcolor{vscstring}{<stdio.h>}}}9
    {<stdlib.h>}{{\textcolor{vscstring}{<stdlib.h>}}}{10}
    {<math.h>}{{\textcolor{vscstring}{<math.h>}}}8
    {<string.h>}{{\textcolor{vscstring}{<string.h>}}}{10}
    {<ctype.h>}{{\textcolor{vscstring}{<ctype.h>}}}9
    {<limits.h>}{{\textcolor{vscstring}{<limits.h>}}}{10}
    {<float.h>}{{\textcolor{vscstring}{<float.h>}}}9
    {<time.h>}{{\textcolor{vscstring}{<time.h>}}}8
    {<stdbool.h>}{{\textcolor{vscstring}{<stdbool.h>}}}{11}
    {<assert.h>}{{\textcolor{vscstring}{<assert.h>}}}{10}
    {<unistd.h>}{{\textcolor{vscstring}{<unistd.h>}}}{10},
}

% file di intestazione (.h): stessi colori, icona diversa
\lstdefinestyle{h}{
  style=c,
  vscicon=h,
}

\lstdefinestyle{shell}{
  language=Shell,
  vscicon=shell,
  keywordstyle=\color{vscfunc},
  stringstyle=\color{vscstring},
  commentstyle=\color{vsccomment},
}

\lstdefinestyle{make}{
  language=Make,
  vscicon=make,
  keywordstyle=[1]\color{vscvar},
  keywordstyle=[2]\color{vscfunc},
  stringstyle=\color{vscstring},
  commentstyle=\color{vsccomment},
}

\lstdefinestyle{plain}{
  language={},
  vscicon=plain,
}

\lstdefinestyle{pseudo}{
  language=Pseudo,
  vscicon=plain,
  keywordstyle=[1]\color{vscpurple},
  keywordstyle=[2]\color{vscblue},
  identifierstyle=\color{vscvar},
  stringstyle=\color{vscstring},
  commentstyle=\color{vsccomment},
}

% ============================================================
%  LA FINESTRA DELL'EDITOR
% ============================================================
\makeatletter
% Chiave vscicon: l'icona da mettere nella scheda (la sceglie lo stile)
\lst@Key{vscicon}{plain}{\def\wtvsc@icon{#1}}

% Il codice e' in corpo 7,5: la colonna dei numeri di riga toglie spazio
% orizzontale, e con il corpo leggermente ridotto le righe stanno in una
% riga anche nelle colonne strette di "affianco".
\newcommand{\wtvsccodefont}{\fontsize{7.5}{9}\selectfont\FiraCodeNL\color{vscfg}}
\newcommand{\wtvsccodefontsmall}{\fontsize{6.6}{7.6}\selectfont\FiraCodeNL\color{vscfg}}
\lstset{
  basicstyle=\linespread{1.15}\wtvsccodefont,
  % literate={} annulla le sostituzioni dello stile FiraCodeStyle (pacchetto
  % lstfiracode, caricato dalla classe): servono a tenere unite le coppie di
  % caratteri che Fira Code disegna come legature, ma cosi' listings non
  % riconosce piu' i delimitatori che le contengono (per esempio "//" e "/*"
  % dei commenti). Le legature sono comunque disattivate nel font.
  literate={},
  backgroundcolor=\color{vscbg},
  identifierstyle={},
  % numeri di riga, come nella "gutter" di VS Code
  numbers=left,
  numberstyle=\fontsize{5.5}{6}\selectfont\FiraCodeNL\color{vsclinenum},
  numbersep=5pt,
  % ogni listato riparte da 1
  firstnumber=1,
  % la cornice e' invisibile (spessore nullo): serve solo a estendere lo
  % sfondo scuro attorno al codice e sotto i numeri di riga
  frame=single,
  framerule=0pt,
  rulecolor=\color{vscbg},
  % il "titolo" del listato diventa la barra delle schede, sopra il codice
  captionpos=t,
  abovecaptionskip=0pt,
  belowcaptionskip=0pt,
  % un TAB vale 4 spazi, come in VS Code
  tabsize=4,
}

% Margini della finestra: la finestra va da #1 dal bordo sinistro della
% riga fino a #2 dal bordo destro. Dentro, a sinistra ci sono 16pt per la
% colonna dei numeri di riga (12pt) e lo spazio prima del codice (4pt);
% a destra 4pt di spazio.
\newcommand{\wtvscL}{0pt}
\newcommand{\wtvscR}{0pt}
\newcommand{\wtvscmargini}[2]{%
  \renewcommand{\wtvscL}{#1}\renewcommand{\wtvscR}{#2}%
  \lstset{xleftmargin=\dimexpr#1+16pt\relax, framexleftmargin=12pt,
          framesep=4pt,
          xrightmargin=\dimexpr#2+4pt\relax, framexrightmargin=0pt}}
\wtvscmargini{0pt}{0pt}

% Icone dei file nella scheda (riprendono il tema di icone Seti di VS Code)
\newcommand{\wtvsc@badge}[2]{%
  {\FiraCodeNL\bfseries\fontsize{6.5}{7}\selectfont\color{#1}#2}}
\newcommand{\wtvsc@iconmark}{%
  \ifcsname wtvsc@icon@\wtvsc@icon\endcsname
    \csname wtvsc@icon@\wtvsc@icon\endcsname
  \else
    \wtvsc@icon@plain
  \fi}
\def\wtvsc@icon@c{\wtvsc@badge{vsciconblue}{C}}
\def\wtvsc@icon@h{\wtvsc@badge{vsciconpurple}{h}}
\def\wtvsc@icon@make{\wtvsc@badge{vsciconorange}{M}}
\def\wtvsc@icon@shell{\wtvsc@badge{vscicongreen}{>\_}}
\def\wtvsc@icon@plain{{\fontsize{7}{7}\selectfont\color{vsctabdim}\ifilecode}}

% La barra delle schede: grigio scuro con gli angoli superiori arrotondati,
% e a sinistra la scheda attiva (dello stesso colore dell'editor, con il
% filetto blu in alto) che contiene icona, nome del file e crocetta.
\newlength{\wtvsc@barw}
\newcommand{\wtvsc@tabbar}{%
  \setlength{\wtvsc@barw}{\dimexpr\linewidth-\wtvscL-\wtvscR\relax}%
  \par\noindent\hspace*{\wtvscL}%
  \begin{tikzpicture}[baseline=0pt]
    \useasboundingbox (0,0) rectangle (\wtvsc@barw,17pt);
    \clip[rounded corners=3pt] (0,-5pt) rectangle (\wtvsc@barw,17pt);
    \fill[vsctabbar] (0,0) rectangle (\wtvsc@barw,17pt);
    \node[anchor=south west, inner xsep=7pt, inner ysep=0pt,
          minimum height=17pt, fill=vscbg, text=vsctabtext,
          font=\Lato\fontsize{7.5}{9}\selectfont] (tab) at (0,0)
      {\wtvsc@iconmark\hspace{4pt}\lstname\hspace{9pt}%
       {\color{vsctabdim}\fontsize{8}{9}\selectfont\texttimes}};
    \fill[vscaccent] (tab.north west) rectangle ([yshift=-1pt]tab.north east);
  \end{tikzpicture}%
  \par\nointerlineskip}

% La barra delle schede compare solo se il listato ha un nome (chiave
% "name"): se manca, il titolo viene svuotato appena prima di disegnarlo.
\lstset{title={\lstname}, everydisplay={\wtvsc@solosenome}}
\newcommand{\wtvsc@solosenome}{%
  \ifx\lst@intname\@empty
    \let\lst@caption\@empty
  \else
    \let\lst@caption\relax
  \fi}
% Al posto della didascalia si disegna la barra delle schede. La ridefinizione
% va fatta a inizio documento perche' il pacchetto caption, caricato prima,
% sostituisce anch'esso \lst@maketitle in quel momento.
\AtBeginDocument{\def\lst@maketitle#1{\wtvsc@tabbar}}
% Nel nome del file listings sostituisce "-" con un trattino lungo: nella
% scheda deve restare il trattino vero, quello del nome del file.
\def\lst@filenamerpl{_\textunderscore $\textdollar}

% Il TAB obbligatorio dei makefile, reso visibile: una freccia larga quanto
% quattro colonne del listato (\lst@width), cioe' quanto il TAB stesso. Si
% scrive (*@\TAB@*) all'inizio della riga di comando.
\newcommand{\TAB}{\makebox[4\lst@width][l]{%
  \color{vsclinenum}\tikz[baseline=-2.6pt]\draw[-{Stealth[length=3pt,width=3pt]},
    line width=0.5pt] (0.4pt,0) -- (3.4*\lst@width,0);}}
\makeatother

% ============================================================
%  CORREZIONE DEI RIQUADRI DELLA CLASSE
% ============================================================
% Nella classe il titolo viene passato a tcolorbox come "title=#1": una
% virgola nel titolo verrebbe interpretata come separatore di chiavi.
% Racchiudendo il titolo in un gruppo il problema sparisce.
\renewcommand{\info}[2]{\begin{infobox}{{#1}}#2\end{infobox}}
\renewcommand{\warning}[2]{\begin{warningbox}{{#1}}#2\end{warningbox}}
\renewcommand{\error}[2]{\begin{errorbox}{{#1}}#2\end{errorbox}}

% ============================================================
%  INDICE: NUMERI DI PAGINA A TRE CIFRE
% ============================================================
% I numeri di pagina a tre cifre non entrano nella casella standard di
% 1,55em; la casella si allarga a una misura fissa, cosi' tutti i numeri
% finiscono sullo stesso margine destro.
\makeatletter
\renewcommand{\@pnumwidth}{22pt}
\makeatother

% ============================================================
%  NUMERAZIONE DI ESEMPI ED ESERCIZI
% ============================================================
% Nella classe i contatori sono agganciati a "section"; con secnumdepth=0
% il contatore di sezione non viene mai incrementato, per cui si otterrebbe
% "Esempio 1.0.3". Li si riaggancia quindi al capitolo.
\counterwithout{esempio}{section}
\counterwithin{esempio}{chapter}
\counterwithout{esercizio}{section}
\counterwithin{esercizio}{chapter}

% L'esempio viene composto in una scatola: se sta in una pagina la si mette
% intera (e passa alla pagina dopo se non entra), altrimenti la si "svuota"
% nella pagina e l'esempio puo' spezzarsi normalmente.
\newbox\wtesbox
\newcommand{\wtemetti}{%
  \ifdim\dimexpr\ht\wtesbox+\dp\wtesbox\relax>\textheight
    \unvbox\wtesbox
  \else
    \box\wtesbox
  \fi}
\renewenvironment{esempio}[1]{%
  \par\addvspace{0.8em}%
  \setbox\wtesbox\vbox\bgroup
  \refstepcounter{esempio}%
  \noindent\rule{\linewidth}{0.5pt}\par\nopagebreak%
  \vspace{0.15em}%
  \noindent{\color{primary!80!black}\textbf{Esempio \theesempio}}%
  \ifx\relax#1\relax\else%
    {\color{primary!80!black}\textbf{ --- #1}}%
  \fi\par\nopagebreak%
  \vspace{0.3em}%
}{%
  \par\vspace{0.3em}%
  \noindent\rule{\linewidth}{0.5pt}\par
  \egroup
  \wtemetti
  \par\vspace{0.8em}%
}

\renewenvironment{esercizio}[1]{%
  \par\addvspace{0.8em}%
  \setbox\wtesbox\vbox\bgroup
  \refstepcounter{esercizio}%
  \noindent\rule{\linewidth}{1.2pt}\vspace{1pt}\par%
  \noindent\rule{\linewidth}{0.4pt}\par%
  \vspace{0.15em}%
  \noindent{\color{accent!80!black}\textbf{Esercizio \theesercizio}}%
  \ifx\relax#1\relax\else%
    {\color{accent!80!black}\textbf{ --- #1}}%
  \fi\par%
  \vspace{0.3em}%
}{%
  \par\vspace{0.3em}%
  \noindent\rule{\linewidth}{0.4pt}\vspace{1pt}\par%
  \noindent\rule{\linewidth}{1.2pt}\par
  \egroup
  \wtemetti
  \par\vspace{0.8em}%
}

% ============================================================
%  MACRO INLINE
% ============================================================
% Codice inline breve (non spezzabile). \leavevmode serve quando il codice
% apre un paragrafo (per esempio all'inizio della cella di una tabella):
% senza, il colore verrebbe applicato prima del paragrafo e la cella
% comincerebbe con una riga vuota.
\DeclareRobustCommand{\code}[1]{\leavevmode{\FiraCode\color{primary!65!black}\mbox{#1}}}
% Codice inline lungo (spezzabile a fine riga)
\DeclareRobustCommand{\codel}[1]{\leavevmode{\FiraCode\color{primary!65!black}#1}}
% Termine tecnico evidenziato (spezzabile, a differenza di \hi)
\def\term#1{{\color{accent!80!black}\LatoBlack #1}}
% URL mostrata come testo
\def\urlx#1{{\FiraCode\small\color{primary!70!black}\url{#1}}}

% \tasto{Ctrl}: un tasto della tastiera; \tasti{Ctrl,Shift,P}: una
% combinazione di tasti da premere insieme
% (le dimensioni minime si azzerano: dentro un nodo di una figura, il tasto
% erediterebbe quelle del nodo che lo contiene)
\newcommand{\tasto}[1]{\tikz[baseline=(k.base)]\node[draw=neutral!55!black,
  fill=neutral!6!white, rounded corners=2pt, line width=0.5pt,
  minimum width=0pt, minimum height=0pt, text width=, align=center,
  inner xsep=3pt, inner ysep=1.6pt, text=black, font=\Lato\scriptsize] (k) {\strut #1};}
\makeatletter
\newcommand{\tasti}[1]{{\let\wt@sep\relax
  \@for\wt@k:=#1\do{\wt@sep\tasto{\wt@k}\def\wt@sep{\,{\Lato\scriptsize+}\,}}}}
% \menu{File,Open Folder...}: un percorso nei menu di un programma
\newcommand{\menu}[1]{{\let\wt@sep\relax\Lato\bfseries\color{primary!65!black}%
  \@for\wt@k:=#1\do{\wt@sep\mbox{\wt@k}\def\wt@sep{\,{\color{neutral!55!black}›}\,}}}}
\makeatother
% \pulsante{Install}: un pulsante di un'interfaccia grafica
\newcommand{\pulsante}[1]{{\setlength{\fboxsep}{1.6pt}%
  \colorbox{dmaccent}{\Lato\scriptsize\color{white}\strut\,#1\,}}}

% Elenco di passi numerati (installazione, procedure): pallini con il numero
\NewDocumentCommand{\wtpassonum}{m}{\tikz[baseline=(n.base)]\node[circle,
  fill=primary!80!black, text=white, font=\Lato\footnotesize\bfseries,
  inner sep=0pt, minimum size=13pt] (n) {#1};}
\newlist{passi}{enumerate}{1}
\setlist[passi]{label=\wtpassonum{\arabic*}, leftmargin=24pt, labelsep=6pt,
  itemsep=3pt, topsep=4pt}

% ============================================================
%  RIQUADRO "IN SINTESI"
% ============================================================
% Chiude le parti piu' dense riassumendo il ragionamento (non una regola
% da imparare a memoria, ma il filo logico da ricordare).
\newtcolorbox{sintesibox}{enhanced, breakable,
  colback=primary!4!white, colframe=primary!70!black, boxrule=0pt,
  leftrule=3pt, arc=0pt, outer arc=0pt, left=8pt, right=8pt, top=5pt,
  bottom=5pt, before skip=0.4cm, after skip=0.5cm, fontupper=\small,
  title={\Lato\bfseries\small In sintesi}, coltitle=primary!70!black,
  colbacktitle=primary!4!white, titlerule=0pt, toptitle=3pt, bottomtitle=0pt}
\newcommand{\insintesi}[1]{\begin{sintesibox}#1\end{sintesibox}}

% ============================================================
%  FINESTRA "BROWSER" (pagine web da cui si scaricano i programmi)
%  Uso: \begin{browser}[indirizzo] ... \end{browser}
% ============================================================
\newcommand{\wtbrowserbefore}{0.35cm}
\newcommand{\wtbrowserafter}{0.6cm}
\newtcolorbox{browser}[1][]{%
  enhanced,
  colback=white,
  colframe=neutral!35!white,
  boxrule=0.8pt,
  arc=3pt,
  drop fuzzy shadow,
  before skip=\wtbrowserbefore, after skip=\wtbrowserafter,
  left=8pt, right=8pt, top=6pt, bottom=6pt,
  fonttitle=\footnotesize\Lato,
  colbacktitle=neutral!15!white,
  coltitle=neutral!35!black,
  title={%
    \tikz[baseline=-2.5pt]{%
      \fill[red-gradient-3] (0,0) circle (2.2pt);
      \fill[yellow-gradient-5] (8pt,0) circle (2.2pt);
      \fill[green-gradient-6] (16pt,0) circle (2.2pt);}%
    \hspace{8pt}\footnotesize #1},
  attach boxed title to top left={yshift*=-\tcboxedtitleheight},
  boxed title style={
    sharp corners=downhill, arc=3pt,
    colframe=neutral!35!white, boxrule=0.8pt,
  },
  fontupper=\Lato\small,
}

% ============================================================
%  TERMINALE, per mostrare comandi e output dei programmi
%  Uso: \begin{console}[titolo] ... \end{console}
% ============================================================
% Il contenuto e' VERBATIM: l'output di un programma contiene spesso
% caratteri speciali per LaTeX (_ \ ^ $ % &) e gli a-capo vanno preservati.
% Le righe che cominciano con il prompt "$" (o "PS C:" in PowerShell) sono
% COMANDI digitati dall'utente e sono piu' chiare; il resto e' cio' che
% stampa il computer. Cio' che l'utente digita mentre un programma e' in
% esecuzione (l'input) si racchiude fra << e >> e compare in giallo.
% gcc cita i nomi fra virgolette tipografiche (‘main’): listings non sa
% gestire i caratteri non ASCII e li stamperebbe fuori posto, quindi nel
% sorgente vanno scritti come (*@‘@*) e (*@’@*).
% Come per lstlisting, questo ambiente non puo' stare dentro l'argomento
% di una macro: nei riquadri va usata la forma ambiente (warningbox...).
\definecolor{consolebg}{HTML}{22252E}
\definecolor{consolefg}{HTML}{C9CDD6}
\definecolor{consolecmd}{HTML}{FFFFFF}
\definecolor{consoleinput}{HTML}{F2D46B}
\definecolor{consolebar}{HTML}{31353F}
\newcommand{\wtconsolebefore}{0.35cm}
\newtcblisting{console}[1][Terminale]{%
  listing only,
  enhanced,
  colback=consolebg,
  colframe=consolebg,
  boxrule=0pt,
  arc=3pt,
  drop fuzzy shadow,
  before skip=\wtconsolebefore, after skip=0.6cm,
  left=8pt, right=8pt, top=6pt, bottom=6pt,
  fonttitle=\footnotesize\Lato,
  colbacktitle=consolebar,
  coltitle=consolefg!80!consolebg,
  title={\iterminal\hspace{6pt}\footnotesize #1},
  attach boxed title to top left={yshift*=-\tcboxedtitleheight},
  boxed title style={sharp corners=downhill, arc=3pt, boxrule=0pt,
                     colframe=consolebar},
  listing options={
    language={},
    basicstyle=\linespread{1.1}\footnotesize\FiraCodeNL\color{consolefg},
    backgroundcolor=\color{consolebg},
    keywordstyle={}, commentstyle={}, stringstyle={},
    moredelim={[l][\color{consolecmd}\bfseries]{\$\ }},
    moredelim={[l][\color{consolecmd}\bfseries]{PS\ C:}},
    moredelim={[is][\color{consoleinput}\bfseries]{<<}{>>}},
    numbers=none, frame=none,
    xleftmargin=0pt, xrightmargin=0pt,
    breaklines=true, columns=fixed,
    mathescape=false, escapeinside={(*@}{@*)},
    showtabs=false, tabsize=8,
  },
}

% ============================================================
%  STILI TIKZ CONDIVISI DALLE FIGURE
% ============================================================
\tikzset{
  % nei riquadri delle figure le parole non vanno spezzate con il trattino
  nosillabe/.style={execute at begin node={\hyphenpenalty=10000\exhyphenpenalty=10000}},
  wtnode/.style={nosillabe,
    rounded corners=3pt, draw=primary!65!black, fill=primary!8!white,
    line width=0.8pt, inner sep=5pt, align=center, font=\Lato\small},
  wtnodeacc/.style={
    wtnode, draw=accent!60!black, fill=accent!8!white},
  wtnodewarn/.style={
    wtnode, draw=orange-gradient-6, fill=orange-gradient-1!45!white},
  wtnodeok/.style={
    wtnode, draw=green-gradient-6, fill=green-gradient-1!45!white},
  wtnodeplain/.style={nosillabe,
    rounded corners=2pt, draw=neutral!60!white, fill=neutral!8!white,
    line width=0.7pt, inner sep=4pt, align=center, font=\Lato\footnotesize},
  wtarrow/.style={
    -{Stealth[length=6pt,width=5pt]}, line width=0.9pt, draw=primary!70!black},
  wtarrowacc/.style={
    -{Stealth[length=6pt,width=5pt]}, line width=0.9pt, draw=accent!70!black},
  wtdash/.style={
    dashed, line width=0.7pt, draw=neutral!50!white},
  wtlab/.style={
    font=\Lato\scriptsize, text=neutral!15!black, inner sep=2.5pt,
    fill=white, align=center},
  wtbrace/.style={
    decorate, decoration={brace, amplitude=4pt, mirror}, draw=neutral!40!black,
    line width=0.7pt},
}

% Figura centrata con didascalia, ridotta alla larghezza indicata
\newcommand{\wtfig}[3][0.95]{%
  \begin{figure}[H]
    \centering
    \resizebox{#1\linewidth}{!}{#2}
    \caption{#3}
  \end{figure}}

% ============================================================
%  IL CODICE NON SI SPEZZA FRA DUE PAGINE
% ============================================================
% Ogni lstlisting viene racchiuso in una minipage. Una minipage e' una
% scatola indivisibile: se il listato non entra nello spazio rimasto in
% fondo alla pagina, passa per intero alla pagina successiva.
\AddToHook{env/lstlisting/before}{\wtattacca\noindent\begin{minipage}{\linewidth}}
\AddToHook{env/lstlisting/after}{\end{minipage}\par}

% ============================================================
%  CODICE E RISULTATO AFFIANCATI
% ============================================================
% Uso:
%   \begin{affianco}[0.55]      % frazione della riga occupata dal codice
%     \begin{lstlisting}...\end{lstlisting}   (uno o piu' listati)
%   \risultato
%     \begin{console} ... \end{console}   (oppure uno schema)
%   \end{affianco}
% Le due colonne sono allineate in alto e formano un blocco unico, che non
% viene mai diviso fra due pagine. Con il secondo argomento facoltativo [c]
% le colonne sono centrate in verticale l'una rispetto all'altra.
\newlength{\wtcodew}
\newlength{\wtresw}
\newlength{\wtgap}
\setlength{\wtgap}{0.5cm}
\NewDocumentEnvironment{affianco}{O{0.56} O{t}}{%
  \par\addvspace{6pt}\noindent
  \setlength{\wtcodew}{#1\linewidth}%
  \setlength{\wtresw}{\linewidth-\wtcodew-\wtgap}%
  \wtvscmargini{0pt}{0pt}%
  \lstset{aboveskip=0pt, belowskip=5pt}%
  \renewcommand{\wtbrowserbefore}{5pt}%
  \renewcommand{\wtbrowserafter}{6pt}%
  \renewcommand{\wtconsolebefore}{5pt}%
  \if c#2%
    \def\risultato{\end{minipage}\hspace{\wtgap}%
      \begin{minipage}[c]{\wtresw}}%
    \begin{minipage}[c]{\wtcodew}%
  \else
    \def\risultato{\end{minipage}\hspace{\wtgap}%
      \begin{minipage}[t]{\wtresw}\vspace{0pt}}%
    \begin{minipage}[t]{\wtcodew}\vspace{0pt}%
  \fi
}{%
  \end{minipage}\par\addvspace{10pt}}

% Listati un po' piu' piccoli, per le righe lunghe dentro una colonna stretta
\newcommand{\codicepiccolo}{%
  \lstset{basicstyle=\linespread{1.12}\wtvsccodefontsmall,
          numberstyle=\fontsize{5}{6}\selectfont\FiraCodeNL\color{vsclinenum}}}

% Didascalia breve sotto un risultato
\newcommand{\wtnota}[1]{%
  {\par\centering\Lato\footnotesize\color{neutral!30!black}#1\par}\vspace{4pt}}

% ============================================================
%  SCHEMA ACCANTO AL CODICE
% ============================================================
% Nella colonna di destra di "affianco":
%   \begin{immagine}[didascalia facoltativa]
%     \begin{tikzpicture} ... \end{tikzpicture}
%   \end{immagine}
% Lo schema viene centrato nella colonna e, se e' piu' largo, ridotto alla
% sua larghezza (nel log compare "Schema ridotto", per accorgersene).
\newsavebox{\wtimgbox}
\NewDocumentEnvironment{immagine}{o}
  {\begin{lrbox}{\wtimgbox}}
  {\end{lrbox}%
   \par\noindent\makebox[\linewidth]{%
     \ifdim\wd\wtimgbox>\linewidth
       \typeout{Schema ridotto: \the\wd\wtimgbox\space -> \the\linewidth
                (riga \the\inputlineno)}%
       \resizebox{\linewidth}{!}{\usebox{\wtimgbox}}%
     \else
       \usebox{\wtimgbox}%
     \fi}\par
   \IfValueT{#1}{\vspace{4pt}\wtnota{#1}}}

% ============================================================
%  SCHEMI DELLA MEMORIA (variabili, puntatori, array, struct)
% ============================================================
% Livelli per i riquadri annidati (per esempio le chiamate di funzione
% sulla pila): il piu' esterno va disegnato sotto agli altri
\pgfdeclarelayer{liv1}
\pgfdeclarelayer{liv2}
\pgfdeclarelayer{liv3}
\pgfsetlayers{liv1,liv2,liv3,background,main}

\colorlet{mref}{accent!70!black}
\colorlet{mok}{green-gradient-6!85!black}
\colorlet{merr}{red-gradient-6!90!black}
\tikzset{
  % casella di una variabile: \node[mvar] (x) {42};
  mvar/.style={draw=neutral!55!black, fill=white, line width=0.7pt,
    minimum height=0.5cm, minimum width=1.4cm, inner sep=2pt,
    font=\FiraCodeNL\scriptsize},
  % nome della variabile, a sinistra della casella
  mnome/.style={font=\FiraCodeNL\scriptsize, text=primary!60!black,
    inner sep=2pt},
  % indirizzo di memoria, accanto a una casella
  mindir/.style={font=\FiraCodeNL\tiny, text=neutral!45!black, inner sep=2pt},
  % intestazione e righe di una struttura (vedi \oggetto)
  mtesta/.style={draw=primary!65!black, fill=primary!16!white,
    line width=0.7pt, minimum height=0.46cm, inner sep=2pt,
    font=\Lato\scriptsize\bfseries, text=primary!35!black},
  mriga/.style={draw=primary!65!black, fill=primary!3!white,
    line width=0.7pt, minimum height=0.46cm, inner sep=3pt,
    font=\FiraCodeNL\scriptsize, align=left},
  % freccia di un puntatore e pallino da cui parte
  mrif/.style={-{Stealth[length=5pt,width=4pt]}, line width=0.85pt,
    draw=mref},
  mpunto/.style={circle, fill=mref, inner sep=1.4pt},
  % freccia tratteggiata: una copia, un percorso, un "va a leggere"
  mcerca/.style={-{Stealth[length=5pt,width=4pt]}, line width=0.8pt,
    draw=green-gradient-6!80!black, dashed},
  % nota esplicativa nello schema
  mnota/.style={font=\Lato\scriptsize, text=neutral!20!black, align=left,
    inner sep=2pt},
  % codice dentro uno schema
  mcod/.style={font=\FiraCodeNL\scriptsize, inner sep=2pt},
  % riquadro di una funzione
  mfun/.style={draw=accent!60!black, fill=accent!8!white, line width=0.7pt,
    rounded corners=6pt, minimum height=0.5cm, inner sep=3pt,
    font=\FiraCodeNL\scriptsize},
  % riquadri annidati: \scopo{livello}{nome}{etichetta}{(nodo1) (nodo2) ...}
  mscopo1/.style={draw=neutral!50!white, fill=neutral!6!white,
    line width=0.7pt, rounded corners=5pt, inner sep=7pt},
  mscopo2/.style={draw=primary!45!white, fill=primary!6!white,
    line width=0.7pt, rounded corners=5pt, inner sep=7pt},
  mscopo3/.style={draw=accent!45!white, fill=accent!7!white,
    line width=0.7pt, rounded corners=5pt, inner sep=7pt},
  mscopolab/.style={font=\Lato\scriptsize\bfseries, text=neutral!25!black,
    inner sep=1.5pt},
}

% \oggetto[stile della testata]{nome}{(posizione)}{larghezza}{titolo}
%         {campo/valore, campo/valore, ...}
% Disegna una struttura (struct): crea i nodi "nome" (testata), "nome-1",
% "nome-2", ... (righe) e "nome-box" (tutta la struttura). La posizione e'
% quella del centro della testata, salvo un'ancora diversa nel primo
% argomento (per esempio anchor=north west). La larghezza indicata e'
% quella minima: se una riga e' piu' lunga, la struttura si allarga. Un
% valore che contiene una virgola va tra graffe: sep/{","}.
\newlength{\mchiavew}
\setlength{\mchiavew}{1.25cm}
\newcommand{\mchiave}[1]{\setlength{\mchiavew}{#1}}
\newlength{\mobjw}
\newlength{\mkeyw}
\newlength{\mrigaw}
\newsavebox{\wtimgmisura}
\NewDocumentCommand{\oggetto}{O{} m m m m m}{%
  \global\mobjw=#4\relax
  \global\mkeyw=\mchiavew
  \ifx\relax#6\relax\else
    \foreach \mk/\mv in {#6} {%
      \pgfinterruptpicture
        \global\setbox\wtimgmisura\hbox{\FiraCodeNL\scriptsize\mk}%
      \endpgfinterruptpicture
      \ifdim\dimexpr\wd\wtimgmisura+6pt\relax>\mkeyw
        \global\mkeyw=\dimexpr\wd\wtimgmisura+6pt\relax
      \fi}%
    \foreach \mk/\mv in {#6} {%
      \pgfinterruptpicture
        \global\setbox\wtimgmisura\hbox{\FiraCodeNL\scriptsize
          \makebox[\mkeyw][l]{\mk}\mv}%
      \endpgfinterruptpicture
      \setlength{\mrigaw}{\wd\wtimgmisura}%
      \ifdim\dimexpr\mrigaw+8pt\relax>\mobjw
        \global\mobjw=\dimexpr\mrigaw+8pt\relax
      \fi}%
  \fi
  \node[mtesta, minimum width=\mobjw, #1] (#2) at #3 {#5};
  \xdef\wtprev{#2}%
  \ifx\relax#6\relax
    \node[mriga, minimum width=\mobjw, anchor=north] (#2-1)
      at ([yshift=0.7pt]\wtprev.south) {\phantom{x}};
    \xdef\wtprev{#2-1}%
  \else
    \foreach \mk/\mv [count=\mi] in {#6} {%
      \node[mriga, minimum width=\mobjw, text width=\dimexpr\mobjw-6pt\relax,
            anchor=north] (#2-\mi) at ([yshift=0.7pt]\wtprev.south)
        {\makebox[\mkeyw][l]{\color{primary!55!black}\mk}\mv};
      \xdef\wtprev{#2-\mi}%
    }%
  \fi
  \node[fit=(#2)(\wtprev), inner sep=0pt] (#2-box) {};}

% \variabile[stile]{nome del nodo}{(posizione)}{nome mostrato}{valore}
% Il nome compare a sinistra della casella (nodo "nome-nome"). Per un
% puntatore si lascia il valore vuoto e si usa \riferimento.
\NewDocumentCommand{\variabile}{O{} m m m m}{%
  \node[mvar, #1] (#2) at #3 {#5};
  \node[mnome, anchor=east] (#2-nome) at ([xshift=-2pt]#2.west) {#4};}

% \riferimento[stile del percorso]{nodo di partenza}{percorso fino alla meta}
% es.: \riferimento{p}{-- (x.west)}  oppure  \riferimento{p}{to[bend left] (x)}
\NewDocumentCommand{\riferimento}{O{} m m}{%
  \draw[mrif, #1] (#2.center) #3;
  \node[mpunto] at (#2.center) {};}

% \scopo[stile]{livello 1-3}{nome}{etichetta}{(nodo) (nodo) ...}
% Il riquadro racchiude i nodi indicati e lascia in alto lo spazio per
% l'etichetta (per esempio il nome della funzione di cui e' la memoria).
\NewDocumentCommand{\scopo}{O{} m m m m}{%
  \node[fit={#5}, inner sep=0pt] (#3-in) {};
  \node[mscopolab, anchor=south west] (#3-lab)
    at ([yshift=1pt]#3-in.north west) {#4};
  \begin{pgfonlayer}{liv#2}
    \node[mscopo#2, fit=(#3-in)(#3-lab), #1] (#3) {};
  \end{pgfonlayer}}

% \celle[stile]{nome}{(posizione)}{v0, v1, ...}: le caselle consecutive di
% un array, con l'indice sotto ciascuna. La posizione e' il bordo sinistro
% della prima casella; crea i nodi nome-0, nome-1, ... e nome-box. Con
% \mvuota si disegna una casella senza valore, con \noindici prima della
% chiamata si omettono gli indici.
\tikzset{
  mcella/.style={draw=primary!65!black, fill=primary!3!white, line width=0.7pt,
    minimum height=0.5cm, minimum width=0.75cm, inner xsep=4pt,
    font=\FiraCodeNL\scriptsize},
  mindice/.style={font=\FiraCodeNL\tiny, text=neutral!45!black, inner sep=1pt},
}
\newcommand{\mvuota}{{\Lato\itshape\color{neutral!50!black}?}}
\newif\ifmindici \mindicitrue
\newcommand{\noindici}{\mindicifalse}
\NewDocumentCommand{\celle}{O{} m m m}{%
  \coordinate (#2-inizio) at #3;
  \xdef\wtprev{#2-inizio}%
  \foreach \v [count=\i from 0] in {#4} {%
    \node[mcella, anchor=west, #1] (#2-\i)
      at ([xshift=-0.7pt]\wtprev.east |- #2-inizio) {\v};
    \ifmindici
      \node[mindice, anchor=north] at ([yshift=-1pt]#2-\i.south) {\i};
    \fi
    \xdef\wtprev{#2-\i}%
  }
  \node[fit=(#2-0)(\wtprev), inner sep=0pt] (#2-box) {};}

% \memoria[stile]{nome}{(posizione)}{indirizzo/valore/etichetta, ...}: una
% colonna di celle di memoria, dall'alto verso il basso, con l'indirizzo a
% sinistra e (facoltativa) un'etichetta a destra. La posizione e' il centro
% della prima cella; crea i nodi nome-1, nome-2, ... e nome-box.
\tikzset{
  mcellam/.style={draw=primary!65!black, fill=primary!3!white, line width=0.7pt,
    minimum height=0.48cm, minimum width=1.7cm, inner sep=2pt,
    font=\FiraCodeNL\scriptsize},
}
\NewDocumentCommand{\memoria}{O{} m m m}{%
  \coordinate (#2-inizio) at #3;
  \foreach \ind/\val/\etic [count=\i] in {#4} {%
    \node[mcellam, #1] (#2-\i) at ($(#2-inizio)+(0,-\i*0.48cm+0.48cm-\i*0.7pt+0.7pt)$) {\val};
    \node[mindir, anchor=east] at ([xshift=-2pt]#2-\i.west) {\ind};
    \node[mnome, anchor=west] at ([xshift=3pt]#2-\i.east) {\etic};
    \xdef\wtultima{#2-\i}%
  }
  \node[fit=(#2-1)(\wtultima), inner sep=0pt] (#2-box) {};}

% ─── Liste collegate ─────────────────────────────────────────
% \nodo[stile]{nome}{(posizione)}{dato}: un nodo di una lista, con il campo
% dei dati e il campo next affiancati. La posizione e' il bordo sinistro;
% crea i nodi nome-dato, nome-next e nome (tutto il nodo). \nodonull{nome}
% barra il campo next: il puntatore vale NULL. \collega{da}{a} disegna la
% freccia dal campo next del primo nodo al secondo.
\tikzset{
  mlnodo/.style={draw=primary!65!black, fill=primary!3!white, line width=0.7pt,
    minimum height=0.55cm, inner sep=2pt, font=\FiraCodeNL\scriptsize},
  mlnuovo/.style={draw=accent!65!black, fill=accent!10!white},
  mlvia/.style={draw=red-gradient-6!70!black, fill=red-gradient-1!25!white},
}
\NewDocumentCommand{\nodo}{O{} m m m}{%
  \node[mlnodo, minimum width=0.95cm, anchor=west, #1] (#2-dato) at #3 {#4};
  \node[mlnodo, minimum width=0.55cm, anchor=west, #1] (#2-next)
    at ([xshift=-0.7pt]#2-dato.east) {};
  \node[fit=(#2-dato)(#2-next), inner sep=0pt] (#2) {};}
\newcommand{\nodonull}[1]{%
  \draw[line width=0.7pt, draw=primary!65!black]
    ([shift={(1.5pt,1.5pt)}]#1-next.south west) -- ([shift={(-1.5pt,-1.5pt)}]#1-next.north east);}
\NewDocumentCommand{\collega}{O{} m m}{%
  \draw[mrif, #1] (#2-next.center) -- (#3.west);
  \node[mpunto] at (#2-next.center) {};}

% Segni di riuscita e di errore negli schemi
\newcommand{\mok}{{\color{mok}\itick}}
\newcommand{\merr}{{\color{merr}\icross}}

% ============================================================
%  DIAGRAMMI DI FLUSSO
% ============================================================
% Forme standard: inizio e fine (ovale), istruzione (rettangolo), input e
% output (parallelogramma), decisione (rombo), sottoprogramma (rettangolo
% con le bande laterali).
\tikzset{
  fcterm/.style={nosillabe, draw=primary!65!black, fill=primary!18!white, line width=0.8pt,
    rounded corners=0.29cm, minimum height=0.58cm, minimum width=1.7cm,
    inner xsep=8pt, font=\Lato\scriptsize\bfseries, text=primary!35!black},
  fcproc/.style={nosillabe, draw=primary!65!black, fill=primary!4!white, line width=0.8pt,
    minimum height=0.58cm, minimum width=2.1cm, inner sep=4pt, align=center,
    font=\Lato\scriptsize},
  fcio/.style={nosillabe, trapezium, trapezium left angle=72, trapezium right angle=108,
    trapezium stretches body, draw=green-gradient-6!80!black,
    fill=green-gradient-1!40!white, line width=0.8pt, minimum height=0.58cm,
    minimum width=2.1cm, inner sep=3pt, align=center, font=\Lato\scriptsize},
  fcdec/.style={nosillabe, diamond, aspect=2.1, draw=accent!60!black, fill=accent!7!white,
    line width=0.8pt, inner sep=1.5pt, align=center, font=\Lato\scriptsize},
  fcsub/.style={fcproc, fill=orange-gradient-1!30!white,
    draw=orange-gradient-6!80!black, inner xsep=9pt,
    path picture={\draw[orange-gradient-6!80!black, line width=0.6pt]
      ([xshift=4pt]path picture bounding box.north west) --
      ([xshift=4pt]path picture bounding box.south west)
      ([xshift=-4pt]path picture bounding box.north east) --
      ([xshift=-4pt]path picture bounding box.south east);}},
  fcfreccia/.style={-{Stealth[length=5pt,width=4.2pt]}, line width=0.8pt,
    draw=neutral!30!black},
  fclinea/.style={line width=0.8pt, draw=neutral!30!black},
  fcsi/.style={font=\Lato\tiny\bfseries, text=green-gradient-6!55!black,
    inner sep=1.5pt},
  fcno/.style={font=\Lato\tiny\bfseries, text=red-gradient-6!75!black,
    inner sep=1.5pt},
  fccod/.style={font=\FiraCodeNL\scriptsize},
  % punto in cui due strade si ricongiungono (o in cui rientra un ciclo)
  fcgiunzione/.style={circle, draw=neutral!30!black, fill=white, line width=0.8pt,
    inner sep=0pt, minimum size=5.5pt},
}

% ============================================================
%  FILE SYSTEM: CARTELLE E FILE
% ============================================================
\tikzset{
  fsdir/.style={inner sep=1.5pt, font=\Lato\scriptsize, text=primary!45!black,
    execute at begin node={{\color{orange-gradient-5!85!black}\ifolder}\,}},
  fsfile/.style={inner sep=1.5pt, font=\Lato\scriptsize, text=neutral!20!black,
    execute at begin node={{\color{neutral!55!black}\ifile}\,}},
  fsqui/.style={fsdir, fill=yellow-gradient-1!70!white, rounded corners=2pt},
  fsramo/.style={draw=neutral!50!black, line width=0.55pt},
}

% ============================================================
%  INTERFACCIA DI VS CODE (tema scuro predefinito "Dark Modern")
% ============================================================
% Colori delle zone della finestra, per le figure che la riproducono
\definecolor{dmbg}{HTML}{1F1F1F}       % editor
\definecolor{dmbar}{HTML}{181818}      % barre: attivita', laterale, stato
\definecolor{dmborder}{HTML}{2B2B2B}   % bordi fra le zone
\definecolor{dmfg}{HTML}{CCCCCC}       % testo
\definecolor{dmdim}{HTML}{9D9D9D}      % testo secondario, icone inattive
\definecolor{dmaccent}{HTML}{0078D4}   % pulsanti, indicatore remoto
\definecolor{dminput}{HTML}{313131}    % caselle di testo
\definecolor{dmsel}{HTML}{2A2D2E}      % voce evidenziata in un elenco
\definecolor{dmyellow}{HTML}{CCA700}   % avvisi

% \vscico[colore]{nome}: le icone della barra delle attivita' di VS Code,
% disegnate in un quadrato 10x10. Nomi: esplora, cerca, git, debug,
% estensioni, remoto, account, impostazioni, gioca.
\makeatletter
\def\vsc@ico@esplora{%
  \draw[line join=round] (2.6,1) rectangle (8.6,7.4);
  \draw[line join=round] (1.2,2.6) -- (1.2,9) -- (6.2,9) -- (7.4,7.8);}
\def\vsc@ico@cerca{%
  \draw (5.9,6.1) circle (2.7);
  \draw[line width=1.1pt] (3.9,4.1) -- (1.2,1.4);}
\def\vsc@ico@git{%
  \draw (3,8.4) circle (1.2); \draw (3,1.6) circle (1.2);
  \draw (7.4,6.6) circle (1.2);
  \draw (3,7.2) -- (3,2.8);
  \draw (7.4,5.4) .. controls (7.4,3.6) and (3,4.6) .. (3,2.8);}
\def\vsc@ico@debug{%
  \draw[line join=round] (1.4,1.2) -- (1.4,8.8) -- (7.4,5) -- cycle;
  \fill (7.6,2.2) circle (1.5);}
\def\vsc@ico@estensioni{%
  \draw (1,1) rectangle (4.6,4.6); \draw (4.6,1) rectangle (8.2,4.6);
  \draw (1,4.6) rectangle (4.6,8.2);
  \draw[rotate around={45:(7.4,7.4)}] (5.6,5.6) rectangle (9.2,9.2);}
\def\vsc@ico@remoto{%
  \draw[line cap=round, line join=round] (1,6.4) -- (3.6,3.8) -- (1,1.2);
  \draw[line cap=round, line join=round] (9,8.8) -- (6.4,6.2) -- (9,3.6);}
\def\vsc@ico@account{%
  \draw (5,6.6) circle (2.2);
  \draw (1.2,0.8) .. controls (1.6,4.2) and (8.4,4.2) .. (8.8,0.8);}
\def\vsc@ico@impostazioni{%
  \draw (5,5) circle (1.6);
  \draw (5,5) circle (3.4);
  \foreach \ang in {0,45,...,315} \draw[line width=1.2pt] (\ang:3.4)++(5,5) -- ++(\ang:1.1);}
\def\vsc@ico@gioca{%
  \fill[line join=round] (2,1) -- (2,9) -- (8.6,5) -- cycle;}
\NewDocumentCommand{\vscico}{O{dmdim} m}{%
  \tikz[baseline=2.4pt, x=1.15pt, y=1.15pt, line width=0.75pt, draw=#1, fill=#1]{%
    \useasboundingbox (0,0) rectangle (10,10);
    \csname vsc@ico@#2\endcsname}}
\makeatother

% Colori del prompt di Ubuntu nel terminale di VS Code (utente@computer in
% verde, cartella in blu)
\definecolor{termgreen}{HTML}{23D18B}
\definecolor{termblue}{HTML}{3B8EEA}

% \vscodelogo[dimensione]: il logo di Visual Studio Code, semplificato
\NewDocumentCommand{\vscodelogo}{O{0.5cm}}{%
  \tikz[baseline=0pt, x=#1, y=#1]{%
    \fill[dmaccent!85!white] (0.70,0.98) -- (0.95,0.86) -- (0.95,0.14) -- (0.70,0.02)
      -- (0.31,0.37) -- (0.13,0.23) -- (0.04,0.28) -- (0.04,0.72) -- (0.13,0.77)
      -- (0.31,0.63) -- cycle;
    \fill[white] (0.70,0.73) -- (0.70,0.27) -- (0.43,0.5) -- cycle;
    \fill[white] (0.13,0.62) -- (0.23,0.5) -- (0.13,0.38) -- cycle;}}

% Icone della barra di stato: errori (cerchio con la croce) e avvisi
% (triangolo con il punto esclamativo)
\NewDocumentCommand{\icoerr}{O{dmfg}}{\tikz[baseline=-0.75ex]{%
  \draw[#1, line width=0.5pt] (0,0) circle (2.7pt);
  \draw[#1, line width=0.5pt] (-1.2pt,-1.2pt) -- (1.2pt,1.2pt) (-1.2pt,1.2pt) -- (1.2pt,-1.2pt);}}
\NewDocumentCommand{\icowarn}{O{dmfg}}{\tikz[baseline=-0.75ex]{%
  \draw[#1, line width=0.5pt, line join=round] (-3pt,-2.5pt) -- (3pt,-2.5pt) -- (0,3pt) -- cycle;
  \draw[#1, line width=0.5pt] (0,-0.2pt) -- (0,1.4pt); \fill[#1] (0,-1.3pt) circle (0.35pt);}}

% Sottolineatura ondulata degli errori nell'editor
\usetikzlibrary{decorations.pathmorphing}
\tikzset{ondina/.style={decorate, decoration={snake, amplitude=0.55pt,
  segment length=2.6pt}, line width=0.5pt}}

% \wtcasellina{0|1}: casella di spunta di un programma di installazione
\newcommand{\wtcasellina}[1]{\tikz[baseline=-2.8pt]{%
  \ifnum#1=1
    \fill[rounded corners=1pt, blue!70!black] (-3.2pt,-3.2pt) rectangle (3.2pt,3.2pt);
    \draw[white, line width=1.1pt] (-1.9pt,-0.2pt) -- (-0.5pt,-1.7pt) -- (2pt,1.8pt);
  \else
    \draw[rounded corners=1pt, neutral!60!black, line width=0.6pt]
      (-3.2pt,-3.2pt) rectangle (3.2pt,3.2pt);
  \fi}}
% Puntatore del mouse, per indicare dove fare clic
\newcommand{\wtpuntatore}{\tikz[baseline=0pt, scale=0.55]{%
  \filldraw[fill=white, draw=black, line width=0.35pt, line join=round]
    (0,0) -- (0,-0.62) -- (0.15,-0.48) -- (0.26,-0.72) -- (0.35,-0.68)
    -- (0.24,-0.44) -- (0.44,-0.44) -- cycle;}}

% ============================================================
%  ICONE DELLA BARRA DI DEBUG DI VS CODE
% ============================================================
% \wtdbgicon{nome} su sfondo chiaro, \wtdbgiconscura{nome} su sfondo scuro.
% Nomi: maniglia, resume, stepover, stepinto, stepout, restart, disconnect.
\definecolor{dbgblu}{HTML}{007ACC}
\definecolor{dbgverde}{HTML}{388A34}
\definecolor{dbgrosso}{HTML}{A1260D}
\definecolor{dbggrigio}{HTML}{8A8A8A}
\definecolor{dbgbluscuro}{HTML}{75BEFF}
\definecolor{dbgverdescuro}{HTML}{89D185}
\definecolor{dbgrossoscuro}{HTML}{F48771}
\definecolor{dbggrigioscuro}{HTML}{8A8A8A}
\makeatletter
\tikzset{dbgfreccia/.style={-{Stealth[length=3pt,width=3.6pt]}}}
\def\wtdbg@colore@maniglia{grigio}
\def\wtdbg@colore@resume{blu}
\def\wtdbg@colore@stepover{blu}
\def\wtdbg@colore@stepinto{blu}
\def\wtdbg@colore@stepout{blu}
\def\wtdbg@colore@restart{verde}
\def\wtdbg@colore@disconnect{rosso}
\def\wtdbg@maniglia{\foreach \x in {3.5,6.5} \foreach \y in {2,5,8} \fill (\x,\y) circle (0.9);}
\def\wtdbg@resume{\draw (1.5,1) -- (1.5,9);
  \draw[line join=round] (4,1.3) -- (4,8.7) -- (9.6,5) -- cycle;}
\def\wtdbg@stepover{\draw[dbgfreccia] (1.2,4.6) arc[start angle=180, end angle=8, radius=3.8];
  \fill (5,1.5) circle (1.5);}
\def\wtdbg@stepinto{\draw[dbgfreccia] (5,10) -- (5,4.2); \fill (5,1.5) circle (1.5);}
\def\wtdbg@stepout{\draw[dbgfreccia] (5,4.4) -- (5,10.2); \fill (5,1.5) circle (1.5);}
\def\wtdbg@restart{\draw[-{Stealth[length=3pt,width=3.6pt]}]
  (2.2,7.6) arc[start angle=135, end angle=-180, radius=3.9];}
\def\wtdbg@disconnect{%
  \begin{scope}[rotate around={45:(5,5)}]
    \draw[rounded corners=0.8] (3.4,3.6) rectangle (6.4,6.4);
    \draw (6.4,4.4) -- (8.4,4.4) (6.4,5.6) -- (8.4,5.6) (0.8,5) -- (3.4,5);
  \end{scope}
  \draw (1.6,8.4) -- (8.4,1.6);}
\newcommand{\wtdbg@icona}[2]{%
  \edef\wtdbg@c{dbg\csname wtdbg@colore@#1\endcsname#2}%
  \tikz[baseline=0pt, x=0.85pt, y=0.85pt, line width=0.85pt, line cap=round,
        draw=\wtdbg@c, fill=\wtdbg@c]{%
    \useasboundingbox (0,0) rectangle (10,10);
    \csname wtdbg@#1\endcsname}}
\newcommand{\wtdbgicon}[1]{\wtdbg@icona{#1}{}}
\newcommand{\wtdbgiconscura}[1]{\wtdbg@icona{#1}{scuro}}
\makeatother

% ============================================================
%  LE SPIEGAZIONI NON SI SPEZZANO FRA DUE PAGINE
% ============================================================
% I riquadri della classe (info, warning, error, definizione) sono
% dichiarati "breakable": qui la chiave "breakable" di tcolorbox viene
% ridefinita in modo che equivalga a "unbreakable": un riquadro che non
% entra passa intero alla pagina successiva, come i listati.
\tcbset{breakable/.style={breakable@false}}

% ============================================================
%  Un concetto non si spezza fra due pagine
% ============================================================
% Il testo che introduce un listato, una figura, una lista o un risultato
% ("... come segue:") resta attaccato a cio' che introduce: se non entrano
% insieme in fondo alla pagina, passano insieme alla successiva.
\raggedbottom
% Un paragrafo non si divide mai fra due pagine
\interlinepenalty=10000
% Un codice inline lungo in fondo alla riga non esce dal margine
\emergencystretch=3em
\makeatletter
% Niente salto pagina fra la frase introduttiva e la lista o il blocco
% centrato (tabelle, figure non flottanti) che la segue
\@beginparpenalty=10000
% Una lista resta tutta sulla stessa pagina
\@itempenalty=10000
\makeatother
\newcommand{\wtattacca}{\par\nopagebreak}
% (per i listati \wtattacca e' gia' nell'hook che li chiude in una minipage)
\AddToHook{env/figure/before}{\wtattacca}
\AddToHook{env/table/before}{\wtattacca}
\AddToHook{env/console/before}{\wtattacca}
\AddToHook{env/browser/before}{\wtattacca}
\AddToHook{env/affianco/before}{\wtattacca}
% Un titolo non resta da solo in fondo alla pagina, e ogni "blocco" (da un
% titolo di sezione o sottosezione al titolo successivo) che sta in una
% pagina non viene spezzato: la posizione di inizio e di fine di ogni blocco
% viene salvata nel file .aux (zref-savepos) e, alla compilazione
% successiva, prima del titolo si chiede tanto spazio quanto e' alto il
% blocco. Se non c'e', il blocco intero passa alla pagina dopo. Servono
% quindi diverse compilazioni perche' le posizioni si stabilizzino.
\usepackage{zref-savepos,zref-abspage}
\newcount\wtblocco
\makeatletter
\zref@addprop{savepos}{abspage}
\newcommand{\wtfineblocco}{%
  \ifnum\wtblocco>\z@ \zsavepos{wtb\the\wtblocco e}\fi}
\newcommand{\wtiniziablocco}{%
  \global\advance\wtblocco\@ne \zsavepos{wtb\the\wtblocco s}}
% \wtspazio{minimo}: Needspace pari all'altezza del prossimo blocco, se sta
% in una pagina, altrimenti pari al minimo. L'altezza misurata quando il
% blocco stava tutto in una pagina viene ricordata (\wtmem nel file .aux):
% cosi' un blocco gia' spostato non torna indietro alla compilazione dopo.
\newdimen\wtaltezza
\newcommand{\wtmem}[2]{\expandafter\gdef\csname wtmem@#1\endcsname{#2}}
% \wtcap{n}: il blocco n apre un capitolo, quindi il blocco n-1 chiude il
% precedente e la sua "fine" viene registrata nella pagina del nuovo capitolo
\newcommand{\wtcap}[1]{\expandafter\gdef\csname wtcap@#1\endcsname{}}
% \wtsegue{n}: il blocco n contiene solo il titolo di una sezione, seguito
% subito da una sottosezione: il titolo viaggia insieme al blocco successivo.
\newcommand{\wtsegue}[1]{\expandafter\gdef\csname wtsegue@#1\endcsname{}}
\newcommand{\wtspazio}[1]{%
  \begingroup
  \edef\wtn{\the\numexpr\wtblocco+1\relax}%
  \edef\wtps{\zref@extractdefault{wtb\wtn s}{abspage}{0}}%
  \edef\wtpe{\zref@extractdefault{wtb\wtn e}{abspage}{0}}%
  \wtaltezza=#1\relax
  \ifcsname wtsegue@\wtn\endcsname
    \edef\wtm{\the\numexpr\wtn+1\relax}%
    \edef\wtqs{\zref@extractdefault{wtb\wtm s}{abspage}{0}}%
    \edef\wtqe{\zref@extractdefault{wtb\wtm e}{abspage}{0}}%
    \ifnum\wtqs>\z@
      \ifnum\wtqe=\wtqs
        \wtaltezza=\dimexpr\zposy{wtb\wtm s}sp-\zposy{wtb\wtm e}sp
                   +4\baselineskip\relax
      \else\ifcsname wtmem@\wtm\endcsname
        \ifdim\csname wtmem@\wtm\endcsname<\maxdimen
          \wtaltezza=\dimexpr\csname wtmem@\wtm\endcsname+4\baselineskip\relax
        \fi
      \fi\fi
    \fi
  \else
  \ifnum\wtps>\z@
    \ifnum\wtpe=\wtps
      \wtaltezza=\dimexpr\zposy{wtb\wtn s}sp-\zposy{wtb\wtn e}sp\relax
      \ifcsname wtmem@\wtn\endcsname
        \ifdim\csname wtmem@\wtn\endcsname>\wtaltezza
          \wtaltezza=\csname wtmem@\wtn\endcsname\relax
        \fi
      \fi
      \immediate\write\@auxout{\string\wtmem{\wtn}{\the\wtaltezza}}%
    \else\ifnum\wtpe=\numexpr\wtps+1\relax
      \ifcsname wtcap@\the\numexpr\wtn+1\relax\endcsname
      \else\ifdim\zposy{wtb\wtn e}sp>\wttesta
        \ifdim\zposy{wtb\wtn s}sp>\wttesta
          \wtaltezza=.96\textheight
          \immediate\write\@auxout{\string\wtmem{\wtn}{\the\wtaltezza}}%
        \fi
      \else\ifcsname wtmem@\wtn\endcsname
        \wtaltezza=\csname wtmem@\wtn\endcsname\relax
        \immediate\write\@auxout{\string\wtmem{\wtn}{\the\wtaltezza}}%
      \else\ifdim\zposy{wtb\wtn s}sp>\wttesta
        \immediate\write\@auxout{\string\wtmem{\wtn}{\the\maxdimen}}%
      \else
        \wtaltezza=\maxdimen
      \fi\fi\fi\fi
    \fi\fi
  \fi
  \fi
  \ifdim\wtaltezza=\maxdimen
    \par\newpage
  \else
    \ifdim\wtaltezza<#1\relax \wtaltezza=#1\relax\fi
    \ifdim\wtaltezza>.97\textheight \wtaltezza=#1\relax\fi
    \Needspace{\dimexpr\wtaltezza+\baselineskip\relax}%
  \fi
  \endgroup}
% posizione verticale oltre la quale un punto e' "in cima alla pagina"
\newdimen\wttesta
\AtBeginDocument{\wttesta=\dimexpr\paperheight-\topmargin-1in-\headheight-\headsep-40pt\relax}
\AddToHook{cmd/chapter/before}{\wtfineblocco\wtiniziablocco\immediate\write\@auxout{\string\wtcap{\the\wtblocco}}}
\AddToHook{cmd/section/before}{\wtfineblocco\wtspazio{22\baselineskip}\wtiniziablocco}
\AddToHook{cmd/subsection/before}{\wtfineblocco
  \if@nobreak\immediate\write\@auxout{\string\wtsegue{\the\wtblocco}}\fi
  \wtspazio{8\baselineskip}\wtiniziablocco}
\AddToHook{enddocument}{\wtfineblocco}
\makeatother


\setcounter{secnumdepth}{0}

\widowpenalty10000
\clubpenalty10000

\graphicspath{ {./images/} }

\begin{document}

  \coverpage{26}{\documentyear}{44}{\documentname}

  % I capitoli partono da 1 (la classe imposta il contatore a -1)
  \setcounter{chapter}{0}

  % ─── Estratto ──────────────────────────────────────────────
  \begin{abstract}
    Questa dispensa raccoglie, riorganizza e spiega il materiale del corso
    di \documentname\ tenuto dal professor \textit{Francesco Isgrò} per il
    corso di laurea triennale in Informatica della Federico II di Napoli
    (slide dell'\documentyear). Il linguaggio del corso è il \textbf{C}.

    È pensata per chi \textbf{parte da zero}: non si dà per scontato di
    aver mai scritto una riga di codice, né di aver mai usato un terminale.
    Il primo capitolo accompagna passo passo nell'installazione di tutto
    ciò che serve (l'editor \textbf{Visual Studio Code}, il compilatore e le
    estensioni) su Windows, macOS e Linux; i successivi seguono l'ordine
    delle lezioni.

    L'obiettivo è \textbf{capire}, non memorizzare: ogni regola è
    accompagnata dal suo \textit{perché}, ogni programma è stato compilato
    ed eseguito e ne viene mostrato l'output reale, e ciò che accade
    \textit{dentro} il computer (le variabili in memoria, i puntatori, le
    chiamate di funzione, i nodi di una lista) è disegnato in uno schema.
    Dove le slide contengono una svista, il testo la corregge e lo segnala.

    Sono inclusi \textbf{soltanto i contenuti tecnici} presentati a
    lezione; le informazioni organizzative (orari, modalità d'esame,
    appelli) restano sulle slide e sui canali del docente, dove vengono
    aggiornate. Questo documento \textbf{non sostituisce} le lezioni, le
    slide e i libri di testo: ne è un complemento.

    \vspace{0.4em}
    Nel testo ricorrono alcuni elementi grafici: riquadri \textbf{verdi}
    con approfondimenti, \textbf{gialli} con avvertenze ed errori tipici,
    \textbf{rossi} con errori gravi da evitare, riquadri \textbf{Definizione}
    e riquadri \textbf{In sintesi} che riassumono il ragionamento. Il codice
    è disegnato come una finestra di VS Code e, \textbf{accanto}, compare
    ciò che produce: il \textbf{terminale}, dove i comandi digitati sono in
    bianco e ciò che si scrive mentre il programma è in esecuzione è in
    \textbf{\color{orange-gradient-6!80!black}giallo}, oppure uno schema.

    \ifx \xauthors \undefined
    \else
      \begin{flushright}
        \textit{Redazione a cura di} \xauthors
      \end{flushright}
    \fi
  \end{abstract}

  \tableofcontents

  \chapter{L'ambiente di lavoro}

Prima di scrivere il primo programma bisogna preparare il proprio
\textit{laboratorio}: qualche programma da installare e un paio di abitudini
da prendere. Questo capitolo accompagna passo passo dall'installazione
alla prima esecuzione, su \textbf{Windows}, \textbf{macOS} e \textbf{Linux}.
Ogni passo è spiegato: non serve sapere nulla di computer oltre ad
accenderli e a scaricare un file da Internet.

\warning{Leggi solo la parte del tuo sistema operativo}{
  Le istruzioni di installazione sono divise per sistema operativo. Segui
  quelle del tuo computer e salta le altre: dalla sezione \textit{Il
  terminale} in poi le istruzioni tornano a essere le stesse per tutti.}

\section{Che cosa ci serve}

Un computer non capisce il C. Capisce soltanto il proprio
\term{linguaggio macchina}: lunghe sequenze di numeri che il processore sa
eseguire, una dopo l'altra. Noi invece scriveremo i programmi in C, un
linguaggio fatto di parole inglesi e simboli, comprensibile a una persona.
Fra noi e il processore servono quindi tre strumenti, che compaiono nella
figura~\ref{fig:catena}.

\begin{figure}[H]
  \centering
  \begin{tikzpicture}[
      card/.style={rounded corners=4pt, draw=#1!65!black, fill=#1!6!white,
        line width=0.8pt, minimum width=4.6cm, minimum height=3.3cm},
      titolo/.style={font=\Lato\small\bfseries, text=#1!55!black},
      schermo/.style={fill=vscbg, rounded corners=2pt, inner sep=5pt,
        font=\FiraCodeNL\scriptsize, text=vscfg, align=left},
      nota/.style={font=\Lato\scriptsize, text=neutral!25!black, align=center}]
    \node[card=primary] (a) at (0,0) {};
    \node[card=accent] (b) at (6.6,0) {};
    \node[card=primary] (c) at (13.2,0) {};
    \node[titolo=primary, anchor=north] at ([yshift=-5pt]a.north) {1. Scrivere};
    \node[titolo=accent, anchor=north] at ([yshift=-5pt]b.north) {2. Tradurre};
    \node[titolo=primary, anchor=north] at ([yshift=-5pt]c.north) {3. Eseguire};
    \node[schermo] at ([yshift=-1pt]a.center) {%
      {\color{vscblue}int} {\color{vscfunc}main}({\color{vscblue}void})\\
      \{\\
      \ \ {\color{vscfunc}printf}({\color{vscstring}"Ciao!\textbackslash n"});\\
      \}};
    \node[schermo, text=vscnum] at ([yshift=-1pt]b.center) {%
      01010101 01001000\\10001001 11100101\\01001000 10001101\\11000011 \dots};
    \node[schermo] at ([yshift=-1pt]c.center) {%
      {\color{white}\bfseries\$ ./ciao}\\Ciao!\\[2pt]\phantom{x}};
    \node[nota, anchor=south] at ([yshift=5pt]a.south) {\FiraCodeNL ciao.c\Lato: un file di testo};
    \node[nota, anchor=south] at ([yshift=5pt]b.south) {\FiraCodeNL ciao\Lato: linguaggio macchina};
    \node[nota, anchor=south] at ([yshift=5pt]c.south) {ciò che il programma stampa};
    \draw[wtarrow] (a.east) -- node[wtlab, above=2pt] {compilare}
      node[below=2pt, font=\FiraCodeNL\scriptsize] {gcc} (b.west);
    \draw[wtarrow] (b.east) -- node[wtlab, above=2pt] {eseguire}
      node[below=2pt, font=\FiraCodeNL\scriptsize] {./ciao} (c.west);
    \node[nota, anchor=north] at ([yshift=-4pt]a.south)
      {\textbf{l'editor}\\Visual Studio Code};
    \node[nota, anchor=north] at ([yshift=-4pt]b.south)
      {\textbf{il compilatore}\\gcc};
    \node[nota, anchor=north] at ([yshift=-4pt]c.south)
      {\textbf{il terminale}\\dove si danno i comandi};
  \end{tikzpicture}
  \caption{Dal testo scritto da noi al programma che gira: i tre strumenti
           da installare.}
  \label{fig:catena}
\end{figure}

\begin{itemize}
  \item L'\term{editor} è il quaderno su cui scriviamo il programma. Il
        programma, per ora, è un normale \textbf{file di testo} con
        estensione \code{.c} (per esempio \code{ciao.c}). Useremo
        \textbf{Visual Studio Code} (o semplicemente \textit{VS Code}), un
        editor gratuito che colora il codice, segnala gli errori mentre si
        scrive e contiene al suo interno anche un terminale.
  \item Il \term{compilatore} è il traduttore: legge il file di testo e
        produce un nuovo file, l'\term{eseguibile}, che contiene lo stesso
        programma tradotto in linguaggio macchina. Useremo \textbf{gcc}, il
        compilatore C del progetto GNU, lo stesso usato a lezione.
  \item Il \term{terminale} è una finestra in cui si \textit{scrivono} i
        comandi al computer invece di cliccare su icone. Da lì
        chiederemo a gcc di compilare e poi lanceremo il programma, che
        stamperà i suoi risultati nella stessa finestra.
\end{itemize}

Il capitolo~\ref{chap:dal-problema} spiegherà in dettaglio che cosa succede
durante la traduzione. Per ora basta l'idea: \textbf{si scrive, si compila,
si esegue}.

\section{Linux, Windows o macOS?}

Il corso \textbf{assume Linux} come sistema operativo: i comandi visti a
lezione (e all'esame) sono quelli di Linux. Non serve però cambiare
computer o cancellare Windows: in ogni sistema c'è un modo semplice per
avere gli stessi strumenti.

\begin{center}
\renewcommand{\arraystretch}{1.45}
\begin{tabular}{@{}>{\bfseries}l>{\raggedright\arraybackslash}p{7.6cm}>{\raggedright\arraybackslash}p{6.4cm}@{}}
  \toprule
  Sistema & Che cosa si installa & Perché \\
  \midrule
  Windows & \textbf{WSL} (\textit{Windows Subsystem for Linux}) con
            \textbf{Ubuntu}, e dentro Ubuntu gcc &
            WSL fa girare un vero Linux dentro Windows, senza riavvii né
            partizioni \\
  macOS   & gli \textbf{strumenti da riga di comando} di Apple
            (\textit{Command Line Tools}) &
            macOS è un parente di Linux: ha già un terminale molto
            simile, manca solo il compilatore \\
  Linux   & il pacchetto \code{build-essential} (Ubuntu, Debian, Mint) &
            è il sistema del corso: basta un comando \\
  \bottomrule
\end{tabular}
\end{center}

\textbf{VS Code} si installa allo stesso modo ovunque (passo 1). Su Windows
c'è una particolarità che conviene capire subito, perché spiega diversi
passaggi successivi: i file del corso e il compilatore stanno
\textit{dentro} Ubuntu, mentre la finestra di VS Code è un normale
programma di Windows. VS Code allora si divide in due parti che collaborano
(figura~\ref{fig:wsl}): l'interfaccia, su Windows, e un piccolo
\textit{server}, dentro Ubuntu, che legge e salva i file e lancia il
compilatore.

\begin{figure}[H]
  \centering
  \begin{tikzpicture}[font=\Lato\scriptsize,
      parte/.style={wtnode, font=\Lato\scriptsize, minimum height=1.3cm,
        text width=3.3cm}]
    % Windows
    \draw[neutral!50!white, fill=neutral!5!white, line width=0.8pt,
          rounded corners=6pt] (-8.6,2.7) rectangle (8.6,-2.6);
    \node[anchor=north west, font=\Lato\small\bfseries, text=neutral!25!black]
      at (-8.4,2.55) {%
      \tikz[baseline=0pt]{\foreach \x/\y in {0/0,0.17/0,0/0.17,0.17/0.17}
        \fill[dmaccent] (\x,\y) rectangle ++(0.15,0.15);}\ Windows};
    % VS Code
    \node[wtnode, text width=3.6cm, minimum height=3.2cm] (vsc) at (-5.9,-0.1)
      {\vscodelogo[0.8cm]\\[4pt]\textbf{Visual Studio Code}\\[2pt]
       la finestra: menu, editor, terminale};
    % Ubuntu
    \draw[orange-gradient-6!80!black, fill=orange-gradient-1!22!white,
          line width=0.8pt, rounded corners=6pt] (-2.9,1.95) rectangle (8.2,-2.25);
    \node[anchor=north west, font=\Lato\small\bfseries,
          text=orange-gradient-6!45!black] at (-2.75,1.85) {Ubuntu (dentro WSL)};
    \node[parte] (srv) at (-0.6,-0.25) {\textbf{VS Code Server}\\con l'estensione C/C++};
    \node[parte, draw=accent!60!black, fill=accent!8!white] (gcc) at (5.6,0.65) {\textbf{gcc}, \textbf{make}, \textbf{gdb}\\gli strumenti del corso};
    \node[parte, draw=green-gradient-6, fill=green-gradient-1!45!white] (file) at (5.6,-1.25) {\textbf{i tuoi file}\\{\FiraCodeNL\tiny /home/fabio/programmazione}};
    \draw[wtarrow, <->] (vsc.east) -- node[wtlab, above=2pt] {si parlano} (srv.west);
    \draw[wtarrow] (srv.east) -- node[wtlab, above, sloped] {lancia} (gcc.west);
    \draw[wtarrow] (srv.east) -- node[wtlab, below, sloped] {legge e salva} (file.west);
  \end{tikzpicture}
  \caption{Su Windows, con WSL: l'interfaccia di VS Code resta su Windows,
           mentre file e compilatore vivono in Ubuntu. Su macOS e su Linux
           tutto sta nello stesso sistema e questa divisione non c'è.}
  \label{fig:wsl}
\end{figure}

\info{Perché non un compilatore ``nativo'' per Windows?}{
  Esistono versioni di gcc per Windows (MinGW, per esempio) e funzionano.
  Ma con esse alcuni dettagli cambiano: l'eseguibile si chiama \code{a.exe}
  invece di \code{a.out}, la fine dell'input si segnala con
  \tasti{Ctrl,Z} invece di \tasti{Ctrl,D}, i comandi del terminale sono
  diversi. Con WSL si usano \textbf{esattamente} i comandi delle lezioni, ed
  è ciò che conta.}

\section{Passo 1: installare Visual Studio Code}

\subsection{Scaricare il programma}

Visual Studio Code si scarica gratuitamente dal sito ufficiale,
\urlx{https://code.visualstudio.com}. La pagina \textit{Download} propone
un pulsante per ogni sistema operativo (figura~\ref{fig:download}): si
sceglie il proprio.

\warning{Visual Studio Code, non Visual Studio}{
  Microsoft produce anche \textbf{Visual Studio} (senza ``Code''): è un
  altro programma, molto più grande, pensato per altri linguaggi. Il nome
  giusto è \textbf{Visual Studio Code}, con l'icona blu della
  figura~\ref{fig:download}.}

\begin{figure}[H]
  \centering
  \begin{minipage}{15.5cm}
  \begin{browser}[code.visualstudio.com/download]
    \centering
    {\Lato\Large\bfseries Download Visual Studio Code}\\[3pt]
    {\Lato\footnotesize\color{neutral!35!black} Free and built on open source.
     Integrated Git, debugging and extensions.}\\[8pt]
    \begin{tikzpicture}[font=\Lato\scriptsize,
        bottone/.style={fill=dmaccent, text=white, rounded corners=2pt,
          minimum width=3.3cm, minimum height=0.95cm, align=center,
          font=\Lato\scriptsize},
        dett/.style={font=\Lato\tiny, text=neutral!35!black, align=center}]
      % Windows
      \foreach \x/\y in {0/0,0.27/0,0/0.27,0.27/0.27}
        \fill[dmaccent] (-0.27+\x,0.55+\y) rectangle ++(0.25,0.25);
      \node[bottone] (w) at (0,-0.15) {{\bfseries\small Windows}\\Windows 10, 11};
      \node[dett, anchor=north] at (w.south) {User Installer \quad x64 \quad Arm64};
      % Linux
      \begin{scope}[shift={(4.6,0.85)}, scale=0.75]
        \fill[black] (0,0) ellipse (0.28 and 0.38);
        \fill[white] (0,-0.08) ellipse (0.18 and 0.26);
        \fill[white] (-0.08,0.2) circle (0.05); \fill[white] (0.08,0.2) circle (0.05);
        \fill[orange-gradient-5] (-0.07,0.12) -- (0.07,0.12) -- (0,0.03) -- cycle;
        \fill[orange-gradient-5] (-0.17,-0.37) ellipse (0.13 and 0.05);
        \fill[orange-gradient-5] (0.17,-0.37) ellipse (0.13 and 0.05);
      \end{scope}
      \node[bottone, minimum width=2.1cm] (deb) at (3.45,-0.15) {{\bfseries\small .deb}\\Debian, Ubuntu};
      \node[bottone, minimum width=2.1cm] (rpm) at (5.75,-0.15) {{\bfseries\small .rpm}\\Red Hat, Fedora};
      \node[dett, anchor=north] at ($(deb.south)!0.5!(rpm.south)$) {x64 \quad Arm64};
      % Mac
      \begin{scope}[shift={(9.6,0.68)}]
        \fill[neutral!20!black] (0,0) circle (0.17);
        \fill[neutral!20!black] (0,-0.05) ellipse (0.16 and 0.14);
        \fill[white] (0.2,0.03) circle (0.07);
        \fill[neutral!20!black, rotate around={-35:(0.03,0.24)}] (0.03,0.24) ellipse (0.035 and 0.08);
      \end{scope}
      \node[bottone] (m) at (9.6,-0.15) {{\bfseries\small Mac}\\macOS};
      \node[dett, anchor=north] at (m.south) {.zip \quad Intel \quad Apple silicon \quad Universal};
    \end{tikzpicture}
  \end{browser}
  \end{minipage}
  \caption{La pagina di download di VS Code. L'aspetto del sito cambia
           spesso; i tre gruppi di pulsanti, invece, restano.}
  \label{fig:download}
\end{figure}

\subsection{Installazione su Windows}

Il pulsante \textit{Windows} scarica un file chiamato più o meno
\code{VSCodeUserSetup-x64-1.xx.x.exe} (i numeri sono la versione). È un
normale programma di installazione:

\begin{passi}
  \item Apri il file scaricato (doppio clic, oppure clic sulla notifica di
        download del browser).
  \item Nella prima schermata scegli \textit{I accept the agreement} (accetti
        la licenza) e fai clic su \textit{Next}.
  \item Nelle due schermate successive (cartella di installazione e menu
        Start) lascia tutto com'è e fai clic su \textit{Next}.
  \item Nella schermata \textit{Select Additional Tasks}
        (figura~\ref{fig:installer}) spunta le due voci \textit{Add ``Open
        with Code'' action\dots}: aggiungono la voce ``Apri con Code'' al
        menu del tasto destro sui file e sulle cartelle, comodissima. La voce
        \textbf{Add to PATH} deve restare spuntata: permette di aprire VS
        Code scrivendo \code{code} nel terminale, e la useremo.
  \item Fai clic su \textit{Next}, poi su \textit{Install}. Al termine,
        lascia spuntato \textit{Launch Visual Studio Code} e fai clic su
        \textit{Finish}.
\end{passi}

\begin{figure}[H]
  \centering
  \begin{tikzpicture}[font=\Lato\scriptsize]
    \def\W{12.4} \def\H{7.6}
    \draw[neutral!45!white, fill=white, line width=0.6pt, rounded corners=4pt]
      (0,0) rectangle (\W,-\H);
    % barra del titolo
    \node[anchor=west] at (0.15,-0.3) {\vscodelogo[0.3cm]};
    \node[anchor=west] at (0.55,-0.3) {Setup - Microsoft Visual Studio Code (User)};
    \node[anchor=east, font=\Lato\footnotesize] at (\W-0.2,-0.3)
      {--\hspace{16pt}$\square$\hspace{16pt}\texttimes};
    % intestazione
    \node[anchor=north west, font=\Lato\footnotesize\bfseries] at (0.45,-0.72) {Select Additional Tasks};
    \node[anchor=north west] at (0.75,-1.1) {Which additional tasks should be performed?};
    \node at (\W-0.85,-1.12) {\vscodelogo[0.75cm]};
    \draw[neutral!35!white] (0,-1.68) -- (\W,-1.68);
    % corpo
    \fill[neutral!3!white] (0.02,-1.69) rectangle (\W-0.02,-6.58);
    \node[anchor=north west, text width=11cm] at (0.45,-1.9)
      {Select the additional tasks you would like Setup to perform while
       installing Visual Studio Code, then click Next.};
    \node[anchor=west] at (0.45,-2.75) {Additional icons:};
    \node[anchor=west] at (0.65,-3.12) {\wtcasellina{0}\hspace{5pt}Create a desktop icon};
    \node[anchor=west] at (0.45,-3.6) {Other:};
    \node[anchor=west] (r1) at (0.65,-3.98) {\wtcasellina{1}\hspace{5pt}Add
      \textquotedbl Open with Code\textquotedbl\ action to Windows Explorer file context menu};
    \node[anchor=west] (r2) at (0.65,-4.38) {\wtcasellina{1}\hspace{5pt}Add
      \textquotedbl Open with Code\textquotedbl\ action to Windows Explorer directory context menu};
    \node[anchor=west] at (0.65,-4.86) {\wtcasellina{1}\hspace{5pt}Register Code as an editor for supported file types};
    \node[anchor=west] (path) at (0.65,-5.34) {\wtcasellina{1}\hspace{5pt}Add to PATH (requires shell restart)};
    \draw[accent!70!black, line width=0.9pt, rounded corners=2pt]
      ([xshift=-3pt,yshift=0.5pt]r1.north west) rectangle ([xshift=3pt,yshift=-0.5pt]r2.south east);
    \draw[accent!70!black, line width=0.9pt, rounded corners=2pt]
      ([xshift=-3pt,yshift=0.5pt]path.north west) rectangle ([xshift=3pt,yshift=-0.5pt]path.south east);
    \node[anchor=west, font=\Lato\scriptsize\bfseries, text=accent!70!black, align=left]
      at ([xshift=12pt]path.east) {deve restare spuntata!};
    \node[anchor=west, font=\Lato\scriptsize\bfseries, text=accent!70!black, align=left]
      at ([xshift=8pt]r2.east |- r1.south) {comode};
    % pulsanti
    \draw[neutral!35!white] (0,-6.6) -- (\W,-6.6);
    \node[draw=neutral!50!white, fill=neutral!8!white, minimum width=1.8cm,
          minimum height=0.5cm] at (\W-5.0,-7.1) {< Back};
    \node[draw=dmaccent, line width=1pt, fill=neutral!8!white, minimum width=1.8cm,
          minimum height=0.5cm] at (\W-2.95,-7.1) {Next >};
    \node[draw=neutral!50!white, fill=neutral!8!white, minimum width=1.8cm,
          minimum height=0.5cm] at (\W-1.05,-7.1) {Cancel};
  \end{tikzpicture}
  \caption{La schermata più importante dell'installazione su Windows.}
  \label{fig:installer}
\end{figure}

\subsection{Installazione su macOS}

Il pulsante \textit{Mac} scarica un archivio \code{.zip}. Safari di solito lo
estrae da solo: nella cartella \textit{Download} compare l'applicazione
\textit{Visual Studio Code} (altrimenti basta un doppio clic sul file
\code{.zip}).

\begin{passi}
  \item Trascina \textit{Visual Studio Code} dalla cartella \textit{Download}
        alla cartella \textit{Applicazioni} (figura~\ref{fig:mac-app}): così
        è installata davvero e la si trova nel Launchpad.
  \item Aprila. Al primo avvio macOS avvisa che è un'applicazione scaricata
        da Internet: conferma con \textit{Apri}.
  \item Per poter aprire VS Code dal terminale con il comando \code{code}:
        premi \tasti{Cmd,Shift,P}, scrivi \textit{shell command} e scegli
        \textit{Shell Command: Install `code' command in PATH}.
\end{passi}

\begin{figure}[H]
  \centering
  \begin{tikzpicture}[font=\Lato\scriptsize]
    \foreach \x/\nome in {0/Download, 7.2/Applicazioni} {
      \draw[neutral!45!white, fill=white, rounded corners=5pt, line width=0.6pt]
        (\x,0) rectangle ++(4.6,-2.6);
      \fill[neutral!10!white, rounded corners=5pt] (\x,0) rectangle ++(4.6,-0.45);
      \fill[neutral!10!white] (\x,-0.25) rectangle ++(4.6,-0.2);
      \foreach \c/\k in {red-gradient-3/0, yellow-gradient-5/1, green-gradient-6/2}
        \fill[\c] (\x+0.25+\k*0.25,-0.22) circle (0.07);
      \node[font=\Lato\scriptsize\bfseries, text=neutral!30!black] at (\x+2.3,-0.22) {\nome};
    }
    \node at (2.3,-1.25) {\vscodelogo[0.9cm]};
    \node[align=center] at (2.3,-2.2) {Visual Studio Code};
    \node[text=neutral!45!black, align=center] at (9.5,-1.45)
      {(le altre applicazioni)};
    \draw[wtarrowacc, line width=1.4pt] (3.4,-1.2) to[bend left=20]
      node[wtlab, above=2pt, fill=none] {trascina} (8.0,-1.2);
  \end{tikzpicture}
  \caption{Su macOS si installa un'applicazione trascinandola nella
           cartella \textit{Applicazioni}.}
  \label{fig:mac-app}
\end{figure}

\subsection{Installazione su Linux}

Su Ubuntu, Debian e derivate si scarica il pacchetto \code{.deb} e lo si
installa dal terminale (il terminale si apre con \tasti{Ctrl,Alt,T}; il
suo uso è spiegato più avanti in questo capitolo):

\begin{console}
$ cd ~/Scaricati
$ sudo apt install ./code_*.deb
\end{console}

Il comando \code{cd} entra nella cartella dei download (su un sistema in
inglese si chiama \code{Downloads}); \code{sudo apt install} installa il
pacchetto, dopo aver chiesto la password dell'utente. In alternativa, VS
Code si trova anche nel negozio di applicazioni della distribuzione (su
Ubuntu, \textit{App Center}).

\section{Passo 2: installare il compilatore}

\subsection{Windows: WSL e Ubuntu}

Su Windows~10 (aggiornato) e Windows~11 WSL si installa con un unico
comando, da scrivere in \textbf{PowerShell} aperta \textbf{come
amministratore}, cioè con il permesso di modificare il sistema.

\begin{passi}
  \item Premi il tasto \tasto{Windows}, scrivi \textit{powershell} e, nel
        riquadro a destra, scegli \textit{Esegui come amministratore}
        (figura~\ref{fig:start}). Windows chiede conferma: rispondi
        \textit{Sì}.
  \item Nella finestra blu (o nera) che si apre, scrivi il comando
        \code{wsl -{}-install} e premi \tasto{Invio}.
  \item Attendi il download (qualche minuto: Ubuntu pesa circa un
        gigabyte) e, quando richiesto, \textbf{riavvia il computer}.
\end{passi}

\begin{figure}[H]
  \centering
  \begin{tikzpicture}[font=\Lato\scriptsize]
    \draw[neutral!40!white, fill=neutral!4!white, rounded corners=7pt,
          line width=0.6pt] (0,0) rectangle (12,-4.5);
    \draw[neutral!40!white, fill=white, rounded corners=9pt] (0.4,-0.3) rectangle (11.6,-0.88);
    \node[anchor=west] at (0.7,-0.59) {\vscico[neutral!40!black]{cerca}\hspace{6pt}powershell};
    \node[anchor=west, font=\Lato\scriptsize\bfseries] at (0.45,-1.3) {Corrispondenza migliore};
    \fill[primary!10!white, rounded corners=3pt] (0.4,-1.55) rectangle (5.6,-2.45);
    \node[fill=primary!80!black, text=white, font=\FiraCodeNL\tiny\bfseries,
          minimum size=0.55cm, rounded corners=1.5pt] at (0.95,-2.0) {>\_};
    \node[anchor=west] at (1.4,-1.82) {\textbf{Windows PowerShell}};
    \node[anchor=west, text=neutral!40!black] at (1.4,-2.15) {App};
    \fill[white, rounded corners=4pt] (6.0,-1.15) rectangle (11.6,-4.2);
    \node[fill=primary!80!black, text=white, font=\FiraCodeNL\scriptsize\bfseries,
          minimum size=0.8cm, rounded corners=2pt] at (8.8,-1.75) {>\_};
    \node[font=\Lato\footnotesize\bfseries] at (8.8,-2.45) {Windows PowerShell};
    \draw[neutral!25!white] (6.3,-2.75) -- (11.3,-2.75);
    \node[anchor=west] at (6.4,-3.0) {Apri};
    \fill[primary!10!white, rounded corners=2pt] (6.25,-3.18) rectangle (11.35,-3.6);
    \node[anchor=west, font=\Lato\scriptsize\bfseries] (adm) at (6.4,-3.39) {Esegui come amministratore};
    \node[anchor=west] at (6.4,-3.85) {Apri percorso file};
    \node at ([xshift=12pt, yshift=-2pt]adm.east) {\wtpuntatore};
  \end{tikzpicture}
  \caption{La ricerca nel menu Start: PowerShell va aperta come
           amministratore.}
  \label{fig:start}
\end{figure}

\begin{console}[Amministratore: Windows PowerShell]
PS C:\WINDOWS\system32> wsl --install
Installazione: Piattaforma macchina virtuale
Piattaforma macchina virtuale è stato installato.
Installazione: Sottosistema Windows per Linux
Sottosistema Windows per Linux è stato installato.
Installazione: Ubuntu
Ubuntu è stato installato.
L'operazione richiesta è riuscita. Le modifiche saranno effettive al riavvio del sistema.
\end{console}

\subsection{Windows: creare l'utente di Ubuntu}

I messaggi possono essere un po' diversi a seconda della versione di
Windows: quello che conta è che non compaiano errori. Dopo il riavvio si
apre da sola una finestra di Ubuntu (se non si apre, cerca \textit{Ubuntu}
nel menu Start) che completa l'installazione e chiede di creare un
\textbf{utente Linux}: un nome (in minuscolo, senza spazi; non deve essere
per forza uguale a quello di Windows) e una password.

\begin{console}[Ubuntu]
Installing, this may take a few minutes...
Please create a default UNIX user account. The username does not need to match your Windows username.
For more information visit: https://aka.ms/wslusers
Enter new UNIX username: <<fabio>>
New password:
Retype new password:
passwd: password updated successfully
Installation successful!
fabio@PC-FABIO:~$
\end{console}

\warning{Mentre scrivi la password non compare nulla}{
  Nel terminale, mentre si digita una password, \textbf{non compaiono
  nemmeno gli asterischi}: sembra che la tastiera non funzioni, invece i
  caratteri vengono letti. Scrivi la password e premi \tasto{Invio}.
  Annotala: servirà ogni volta che installerai qualcosa.}

\subsection{Windows: installare il compilatore in Ubuntu}

L'ultima riga stampata da Ubuntu, \code{fabio@PC-FABIO:\textasciitilde\$}, è il
\term{prompt}: Ubuntu è pronta a ricevere comandi (la spiegheremo pezzo per
pezzo più avanti). Restano da installare il compilatore e gli altri
strumenti del corso. Scrivi questi due comandi, uno alla volta, premendo
\tasto{Invio} dopo ciascuno:

\begin{console}[Ubuntu]
$ sudo apt update
[sudo] password for fabio:
...
$ sudo apt install build-essential gdb
...
Do you want to continue? [Y/n] <<y>>
...
\end{console}

\begin{itemize}
  \item \code{sudo} (\textit{superuser do}) esegue il comando che segue
        con i permessi di amministratore, e per questo chiede la password
        appena creata.
  \item \code{apt} è il \term{gestore di pacchetti} di Ubuntu: una specie
        di negozio di applicazioni che si usa dal terminale. Con
        \code{update} scarica l'elenco aggiornato dei programmi, con
        \code{install} installa quelli indicati (e chiede conferma: si
        risponde \code{y}, \textit{yes}).
  \item \code{build-essential} è un pacchetto che ne raccoglie altri:
        il compilatore \textbf{gcc}, il programma \textbf{make} (che useremo
        dal capitolo~\ref{chap:make}) e le \textit{librerie} standard del C.
        \textbf{gdb} è il \textit{debugger}, utile per eseguire un programma
        un'istruzione alla volta.
\end{itemize}

Per verificare che tutto sia andato a buon fine, si chiede a gcc la sua
versione:

\begin{console}[Ubuntu]
$ gcc --version
gcc (Ubuntu 13.3.0-6ubuntu2~24.04.1) 13.3.0
Copyright (C) 2023 Free Software Foundation, Inc.
This is free software; see the source for copying conditions.  There is NO
warranty; not even for MERCHANTABILITY or FITNESS FOR A PARTICULAR PURPOSE.
\end{console}

I numeri di versione possono essere diversi: va bene qualunque versione.
Se invece compare \code{gcc: command not found}, l'installazione non è
riuscita: ripeti i due comandi \code{apt} e controlla che non stampino
errori.

\warning{Se \code{wsl -{}-install} non funziona}{
  L'errore più comune all'avvio di Ubuntu è un messaggio che parla di
  \textit{virtualizzazione} (per esempio il codice \code{0x80370102}): WSL
  usa una funzione del processore che su alcuni computer è disattivata.
  Si attiva dal BIOS/UEFI del computer (la voce si chiama
  \textit{Virtualization Technology}, \textit{Intel VT-x} o \textit{AMD-V});
  la procedura cambia da modello a modello, quindi conviene cercare
  ``attivare virtualizzazione'' seguito dal modello del proprio PC. Le
  istruzioni ufficiali di Microsoft sono indicate nei riferimenti a fine
  capitolo.}

\subsection{macOS: gli strumenti da riga di comando}

Su macOS il compilatore fa parte degli \textit{strumenti da riga di
comando} di Apple, gratuiti. Non serve installare Xcode (l'ambiente di
sviluppo completo, molto più grande): bastano gli strumenti.

\begin{passi}
  \item Apri il \textbf{Terminale}: premi \tasti{Cmd,Spazio} (si apre
        Spotlight), scrivi \textit{terminale} e premi \tasto{Invio}.
  \item Scrivi il comando \code{xcode-select -{}-install} e premi
        \tasto{Invio}.
  \item Compare una finestra come quella della figura~\ref{fig:xcode}: fai
        clic su \textit{Installa} (non su \textit{Ottieni Xcode}) e accetta
        la licenza. Il download richiede qualche minuto.
\end{passi}

\begin{figure}[H]
  \centering
  \begin{tikzpicture}[font=\Lato\scriptsize]
    \draw[neutral!35!white, fill=neutral!3!white, rounded corners=9pt,
          line width=0.6pt] (0,0) rectangle (10.2,-3.3);
    \node[fill=neutral!25!white, rounded corners=4pt, minimum size=1cm,
          font=\FiraCodeNL\small\bfseries, text=neutral!25!black] at (0.95,-1.05) {>\_};
    \node[anchor=north west, text width=8.3cm, align=left] at (1.65,-0.35)
      {\textbf{Il comando ``xcode-select'' richiede gli strumenti per
       sviluppatori della riga di comando. Vuoi installarli adesso?}\\[4pt]
       Scegli Installa per scaricare e installare ora gli strumenti per
       sviluppatori della riga di comando.};
    \node[draw=neutral!45!white, fill=white, rounded corners=3pt,
          minimum width=2.1cm, minimum height=0.48cm] at (3.0,-2.75) {Ottieni Xcode};
    \node[draw=neutral!45!white, fill=white, rounded corners=3pt,
          minimum width=1.8cm, minimum height=0.48cm] at (7.0,-2.75) {Non ora};
    \node[fill=dmaccent, text=white, rounded corners=3pt,
          minimum width=1.8cm, minimum height=0.48cm] (ok) at (9.0,-2.75) {Installa};
    \node at ([xshift=6pt, yshift=-5pt]ok.center) {\wtpuntatore};
  \end{tikzpicture}
  \caption{La richiesta di installazione degli strumenti da riga di
           comando (il testo esatto dipende dalla versione di macOS).}
  \label{fig:xcode}
\end{figure}

Al termine, la verifica:

\begin{console}[Terminale di macOS]
$ gcc --version
Apple clang version 16.0.0 (clang-1600.0.26.6)
Target: arm64-apple-darwin24.1.0
Thread model: posix
InstalledDir: /Library/Developer/CommandLineTools/usr/bin
\end{console}

Sorpresa: il comando \code{gcc} risponde \textit{Apple clang}. Su macOS il
compilatore è \textbf{clang}, un altro compilatore C, che per comodità
risponde anche al nome \code{gcc} e accetta le stesse opzioni. Per tutto
ciò che si fa nel corso il comportamento è lo stesso.

\info{Il prompt di macOS}{
  Nel Terminale di macOS il prompt finisce con il simbolo \code{\%} invece
  che con \code{\$} (per esempio \code{fabio@MacBook-Air \textasciitilde\
  \%}). In questa dispensa i comandi sono sempre preceduti da \code{\$}: su
  Mac si scrivono allo stesso modo.}

\subsection{Linux}

Su Ubuntu, Debian e Linux Mint bastano gli stessi due comandi visti per
Windows (in Ubuntu, il terminale si apre con \tasti{Ctrl,Alt,T}):

\begin{console}
$ sudo apt update
$ sudo apt install build-essential gdb
\end{console}

Su Fedora il comando equivalente è
\code{sudo dnf install gcc make gdb}. La verifica, anche qui, è
\code{gcc -{}-version}.

\section{Passo 3: le estensioni di VS Code}

VS Code, appena installato, sa già colorare il codice C. Le
\term{estensioni} sono componenti aggiuntivi gratuiti che gli insegnano a
fare di più. Si installano dalla vista \textit{Extensions}: si apre con
l'icona dei quattro quadratini nella barra a sinistra (oppure con
\tasti{Ctrl,Shift,X}, su Mac \tasti{Cmd,Shift,X}); si scrive il nome
nella casella di ricerca e si fa clic su \pulsante{Install}
(figura~\ref{fig:estensioni}).

\begin{figure}[H]
  \centering
  \begin{tikzpicture}[font=\Lato\scriptsize,
      esten/.style={regular polygon, regular polygon sides=6, minimum size=0.95cm,
        inner sep=0pt, font=\Lato\tiny\bfseries, text=white},
      nome/.style={anchor=west, text=dmfg, font=\Lato\scriptsize\bfseries},
      desc/.style={anchor=west, text=dmdim, font=\Lato\tiny},
      edit/.style={anchor=west, text=dmfg, font=\Lato\tiny}]
    \def\W{10.2}
    \fill[dmbar] (0,0) rectangle (\W,-5.4);
    \draw[dmborder] (0.75,0) -- (0.75,-5.4);
    % barra delle attivita'
    \foreach \ico/\y in {esplora/-0.55, cerca/-1.2, git/-1.85, debug/-2.5} {
      \node at (0.375,\y) {\vscico{\ico}};
    }
    \node at (0.375,-3.15) {\vscico[white]{estensioni}};
    \fill[white] (0,-2.92) rectangle (0.04,-3.38);
    % intestazione e ricerca
    \node[anchor=west, text=dmfg, font=\Lato\tiny] at (0.95,-0.3) {EXTENSIONS: MARKETPLACE};
    \draw[dmaccent, fill=dminput, line width=0.6pt] (0.95,-0.58) rectangle (\W-0.25,-1.02);
    \node[anchor=west, text=dmfg] at (1.05,-0.8) {c/c++};
    % risultati
    \fill[dmsel] (0.76,-1.2) rectangle (\W,-2.5);
    \foreach \y/\c in {-1.2/purple-gradient-6!80!black, -2.5/purple-gradient-5!70!black,
                       -3.8/purple-gradient-4!60!black}
      \node[esten, fill=\c] at (1.5,\y-0.62) {C/C++};
    \node[nome] at (2.2,-1.45) {C/C++};
    \node[desc] at (2.2,-1.78) {C/C++ IntelliSense, debugging, and code browsing.};
    \node[edit] (p1) at (2.2,-2.13) {\tikz[baseline=-0.6ex]\node[circle, fill=dmaccent, inner sep=0.6pt,
      text=white, font=\Lato\tiny]{\itick};\ Microsoft};
    \node[nome] at (2.2,-2.75) {C/C++ Extension Pack};
    \node[desc] at (2.2,-3.08) {Popular extensions for C++ development in Visual Studio Code.};
    \node[edit] at (2.2,-3.43) {\tikz[baseline=-0.6ex]\node[circle, fill=dmaccent, inner sep=0.6pt,
      text=white, font=\Lato\tiny]{\itick};\ Microsoft};
    \node[nome] at (2.2,-4.05) {C/C++ Themes};
    \node[desc] at (2.2,-4.38) {UI Themes for C/C++ extension.};
    \node[edit] at (2.2,-4.73) {\tikz[baseline=-0.6ex]\node[circle, fill=dmaccent, inner sep=0.6pt,
      text=white, font=\Lato\tiny]{\itick};\ Microsoft};
    \foreach \y in {-2.13,-3.43,-4.73}
      \node[fill=dmaccent, text=white, font=\Lato\tiny, inner xsep=4pt, inner ysep=2pt,
            rounded corners=1pt] at (\W-0.75,\y) {Install};
    \node at (\W-0.45,-2.25) {\wtpuntatore};
    % didascalia a destra
    \node[anchor=west, font=\Lato\scriptsize, text=accent!70!black, align=left]
      at (\W+0.3,-1.85) {\textbf{questa}: l'estensione\\ufficiale di Microsoft};
    \draw[wtarrowacc] (\W+0.3,-1.85) -- (\W+0.05,-1.85);
  \end{tikzpicture}
  \caption{La vista \textit{Extensions} con la ricerca ``c/c++''. Basta la
           prima: \textit{C/C++} di Microsoft (si riconosce dal segno blu di
           ``editore verificato'').}
  \label{fig:estensioni}
\end{figure}

Le estensioni da installare sono poche:

\begin{center}
\renewcommand{\arraystretch}{1.45}
\begin{tabular}{@{}>{\raggedright\arraybackslash}p{3.2cm}>{\raggedright\arraybackslash}p{2.4cm}>{\raggedright\arraybackslash}p{10.6cm}@{}}
  \toprule
  \textbf{Estensione} & \textbf{Per chi} & \textbf{A che cosa serve} \\
  \midrule
  \textbf{C/C++} (Microsoft) & tutti &
    \textit{IntelliSense}: suggerimenti mentre si scrive (nomi di funzione,
    parametri), errori sottolineati in rosso ancora prima di compilare,
    navigazione nel codice (clic destro, \textit{Go to Definition}),
    debugger. \\
  \textbf{WSL} (Microsoft) & solo Windows &
    collega VS Code a Ubuntu, come nella figura~\ref{fig:wsl}. \\
  \bottomrule
\end{tabular}
\end{center}

Spesso non serve nemmeno cercarle: aprendo per la prima volta un file
\code{.c}, VS Code propone da solo, con una notifica in basso a destra, di
installare l'estensione consigliata per il linguaggio C.

\info{E la lingua italiana?}{
  VS Code può essere tradotto in italiano con l'estensione
  \textit{Italian Language Pack}. In questa dispensa però i menu sono
  citati \textbf{in inglese}, come compaiono dopo l'installazione: se
  scegli l'italiano dovrai tradurre i nomi delle voci (\textit{File} resta
  \textit{File}, \textit{Terminal} diventa \textit{Terminale}, e così
  via).}

\subsection{Windows: collegare VS Code a Ubuntu}

Su Windows l'ordine conta: prima si installa l'estensione \textbf{WSL}, poi
si apre VS Code \textit{dentro} Ubuntu, e solo allora si installa
\textbf{C/C++}. Il modo più semplice per aprire VS Code dentro Ubuntu è
farlo dal terminale di Ubuntu, con il comando \code{code} seguito da un
punto (il punto significa ``la cartella in cui mi trovo''):

\begin{console}[Ubuntu]
$ code .
Installing VS Code Server for Linux x64 (...)
Downloading: 100%
Unpacking: 100%
\end{console}

La prima volta VS Code installa il suo \textit{server} dentro Ubuntu (lo
fa da solo, in pochi secondi), poi si apre la finestra. Come si capisce che
è collegata a Ubuntu? Dall'\term{indicatore remoto}, il riquadro blu
nell'angolo in basso a sinistra (figura~\ref{fig:remoto}).

\begin{figure}[H]
  \centering
  \begin{tikzpicture}[font=\Lato\scriptsize]
    \fill[dmbar] (0,0) rectangle (13,-0.44);
    \draw[dmborder] (0,0) -- (13,0);
    \fill[dmaccent] (0,0) rectangle (2.5,-0.44);
    \node[anchor=west, text=white] (rem) at (0.15,-0.22)
      {\vscico[white]{remoto}\hspace{5pt}WSL: Ubuntu};
    \node[anchor=west, text=dmfg] at (2.75,-0.22) {\icoerr\ 0\hspace{6pt}\icowarn\ 0};
    \node[anchor=east, text=dmfg] at (12.8,-0.22)
      {Ln 5, Col 31\hspace{14pt}Spaces: 4\hspace{14pt}UTF-8\hspace{14pt}LF\hspace{14pt}\{\} C\hspace{14pt}Linux};
    \draw[wtarrowacc] (1.25,-1.15) -- (1.25,-0.5);
    \node[anchor=north, font=\Lato\scriptsize, text=accent!70!black, align=center]
      at (1.25,-1.15) {\textbf{collegato a Ubuntu}};
  \end{tikzpicture}
  \caption{La barra di stato di VS Code collegato a Ubuntu. Se il riquadro
           in basso a sinistra non dice \textit{WSL: Ubuntu}, VS Code sta
           lavorando su Windows e non troverà il compilatore.}
  \label{fig:remoto}
\end{figure}

Ora, nella vista \textit{Extensions}, l'estensione C/C++ mostra il
pulsante \pulsante{Install in WSL: Ubuntu}: è quello giusto, perché
l'estensione deve lavorare dove stanno i file e il compilatore. Le volte
successive, VS Code ricorda le cartelle aperte: dal menu \menu{File,Open
Recent} si riapre direttamente una cartella di Ubuntu, già collegata.

\info{Dove sono i file di Ubuntu?}{
  La cartella personale di Ubuntu (\code{/home/fabio}) è visibile anche da
  Esplora file di Windows, alla voce \textit{Linux} nel riquadro di
  sinistra. Viceversa, da Ubuntu il disco \code{C:} di Windows si trova in
  \code{/mnt/c}. Conviene però tenere i file del corso \textbf{dentro
  Ubuntu}, nella cartella personale: lì compilare è più veloce e non ci
  sono problemi di permessi.}

\section{Passo 4: la cartella del corso}

Conviene tenere tutti i programmi del corso in una cartella dedicata. La
creiamo dal terminale (su Windows, quello di \textbf{Ubuntu}) e la apriamo
in VS Code:

\begin{console}
$ mkdir programmazione
$ cd programmazione
$ code .
\end{console}

\code{mkdir} (\textit{make directory}) crea la cartella, \code{cd}
(\textit{change directory}) ci entra, \code{code .} apre VS Code sulla
cartella corrente. In alternativa si può fare tutto con il mouse: si crea
la cartella come al solito e la si apre in VS Code da \menu{File,Open
Folder...}.

La prima volta che si apre una cartella, VS Code chiede se ci si fida
degli autori dei file che contiene (\textit{Do you trust the authors of the
files in this folder?}): sono i nostri file, quindi si risponde
\pulsante{Yes, I trust the authors}. Senza questa conferma VS Code apre la
cartella in \textit{modalità con restrizioni} e disattiva le estensioni.

\subsection{Com'è fatta la finestra di VS Code}

La figura~\ref{fig:vscode} mostra VS Code con la cartella del corso
aperta, un file C nell'editor e il terminale in basso. Vale la pena
guardarla con calma: sono le zone che useremo sempre.

% L'editor della figura contiene un listato vero (senza barra delle schede)
\newsavebox{\vsceditor}
\begin{lrbox}{\vsceditor}
\begin{minipage}{12.6cm}
\lstset{backgroundcolor=\color{dmbg}, rulecolor=\color{dmbg}}
\begin{lstlisting}[style=c, aboveskip=0pt, belowskip=0pt]
#include <stdio.h>

int main(void)
{
    printf("Ciao, mondo!\n");
    return 0;
}
\end{lstlisting}
\end{minipage}
\end{lrbox}

\begin{figure}[H]
  \centering
  \begin{tikzpicture}[font=\Lato\scriptsize,
      num/.style={circle, fill=accent!75!black, text=white, draw=white,
        line width=0.8pt, font=\Lato\scriptsize\bfseries, inner sep=0pt,
        minimum size=0.42cm},
      piccolo/.style={font=\Lato\tiny, text=dmfg, inner sep=0pt},
      term/.style={anchor=west, font=\FiraCodeNL\scriptsize, inner sep=0pt}]
    \def\W{17.4} \def\H{10.3}
    % fondo e barra del titolo
    \fill[dmbar, rounded corners=4pt] (0,0) rectangle (\W,-\H);
    \node at (0.38,-0.25) {\vscodelogo[0.3cm]};
    \node[piccolo, anchor=west] at (0.8,-0.25)
      {File\hspace{9pt}Edit\hspace{9pt}Selection\hspace{9pt}View\hspace{9pt}Go\hspace{9pt}Run\hspace{9pt}Terminal\hspace{9pt}Help};
    \fill[dminput, rounded corners=3pt] (6.6,-0.1) rectangle (11.0,-0.41);
    \node[piccolo, text=dmdim] at (8.8,-0.255) {\scalebox{0.7}{\vscico{cerca}}\ programmazione [WSL: Ubuntu]};
    \node[piccolo, anchor=east, font=\Lato\scriptsize] at (\W-0.25,-0.25)
      {--\hspace{15pt}$\square$\hspace{15pt}\texttimes};
    \draw[dmborder] (0,-0.5) -- (\W,-0.5);
    % barra delle attivita'
    \draw[dmborder] (0.75,-0.5) -- (0.75,-\H+0.35);
    \node at (0.375,-0.95) {\vscico[white]{esplora}};
    \fill[dmaccent] (0,-0.72) rectangle (0.045,-1.18);
    \foreach \ico/\y in {cerca/-1.6, git/-2.25, debug/-2.9, estensioni/-3.55}
      \node at (0.375,\y) {\vscico{\ico}};
    \node at (0.375,-\H+1.4) {\vscico{account}};
    \node at (0.375,-\H+0.8) {\vscico{impostazioni}};
    % barra laterale (Explorer)
    \draw[dmborder] (4.1,-0.5) -- (4.1,-\H+0.35);
    \node[piccolo, anchor=west] at (0.95,-0.78) {EXPLORER};
    \node[piccolo, anchor=east] at (3.95,-0.78) {$\cdots$};
    \node[piccolo, anchor=west, font=\Lato\tiny\bfseries] at (0.9,-1.15)
      {\rotatebox[origin=c]{-90}{›}\ PROGRAMMAZIONE [WSL: UBUNTU]};
    \fill[dmsel] (0.76,-1.38) rectangle (4.09,-1.72);
    \draw[dmaccent, line width=0.5pt] (0.77,-1.38) rectangle (4.08,-1.72);
    \node[piccolo, anchor=west, font=\Lato\scriptsize] at (1.15,-1.55)
      {{\FiraCodeNL\bfseries\color{vsciconblue}C}\ \ ciao.c};
    \node[piccolo, anchor=west, font=\Lato\scriptsize] at (1.15,-1.92)
      {{\color{vsctabdim}\ifilecode}\ \ ciao};
    % editor: barra delle schede
    \fill[dmbg] (4.11,-0.51) rectangle (\W,-\H+3.75);
    \fill[dmbar] (4.11,-0.51) rectangle (\W,-0.98);
    \fill[dmbg] (4.11,-0.51) rectangle (6.3,-0.98);
    \fill[dmaccent] (4.11,-0.51) rectangle (6.3,-0.55);
    \node[piccolo, anchor=west, font=\Lato\scriptsize] at (4.3,-0.75)
      {{\FiraCodeNL\bfseries\color{vsciconblue}C}\ \ ciao.c\hspace{9pt}{\color{dmdim}\texttimes}};
    \node at (\W-1.05,-0.75) {\vscico[dmfg]{gioca}};
    \node[piccolo] at (\W-0.7,-0.75) {\rotatebox[origin=c]{-90}{›}};
    \node[piccolo, anchor=west, text=dmdim] at (4.3,-1.2)
      {{\FiraCodeNL\color{vsciconblue}C}\ ciao.c\ \ ›\ \ main};
    \node[anchor=north west, inner sep=0pt] at (4.15,-1.45) {\usebox{\vsceditor}};
    % pannello del terminale
    \draw[dmborder] (4.1,-\H+3.75) -- (\W,-\H+3.75);
    \fill[dmbar] (4.11,-\H+3.74) rectangle (\W,-\H+0.36);
    \node[piccolo, anchor=west, text=dmdim] (t1) at (4.35,-\H+3.42) {PROBLEMS};
    \node[piccolo, anchor=west, text=dmdim] (t2) at ([xshift=12pt]t1.east) {OUTPUT};
    \node[piccolo, anchor=west, text=dmdim] (t3) at ([xshift=12pt]t2.east) {DEBUG CONSOLE};
    \node[piccolo, anchor=west, text=white] (t4) at ([xshift=12pt]t3.east) {TERMINAL};
    \node[piccolo, anchor=west, text=dmdim] at ([xshift=12pt]t4.east) {PORTS};
    \fill[dmaccent] ([yshift=-3pt]t4.south west) rectangle ([yshift=-3.6pt]t4.south east);
    \node[piccolo, anchor=east] at (\W-0.3,-\H+3.42)
      {\iterminal\hspace{3pt}bash\hspace{10pt}+\hspace{10pt}\texttimes};
    \node[term] at (4.35,-\H+2.85) {\textcolor{termgreen}{\bfseries fabio@PC-FABIO}\textcolor{dmfg}{:}\textcolor{termblue}{\bfseries\textasciitilde/programmazione}\textcolor{dmfg}{\$ }\textcolor{white}{gcc -Wall -o ciao ciao.c}};
    \node[term] at (4.35,-\H+2.45) {\textcolor{termgreen}{\bfseries fabio@PC-FABIO}\textcolor{dmfg}{:}\textcolor{termblue}{\bfseries\textasciitilde/programmazione}\textcolor{dmfg}{\$ }\textcolor{white}{./ciao}};
    \node[term, text=dmfg] at (4.35,-\H+2.05) {Ciao, mondo!};
    \node[term] at (4.35,-\H+1.65) {\textcolor{termgreen}{\bfseries fabio@PC-FABIO}\textcolor{dmfg}{:}\textcolor{termblue}{\bfseries\textasciitilde/programmazione}\textcolor{dmfg}{\$ }\textcolor{white}{\rule[-1pt]{4pt}{7pt}}};
    % barra di stato
    \draw[dmborder] (0,-\H+0.35) -- (\W,-\H+0.35);
    \begin{scope}
      \clip[rounded corners=4pt] (0,0) rectangle (\W,-\H);
      \fill[dmaccent] (0,-\H+0.35) rectangle (2.3,-\H);
    \end{scope}
    \node[piccolo, anchor=west, text=white, font=\Lato\tiny] at (0.12,-\H+0.175)
      {\scalebox{0.6}{\vscico[white]{remoto}}\ \ WSL: Ubuntu};
    \node[piccolo, anchor=west] at (2.5,-\H+0.175) {\icoerr\ 0\hspace{5pt}\icowarn\ 0};
    \node[piccolo, anchor=east] at (\W-0.25,-\H+0.175)
      {Ln 5, Col 31\hspace{10pt}Spaces: 4\hspace{10pt}UTF-8\hspace{10pt}LF\hspace{10pt}\{\} C\hspace{10pt}Linux};
    % numeri
    \node[num] at (0.375,-4.4) {1};
    \node[num] at (2.45,-3.0) {2};
    \node[num] at (13.6,-3.0) {3};
    \node[num] at (\W-1.75,-0.75) {4};
    \node[num] at (14.8,-\H+2.2) {5};
    \node[num] at (2.6,-\H+0.6) {6};
  \end{tikzpicture}
  \caption{Visual Studio Code con la cartella \code{programmazione} aperta
           (su Windows, collegato a Ubuntu).}
  \label{fig:vscode}
\end{figure}

\begin{passi}
  \item La \term{barra delle attività}: cambia ciò che compare nella barra
        laterale. Le icone, dall'alto: \textit{Explorer} (i file),
        \textit{Search} (cerca nei file), \textit{Source Control} (Git),
        \textit{Run and Debug} (il debugger), \textit{Extensions}.
  \item La \term{barra laterale}, qui con l'\textit{Explorer}: i file della
        cartella aperta. Un clic apre un file nell'editor; il pulsante con
        il foglio e il segno ``+'', che compare passando con il mouse sul
        nome della cartella, crea un nuovo file.
  \item L'\term{editor}: il file aperto, con i numeri di riga a sinistra. In
        alto ci sono le schede dei file aperti e il ``percorso'' (\textit{ciao.c
        › main}: il cursore è dentro la funzione \code{main}).
  \item Il pulsante \textit{Run}, aggiunto dall'estensione C/C++ (ne
        parleremo alla fine del capitolo).
  \item Il \term{pannello} in basso, qui con il \textit{terminale}
        integrato: un vero terminale (su Windows, di Ubuntu) già posizionato
        nella cartella aperta. Si apre dal menu \menu{Terminal,New
        Terminal}; accanto alla voce del menu è indicata anche la
        scorciatoia da tastiera, che dipende dalla tastiera in uso.
  \item La \term{barra di stato}: a sinistra l'indicatore remoto e il
        numero di errori e avvisi; a destra la posizione del cursore
        (\textit{Ln} riga, \textit{Col} colonna) e il linguaggio
        riconosciuto (\textit{C}).
\end{passi}

\subsection{Creare un file}

Un nuovo file si crea dall'\textit{Explorer} (pulsante ``nuovo file''
accanto al nome della cartella) oppure da \menu{File,New File...}. Il nome
deve finire con \textbf{\code{.c}}: è l'\term{estensione} del nome a dire a
VS Code (e a gcc) che il file contiene codice C. Con un nome come
\code{ciao.txt}, o senza estensione, niente colori, niente suggerimenti e
gcc si rifiuterebbe di compilarlo.

\warning{Nomi di file semplici}{
  Usa nomi brevi, in minuscolo, \textbf{senza spazi} e senza lettere
  accentate: \code{esercizio1.c}, \code{media\_voti.c}. Gli spazi nei nomi
  costringono a usare le virgolette nei comandi del terminale e sono una
  fonte continua di errori.}

\section{Il terminale}

Il terminale è un dialogo scritto con il computer: noi scriviamo un
\textbf{comando} e premiamo \tasto{Invio}, il computer lo esegue e risponde
con del testo. Il programma che legge i comandi si chiama \term{shell}; in
Ubuntu è \textit{bash}, su macOS \textit{zsh}. Per il corso servono pochi
comandi, ma vanno capiti bene: sono il modo in cui compileremo ed eseguiremo
ogni programma.

\subsection{Il prompt}

Quando la shell è pronta mostra il \term{prompt}, una riga che finisce con
\code{\$} (su Mac \code{\%}) e dopo la quale si scrive. Il prompt di
Ubuntu non è decorativo: dice \textit{chi} siamo, \textit{dove} siamo e
che la shell \textit{aspetta} un comando (figura~\ref{fig:prompt}).

\begin{figure}[H]
  \centering
  \begin{tikzpicture}[font=\FiraCodeNL\normalsize,
      pezzo/.style={inner sep=0pt, anchor=base west},
      eti/.style={font=\Lato\scriptsize, text=neutral!20!black, align=center}]
    \fill[consolebg, rounded corners=4pt] (-0.4,0.55) rectangle (8.6,-0.4);
    \node[pezzo, text=termgreen, font=\FiraCodeNL\normalsize\bfseries] (u) at (0,0) {fabio};
    \node[pezzo, text=termgreen, font=\FiraCodeNL\normalsize\bfseries] (at) at (u.base east) {@};
    \node[pezzo, text=termgreen, font=\FiraCodeNL\normalsize\bfseries] (h) at (at.base east) {PC-FABIO};
    \node[pezzo, text=consolefg] (due) at (h.base east) {:};
    \node[pezzo, text=termblue, font=\FiraCodeNL\normalsize\bfseries] (d) at (due.base east) {\textasciitilde/programmazione};
    \node[pezzo, text=consolefg] (s) at (d.base east) {\$};
    \node[pezzo, text=white, font=\FiraCodeNL\normalsize\bfseries] (c) at ([xshift=8pt]s.base east) {ls -l};
    % graffe sotto (punta in basso) per le parti del prompt
    \foreach \p in {u, h, d, s}
      \draw[wtbrace] ([yshift=-0.55cm]\p.base west) -- ([yshift=-0.55cm]\p.base east);
    \node[eti, anchor=north] at ([yshift=-0.75cm]u.base) {l'utente};
    \node[eti, anchor=north] at ([yshift=-0.75cm]h.base) {il nome del\\computer};
    \node[eti, anchor=north] at ([yshift=-0.75cm]d.base) {la \textbf{cartella corrente}\\(\textasciitilde\ è la tua cartella\\personale)};
    \node[eti, anchor=north] at ([yshift=-0.75cm]s.base) {``pronto''};
    % graffa sopra (punta in alto) per il comando
    \draw[wtbrace] ([yshift=0.75cm]c.base east) -- ([yshift=0.75cm]c.base west);
    \node[eti, anchor=south] at ([yshift=0.95cm]c.base) {il comando, scritto da te};
  \end{tikzpicture}
  \caption{Le parti del prompt di Ubuntu.}
  \label{fig:prompt}
\end{figure}

La parte più importante è la \term{cartella corrente} (o \textit{directory
di lavoro}): ogni comando lavora \textit{lì}, a meno di indicare
diversamente. Se scriviamo \code{gcc ciao.c} e il file \code{ciao.c} non
sta nella cartella corrente, gcc non lo trova. Moltissimi problemi dei
principianti si risolvono guardando il prompt.

\info{Il prompt, nella dispensa}{
  D'ora in poi i comandi sono mostrati con il solo \code{\$} come prompt e
  in \textbf{grassetto bianco}; le righe normali sono ciò che il computer
  risponde, e il testo in \textbf{\color{orange-gradient-6!80!black}giallo}
  è ciò che si scrive mentre un programma è in esecuzione. Il \code{\$}
  \textbf{non va digitato}: si scrive solo ciò che lo segue.}

\subsection{Cartelle e percorsi}

I file sono organizzati in cartelle, che contengono altre cartelle, che
contengono altri file: un \term{albero}, che in Linux parte da un'unica
radice chiamata \code{/} (figura~\ref{fig:albero}).

\begin{figure}[H]
  \centering
  \begin{tikzpicture}[font=\Lato\scriptsize]
    \node[fsdir] (root) at (0,0) {/};
    \node[fsdir] (home) at (-4.2,-1.1) {home};
    \node[fsdir] (usr) at (0.6,-1.1) {usr};
    \node[fsdir] (mnt) at (4.4,-1.1) {mnt};
    \node[fsdir] (fabio) at (-4.2,-2.2) {fabio};
    \node[fsqui] (prog) at (-4.2,-3.3) {programmazione};
    \node[fsfile] (ciao) at (-5.6,-4.4) {ciao.c};
    \node[fsfile] (eseg) at (-2.8,-4.4) {ciao};
    \node[fsdir] (bin) at (0.6,-2.2) {bin};
    \node[fsfile] (gcc) at (0.6,-3.3) {gcc};
    \node[fsdir] (c) at (4.4,-2.2) {c};
    \foreach \a/\b in {root/home, root/usr, root/mnt, home/fabio, fabio/prog,
                       prog/ciao, prog/eseg, usr/bin, bin/gcc, mnt/c}
      \draw[fsramo] (\a.south) -- ++(0,-0.2) -| (\b.north);
    \node[mnota, anchor=west, text=primary!50!black] at ([xshift=4pt]fabio.east)
      {\FiraCodeNL\textasciitilde\Lato: la tua cartella personale};
    \node[mnota, anchor=west, text=orange-gradient-6!50!black] at ([xshift=4pt]prog.east)
      {la cartella corrente};
    \node[mnota, anchor=north, align=center] at ([yshift=-2pt]gcc.south)
      {il programma gcc\\è un file anche lui};
    \node[mnota, anchor=north, align=center] at ([yshift=-2pt]c.south)
      {il disco \FiraCodeNL C:\Lato\ di Windows\\(solo con WSL)};
  \end{tikzpicture}
  \caption{Una parte dell'albero delle cartelle di Ubuntu.}
  \label{fig:albero}
\end{figure}

Un file si indica con il suo \term{percorso}, cioè la strada per
raggiungerlo, con i nomi delle cartelle separati da \code{/}:

\begin{itemize}
  \item un percorso \term{assoluto} parte dalla radice e vale ovunque ci si
        trovi: \code{/home/fabio/programmazione/ciao.c};
  \item un percorso \term{relativo} parte dalla cartella corrente: se si è
        in \code{/home/fabio}, lo stesso file è
        \code{programmazione/ciao.c}; se si è già in \code{programmazione},
        è semplicemente \code{ciao.c}.
\end{itemize}

Tre nomi speciali completano il quadro: \code{.} è la cartella corrente,
\code{..} è la cartella che la contiene (``sali di un livello'') e
\code{\textasciitilde} è la cartella personale. Dalla cartella
\code{programmazione}, quindi, \code{..} è \code{/home/fabio}.

\subsection{I comandi essenziali}

\begin{center}
\renewcommand{\arraystretch}{1.4}
\begin{tabular}{@{}>{\FiraCodeNL\small}l>{\raggedright\arraybackslash}p{6.0cm}>{\raggedright\arraybackslash}p{5.6cm}@{}}
  \toprule
  \normalfont\textbf{Comando} & \textbf{Che cosa fa} & \textbf{Da dove viene il nome} \\
  \midrule
  pwd           & stampa il percorso della cartella corrente & \textit{print working directory} \\
  ls            & elenca il contenuto della cartella corrente & \textit{list} \\
  ls -l         & come sopra, con dimensione e data dei file & \textit{long} \\
  cd nome       & entra nella cartella \code{nome} & \textit{change directory} \\
  cd ..         & sale nella cartella superiore & \\
  cd            & torna nella cartella personale & \\
  mkdir nome    & crea la cartella \code{nome} & \textit{make directory} \\
  cp da a       & copia il file \code{da} nel file \code{a} & \textit{copy} \\
  mv da a       & sposta o rinomina un file & \textit{move} \\
  rm file       & cancella un file (\textbf{senza cestino!}) & \textit{remove} \\
  cat file      & mostra il contenuto di un file di testo & \textit{concatenate} \\
  clear         & pulisce lo schermo del terminale & \\
  \bottomrule
\end{tabular}
\end{center}

Ecco una breve sessione di lavoro; la spiegazione segue subito dopo.

\begin{console}
$ pwd
/home/fabio
$ ls
programmazione
$ cd programmazione
$ ls
ciao.c
$ mkdir lezione2
$ ls
ciao.c  lezione2
$ cp ciao.c lezione2/prova.c
$ cd lezione2
$ ls
prova.c
$ cd ..
$ pwd
/home/fabio/programmazione
\end{console}

Si parte dalla cartella personale (\code{pwd} lo conferma), si entra in
\code{programmazione}, si crea la sottocartella \code{lezione2} e ci si
copia il file \code{ciao.c} con il nuovo nome \code{prova.c} (si noti il
percorso relativo \code{lezione2/prova.c}). Poi si entra in
\code{lezione2}, si controlla che la copia ci sia, e con \code{cd ..} si
torna su di un livello. Molti comandi, quando riescono, \textbf{non
stampano niente}: è la convenzione di Linux, ``nessuna notizia, buona
notizia''.

\error{\code{rm} non perdona}{
  Il comando \code{rm} cancella definitivamente: il file non passa per il
  cestino e non si recupera. Controlla sempre due volte il nome prima di
  premere \tasto{Invio}.}

\subsection{Trucchi che fanno risparmiare tempo}

\begin{itemize}
  \item \tasto{Tab} \textbf{completa} i nomi: scrivendo \code{cd prog} e
        premendo \tasto{Tab}, la shell completa da sola in
        \code{cd programmazione/}. Oltre a risparmiare tempo evita gli
        errori di battitura: se il completamento non funziona, il file non
        si trova dove pensiamo.
  \item \tasto{$\uparrow$} e \tasto{$\downarrow$} scorrono i \textbf{comandi già dati}: per
        ricompilare dopo una modifica basta premere \tasto{$\uparrow$} fino a
        ritrovare il comando \code{gcc} e poi \tasto{Invio}.
  \item \tasti{Ctrl,C} \textbf{interrompe} il programma in esecuzione. Se un
        programma non termina mai (capiterà, con i cicli infiniti del
        capitolo~\ref{chap:flusso}), è il modo per fermarlo.
  \item \tasti{Ctrl,D}, scritto come input a un programma, segnala la
        \textbf{fine dell'input}; lo useremo dal
        capitolo~\ref{chap:primi-programmi}.
  \item Nel terminale \tasti{Ctrl,C} non copia: per copiare e incollare,
        nel terminale di VS Code si usano il tasto destro del mouse o
        \tasti{Ctrl,Shift,C} e \tasti{Ctrl,Shift,V}.
\end{itemize}

\section{Il primo programma}

È il momento di mettere insieme i pezzi. Nella cartella
\code{programmazione} crea il file \code{ciao.c} e scrivici, esattamente
così, il programma seguente.

\begin{lstlisting}[style=c, name=ciao.c]
#include <stdio.h>

int main(void)
{
    printf("Ciao, mondo!\n");
    return 0;
}
\end{lstlisting}

Il significato di ogni riga sarà spiegato nel
capitolo~\ref{chap:primi-programmi}; per ora basta sapere che
\code{printf} stampa il testo fra virgolette, e che \code{\textbackslash n}
significa ``vai a capo''. Prima di proseguire, \textbf{salva} il file con
\tasti{Ctrl,S} (su Mac \tasti{Cmd,S}).

\begin{figure}[H]
  \centering
  \begin{tikzpicture}[font=\Lato\scriptsize]
    \foreach \x/\segno in {0/{\color{white}$\bullet$}, 7/{\color{dmdim}\texttimes}} {
      \fill[dmbar] (\x,0) rectangle ++(5,-0.5);
      \fill[dmbg] (\x,0) rectangle ++(2.3,-0.5);
      \fill[dmaccent] (\x,0) rectangle ++(2.3,-0.04);
      \node[anchor=west, text=white, font=\Lato\scriptsize] at (\x+0.2,-0.26)
        {{\FiraCodeNL\bfseries\color{vsciconblue}C}\ \ ciao.c\hspace{9pt}\segno};
    }
    \node[anchor=north, align=center] at (2.5,-0.65) {un pallino: modifiche \textbf{non salvate}};
    \node[anchor=north, align=center] at (9.5,-0.65) {una crocetta: il file è \textbf{salvato}};
  \end{tikzpicture}
  \caption{La scheda del file dice se ci sono modifiche non salvate.}
  \label{fig:salvato}
\end{figure}

gcc legge il file \textbf{dal disco}: se le modifiche non sono salvate,
compila la versione vecchia. È l'errore più frequente all'inizio
(``ho corretto ma non cambia niente!''). Per evitarlo una volta per tutte,
si può attivare il salvataggio automatico dal menu \menu{File,Auto Save}.

\subsection{Compilare ed eseguire}

Ora, nel terminale, compiliamo ed eseguiamo:

\begin{console}
$ gcc ciao.c
$ ls
a.out  ciao.c
$ ./a.out
Ciao, mondo!
\end{console}

\begin{itemize}
  \item \code{gcc ciao.c} compila il file. Non stampa nulla: nessun
        messaggio significa \textbf{nessun errore}. Il risultato è un nuovo
        file, l'eseguibile, che per tradizione si chiama \code{a.out}.
  \item \code{./a.out} esegue il programma, che stampa la sua riga.
\end{itemize}

Il nome \code{a.out} va bene per una prova, ma con più programmi nella
stessa cartella ogni compilazione sovrascriverebbe la precedente. Con
l'opzione \code{-o} (\textit{output}) si sceglie il nome dell'eseguibile; e
con \code{-Wall} (\textit{warnings: all}) si chiede a gcc di segnalare
anche le cose sospette, che non impediscono la compilazione ma spesso
nascondono un errore. È il modo in cui conviene compilare sempre:

\begin{console}
$ gcc -Wall -o ciao ciao.c
$ ./ciao
Ciao, mondo!
\end{console}

Ogni volta che si cambia il programma bisogna ripetere \textbf{tutto il
giro} (figura~\ref{fig:ciclo}): l'eseguibile è una traduzione fatta in un
certo momento, e non si aggiorna da solo quando cambia il testo.

\begin{figure}[H]
  \centering
  \begin{tikzpicture}[font=\Lato\footnotesize,
      fase/.style={wtnode, minimum width=3.0cm, minimum height=1.05cm}]
    \node[fase] (m) at (0,0) {\textbf{modifica}\\[1pt]{\scriptsize nell'editor}};
    \node[fase] (s) at (4.6,0) {\textbf{salva}\\[1pt]\tasti{Ctrl,S}};
    \node[fase, wtnodeacc] (c) at (9.2,0) {\textbf{compila}\\[1pt]{\FiraCodeNL\scriptsize gcc -Wall -o ciao ciao.c}};
    \node[fase, wtnodeok] (e) at (13.8,0) {\textbf{esegui}\\[1pt]{\FiraCodeNL\scriptsize ./ciao}};
    \draw[wtarrow] (m) -- (s);
    \draw[wtarrow] (s) -- (c);
    \draw[wtarrow] (c) -- (e);
    \draw[wtarrow] (e.south) -- ++(0,-0.55) -| node[wtlab, pos=0.25] {il risultato non è quello atteso? si ricomincia} (m.south);
    \draw[wtarrowacc] (c.north) -- ++(0,0.5) -| node[wtlab, pos=0.25] {errori di compilazione: si corregge il testo} (m.north);
  \end{tikzpicture}
  \caption{Il ciclo di lavoro: modifica, salva, compila, esegui.}
  \label{fig:ciclo}
\end{figure}

\subsection{Perché \code{./}?}

Perché per eseguire il programma si scrive \code{./ciao} e non solo
\code{ciao}, mentre per gcc basta \code{gcc}? Quando si scrive il nome di un
comando, la shell non lo cerca dappertutto: lo cerca solo in un elenco di
cartelle prestabilite, quelle dove stanno i programmi installati (come
\code{/usr/bin}, dove si trova \code{gcc}). L'elenco è una
\textit{variabile d'ambiente} chiamata \code{PATH}, e per sicurezza non
contiene la cartella corrente. Scrivendo \code{./ciao} si dà il
\textbf{percorso} del file invece del nome di un comando: ``il programma
\code{ciao} che sta \textit{qui}'' (\code{.} è la cartella corrente).

\begin{console}
$ ciao
ciao: command not found
$ ./ciao
Ciao, mondo!
\end{console}

\subsection{Gli errori di compilazione}

Che cosa succede se il programma contiene un errore? Togliamo il punto e
virgola alla fine della riga di \code{printf} e ricompiliamo:

\begin{affianco}[0.42][c]
\begin{lstlisting}[style=c, name=errore.c]
#include <stdio.h>

int main(void)
{
    printf("Ciao, mondo!\n")
    return 0;
}
\end{lstlisting}
\risultato
\begin{console}
$ gcc -Wall -o errore errore.c
errore.c: In function (*@‘@*)main(*@’@*):
errore.c:5:29: error: expected (*@‘@*);(*@’@*) before (*@‘@*)return(*@’@*)
    5 |     printf("Ciao, mondo!\n")
      |                             ^
      |                             ;
    6 |     return 0;
      |     ~~~~~~
\end{console}
\end{affianco}

Il messaggio va letto con calma, perché dice quasi tutto
(figura~\ref{fig:errore}): \textbf{dove} è l'errore (file, riga e colonna),
\textbf{di che tipo} è e \textbf{che cosa} si aspettava il compilatore. Le
righe successive ricopiano il codice e indicano con \code{\^{}} il punto
esatto; gcc arriva perfino a suggerire il \code{;} mancante. Quando c'è un
\code{error} l'eseguibile \textbf{non viene prodotto}: bisogna correggere e
ricompilare.

\begin{figure}[H]
  \centering
  \begin{tikzpicture}[font=\FiraCodeNL\normalsize,
      pezzo/.style={inner sep=0pt, anchor=base west, text=consolefg},
      eti/.style={font=\Lato\scriptsize, text=neutral!20!black, align=center}]
    \fill[consolebg, rounded corners=4pt] (-0.4,0.55) rectangle (11.6,-0.4);
    \node[pezzo] (f) at (0,0) {errore.c};
    \node[pezzo] (d1) at (f.base east) {:};
    \node[pezzo] (r) at (d1.base east) {5:29};
    \node[pezzo] (d3) at (r.base east) {:};
    \node[pezzo, text=red-gradient-2, font=\FiraCodeNL\normalsize\bfseries] (t) at ([xshift=6pt]d3.base east) {error:};
    \node[pezzo] (m) at ([xshift=6pt]t.base east) {expected ‘;’ before ‘return’};
    % sotto (punta in basso): file, riga e colonna, messaggio
    \foreach \p in {f, r, m}
      \draw[wtbrace] ([yshift=-0.55cm]\p.base west) -- ([yshift=-0.55cm]\p.base east);
    \node[eti, anchor=north] at ([yshift=-0.75cm]f.base) {il file};
    \node[eti, anchor=north] at ([yshift=-0.75cm]r.base) {riga 5,\\colonna 29};
    \node[eti, anchor=north] at ([yshift=-0.75cm]m.base) {che cosa non va: ``mi aspettavo\\un \FiraCodeNL;\Lato\ prima di \FiraCodeNL return\Lato''};
    % sopra (punta in alto): il tipo di messaggio
    \draw[wtbrace] ([yshift=0.75cm]t.base east) -- ([yshift=0.75cm]t.base west);
    \node[eti, anchor=south] at ([yshift=0.95cm]t.base) {il tipo: \textbf{error} (errore) o \textbf{warning} (avviso)};
  \end{tikzpicture}
  \caption{Anatomia di un messaggio di errore di gcc.}
  \label{fig:errore}
\end{figure}

Si noti che l'errore è segnalato alla riga \textbf{5}, alla fine della
riga di \code{printf}, anche se il compilatore se ne accorge solo leggendo
\code{return} alla riga 6: il punto e virgola poteva mancare solo lì. In
generale l'errore sta nella riga indicata \textbf{o poco prima}.

Esistono due tipi di messaggi:

\begin{itemize}
  \item un \code{error} è un errore vero: il testo non è C valido e
        l'eseguibile non viene creato;
  \item un \code{warning} (avviso) segnala qualcosa di lecito ma sospetto:
        l'eseguibile viene creato, ma quasi sempre il warning indica un
        errore di ragionamento. \textbf{I warning vanno trattati come
        errori}: un programma compilato con \code{-Wall} non deve
        produrne nessuno.
\end{itemize}

\info{Leggere solo il primo errore}{
  Un solo errore può confondere il compilatore e generarne a cascata molti
  altri, alcuni dei quali senza senso. La regola pratica: si corregge
  \textbf{il primo} messaggio dell'elenco, si ricompila, e si ripete.}

Con l'estensione C/C++ installata, VS Code si accorge di molti errori
\textbf{prima} della compilazione e li sottolinea con un'ondina rossa;
passando con il mouse sopra la sottolineatura compare la spiegazione
(figura~\ref{fig:ondina}). Il numero di errori compare anche nella barra di
stato, accanto a \icoerr[black].

\begin{figure}[H]
  \centering
  \begin{tikzpicture}[font=\FiraCodeNL\scriptsize,
      riga/.style={anchor=base west, inner sep=0pt},
      n/.style={anchor=base east, inner sep=0pt, text=vsclinenum, font=\FiraCodeNL\tiny}]
    \fill[dmbg, rounded corners=3pt] (0,0.35) rectangle (11.5,-1.95);
    \node[n] at (0.45,0) {4};
    \node[riga, text=vscfg] at (0.7,0) {\{};
    \node[n] at (0.45,-0.42) {5};
    \node[riga] at (0.7,-0.42) {\ \ \ \ {\color{vscfunc}printf}{\color{vscfg}(}{\color{vscstring}"Ciao, mondo!\textbackslash n"}{\color{vscfg})}};
    \node[n] at (0.45,-0.84) {6};
    \node[riga] (sp) at (0.7,-0.84) {\ \ \ \ };
    \node[riga, text=vscpurple] (ret) at (sp.base east) {return};
    \node[riga, text=vscfg] at (ret.base east) {\ 0;};
    \draw[red-gradient-3, ondina] ([yshift=-1.6pt]ret.base west) -- ([yshift=-1.6pt]ret.base east);
    \node[anchor=north west, fill=dmbar, draw=dmborder, text=dmfg, font=\Lato\scriptsize,
          inner sep=5pt, rounded corners=2pt, align=left] at ([yshift=-6pt]ret.base west)
      {expected a ';' \ \ {\color{dmdim}C/C++(65)}};
  \end{tikzpicture}
  \caption{L'estensione C/C++ sottolinea l'errore mentre si scrive.}
  \label{fig:ondina}
\end{figure}

\subsection{E il pulsante “Run”?}

L'estensione C/C++ aggiunge in alto a destra nell'editor un pulsante
\vscico[black]{gioca} (\textit{Run C/C++ File}) che compila ed esegue il
file aperto con un clic; la prima volta chiede quale compilatore usare (si
sceglie \textit{gcc}). È comodo, ma in questa dispensa compileremo
\textbf{sempre dal terminale}: è ciò che si fa a lezione e all'esame,
permette di scegliere le opzioni di gcc, e soprattutto fa vedere che cosa
succede davvero, passo per passo. Dal capitolo~\ref{chap:make} useremo
\code{make}, che si lancia anch'esso dal terminale.

\section{Quando qualcosa non funziona}

\begin{center}
\renewcommand{\arraystretch}{1.45}
\begin{tabular}{@{}>{\raggedright\arraybackslash}p{5.6cm}>{\raggedright\arraybackslash}p{10.6cm}@{}}
  \toprule
  \textbf{Sintomo} & \textbf{Causa probabile e rimedio} \\
  \midrule
  \code{gcc: command not found} &
    Il compilatore non è installato (ripeti il passo 2). Su Windows,
    controlla di essere nel terminale di \textbf{Ubuntu} (prompt con
    \code{\$}) e non in PowerShell (prompt \code{PS C:\textbackslash...}). \\
  \code{code: command not found} &
    Su Windows chiudi e riapri il terminale di Ubuntu dopo aver installato
    VS Code (il \code{PATH} si aggiorna solo all'apertura); su Mac esegui
    \textit{Shell Command: Install `code' command in PATH}. \\
  \codel{ciao.c: No such file or directory} &
    Il file non è nella cartella corrente: controlla con \code{pwd} e
    \code{ls} dove ti trovi e che cosa c'è, oppure il file non è stato
    salvato con quel nome. \\
  \codel{./ciao: No such file or directory} &
    Il programma non è stato compilato, oppure la compilazione è fallita
    (c'erano degli \code{error}), oppure l'eseguibile ha un altro nome
    (\code{a.out}?). \\
  Ho corretto ma il programma fa ancora la cosa sbagliata &
    Hai salvato? Hai ricompilato? Stai eseguendo l'eseguibile giusto? \\
  Il programma non termina &
    Probabilmente un ciclo infinito, oppure il programma aspetta un input
    da tastiera. \tasti{Ctrl,C} lo interrompe. \\
  In VS Code \code{\#include <stdio.h>} è sottolineato in rosso &
    L'estensione C/C++ non trova il compilatore. Su Windows: VS Code non
    è collegato a Ubuntu (guarda l'indicatore remoto) o l'estensione non è
    installata \textit{in WSL}. Altrove: verifica che \code{gcc -{}-version}
    funzioni. \\
  \bottomrule
\end{tabular}
\end{center}

\insintesi{
  Per scrivere ed eseguire un programma C servono tre strumenti: un
  \textbf{editor} (VS Code) per scrivere il testo del programma in un file
  \code{.c}, un \textbf{compilatore} (gcc) che lo traduce in linguaggio
  macchina, e un \textbf{terminale} da cui lanciare il compilatore e il
  programma. Su Windows tutto ciò vive in Ubuntu, grazie a WSL. Il ciclo di
  lavoro è sempre lo stesso: \textit{modifica, salva, compila, esegui}; e
  i messaggi di gcc indicano file, riga e colonna del problema.}

\section{Riferimenti}

\begin{itemize}
  \item \textbf{Installare Linux in Windows con WSL} --- Microsoft Learn. \\
        \urlx{https://learn.microsoft.com/it-it/windows/wsl/install}
  \item \textbf{Using C++ and WSL in VS Code} --- documentazione di VS Code
        (vale anche per il C). \\
        \urlx{https://code.visualstudio.com/docs/cpp/config-wsl}
  \item \textbf{Using Clang in Visual Studio Code} (macOS) e \textbf{Using
        C++ on Linux in VS Code} --- documentazione di VS Code. \\
        \urlx{https://code.visualstudio.com/docs/cpp/config-clang-mac} \\
        \urlx{https://code.visualstudio.com/docs/cpp/config-linux}
  \item \textbf{User Interface} --- la guida alle zone della finestra di
        VS Code. \\
        \urlx{https://code.visualstudio.com/docs/getstarted/userinterface}
\end{itemize}


  \chapter{Il calcolatore}
\label{chap:calcolatore}

Per imparare a guidare non serve saper costruire un motore, ma aiuta molto
sapere che cosa fanno il volante, i pedali e il cambio. Allo stesso modo, per
programmare non serve conoscere l'elettronica di un computer, ma bisogna
avere un'idea chiara di come è fatto e di come esegue un programma: molte
regole del C, che altrimenti sembrerebbero capricci, diventano ovvie. Questo
capitolo costruisce quell'idea, partendo da che cosa significa
\textit{programmare}.

\section{Che cos'è la programmazione}

Il corso si apre con la definizione del dizionario Treccani:

\begin{center}
\begin{minipage}{0.86\linewidth}
\itshape “In informatica, la stesura di un programma in una forma
interpretabile dal calcolatore elettronico (o direttamente o dopo una
rielaborazione opportuna), compiuta dal programmatore sulla base delle
specifiche di programma fornite dall'analista; la programmazione comprende
la definizione del problema da risolvere, la preparazione del diagramma di
flusso del programma, la sua codifica (cioè la trascrizione, in un
linguaggio opportuno, delle varie istruzioni), la messa a punto e la
preparazione di una documentazione a uso dell'operatore.”
\end{minipage}
\end{center}

La definizione contiene già tutto il percorso del corso. Scrivere codice
(la \textit{codifica}) è solo una delle fasi, e non la prima: prima bisogna
capire il problema e progettare la soluzione, dopo bisogna verificare che
funzioni (la \textit{messa a punto}) e spiegarla (figura~\ref{fig:fasi}).
La parentesi “o direttamente o dopo una rielaborazione opportuna” è
importante: il programma che scriviamo non è ancora ciò che il computer
esegue, va \textit{rielaborato}, cioè tradotto. È il compito del
compilatore incontrato nel capitolo precedente.

\begin{figure}[H]
  \centering
  \begin{tikzpicture}[font=\Lato\scriptsize,
      fase/.style={wtnode, font=\Lato\scriptsize, text width=2.55cm,
        minimum height=1.5cm}]
    \node[fase] (a) at (0,0) {\textbf{definire il problema}\\[2pt]che cosa si deve ottenere, a partire da quali dati};
    \node[fase] (b) at (3.6,0) {\textbf{progettare}\\[2pt]l'algoritmo, per esempio con un diagramma di flusso};
    \node[fase, draw=accent!60!black, fill=accent!8!white] (c) at (7.2,0) {\textbf{codificare}\\[2pt]scrivere il programma in un linguaggio (il C)};
    \node[fase] (d) at (10.8,0) {\textbf{mettere a punto}\\[2pt]provare, trovare e correggere gli errori};
    \node[fase] (e) at (14.4,0) {\textbf{documentare}\\[2pt]spiegare come si usa e come funziona};
    \foreach \x/\y in {a/b, b/c, c/d, d/e} \draw[wtarrow] (\x) -- (\y);
    \draw[wtarrowacc] (d.south) -- ++(0,-0.4) -| node[wtlab, pos=0.25] {gli errori trovati riportano indietro} (b.south);
  \end{tikzpicture}
  \caption{Le fasi della programmazione secondo la definizione: scrivere il
           codice è solo una di esse.}
  \label{fig:fasi}
\end{figure}

\info{Come si impara a programmare}{
  A lezione il docente avverte che il corso parte da zero e all'inizio va
  piano, ma che velocità e difficoltà crescono presto, e che entro il primo
  mese quasi tutti incontrano difficoltà serie, soprattutto negli esercizi.
  È normale e si supera, con alcune abitudini:
  \begin{itemize}
    \item non dare per scontato di saper già programmare, e non saltare le
          lezioni “facili”;
    \item studiare con costanza: la programmazione si accumula, ogni
          argomento usa i precedenti;
    \item \textbf{scrivere ed eseguire} ogni esempio, poi modificarlo per
          vedere che cosa succede: leggere un programma non basta, come non
          basta guardare qualcuno nuotare;
    \item chiedere appena qualcosa non è chiaro (“non siete gli unici”);
    \item prendere appunti e fare sintesi: le slide sono, per loro natura,
          sintetiche; questa dispensa prova a riempire gli spazi, ma non
          sostituisce le lezioni.
  \end{itemize}}

\section{Lo schema di un calcolatore}

Un computer moderno è un oggetto complicatissimo, ma il suo funzionamento si
capisce con uno schema di tre soli blocchi (figura~\ref{fig:vonneumann}):

\begin{itemize}
  \item la \term{CPU} (\textit{Central Processing Unit}, il
        \textit{processore}), che esegue le istruzioni;
  \item la \term{memoria centrale} (o \textit{principale}, la RAM), che
        contiene le istruzioni da eseguire e i dati su cui lavorano;
  \item i \term{dispositivi di ingresso e uscita} (\textit{I/O},
        \textit{Input/Output}): tastiera, mouse, schermo, disco, rete\dots,
        collegati attraverso delle \textit{interfacce}.
\end{itemize}

I blocchi comunicano attraverso il \term{bus}, un insieme di collegamenti
condiviso su cui viaggiano dati, indirizzi e comandi. Fisicamente, è la
\textit{scheda madre} (\textit{motherboard}) a collegare fra loro processore,
memoria e periferiche.

\begin{figure}[H]
  \centering
  \begin{tikzpicture}[font=\Lato\scriptsize,
      dev/.style={wtnodeplain, minimum width=2.3cm, minimum height=0.75cm},
      intf/.style={draw=green-gradient-6, fill=green-gradient-1!45!white,
        line width=0.7pt, minimum width=2.3cm, minimum height=0.45cm,
        font=\Lato\tiny},
      blk/.style={wtnode, minimum width=3.4cm, minimum height=1.2cm}]
    \foreach \x/\n in {0/tastiera, 3/mouse, 6/schermo, 9/disco (memoria secondaria)} {
      \node[dev, text width=2.1cm] (d\x) at (\x,2.2) {\n};
      \node[intf] (i\x) at (\x,1.1) {interfaccia di I/O};
      \draw[wtarrow, <->] (d\x) -- (i\x);
      \draw[wtarrow, <->] (i\x) -- (\x,0.22);
    }
    \fill[orange-gradient-4] (-1.4,0.2) rectangle (10.4,-0.2);
    \node[font=\Lato\scriptsize\bfseries, text=orange-gradient-6!35!black] at (4.5,0) {bus};
    \node[blk, draw=red-gradient-6!70!black, fill=red-gradient-1!25!white] (mem) at (2.2,-1.6) {\textbf{memoria centrale}\\istruzioni e dati};
    \node[blk, draw=accent!60!black, fill=accent!8!white] (cpu) at (6.8,-1.6) {\textbf{CPU}\\esegue le istruzioni};
    \draw[wtarrow, <->] (mem) -- (2.2,-0.22);
    \draw[wtarrow, <->] (cpu) -- (6.8,-0.22);
  \end{tikzpicture}
  \caption{Schema semplificato di un calcolatore (architettura di von
           Neumann): CPU, memoria e dispositivi di I/O collegati dal bus.}
  \label{fig:vonneumann}
\end{figure}

Questo schema porta il nome di \textit{John von Neumann}, che nel 1945 ne
descrisse l'idea centrale: le \textbf{istruzioni} del programma stanno
nella \textbf{stessa memoria} dei \textbf{dati}. Un programma non è cablato
nei circuiti: è una sequenza di numeri in memoria, che si può caricare,
sostituire, copiare. È questo che rende un computer una macchina
\textit{generale}, capace di fare cose diversissime a seconda del programma
caricato.

\section{La CPU e il linguaggio macchina}

\subsection{Che cosa sa fare un processore}

Tutte le elaborazioni di un calcolatore sono svolte dalla CPU, oppure da
altri componenti su suo comando. Sorprendentemente, un processore sa
compiere solo operazioni \textbf{molto semplici}:

\begin{itemize}
  \item leggere, scrivere o copiare una o più celle di memoria;
  \item sommare, sottrarre, moltiplicare, dividere il contenuto di celle di
        memoria;
  \item leggere e scrivere in zone di memoria “speciali” per pilotare i
        dispositivi di I/O (per esempio la scheda video);
  \item poche altre operazioni elementari, come confrontare due numeri e
        \textbf{saltare} a un'altra istruzione a seconda del risultato.
\end{itemize}

Il segreto è la velocità: un processore moderno esegue \textbf{miliardi} di
queste operazioni al secondo. Ogni cosa che un computer fa, dal mostrare
una pagina web al riprodurre un video, è scomposta in una lunghissima
sequenza di passi elementari.

\subsection{Il linguaggio macchina}

Ogni processore ha un proprio insieme di istruzioni, e ogni istruzione è
identificata da una certa configurazione di bit (di zeri e uni). L'insieme
delle istruzioni di un processore si chiama \term{linguaggio macchina} di
quel processore. Processori di famiglie diverse hanno linguaggi macchina
diversi: un programma tradotto per i processori Intel e AMD dei PC (famiglia
\textit{x86-64}) non gira su un Mac con processore Apple (famiglia
\textit{ARM}), e viceversa.

Per eseguire un programma, allora, bisogna fare due cose: mettere in
memoria la sequenza di configurazioni di bit che rappresentano le sue
istruzioni, e dire al processore a quale indirizzo si trova la prima. Da lì
in poi il processore ripete sempre lo stesso ciclo
(figura~\ref{fig:ciclo-cpu}): \textbf{preleva} l'istruzione all'indirizzo
indicato da un suo registro speciale, il \term{program counter};
\textbf{decodifica} l'istruzione, cioè capisce che cosa chiede;
la \textbf{esegue}; e passa all'istruzione successiva, a meno che
l'istruzione appena eseguita non fosse un \textit{salto}, che cambia il
program counter.

\begin{figure}[H]
  \centering
  \begin{tikzpicture}[font=\Lato\scriptsize]
    % memoria
    \memoria{m}{(0,0)}{1000/MOV \ldots/, 1003/ADD \ldots/, 1005/JMP 1000/, 1008/HLT/, \ldots/\ldots/}
    \node[mnota, anchor=south] at ([yshift=2pt]m-1.north) {\textbf{memoria}};
    \node[mcellam, fill=yellow-gradient-1!70!white] at (m-2) {ADD \ldots};
    % CPU
    \node[wtnode, draw=accent!60!black, fill=accent!6!white, minimum width=4.6cm,
          minimum height=2.75cm, minimum width=5.3cm, anchor=north west] (cpu) at (3.0,0.55) {};
    \node[font=\Lato\scriptsize\bfseries, anchor=north west] at ([shift={(5pt,-4pt)}]cpu.north west) {CPU};
    \node[mvar] (pc) at (4.55,-0.49) {1003};
    \node[mnome, anchor=west] at ([xshift=3pt]pc.east) {program counter};
    \node[mvar] (ax) at (4.55,-1.2) {42};
    \node[mnome, anchor=west] at ([xshift=3pt]ax.east) {registro AX};
    \node[mvar] (bx) at (4.55,-1.8) {8};
    \node[mnome, anchor=west] at ([xshift=3pt]bx.east) {registro BX};
    \draw[mrif] (pc.west) -- ([xshift=8pt]m-2.east);
    \node[mpunto] at (pc.west) {};
    % ciclo
    \begin{scope}[shift={(10.9,-0.95)}]
      \node[wtnodeplain] (f) at (90:1.25) {1. preleva};
      \node[wtnodeplain] (d) at (-30:1.25) {2. decodifica};
      \node[wtnodeplain] (e) at (210:1.25) {3. esegue};
      \draw[wtarrow] (f) to[bend left=30] (d);
      \draw[wtarrow] (d) to[bend left=30] (e);
      \draw[wtarrow] (e) to[bend left=30] (f);
    \end{scope}
  \end{tikzpicture}
  \caption{Il program counter indica l'indirizzo della prossima istruzione;
           la CPU la preleva, la decodifica e la esegue, poi ricomincia. I
           registri sono piccole celle di memoria interne alla CPU.}
  \label{fig:ciclo-cpu}
\end{figure}

\section{L'assembly}

Il linguaggio macchina è una sequenza di bit, praticamente illeggibile per
una persona. Il primo passo per renderlo comprensibile è l'\term{assembly}:
un linguaggio quasi identico al linguaggio macchina, in cui però ogni codice
operativo è sostituito da un'abbreviazione (un \textit{codice mnemonico}):
\code{MOV} per \textit{move} (copia), \code{ADD} per \textit{add} (somma),
\code{JMP} per \textit{jump} (salta). Un programma, l'\term{assembler},
traduce l'assembly in linguaggio macchina, un'istruzione alla volta.

\subsection{Un esempio: contare i bit a 1}

Il programma seguente, visto a lezione, conta quanti bit valgono 1 in un
dato di 16 bit. A sinistra di ogni istruzione c'è il suo indirizzo in
memoria (in esadecimale, cioè in base 16: per ora basta considerarlo un
“numero civico”), a destra la spiegazione.

\begin{center}
\renewcommand{\arraystretch}{1.25}
{\FiraCodeNL\scriptsize
\begin{tabular}{@{}l l l >{\Lato\scriptsize}l@{}}
  \toprule
  \normalfont\Lato\bfseries indirizzo & \normalfont\Lato\bfseries istruzione & & \normalfont\Lato\bfseries significato \\
  \midrule
  3F0D0100 & 00000000 00000000 & & il dato da elaborare (16 bit) \\
  3F0D0102 & 00000000 & & qui andrà il risultato (8 bit) \\
  \midrule
  40B1A200 & MOV & \$0, \%AL         & azzera il contatore, il registro AL \\
  40B1A203 & MOV & (3F0D0100), \%DX  & copia nel registro DX il dato da elaborare \\
  40B1A207 & CMP & \%DX, \$0         & confronta DX con 0\dots \\
  40B1A211 & JE  & 40B1A232          & \dots e se sono uguali salta alla fine \\
  40B1A215 & SHL & \%DX              & sposta i bit di DX di un posto a sinistra: \\
           &     &                   & il bit che “esce” finisce nel \textit{carry flag} CF \\
  40B1A217 & JC  & 40B1A225          & se CF vale 1 (è uscito un 1) salta all'incremento \\
  40B1A221 & JMP & 40B1A207          & altrimenti torna al confronto \\
  40B1A225 & ADD & \$1, \%AL         & incrementa il contatore AL \\
  40B1A228 & JMP & 40B1A207          & torna al confronto \\
  40B1A232 & MOV & \%AL, (3F0D0102)  & scrive il risultato in memoria \\
  40B1A236 & HLT &                   & fine (\textit{halt}) \\
  \bottomrule
\end{tabular}}
\end{center}

L'idea è questa: a ogni giro si fa “scivolare” il dato di un posto a
sinistra, e si guarda il bit che cade fuori dal bordo; se è un 1, si
incrementa il contatore. Quando nel dato non resta nessun 1 (cioè vale 0),
si è finito. Seguiamo l'esecuzione con un dato che contiene due bit a 1,
\code{1010000000000000}:

\begin{center}
\renewcommand{\arraystretch}{1.3}
\begin{tabular}{@{}c>{\FiraCodeNL\small}c>{\FiraCodeNL\small}c>{\FiraCodeNL\small}c>{\FiraCodeNL\small}c>{\raggedright\arraybackslash}p{4.4cm}@{}}
  \toprule
  \textbf{giro} & \normalfont\textbf{DX prima} & \normalfont\textbf{bit uscito (CF)} & \normalfont\textbf{DX dopo} & \normalfont\textbf{AL} & \textbf{che cosa succede} \\
  \midrule
  1 & 1010000000000000 & 1 & 0100000000000000 & 1 & è uscito un 1: \code{JC} salta all'\code{ADD} \\
  2 & 0100000000000000 & 0 & 1000000000000000 & 1 & è uscito uno 0: \code{JMP} torna al confronto \\
  3 & 1000000000000000 & 1 & 0000000000000000 & 2 & è uscito un 1: si incrementa \\
  4 & 0000000000000000 & & & 2 & DX vale 0: \code{JE} salta alla fine \\
  \bottomrule
\end{tabular}
\end{center}

Il risultato, 2, viene scritto in memoria e il programma si ferma. Si noti
che in assembly non esiste un'istruzione “ripeti finché”: la ripetizione
si ottiene \textbf{saltando indietro} (\code{JMP 40B1A207}) e la scelta
\textbf{saltando avanti} se una condizione è vera (\code{JE}, \code{JC}).

\subsection{I limiti dell'assembly}

Con l'assembly si può scrivere qualunque programma, anche un sistema
operativo o un'interfaccia grafica, ma è un lavoro enorme: le istruzioni
sono di \textbf{basso livello}, cioè vicine alla macchina e lontane dal
modo in cui ragioniamo. Per “contare gli uni” abbiamo dovuto pensare in
termini di registri, flag e indirizzi di salto; e il programma funziona
solo sul processore per cui è scritto.

\section{I livelli di astrazione}

Servono istruzioni più vicine al modo di ragionare umano. La soluzione è
l'\term{astrazione}: costruire un linguaggio che “nasconde” la macchina
sottostante e permette di esprimere gli algoritmi in modo simbolico. Sono i
\term{linguaggi di alto livello}, come il C, che hanno due proprietà
fondamentali:

\begin{itemize}
  \item sono più vicini al linguaggio umano (e matematico) che al
        linguaggio macchina;
  \item sono \textbf{generali}: non dipendono dal processore. Lo stesso
        programma C si traduce per un PC, per un Mac o per un telefono.
\end{itemize}

La figura~\ref{fig:livelli} mostra la stessa istruzione ai tre livelli. È
la traduzione reale prodotta da gcc su un PC (processore x86-64) per
l'istruzione C \code{c = a + b;}, in cui \code{a}, \code{b} e \code{c} sono
tre variabili intere.

\begin{figure}[H]
  \centering
  \begin{tikzpicture}[font=\Lato\scriptsize,
      liv/.style={rounded corners=4pt, line width=0.8pt, minimum width=17.2cm,
        inner sep=6pt, align=left, anchor=north west},
      tit/.style={font=\Lato\scriptsize\bfseries, anchor=north west}]
    \node[liv, draw=green-gradient-6, fill=green-gradient-1!35!white] (alto) at (0,0)
      {\begin{tabular}{@{}p{3.5cm}l@{}}
         \textbf{linguaggio di alto livello} & {\FiraCodeNL\small c = a + b;} \\
         \textit{si ragiona sul problema} & \\
       \end{tabular}};
    \node[liv, draw=primary!65!black, fill=primary!6!white] (asm) at ([yshift=-0.35cm]alto.south west)
      {\begin{tabular}{@{}p{3.5cm}>{\FiraCodeNL\small}l>{\Lato\scriptsize}l@{}}
         \textbf{assembly} & mov edx, [rbp-0xc] & copia \code{a} nel registro edx \\
         \textit{si ragiona sulla macchina:} & mov eax, [rbp-0x8] & copia \code{b} nel registro eax \\
         \textit{registri e indirizzi} & add eax, edx & eax = eax + edx \\
         & mov [rbp-0x4], eax & copia il risultato in \code{c} \\
       \end{tabular}};
    \node[liv, draw=accent!60!black, fill=accent!6!white] (mac) at ([yshift=-0.35cm]asm.south west)
      {\begin{tabular}{@{}p{3.5cm}>{\FiraCodeNL\small}l@{}}
         \textbf{linguaggio macchina} & 10001011 01010101 11110100 \\
         \textit{bit, cioè numeri} & 10001011 01000101 11111000 \\
         & 00000001 11010000 \\
         & 10001001 01000101 11111100 \\
       \end{tabular}};
    \draw[wtarrowacc, line width=1.2pt] ([xshift=-0.6cm, yshift=-0.2cm]alto.north east) -- node[wtlab, left=4pt, fill=none, align=right] {traduzione\\(compilatore)} ([xshift=-0.6cm, yshift=0.2cm]mac.south east);
  \end{tikzpicture}
  \caption{La stessa istruzione a tre livelli di astrazione. Un'istruzione
           di alto livello diventa più istruzioni macchina; ogni istruzione
           assembly corrisponde a un'istruzione macchina.}
  \label{fig:livelli}
\end{figure}

Ogni livello richiede di conoscere cose diverse:

\begin{itemize}
  \item per il \textbf{linguaggio macchina}, come sono rappresentate le
        informazioni in bit;
  \item per l'\textbf{assembly}, le caratteristiche della macchina:
        registri, insieme di istruzioni, dimensione dei dati;
  \item per un \textbf{linguaggio di alto livello}, solo il linguaggio: ci
        si astrae dai dettagli della macchina.
\end{itemize}

Il prezzo da pagare è che un programma di alto livello non può essere
eseguito direttamente: va \textbf{tradotto} in linguaggio macchina. È il
compito del \term{compilatore}, l'argomento del
capitolo~\ref{chap:dal-problema}.

\subsection{Algoritmo e programma}

Due parole che useremo continuamente vanno distinte con cura.

\dfn{Algoritmo}{
  Una sequenza di passi elementari che risolve un problema. La sequenza
  deve essere \textbf{finita}, \textbf{non ambigua},
  \textbf{deterministica} e \textbf{generale}.}

\dfn{Programma}{
  La formalizzazione di un algoritmo in un linguaggio di programmazione.}

Le quattro proprietà dell'algoritmo non sono pignolerie: ciascuna esclude un
modo concreto in cui un procedimento può fallire. Prendiamo una ricetta di
cucina, che è quasi un algoritmo:

\begin{itemize}
  \item \textbf{finita}: deve terminare. “Mescola finché non è pronto” è
        accettabile solo se prima o poi “è pronto”;
  \item \textbf{non ambigua}: ogni passo deve avere un solo significato.
        “Sale quanto basta” va bene per un cuoco, non per un computer, che
        non ha buon senso: bisogna dire quanti grammi;
  \item \textbf{deterministica}: in ogni momento il passo successivo è
        determinato; con gli stessi dati si ottiene sempre lo stesso
        risultato;
  \item \textbf{generale}: risolve una \textit{classe} di problemi, non un
        solo caso. Una ricetta “per 4 persone” è meno generale di una “per
        $n$ persone”; un algoritmo che calcola l'area di un cerchio di
        raggio 2 è inutile, ne serve uno che funzioni per qualunque raggio.
\end{itemize}

\section{La memoria}

\subsection{Celle e indirizzi}

La memoria centrale è il contenitore dei dati su cui lavora il processore.
Si può immaginare come una lunghissima fila di caselle consecutive, le
\term{locazioni} (o \textit{celle}) di memoria. Ogni cella:

\begin{itemize}
  \item è l'unità minima di memorizzazione: si legge e si scrive una cella
        intera, non un pezzo;
  \item contiene tipicamente \textbf{8 bit}, cioè un \term{byte};
  \item è identificata da un numero, il suo \term{indirizzo}, come le case
        di una strada dal numero civico.
\end{itemize}

\begin{affianco}[0.5][c]
\begin{immagine}
\begin{tikzpicture}
  \memoria{m}{(0,0)}{1000/00000101/, 1001/01000001/, 1002/11111111/, 1003/00000000/, 1004/\dots/}
  \node[mnota, anchor=west] at ([xshift=12pt]m-1.east) {il numero 5};
  \node[mnota, anchor=west] at ([xshift=12pt]m-2.east) {il numero 65, o la lettera A};
  \node[mnota, anchor=west] at ([xshift=12pt]m-3.east) {il numero 255};
  \node[mnota, anchor=south] at ([yshift=3pt]m-1.north) {\textbf{contenuto} (8 bit)};
  \node[mnota, anchor=south east] at ([yshift=3pt, xshift=-4pt]m-1.north west) {\textbf{indirizzo}};
\end{tikzpicture}
\end{immagine}
\risultato
Un \term{bit} (\textit{binary digit}, cifra binaria) è la più piccola unità
di informazione: vale 0 oppure 1. Elettronicamente corrisponde a qualcosa
che ha due stati, come un interruttore acceso o spento. Con un bit solo si
distinguono due situazioni; mettendone insieme 8 se ne distinguono molte di
più: vediamo quante, e come si usano per rappresentare i numeri.
\end{affianco}

\subsection{La notazione binaria}

Siamo abituati alla notazione \textbf{decimale}: dieci cifre, da 0 a 9, e
un sistema \textbf{posizionale}, in cui il valore di una cifra dipende dalla
posizione. In \code{352} il 3 vale trecento perché sta nella posizione delle
centinaia:
\[
  352 = 3 \cdot 100 + 5 \cdot 10 + 2 \cdot 1 = 3 \cdot 10^2 + 5 \cdot 10^1 + 2 \cdot 10^0 .
\]
La notazione \textbf{binaria} funziona allo stesso modo, ma con due sole
cifre, 0 e 1, e con le potenze di \textbf{2} al posto delle potenze di 10
(figura~\ref{fig:binario}).

\begin{figure}[H]
  \centering
  \begin{tikzpicture}[font=\FiraCodeNL\small]
    \foreach \p [count=\i from 0] in {128,64,32,16,8,4,2,1} {
      \node[mcella, minimum width=1.25cm, minimum height=0.7cm, font=\FiraCodeNL\normalsize]
        (b\i) at (\i*1.25,0) {};
      \node[font=\FiraCodeNL\scriptsize, text=neutral!40!black] at (\i*1.25,0.62) {\p};
      \pgfmathtruncatemacro{\e}{7-\i}
      \node[font=\Lato\tiny, text=neutral!50!black] at (\i*1.25,1.0) {$2^{\e}$};
    }
    \foreach \val [count=\i from 0] in {0,0,1,0,1,1,0,1} {
      \ifnum\val=1
        \node[font=\FiraCodeNL\normalsize\bfseries, text=accent!80!black] at (b\i) {1};
      \else
        \node[font=\FiraCodeNL\normalsize, text=neutral!55!black] at (b\i) {0};
      \fi
    }
    \node[anchor=west, font=\Lato\small] at (10.1,0)
      {$= 32 + 8 + 4 + 1 = \mathbf{45}$};
    \node[font=\Lato\tiny, anchor=east, text=neutral!50!black] at (-0.75,1.0) {potenza};
    \node[font=\Lato\tiny, anchor=east, text=neutral!50!black] at (-0.75,0.62) {peso};
    \node[font=\Lato\tiny, anchor=east, text=neutral!50!black] at (-0.75,0) {bit};
  \end{tikzpicture}
  \caption{Il byte \code{00101101}: ogni bit a 1 aggiunge il peso della sua
           posizione.}
  \label{fig:binario}
\end{figure}

Per passare dal binario al decimale, quindi, si sommano i pesi delle
posizioni in cui c'è un 1. Il passaggio inverso si fa con le
\textbf{divisioni successive per 2}: si divide il numero per 2, si annota il
resto (0 o 1), si ripete sul quoziente fino ad arrivare a 0; i resti, letti
\textbf{dal basso verso l'alto}, sono le cifre binarie.

\begin{affianco}[0.36][c]
\begin{immagine}
\begin{tikzpicture}[font=\FiraCodeNL\small]
  \foreach \num/\quoz/\resto [count=\i from 0] in {13/6/1, 6/3/0, 3/1/1, 1/0/1} {
    \node[anchor=east] at (0,-\i*0.55) {\num\ : 2 = \quoz};
    \node[anchor=west, text=accent!80!black] at (0.3,-\i*0.55) {resto \textbf{\resto}};
  }
  \draw[wtarrowacc] (2.75,-1.75) -- node[right, font=\Lato\scriptsize, text=accent!70!black, align=left] {leggere\\dal basso} (2.75,0.15);
\end{tikzpicture}
\end{immagine}
\risultato
Per esempio, 13 in binario è \code{1101}: lo si verifica sommando i pesi,
$8 + 4 + 1 = 13$. Perché funziona? Ogni divisione per 2 “stacca” l'ultima
cifra binaria del numero, esattamente come dividere per 10 un numero
decimale ne stacca l'ultima cifra (il resto di $352 : 10$ è 2, l'ultima
cifra). Per questo il primo resto trovato è la cifra più a destra, e i resti
vanno letti al contrario.
\end{affianco}

\subsection{Quanti numeri in un byte}

Con $n$ posizioni e 2 simboli possibili per ciascuna si formano $2^n$
combinazioni diverse: ogni bit in più \textbf{raddoppia} le possibilità.
Un byte ne ha $2^8 = 256$, quindi rappresenta i numeri naturali da 0 a 255:

\begin{center}
\renewcommand{\arraystretch}{1.15}
\begin{tabular}{@{}>{\FiraCodeNL}c c@{\hspace{1.4cm}}>{\FiraCodeNL}c c@{}}
  \toprule
  \normalfont\textbf{bit} & \textbf{valore} & \normalfont\textbf{bit} & \textbf{valore} \\
  \midrule
  00000000 & 0 & 00000100 & 4 \\
  00000001 & 1 & 00000101 & 5 \\
  00000010 & 2 & \dots    & \dots \\
  00000011 & 3 & 11111111 & 255 \\
  \bottomrule
\end{tabular}
\end{center}

\subsection{I numeri negativi}

E i numeri negativi? Bisogna usare una parte dell'informazione per il
\textbf{segno}. L'idea più semplice, vista a lezione, è la
\term{rappresentazione in modulo e segno}: il primo bit indica il segno (0
positivo, 1 negativo) e gli altri 7 il valore assoluto. Con 7 bit si
rappresentano $2^7 = 128$ valori assoluti, da 0 a 127, quindi l'intervallo
è $[-127,\ 127]$.

\warning{Una svista nelle slide}{
  Le slide indicano l'intervallo $[-127 .. 128]$, ma in modulo e segno il
  massimo è $+127$ (\code{01111111}), non $128$: con 7 bit il valore
  assoluto più grande è $2^7 - 1 = 127$. L'intervallo è simmetrico,
  $[-127,\ 127]$, e lo zero compare due volte: \code{00000000} ($+0$) e
  \code{10000000} ($-0$).}

Proprio quel “doppio zero”, e il fatto che le somme con numeri negativi
richiederebbero circuiti speciali, fanno sì che i computer reali usino
un'altra rappresentazione, il \term{complemento a due}, in cui con un byte
si rappresentano i numeri da $-128$ a $+127$. L'idea si vede bene con 3
bit (figura~\ref{fig:ruota}): i numeri sono disposti su una ruota, come le
cifre del contachilometri di un'auto; partendo da \code{000} e andando
\textbf{indietro} di un passo si arriva a \code{111}, che quindi vale $-1$.

\begin{figure}[H]
  \centering
  \begin{tikzpicture}[font=\FiraCodeNL\scriptsize]
    \draw[neutral!40!white, line width=6pt] (0,0) circle (1.75);
    \foreach \bits/\val/\pos in {000/0/90, 001/1/45, 010/2/0, 011/3/-45,
                                  100/-4/-90, 101/-3/-135, 110/-2/180, 111/-1/135} {
      \node[mcella, minimum width=0.85cm] at (\pos:1.75) {\bits};
      \node[font=\Lato\small\bfseries, text=primary!60!black] at (\pos:2.5) {\val};
    }
    \draw[wtarrowacc, line width=1pt] (100:1.25) arc[start angle=100, end angle=125, radius=1.25];
    \node[font=\Lato\tiny, text=accent!70!black, align=center] at (0,0.3) {un passo indietro\\da 0 è $-1$};
    \node[font=\Lato\tiny, text=neutral!35!black, align=center] at (0,-0.45) {il primo bit\\è 1 per i negativi};
  \end{tikzpicture}
  \caption{Il complemento a due su 3 bit: 8 combinazioni, da $-4$ a $+3$,
           un solo zero. Su 8 bit si va da $-128$ a $+127$.}
  \label{fig:ruota}
\end{figure}

Il complemento a due ha conseguenze molto concrete in C (per esempio su
che cosa succede quando un numero “trabocca” oltre il massimo): le
vedremo nel capitolo~\ref{chap:tipi}.

\section{I tipi di dato}

\subsection{Gli stessi bit, significati diversi}

Il byte \code{01000001} contiene semplicemente otto bit. È il numero 65?
È la lettera \textit{A}? Dipende da come lo si \textbf{interpreta}: la
stessa quantità di memoria può rappresentare un numero naturale, un intero
con segno, un numero reale o un carattere. La memoria non sa che cosa
contiene; lo sa il programma.

\begin{figure}[H]
  \centering
  \begin{tikzpicture}[font=\Lato\scriptsize]
    \node[mcella, minimum width=2.6cm, minimum height=0.7cm, font=\FiraCodeNL\normalsize] (b) at (0,0) {01000001};
    \foreach \ang/\testo in {150/{come naturale:\\\textbf{65}}, 30/{come carattere:\\\textbf{A} (codice ASCII 65)},
                        -150/{come intero con segno:\\\textbf{+65}}, -30/{come parte di un \code{float}:\\dipende dagli altri 3 byte}} {
      \node[wtnodeplain, text width=3.3cm] (n) at (\ang:3.0) {\testo};
      \draw[wtarrow] (b) -- (n);
    }
  \end{tikzpicture}
  \caption{Lo stesso byte, letto con occhiali diversi.}
  \label{fig:occhiali}
\end{figure}

\subsection{Il tipo come astrazione}

I linguaggi di alto livello risolvono il problema con i \term{tipi di
dato}: si scrive il programma in termini di \textit{che cosa} sono i dati,
senza preoccuparsi di come e dove saranno davvero memorizzati. In C si può
scrivere

\begin{lstlisting}[style=c]
int variabile = 5;
\end{lstlisting}

e non ci si chiede in quali celle finirà il 5, né come sarà scritto in bit:
si usa \code{int} nel suo significato astratto di “contenitore di numeri
interi”. Eppure un \code{int} non sta in una cella sola: ne occupa
\textbf{quattro consecutive} (32 bit, quindi valori ben oltre 255). È il
tipo a stabilire che quelle quattro celle vanno lette insieme, come un
unico numero (figura~\ref{fig:int4}).

\begin{figure}[H]
  \centering
  \begin{tikzpicture}[font=\FiraCodeNL\scriptsize]
    \noindici
    \celle{c}{(0,0)}{00000101, 00000000, 00000000, 00000000}
    \foreach \i/\a in {0/2000, 1/2001, 2/2002, 3/2003}
      \node[mindir, anchor=north] at ([yshift=-1pt]c-\i.south) {\a};
    \draw[wtbrace] ([yshift=0.15cm]c-3.north east) -- ([yshift=0.15cm]c-0.north west);
    \node[font=\Lato\scriptsize, anchor=south, align=center] at ([yshift=0.45cm]$(c-0.north)!0.5!(c-3.north)$)
      {\FiraCodeNL variabile\Lato\ (\FiraCodeNL int\Lato): \textbf{4 byte} letti insieme come il numero 5};
  \end{tikzpicture}
  \caption{Un \code{int} occupa 4 celle consecutive (sui PC il byte meno
           significativo sta all'indirizzo più basso).}
  \label{fig:int4}
\end{figure}

Ogni linguaggio definisce i propri tipi. In C, i principali sono:

\begin{center}
\renewcommand{\arraystretch}{1.3}
\begin{tabular}{@{}>{\FiraCodeNL}l l l@{}}
  \toprule
  \normalfont\textbf{tipo} & \textbf{che cosa rappresenta} & \textbf{dimensione tipica} \\
  \midrule
  char   & un carattere (codice ASCII)     & 8 bit (1 byte) \\
  int    & un numero intero                & 32 bit (4 byte) \\
  float  & un numero reale                 & 32 bit (4 byte) \\
  double & un numero reale, più preciso    & 64 bit (8 byte) \\
  \bottomrule
\end{tabular}
\end{center}

Le dimensioni sono quelle dei computer di oggi; lo standard del C fissa solo
dei minimi, come vedremo nel capitolo~\ref{chap:tipi}.

\section{La gerarchia di memoria}

La memoria centrale può contenere molti programmi e dati, ma ha due difetti:

\begin{itemize}
  \item è \textbf{lenta} rispetto al processore: leggere una cella richiede
        il tempo di un centinaio di operazioni della CPU;
  \item è \textbf{volatile}: quando si spegne il computer, il contenuto si
        perde.
\end{itemize}

Contro la lentezza si costruisce una \term{gerarchia}: piccole memorie
velocissime vicino al processore, memorie via via più grandi e lente man
mano che ci si allontana. Le memorie veloci costano molto per ogni bit,
quindi sono piccole. Contro la volatilità ci sono le \term{memorie di massa}
(dischi, SSD, chiavette), che conservano i dati anche a computer spento
(figura~\ref{fig:gerarchia}).

\begin{figure}[H]
  \centering
  \begin{tikzpicture}[font=\Lato\scriptsize]
    \def\h{0.78}
    \foreach \n/\nome/\vel/\dim/\col [count=\i from 0] in {
        0/registri (dentro la CPU)/circa 1 ns/centinaia di byte/accent,
        1/cache/circa 10 ns/kilobyte\,--\,megabyte/accent,
        2/memoria centrale (RAM)/circa 100 ns/gigabyte/primary,
        3/memoria di massa (disco{,} SSD)/da 0{,}1 a 10 ms/terabyte/green-gradient-6,
        4/archivio (nastri{,} cloud)/secondi/molti terabyte/neutral} {
      \pgfmathsetmacro{\w}{2.9+\i*1.3}
      \fill[\col!14!white, draw=\col!65!black, line width=0.7pt]
        (-\w/2,-\i*\h) -- (\w/2,-\i*\h) -- (\w/2+0.65,-\i*\h-\h) -- (-\w/2-0.65,-\i*\h-\h) -- cycle;
      \node[font=\Lato\scriptsize\bfseries] at (0,-\i*\h-\h/2) {\nome};
      \node[anchor=west] at (5.2,-\i*\h-\h/2) {\vel};
      \node[anchor=west] at (7.6,-\i*\h-\h/2) {\dim};
    }
    \node[anchor=west, font=\Lato\scriptsize\bfseries] at (5.2,0.3) {tempo di accesso};
    \node[anchor=west, font=\Lato\scriptsize\bfseries] at (7.6,0.3) {dimensione};
    \draw[wtarrow, <->] (-5.3,-0.1) -- node[wtlab, rotate=90, fill=white] {più veloci, più costose} (-5.3,-3.8);
    \draw[neutral!50!white] (10.2,-0.25) -- (10.2,-2.25);
    \node[font=\Lato\tiny, text=neutral!35!black, anchor=west] at (10.3,-1.25) {volatili};
    \draw[neutral!50!white] (10.2,-2.45) -- (10.2,-3.85);
    \node[font=\Lato\tiny, text=neutral!35!black, anchor=west] at (10.3,-3.15) {permanenti};
  \end{tikzpicture}
  \caption{La gerarchia di memoria (1 ns è un miliardesimo di secondo). In
           alto memorie piccole, veloci, costose e volatili; in basso grandi,
           lente, economiche e permanenti.}
  \label{fig:gerarchia}
\end{figure}

Un'immagine aiuta: la CPU è un impiegato alla scrivania. I registri sono ciò
che tiene in mano; la cache, i fogli sulla scrivania; la RAM, lo scaffale
nella stanza; il disco, l'archivio in cantina. L'impiegato lavora solo su
ciò che ha in mano, e se un foglio non è sulla scrivania deve alzarsi a
prenderlo: per questo i dati usati spesso vengono tenuti vicini.

\section{Il software}

\subsection{Gli strati del software}

Fra l'hardware e i programmi che usiamo c'è un sistema software a strati
(figura~\ref{fig:strati}). Ogni strato usa i servizi di quello sotto e ne
offre di più comodi a quello sopra: è ancora l'idea di astrazione.

\begin{figure}[H]
  \centering
  \begin{tikzpicture}[font=\Lato\scriptsize,
      strato/.style={draw=#1!65!black, fill=#1!12!white, line width=0.7pt,
        minimum width=9.5cm, minimum height=0.68cm, anchor=north},
      es/.style={anchor=west, font=\Lato\scriptsize, text=neutral!25!black}]
    \node[strato=accent] (app) at (0,0) {\textbf{applicazioni}};
    \node[strato=primary] (lib) at ([yshift=-2pt]app.south) {\textbf{librerie} (per esempio la libreria standard del C)};
    \node[strato=green-gradient-6] (int) at ([yshift=-2pt]lib.south) {\textbf{interfaccia}: a riga di comando (la shell) o grafica};
    \node[strato=orange-gradient-6] (so) at ([yshift=-2pt]int.south) {\textbf{sistema operativo}: gestione di processi, memoria, periferiche, file};
    \node[strato=red-gradient-6] (ker) at ([yshift=-2pt]so.south) {\textbf{nucleo} (\textit{kernel}) del sistema operativo};
    \node[strato=neutral] (hw) at ([yshift=-2pt]ker.south) {\textbf{hardware}: CPU, memoria, dispositivi};
    \node[es] at ([xshift=0.4cm]app.east) {browser, VS Code, i nostri programmi};
    \node[es] at ([xshift=0.4cm]lib.east) {\code{printf}, \code{scanf}\dots};
    \node[es] at ([xshift=0.4cm]int.east) {bash, il desktop};
    \node[es] at ([xshift=0.4cm]so.east) {Linux, Windows, macOS};
  \end{tikzpicture}
  \caption{Fra l'hardware e l'ambiente in cui si sviluppano i programmi
           c'è un sistema software stratificato.}
  \label{fig:strati}
\end{figure}

\subsection{L'avvio del computer}

Il primo programma eseguito all'accensione è il \term{sistema operativo}:
è lui che, quando chiediamo di eseguire un programma, lo carica in memoria
e dice al processore dove si trova la sua prima istruzione. Ma chi dice al
processore dove si trova la prima istruzione del sistema operativo? È un
problema dell'uovo e della gallina, risolto così
(figura~\ref{fig:avvio}): all'accensione il computer esegue un piccolo
programma che sta in una memoria permanente della scheda madre (il
\textit{firmware}); questo carica in memoria il contenuto di un blocco
predefinito di un disco, il \term{Master Boot Record} (MBR). Il codice del
MBR, chiamato \term{bootloader}, è una sequenza di byte che si assume essere
un programma in linguaggio macchina: il suo compito è caricare il sistema
operativo. Il MBR è configurabile: se sul disco ci sono più sistemi
operativi, si può scegliere quale avviare.

\begin{figure}[H]
  \centering
  \begin{tikzpicture}[font=\Lato\scriptsize,
      passo/.style={wtnode, font=\Lato\scriptsize, text width=2.8cm, minimum height=1.3cm}]
    \node[passo, draw=neutral!60!white, fill=neutral!8!white, text width=2.4cm] (a) at (0,0) {\textbf{accensione}};
    \node[passo] (b) at (3.6,0) {\textbf{firmware}\\(BIOS o UEFI)\\nella scheda madre};
    \node[passo] (c) at (7.4,0) {\textbf{bootloader}\\letto dal MBR del disco};
    \node[passo, draw=accent!60!black, fill=accent!8!white] (d) at (11.2,0) {\textbf{sistema operativo}\\caricato in memoria};
    \node[passo, draw=green-gradient-6, fill=green-gradient-1!45!white, text width=2.4cm] (e) at (14.8,0) {i \textbf{programmi} dell'utente};
    \foreach \x/\y in {a/b, b/c, c/d, d/e} \draw[wtarrow] (\x) -- (\y);
  \end{tikzpicture}
  \caption{L'avvio: ogni programma carica il successivo, più grande e più
           capace.}
  \label{fig:avvio}
\end{figure}

\info{MBR e UEFI}{
  Il Master Boot Record è il meccanismo storico dei PC. I computer recenti
  usano il firmware \textit{UEFI}, che invece di un singolo blocco legge il
  bootloader da una piccola partizione dedicata del disco. L'idea non
  cambia: un programma minuscolo, che sta in un posto noto in anticipo,
  avvia un programma più grande.}

\subsection{Che cosa fa il sistema operativo}

Il sistema operativo \textbf{alza il livello di astrazione} della macchina:
è più facile da programmare dell'hardware sottostante e semplifica la
scrittura di ogni altro programma. Le sue funzioni principali:

\begin{itemize}
  \item \textbf{gestione dei processi}, cioè dei programmi in esecuzione:
        li crea e li termina, li sospende e li riattiva, li fa comunicare.
        Il processore esegue un'istruzione alla volta, ma il sistema
        operativo lo assegna a turno ai vari programmi per pochi millesimi
        di secondo ciascuno (la \textit{schedulazione}): così sembrano
        procedere tutti insieme;
  \item \textbf{gestione della memoria centrale}: assegna la memoria ai
        processi, li carica e li rimuove, e la \textbf{protegge}, cioè
        impedisce a un programma di scrivere nella memoria di un altro;
  \item \textbf{gestione delle periferiche}: le configura e le pilota, con
        il \textit{buffering} (tiene i dati in memoria finché non vengono
        elaborati) e il \textit{caching} (tiene vicini i dati usati spesso);
  \item \textbf{gestione del file system}: nasconde all'utente come sono
        organizzati i dati sul disco, offrendo il concetto di \textit{file}
        e di \textit{cartella} (l'albero di cartelle del
        capitolo~\ref{chap:ambiente}), con le operazioni di creazione,
        lettura, scrittura, copia, ricerca e i diritti di accesso;
  \item \textbf{interfaccia utente}: a riga di comando (l'interprete dei
        comandi, la shell che usiamo nel terminale) o grafica;
  \item \textbf{rilevamento degli errori}: divisione per zero, carta
        finita, rete caduta\dots; per ogni errore serve un'azione di
        recupero appropriata.
\end{itemize}

L'ultima voce ci riguarderà da vicino. Un programma C che divide un intero
per zero chiede al processore un'operazione impossibile; il processore
avvisa il sistema operativo, che \textbf{termina il programma}:

\begin{affianco}[0.45][c]
\begin{lstlisting}[style=c, name=div0.c]
#include <stdio.h>

int main(void)
{
    int a = 7, b = 0;
    printf("%d\n", a / b);
    return 0;
}
\end{lstlisting}
\risultato
\begin{console}
$ gcc -o div0 div0.c
$ ./div0
Floating point exception
\end{console}
\wtnota{Il messaggio parla di “virgola mobile” anche se la divisione è
fra interi: è un nome storico.}
\end{affianco}

\subsection{I servizi per i programmi}

Il sistema operativo offre servizi ai programmi: caricarli in memoria ed
eseguirli, compiere le operazioni di I/O (un programma normale non può
pilotare direttamente i dispositivi), leggere e scrivere i file, far
comunicare processi sullo stesso computer o su computer diversi (con memoria
condivisa o scambio di messaggi), rilevare gli errori. Ci sono poi servizi
per gestire il sistema nel suo insieme: allocazione delle risorse fra utenti
e processi, contabilità del loro uso, protezione e sicurezza.

Anche il più semplice dei nostri programmi passa da questi servizi. Quando
il programma del capitolo precedente esegue \code{printf}, il testo non va
direttamente sullo schermo (figura~\ref{fig:printf-so}).

\begin{figure}[H]
  \centering
  \begin{tikzpicture}[font=\Lato\scriptsize,
      p/.style={wtnode, font=\Lato\scriptsize, text width=3.5cm, minimum height=1.15cm}]
    \node[p] (a) at (0,0) {il nostro programma chiama \code{printf}};
    \node[p] (b) at (4.5,0) {la \textbf{libreria} del C prepara il testo};
    \node[p, draw=orange-gradient-6!80!black, fill=orange-gradient-1!30!white] (c) at (9.0,0) {chiede al \textbf{sistema operativo} di scriverlo};
    \node[p, draw=green-gradient-6, fill=green-gradient-1!45!white] (d) at (13.5,0) {il testo compare nel \textbf{terminale}};
    \foreach \x/\y in {a/b, b/c, c/d} \draw[wtarrow] (\x) -- (\y);
  \end{tikzpicture}
  \caption{Il percorso di una \code{printf}: dal programma alla libreria,
           dalla libreria al sistema operativo, e infine al dispositivo.}
  \label{fig:printf-so}
\end{figure}

\subsection{Installare e scrivere un sistema operativo}

Lo stesso sistema operativo gira su macchine diverse, quindi va
\textbf{configurato} in fase di installazione: processore, memoria,
suddivisione del disco (partizioni), dispositivi presenti, opzioni come
l'algoritmo di schedulazione o il numero massimo di processi.

I primi sistemi operativi erano scritti in assembly; oggi si scrivono in
linguaggi di alto livello, perché il codice si scrive più in fretta, è più
compatto, più facile da capire e da correggere. Un sistema operativo scritto
in un linguaggio di alto livello è molto più flessibile e manutenibile. Non
a caso, il C nacque proprio per scrivere un sistema operativo, Unix, come
racconta il capitolo~\ref{chap:strutturata}.

\insintesi{
  Un calcolatore è fatto di CPU, memoria e dispositivi di I/O, collegati
  dal bus. La memoria è una fila di celle di un byte, ognuna con un
  indirizzo, che contengono sia i dati sia le istruzioni dei programmi. La
  CPU sa eseguire solo istruzioni elementari del proprio
  \textbf{linguaggio macchina}, una alla volta, guidata dal program counter;
  ripetizioni e scelte si ottengono con i salti. I linguaggi di alto
  livello come il C nascondono questi dettagli (astrazione), ma vanno
  \textbf{tradotti}. I bit di per sé non hanno significato: sono i
  \textbf{tipi} a stabilire come interpretarli. Il sistema operativo, il
  primo programma caricato all'avvio, gestisce processi, memoria, periferiche
  e file, e offre i suoi servizi a tutti gli altri programmi.}

\section{Riferimenti}

\begin{itemize}
  \item \textbf{Deitel \& Deitel}, \textit{Il linguaggio C --- Fondamenti e
        tecniche di programmazione}, Pearson: capitolo~1 (computer,
        hardware e software, linguaggi macchina, assembly e di alto
        livello).
  \item \textbf{Treccani}, voce \textit{programmazione}. \\
        \urlx{https://www.treccani.it/vocabolario/programmazione/}
\end{itemize}


  \chapter{Dal problema al programma}
\label{chap:dal-problema}

Il capitolo precedente ha mostrato che cosa sa fare la macchina: eseguire,
velocissima, istruzioni elementari scritte in linguaggio macchina. Questo
capitolo percorre la strada nell'altro senso, partendo da un
\textbf{problema} da risolvere: come si descrive la soluzione (l'algoritmo),
come la si scrive in un linguaggio di alto livello (il programma) e come il
programma diventa qualcosa che la CPU può eseguire.

\section{Programmare}

La programmazione è l'attività che porta a sviluppare un software. Non ogni
procedimento, però, può diventare un programma: deve poter essere
\textbf{descritto in modo non ambiguo}, deve \textbf{raggiungere l'obiettivo
in un tempo finito} e deve essere \textbf{deterministico}. Un procedimento
con queste caratteristiche si chiama \term{algoritmo}: una ricetta di cucina
ben scritta, o il metodo per risolvere un'equazione di secondo grado.

Programmare significa descrivere al calcolatore \textit{come} risolvere un
problema. Il processore mette a disposizione molte operazioni di base
(somma, sottrazione, operazioni logiche, spostamenti di valori, letture e
scritture verso le periferiche), ma, come abbiamo visto, conviene esprimersi
in un linguaggio di alto livello e lasciare a un programma il compito di
tradurre.

\section{I linguaggi di alto livello}

Un linguaggio di alto livello (\textit{HLL}, \textit{High Level Language})
ha tre caratteristiche:

\begin{itemize}
  \item un \textbf{alto livello di astrazione}: le istruzioni hanno un
        significato più complesso di quelle della macchina, e non serve
        sapere come il processore le realizza (\code{A + B} somma due
        valori, punto);
  \item viene \textbf{tradotto} in linguaggio macchina (o in assembly) da
        un programma traduttore, ed è quindi indipendente dal processore e
        dal sistema operativo;
  \item è più \textbf{simile al linguaggio umano}, quindi più facile da
        scrivere e da ricordare (in molti linguaggi \code{print X} mostra il
        valore di \code{X}).
\end{itemize}

Comunque avvenga la traduzione, il programmatore scrive le istruzioni in un
file di testo: il \term{codice sorgente} (o \textit{programma sorgente}, o
semplicemente \textit{programma}). Per tradurlo ci sono due strategie,
l'\textbf{interprete} e il \textbf{compilatore}.

\subsection{Interprete e compilatore}

Un \term{interprete} traduce le istruzioni del codice sorgente \textbf{una
alla volta} in linguaggio macchina, e ciascuna viene eseguita subito dalla
CPU, prima di passare alla successiva. Un \term{compilatore}, invece,
traduce \textbf{tutto} il codice sorgente in linguaggio macchina e lo salva
in un file, il \term{programma eseguibile}; l'esecuzione avviene dopo, ed è
un'operazione separata (figura~\ref{fig:interprete}).

\begin{figure}[H]
  \centering
  \begin{tikzpicture}[font=\Lato\scriptsize,
      src/.style={draw=neutral!55!white, fill=neutral!8!white, minimum width=2.1cm,
        minimum height=0.42cm, font=\Lato\tiny},
      sw/.style={fill=primary!65!white, text=white, minimum width=2.1cm,
        minimum height=0.85cm, font=\Lato\scriptsize\bfseries, align=center},
      ese/.style={draw=neutral!55!white, fill=neutral!14!white, minimum width=2.2cm,
        minimum height=0.85cm, font=\Lato\tiny, align=center}]
    % interprete
    \node[font=\Lato\small\bfseries, anchor=west] at (-1.2,1.1) {Interprete};
    \draw[neutral!40!white, fill=neutral!3!white] (-1.2,0.7) rectangle (1.2,-2.6);
    \node[font=\Lato\tiny\bfseries] at (0,0.5) {FILE SORGENTE};
    \node[src] (i1) at (0,0) {istruzione 1};
    \node[src] (i2) at (0,-0.6) {istruzione 2};
    \node at (0,-1.25) {\vdots};
    \node[src] (in) at (0,-1.9) {istruzione N};
    \node[font=\FiraCodeNL\tiny] at (0,-2.4) {programma.py};
    \node[sw] (s1) at (3.5,0) {interprete};
    \node[sw] (s2) at (3.5,-1.2) {interprete};
    \node[ese] (e1) at (6.3,0) {esecuzione\\istruzione 1};
    \node[ese] (e2) at (6.3,-1.2) {esecuzione\\istruzione 2};
    \node at (3.5,-2.0) {\vdots};
    \draw[wtarrow] (i1.east) -- (s1.west);
    \draw[wtarrow] (i2.east) -- ++(0.6,0) |- (s2.west);
    \draw[wtarrow] (s1) -- (e1);
    \draw[wtarrow] (s2) -- (e2);
    \node[mnota, text width=4.6cm, anchor=north west] at (-1.2,-2.8)
      {la traduzione avviene \textbf{a ogni esecuzione}, un'istruzione alla volta};
    % compilatore
    \begin{scope}[xshift=9.0cm]
      \node[font=\Lato\small\bfseries, anchor=west] at (-1.2,1.1) {Compilatore};
      \draw[neutral!40!white, fill=neutral!3!white] (-1.2,0.7) rectangle (1.2,-2.6);
      \node[font=\Lato\tiny\bfseries] at (0,0.5) {FILE SORGENTE};
      \node[src] at (0,0) {istruzione 1};
      \node[src] at (0,-0.6) {istruzione 2};
      \node at (0,-1.25) {\vdots};
      \node[src] at (0,-1.9) {istruzione N};
      \node[font=\FiraCodeNL\tiny] at (0,-2.4) {programma.c};
      \node[sw] (c) at (2.95,-0.95) {compilatore};
      \node[ese, fill=accent!12!white, draw=accent!50!black, minimum width=1.8cm] (x) at (5.3,-0.95) {file\\eseguibile};
      \node[ese, minimum width=1.8cm] (r) at (7.6,-0.95) {esecuzione};
      \draw[wtarrow] (1.2,-0.95) -- (c);
      \draw[wtarrow] (c) -- (x);
      \draw[wtarrow] (x) -- (r);
      \node[mnota, text width=6.0cm, anchor=north west] at (-1.2,-2.8)
        {la traduzione avviene \textbf{una volta sola}, prima; poi il file
         eseguibile si esegue quante volte si vuole};
    \end{scope}
  \end{tikzpicture}
  \caption{Interprete (a sinistra) e compilatore (a destra).}
  \label{fig:interprete}
\end{figure}

Il C è un linguaggio \textbf{compilato}; Python e JavaScript, per esempio,
sono di solito interpretati. Le due strategie hanno pregi e difetti
opposti:

\begin{center}
\renewcommand{\arraystretch}{1.45}
\begin{tabular}{@{}>{\raggedright\arraybackslash}p{2.8cm}>{\raggedright\arraybackslash}p{6.6cm}>{\raggedright\arraybackslash}p{6.6cm}@{}}
  \toprule
  & \textbf{Programma interpretato} & \textbf{Programma compilato} \\
  \midrule
  \textbf{velocità} & più lento: a ogni esecuzione l'interprete deve ritradurre
    le istruzioni & più veloce: è già tradotto; e il compilatore sa
    \textit{ottimizzare} il codice, per renderlo più veloce o più piccolo \\
  \textbf{chi deve avere che cosa} & l'utente deve avere l'interprete e saperlo
    usare (si può nascondere, ma c'è) & all'utente basta l'eseguibile (un
    doppio clic); solo il programmatore ha bisogno del compilatore \\
  \textbf{codice sorgente} & serve all'utente, quindi è visibile e
    modificabile da chiunque & non serve per eseguire: si può non
    distribuirlo, proteggendo il lavoro (risalire al sorgente è possibile, ma
    difficile) \\
  \bottomrule
\end{tabular}
\end{center}

\subsection{Una via di mezzo: il bytecode}

Alcuni linguaggi, come Java e Python, compilano il sorgente in un file
intermedio, non eseguibile direttamente, scritto in una sorta di assembly
generico chiamato \term{bytecode}, che si traduce molto rapidamente nel
codice macchina di un processore reale. Su ogni computer si installa un
interprete del bytecode (per Java, la \textit{Java Virtual Machine}): così
lo stesso file gira su macchine con linguaggi macchina diversi
(figura~\ref{fig:bytecode}).

\begin{figure}[H]
  \centering
  \begin{tikzpicture}[font=\Lato\scriptsize,
      b/.style={wtnode, font=\Lato\scriptsize, minimum height=0.9cm, text width=2.6cm}]
    \node[b] (s) at (0,0) {sorgente\\{\FiraCodeNL Programma.java}};
    \node[b, draw=accent!60!black, fill=accent!8!white] (c) at (3.6,0) {compilatore\\{\FiraCodeNL javac}};
    \node[b, draw=orange-gradient-6!80!black, fill=orange-gradient-1!30!white] (bc) at (7.2,0) {bytecode\\{\FiraCodeNL Programma.class}};
    \foreach \y/\m in {1.15/{JVM per Windows (x86-64)}, 0/{JVM per macOS (ARM)}, -1.15/{JVM per Linux}} {
      \node[b, draw=green-gradient-6, fill=green-gradient-1!45!white, text width=3.4cm] (j) at (11.6,\y) {\m};
      \draw[wtarrow] (bc.east) -- (j.west);
    }
    \draw[wtarrow] (s) -- (c);
    \draw[wtarrow] (c) -- (bc);
  \end{tikzpicture}
  \caption{Un solo file di bytecode, un interprete diverso per ogni tipo di
           macchina.}
  \label{fig:bytecode}
\end{figure}

\section{Le librerie}

In un linguaggio di alto livello il programmatore non deve programmare da sé
le funzioni di base: calcolare una radice quadrata o una tangente, stampare
un numero sullo schermo, leggere un valore dalla tastiera. Queste operazioni
sono state programmate e compilate una volta per tutte da chi produce il
compilatore, e sono a disposizione sotto forma di \term{funzioni} (per
esempio \code{tan(x)} calcola la tangente di \code{x}, \code{printf} stampa).

Il loro codice, già tradotto in linguaggio macchina, è raccolto in file
chiamati \term{librerie}: collezioni di funzioni pronte all'uso. Un
programma C “vero” è quindi fatto di due parti: il codice scritto da noi
e il codice delle librerie che usa.

\section{Come nasce un eseguibile}

Torniamo al compilatore. Finora lo abbiamo trattato come una scatola
chiusa, che riceve il sorgente e restituisce l'eseguibile; in realtà la
traduzione di un programma C avviene in più passaggi, ciascuno svolto da uno
strumento diverso. Conoscerli aiuta a capire i messaggi di errore, che
cambiano a seconda del passaggio in cui qualcosa è andato storto.

\subsection{Compilazione e linking}

La creazione di un eseguibile (in inglese \textit{build}) avviene in due
fasi (figura~\ref{fig:build}):

\begin{passi}
  \item \textbf{compilazione}: il sorgente viene tradotto in codice
        macchina, ma le funzioni di libreria che usa sono ancora
        \textbf{mancanti}. Il risultato è un file intermedio, il
        \term{file oggetto} (estensione \code{.o}), che non si può
        eseguire;
  \item \textbf{linking} (collegamento): uno o più file oggetto e una o più
        librerie (già in codice macchina) vengono \textbf{uniti} dal
        \term{linker}, che aggiunge le parti mancanti e produce un unico
        file eseguibile.
\end{passi}

\begin{figure}[H]
  \centering
  \begin{tikzpicture}[font=\Lato\scriptsize,
      file/.style={wtnodeplain, minimum width=2.6cm, minimum height=0.8cm},
      comp/.style={wtnode, draw=accent!60!black, fill=accent!8!white,
        font=\Lato\scriptsize\bfseries, minimum width=2.6cm},
      ogg/.style={wtnode, draw=primary!60!black, fill=primary!6!white,
        font=\Lato\scriptsize, minimum width=2.6cm, minimum height=0.8cm}]
    \node[file] (s1) at (0,0) {sorgente {\FiraCodeNL main.c}\\(contiene \code{main})};
    \node[file] (s2) at (3.6,0) {sorgente {\FiraCodeNL funzioni.c}\\(altre funzioni)};
    \node[comp] (c1) at (0,-1.3) {compilatore};
    \node[comp] (c2) at (3.6,-1.3) {compilatore};
    \node[ogg] (o1) at (0,-2.6) {{\FiraCodeNL main.o}\\oggetto, non eseguibile};
    \node[ogg] (o2) at (3.6,-2.6) {{\FiraCodeNL funzioni.o}\\oggetto, non eseguibile};
    \node[ogg, draw=orange-gradient-6!80!black, fill=orange-gradient-1!30!white] (lib) at (7.6,-2.6) {\textbf{libreria}\\funzioni precompilate};
    \node[wtnode, draw=green-gradient-6, fill=green-gradient-1!45!white, minimum width=9cm,
          font=\Lato\scriptsize\bfseries] (lk) at (3.8,-4.0) {linker};
    \node[wtnode, draw=red-gradient-6!70!black, fill=red-gradient-1!25!white,
          font=\Lato\scriptsize, minimum width=2.6cm] (ex) at (3.8,-5.3) {\textbf{eseguibile}\\{\FiraCodeNL programma}};
    \foreach \a/\b in {s1/c1, s2/c2, c1/o1, c2/o2} \draw[wtarrow] (\a) -- (\b);
    \draw[wtarrow, draw=green-gradient-6!70!black] (o1.south) -- (o1.south |- lk.north);
    \draw[wtarrow, draw=green-gradient-6!70!black] (o2.south) -- (o2.south |- lk.north);
    \draw[wtarrow, draw=green-gradient-6!70!black] (lib.south) -- (lib.south |- lk.north);
    \draw[wtarrow, line width=1.3pt] (lk) -- (ex);
  \end{tikzpicture}
  \caption{La compilazione produce i file oggetto, il linking li unisce fra
           loro e con le librerie.}
  \label{fig:build}
\end{figure}

Quando nel capitolo~\ref{chap:ambiente} abbiamo scritto
\code{gcc ciao.c}, gcc ha eseguito entrambe le fasi in un colpo solo. Le si
può però separare, ed è istruttivo farlo con il primo programma di esempio
del corso, \code{area\_cerchio.c} (lo vedremo per intero fra poco).
L'opzione \code{-c} (\textit{compile}) chiede a gcc di fermarsi dopo la
compilazione:

\begin{console}
$ gcc -c area_cerchio.c
$ ls -l area_cerchio.c area_cerchio.o
-rw-r--r-- 1 fabio fabio  255 ott 31 10:09 area_cerchio.c
-rw-r--r-- 1 fabio fabio 1960 ott 31 10:09 area_cerchio.o
$ nm area_cerchio.o
                 U __isoc99_scanf
                 U __stack_chk_fail
0000000000000000 T main
                 U printf
\end{console}

Il comando \code{nm} elenca i nomi che compaiono in un file oggetto. La
\code{T} accanto a \code{main} significa che il suo codice è
\textbf{presente} nel file (\textit{text}, cioè codice); la \code{U} accanto
a \code{printf} e \code{scanf} (che compare con il nome interno
\code{\_\_isoc99\_scanf}) significa \textbf{undefined}: il programma le usa,
ma il loro codice non c'è. Sono proprio le “funzioni di libreria ancora
mancanti”. Il linking le aggiunge:

\begin{console}
$ gcc -o area_cerchio area_cerchio.o
$ ./area_cerchio
Inserisci il raggio del cerchio: <<5>>
L'area del cerchio è: 78.54
\end{console}

Dato un file \code{.o} invece di un \code{.c}, gcc salta la compilazione e
fa solo il linking. Separare le due fasi diventa utile quando un programma è
diviso in più file: se ne cambia uno, basta ricompilare quello e rifare il
linking. Lo vedremo nel capitolo~\ref{chap:make}.

\subsection{La catena completa}

La figura~\ref{fig:catena-completa} mostra tutto il percorso di un programma
C, dall'editor fino alla CPU, con gli strumenti coinvolti.

\begin{figure}[H]
  \centering
  \begin{tikzpicture}[font=\Lato\scriptsize,
      st/.style={wtnode, font=\Lato\scriptsize\bfseries, minimum width=2.5cm, minimum height=0.85cm,
        draw=accent!60!black, fill=accent!8!white},
      fl/.style={wtnodeplain, font=\Lato\scriptsize, minimum width=2.5cm, minimum height=0.85cm}]
    \node[st] (ed) at (0,0) {editor\\{\mdseries (VS Code)}};
    \node[fl] (src) at (3.4,0) {programma sorgente\\{\FiraCodeNL prova.c}};
    \node[st] (pre) at (6.8,0) {preprocessore};
    \node[st] (cmp) at (10.2,0) {compilatore};
    \node[fl] (obj) at (13.6,0) {programma oggetto\\{\FiraCodeNL prova.o}};
    \node[st] (lnk) at (13.6,-2.0) {linker};
    \node[fl, draw=orange-gradient-6!80!black, fill=orange-gradient-1!30!white] (lib) at (13.6,-3.6) {librerie\\{\FiraCodeNL libc, libm\dots}};
    \node[fl] (exe) at (10.2,-2.0) {programma eseguibile\\{\FiraCodeNL prova}};
    \node[st] (ld) at (6.8,-2.0) {loader\\{\mdseries (sistema operativo)}};
    \node[fl, draw=red-gradient-6!70!black, fill=red-gradient-1!20!white] (mem) at (3.4,-2.0) {programma\\in memoria};
    \node[st, draw=primary!65!black, fill=primary!12!white] (cpu) at (0,-2.0) {CPU};
    \node[st, draw=green-gradient-6, fill=green-gradient-1!45!white] (dbg) at (3.4,-3.6) {debugger\\{\mdseries (gdb)}};
    \foreach \a/\b in {ed/src, src/pre, pre/cmp, cmp/obj, obj/lnk, lnk/exe, exe/ld, ld/mem, mem/cpu}
      \draw[wtarrow] (\a) -- (\b);
    \draw[wtarrow] (lib) -- (lnk);
    \draw[wtarrowacc, dashed, <->] (dbg) -- node[wtlab, right=2pt, fill=none] {osserva e controlla} (mem);
  \end{tikzpicture}
  \caption{Dal sorgente al processore: gli strumenti (in viola) e i file che
           producono.}
  \label{fig:catena-completa}
\end{figure}

\begin{itemize}
  \item L'\textbf{editor} serve a scrivere il sorgente.
  \item Il \term{preprocessore} lavora sul testo prima della compilazione
        vera e propria: esegue le righe che cominciano con \code{\#}.
        \code{\#include <stdio.h>} viene sostituita dall'intero contenuto
        del file \code{stdio.h}; \code{\#define PI 3.14159} fa sì che ogni
        \code{PI} del programma diventi \code{3.14159}.
  \item Il \textbf{compilatore} traduce in codice macchina e produce il
        file oggetto; il \textbf{linker} lo unisce alle librerie.
  \item Il \term{loader} è la parte del sistema operativo che, quando
        scriviamo \code{./prova}, copia l'eseguibile in memoria e fa partire
        la CPU dalla sua prima istruzione (capitolo~\ref{chap:calcolatore}).
  \item Il \term{debugger} permette di eseguire il programma un passo alla
        volta, fermarlo e guardare il contenuto delle variabili: è lo
        strumento per la \textit{messa a punto}.
\end{itemize}

Il lavoro del preprocessore si può vedere con l'opzione \code{-E}, che
ferma gcc subito dopo questa fase e stampa il testo risultante. Il nostro
programma di 15 righe diventa di oltre 800, quasi tutte provenienti da
\code{stdio.h}; in fondo c'è il nostro codice, con \code{PI} già sostituito:

\begin{console}
$ gcc -E area_cerchio.c | wc -l
828
$ gcc -E area_cerchio.c | tail -6
    scanf("%f", &raggio);

    area = 3.14159 * raggio * raggio;

    printf("L'area del cerchio è: %.2f\n", area);
\end{console}

(Il simbolo \code{|} è una \textit{pipe}: passa l'output di un comando a un
altro. \code{wc -l} conta le righe, \code{tail -6} mostra le ultime 6.)

\subsection{Librerie statiche e dinamiche}

Il linking può collegare una libreria in due modi.

Con una \term{libreria statica}, il codice delle funzioni usate viene
\textbf{copiato} dentro il file eseguibile. Le stesse funzioni possono
essere usate da più programmi, ma ognuno ne contiene una copia
(figura~\ref{fig:librerie}, a sinistra).

Con una \term{libreria dinamica}, nell'eseguibile \textbf{non} viene inserito
il codice delle funzioni, ma solo un \textbf{riferimento}: il nome del file
della libreria e il nome della funzione. Su Windows sono le \textit{DLL}
(\textit{Dynamic Link Libraries}), su Linux e macOS le \textit{shared
libraries} (file \code{.so} e \code{.dylib}). Quando parte il primo
programma che usa una libreria dinamica, il sistema operativo la carica in
memoria; se poi parte un altro programma che usa la stessa libreria, la
trova già in memoria e la \textbf{condivide} (figura~\ref{fig:librerie}, a
destra).

\begin{figure}[H]
  \centering
  \begin{tikzpicture}[font=\Lato\scriptsize,
      app/.style={draw=primary!65!black, fill=primary!10!white, line width=0.7pt,
        minimum width=1.6cm, minimum height=0.55cm, font=\Lato\tiny\bfseries},
      lib/.style={draw=orange-gradient-6!80!black, fill=orange-gradient-1!35!white,
        line width=0.7pt, minimum width=1.6cm, minimum height=1.3cm, font=\Lato\tiny,
        align=center}]
    % statiche
    \node[font=\Lato\scriptsize\bfseries] at (1.85,1.2) {librerie statiche};
    \node[app] (a1) at (0.9,0.45) {applicazione A};
    \node[lib, anchor=north] at ([yshift=0.7pt]a1.south) {copia della\\libreria};
    \node[app] (a2) at (2.8,0.45) {applicazione B};
    \node[lib, anchor=north] at ([yshift=0.7pt]a2.south) {copia della\\libreria};
    % dinamiche
    \node[font=\Lato\scriptsize\bfseries] at (8.1,1.2) {librerie dinamiche};
    \node[app] (b1) at (6.9,0.45) {applicazione A};
    \node[draw=primary!65!black, fill=white, line width=0.7pt, minimum width=1.6cm,
          minimum height=0.45cm, anchor=north, font=\Lato\tiny] (r1) at ([yshift=0.7pt]b1.south) {riferimento};
    \node[app] (b2) at (9.3,0.45) {applicazione B};
    \node[draw=primary!65!black, fill=white, line width=0.7pt, minimum width=1.6cm,
          minimum height=0.45cm, anchor=north, font=\Lato\tiny] (r2) at ([yshift=0.7pt]b2.south) {riferimento};
    \node[lib] (sh) at (8.1,-1.75) {libreria\\condivisa\\(una sola,\\in memoria)};
    \draw[mrif] (r1.south) -- (sh.north west);
    \draw[mrif] (r2.south) -- (sh.north east);
  \end{tikzpicture}
  \caption{A sinistra ogni eseguibile contiene la sua copia della libreria;
           a destra gli eseguibili contengono solo un riferimento a un'unica
           copia, condivisa.}
  \label{fig:librerie}
\end{figure}

I vantaggi delle librerie dinamiche sono evidenti: gli eseguibili sono più
piccoli, un programma parte più in fretta se la libreria è già in memoria e,
se una funzione va aggiornata (una nuova versione, la correzione di un
errore), basta sostituire la libreria, senza ritoccare i programmi. Il
rovescio della medaglia: se la libreria manca o è danneggiata, il programma
non funziona.

gcc usa le librerie dinamiche, se non si chiede diversamente. Il comando
\code{ldd} mostra da quali dipende un eseguibile: il nostro dipende da
\code{libc.so.6}, la libreria standard del C. Con l'opzione \code{-static}
si ottiene invece un collegamento statico, e la differenza di dimensione
parla da sola:

\begin{console}
$ ldd area_cerchio
	linux-vdso.so.1 (0x00007f1b90918000)
	libc.so.6 => /lib/x86_64-linux-gnu/libc.so.6 (0x00007f1b90600000)
	/lib64/ld-linux-x86-64.so.2 (0x00007f1b9091a000)
$ gcc -static -o area_statico area_cerchio.o
$ ls -l area_cerchio area_statico
-rwxr-xr-x 1 fabio fabio  16072 ott 31 10:09 area_cerchio
-rwxr-xr-x 1 fabio fabio 915784 ott 31 10:09 area_statico
\end{console}

Lo stesso programma pesa 16 kilobyte con la libreria dinamica e quasi un
megabyte con quella statica, che contiene dentro di sé tutto il codice di
\code{printf}, \code{scanf} e delle funzioni che queste usano. (Su macOS
\code{-static} non è disponibile per i programmi normali, e al posto di
\code{ldd} si usa \code{otool -L}.)

\section{Gli algoritmi}

\subsection{Che cos'è un algoritmo}

Possiamo ora precisare la definizione del capitolo precedente. Un
algoritmo è la \textbf{descrizione operativa} di come risolvere un problema:
un \textbf{elenco finito di istruzioni elementari}, da eseguire in
successione. Le singole istruzioni devono:

\begin{itemize}
  \item richiedere un \textbf{tempo finito};
  \item \textbf{non essere ambigue};
  \item avere un \textbf{effetto osservabile}.
\end{itemize}

Inoltre l'algoritmo deve produrre un \textbf{risultato}, sempre lo stesso a
partire dalle stesse condizioni iniziali, e le risorse che richiede devono
essere \textbf{ragionevoli} per la tecnologia attuale: il tempo complessivo
di esecuzione, la memoria massima necessaria, eventuali dispositivi
periferici particolari.

Il termine nasce in informatica, ma si applica a qualunque procedimento.
Cambiano le istruzioni elementari, non le caratteristiche:

\begin{center}
\renewcommand{\arraystretch}{1.3}
\begin{tabular}{@{}ll@{}}
  \toprule
  \textbf{Algoritmo per\dots} & \textbf{un'istruzione elementare} \\
  \midrule
  risolvere un'equazione di secondo grado & calcola $b^2 - 4ac$ \\
  preparare una ricetta                   & sbatti 4 albumi \\
  montare un armadio IKEA                 & unisci i pannelli 12 e 13 con una vite C \\
  avviare il motore di un aereo           & inserisci i magneti \\
  installare un'app sul cellulare         & tocca il pulsante Installa \\
  \bottomrule
\end{tabular}
\end{center}

Sono invece istruzioni \textbf{non adeguate}:

\begin{itemize}
  \item “calcola tutte le cifre di $\pi$”: le cifre sono infinite, quindi
        l'istruzione non termina mai (tempo non finito);
  \item “aggiungi alcuni degli ingredienti”: quali? quanti? (ambigua).
\end{itemize}

E un algoritmo che con gli stessi dati iniziali produce risultati diversi
\textbf{non è corretto}, perché non è deterministico: un'equazione di
secondo grado deve dare sempre le stesse soluzioni per gli stessi valori di
$a$, $b$ e $c$.

\subsection{Quando le risorse non bastano}

Un algoritmo che richiede un tempo irragionevole per la tecnologia attuale,
considerati costi e benefici, \textbf{non risolve il problema}, anche se è
corretto. L'esempio delle lezioni: aprire un file zip protetto da password
senza conoscerla. Un algoritmo c'è, ed è semplicissimo: provare tutte le
combinazioni di caratteri finché non si trova quella giusta (il cosiddetto
attacco a \textit{forza bruta}).

Supponiamo che la password sia di 8 caratteri, scelti fra i 96 caratteri
stampabili del codice ASCII standard. Per il primo carattere ci sono 96
possibilità; per ciascuna, 96 per il secondo, e così via: in tutto
$96^8 \approx 7{,}2 \cdot 10^{15}$ password possibili, cioè circa 7,2
milioni di miliardi. Con un computer che ne prova 34 milioni al secondo,
senza mai fermarsi:

\[
  \frac{7{,}2 \cdot 10^{15}}{3{,}4 \cdot 10^{7}\ \text{al secondo}}
  \approx 2{,}1 \cdot 10^{8}\ \text{secondi} \approx 6{,}7\ \text{anni}.
\]

\warning{Un conto da rifare}{
  Le slide parlano di “quasi 5 anni” per provare tutte le password; il
  calcolo esatto dà circa \textbf{6,7 anni}. In media, però, la password
  giusta si trova dopo averne provate la metà, cioè in circa 3,4 anni. La
  conclusione non cambia: per chi ha bisogno del file oggi, l'algoritmo è
  inutile.}

La cosa più istruttiva è che cosa succede allungando la password: ogni
carattere in più \textbf{moltiplica per 96} il numero di tentativi.

\begin{center}
\renewcommand{\arraystretch}{1.3}
\begin{tabular}{@{}crr@{}}
  \toprule
  \textbf{lunghezza} & \textbf{password possibili} & \textbf{tempo per provarle tutte} \\
  \midrule
  8  & $7{,}2 \cdot 10^{15}$ & circa 6,7 anni \\
  9  & $6{,}9 \cdot 10^{17}$ & circa 645 anni \\
  10 & $6{,}6 \cdot 10^{19}$ & circa 62\,000 anni \\
  \bottomrule
\end{tabular}
\end{center}

Una crescita di questo tipo si chiama \textit{esponenziale}, ed è il motivo
per cui le password lunghe sono sicure: non serve un algoritmo più furbo per
violarle, servirebbe un tempo che nessuno ha.

\section{Come si descrive un algoritmo}

Prima di scrivere un programma conviene descrivere l'algoritmo in una forma
intermedia, più libera di un linguaggio di programmazione. Le forme più
usate sono:

\begin{itemize}
  \item la \textbf{descrizione testuale}, in italiano;
  \item lo \term{pseudo-codice}: una descrizione testuale con un certo
        livello di formalismo, con parole chiave fisse (\code{LEGGI},
        \code{STAMPA}, \code{SE}\dots) e una struttura simile a quella di un
        programma, ma senza le regole rigide di un linguaggio vero;
  \item il \term{diagramma di flusso} (\textit{flow-chart}): una
        rappresentazione grafica, in cui ogni forma ha un significato
        (figura~\ref{fig:forme});
  \item il \textbf{linguaggio di alto livello}, cioè il programma vero e
        proprio.
\end{itemize}

\begin{figure}[H]
  \centering
  \begin{tikzpicture}[font=\Lato\scriptsize]
    \node[fcterm] (a) at (0,0) {Inizio};
    \node[fcio] (b) at (3.4,0) {leggi \code{x}};
    \node[fcproc] (c) at (6.8,0) {\code{y = 2 * x}};
    \node[fcdec] (d) at (10.2,0) {\code{y > 10}?};
    \node[fcsub] (e) at (13.6,0) {ordina};
    \foreach \nodo/\testo in {a/{inizio e fine}, b/{input e output\\(leggere, stampare)},
                       c/{un'istruzione\\(un calcolo)}, d/{una decisione:\\due uscite, sì e no},
                       e/{un sottoprogramma\\(cap.~\ref{chap:strutturata})}}
      \node[mnota, align=center, anchor=north] at ([yshift=-0.3cm]\nodo.center |- d.south) {\testo};
  \end{tikzpicture}
  \caption{Le forme dei diagrammi di flusso.}
  \label{fig:forme}
\end{figure}

Conviene anche organizzare la descrizione in tre parti: l'\textbf{input}
(i dati di partenza), l'\textbf{elaborazione} (i calcoli) e
l'\textbf{output} (i risultati). Vediamo i tre esempi delle lezioni, in
tutte le forme. Il codice C è commentato solo a grandi linee: ogni
dettaglio sarà spiegato nel capitolo~\ref{chap:primi-programmi}. Qui conta
vedere come ogni riga di pseudo-codice diventa una riga di programma.

\subsection{Esempio 1: l'area di un cerchio}

\begin{itemize}
  \item \textbf{Input}: il raggio del cerchio $r$, un numero reale positivo.
  \item \textbf{Elaborazione}: l'area si calcola con la formula
        $A = \pi \cdot r^2$, dove $\pi$ è la costante che rappresenta il
        rapporto fra la circonferenza di un cerchio e il suo diametro.
  \item \textbf{Output}: il valore numerico dell'area.
\end{itemize}

In forma testuale: \textit{acquisisci dall'utente il valore del raggio;
calcola l'area con la formula $\pi \cdot r \cdot r$; visualizza il valore
ottenuto.} In pseudo-codice e come diagramma di flusso:

\begin{affianco}[0.42][c]
\begin{lstlisting}[style=pseudo, name=pseudocodice]
BEGIN
  LEGGI raggio
  AREA = PI * raggio * raggio
  STAMPA AREA
END
\end{lstlisting}
\risultato
\begin{immagine}
\begin{tikzpicture}[font=\Lato\scriptsize]
  \node[fcterm] (a) at (0,0) {Inizio};
  \node[fcio] (b) at (0,-1.0) {leggi \code{raggio}};
  \node[fcproc] (c) at (0,-2.0) {\code{area = PI * raggio * raggio}};
  \node[fcio] (d) at (0,-3.0) {stampa \code{area}};
  \node[fcterm] (e) at (0,-4.0) {Fine};
  \foreach \x/\y in {a/b, b/c, c/d, d/e} \draw[fcfreccia] (\x) -- (\y);
\end{tikzpicture}
\end{immagine}
\end{affianco}

Ed ecco il programma C, con l'esecuzione accanto:

\begin{affianco}[0.55]
\begin{lstlisting}[style=c, name=area_cerchio.c]
#include <stdio.h>
#define PI 3.14159

int main() {
    float raggio, area;

    printf("Inserisci il raggio del cerchio: ");
    scanf("%f", &raggio);

    area = PI * raggio * raggio;

    printf("L'area del cerchio è: %.2f\n", area);

    return 0;
}
\end{lstlisting}
\risultato
\begin{console}
$ gcc -Wall -o area_cerchio area_cerchio.c
$ ./area_cerchio
Inserisci il raggio del cerchio: <<5>>
L'area del cerchio è: 78.54
\end{console}
\end{affianco}

La corrispondenza con lo pseudo-codice è diretta:

\begin{center}
\renewcommand{\arraystretch}{1.35}
\begin{tabular}{@{}>{\FiraCodeNL\small}l>{\FiraCodeNL\small}l>{\raggedright\arraybackslash}p{6.4cm}@{}}
  \toprule
  \normalfont\textbf{pseudo-codice} & \normalfont\textbf{C} & \textbf{nota} \\
  \midrule
  & float raggio, area; & in C bisogna dichiarare le variabili prima di
                          usarle, dicendone il tipo (\code{float}: numero
                          reale) \\
  LEGGI raggio & scanf("\%f", \&raggio); & prima si stampa un messaggio
                          che dice all'utente che cosa scrivere \\
  AREA = PI * raggio * raggio & area = PI * raggio * raggio; & identico;
                          \code{PI} è definito dalla riga \code{\#define} \\
  STAMPA AREA & printf("\dots \%.2f\textbackslash n", area); & \code{\%.2f}
                          significa: un numero reale con 2 decimali \\
  \bottomrule
\end{tabular}
\end{center}

\subsection{Esempio 2: da Fahrenheit a Celsius}

\begin{itemize}
  \item \textbf{Input}: una temperatura in gradi Fahrenheit (un numero
        reale).
  \item \textbf{Elaborazione}: la formula di conversione è
        $\mathit{Celsius} = (\mathit{Fahrenheit} - 32) \cdot 5 / 9$.
  \item \textbf{Output}: la temperatura in gradi Celsius.
\end{itemize}

Lo pseudo-codice ha la stessa forma del primo esempio, e così il programma:

\begin{affianco}[0.5][c]
\begin{lstlisting}[style=pseudo, name=pseudocodice]
BEGIN
  LEGGI fahrenheit
  CELSIUS = (fahrenheit - 32) * 5 / 9
  STAMPA CELSIUS
END
\end{lstlisting}
\risultato
\wtnota{LEGGI diventa \code{scanf}, il calcolo resta identico, STAMPA
diventa \code{printf}.}
\end{affianco}

\begin{affianco}[0.55]
\begin{lstlisting}[style=c, name=conversione.c]
#include <stdio.h>

int main() {
    float fahrenheit, celsius;

    printf("Inserisci la temperatura in Fahrenheit: ");
    scanf("%f", &fahrenheit);

    celsius = (fahrenheit - 32) * 5 / 9;

    printf("La temperatura in Celsius è: %.2f\n", celsius);

    return 0;
}
\end{lstlisting}
\risultato
\begin{console}
$ gcc -Wall -o conversione conversione.c
$ ./conversione
Inserisci la temperatura in Fahrenheit: <<98.6>>
La temperatura in Celsius è: 37.00
\end{console}
\end{affianco}

\warning{L'ordine delle operazioni conta}{
  Matematicamente $(F - 32) \cdot 5 / 9$ e $5 / 9 \cdot (F - 32)$ sono la
  stessa cosa. In C no: scrivendo \code{celsius = 5 / 9 * (fahrenheit - 32);}
  il risultato è sempre \textbf{0}. Il motivo è che \code{5} e \code{9} sono
  numeri interi, e in C la divisione fra interi dà un intero, troncando i
  decimali: \code{5 / 9} vale 0, e 0 per qualunque cosa fa 0. Nella versione
  del programma, invece, \code{(fahrenheit - 32)} è un \code{float}, quindi
  anche le operazioni successive avvengono fra reali. È una delle trappole
  più comuni del C, e la ritroveremo nel capitolo~\ref{chap:primi-programmi}.}

\subsection{Esempio 3: il segno di un numero}

Il terzo esempio introduce una \textbf{scelta}: il programma deve fare cose
diverse a seconda del dato.

\begin{itemize}
  \item \textbf{Input}: un numero intero.
  \item \textbf{Elaborazione}: se il numero è maggiore di zero è positivo;
        se è minore di zero è negativo; altrimenti è nullo.
  \item \textbf{Output}: un messaggio che indica se il numero è positivo,
        negativo o nullo.
\end{itemize}

\begin{affianco}[0.46][c]
\begin{lstlisting}[style=pseudo, name=pseudocodice]
BEGIN
  LEGGI numero
  SE numero > 0 ALLORA
    STAMPA "Il numero è positivo"
  ALTRIMENTI SE numero < 0 ALLORA
    STAMPA "Il numero è negativo"
  ALTRIMENTI
    STAMPA "Il numero è nullo"
  END SE
END
\end{lstlisting}
\risultato
\begin{immagine}
\begin{tikzpicture}[font=\Lato\scriptsize]
  \node[fcterm] (a) at (0,0) {Inizio};
  \node[fcio] (b) at (0,-0.95) {leggi \code{numero}};
  \node[fcdec] (c) at (0,-2.15) {\code{numero > 0}?};
  \node[fcio] (p) at (3.1,-2.15) {stampa “positivo”};
  \node[fcdec] (d) at (0,-3.75) {\code{numero < 0}?};
  \node[fcio] (n) at (3.1,-3.75) {stampa “negativo”};
  \node[fcio] (z) at (0,-5.05) {stampa “nullo”};
  \node[fcterm] (e) at (1.55,-6.05) {Fine};
  \draw[fcfreccia] (a) -- (b);
  \draw[fcfreccia] (b) -- (c);
  \draw[fcfreccia] (c) -- node[fcsi, above] {sì} (p);
  \draw[fcfreccia] (c) -- node[fcno, right] {no} (d);
  \draw[fcfreccia] (d) -- node[fcsi, above] {sì} (n);
  \draw[fcfreccia] (d) -- node[fcno, right] {no} (z);
  \draw[fcfreccia] (p.east) -- ++(0.35,0) |- (e.east);
  \draw[fcfreccia] (n.east) -- ++(0.35,0) |- (e.east);
  \draw[fcfreccia] (z.south) |- (e.west);
\end{tikzpicture}
\end{immagine}
\end{affianco}

Nel diagramma di flusso si vede bene che, qualunque sia il numero, viene
eseguita \textbf{una sola} delle tre stampe: le strade si separano ai rombi
e si ricongiungono prima della fine. In C la scelta si scrive con
\code{if} (“se”) ed \code{else} (“altrimenti”):

\begin{affianco}[0.55]
\begin{lstlisting}[style=c, name=positivo.c]
#include <stdio.h>

int main() {
    int numero;

    printf("Inserisci un numero: ");
    scanf("%d", &numero);

    if (numero > 0) {
        printf("Il numero è positivo.\n");
    } else if (numero < 0) {
        printf("Il numero è negativo.\n");
    } else {
        printf("Il numero è zero.\n");
    }

    return 0;
}
\end{lstlisting}
\risultato
\begin{console}
$ ./positivo
Inserisci un numero: <<7>>
Il numero è positivo.
$ ./positivo
Inserisci un numero: <<-3>>
Il numero è negativo.
$ ./positivo
Inserisci un numero: <<0>>
Il numero è zero.
\end{console}
\wtnota{Per verificare un programma con una scelta, lo si prova con un
dato per ogni strada possibile.}
\end{affianco}

Si noti l'ultima nota: un programma con tre strade va provato \textbf{tre
volte}, una per strada. Un errore nascosto in un ramo che non si prova mai
resta lì, in attesa del giorno sbagliato.

\insintesi{
  Un algoritmo è un elenco finito di istruzioni elementari, non ambigue,
  ciascuna eseguibile in un tempo finito, che produce sempre lo stesso
  risultato dagli stessi dati con risorse ragionevoli. Lo si descrive in
  italiano, in pseudo-codice, con un diagramma di flusso e infine in un
  linguaggio di alto livello. Il programma sorgente va tradotto: un
  \textbf{interprete} lo fa un'istruzione alla volta a ogni esecuzione, un
  \textbf{compilatore} una volta sola. Per il C la traduzione passa per il
  preprocessore, il compilatore (che produce i file oggetto) e il linker
  (che li unisce alle librerie, statiche o dinamiche); al momento
  dell'esecuzione il loader carica il programma in memoria.}

\section{Riferimenti}

\begin{itemize}
  \item \textbf{Deitel \& Deitel}, \textit{Il linguaggio C}: capitoli 1 e 3
        (ambiente di sviluppo, algoritmi, pseudo-codice).
  \item \textbf{GCC --- Overall Options}: le opzioni \code{-E}, \code{-S},
        \code{-c} e \code{-o}. \\
        \urlx{https://gcc.gnu.org/onlinedocs/gcc/Overall-Options.html}
\end{itemize}


  \chapter{La programmazione strutturata}
\label{chap:strutturata}

Prima di scrivere programmi in C conviene sapere da dove viene il linguaggio
e, soprattutto, quale \textbf{disciplina} segue: la programmazione
strutturata. È un'idea semplice e potentissima: \textit{ogni} programma si
può costruire combinando tre soli schemi, la sequenza, la selezione e
l'iterazione. Questo capitolo li presenta con i diagrammi di flusso, prima
con esempi della vita di tutti i giorni e poi con problemi di informatica;
nei capitoli successivi ciascuno diventerà un'istruzione del C.

\section{Il linguaggio C}

\subsection{Da dove viene}

Il C è stato creato da \textbf{Dennis Ritchie}, ai Bell Labs, nel 1972,
come linguaggio per scrivere il sistema operativo \textbf{Unix}. La prima
versione di Unix, opera di \textbf{Ken Thompson}, era scritta in assembly e
in un linguaggio chiamato \textbf{B}. Il B aveva dei limiti seri:

\begin{itemize}
  \item non era \textbf{tipizzato}: conosceva un solo tipo di dato, la
        \textit{parola macchina} (il gruppo di bit su cui lavora il
        processore), e non distingueva fra numeri interi, caratteri o
        indirizzi;
  \item faceva un uso pesante di puntatori e indirizzi di memoria.
\end{itemize}

Il C nasce per superare questi limiti: è un'evoluzione del B che incorpora
i \textbf{tipi} (ricordate gli “occhiali” del
capitolo~\ref{chap:calcolatore}: è il tipo a dire come interpretare i
bit), e rispetta i canoni della programmazione strutturata. Nel tempo il
linguaggio si è evoluto (figura~\ref{fig:storia-c}):

\begin{itemize}
  \item il \textbf{C tradizionale} (\textit{Traditional C}), descritto nel
        1978 nel libro di Kernighan e Ritchie e diffuso all'inizio degli anni
        '80, introduce fra l'altro i tipi \code{void} ed \code{enum};
  \item l'\textbf{ANSI C}, della fine degli anni '80, aggiunge il tipo
        \code{void *}, i \textbf{prototipi} delle funzioni e una definizione
        del linguaggio molto più precisa.
\end{itemize}

\begin{figure}[H]
  \centering
  \begin{tikzpicture}[font=\Lato\scriptsize]
    \draw[neutral!45!black, line width=1pt, -{Stealth[length=6pt]}] (0,0) -- (17.6,0);
    \foreach \x/\anno/\cosa/\h in {
        0.6/1967/{BCPL\\(Martin Richards)}/0.45,
        2.6/1969/{B\\(Ken Thompson)}/1.3,
        4.6/1972/{\textbf{C}\\(Dennis Ritchie)}/0.45,
        6.6/1978/{libro di Kernighan\\e Ritchie: il C\\tradizionale}/1.3,
        8.6/1989/{ANSI C\\(ISO nel 1990)}/0.45,
        10.6/1999/{C99}/1.3,
        12.6/2011/{C11}/0.45,
        14.6/2018/{C17}/1.3,
        16.6/2024/{C23}/0.45} {
      \fill[primary!70!black] (\x,0) circle (2.3pt);
      \node[anchor=north, font=\Lato\scriptsize\bfseries, text=primary!55!black] at (\x,-0.12) {\anno};
      \draw[neutral!55!white] (\x,0.1) -- (\x,\h);
      \node[anchor=south, align=center, inner sep=1.5pt] at (\x,\h) {\cosa};
    }
  \end{tikzpicture}
  \caption{Le tappe del C, dai suoi antenati alle versioni dello standard.}
  \label{fig:storia-c}
\end{figure}

\warning{Una precisazione storica}{
  Le slide datano il B al 1967. Il 1967 è l'anno del suo antenato, il
  \textit{BCPL} di Martin Richards; il B fu sviluppato da Ken Thompson
  intorno al 1969, partendo proprio dal BCPL.}

\subsection{Lo standard}

L'\textbf{ANSI} (\textit{American National Standards Institute}) è l'ente
statunitense che definisce gli standard per molti tipi di sistemi, compresi
i linguaggi di programmazione. Lo standard ANSI del C (1989) fu approvato
nel 1990 dall'\textbf{ISO} (\textit{International Organization for
Standardization}), e l'ANSI C divenne così lo standard internazionale.

Uno standard specifica il \textbf{formalismo} del linguaggio (quali testi
sono programmi C corretti, cioè la sua \textit{sintassi}) e la sua
\textbf{semantica} (che cosa significa ciascun costrutto, cioè che cosa
deve fare il programma). Lo scopo è triplice:

\begin{itemize}
  \item \textbf{portabilità}: lo stesso programma si compila ed esegue allo
        stesso modo con compilatori diversi, su macchine diverse;
  \item \textbf{affidabilità}: il comportamento è definito, non lasciato
        all'interpretazione di chi scrive il compilatore;
  \item \textbf{manutenibilità}: chiunque conosca lo standard può leggere e
        modificare il programma.
\end{itemize}

Dopo il 1990 lo standard è stato aggiornato più volte (C99, C11, C17, C23).
Le versioni recenti di gcc seguono per default una di queste, con alcune
estensioni proprie; i programmi di questa dispensa funzionano con tutte.

\section{La programmazione strutturata}

\subsection{Il problema: gli spaghetti}

Nel capitolo~\ref{chap:calcolatore} abbiamo visto che in assembly
ripetizioni e scelte si ottengono con i \textbf{salti}: “vai
all'istruzione all'indirizzo tale”. I primi linguaggi di alto livello
ereditarono questa abitudine con l'istruzione \code{GOTO} (“vai a”), che
permette di saltare in qualunque punto del programma. In un programma
grande, pieno di salti in avanti e all'indietro, seguire l'esecuzione
diventa come seguire uno spaghetto nel piatto: è il cosiddetto
\term{spaghetti code} (figura~\ref{fig:spaghetti}, a sinistra), di bassa
qualità e leggibilità.

Nel 1968--1969 \textbf{Edsger W. Dijkstra} criticò duramente l'uso del
\code{GOTO} e propose di regolamentare e standardizzare le metodologie di
programmazione: è la \term{programmazione strutturata}. L'obiettivo è
rendere più facile la \textbf{lettura}, la \textbf{modifica} e la
\textbf{manutenzione} dei programmi.

\begin{figure}[H]
  \centering
  \begin{tikzpicture}[font=\Lato\scriptsize]
    % spaghetti
    \node[font=\Lato\small\bfseries, text=red-gradient-6!60!black] at (0,0.9) {con i GOTO};
    \foreach \i in {0,...,5}
      \node[fcproc, minimum width=2.2cm] (s\i) at (0,-\i*0.85) {istruzione \pgfmathparse{int(\i+1)}\pgfmathresult};
    \foreach \i/\j in {0/1, 1/2, 2/3, 3/4, 4/5}
      \draw[fcfreccia] (s\i) -- (s\j);
    \begin{scope}[red-gradient-6!70!black]
      \draw[fcfreccia, draw=red-gradient-6!70!black] (s3.east) to[out=0, in=0, looseness=1.7] (s0.east);
      \draw[fcfreccia, draw=red-gradient-6!70!black] (s1.west) to[out=180, in=180, looseness=1.5] (s4.west);
      \draw[fcfreccia, draw=red-gradient-6!70!black] (s5.east) to[out=0, in=0, looseness=1.3] (s2.east);
      \draw[fcfreccia, draw=red-gradient-6!70!black] (s2.west) to[out=180, in=180, looseness=1.3] (s3.west);
      \draw[fcfreccia, draw=red-gradient-6!70!black] (s4.east) to[out=-10, in=0, looseness=1.6] (s1.east);
    \end{scope}
    \node[mnota, align=center, anchor=north] at (0,-4.75) {per sapere che cosa viene\\eseguito dopo, bisogna\\seguire ogni freccia};
    % strutturato
    \begin{scope}[xshift=8.6cm]
      \node[font=\Lato\small\bfseries, text=green-gradient-6!50!black] at (0,0.9) {strutturato};
      \draw[green-gradient-6!80!black, line width=0.8pt, rounded corners=5pt, fill=green-gradient-1!25!white]
        (-3.6,0.45) rectangle (3.6,-4.4);
      \node[font=\Lato\scriptsize\bfseries, anchor=north west] at (-3.5,0.4) {sequenza};
      \node[fcproc, minimum width=2.4cm] (a) at (0,0.05) {istruzione};
      \draw[accent!60!black, line width=0.8pt, rounded corners=5pt, fill=accent!7!white]
        (-3.1,-0.5) rectangle (3.1,-3.35);
      \node[font=\Lato\scriptsize\bfseries, anchor=north west] at (-3.0,-0.55) {iterazione: ripeti finché\dots};
      \draw[primary!60!black, line width=0.8pt, rounded corners=5pt, fill=primary!6!white]
        (-2.6,-1.05) rectangle (2.6,-3.05);
      \node[font=\Lato\scriptsize\bfseries, anchor=north west] at (-2.5,-1.1) {selezione: se\dots altrimenti\dots};
      \node[fcproc, minimum width=1.9cm] at (-1.2,-2.2) {istruzione};
      \node[fcproc, minimum width=1.9cm] at (1.2,-2.2) {istruzione};
      \node[fcproc, minimum width=2.4cm] (z) at (0,-3.85) {istruzione};
      \node[mnota, align=center, anchor=north] at (0,-4.75) {blocchi annidati, ciascuno con\\\textbf{una sola entrata e una sola uscita}:\\si legge dall'alto in basso};
    \end{scope}
  \end{tikzpicture}
  \caption{Lo stesso tipo di logica scritto con salti arbitrari (a
           sinistra) e con le strutture di controllo (a destra).}
  \label{fig:spaghetti}
\end{figure}

\subsection{Il teorema di Böhm e Jacopini}

Si può davvero fare a meno dei salti arbitrari? La risposta è un teorema,
dimostrato nel 1966 da due informatici italiani, \textbf{Corrado Böhm} e
\textbf{Giuseppe Jacopini}.

\dfn{Teorema di Böhm-Jacopini}{
  Qualsiasi programma scritto usando il \code{GOTO} può essere riscritto
  senza, a patto di avere a disposizione tre tipi di \textbf{strutture di
  controllo}: la \textbf{sequenza}, la \textbf{selezione} e
  l'\textbf{iterazione}.}

Le tre strutture bastano, quindi, per \textit{qualunque} algoritmo. Hanno
una caratteristica comune decisiva: ciascuna ha \textbf{una sola entrata e
una sola uscita}. Per questo si possono incastrare l'una dentro l'altra
come scatole (figura~\ref{fig:spaghetti}, a destra), e un programma si
legge dall'alto in basso, capendo ogni blocco senza dover guardare il resto.

\info{Che fine ha fatto il \code{GOTO}?}{
  Il C ha ancora un'istruzione \code{goto}, per ragioni storiche e per pochi
  casi molto particolari. In questo corso \textbf{non si usa}: tutto ciò che
  serve si scrive con le strutture di controllo.}

\section{Le tre strutture di controllo}

Le strutture di controllo del flusso sono dunque tre
(figura~\ref{fig:tre-strutture}). Vediamole nel dettaglio, ciascuna con un
problema del mondo reale e un problema di informatica.

\begin{figure}[H]
  \centering
  \begin{tikzpicture}[font=\Lato\scriptsize]
    % sequenza
    \node[font=\Lato\small\bfseries, text=primary!55!black] at (0,0.75) {sequenza};
    \node[fcproc] (a1) at (0,0) {passo};
    \node[fcproc] (a2) at (0,-0.9) {passo};
    \node[fcproc] (a3) at (0,-1.8) {passo};
    \draw[fcfreccia] (0,0.5) -- (a1); \draw[fcfreccia] (a1) -- (a2);
    \draw[fcfreccia] (a2) -- (a3); \draw[fcfreccia] (a3) -- (0,-2.5);
    \node[mnota, align=center, anchor=north] at (0,-2.75) {i passi vengono eseguiti\\uno dopo l'altro,\\incondizionatamente};
    % selezione
    \begin{scope}[xshift=6cm]
      \node[font=\Lato\small\bfseries, text=primary!55!black] at (0,0.75) {selezione};
      \node[fcdec] (c) at (0,-0.15) {condizione};
      \node[fcproc, minimum width=1.6cm] (v) at (-1.5,-1.2) {passo};
      \node[fcproc, minimum width=1.6cm] (f) at (1.5,-1.2) {passo};
      \node[fcgiunzione] (g) at (0,-2.0) {};
      \draw[fcfreccia] (0,0.5) -- (c);
      \draw[fcfreccia] (c.west) -| node[fcsi, above, pos=0.3] {vero} (v.north);
      \draw[fcfreccia] (c.east) -| node[fcno, above, pos=0.3] {falso} (f.north);
      \draw[fcfreccia] (v.south) |- (g.west);
      \draw[fcfreccia] (f.south) |- (g.east);
      \draw[fcfreccia] (g) -- (0,-2.5);
      \node[mnota, align=center, anchor=north] at (0,-2.75) {i passi vengono eseguiti\\solo dopo aver verificato\\una condizione};
    \end{scope}
    % iterazione
    \begin{scope}[xshift=12cm]
      \node[font=\Lato\small\bfseries, text=primary!55!black] at (0,0.75) {iterazione};
      \node[fcgiunzione] (j) at (0,0.3) {};
      \node[fcdec] (t) at (0,-0.5) {condizione};
      \node[fcproc] (b) at (0,-1.6) {passo};
      \draw[fcfreccia] (0,0.6) -- (j);
      \draw[fcfreccia] (j) -- (t);
      \draw[fcfreccia] (t) -- node[fcsi, right] {vero} (b);
      \draw[fcfreccia] (b.south) -- (0,-2.15) -- (1.7,-2.15) -- (1.7,0.3) -- (j.east);
      \draw[fcfreccia] (t.west) -- node[fcno, above] {falso} (-1.6,-0.5) -- (-1.6,-2.5);
      \node[mnota, align=center, anchor=north] at (0,-2.75) {i passi elementari\\vengono eseguiti\\più volte};
    \end{scope}
  \end{tikzpicture}
  \caption{Le tre strutture di controllo del flusso.}
  \label{fig:tre-strutture}
\end{figure}

\subsection{La sequenza}

La \term{sequenza} è la struttura più banale: i passi vengono eseguiti
\textbf{uno dopo l'altro}, incondizionatamente, nell'ordine in cui sono
scritti. Ma l'ordine conta: per alzarsi dal letto bisogna prima mettersi a
sedere, e per stampare un risultato bisogna prima averlo calcolato
(figura~\ref{fig:sequenza}).

\begin{figure}[H]
  \centering
  \begin{tikzpicture}[font=\Lato\scriptsize,
      tit/.style={font=\Lato\scriptsize\bfseries, text=neutral!25!black, align=center}]
    \node[tit] at (0,0.98) {Svegliarsi\\{\mdseries\itshape problema del mondo reale}};
    \foreach \t [count=\i from 0] in {apri gli occhi, stiracchiati, mettiti a sedere sul letto,
                                       infila le ciabatte, alzati}
      \node[fcproc, minimum width=3.0cm] (s\i) at (0,-\i*0.85) {\t};
    \foreach \i/\j in {0/1, 1/2, 2/3, 3/4} \draw[fcfreccia] (s\i) -- (s\j);
    \begin{scope}[xshift=8.4cm]
      \node[tit] at (0,0.98) {Calcolare l'area di un quadrato\\{\mdseries\itshape problema di informatica}};
      \node[fcterm] (a) at (0,0) {Inizio};
      \node[fcio, text width=3.6cm] (b) at (0,-0.95) {prendi in input il valore della variabile \code{lato}};
      \node[fcproc, text width=3.6cm] (c) at (0,-2.05) {calcola il prodotto di \code{lato} per \code{lato} e metti il valore nella variabile \code{area}};
      \node[fcio, text width=3.6cm] (d) at (0,-3.15) {metti in output il valore della variabile \code{area}};
      \node[fcterm] (e) at (0,-4.1) {Fine};
      \foreach \x/\y in {a/b, b/c, c/d, d/e} \draw[fcfreccia] (\x) -- (\y);
    \end{scope}
  \end{tikzpicture}
  \caption{Due sequenze.}
  \label{fig:sequenza}
\end{figure}

Nel quadrato compare per la prima volta la parola \term{variabile}: un
contenitore con un nome (\code{lato}, \code{area}) in cui si mette un
valore. È una cella di memoria (o un gruppo di celle) a cui diamo un nome,
così da non doverci occupare del suo indirizzo.

\subsection{La selezione a una via}

La \term{selezione} esegue dei passi \textbf{solo dopo aver verificato una
condizione}. Ne esistono tre tipi; il più semplice è la \textbf{selezione a
una via}: se la condizione è vera si eseguono dei passi, se è falsa non si
fa niente di particolare e si prosegue (figura~\ref{fig:una-via}).

\begin{figure}[H]
  \centering
  \begin{tikzpicture}[font=\Lato\scriptsize,
      tit/.style={font=\Lato\scriptsize\bfseries, text=neutral!25!black, align=center}]
    \node[tit] at (0,0.98) {Fare una telefonata\\{\mdseries\itshape problema del mondo reale}};
    \node[fcproc] (a) at (0,0) {componi il numero};
    \node[fcdec] (c) at (0,-1.1) {risponde?};
    \node[fcproc, minimum width=1.6cm] (p) at (-1.6,-2.15) {parla};
    \node[fcgiunzione] (g) at (0,-2.85) {};
    \node[fcproc] (r) at (0,-3.6) {riaggancia};
    \draw[fcfreccia] (a) -- (c);
    \draw[fcfreccia] (c.west) -| node[fcsi, above, pos=0.3] {vero} (p.north);
    \draw[fcfreccia] (p.south) |- (g.west);
    \draw[fcfreccia] (c.east) -- node[fcno, above] {falso} ++(0.55,0) |- (g.east);
    \draw[fcfreccia] (g) -- (r);
    \begin{scope}[xshift=8.4cm]
      \node[tit] at (0,0.98) {Simulare la roulette\\{\mdseries\itshape problema di informatica}};
      \node[fcproc, text width=4.2cm] (a) at (0,-0.05) {estrai un numero casuale fra 0 e 36 e assegnalo alla variabile \code{x}};
      \node[fcproc, text width=4.2cm] (b) at (0,-1.15) {assegna alla variabile \code{msg} il testo “È uscito il” seguito dal valore di \code{x}};
      \node[fcdec] (c) at (0,-2.3) {\code{x} vale 0?};
      \node[fcproc, text width=2.4cm] (d) at (-2.3,-3.3) {aggiungi a \code{msg} “Vince il banco”};
      \node[fcgiunzione] (g) at (0,-4.05) {};
      \node[fcio, text width=3.0cm] (e) at (0,-4.8) {metti in output la variabile \code{msg}};
      \draw[fcfreccia] (a) -- (b);
      \draw[fcfreccia] (b) -- (c);
      \draw[fcfreccia] (c.west) -| node[fcsi, above, pos=0.3] {vero} (d.north);
      \draw[fcfreccia] (d.south) |- (g.west);
      \draw[fcfreccia] (c.east) -- node[fcno, above] {falso} ++(0.55,0) |- (g.east);
      \draw[fcfreccia] (g) -- (e);
    \end{scope}
  \end{tikzpicture}
  \caption{Due selezioni a una via: il ramo “falso” non contiene passi.}
  \label{fig:una-via}
\end{figure}

Il pallino dove le strade si ricongiungono è importante: qualunque cosa
succeda alla domanda, dopo si riaggancia. La selezione ha una sola entrata
(la domanda) e una sola uscita (il pallino), proprio come una singola
istruzione.

\subsection{La selezione a due vie}

La \textbf{selezione a due vie} prevede delle operazioni da svolgere sia nel
caso in cui la condizione risulti \textbf{vera}, sia nel caso in cui risulti
\textbf{falsa}: viene eseguito sempre \textit{uno} dei due rami, mai
entrambi, mai nessuno (figura~\ref{fig:due-vie}).

\begin{figure}[H]
  \centering
  \begin{tikzpicture}[font=\Lato\scriptsize,
      tit/.style={font=\Lato\scriptsize\bfseries, text=neutral!25!black, align=center}]
    \node[tit] at (0,0.98) {Fare un viaggio in autobus\\{\mdseries\itshape problema del mondo reale}};
    \node[fcproc] (a) at (0,0) {sali sull'autobus};
    \node[fcproc] (b) at (0,-0.85) {oblitera};
    \node[fcdec] (c) at (0,-1.95) {c'è posto?};
    \node[fcproc, minimum width=2.4cm] (v) at (-1.85,-3.05) {fai il viaggio seduto};
    \node[fcproc, minimum width=2.4cm] (f) at (1.85,-3.05) {fai il viaggio in piedi};
    \node[fcgiunzione] (g) at (0,-3.8) {};
    \node[fcproc] (s) at (0,-4.55) {scendi dall'autobus};
    \draw[fcfreccia] (a) -- (b); \draw[fcfreccia] (b) -- (c);
    \draw[fcfreccia] (c.west) -| node[fcsi, above, pos=0.3] {vero} (v.north);
    \draw[fcfreccia] (c.east) -| node[fcno, above, pos=0.3] {falso} (f.north);
    \draw[fcfreccia] (v.south) |- (g.west);
    \draw[fcfreccia] (f.south) |- (g.east);
    \draw[fcfreccia] (g) -- (s);
    \begin{scope}[xshift=8.6cm]
      \node[tit] at (0,0.98) {Il segno di un numero casuale\\{\mdseries\itshape problema di informatica}};
      \node[fcproc, text width=4.2cm] (a) at (0,-0.05) {estrai un intero casuale fra $-50$ e $50$ e assegnalo alla variabile \code{x}};
      \node[fcdec] (c) at (0,-1.3) {\code{x < 0}?};
      \node[fcio, text width=2.5cm] (v) at (-1.95,-2.5) {metti in output il valore di \code{x} e “è negativo”};
      \node[fcio, text width=2.5cm] (f) at (1.95,-2.5) {metti in output il valore di \code{x} e “è positivo”};
      \node[fcgiunzione] (g) at (0,-3.45) {};
      \draw[fcfreccia] (a) -- (c);
      \draw[fcfreccia] (c.west) -| node[fcsi, above, pos=0.3] {vero} (v.north);
      \draw[fcfreccia] (c.east) -| node[fcno, above, pos=0.3] {falso} (f.north);
      \draw[fcfreccia] (v.south) |- (g.west);
      \draw[fcfreccia] (f.south) |- (g.east);
      \draw[fcfreccia] (g) -- (0,-4.0);
    \end{scope}
  \end{tikzpicture}
  \caption{Due selezioni a due vie.}
  \label{fig:due-vie}
\end{figure}

\warning{E se esce lo zero?}{
  L'algoritmo del segno, così com'è, ha un caso particolare: se il numero
  estratto è 0, la condizione \code{x < 0} è falsa e viene stampato “0 è
  positivo”. Lo zero però non è né positivo né negativo. Per trattarlo
  correttamente servono \textbf{tre} strade, come nell'esempio 3 del
  capitolo~\ref{chap:dal-problema}. È un buon esempio di come si verifica un
  algoritmo: si prova ogni ramo, e si cercano i \textbf{casi limite} (qui,
  il confine fra le due strade).}

\subsection{La selezione a n vie}

La \textbf{selezione a n vie} prevede operazioni diverse in \textbf{diversi
casi}, a seconda del valore di qualcosa: invece di una domanda con risposta
vero o falso, c'è una domanda con molte risposte possibili, ciascuna con la
sua strada (figura~\ref{fig:n-vie}).

\begin{figure}[H]
  \centering
  \begin{tikzpicture}[font=\Lato\scriptsize,
      tit/.style={font=\Lato\scriptsize\bfseries, text=neutral!25!black, align=center},
      caso/.style={fcproc, minimum width=2.25cm, text width=2.05cm, minimum height=0.9cm},
      val/.style={font=\Lato\tiny\bfseries, text=accent!65!black, fill=white, inner sep=1pt}]
    \node[tit] at (-5.0,0.1) {Decidere cosa indossare\\{\mdseries\itshape problema del mondo reale}};
    \node[fcproc] (a) at (0,0.35) {guarda fuori dalla finestra};
    \node[fcdec] (c) at (0,-0.75) {che tempo fa?};
    \foreach \x/\valore/\testo [count=\i] in {-3.9/30 gradi/esci a mezze maniche, -1.3/15 gradi/esci con la giacca a vento,
                                     1.3/0 gradi/esci con il piumino, 3.9/$-15$ gradi/non uscire} {
      \node[caso] (k\i) at (\x,-2.25) {\testo};
      \draw[fcfreccia] (c.south) -- node[val, pos=0.62] {\valore} (k\i.north);
      \draw[fcfreccia] (k\i.south) -- (0,-3.3);
    }
    \node[fcgiunzione] at (0,-3.3) {};
    \draw[fcfreccia] (a) -- (c);
    \draw[fcfreccia] (0,-3.38) -- (0,-3.8);
    \begin{scope}[yshift=-5.1cm]
      \node[tit] at (-5.6,0.1) {Estrarre un giorno\\lavorativo a caso\\{\mdseries\itshape problema di informatica}};
      \node[fcproc, text width=4.4cm] (a) at (0,0.35) {estrai un intero casuale fra 1 e 5 e assegnalo alla variabile \code{x}};
      \node[fcdec, minimum width=1.6cm, minimum height=0.8cm] (c) at (0,-0.85) {valore di \code{x}};
      \foreach \x/\valore/\testo [count=\i] in {-5.0/1/Lunedì, -2.5/2/Martedì, 0/3/Mercoledì, 2.5/4/Giovedì, 5.0/5/Venerdì} {
        \node[fcio, minimum width=1.9cm, text width=1.7cm, minimum height=0.9cm] (k\i) at (\x,-2.45) {metti in output “\testo”};
        \draw[fcfreccia] (c.south) -- node[val, pos=0.6] {\valore} (k\i.north);
        \draw[fcfreccia] (k\i.south) -- (0,-3.5);
      }
      \node[fcgiunzione] at (0,-3.5) {};
      \draw[fcfreccia] (a) -- (c);
      \draw[fcfreccia] (0,-3.58) -- (0,-4.0);
    \end{scope}
  \end{tikzpicture}
  \caption{Due selezioni a n vie: le etichette sulle frecce sono i valori
           possibili.}
  \label{fig:n-vie}
\end{figure}

Si potrebbe ottenere lo stesso risultato con una catena di selezioni a due
vie (“\code{x} vale 1? altrimenti, vale 2? altrimenti\dots”), ma la
selezione a n vie dice in modo più diretto che \textbf{una sola} domanda
decide fra molte strade. In C esistono entrambe le forme
(capitolo~\ref{chap:flusso}).

\subsection{L'iterazione}

Il concetto di \term{ciclo} (o \textit{iterazione}) è stato introdotto per
alleggerire situazioni come questa: dipingere le statuette dei sette nani.
Scritto come sequenza, l'algoritmo ripete sette volte lo stesso passo
(figura~\ref{fig:nani}, a sinistra). Una volta isolata l'istruzione da
eseguire \textbf{più volte}, si può scrivere un algoritmo che la esegue
semplicemente quante volte serve (a destra).

\begin{figure}[H]
  \centering
  \begin{tikzpicture}[font=\Lato\scriptsize,
      tit/.style={font=\Lato\scriptsize\bfseries, text=neutral!25!black, align=center}]
    \node[tit] at (0,0.98) {come sequenza};
    \foreach \n [count=\i from 0] in {Mammolo, Cucciolo, Dotto, Eolo, Brontolo, Gongolo, Pisolo}
      \node[fcproc, minimum width=3.4cm] (s\i) at (0,-\i*0.72) {dipingi la statuetta di \n};
    \foreach \i in {0,...,5} {\pgfmathtruncatemacro{\j}{\i+1}\draw[fcfreccia] (s\i) -- (s\j);}
    \draw[wtarrowacc, line width=1.4pt] (2.3,-2.2) -- node[wtlab, above=3pt, fill=none] {diventa} (4.6,-2.2);
    \begin{scope}[xshift=7.6cm]
      \node[tit] at (0,0.98) {come ciclo};
      \node[fcproc] (i0) at (0,0.1) {\code{i = 0}};
      \node[fcgiunzione] (j) at (0,-0.6) {};
      \node[fcdec] (t) at (0,-1.45) {\code{i < 7}?};
      \node[fcproc] (b1) at (0,-2.6) {dipingi una statuetta};
      \node[fcproc] (b2) at (0,-3.45) {incrementa \code{i}};
      \node[fcgiunzione] (u) at (0,-4.55) {};
      \draw[fcfreccia] (i0) -- (j);
      \draw[fcfreccia] (j) -- (t);
      \draw[fcfreccia] (t) -- node[fcsi, right] {vero} (b1);
      \draw[fcfreccia] (b1) -- (b2);
      \draw[fcfreccia, draw=accent!70!black] (b2.south) -- (0,-3.95) -- (1.85,-3.95) -- (1.85,-0.6) -- (j.east);
      \draw[fcfreccia] (t.west) -- node[fcno, above] {falso} (-1.75,-1.45) -- (-1.75,-4.55) -- (u.west);
      \draw[fcfreccia] (u) -- (0,-5.0);
    \end{scope}
  \end{tikzpicture}
  \caption{Dipingere le statuette dei sette nani: la ripetizione viene
           isolata in un ciclo.}
  \label{fig:nani}
\end{figure}

Per sapere quando fermarsi, il ciclo usa una variabile, \code{i}, che
\textbf{conta} quante volte il corpo del ciclo è già stato eseguito. Una
variabile di questo tipo si chiama \term{contatore} e per tradizione si
chiama \code{i}. Esistono però due grandi famiglie di cicli, e la difficoltà
sta nel capire in quale circostanza conviene usare l'una o l'altra:

\begin{itemize}
  \item il \textbf{ciclo su contatore}, quando si sa \textit{quante volte}
        ripetere;
  \item il \textbf{ciclo su condizione}, quando non lo si sa, ma si sa
        \textit{quando} continuare e quando smettere.
\end{itemize}

\subsection{Il ciclo su contatore}

Nel ciclo su contatore la variabile contatore viene
\textbf{inizializzata} prima di fare il ciclo, \textbf{testata} prima di
ripetere il ciclo e \textbf{modificata} ogni volta che il ciclo è stato
eseguito (figura~\ref{fig:contatore}, a sinistra). Dimenticare uno dei tre
passi è un errore classico: senza inizializzazione il conto parte da un
valore a caso, senza modifica il test resta sempre vero e il ciclo non
finisce mai.

\begin{figure}[H]
  \centering
  \begin{tikzpicture}[font=\Lato\scriptsize,
      tit/.style={font=\Lato\scriptsize\bfseries, text=neutral!25!black, align=center}]
    \node[tit] at (0,0.98) {lo schema generale};
    \node[fcproc, fill=green-gradient-1!40!white] (i0) at (0,0.1) {inizializzazione};
    \node[fcgiunzione] (j) at (0,-0.6) {};
    \node[fcdec, fill=yellow-gradient-1!60!white] (t) at (0,-1.45) {test};
    \node[fcproc] (b1) at (0,-2.6) {passi da ripetere};
    \node[fcproc, fill=accent!10!white] (b2) at (0,-3.45) {modifica};
    \node[fcgiunzione] (u) at (0,-4.55) {};
    \draw[fcfreccia] (i0) -- (j); \draw[fcfreccia] (j) -- (t);
    \draw[fcfreccia] (t) -- node[fcsi, right] {vero} (b1);
    \draw[fcfreccia] (b1) -- (b2);
    \draw[fcfreccia, draw=accent!70!black] (b2.south) -- (0,-3.95) -- (1.85,-3.95) -- (1.85,-0.6) -- (j.east);
    \draw[fcfreccia] (t.west) -- node[fcno, above] {falso} (-1.75,-1.45) -- (-1.75,-4.55) -- (u.west);
    \draw[fcfreccia] (u) -- (0,-5.0);
    \begin{scope}[xshift=7.6cm]
      \node[tit] at (0,0.98) {Stampare la tabellina del 3\\{\mdseries\itshape problema di informatica}};
      \node[fcproc] (i0) at (0,0.1) {\code{i = 0}};
      \node[fcgiunzione] (j) at (0,-0.6) {};
      \node[fcdec] (t) at (0,-1.45) {\code{i < 11}?};
      \node[fcio, text width=2.8cm] (b1) at (0,-2.6) {metti in output il prodotto di 3 per \code{i}};
      \node[fcproc] (b2) at (0,-3.6) {incrementa \code{i}};
      \node[fcgiunzione] (u) at (0,-4.55) {};
      \draw[fcfreccia] (i0) -- (j); \draw[fcfreccia] (j) -- (t);
      \draw[fcfreccia] (t) -- node[fcsi, right] {vero} (b1);
      \draw[fcfreccia] (b1) -- (b2);
      \draw[fcfreccia, draw=accent!70!black] (b2.south) -- (0,-4.05) -- (2.1,-4.05) -- (2.1,-0.6) -- (j.east);
      \draw[fcfreccia] (t.west) -- node[fcno, above] {falso} (-1.9,-1.45) -- (-1.9,-4.55) -- (u.west);
      \draw[fcfreccia] (u) -- (0,-5.0);
    \end{scope}
  \end{tikzpicture}
  \caption{Il ciclo su contatore: lo schema e un esempio.}
  \label{fig:contatore}
\end{figure}

Seguiamo l'esecuzione della tabellina, scrivendo a ogni giro il valore di
\code{i}, l'esito del test e l'output. È la tecnica più utile per capire un
ciclo (e per trovarne gli errori): \textbf{simularlo a mano}, una riga per
giro.

\begin{center}
\renewcommand{\arraystretch}{1.3}
\setlength{\tabcolsep}{5.2pt}
\begin{tabular}{@{}l*{12}{c}@{}}
  \toprule
  \textbf{valore di \code{i}} & 0 & 1 & 2 & 3 & 4 & 5 & 6 & 7 & 8 & 9 & 10 & 11 \\
  \textbf{\code{i < 11}?} & vero & vero & vero & vero & vero & vero & vero & vero & vero & vero & vero & \textbf{falso} \\
  \textbf{output} & 0 & 3 & 6 & 9 & 12 & 15 & 18 & 21 & 24 & 27 & 30 & --- \\
  \bottomrule
\end{tabular}
\end{center}

Il corpo viene eseguito 11 volte (per \code{i} da 0 a 10), e il test viene
valutato 12 volte: l'ultima, con \code{i} che vale 11, è quella che fa
uscire dal ciclo.

\subsection{Il ciclo su condizione}

Se non sappiamo a priori \textbf{quante} volte dobbiamo eseguire il ciclo,
sicuramente sappiamo \textbf{quando} eseguirlo e quando no: serve allora un
\term{ciclo su condizione}, che si ripete finché una condizione è vera.
Anche qui ci sono due casi, a seconda di quando si fa il test:

\begin{itemize}
  \item prima eseguiamo il ciclo e poi facciamo il test (\textbf{condizione
        in coda});
  \item oppure prima facciamo il test e poi eseguiamo il ciclo
        (\textbf{condizione in testa}).
\end{itemize}

La differenza sembra piccola, ma non lo è: con la condizione in coda il
corpo del ciclo viene eseguito \textbf{almeno una volta}; con la condizione
in testa può essere eseguito \textbf{zero volte}.

\subsection{Condizione in coda}

Si pensi a mangiare un piatto di spaghetti: si mangia una forchettata, poi
ci si chiede se ce n'è ancora; se sì, si mangia un'altra forchettata, e
così via. Chiedersi “ce n'è ancora?” ha senso solo \textbf{dopo} aver
cominciato (figura~\ref{fig:coda}, a sinistra).

\begin{figure}[H]
  \centering
  \begin{tikzpicture}[font=\Lato\scriptsize,
      tit/.style={font=\Lato\scriptsize\bfseries, text=neutral!25!black, align=center}]
    \node[tit] at (0,0.98) {Mangiare un piatto di spaghetti\\{\mdseries\itshape problema del mondo reale}};
    \node[fcgiunzione] (j) at (0,0.0) {};
    \node[fcproc] (b) at (0,-0.75) {mangia una forchettata};
    \node[fcdec] (t) at (0,-1.95) {ce n'è ancora?};
    \draw[fcfreccia] (0,0.35) -- (j);
    \draw[fcfreccia] (j) -- (b); \draw[fcfreccia] (b) -- (t);
    \draw[fcfreccia, draw=accent!70!black] (t.east) -- node[fcsi, above] {vero} (2.25,-1.95) -- (2.25,0) -- (j.east);
    \draw[fcfreccia] (t.south) -- node[fcno, right] {falso} (0,-3.0);
    \begin{scope}[xshift=8.0cm]
      \node[tit] at (0,0.98) {Tenere sotto controllo le calorie\\{\mdseries\itshape problema di informatica}};
      \node[fcio, text width=3.9cm] (a) at (0,0.05) {chiedi in input quante calorie si possono mangiare e acquisiscile nella variabile \code{tot}};
      \node[fcgiunzione] (j) at (0,-0.85) {};
      \node[fcio, text width=3.9cm] (b) at (0,-1.7) {chiedi in input le calorie della pietanza mangiata e acquisiscile nella variabile \code{x}};
      \node[fcproc] (c) at (0,-2.75) {decrementa \code{tot} di \code{x}};
      \node[fcdec] (t) at (0,-3.8) {\code{tot > 0}?};
      \draw[fcfreccia] (a) -- (j); \draw[fcfreccia] (j) -- (b);
      \draw[fcfreccia] (b) -- (c); \draw[fcfreccia] (c) -- (t);
      \draw[fcfreccia, draw=accent!70!black] (t.east) -- node[fcsi, above] {vero} (2.75,-3.8) -- (2.75,-0.85) -- (j.east);
      \draw[fcfreccia] (t.south) -- node[fcno, right] {falso} (0,-4.8);
    \end{scope}
  \end{tikzpicture}
  \caption{Cicli con la condizione in coda: il corpo viene eseguito almeno
           una volta.}
  \label{fig:coda}
\end{figure}

Il ciclo su condizione con verifica in coda si presta a due casi:

\begin{itemize}
  \item quando ha senso porsi la domanda \textbf{solo dopo} aver già
        eseguito una prima volta il ciclo, come per gli spaghetti e le
        calorie;
  \item quando c'è un obiettivo da raggiungere mediante uno o più
        \textbf{tentativi}: si tenta, e si ripete finché il tentativo non
        riesce (figura~\ref{fig:tentativi}).
\end{itemize}

\begin{figure}[H]
  \centering
  \begin{tikzpicture}[font=\Lato\scriptsize,
      tit/.style={font=\Lato\scriptsize\bfseries, text=neutral!25!black, align=center}]
    \node[tit] at (0,0.98) {Richiamare qualcuno che se ne va\\{\mdseries\itshape problema del mondo reale}};
    \node[fcgiunzione] (j) at (0,0.0) {};
    \node[fcproc, text width=2.8cm] (b) at (0,-0.8) {chiama a gran voce chi se ne sta andando via};
    \node[fcdec] (t) at (0,-2.05) {continua ad\\andare via?};
    \draw[fcfreccia] (0,0.35) -- (j);
    \draw[fcfreccia] (j) -- (b); \draw[fcfreccia] (b) -- (t);
    \draw[fcfreccia, draw=accent!70!black] (t.east) -- node[fcsi, above] {vero} (2.35,-2.05) -- (2.35,0) -- (j.east);
    \draw[fcfreccia] (t.south) -- node[fcno, right] {falso} (0,-3.2);
    \begin{scope}[xshift=8.0cm]
      \node[tit] at (0,0.98) {Input controllato\\{\mdseries\itshape problema di informatica}};
      \node[fcgiunzione] (j) at (0,0.0) {};
      \node[fcio, text width=3.9cm] (a) at (0,-0.8) {mostra in output un messaggio che chiede di inserire un intero di 3 cifre};
      \node[fcio, text width=3.9cm] (b) at (0,-1.85) {acquisisci l'input da tastiera nella variabile \code{x}};
      \node[fcdec] (t) at (0,-3.15) {\code{x < 100} oppure\\\code{x > 999}?};
      \draw[fcfreccia] (0,0.35) -- (j);
      \draw[fcfreccia] (j) -- (a); \draw[fcfreccia] (a) -- (b); \draw[fcfreccia] (b) -- (t);
      \draw[fcfreccia, draw=accent!70!black] (t.east) -- node[fcsi, above] {vero} (2.85,-3.15) -- (2.85,0) -- (j.east);
      \draw[fcfreccia] (t.south) -- node[fcno, right] {falso} (0,-4.35);
    \end{scope}
  \end{tikzpicture}
  \caption{Cicli con la condizione in coda per raggiungere un obiettivo con
           più tentativi.}
  \label{fig:tentativi}
\end{figure}

L'input controllato è un caso che ritroveremo spesso: un programma non può
fidarsi di ciò che scrive l'utente. Chiede un numero di tre cifre e, finché
il numero ricevuto è fuori dall'intervallo da 100 a 999, lo richiede.

\subsection{Condizione in testa}

Quando invece ha senso porsi la domanda \textbf{prima} di eseguire la
prima volta il ciclo, allora si pone la verifica della condizione in testa:
è il \textbf{ciclo su condizione con verifica in testa}. Per alzare tutte
le tapparelle di casa ci si chiede prima se ce ne sono ancora da alzare:
potrebbero essere già tutte aperte, e allora non si fa nulla
(figura~\ref{fig:testa}). Si noti che il ciclo su contatore visto prima è
un caso particolare di questo schema, con il test in testa.

\begin{figure}[H]
  \centering
  \begin{tikzpicture}[font=\Lato\scriptsize,
      tit/.style={font=\Lato\scriptsize\bfseries, text=neutral!25!black, align=center}]
    \node[tit] at (0,0.98) {Alzare tutte le tapparelle di casa\\{\mdseries\itshape problema del mondo reale}};
    \node[fcgiunzione] (j) at (0,0.0) {};
    \node[fcdec] (t) at (0,-0.95) {ci sono tapparelle\\da alzare?};
    \node[fcproc] (b) at (0,-2.25) {alza una tapparella};
    \node[fcgiunzione] (u) at (0,-3.25) {};
    \draw[fcfreccia] (0,0.35) -- (j); \draw[fcfreccia] (j) -- (t);
    \draw[fcfreccia] (t) -- node[fcsi, right] {vero} (b);
    \draw[fcfreccia, draw=accent!70!black] (b.south) -- (0,-2.75) -- (2.3,-2.75) -- (2.3,0) -- (j.east);
    \draw[fcfreccia] (t.west) -- node[fcno, above] {falso} (-2.2,-0.95) -- (-2.2,-3.25) -- (u.west);
    \draw[fcfreccia] (u) -- (0,-3.7);
    \begin{scope}[xshift=8.0cm]
      \node[tit] at (0,0.98) {Contare le cifre di un numero\\{\mdseries\itshape problema di informatica}};
      \node[fcio, text width=3.9cm] (a) at (0,0.05) {chiedi in input un intero positivo e acquisiscilo nella variabile \code{num}};
      \node[fcproc] (z) at (0,-0.95) {azzera la variabile \code{tot}};
      \node[fcgiunzione] (j) at (0,-1.6) {};
      \node[fcdec] (t) at (0,-2.45) {\code{num > 0}?};
      \node[fcproc, text width=3.6cm] (b1) at (0,-3.7) {calcola il quoziente di \code{num} diviso 10 e metti la parte intera in \code{num}};
      \node[fcproc] (b2) at (0,-4.75) {incrementa \code{tot} di 1};
      \node[fcgiunzione] (u) at (0,-5.7) {};
      \draw[fcfreccia] (a) -- (z); \draw[fcfreccia] (z) -- (j); \draw[fcfreccia] (j) -- (t);
      \draw[fcfreccia] (t) -- node[fcsi, right] {vero} (b1);
      \draw[fcfreccia] (b1) -- (b2);
      \draw[fcfreccia, draw=accent!70!black] (b2.south) -- (0,-5.2) -- (2.55,-5.2) -- (2.55,-1.6) -- (j.east);
      \draw[fcfreccia] (t.west) -- node[fcno, above] {falso} (-2.45,-2.45) -- (-2.45,-5.7) -- (u.west);
      \draw[fcfreccia] (u) -- (0,-6.15);
    \end{scope}
  \end{tikzpicture}
  \caption{Cicli con la condizione in testa: il corpo può non essere mai
           eseguito.}
  \label{fig:testa}
\end{figure}

Perché il secondo algoritmo conta le cifre? Dividere per 10 un numero
intero, tenendo solo la parte intera, \textbf{toglie l'ultima cifra}: da
4072 si passa a 407. Quindi le cifre sono tante quante le divisioni che
servono per arrivare a 0. Simuliamo con \code{num} che vale 4072:

\begin{center}
\renewcommand{\arraystretch}{1.3}
\begin{tabular}{@{}cccc@{}}
  \toprule
  \textbf{giro} & \textbf{\code{num} prima} & \textbf{\code{num > 0}?} & \textbf{dopo il giro} \\
  \midrule
  1 & 4072 & vero  & \code{num} = 407, \code{tot} = 1 \\
  2 & 407  & vero  & \code{num} = 40, \code{tot} = 2 \\
  3 & 40   & vero  & \code{num} = 4, \code{tot} = 3 \\
  4 & 4    & vero  & \code{num} = 0, \code{tot} = 4 \\
  5 & 0    & \textbf{falso} & si esce: le cifre sono \textbf{4} \\
  \bottomrule
\end{tabular}
\end{center}

L'algoritmo chiede un intero \textbf{positivo}, ed è importante: con
\code{num} uguale a 0 il test sarebbe subito falso e \code{tot} resterebbe
0, mentre il numero 0 ha una cifra. Ancora una volta, i casi limite.

\section{Combinare le strutture}

Le strutture di controllo possono essere \textbf{combinate} tra loro:
ovunque un diagramma contiene un passo, quel passo può essere sostituito da
un'intera struttura, perché ogni struttura ha una sola entrata e una sola
uscita, esattamente come un passo.

Riprendiamo la telefonata. In realtà, il passo “componi il numero” non è
esattamente un passo elementare, in quanto è scomponibile: si digita una
cifra, e si continua finché ci sono ancora cifre (un ciclo con la
condizione in coda). Sostituendo il passo con il suo ciclo si ottiene il
diagramma della figura~\ref{fig:combinazione}.

\begin{figure}[H]
  \centering
  \begin{tikzpicture}[font=\Lato\scriptsize,
      tit/.style={font=\Lato\scriptsize\bfseries, text=neutral!25!black, align=center}]
    % prima
    \node[tit] at (0,0.98) {prima};
    \node[fcproc, fill=yellow-gradient-1!70!white] (a) at (0,0) {componi il numero};
    \node[fcdec] (c) at (0,-1.1) {risponde?};
    \node[fcproc, minimum width=1.6cm] (p) at (-1.6,-2.15) {parla};
    \node[fcgiunzione] (g) at (0,-2.85) {};
    \node[fcproc] (r) at (0,-3.6) {riaggancia};
    \draw[fcfreccia] (a) -- (c);
    \draw[fcfreccia] (c.west) -| node[fcsi, above, pos=0.3] {vero} (p.north);
    \draw[fcfreccia] (p.south) |- (g.west);
    \draw[fcfreccia] (c.east) -- node[fcno, above] {falso} ++(0.55,0) |- (g.east);
    \draw[fcfreccia] (g) -- (r);
    \draw[wtarrowacc, line width=1.4pt] (2.6,-1.6) -- node[wtlab, above=3pt, fill=none] {si espande} (4.6,-1.6);
    % dopo
    \begin{scope}[xshift=8.2cm, yshift=0.6cm]
      \node[tit] at (0,0.15) {dopo};
      \draw[yellow-gradient-6!90!black, dashed, line width=0.8pt, fill=yellow-gradient-1!35!white, rounded corners=4pt]
        (-2.6,-0.2) rectangle (2.6,-2.75);
      \node[font=\Lato\tiny\bfseries, anchor=north west, text=yellow-gradient-6!40!black] at (-2.55,-0.25) {componi il numero};
      \node[fcgiunzione] (j) at (0,-0.5) {};
      \node[fcproc] (d) at (0,-1.15) {digita una cifra};
      \node[fcdec] (t) at (0,-2.15) {ci sono ancora\\cifre?};
      \draw[fcfreccia] (0,0.0) -- (j); \draw[fcfreccia] (j) -- (d); \draw[fcfreccia] (d) -- (t);
      \draw[fcfreccia, draw=accent!70!black] (t.east) -- node[fcsi, above] {vero} (2.25,-2.15) -- (2.25,-0.5) -- (j.east);
      \node[fcdec] (c) at (0,-3.55) {risponde?};
      \node[fcproc, minimum width=1.6cm] (p) at (-1.6,-4.6) {parla};
      \node[fcgiunzione] (g) at (0,-5.3) {};
      \node[fcproc] (r) at (0,-6.05) {riaggancia};
      \draw[fcfreccia] (t.south) -- node[fcno, right] {falso} (c);
      \draw[fcfreccia] (c.west) -| node[fcsi, above, pos=0.3] {vero} (p.north);
      \draw[fcfreccia] (p.south) |- (g.west);
      \draw[fcfreccia] (c.east) -- node[fcno, above] {falso} ++(0.55,0) |- (g.east);
      \draw[fcfreccia] (g) -- (r);
    \end{scope}
  \end{tikzpicture}
  \caption{Il passo “componi il numero” sostituito dal ciclo che lo
           realizza: una selezione che segue un ciclo.}
  \label{fig:combinazione}
\end{figure}

Questo principio, sostituire un passo con la struttura che lo realizza, è
alla base dell'approccio \textbf{top-down}.

\section{Lo sviluppo top-down}

\subsection{Dividere il problema}

Per \term{sviluppo top-down} (“dall'alto verso il basso”) si intende un
approccio alla soluzione di un problema che prevede la sua
\textbf{scomposizione in sotto-problemi}, più piccoli e quindi più semplici
da risolvere. Ogni sotto-problema, a sua volta, si può scomporre ancora,
finché non si arriva a passi abbastanza semplici da scrivere direttamente.

Ogni sotto-problema sarà risolto con la realizzazione di un
\term{sotto-programma} (in C, una \textit{funzione}), e uno dei passi
elementari del nostro algoritmo potrebbe essere proprio l'esecuzione di
uno dei sotto-programmi. Lo svolgimento di un compito ben preciso (anche
piccolo) può avvenire in autonomia, ma spesso il sotto-programma ha bisogno
di scambiare dei dati con chi lo usa. Per esempio, “componi il numero”
richiede che gli si fornisca il numero da comporre.

\begin{itemize}
  \item I dati che un sotto-programma riceve da chi lo invoca si chiamano
        \term{parametri}.
  \item Il dato che un sotto-programma fornisce in output a chi lo ha
        invocato si chiama \term{valore restituito}.
\end{itemize}

Un'analogia: in un libro di cucina, la ricetta delle lasagne dice
“prepara un litro di besciamella (vedi pagina 112)”. La ricetta della
besciamella è un sotto-programma; “un litro” è il parametro, e la
besciamella pronta è il valore restituito. La ricetta delle lasagne non deve
sapere come si fa la besciamella: le basta sapere che cosa darle e che cosa
riceverà.

\subsection{Come si invoca un sotto-programma}

Supponiamo di avere un programma (che questo sia il programma principale o
un sotto-programma non ha alcuna importanza) e un sotto-programma qualsiasi.
Quando progettiamo il sotto-programma dobbiamo, fra l'altro, indicare i
parametri di cui necessita: sono i \term{parametri formali}, i nomi con cui
il sotto-programma si riferisce ai dati ricevuti. Supponiamo che siano
\code{a}, \code{b}, \code{c} e \code{d}.

L'\textbf{invocazione} (da parte del programma chiamante) avviene indicando
il nome del sotto-programma, fornendo i parametri con cui lavorare (i
\term{parametri attuali}) e indicando che cosa fare del valore restituito.
La figura~\ref{fig:sottoprogramma} segue l'invocazione passo per passo.

\begin{figure}[H]
  \centering
  \begin{tikzpicture}[font=\Lato\scriptsize,
      passo/.style={font=\Lato\tiny, anchor=west, inner sep=1pt, text=neutral!35!black},
      par/.style={draw=neutral!55!black, fill=neutral!25!white, minimum size=0.36cm,
        inner sep=0pt, font=\FiraCodeNL\tiny\bfseries},
      num/.style={circle, fill=accent!75!black, text=white, font=\Lato\scriptsize\bfseries,
        inner sep=0pt, minimum size=0.4cm}]
    % programma chiamante
    \node[font=\Lato\scriptsize\bfseries, anchor=south west] at (0,0.05) {programma chiamante};
    \draw[neutral!50!black, fill=white, line width=0.7pt] (0,0) rectangle (4.4,-6.2);
    \foreach \k [count=\i] in {1,2,3} \node[passo] at (0.1,-\i*0.32) {\k.\ \ passo elementare};
    \node[passo] at (0.1,-1.25) {\dots};
    \node[passo] at (0.1,-1.6) {138.\ \ passo elementare};
    \node[draw=accent!70!black, fill=accent!8!white, line width=0.8pt, text width=4.0cm,
          font=\Lato\tiny, align=left, anchor=north west] (inv) at (0.1,-1.85)
      {139.\ \ invoca il sotto-programma qui a lato con i parametri \code{n1}, \code{n2},
       \code{n3}, \code{n4} e metti il valore restituito nella variabile \code{x}};
    \node[passo] (p140) at (0.1,-3.2) {140.\ \ passo elementare};
    \node[passo] at (0.1,-3.55) {\dots};
    \node[passo] at (0.1,-5.9) {1865.\ \ passo elementare};
    \foreach \n [count=\i] in {n1,n2,n3,n4}
      \node[par] (q\i) at (4.75,-1.85-\i*0.38) {\n};
    \node[par, fill=accent!25!white] (qx) at (4.75,-3.85) {x};
    % sotto-programma
    \node[font=\Lato\scriptsize\bfseries, anchor=south west] at (10.6,-0.95) {sotto-programma};
    \draw[neutral!50!black, fill=white, line width=0.7pt] (10.6,-1.0) rectangle (14.6,-5.4);
    \foreach \k [count=\i] in {1,2,3} \node[passo] at (10.7,-1.0-\i*0.32) {\k.\ \ passo elementare};
    \node[passo] at (10.7,-2.3) {\dots};
    \node[passo] at (10.7,-4.75) {130.\ \ passo elementare};
    \node[passo] at (10.7,-5.1) {131.\ \ restituisce il risultato};
    \foreach \n [count=\i] in {a,b,c,d}
      \node[par, fill=primary!25!white] (f\i) at (10.2,-1.0-\i*0.38) {\n};
    % frecce
    \foreach \i in {1,...,4}
      \draw[wtarrow, draw=dbgblu] (q\i.east) -- (f\i.west);
    \draw[wtarrow, draw=red-gradient-6!75!black] (10.6,-5.1) -- (qx.east);
    \node[num] at (7.4,-1.35) {1};
    \node[num] at (12.6,-3.5) {2};
    \node[num] at (7.6,-4.75) {3};
    \draw[wtarrow, draw=green-gradient-6!70!black, dashed] (inv.south) ++(-1.0,0) -- ([xshift=0.9cm]p140.north west);
    \node[num] at (-0.45,-3.0) {4};
  \end{tikzpicture}
  \caption{L'invocazione di un sotto-programma con quattro parametri.}
  \label{fig:sottoprogramma}
\end{figure}

\begin{passi}
  \item All'invocazione, il controllo del flusso passa al sotto-programma,
        e ogni parametro formale riceve una \textbf{copia} del valore del
        parametro attuale corrispondente: in \code{a} viene copiato il
        valore di \code{n1}, in \code{b} quello di \code{n2}, e così via.
  \item Il sotto-programma esegue i suoi passi, lavorando con \code{a},
        \code{b}, \code{c} e \code{d}.
  \item Quando ha eseguito la sua ultima istruzione, fornisce al programma
        chiamante il suo \textbf{valore di restituzione}, che finisce in
        \code{x}\dots
  \item \dots e l'esecuzione riprende dall'istruzione successiva a quella
        dell'invocazione (la 140) del programma chiamante.
\end{passi}

Due osservazioni importanti. I nomi dei parametri attuali e formali
\textbf{possono anche non coincidere} (qui \code{n1} e \code{a}): il
sotto-programma è scritto una volta sola e deve funzionare con qualunque
dato gli venga passato, quindi usa nomi suoi. Invece devono
\textbf{assolutamente coincidere} il numero, il \textbf{tipo} di ogni
parametro e l'\textbf{ordine}: il primo parametro attuale va nel primo
formale, il secondo nel secondo. Scambiare l'ordine non produce un errore
visibile, produce un risultato sbagliato.

\section{I tipi di istruzione}

Riassumendo, ogni passo elementare di un algoritmo sarà un'istruzione di uno
di questi tipi: \textbf{input}, \textbf{output}, \textbf{assegnazione}
(mettere un valore in una variabile), \textbf{selezione},
\textbf{iterazione}, \textbf{invocazione di funzione}. Selezione e
iterazione si dividono a loro volta come nella figura~\ref{fig:istruzioni}.

\begin{figure}[H]
  \centering
  \begin{tikzpicture}[font=\Lato\scriptsize,
      r/.style={wtnode, font=\Lato\scriptsize\bfseries, minimum width=2.6cm},
      f/.style={wtnodeplain, minimum width=1.8cm}]
    \node[r] (s) at (0,0) {istruzioni di selezione};
    \node[f] (s1) at (-2.2,-1.2) {a 1 via};
    \node[f] (s2) at (0,-1.2) {a 2 vie};
    \node[f] (s3) at (2.2,-1.2) {a n vie};
    \foreach \x in {s1,s2,s3} \draw[fsramo] (s.south) -- ++(0,-0.3) -| (\x.north);
    \begin{scope}[xshift=8.6cm]
      \node[r] (i) at (0,0) {istruzioni di iterazione};
      \node[f] (i1) at (-1.6,-1.2) {su contatore};
      \node[f] (i2) at (1.6,-1.2) {su condizione};
      \node[f] (i3) at (0.6,-2.4) {in testa};
      \node[f] (i4) at (2.6,-2.4) {in coda};
      \foreach \x in {i1,i2} \draw[fsramo] (i.south) -- ++(0,-0.3) -| (\x.north);
      \foreach \x in {i3,i4} \draw[fsramo] (i2.south) -- ++(0,-0.3) -| (\x.north);
    \end{scope}
  \end{tikzpicture}
  \caption{I tipi di istruzioni di selezione e di iterazione.}
  \label{fig:istruzioni}
\end{figure}

Ognuno di questi concetti diventerà un costrutto preciso del C. La tabella
seguente anticipa la corrispondenza, per avere fin d'ora una mappa del
percorso: non serve ricordarla, la ritroveremo pezzo per pezzo.

\begin{center}
\renewcommand{\arraystretch}{1.35}
\begin{tabular}{@{}ll>{\FiraCodeNL\small}lc@{}}
  \toprule
  \textbf{istruzione} & & \normalfont\textbf{in C} & \textbf{capitolo} \\
  \midrule
  input                   & & scanf, getchar & \ref{chap:primi-programmi}, \ref{chap:tipi} \\
  output                  & & printf, putchar & \ref{chap:primi-programmi}, \ref{chap:tipi} \\
  assegnazione            & & x = valore; & \ref{chap:primi-programmi} \\
  selezione               & a 1 via & if & \ref{chap:primi-programmi}, \ref{chap:flusso} \\
                          & a 2 vie & if \dots\ else & \ref{chap:primi-programmi}, \ref{chap:flusso} \\
                          & a n vie & switch & \ref{chap:flusso} \\
  iterazione              & su contatore & for & \ref{chap:primi-programmi}, \ref{chap:flusso} \\
                          & su condizione, in testa & while & \ref{chap:primi-programmi}, \ref{chap:flusso} \\
                          & su condizione, in coda & do \dots\ while & \ref{chap:flusso} \\
  invocazione di funzione & & nome(parametri) & \ref{chap:funzioni} \\
  \bottomrule
\end{tabular}
\end{center}

\insintesi{
  Il C nasce nel 1972 per scrivere Unix, come evoluzione tipizzata del B;
  lo standard (ANSI e ISO, dal 1989--90) ne garantisce portabilità,
  affidabilità e manutenibilità. La programmazione strutturata rinuncia ai
  salti arbitrari: per il teorema di Böhm-Jacopini bastano
  \textbf{sequenza}, \textbf{selezione} (a una, due o n vie) e
  \textbf{iterazione} (su contatore, o su condizione con il test in testa o
  in coda). Ogni struttura ha una sola entrata e una sola uscita, quindi le
  strutture si annidano come scatole; sostituire un passo con la struttura
  che lo realizza è l'idea dello sviluppo \textbf{top-down}, in cui i
  sotto-problemi diventano sotto-programmi che ricevono \textbf{parametri}
  (copiati, nell'ordine) e restituiscono un \textbf{valore}.}

\section{Riferimenti}

\begin{itemize}
  \item \textbf{Kernighan \& Ritchie}, \textit{Il linguaggio C}: prefazione e
        introduzione (storia del linguaggio).
  \item \textbf{E. W. Dijkstra}, \textit{Go To Statement Considered Harmful},
        Communications of the ACM, 1968.
  \item \textbf{C. Böhm, G. Jacopini}, \textit{Flow diagrams, Turing machines
        and languages with only two formation rules}, Communications of the
        ACM, 1966.
\end{itemize}


  \chapter{Primi programmi in C}
\label{chap:primi-programmi}

Questo capitolo è una panoramica del C attraverso una serie di programmi di
esempio, ciascuno spiegato elemento per elemento. Si introducono
brevemente i costrutti principali (l'output, le variabili, l'input, le
scelte, i cicli, le funzioni) per avere subito un'idea di come è fatto un
programma vero; ogni argomento verrà poi approfondito nei capitoli
successivi. Il consiglio è uno solo: \textbf{scrivere ed eseguire ogni
esempio}, poi modificarlo e osservare che cosa cambia.

\section{Il primo programma, riga per riga}

\subsection{Scrivere, compilare, eseguire}

Un programma deve comunicare per essere utile: produce sempre un output di
qualche tipo. Il primo esempio stampa una frase sullo schermo. Lo scriviamo
nel file \code{sea.c}, lo compiliamo e lo eseguiamo, come abbiamo imparato
nel capitolo~\ref{chap:ambiente}:

\begin{affianco}[0.5]
\begin{lstlisting}[style=c, name=sea.c]
#include <stdio.h>

int main(void)
{
    printf("from sea to shining C\n");
    return 0;
}
\end{lstlisting}
\risultato
\begin{console}
$ gcc sea.c
$ ./a.out
from sea to shining C
\end{console}
\wtnota{Se non ci sono errori, gcc crea il file eseguibile \code{a.out}.}
\end{affianco}

Sono sette righe, ma ognuna ha un ruolo preciso. La
figura~\ref{fig:anatomia} le numera, e l'elenco che la segue le spiega.

\begin{figure}[H]
  \centering
  \begin{tikzpicture}[font=\FiraCodeNL\small,
      riga/.style={anchor=base west, inner sep=0pt},
      num/.style={circle, fill=accent!75!black, text=white, font=\Lato\scriptsize\bfseries,
        inner sep=0pt, minimum size=0.42cm}]
    \fill[vscbg, rounded corners=4pt] (-0.3,0.55) rectangle (9.2,-3.4);
    \node[riga] (r1) at (0,0) {\textcolor{vscpurple}{\#include} \textcolor{vscstring}{<stdio.h>}};
    \node[riga] (r2) at (0,-0.6) {\textcolor{vscblue}{int} \textcolor{vscfunc}{main}\textcolor{vscfg}{(}\textcolor{vscblue}{void}\textcolor{vscfg}{)}};
    \node[riga, text=vscfg] (r3) at (0,-1.2) {\{};
    \node[riga] (r4) at (0,-1.8) {\textcolor{vscfg}{\ \ \ \ }\textcolor{vscfunc}{printf}\textcolor{vscfg}{(}\textcolor{vscstring}{"from sea to shining C\textbackslash n"}\textcolor{vscfg}{);}};
    \node[riga] (r5) at (0,-2.4) {\textcolor{vscfg}{\ \ \ \ }\textcolor{vscpurple}{return} \textcolor{vscfg}{0;}};
    \node[riga, text=vscfg] (r6) at (0,-3.0) {\}};
    \foreach \r/\n in {r1/1, r2/2, r3/3, r4/4, r5/5, r6/3}
      \node[num] at ([yshift=3.5pt]8.75,0 |- \r.base) {\n};
  \end{tikzpicture}
  \caption{Anatomia del primo programma: i numeri rimandano all'elenco
           seguente.}
  \label{fig:anatomia}
\end{figure}

\begin{passi}
  \item \textbf{Direttiva per il preprocessore}: copia qui il contenuto del
        file di intestazione \code{stdio.h}, che descrive \code{printf}.
  \item \textbf{Testa della funzione} \code{main}, da cui comincia
        l'esecuzione: restituisce un \code{int} e non riceve parametri
        (\code{void}).
  \item Le \textbf{graffe} delimitano il \textbf{corpo} della funzione.
  \item Un'\textbf{istruzione}: chiama \code{printf} con un argomento, una
        \textbf{stringa costante}; \code{\textbackslash n} va a capo; il
        \code{;} chiude l'istruzione.
  \item Restituisce 0 al sistema operativo: “tutto bene”.
\end{passi}

\subsection{Il preprocessore e \code{\#include}}

Prima della compilazione vera e propria viene eseguito, automaticamente, il
\term{preprocessore} (capitolo~\ref{chap:dal-problema}). Le righe che
cominciano con \code{\#} sono \textbf{direttive} per il preprocessore. La
direttiva \code{\#include <stdio.h>} chiede di includere in quel punto una
copia del file di intestazione (\textit{header file}) \code{stdio.h}, che
contiene le informazioni necessarie per usare le funzioni di input e output,
come \code{printf}.

\warning{\code{\#include} non include la libreria}{
  A lezione il docente è categorico: guai a dire che \code{\#include}
  “include la libreria”. Include un \textbf{file di testo}, l'header, che
  \textbf{descrive} le funzioni (come si chiamano, che parametri vogliono,
  che cosa restituiscono). Il \textbf{codice} delle funzioni sta nella
  libreria, che viene collegata dal \textbf{linker}, dopo. Sono due cose
  diverse, che intervengono in due fasi diverse: lo vedremo in modo molto
  concreto con la funzione \code{pow} più avanti in questo capitolo.}

\subsection{La funzione \code{main}}

\code{int main(void)} è la prima riga della definizione di una
\term{funzione} chiamata \code{main}. Le parentesi dopo il nome dicono che
\code{main} è una funzione; \code{int} e \code{void} sono parole speciali
che danno informazioni al compilatore: \code{void} fra le parentesi
significa che la funzione \textbf{non ha parametri}, \code{int} davanti al
nome significa che \textbf{restituisce un intero}.

Ogni programma C ha \textbf{una} funzione \code{main}, e l'esecuzione
comincia da lei. Il corpo della funzione è racchiuso fra le graffe
\code{\{} e \code{\}}, che in C servono in generale a raggruppare
istruzioni.

\subsection{\code{printf}, le stringhe e il punto e virgola}

\code{printf} è una funzione che stampa sullo schermo (più precisamente
sullo \textit{standard output}, come vedremo). Non fa parte del
linguaggio: appartiene alla \term{libreria standard}, un insieme di
librerie che forniscono al programmatore funzioni pronte; gli header file
permettono di usarle, e \code{stdio.h} fornisce le informazioni per
\code{printf}.

Il testo fra virgolette, \code{"from sea to shining C\textbackslash n"}, è
una \term{stringa costante}: è l'\textbf{argomento} di \code{printf} e
stabilisce che cosa viene stampato. Al suo interno, \code{\textbackslash n}
(una barra rovesciata seguita da \code{n}) è un \textbf{singolo carattere},
il \textit{newline}: fa avanzare il cursore all'inizio della riga
successiva. L'intera riga è un'\term{istruzione}, e in C le istruzioni
\textbf{terminano con un punto e virgola}.

\subsection{\code{return} e il valore restituito al sistema}

\code{return 0;} termina la funzione \code{main} e restituisce il valore 0
al sistema operativo. Per convenzione 0 significa “tutto è andato bene”,
un valore diverso da 0 segnala un problema. Il valore non si vede
nell'output, ma la shell lo conserva, e lo si può leggere con il comando
\code{echo \$?} subito dopo l'esecuzione:

\begin{console}
$ ./a.out
from sea to shining C
$ echo $?
0
\end{console}

\code{return} \textbf{termina} la funzione: ciò che viene dopo non è
eseguito. Il file \code{sea.c} distribuito con le lezioni ne dà una prova,
perché dopo il \code{return} contiene un'altra \code{printf}:

\begin{affianco}[0.5]
\begin{lstlisting}[style=c, name=sea.c]
#include <stdio.h>

int main()
{
    printf("from sea to shining C\n");

    return 0;
    printf("Pippo\n");
}
\end{lstlisting}
\risultato
\begin{console}
$ gcc -o sea sea.c
$ ./sea
from sea to shining C
\end{console}
\wtnota{“Pippo” non viene mai stampato: l'esecuzione esce da \code{main}
al \code{return}.}
\end{affianco}

(Qui \code{main} è scritto con le parentesi vuote, \code{main()}: per una
definizione di funzione è equivalente a \code{main(void)}; in questa
dispensa useremo la forma con \code{void}, più esplicita.)

\section{L'output con \code{printf}}

\subsection{Il cursore}

\code{printf} scrive sullo schermo da sinistra a destra e dall'alto verso
il basso, come si scrive su un quaderno; va a capo \textbf{solo} quando
incontra il carattere \code{\textbackslash n}. Ogni chiamata di
\code{printf} comincia da dove la precedente ha lasciato il cursore. Per
questo il programma seguente, diverso dal primo, stampa \textbf{lo stesso}
output:

\begin{affianco}[0.5][c]
\begin{lstlisting}[style=c]
printf("from sea to ");
printf("shining ");
printf("C\n");
\end{lstlisting}
\risultato
\begin{console}
from sea to shining C
\end{console}
\end{affianco}

Se vogliamo scrivere su tre righe dobbiamo inserire dei \code{\textbackslash
n}; e se ne possono mettere più di uno nella stessa stringa:

\begin{affianco}[0.5][c]
\begin{lstlisting}[style=c]
printf("from sea to\n");
printf("shining\nC\n");
\end{lstlisting}
\risultato
\begin{console}
from sea to
shining
C
\end{console}
\end{affianco}

\subsection{Disegnare con i caratteri}

Per essere leggibile, l'output va spaziato opportunamente sullo schermo.
Il programma seguente stampa la stessa frase in una cornice di asterischi,
preceduta e seguita da dieci righe vuote (ciascun \code{\textbackslash n}
iniziale produce una riga vuota):

\begin{affianco}[0.58][c]
\begin{lstlisting}[style=c, name=sea_box.c]
#include <stdio.h>

int main(void)
{
    printf("\n\n\n\n\n\n\n\n\n\n");
    printf("      ***********************\n");
    printf("      *  from sea           *\n");
    printf("      *  to shining C       *\n");
    printf("      ***********************\n");
    printf("\n\n\n\n\n\n\n\n\n\n");
    return 0;
}
\end{lstlisting}
\risultato
\begin{console}
$ ./sea_box


      ***********************
      *  from sea           *
      *  to shining C       *
      ***********************


$
\end{console}
\wtnota{Qui sono mostrate solo due delle dieci righe vuote sopra e sotto.}
\end{affianco}

\section{Variabili, espressioni e assegnazione}

\subsection{Il problema della maratona}

Il secondo esempio converte una distanza: la lunghezza di una maratona è di
26 miglia e 385 iarde, e vogliamo sapere quanti chilometri sono. Servono due
regole: in un miglio ci sono 1760 iarde, e un miglio è lungo 1,609 km.

\begin{affianco}[0.56]
\begin{lstlisting}[style=c, name=marathon.c]
/* The distance of a marathon in kilometers. */

#include <stdio.h>

int main(void)
{
    int     miles, yards;
    float   kilometers;

    miles = 26;
    yards = 385;
    kilometers = 1.609 * (miles + yards / 1760.0);
    printf("\nA marathon is %f kilometers.\n\n", kilometers);
    return 0;
}
\end{lstlisting}
\risultato
\begin{console}
$ gcc -o marathon marathon.c
$ ./marathon

A marathon is 42.185970 kilometers.

\end{console}
\wtnota{L'opzione \code{-o} sceglie il nome del file eseguibile.}
\end{affianco}

\subsection{I commenti}

Tutto ciò che è compreso fra \code{/*} e \code{*/} è un
\term{commento}: viene ignorato dal compilatore (anzi, eliminato già dal
preprocessore) e serve solo a chi legge. Si possono anche scrivere commenti
su una sola riga, che cominciano con \code{//} e arrivano fino alla fine
della riga. Un commento all'inizio del file spiega lo scopo del programma;
vi si possono aggiungere informazioni sugli autori e sulle versioni.

\subsection{Le dichiarazioni}

\begin{lstlisting}[style=c]
int     miles, yards;
float   kilometers;
\end{lstlisting}

Queste righe sono \term{dichiarazioni}. \code{int} e \code{float} sono
\textbf{parole chiave} del C che indicano due \textbf{tipi di dato
fondamentali}: la prima riga dice al compilatore che \code{miles} e
\code{yards} sono \term{variabili} di tipo intero, che assumono valori
interi; la seconda che \code{kilometers} è una variabile di tipo
\code{float}, che assume valori reali. Anche le dichiarazioni terminano con
un punto e virgola.

Dichiarare una variabile significa chiedere al compilatore di riservarle
dello spazio in memoria e di ricordarne il tipo. In C \textbf{ogni
variabile va dichiarata prima di essere usata}.

\subsection{L'assegnazione}

\code{miles = 26;} è un'\term{istruzione di assegnazione}: il simbolo
\code{=} è l'\textbf{operatore di assegnazione}, \code{26} è una costante
intera, e l'effetto è che il valore 26 viene scritto nella variabile
\code{miles} (figura~\ref{fig:maratona-memoria}). Allo stesso modo 385 viene
assegnato a \code{yards}.

\warning{\code{=} non è l'uguale della matematica}{
  In matematica $x = x + 1$ è un'equazione impossibile. In C
  \code{x = x + 1;} è perfettamente sensato e significa: \textbf{calcola} il
  valore dell'espressione a destra (il valore attuale di \code{x} più 1) e
  \textbf{scrivilo} nella variabile a sinistra. L'assegnazione va letta da
  destra a sinistra: “\code{x} riceve\dots”.}

\begin{figure}[H]
  \centering
  \begin{tikzpicture}[font=\Lato\scriptsize]
    \foreach \titolo/\m/\y/\km [count=\i from 0] in {
        {dopo le dichiarazioni}/\mvuota/\mvuota/\mvuota,
        {dopo \code{miles = 26;}}/26/\mvuota/\mvuota,
        {dopo \code{yards = 385;}}/26/385/\mvuota,
        {dopo il calcolo}/26/385/42.18597} {
      \begin{scope}[xshift=\i*4.4cm]
        \node[font=\Lato\scriptsize\bfseries] at (0.3,0.6) {\titolo};
        \variabile{m\i}{(0.6,0)}{miles}{\m}
        \variabile{y\i}{(0.6,-0.65)}{yards}{\y}
        \variabile[minimum width=1.4cm]{k\i}{(0.6,-1.3)}{kilometers}{\km}
      \end{scope}
    }
  \end{tikzpicture}
  \caption{Le variabili della maratona, istruzione dopo istruzione. Appena
           dichiarata, una variabile contiene un valore a caso (“?”):
           finché non le si assegna qualcosa, non va usata.}
  \label{fig:maratona-memoria}
\end{figure}

\subsection{Come si calcola un'espressione}

\code{kilometers = 1.609 * (miles + yards / 1760.0);} assegna un valore
calcolato da un'\term{espressione}. Le regole sono quelle dell'aritmetica:
si calcola prima ciò che sta fra parentesi, e dentro le parentesi la
divisione ha la precedenza sulla somma. Quindi prima \code{yards / 1760.0},
poi il risultato viene sommato a \code{miles}, e infine moltiplicato per
\code{1.609} (figura~\ref{fig:albero-maratona}).

\begin{figure}[H]
  \centering
  \begin{tikzpicture}[font=\FiraCodeNL\small,
      op/.style={circle, draw=accent!65!black, fill=accent!8!white, line width=0.8pt,
        minimum size=0.62cm, inner sep=0pt, font=\FiraCodeNL\small\bfseries},
      fo/.style={wtnodeplain, font=\FiraCodeNL\scriptsize, minimum width=1.2cm},
      val/.style={font=\Lato\scriptsize\bfseries, text=green-gradient-6!50!black, anchor=west},
      ord/.style={circle, fill=primary!75!black, text=white, font=\Lato\tiny\bfseries,
        inner sep=0pt, minimum size=0.36cm}]
    \node[op] (mul) at (0,0) {*};
    \node[fo] (k) at (-1.6,-1.1) {1.609};
    \node[op] (add) at (1.6,-1.1) {+};
    \node[fo] (m) at (0.4,-2.2) {miles\\[-1pt]{\color{neutral!50!black}(26)}};
    \node[op] (div) at (2.8,-2.2) {/};
    \node[fo] (y) at (1.8,-3.35) {yards\\[-1pt]{\color{neutral!50!black}(385)}};
    \node[fo] (c) at (3.8,-3.35) {1760.0};
    \foreach \a/\b in {mul/k, mul/add, add/m, add/div, div/y, div/c} \draw[fsramo] (\a) -- (\b);
    \node[val] at ([xshift=4pt]div.east) {= 0.21875};
    \node[val] at ([xshift=4pt]add.east) {= 26.21875};
    \node[val] at ([xshift=4pt]mul.east) {= 42.18597};
    \node[ord] at ([shift={(-0.42,0.32)}]div.center) {1};
    \node[ord] at ([shift={(-0.42,0.32)}]add.center) {2};
    \node[ord] at ([shift={(-0.42,0.32)}]mul.center) {3};
  \end{tikzpicture}
  \caption{L'albero dell'espressione: si calcola dal basso verso l'alto,
           nell'ordine indicato dai numeri.}
  \label{fig:albero-maratona}
\end{figure}

\code{1.609} e \code{1760.0} sono \textbf{costanti reali} (in C, di tipo
\code{double}): per questo \code{yards / 1760.0} è una divisione fra reali,
con risultato reale. Infine \code{printf("\dots\%f\dots", kilometers)}
stampa il valore: la stringa di controllo contiene una
\term{specifica di conversione}, \code{\%f}, che dice “qui va stampato un
numero reale”; la variabile corrispondente, il secondo argomento, deve
essere di quel tipo.

\subsection{La trappola della divisione intera}

Basta togliere il \code{.0} alla costante per cambiare completamente il
risultato:

\begin{affianco}[0.6][c]
\begin{lstlisting}[style=c]
kilometers = 1.609 * (miles + yards / 1760);
\end{lstlisting}
\risultato
\begin{console}
A marathon is 41.834000 kilometers.
\end{console}
\end{affianco}

Ora \code{1760} è una costante \textbf{intera}, e \code{yards / 1760} è
una divisione fra interi: in C il risultato di una divisione fra interi è
un intero, e il resto si perde. 385 diviso 1760 fa 0 con resto 385, quindi
\code{yards / 1760} vale \textbf{0} e i 385 metri e rotti\dots\ spariscono:
$1{,}609 \cdot 26 = 41{,}834$. Il compilatore non segnala niente, perché
non c'è nessun errore di linguaggio: c'è un errore di ragionamento.

\section{Parole chiave e nomi}

Parole come \code{int} e \code{float} sono \term{parole chiave}
(\textit{keywords}): sono \textbf{riservate} e non possono essere usate come
nomi di variabili. Una variabile non può chiamarsi \code{int}, né
\code{return}. Per lo stesso motivo di chiarezza, non si usano come nomi di
variabili i nomi delle funzioni: niente variabili chiamate \code{printf} o
\code{main}. L'elenco completo delle parole chiave e le regole per i nomi
sono nel capitolo~\ref{chap:sintassi}.

\section{Il preprocessore: \code{\#define} e \code{\#include}}

\subsection{Le costanti simboliche}

Il preprocessore scandisce tutto il file e individua le direttive, cioè le
righe che cominciano con \code{\#}. Due delle più usate sono
\code{\#define} e \code{\#include}. Consideriamo le righe

\begin{lstlisting}[style=c]
#define LIMIT 100
#define PI    3.14159
\end{lstlisting}

Il preprocessore sostituirà nel file sorgente tutte le occorrenze di
\code{LIMIT} con \code{100} e tutte quelle di \code{PI} con
\code{3.14159}, tranne quelle che compaiono fra virgolette (nelle stringhe)
e nei commenti. Un \code{\#define} può stare ovunque nel file e ha effetto
solo sulle righe successive; di solito si mette all'inizio, e per
convenzione il nome si scrive in maiuscolo.

Le \term{costanti simboliche} rendono il programma più \textbf{leggibile}
(\code{LIMIT} dice che cosa rappresenta, \code{100} no) e più facile e
sicuro da \textbf{modificare}. Se la velocità della luce, circa 299\,792,458
km/s, compare in cento punti di un programma, scrivendo

\begin{lstlisting}[style=c]
#define C 299792.458   /* velocità della luce in km/s */
\end{lstlisting}

e usando \code{C} ovunque, una stima più accurata si inserisce cambiando
una sola riga, senza rischiare di dimenticare qualche occorrenza.

\subsection{I file di intestazione}

Una riga come

\begin{lstlisting}[style=c]
#include "my_file.h"
\end{lstlisting}

copia il contenuto di \code{my\_file.h} in quel punto del sorgente. Anche
\code{\#include} può stare ovunque, ma tipicamente si mette all'inizio. I
file che si includono sono gli \term{header file} e per convenzione hanno
estensione \code{.h}. Il sistema C ne fornisce alcuni, come
\code{stdio.h}, \code{stdlib.h}, \code{string.h}, \code{math.h}. La
differenza fra le due forme:

\begin{itemize}
  \item \code{\#include <stdio.h>}, con le parentesi angolari, cerca il file
        fra gli header del sistema;
  \item \code{\#include "my\_file.h"}, con le virgolette, lo cerca prima
        nella cartella del sorgente: si usa per i nostri header.
\end{itemize}

Ogni volta che si usa una funzione di input/output come \code{printf} si
deve includere \code{stdio.h}, che ne contiene il \textbf{prototipo}: la
descrizione che serve al compilatore per controllare che la chiamata sia
corretta. L'header serve alla \textbf{compilazione}; il codice della
funzione arriva dalla libreria, durante il \textbf{linking}.

\subsection{Header e libreria: la prova}

La funzione \code{pow(x, y)} calcola $x^y$; il suo prototipo sta in
\code{math.h}, il suo codice nella \textbf{libreria matematica}, che su
Linux è separata dalla libreria standard e va chiesta esplicitamente al
linker con l'opzione \code{-lm} (\textit{library m}):

\begin{affianco}[0.48]
\begin{lstlisting}[style=c, name=potenza.c]
#include <stdio.h>
#include <math.h>

int main(void)
{
    double base, esponente;

    printf("Base ed esponente: ");
    scanf("%lf%lf", &base, &esponente);
    printf("%g elevato a %g fa %g\n",
           base, esponente, pow(base, esponente));
    return 0;
}
\end{lstlisting}
\risultato
\begin{console}
$ gcc -Wall -o potenza potenza.c
/usr/bin/ld: /tmp/ccVCbC9C.o: in function `main':
potenza.c:(.text+0x61): undefined reference to `pow'
collect2: error: ld returned 1 exit status
$ gcc -Wall -o potenza potenza.c -lm
$ ./potenza
Base ed esponente: <<2 10>>
2 elevato a 10 fa 1024
\end{console}
\end{affianco}

Il primo tentativo fallisce, ma si guardi \textbf{chi} protesta:
\code{ld}, il linker, non il compilatore. Grazie a \code{math.h} il
compilatore sapeva tutto ciò che gli serviva e ha prodotto il file oggetto;
è il linker che non ha trovato il \textbf{codice} di \code{pow}
(\textit{undefined reference}). Con \code{-lm} gli si indica la libreria, e
tutto funziona. (Su macOS la libreria matematica è già compresa in quella
di sistema e \code{-lm} non serve, ma non fa danni.)

\section{Un esempio con le costanti: Pacific Sea}

\subsection{Il programma}

\textit{Pacific Sea} è il nome dato alla parte di oceano Pacifico visibile
dal campus di Santa Cruz, sulla baia di Monterey. Il programma converte la
sua area da chilometri quadrati in acri, miglia quadrate, piedi quadrati e
pollici quadrati. Le costanti stanno in un header scritto da noi:

\begin{lstlisting}[style=h, name=pacific_sea.h]
#include <stdio.h>

#define   AREA                        2337
#define   SQ_MILES_PER_SQ_KILOMETER   0.3861021585424458
#define   SQ_FEET_PER_SQ_MILE         (5280 * 5280)
#define   SQ_INCHES_PER_SQ_FOOT       144
#define   ACRES_PER_SQ_MILE           640
\end{lstlisting}

\begin{lstlisting}[style=c, name=pacific_sea.c]
/* Measuring the Pacific Sea. */

#include "pacific_sea.h"

int main(void)
{
    const int   pacific_sea = AREA;     /* in sq kilometers */
    double      acres, sq_miles, sq_feet, sq_inches;

    printf("\nThe Pacific Sea covers an area");
    printf(" of %d square kilometers.\n", pacific_sea);
    sq_miles = SQ_MILES_PER_SQ_KILOMETER * pacific_sea;
    sq_feet = SQ_FEET_PER_SQ_MILE * sq_miles;
    sq_inches = SQ_INCHES_PER_SQ_FOOT * sq_feet;
    acres = ACRES_PER_SQ_MILE * sq_miles;
    printf("In other units of measure this is:\n\n");
    printf("%22.7e acres\n", acres);
    printf("%22.7e square miles\n", sq_miles);
    printf("%22.7e square feet\n", sq_feet);
    printf("%22.7e square inches\n\n", sq_inches);
    return 0;
}
\end{lstlisting}

\begin{console}
$ ./pacific_sea

The Pacific Sea covers an area of 2337 square kilometers.
In other units of measure this is:

         5.7748528e+05 acres
         9.0232074e+02 square miles
         2.5155259e+10 square feet
         3.6223572e+12 square inches

\end{console}

\subsection{Gli \code{\#include} a catena}

La prima riga include una copia di \code{pacific\_sea.h}, e siccome questo
contiene a sua volta \code{\#include <stdio.h>}, il preprocessore include
anche una copia di \code{stdio.h}, che serve per \code{printf}
(figura~\ref{fig:include}). Il risultato è un unico testo, che è quello che
il compilatore riceve.

\begin{figure}[H]
  \centering
  \begin{tikzpicture}[font=\Lato\scriptsize,
      f/.style={wtnodeplain, minimum width=2.9cm, minimum height=1.2cm, align=center}]
    \node[f] (a) at (0,0) {{\FiraCodeNL pacific\_sea.c}\\{\FiraCodeNL\tiny \#include "pacific\_sea.h"}\\il programma};
    \node[f] (b) at (4.6,0) {{\FiraCodeNL pacific\_sea.h}\\{\FiraCodeNL\tiny \#include <stdio.h>}\\le costanti};
    \node[f] (c) at (9.2,0) {{\FiraCodeNL stdio.h}\\{\tiny (header di sistema)}\\i prototipi di I/O};
    \node[wtnode, minimum width=3.0cm, minimum height=1.2cm] (d) at (13.8,0) {un unico testo\\per il compilatore};
    \draw[wtarrowacc] (b) -- node[wtlab, above=3pt, fill=none] {copiato in} (a);
    \draw[wtarrowacc] (c) -- node[wtlab, above=3pt, fill=none] {copiato in} (b);
    \draw[wtarrow] (a.south) -- ++(0,-0.45) -| node[wtlab, pos=0.25] {dopo il preprocessore} (d.south);
  \end{tikzpicture}
  \caption{Il preprocessore risolve gli \code{\#include} uno dentro l'altro.}
  \label{fig:include}
\end{figure}

\subsection{Le costanti dell'header}

Le costanti simboliche dell'header (una reale, le altre intere) rendono il
codice più chiaro: se ci fosse una nuova stima dell'area basterebbe
cambiare \code{AREA}, e le altre danno un nome significativo ai numeri. Su
\code{SQ\_FEET\_PER\_SQ\_MILE} vale la pena fermarsi:

\begin{itemize}
  \item scrivere \code{(5280 * 5280)} invece di \code{27878400} è più
        chiaro: dice che in un miglio ci sono 5280 piedi, e che un miglio
        quadrato ne contiene il quadrato. Non si perde efficienza: il
        compilatore calcola da sé le espressioni fatte solo di costanti;
  \item le parentesi non sono sintatticamente necessarie, ma sono un'ottima
        abitudine. Il preprocessore fa una sostituzione di
        \textbf{testo}, senza capire nulla: se fosse
        \code{\#define SQ\_FEET\_PER\_SQ\_MILE 5280 * 5280}, l'espressione
        \code{x / SQ\_FEET\_PER\_SQ\_MILE} diventerebbe
        \code{x / 5280 * 5280}, cioè (divisione e moltiplicazione si
        eseguono da sinistra a destra) $(x / 5280) \cdot 5280$, che vale di
        nuovo circa $x$. Con le parentesi diventa
        \code{x / (5280 * 5280)}, come si voleva.
\end{itemize}

\subsection{\code{const}}

\code{const int pacific\_sea = AREA;} è la dichiarazione di una variabile
\code{pacific\_sea} di tipo \code{int}, \textbf{inizializzata} (cioè con un
valore assegnato già nella dichiarazione) a 2337: prima della compilazione
il preprocessore ha sostituito \code{AREA} con il suo valore. La parola
chiave \code{const} qualifica il tipo: la variabile riceve il suo valore
all'inizializzazione, ma \textbf{non può essere modificata}. Il compilatore
lo fa rispettare:

\begin{affianco}[0.5]
\begin{lstlisting}[style=c, name=costante.c]
#include <stdio.h>

int main(void)
{
    const int pacific_sea = 2337;

    pacific_sea = 3000;
    printf("%d\n", pacific_sea);
    return 0;
}
\end{lstlisting}
\risultato
\begin{console}
$ gcc -Wall -o costante costante.c
costante.c: In function (*@‘@*)main(*@’@*):
costante.c:7:17: error: assignment of read-only variable (*@‘@*)pacific_sea(*@’@*)
    7 |     pacific_sea = 3000;
      |                 ^
\end{console}
\end{affianco}

\subsection{I tipi reali e la notazione scientifica}

\code{acres}, \code{sq\_miles} e gli altri sono dichiarati \code{double}. In
C i tipi reali (detti \textit{floating point}, “a virgola mobile”) sono
tre: \code{float} (tipicamente 6 cifre significative), \code{double}
(tipicamente 15) e \code{long double} (almeno quanto \code{double}; sui PC
di solito 18). I valori dipendono dall'implementazione; il
capitolo~\ref{chap:tipi} li approfondisce.

L'istruzione \code{printf("\%22.7e acres\textbackslash n", acres);} stampa la
riga \code{5.7748528e+05 acres}. Il numero è in \textbf{notazione
scientifica}: \code{e+05} significa “per $10^5$”, quindi il valore è
$5{,}7748528 \cdot 10^5 = 577\,485{,}28$ acri. La conversione \code{\%e}
chiede proprio questo formato (detto \textit{e-format}), e la forma
\code{\%$m$.$n$e}, con $m$ e $n$ interi positivi, indica che il numero
occupa in tutto $m$ caratteri (allineato a destra) con $n$ cifre dopo il
punto: qui 22 caratteri e 7 cifre decimali. Per questo nell'output i numeri
sono preceduti da spazi e allineati in colonna.

\section{\code{printf} e \code{scanf}}

\subsection{Funzioni di input e output formattato}

\code{printf} si usa per l'output, \code{scanf} per l'input; la \code{f}
finale sta per \textit{formatted} (formattato). Non fanno parte del
linguaggio, ma della libreria standard, che viene collegata implicitamente
alla fine della compilazione. Il programmatore deve comunque dichiarare al
compilatore le funzioni che usa: il meccanismo dell'ANSI C è il
\textbf{prototipo}, e i prototipi delle funzioni della libreria standard
sono negli header standard (per \code{printf} e \code{scanf},
\code{stdio.h}).

Entrambe accettano una lista di argomenti:

\begin{itemize}
  \item la \term{stringa di controllo}, una stringa che contiene le
        \textbf{specifiche di conversione}: ciascuna è il carattere \code{\%}
        seguito da un \textbf{carattere di conversione} (per esempio
        \code{\%d} per un intero);
  \item gli altri argomenti, \textbf{uno per ogni conversione} della stringa
        di controllo, nello stesso ordine.
\end{itemize}

\subsection{Le conversioni di \code{printf}}

Per stampare le prime tre lettere dell'alfabeto ci sono molti modi: senza
conversioni, con tre conversioni per caratteri, con una conversione per
stringhe. Si può anche decidere quanto spazio dare a ciascuna conversione:
\code{\%3c} significa “un carattere, in un campo largo 3”, allineato a
destra.

\begin{affianco}[0.55][c]
\begin{lstlisting}[style=c]
printf("abc\n");
printf("%c%c%c\n", 'a', 'b', 'c');
printf("%s\n", "abc");
printf("%c%3c%5c\n", 'A', 'B', 'C');
\end{lstlisting}
\risultato
\begin{console}
abc
abc
abc
A  B    C
\end{console}
\end{affianco}

Nel codice, \code{'a'} fra apici singoli è una \textbf{costante carattere}
(un solo carattere), mentre \code{"abc"} fra virgolette è una
\textbf{stringa}. I caratteri di conversione principali di \code{printf}
sono:

\begin{center}
\renewcommand{\arraystretch}{1.3}
\begin{tabular}{@{}>{\FiraCodeNL}c l l@{}}
  \toprule
  \normalfont\textbf{conversione} & \textbf{stampa} & \textbf{esempio} \\
  \midrule
  \%c & un carattere & \code{printf("\%c", 'A')} $\to$ \code{A} \\
  \%d & un intero, in decimale & \code{printf("\%d", 42)} $\to$ \code{42} \\
  \%e & un reale, in notazione scientifica & \code{printf("\%e", 314159.0)} $\to$ \code{3.141590e+05} \\
  \%f & un reale & \code{printf("\%f", 3.14159)} $\to$ \code{3.141590} \\
  \%g & un reale, nel più corto fra \code{\%e} e \code{\%f} & \code{printf("\%g", 314159.0)} $\to$ \code{314159} \\
  \%s & una stringa & \code{printf("\%s", "abc")} $\to$ \code{abc} \\
  \bottomrule
\end{tabular}
\end{center}

Fra \code{\%} e il carattere di conversione si possono indicare l'ampiezza
del campo e, per i reali, il numero di cifre decimali (la
\textit{precisione}); un segno meno allinea a sinistra. Le parentesi quadre
nell'esempio servono solo a mostrare dove comincia e finisce il campo:

\begin{affianco}[0.55][c]
\begin{lstlisting}[style=c]
printf("[%d] [%5d] [%-5d]\n", 42, 42, 42);
printf("[%f] [%.2f] [%10.3f]\n", 3.14159, 3.14159, 3.14159);
printf("[%e] [%g] [%g]\n", 314159.0, 314159.0, 0.0000314);
\end{lstlisting}
\risultato
\begin{console}[Terminale]
[42] [   42] [42   ]
[3.141590] [3.14] [     3.142]
[3.141590e+05] [314159] [3.14e-05]
\end{console}
\end{affianco}

Si noti che \code{\%f} stampa sempre 6 cifre decimali, se non si indica
altro, e che la precisione \textbf{arrotonda}: 3.14159 con 3 decimali
diventa 3.142.

\subsection{\code{scanf} e l'operatore \code{\&}}

\code{scanf} è l'analoga di \code{printf} per l'input. Il primo parametro è
la stringa di controllo; l'input è sempre una sequenza di caratteri (ciò che
si digita sulla tastiera), e le conversioni specificano come i caratteri
devono essere interpretati. I parametri successivi, però, non sono variabili
ma \textbf{indirizzi}: gli indirizzi delle variabili in cui scrivere i
valori letti.

\begin{lstlisting}[style=c]
scanf("%d", &x);
\end{lstlisting}

\code{\&} è l'\term{operatore di indirizzo}: \code{\&x} è l'indirizzo della
variabile \code{x}, il “numero civico” della sua cella di memoria
(figura~\ref{fig:scanf-indirizzo}). Perché serve? Perché \code{scanf} deve
\textbf{scrivere} nella variabile, e per farlo deve sapere dove si trova.
Scrivendo \code{scanf("\%d", x)} le si passerebbe il \textit{valore} di
\code{x} (un numero qualunque, ancora indefinito), e \code{scanf} non
saprebbe dove mettere il dato letto; il motivo profondo lo vedremo alla
fine del capitolo, con il passaggio dei parametri per valore.

\begin{figure}[H]
  \centering
  \begin{tikzpicture}[font=\Lato\scriptsize]
    \node[fill=vscbg, rounded corners=3pt, inner sep=6pt, font=\FiraCodeNL\small] (cod) at (0,0)
      {\textcolor{vscfunc}{scanf}\textcolor{vscfg}{(}\textcolor{vscstring}{"\%d"}\textcolor{vscfg}{, }\textcolor{yellow-gradient-2}{\bfseries \&x}\textcolor{vscfg}{);}};
    \node[wtnode, font=\Lato\scriptsize, text width=3.4cm] (sc) at (6.0,0) {\textbf{scanf}\\legge i caratteri digitati, li converte in un intero e lo scrive all'indirizzo ricevuto};
    \memoria{m}{(11.8,0.72)}{996/\ldots/, 1000/42/x, 1004/\ldots/}
    \draw[wtarrow] (cod.east) -- node[wtlab, above=2pt, fill=none] {riceve \code{\&x}, cioè 1000} (sc.west);
    \draw[wtarrowacc] (sc.east) -- node[wtlab, above=2pt, fill=none] {scrive 42} ([xshift=-30pt]m-2.west);
    \node[mnota, anchor=north, align=center] at ([yshift=-4pt]m-3.south) {la variabile \code{x}\\all'indirizzo 1000};
  \end{tikzpicture}
  \caption{\code{scanf} riceve l'indirizzo di \code{x} e scrive lì il
           valore letto (l'indirizzo 1000 è solo un esempio).}
  \label{fig:scanf-indirizzo}
\end{figure}

Le conversioni principali di \code{scanf}:

\begin{center}
\renewcommand{\arraystretch}{1.3}
\begin{tabular}{@{}>{\FiraCodeNL}c l@{}}
  \toprule
  \normalfont\textbf{conversione} & \textbf{legge e scrive in una variabile di tipo} \\
  \midrule
  \%c  & \code{char} (un carattere) \\
  \%d  & \code{int} (un intero decimale) \\
  \%f  & \code{float} \\
  \%lf & \code{double} \\
  \%s  & una stringa (capitolo~\ref{chap:stringhe}) \\
  \bottomrule
\end{tabular}
\end{center}

\warning{\code{\%f} o \code{\%lf}?}{
  Con \code{printf}, \code{\%f} va bene sia per i \code{float} sia per i
  \code{double}. Con \code{scanf} no: per un \code{double} serve
  \code{\%lf}. Il motivo: \code{scanf} riceve solo un indirizzo, e dalla
  conversione deve capire quanti byte scrivere a quell'indirizzo (4 per un
  \code{float}, 8 per un \code{double}). Con la conversione sbagliata
  scrive un numero di byte sbagliato, e il valore letto è senza senso.}

\subsection{Un esempio: echo}

Il programma seguente legge tre caratteri, un intero, un \code{float} e un
\code{double}, e li ristampa (fa da “eco”):

\begin{lstlisting}[style=c, name=echo.c]
#include <stdio.h>

int main(void)
{
    char     c1, c2, c3;
    int      i;
    float    x;
    double   y;

    printf("\n%s\n%s", "Input three characters ", "an int, a float, and a double \n");
    scanf("%c%c%c%d%f%lf", &c1, &c2, &c3, &i, &x, &y);
    printf("\nHere is the data that you typed in:\n");
    printf("%3c%3c%3c\n%17d%17f%17f\n\n", c1, c2, c3, i, x, y);
    return 0;
}
\end{lstlisting}

Proviamolo due volte, con input quasi uguali: nel secondo c'è solo uno
spazio fra \code{B} e \code{C}.

\begin{console}
$ ./echo

Input three characters
an int, a float, and a double
<<ABC 3 55 77.7>>

Here is the data that you typed in:
  A  B  C
                3        55.000000        77.700000

\end{console}

\begin{console}
$ ./echo

Input three characters
an int, a float, and a double
<<AB C 3 55 77.7>>

Here is the data that you typed in:
  A  B
       1386527472         0.000000         0.000000

\end{console}

Il secondo risultato è disastroso, e il motivo è una regola di
\code{scanf}: quando legge \textbf{numeri} salta gli spazi che li
precedono, ma quando legge \textbf{caratteri} no, perché lo spazio
\textit{è} un carattere. La figura~\ref{fig:scanf-buffer} mostra quali
caratteri “mangia” ciascuna conversione.

\begin{figure}[H]
  \centering
  \begin{tikzpicture}[font=\FiraCodeNL\scriptsize,
      ch/.style={mcella, minimum width=0.46cm, minimum height=0.5cm, inner xsep=1pt},
      conv/.style={font=\FiraCodeNL\scriptsize\bfseries, text=accent!70!black},
      tit/.style={font=\Lato\scriptsize\bfseries, anchor=east}]
    \noindici
    \node[tit] at (-0.2,0) {input 1};
    \celle[ch]{a}{(0,0)}{A,B,C,\char"2423,3,\char"2423,5,5,\char"2423,7,7,.,7,\textbackslash n}
    \foreach \ia/\ib/\conv in {0/0/\%c, 1/1/\%c, 2/2/\%c, 3/4/\%d, 5/7/\%f, 8/12/\%lf} {
      \draw[wtbrace] ([yshift=-3pt]a-\ia.south west) -- ([yshift=-3pt]a-\ib.south east);
      \node[conv, anchor=north] at ([yshift=-0.32cm]$(a-\ia.south)!0.5!(a-\ib.south)$) {\conv};
    }
    \node[mnota, anchor=west] at ([xshift=0.4cm]a-13.east) {6 conversioni riuscite: \textbf{A}, \textbf{B}, \textbf{C}, \textbf{3}, \textbf{55}, \textbf{77.7}};
    \node[tit] at (-0.2,-2.0) {input 2};
    \celle[ch]{b}{(0,-2.0)}{A,B,\char"2423,C,\char"2423,3,\char"2423,5,5,\char"2423,7,7,.,7}
    \foreach \ia/\ib/\conv in {0/0/\%c, 1/1/\%c, 2/2/\%c} {
      \draw[wtbrace] ([yshift=-3pt]b-\ia.south west) -- ([yshift=-3pt]b-\ib.south east);
      \node[conv, anchor=north] at ([yshift=-0.32cm]$(b-\ia.south)!0.5!(b-\ib.south)$) {\conv};
    }
    \draw[wtbrace, draw=merr] ([yshift=-3pt]b-3.south west) -- ([yshift=-3pt]b-3.south east);
    \node[conv, anchor=north, text=merr] at ([yshift=-0.32cm]b-3.south) {\%d \merr};
    \node[mnota, anchor=west, text width=6.4cm] at ([xshift=0.4cm]b-13.east) {la terza \code{\%c} legge lo \textbf{spazio}; poi \code{\%d} trova
      una \code{C}, che non è un numero, e \code{scanf} si ferma: \code{i}, \code{x} e \code{y}
      restano con i valori a caso che avevano};
  \end{tikzpicture}
  \caption{I caratteri dell'input (\code{\char"2423} indica uno spazio) e la
           conversione che li consuma.}
  \label{fig:scanf-buffer}
\end{figure}

Nel secondo caso \code{i}, \code{x} e \code{y} non ricevono nulla, e
\code{printf} stampa ciò che contenevano: valori casuali, diversi da un
computer all'altro (e da un'esecuzione all'altra). Come accorgersene dal
programma? \code{scanf} \textbf{restituisce il numero di conversioni
riuscite}. Salvandolo in una variabile e stampandolo (è il programma
\code{echo\_conversions.c} delle lezioni) si ottiene 6 nel primo caso e 3 nel
secondo:

\begin{affianco}[0.55][c]
\begin{lstlisting}[style=c]
int conversions;
...
conversions = scanf("%c%c%c%d%f%lf", &c1, &c2, &c3, &i, &x, &y);
...
printf("Correct conversions. %d\n", conversions);
\end{lstlisting}
\risultato
\begin{console}
Correct conversions. 6
...
Correct conversions. 3
\end{console}
\end{affianco}

Controllare il valore restituito da \code{scanf} è il modo corretto di
verificare che l'input sia stato letto: lo useremo fra poco.

\section{Il controllo del flusso}

\subsection{Vero e falso in C}

Le istruzioni sono normalmente eseguite in sequenza, ma la maggior parte
degli algoritmi richiede di alterare il flusso sequenziale
(capitolo~\ref{chap:strutturata}). Il C offre \code{if} e
\code{if}-\code{else} per scegliere fra azioni alternative, \code{while} e
\code{for} per ripetere. Tutti richiedono di valutare un'espressione
\textbf{logica}, cioè vera o falsa, e in C la regola è semplicissima:

\begin{itemize}
  \item il \textbf{falso} è rappresentato dal valore \textbf{0};
  \item \textbf{ogni valore diverso da zero} è interpretato come
        \textbf{vero}.
\end{itemize}

\subsection{\code{if}}

La forma generale è

\begin{lstlisting}[style=c]
if (espressione) {
    istruzioni_if;
}
istruzione_successiva;
\end{lstlisting}

Se \code{espressione} è diversa da zero (vera) si eseguono
\code{istruzioni\_if} e poi \code{istruzione\_successiva}; se vale 0 (falsa)
si salta direttamente a \code{istruzione\_successiva}. È la selezione a una
via del capitolo precedente. Un esempio:

\begin{affianco}[0.4][c]
\begin{lstlisting}[style=c]
a = 1;
if (b == 3) {
    a = 5;
}
printf("%d", a);
\end{lstlisting}
\risultato
Il simbolo \code{==} è l'operatore di \textbf{uguaglianza}: verifica se
\code{b} è uguale a 3. Se è vero, si assegna 5 ad \code{a} e viene stampato
\textbf{5}; se è falso si esegue subito la stampa, e viene stampato
\textbf{1}.
\end{affianco}

\error{\code{==} e \code{=} sono due cose diverse}{
  \code{==} \textbf{confronta}, \code{=} \textbf{assegna}. Scrivere
  \code{if (b = 3)} non è un errore per il compilatore: assegna 3 a
  \code{b}, e siccome il valore dell'assegnazione è 3 (diverso da zero, cioè
  vero), il ramo viene eseguito \textbf{sempre}. È uno degli errori più
  frequenti in C; il capitolo~\ref{chap:flusso} ci torna in dettaglio.
  Compilando con \code{-Wall}, gcc avvisa.}

Fra le graffe può esserci una sequenza di istruzioni; con una sola
istruzione le graffe si possono omettere, ma conviene \textbf{metterle
sempre}: se in seguito si aggiunge un'istruzione al blocco, senza graffe
resterebbe fuori dall'\code{if}, con un errore difficile da vedere.

\subsection{\code{if}-\code{else}}

\begin{affianco}[0.4][c]
\begin{lstlisting}[style=c]
if (espressione) {
    istruzioni_if;
} else {
    istruzioni_else;
}
\end{lstlisting}
\risultato
Se \code{espressione} è diversa da zero si esegue il primo blocco, se è zero
il secondo: è la selezione a due vie. Uno dei due blocchi viene sempre
eseguito, mai entrambi.
\end{affianco}

Nell'esempio seguente, se \code{cnt} vale 0 il programma stampa
$2 + 3 + 5 = 10$, altrimenti $-1 - 2 - 3 = -6$. A destra la variante con
\code{if (cnt)}: non c'è nessun confronto, la decisione si basa solo sul
valore di \code{cnt} (0 è falso, il resto vero), e i risultati si
invertono.

\begin{affianco}[0.5]
\begin{lstlisting}[style=c]
if (cnt == 0) {
    a = 2;
    b = 3;
    c = 5;
} else {
    a = -1;
    b = -2;
    c = -3;
}
printf("%d", a + b + c);
\end{lstlisting}
\risultato
\begin{lstlisting}[style=c]
if (cnt) {
    a = 2;
    b = 3;
    c = 5;
} else {
    a = -1;
    b = -2;
    c = -3;
}
printf("%d", a + b + c);
\end{lstlisting}
\end{affianco}

\begin{center}
\renewcommand{\arraystretch}{1.3}
\begin{tabular}{@{}lcc@{}}
  \toprule
  & \textbf{con \code{cnt == 0}} & \textbf{con \code{cnt}} \\
  \midrule
  se \code{cnt} vale 0           & stampa 10 & stampa $-6$ \\
  se \code{cnt} vale, per esempio, 4 & stampa $-6$ & stampa 10 \\
  \bottomrule
\end{tabular}
\end{center}

\subsection{\code{while}}

\begin{affianco}[0.42]
\begin{lstlisting}[style=c, name=sum.c]
#include <stdio.h>

int main(void)
{
    int   i = 1, sum = 0;

    while (i <= 5) {
        sum = sum + i;
        i = i + 1;
    }
    printf("sum = %d\n", sum);
    return 0;
}
\end{lstlisting}
\risultato
\begin{console}
$ ./sum
sum = 15
\end{console}
Il programma somma i numeri da 1 a 5. Nella dichiarazione,
\code{int i = 1, sum = 0;} le due variabili vengono anche
\textbf{inizializzate}. \code{<=} significa “minore o uguale”.
\end{affianco}

Il \code{while} verifica se \code{i <= 5}: se la condizione è vera esegue
le istruzioni del blocco fino alla \code{\}}, poi il controllo
\textbf{torna all'inizio} del ciclo per ripetere il test. L'iterazione
continua finché il test non fallisce. È il ciclo con la condizione in testa
del capitolo~\ref{chap:strutturata}. Simuliamolo:

\begin{center}
\renewcommand{\arraystretch}{1.3}
\begin{tabular}{@{}cccc@{}}
  \toprule
  \textbf{giro} & \textbf{\code{i <= 5}?} & \textbf{\code{sum} dopo il giro} & \textbf{\code{i} dopo il giro} \\
  \midrule
  1 & $1 \le 5$, vero & 0 + 1 = 1 & 2 \\
  2 & $2 \le 5$, vero & 1 + 2 = 3 & 3 \\
  3 & $3 \le 5$, vero & 3 + 3 = 6 & 4 \\
  4 & $4 \le 5$, vero & 6 + 4 = 10 & 5 \\
  5 & $5 \le 5$, vero & 10 + 5 = 15 & 6 \\
  6 & $6 \le 5$, \textbf{falso} & si esce: \code{sum} vale 15 & \\
  \bottomrule
\end{tabular}
\end{center}

La forma generale è

\begin{lstlisting}[style=c]
while (espressione) {
    istruzioni_while;
}
istruzioni;
\end{lstlisting}

Se \code{espressione} è diversa da zero si eseguono
\code{istruzioni\_while} e poi si verifica di nuovo \code{espressione}; il
processo continua finché \code{espressione} vale 0. Perché il ciclo
finisca, è \textbf{necessaria} un'istruzione dentro il ciclo che possa
cambiare il valore dell'espressione (qui \code{i = i + 1}): senza, il test
resterebbe vero per sempre.

\subsection{Assegnazioni abbreviate e incremento}

Istruzioni come \code{sum = sum + i;} sono così frequenti che il C ha una
forma abbreviata: \code{sum += i;}. In generale, un'istruzione della forma
\codel{variabile op= espressione;} (con \code{op} uno fra \code{+}, \code{-},
\code{*}, \code{/}) equivale a \codel{variabile = variabile op espressione;}.

Per aggiungere o togliere 1 ci sono gli \textbf{operatori di incremento e
decremento}:

\begin{center}
\renewcommand{\arraystretch}{1.3}
\begin{tabular}{@{}>{\FiraCodeNL}l l@{}}
  \toprule
  \normalfont\textbf{forma} & \textbf{come istruzione a sé, equivale a} \\
  \midrule
  ++i \textnormal{e} i++ & \code{i = i + 1;} \\
  --i \textnormal{e} i-- & \code{i = i - 1;} \\
  \bottomrule
\end{tabular}
\end{center}

Il ciclo di \code{sum.c} si scrive quindi anche così:

\begin{lstlisting}[style=c]
while (i <= 5) {
    sum += i;
    ++i;
}
\end{lstlisting}

\warning{\code{++} dentro le espressioni}{
  Come istruzioni isolate, \code{++i} e \code{i++} fanno la stessa cosa. Usati
  \textbf{dentro} un'espressione più grande (per esempio \code{a = i++;}) si
  comportano in modo diverso, e il loro funzionamento è più complicato: lo
  vedremo con calma nel capitolo~\ref{chap:sintassi}. Per ora li useremo
  solo da soli.}

\subsection{\code{for}}

Il ciclo \code{for} ha la forma

\begin{affianco}[0.48][c]
\begin{lstlisting}[style=c]
for (expr1; expr2; expr3) {
    istruzioni_for;
}
\end{lstlisting}
\risultato
\begin{lstlisting}[style=c]
expr1;
while (expr2) {
    istruzioni_for;
    expr3;
}
\end{lstlisting}
\end{affianco}

ed è equivalente al \code{while} a destra. Tipicamente \code{expr1} fa
un'assegnazione iniziale, \code{expr2} un test, \code{expr3} incrementa una
variabile: sono i tre passi del ciclo su contatore del
capitolo~\ref{chap:strutturata} (inizializzazione, test, modifica), riuniti
in un'unica riga. La somma da 1 a 5 diventa:

\begin{affianco}[0.48][c]
\begin{lstlisting}[style=c]
for (i = 1; i <= 5; i++) {
    sum += i;
}
\end{lstlisting}
\risultato
\begin{immagine}
\begin{tikzpicture}[font=\Lato\scriptsize]
  \node[fcproc, fill=green-gradient-1!40!white] (a) at (0,0) {\code{i = 1}};
  \node[fcgiunzione] (j) at (0,-0.6) {};
  \node[fcdec, fill=yellow-gradient-1!60!white] (t) at (0,-1.35) {\code{i <= 5}?};
  \node[fcproc] (b) at (0,-2.35) {\code{sum += i}};
  \node[fcproc, fill=accent!10!white] (c) at (0,-3.15) {\code{i++}};
  \draw[fcfreccia] (a) -- (j); \draw[fcfreccia] (j) -- (t);
  \draw[fcfreccia] (t) -- node[fcsi, right] {vero} (b);
  \draw[fcfreccia] (b) -- (c);
  \draw[fcfreccia, draw=accent!70!black] (c.south) -- (0,-3.6) -- (1.6,-3.6) -- (1.6,-0.6) -- (j.east);
  \draw[fcfreccia] (t.west) -- node[fcno, above] {falso} (-1.6,-1.35) -- (-1.6,-3.9);
\end{tikzpicture}
\end{immagine}
\end{affianco}

\section{Un esempio completo: \code{running\_sum}}

\subsection{Il programma}

Il programma seguente legge una sequenza di numeri e, per ciascuno, stampa
una riga di tabella con il conteggio, il numero letto, il minimo, il
massimo, la somma e la media dei numeri letti fino a quel momento.

\begin{lstlisting}[style=c, name=running_sum.c]
/* Compute the minimum, maximum, sum, and average. */

#include <stdio.h>
#include <stdlib.h>

int main(void)
{
    int      i;
    double   x, min, max, sum, avg;

    if (scanf("%lf", &x) != 1) {
        printf("No data found - bye!\n");
        exit(1);
    }
    min = max = sum = avg = x;
    printf("%5s%9s%9s%9s%12s%12s\n%5s%9s%9s%9s%12s%12s\n\n",
           "Count", "Item", "Min", "Max", "Sum", "Average",
           "-----", "----", "---", "---", "---", "-------");
    printf("%5d%9.1f%9.1f%9.1f%12.3f%12.3f\n",
           1, x, min, max, sum, avg);
    for (i = 2; scanf("%lf", &x) == 1; ++i) {
        if (x < min)
            min = x;
        else if (x > max)
            max = x;
        sum += x;
        avg = sum / i;
        printf("%5d%9.1f%9.1f%9.1f%12.3f%12.3f\n",
               i, x, min, max, sum, avg);
    }
    return 0;
}
\end{lstlisting}

\subsection{Riga per riga}

\begin{itemize}
  \item \code{if (scanf("\%lf", \&x) != 1)}: \code{scanf} restituisce il
        numero di conversioni riuscite; qui ne è richiesta una sola, quindi
        un valore diverso da 1 (\code{!=} significa “diverso”) vuol dire
        che non è stato letto niente. In quel caso il programma stampa un
        messaggio ed esce con \code{exit(1)}: \code{exit}, il cui prototipo
        sta in \code{stdlib.h}, termina il programma restituendo al sistema
        operativo il valore indicato (qui 1, “qualcosa è andato storto”).
  \item \code{min = max = sum = avg = x;}: un'assegnazione in C è anche
        un'espressione, che \textbf{vale} il valore assegnato. Le
        assegnazioni a catena si eseguono da destra a sinistra:
        \code{avg = x}, poi \code{sum = (avg = x)}, e così via. Tutte e
        quattro le variabili ricevono il primo numero letto.
  \item Le due \code{printf} con \code{\%5s}, \code{\%9s}\dots\ stampano
        l'intestazione della tabella: stringhe costanti in campi di
        larghezza fissa, così le colonne restano allineate.
  \item \code{for (i = 2; scanf("\%lf", \&x) == 1; ++i)}: un ciclo \code{for}
        in cui il \textbf{test è la lettura stessa}. Il ciclo continua finché
        \code{scanf} riesce a leggere un numero, cioè finché c'è input; non
        si sa in anticipo quante volte girerà. \code{i} conta i numeri letti
        e parte da 2 perché il primo è già stato letto prima del ciclo.
  \item \code{if (x < min) min = x; else if (x > max) max = x;}: aggiorna, se
        necessario, il minimo o il massimo. L'\code{else} contiene a sua
        volta un \code{if}. (Si poteva anche scrivere con due \code{if}
        indipendenti.)
  \item \code{avg = sum / i;}: qui la divisione è fra un \code{double} e un
        \code{int}, quindi è reale: nessuna trappola.
\end{itemize}

\subsection{Tre modi di dare l'input}

Il programma legge i numeri dallo \term{standard input}, il “canale” da cui
un programma riceve i dati. Normalmente è la tastiera, e allora i numeri si
possono scrivere uno alla volta, premendo \tasto{Invio} dopo ciascuno, oppure
tutti su una riga; ma come fa il programma a sapere che l'input è finito?
Bisogna segnalarlo con \tasti{Ctrl,D} all'inizio di una riga: \code{scanf}
non riesce più a leggere, restituisce un valore diverso da 1, e il ciclo
finisce.

\begin{console}
$ ./running_sum
<<3>>
Count     Item      Min      Max         Sum     Average
-----     ----      ---      ---         ---     -------

    1      3.0      3.0      3.0       3.000       3.000
<<-5 7>>
    2     -5.0     -5.0      3.0      -2.000      -1.000
    3      7.0     -5.0      7.0       5.000       1.667
<<Ctrl+D>>
$
\end{console}

(\tasti{Ctrl,D} in realtà non compare sullo schermo: qui è indicato per
chiarezza.) Il terzo modo è il più comodo per provare un programma: scrivere i dati in
un file e usarlo \textbf{al posto} della tastiera, con la
\term{redirezione} \code{<} (figura~\ref{fig:redirezione}). Il file
\code{running\_sum.in} delle lezioni contiene i dieci numeri
\code{3 -5 7 -9 11 -13 15 -17 19 -21}.

\begin{console}
$ ./running_sum < running_sum.in
Count     Item      Min      Max         Sum     Average
-----     ----      ---      ---         ---     -------

    1      3.0      3.0      3.0       3.000       3.000
    2     -5.0     -5.0      3.0      -2.000      -1.000
    3      7.0     -5.0      7.0       5.000       1.667
    4     -9.0     -9.0      7.0      -4.000      -1.000
    5     11.0     -9.0     11.0       7.000       1.400
    6    -13.0    -13.0     11.0      -6.000      -1.000
    7     15.0    -13.0     15.0       9.000       1.286
    8    -17.0    -17.0     15.0      -8.000      -1.000
    9     19.0    -17.0     19.0      11.000       1.222
   10    -21.0    -21.0     19.0     -10.000      -1.000
\end{console}

\begin{figure}[H]
  \centering
  \begin{tikzpicture}[font=\Lato\scriptsize,
      prog/.style={wtnode, font=\Lato\scriptsize\bfseries, minimum width=2.6cm, minimum height=0.9cm},
      dev/.style={wtnodeplain, minimum width=2.3cm, minimum height=0.8cm},
      can/.style={font=\FiraCodeNL\scriptsize, text=accent!70!black}]
    \node[dev] (k) at (0,0) {tastiera};
    \node[prog] (p) at (4.9,0) {running\_sum};
    \node[dev] (s) at (9.8,0) {schermo};
    \draw[wtarrow] (k) -- node[can, above] {standard input} (p);
    \draw[wtarrow] (p) -- node[can, above] {standard output} (s);
    \node[font=\Lato\scriptsize\bfseries, anchor=west] at (11.4,0) {normalmente};
    \node[dev, draw=orange-gradient-6!80!black, fill=orange-gradient-1!30!white] (f) at (0,-1.5) {\FiraCodeNL running\_sum.in};
    \node[prog] (p2) at (4.9,-1.5) {running\_sum};
    \node[dev] (s2) at (9.8,-1.5) {schermo};
    \draw[wtarrow] (f) -- node[can, above] {<} (p2);
    \draw[wtarrow] (p2) -- (s2);
    \node[font=\Lato\scriptsize\bfseries, anchor=west, align=left] at (11.4,-1.5) {\FiraCodeNL ./running\_sum < running\_sum.in};
    \node[dev] (k3) at (0,-3.0) {tastiera};
    \node[prog] (p3) at (4.9,-3.0) {running\_sum};
    \node[dev, draw=orange-gradient-6!80!black, fill=orange-gradient-1!30!white] (s3) at (9.8,-3.0) {\FiraCodeNL risultati.txt};
    \draw[wtarrow] (k3) -- (p3);
    \draw[wtarrow] (p3) -- node[can, above] {>} (s3);
    \node[font=\Lato\scriptsize\bfseries, anchor=west, align=left] at (11.4,-3.0) {\FiraCodeNL ./running\_sum > risultati.txt};
  \end{tikzpicture}
  \caption{La redirezione: con \code{<} un file prende il posto della
           tastiera, con \code{>} un file prende il posto dello schermo.}
  \label{fig:redirezione}
\end{figure}

È la shell a collegare il file al programma: il programma non se ne accorge
nemmeno, continua a leggere dallo standard input. Simmetricamente, con
\code{>} si può mandare l'output in un file invece che sullo schermo, e i due
si possono combinare: \code{./running\_sum < running\_sum.in > risultati.txt}.
Infine, se non c'è nessun dato, il programma esce con il suo messaggio, e il
valore restituito al sistema è 1:

\begin{console}
$ ./running_sum < /dev/null
No data found - bye!
$ echo $?
1
\end{console}

(\code{/dev/null} è un file speciale di Linux, sempre vuoto.)

\subsection{Una versione senza la lettura iniziale}

Nel programma si fa una prima lettura \textbf{prima} del ciclo, poi le altre
\textbf{nel} ciclo: lo stesso codice di lettura compare due volte. Quando è
possibile, è più elegante mettere tutto dentro il ciclo. Ecco la versione
delle lezioni con un \code{while}, in cui \code{cont} conta i numeri letti e
il caso “nessun dato” si controlla alla fine:

\begin{lstlisting}[style=c, name=running_sum2.c]
/* Compute the minimum, maximum, sum, and average. */

#include <stdio.h>

int main(void)
{
    int      cont = 0;
    double   x, min, max, sum, avg;

    while (scanf("%lf", &x) == 1) {
        cont = cont + 1;
        if (cont == 1) {
            sum = avg = min = max = x;
            printf("%5s%9s%9s%9s%12s%12s\n%5s%9s%9s%9s%12s%12s\n\n",
                   "Count", "Item", "Min", "Max", "Sum", "Average",
                   "-----", "----", "---", "---", "---", "-------");
        } else {
            if (x < min) min = x;
            else if (x > max) max = x;
            sum += x;
            avg = sum / cont;
        }
        printf("%5d%9.1f%9.1f%9.1f%12.3f%12.3f\n",
               cont, x, min, max, sum, avg);
    }
    if (cont == 0) {
        printf("No data found - bye!\n");
    }
    return 0;
}
\end{lstlisting}

L'output è identico. Di solito una variabile come \code{i} è la variabile di
un ciclo su contatore; qui \code{cont} conta gli elementi letti, e il primo
elemento è trattato a parte dentro il ciclo (\code{cont == 1}). Non serve
più \code{stdlib.h}, perché non si usa \code{exit}.

\section{Quando la compilazione fallisce}

I messaggi del compilatore sono importanti e vanno letti; a volte però
possono creare confusione. A lezione si compila il file
\code{running\_sum2Errore.c}, identico a quello appena visto tranne che per
un dettaglio: nella prima riga manca la barra iniziale del commento.

\begin{lstlisting}[style=c, name=running_sum2Errore.c]
* Compute the minimum, maximum, sum, and average. */

#include <stdio.h>
...
\end{lstlisting}

Il risultato è una lunga, deprimente lista di \textbf{51 errori}, molti dei
quali addirittura negli header file di sistema:

\begin{console}
$ gcc -Wall -o running_sum2 running_sum2Errore.c
running_sum2Errore.c:1:11: error: expected '=', ',', ';', 'asm' or '__attribute__' before 'the'
    1 | * Compute the minimum, maximum, sum, and average. */
      |           ^~~
running_sum2Errore.c:1:11: error: unknown type name 'the'
In file included from /usr/include/stdio.h:44,
                 from running_sum2Errore.c:3:
/usr/include/x86_64-linux-gnu/bits/types/struct_FILE.h:95:3: error: unknown type name 'size_t'
   95 |   size_t __pad5;
      |   ^~~~~~
...
\end{console}

Che cosa è successo? Senza \code{/*} la prima riga non è più un commento: il
compilatore la legge come codice C e cerca di darle un senso
(\code{* Compute the\dots}), fallendo. Ma non sa dove finisce l'errore, e
continua a leggere il resto del file “fuori fase”: anche il contenuto di
\code{stdio.h}, perfettamente corretto, viene letto nel contesto sbagliato e
produce errori a cascata. La regola del capitolo~\ref{chap:ambiente} vale
più che mai: si guarda \textbf{il primo} errore, che qui indica
correttamente la riga 1, e si ignorano gli altri finché non lo si è
corretto.

\info{Una curiosità}{
  Il file \code{running\_sum2.c} distribuito con le lezioni ha lo stesso
  difetto di \code{running\_sum2Errore.c} (i due file sono identici):
  per compilarlo bisogna aggiungere la \code{/} mancante all'inizio.}

\section{Le funzioni}

\subsection{Dichiarare prima di usare}

Una \term{funzione} è un pezzo di codice che rappresenta un blocco del
processo di soluzione del problema: è il sotto-programma del
capitolo~\ref{chap:strutturata}. Un programma C è costituito da una o più
funzioni, in uno o più file sorgente. Una di esse è necessariamente
\code{main}, da cui comincia l'esecuzione; le altre sono richiamate da
\code{main} o da altre funzioni.

In C le funzioni devono essere \textbf{dichiarate prima di usarle}.
Supponiamo di voler usare la funzione \code{pow} della libreria matematica:
la chiamata \code{pow(x, y)} restituisce $x^y$. La sua dichiarazione è

\begin{lstlisting}[style=c]
double pow(double x, double y);
\end{lstlisting}

Questo tipo di dichiarazione si chiama \term{prototipo} della funzione: dice
il tipo del valore restituito (\code{double}), il nome e il tipo dei
parametri. I nomi dei parametri, nel prototipo, sono ignorati dal
compilatore: \code{double pow(double, double);} è equivalente, perché al
compilatore interessano solo i tipi. È comunque buona norma metterli,
soprattutto se sono significativi.

La forma generale di un prototipo è
\code{tipo nome\_funzione(lista dei tipi dei parametri);}. La parola chiave
\code{void} si usa quando la funzione non ha parametri
(\code{int f(void);}) o non restituisce niente (\code{void f(int);}). Se
necessario, al momento della chiamata gli argomenti vengono convertiti al
tipo indicato nel prototipo: chiamando \code{pow(2, 10)} con due interi,
questi diventano \code{2.0} e \code{10.0}.

\subsection{L'esempio maxmin}

Il programma seguente stampa delle informazioni su se stesso, legge un intero
\code{n}, legge \code{n} numeri reali e ne trova il massimo e il minimo.

\begin{lstlisting}[style=c, name=maxmin.c]
#include <stdio.h>

float   maximum(float x, float y);
float   minimum(float x, float y);
void    prn_info(void);

int main(void)
{
    int     i, n;
    float   max, min, x;

    prn_info();
    printf("Input n:  ");
    scanf("%d", &n);
    printf("\nInput %d real numbers:  ", n);
    scanf("%f", &x);
    max = min = x;
    for (i = 2; i <= n; ++i) {
        scanf("%f", &x);
        max = maximum(max, x);
        min = minimum(min, x);
    }
    printf("\n%s%11.3f\n%s%11.3f\n\n",
           "Maximum value:", max,
           "Minimum value:", min);
    return 0;
}

float maximum(float x, float y)
{
    if (x > y)
        return x;
    else
        return y;
}

float minimum(float x, float y)
{
    if (x < y)
        return x;
    else
        return y;
}

void prn_info(void)
{
    printf("\n%s\n%s\n\n",
           "This program reads an integer value for n, and then",
           "processes n real numbers to find max and min values.");
}
\end{lstlisting}

\begin{console}
$ ./maxmin

This program reads an integer value for n, and then
processes n real numbers to find max and min values.

Input n:  <<5>>

Input 5 real numbers:  <<737.7799  -11.2e+3  -777  0.001  3.14159>>

Maximum value:    737.780
Minimum value: -11200.000

\end{console}

Si noti che \code{scanf} con \code{\%f} accetta anche la notazione
scientifica: \code{-11.2e+3} è $-11{,}2 \cdot 10^3 = -11\,200$.

\subsection{Come funziona}

\begin{itemize}
  \item In cima al file ci sono i \textbf{prototipi} di \code{maximum},
        \code{minimum} e \code{prn\_info}. I primi due dicono che le funzioni
        hanno due parametri di tipo \code{float} e restituiscono un
        \code{float}; il terzo che \code{prn\_info} non ha parametri e non
        restituisce alcun valore.
  \item \code{prn\_info();} chiama la funzione, che a sua volta chiama
        \code{printf} per stampare le informazioni sul programma e sul suo
        uso.
  \item Il programma chiede e legge \code{n} (\code{\%d}: un intero
        decimale), poi legge il primo numero reale in \code{x}, e
        \code{max = min = x;} inizializza massimo e minimo con quel valore.
  \item A ogni giro del ciclo si legge un nuovo valore di \code{x};
        \code{max = maximum(max, x);} calcola il massimo fra \code{max} e
        \code{x} e lo assegna a \code{max}, e lo stesso per il minimo.
\end{itemize}

La \textbf{definizione} di \code{maximum} specifica che cosa la funzione fa
quando viene chiamata. È fatta di una \textbf{testa}, la parte prima della
prima graffa (\code{float maximum(float x, float y)}: il primo \code{float}
è il tipo del valore restituito, gli altri due i tipi dei parametri), e di un
\textbf{corpo}, la parte fra le graffe.

L'istruzione \code{return} \textbf{interrompe} la funzione e restituisce al
chiamante il valore della sua espressione; la forma generale è
\code{return espressione;}. In una funzione \code{void} si può scrivere
\code{return;} senza espressione, o lasciare che la funzione arrivi da sola
alla graffa finale, come fa \code{prn\_info}.

\subsection{Due \code{x} diverse}

Le funzioni \code{maximum} e \code{minimum} hanno entrambe i parametri
\code{x} e \code{y}, e una variabile \code{x} c'è anche in \code{main}. Non
c'è nessuna relazione fra loro: sono variabili \textbf{logicamente e
fisicamente diverse}, in celle di memoria diverse. Usare lo stesso nome
informa solo chi legge che il loro uso è simile. La
figura~\ref{fig:due-x} mostra la memoria durante la chiamata
\code{maximum(max, x)} al secondo giro.

\begin{figure}[H]
  \centering
  \begin{tikzpicture}[font=\Lato\scriptsize]
    \variabile{mi}{(1.4,0)}{i}{2}
    \variabile{mn}{(1.4,-0.6)}{n}{5}
    \variabile[minimum width=1.7cm]{mmax}{(1.55,-1.2)}{max}{737.7799}
    \variabile[minimum width=1.7cm]{mmin}{(1.55,-1.8)}{min}{737.7799}
    \variabile[minimum width=1.7cm]{mx}{(1.55,-2.4)}{x}{-11200}
    \scopo{2}{main}{main}{(mi-nome) (mn) (mmax) (mmin) (mx) (mx-nome)}
    \variabile[minimum width=1.7cm]{fx}{(8.0,-1.2)}{x}{737.7799}
    \variabile[minimum width=1.7cm]{fy}{(8.0,-1.8)}{y}{-11200}
    \scopo{3}{maximum}{maximum}{(fx-nome) (fx) (fy) (fy-nome)}
    \draw[mcerca] (mmax.east) -- node[wtlab, above, fill=none] {copia} (fx.west -| fx-nome.west);
    \draw[mcerca] (mx.east) -- node[wtlab, below, fill=none] {copia} (fy.west -| fy-nome.west);
  \end{tikzpicture}
  \caption{La \code{x} di \code{main} e la \code{x} di \code{maximum} sono
           due variabili distinte: \code{maximum} riceve una
           \textbf{copia} dei valori.}
  \label{fig:due-x}
\end{figure}

\subsection{Che cosa succede senza prototipi}

Le slide mostrano anche una versione di \code{maxmin} senza i prototipi in
cima al file. Con gcc il risultato è questo:

\begin{console}
$ gcc -Wall -o maxmin maxmin.c
maxmin.c:9:4: warning: implicit declaration of function (*@‘@*)prn_info(*@’@*)
maxmin.c:17:13: warning: implicit declaration of function (*@‘@*)maximum(*@’@*)
maxmin.c:18:13: warning: implicit declaration of function (*@‘@*)minimum(*@’@*)
maxmin.c:26:7: error: conflicting types for (*@‘@*)maximum(*@’@*); have (*@‘@*)float(float,  float)(*@’@*)
maxmin.c:34:7: error: conflicting types for (*@‘@*)minimum(*@’@*); have (*@‘@*)float(float,  float)(*@’@*)
\end{console}

Il compilatore legge il file dall'alto in basso. Quando incontra
\code{maximum(max, x)} non sa ancora nulla di \code{maximum}: per una vecchia
regola del C, “immagina” che sia una funzione che restituisce un
\code{int} (è la \textit{dichiarazione implicita} di cui avvisa il
warning). Più avanti trova la definizione vera, che restituisce un
\code{float}, e la contraddizione diventa un errore. Le versioni più recenti
di gcc (dalla 14) non accettano più le dichiarazioni implicite nemmeno come
avviso: la chiamata di una funzione non dichiarata è direttamente un errore.
La lezione è semplice: \textbf{ogni funzione va dichiarata prima dell'uso}.

Un modo alternativo, che va bene per i file piccoli, è mettere le
\textbf{definizioni} delle funzioni \textbf{prima} di \code{main}: quando il
compilatore arriva alla chiamata ha già letto la definizione, e la testa
della funzione fa da prototipo.

\subsection{Un punto e virgola di troppo}

Nelle prime slide su \code{maxmin} la funzione \code{maximum} contiene un
errore di battitura istruttivo:

\begin{affianco}[0.4]
\begin{lstlisting}[style=c]
float maximum(float x, float y)
{
    if (x > y) ;{
        return x;
    }
    else {
        return y;
    }
}
\end{lstlisting}
\risultato
\begin{console}
$ gcc -Wall -o maxmin maxmin.c
maxmin.c: In function (*@‘@*)maximum(*@’@*):
maxmin.c:31:4: warning: this (*@‘@*)if(*@’@*) clause does not guard... [-Wmisleading-indentation]
   31 |    if (x > y) ;{
      |    ^~
maxmin.c:34:4: error: (*@‘@*)else(*@’@*) without a previous (*@‘@*)if(*@’@*)
   34 |    else {
      |    ^~~~
\end{console}
\end{affianco}

Il \code{;} dopo la condizione è un'\textbf{istruzione vuota}: l'\code{if}
controlla quella, cioè nulla. Il blocco \code{\{ return x; \}} che segue non
fa più parte dell'\code{if}: è un blocco qualunque, eseguito sempre. E
quando arriva l'\code{else}, l'\code{if} a cui dovrebbe appartenere è già
concluso: “\code{else} senza un \code{if} precedente”. gcc avvisa anche
che l'indentazione è \textit{ingannevole}: il codice sembra dire una cosa e
ne fa un'altra.

\section{Il passaggio dei parametri per valore}

Esistono due modi di passare i parametri a una funzione:

\begin{itemize}
  \item il \textbf{passaggio per valore}: viene creata una \textbf{copia}
        della variabile passata come parametro;
  \item il \textbf{passaggio per riferimento}: viene passata la variabile
        stessa.
\end{itemize}

\textbf{In C i parametri sono passati solo per valore.} Il programma
seguente lo dimostra: la funzione prova a cambiare il valore di \code{a},
ma \code{main} non se ne accorge.

\begin{affianco}[0.5]
\begin{lstlisting}[style=c, name=try_to.c]
#include <stdio.h>

int main(void)
{
    int    a = 1;
    void   try_to_change_it(int);

    printf("%d\n", a);     /* 1 is printed */
    try_to_change_it(a);
    printf("%d\n", a);     /* 1 is printed again! */
    return 0;
}

void try_to_change_it(int a)
{
    a = 777;
}
\end{lstlisting}
\risultato
\begin{console}
$ ./try_to
1
1
\end{console}
\wtnota{Il prototipo è scritto dentro \code{main}: è lecito, e vale solo
lì. Di solito lo si mette in cima al file.}
\end{affianco}

\begin{figure}[H]
  \centering
  \begin{tikzpicture}[font=\Lato\scriptsize]
    \foreach \col/\testo in {0/{1. prima della chiamata}, 1/{2. dentro la funzione, dopo \code{a = 777}}, 2/{3. dopo il ritorno}} {
      \node[font=\Lato\scriptsize\bfseries] at (\col*5.6+1.6,1.25) {\testo};
    }
    % 1
    \variabile{a1}{(1.6,0)}{a}{1}
    \scopo{2}{m1}{main}{(a1-nome) (a1)}
    % 2
    \begin{scope}[xshift=5.6cm]
      \variabile{a2}{(1.6,0)}{a}{1}
      \scopo{2}{m2}{main}{(a2-nome) (a2)}
      \variabile[fill=accent!10!white]{b2}{(1.6,-1.55)}{a}{777}
      \scopo{3}{f2}{try\_to\_change\_it}{(b2-nome) (b2)}
      \draw[mcerca] (a2.east) to[out=-20, in=20] node[wtlab, right, fill=none] {copia di 1} (b2.east);
    \end{scope}
    % 3
    \begin{scope}[xshift=11.2cm]
      \variabile{a3}{(1.6,0)}{a}{1}
      \scopo{2}{m3}{main}{(a3-nome) (a3)}
      \node[mnota, align=center] at (1.3,-1.55) {la copia non esiste più:\\\code{a} di \code{main} vale ancora 1};
    \end{scope}
  \end{tikzpicture}
  \caption{Il passaggio per valore: la funzione lavora su una copia, che
           scompare quando la funzione termina.}
  \label{fig:per-valore}
\end{figure}

Alla chiamata, il parametro \code{a} di \code{try\_to\_change\_it} (una
variabile diversa da quella di \code{main}, anche se ha lo stesso nome)
riceve una \textbf{copia} del valore 1. L'assegnazione \code{a = 777}
modifica la copia; quando la funzione termina la copia scompare, e
\code{main} ritrova la sua \code{a} intatta (figura~\ref{fig:per-valore}).

Ecco perché \code{scanf} ha bisogno di \code{\&x}: se ricevesse \code{x},
riceverebbe solo una copia del suo valore e non potrebbe modificare la
variabile. Ricevendo l'\textbf{indirizzo}, sa dove andare a scrivere. Come
una funzione possa usare un indirizzo per modificare una variabile del
chiamante (i \textit{puntatori}) è l'argomento del
capitolo~\ref{chap:puntatori}.

\insintesi{
  Un programma C è fatto di funzioni, e l'esecuzione comincia da
  \code{main}. Le righe che cominciano con \code{\#} sono per il
  preprocessore: \code{\#include} copia un header (che \textbf{descrive} le
  funzioni; il loro codice arriva dalla libreria al momento del linking),
  \code{\#define} definisce costanti simboliche. Le variabili si dichiarano
  con un tipo prima di usarle; \code{=} assegna (si legge da destra a
  sinistra); la divisione fra interi tronca. \code{printf} e \code{scanf}
  usano una stringa di controllo con le conversioni (\code{\%d}, \code{\%f},
  \code{\%lf}\dots); \code{scanf} vuole gli \textbf{indirizzi} delle
  variabili e restituisce il numero di conversioni riuscite. In C 0 è falso
  e ogni altro valore è vero; \code{if}, \code{while} e \code{for}
  realizzano selezione e iterazione. Le funzioni vanno dichiarate con un
  prototipo prima dell'uso, e ricevono sempre una \textbf{copia} dei
  parametri.}

\section{Riferimenti}

\begin{itemize}
  \item \textbf{A. Kelley, I. Pohl}, \textit{A Book on C}, Addison-Wesley:
        capitolo~1, da cui provengono gli esempi di questo capitolo.
  \item \textbf{Deitel \& Deitel}, \textit{Il linguaggio C}: capitoli 2--5.
  \item \textbf{printf} e \textbf{scanf} --- cppreference.com. \\
        \urlx{https://en.cppreference.com/w/c/io/fprintf} \\
        \urlx{https://en.cppreference.com/w/c/io/fscanf}
\end{itemize}


  \chapter{Elementi lessicali e operatori}
\label{chap:sintassi}

% Un token disegnato come una casella colorata secondo la sua categoria
\tikzset{
  tokbase/.style={rounded corners=2pt, line width=0.6pt, inner xsep=2.5pt,
    inner ysep=2pt, font=\FiraCodeNL\small},
  tokkw/.style={tokbase, draw=primary!70!black, fill=primary!18!white},
  tokid/.style={tokbase, draw=green-gradient-6!85!black, fill=green-gradient-1!45!white},
  tokop/.style={tokbase, draw=accent!65!black, fill=accent!12!white},
  tokpu/.style={tokbase, draw=neutral!60!black, fill=neutral!12!white},
  tokco/.style={tokbase, draw=orange-gradient-6!85!black, fill=orange-gradient-1!35!white},
}
\newcommand{\tk}[2]{\tikz[baseline=(t.base)]\node[#1] (t) {\strut #2};}

Il C è un linguaggio, e come ogni linguaggio ha un \textbf{alfabeto}, delle
regole per formare le \textbf{parole} e delle regole per comporre le parole
in \textbf{frasi} corrette, la \term{sintassi}. Il compilatore, per prima
cosa, verifica che il programma sia sintatticamente corretto: la
correttezza sintattica è necessaria per la traduzione (un programma con un
errore di sintassi non viene tradotto). Non è però sufficiente: un programma
sintatticamente perfetto può comunque fare la cosa sbagliata, come la
divisione intera della maratona nel capitolo~\ref{chap:primi-programmi}.

Questo capitolo guarda il C “al microscopio”: come il compilatore legge
il testo, quali parole sono riservate, come si scelgono i nomi, in che
ordine si calcolano le espressioni.

\section{Dai caratteri ai token}

\subsection{I token}

Un programma è una sequenza di caratteri leciti (lettere, cifre, simboli,
spazi). Il compilatore, per prima cosa, raggruppa i caratteri in unità
dotate di significato, i \term{token}, come noi raggruppiamo le lettere in
parole. Prendiamo un programma che legge due interi e ne stampa la somma:

\begin{affianco}[0.58]
\begin{lstlisting}[style=c, name=somma.c]
/* Read in two integers and print their sum. */

#include <stdio.h>

int main(void)
{
    int   a, b, sum;

    printf("Input two integers:  ");
    scanf("%d%d", &a, &b);
    sum = a + b;
    printf("%d + %d = %d\n", a, b, sum);
    return 0;
}
\end{lstlisting}
\risultato
\begin{console}
$ ./somma
Input two integers:  <<7 35>>
7 + 35 = 42
\end{console}
\end{affianco}

I token sono di alcune categorie (figura~\ref{fig:token}):

\begin{itemize}
  \item gli \textbf{identificatori}, i nomi: \code{main}, \code{a},
        \code{b}, \code{sum}, \code{printf}, \code{scanf};
  \item le \textbf{parole chiave}: \code{int}, \code{void}, \code{return};
  \item gli \textbf{operatori}: \code{=}, \code{+}, \code{\&}, e le
        parentesi \code{( )} dopo \code{main}, che dicono che \code{main} è
        una funzione;
  \item la \textbf{punteggiatura}: \code{\{ \}}, \code{;} e \code{,};
  \item le \textbf{costanti}: i numeri come \code{0} e le stringhe come
        \code{"Input two integers:  "}, che il compilatore tratta come un
        \textbf{unico} token, spazi compresi.
\end{itemize}

\begin{figure}[H]
  \centering
  \begin{minipage}{15cm}
    \centering
    \tk{tokkw}{int}\ \ \tk{tokid}{a}\tk{tokpu}{,}\ \tk{tokid}{b}\tk{tokpu}{,}\ \tk{tokid}{sum}\tk{tokpu}{;}
    \\[7pt]
    \tk{tokid}{scanf}\tk{tokop}{(}\tk{tokco}{"\%d\%d"}\tk{tokpu}{,}\ \tk{tokop}{\&}\tk{tokid}{a}\tk{tokpu}{,}\ \tk{tokop}{\&}\tk{tokid}{b}\tk{tokop}{)}\tk{tokpu}{;}
    \\[7pt]
    \tk{tokid}{sum}\ \tk{tokop}{=}\ \tk{tokid}{a}\ \tk{tokop}{+}\ \tk{tokid}{b}\tk{tokpu}{;}
    \\[7pt]
    \tk{tokkw}{return}\ \tk{tokco}{0}\tk{tokpu}{;}
    \\[12pt]
    {\Lato\scriptsize
     \tk{tokkw}{\Lato\scriptsize parola chiave}\quad
     \tk{tokid}{\Lato\scriptsize identificatore}\quad
     \tk{tokop}{\Lato\scriptsize operatore}\quad
     \tk{tokpu}{\Lato\scriptsize punteggiatura}\quad
     \tk{tokco}{\Lato\scriptsize costante}}
  \end{minipage}
  \caption{Alcune righe di \code{somma.c} scomposte in token, colorati per
           categoria.}
  \label{fig:token}
\end{figure}

\subsection{Gli spazi}

Il compilatore usa gli spazi bianchi (spazi, tabulazioni, a capo) per
distinguere i token. Dove due token si “toccherebbero” formando una parola
diversa, lo spazio è \textbf{necessario}: non si può scrivere
\code{inta, b, sum;}, perché il compilatore non riesce a staccare
\code{int} da \code{a} e vede un unico identificatore, \code{inta}, mai
dichiarato:

\begin{console}
$ gcc -Wall -c somma.c
somma.c: In function (*@‘@*)main(*@’@*):
somma.c:7:5: error: (*@‘@*)inta(*@’@*) undeclared (first use in this function); did you mean (*@‘@*)int(*@’@*)?
    7 |     inta, b, sum;
      |     ^~~~
      |     int
\end{console}

Non tutti gli spazi però sono necessari: \code{int a,b,sum;} va benissimo,
perché le virgole separano già i token. E dove servono, servono davvero:
\code{int absum;} dichiara un'unica variabile chiamata \code{absum}.

Al contrario, gli spazi \textbf{in più} sono ignorati: \code{sum=a+b;} e
\code{sum   =     a     +      b   ;} sono la stessa istruzione. Ma uno
spazio \textbf{dentro} un token lo spezza: \code{s u m = a + b;} non è
lecito, perché il compilatore tratta \code{s}, \code{u} e \code{m} come tre
token diversi.

\begin{console}
$ gcc -Wall -c somma.c
somma.c:11:5: error: unknown type name (*@‘@*)s(*@’@*)
   11 |     s u m = a + b;
      |     ^
\end{console}

Lo stesso vale per l'operatore di indirizzo: in \code{\&a} il carattere
\code{\&} è un token a sé, separato dalla variabile, quindi \code{\& a} o
\code{\&  b} sono scritture lecite (anche se insolite). Però i due argomenti
devono comunque essere separati dalla virgola: \code{\& a \& b} non è lecito,
perché viene letto come \code{\&a \& b}, cioè l'operatore \code{\&} (che fra
due operandi ha un altro significato) applicato all'indirizzo di \code{a} e
a \code{b}, e il compilatore protesta:

\begin{console}
$ gcc -Wall -c somma.c
somma.c:10:23: error: invalid operands to binary & (have (*@‘@*)int *(*@’@*) and (*@‘@*)int(*@’@*))
   10 |     scanf("%d%d", & a & b);
      |                   ~~~ ^
\end{console}

\subsection{Identificatori della libreria}

\code{printf} e \code{scanf} sono identificatori come gli altri; le
parentesi che li seguono indicano che sono funzioni. Il compilatore non ne
conosce il codice: dopo la compilazione, il linker crea l'eseguibile
prendendo la loro definizione dalla libreria standard. Non sono parole
chiave, ma andrebbero considerati \textbf{riservati}: usarli per altri scopi
è una pessima idea, come mostra l'ultimo esempio di questo capitolo.

\section{I commenti}

I \term{commenti} sono sequenze arbitrarie di testo fra i delimitatori
\code{/*} e \code{*/}, oppure da \code{//} alla fine della riga. I commenti
\textbf{non sono token}: il compilatore trasforma ogni commento in uno
spazio. Lo si può verificare con un esempio un po' estremo, in cui un
commento separa \code{int} da \code{x}:

\begin{affianco}[0.6][c]
\begin{lstlisting}[style=c, name=commenti.c]
#include <stdio.h>

int main(void)
{
    int/* un commento vale come uno spazio */x = 5;

    printf("%d\n", x);   // un commento su una riga
    return 0;
}
\end{lstlisting}
\risultato
\begin{console}
$ ./commenti
5
\end{console}
\end{affianco}

I commenti \code{/* */} \textbf{non si annidano}: il primo \code{*/}
chiude il commento, qualunque cosa ci sia prima. Con
\codel{/* esterno /* interno */ ancora esterno? */} il commento finisce dopo
\code{interno}, e “ancora esterno?” viene letto come codice (gcc avvisa:
\textit{"/*" within comment}).

I commenti sono la documentazione del codice. Andrebbero scritti
\textbf{insieme} al codice: una volta che il programma funziona si tende a
scriverne di brevissimi, o a non scriverne affatto. Un buon commento spiega
il \textit{perché}, non ripete il \textit{che cosa}: \code{i = i + 1;
/* incrementa i */} non aiuta nessuno, \code{/* il primo elemento è già
stato letto */} sì.

\section{Le parole chiave}

Le \term{parole chiave} sono parole riservate, che non possono essere usate
in altri contesti (per esempio come nomi di variabili). Il C ne ha, nel
complesso, poche; l'ANSI C ne definisce 32:

\begin{center}
\renewcommand{\arraystretch}{1.25}
{\FiraCodeNL\small
\begin{tabular}{@{}*{8}{l}@{}}
  \toprule
  auto & break & case & char & const & continue & default & do \\
  double & else & enum & extern & float & for & goto & if \\
  int & long & register & return & short & signed & sizeof & static \\
  struct & switch & typedef & union & unsigned & void & volatile & while \\
  \bottomrule
\end{tabular}}
\end{center}

\warning{Le parole chiave sono minuscole}{
  Nella tabella delle slide compaiono \code{Volatile} e \code{Unsigned} con
  l'iniziale maiuscola: è un effetto della correzione automatica. Le parole
  chiave sono tutte \textbf{minuscole}, e siccome il C distingue maiuscole e
  minuscole, \code{Unsigned} non sarebbe una parola chiave ma un
  identificatore qualunque.}

Alcune implementazioni hanno parole chiave aggiuntive (gcc, per esempio,
riconosce \code{asm} e \code{typeof}), e le versioni successive dello
standard ne hanno aggiunte poche altre (fra cui \code{inline}, e nel C23
\code{bool}, \code{true} e \code{false}). Non serve impararle a memoria:
VS Code le colora in blu o in viola, e l'errore di usarne una come nome si
vede subito.

\section{Gli identificatori}

\subsection{La regola}

Un \term{identificatore} (un nome di variabile, di funzione, di tipo\dots)
è una sequenza di caratteri composta da \textbf{lettere}, \textbf{cifre} e
dal carattere speciale \code{\_} (\textit{underscore}), che \textbf{non
comincia con una cifra}. Formalmente, nella notazione delle slide:

\begin{center}
  \FiraCodeNL identifier ::= \{letter | underscore\}$_1$ \{letter | underscore | digit\}$_{0+}$
\end{center}

dove \code{|} significa “oppure”, il pedice $_1$ “esattamente una volta”
e il pedice $_{0+}$ “zero o più volte”. La stessa regola si può disegnare
come un percorso da seguire, un \textit{diagramma sintattico}
(figura~\ref{fig:sintattico}): un nome è lecito se lo si può ottenere
percorrendo le frecce dall'inizio alla fine e raccogliendo un carattere a
ogni casella attraversata.

\begin{figure}[H]
  \centering
  \begin{tikzpicture}[font=\Lato\scriptsize,
      cas/.style={rounded corners=6pt, draw=primary!65!black, fill=primary!8!white,
        line width=0.8pt, minimum width=1.6cm, minimum height=0.5cm},
      bin/.style={line width=0.9pt, draw=neutral!35!black},
      fr/.style={bin, -{Stealth[length=5pt,width=4pt]}}]
    % binario principale
    \draw[fr] (0,0) -- (1.0,0);
    \node[anchor=east, font=\Lato\scriptsize\bfseries] at (0,0) {inizio};
    % prima casella: lettera oppure _
    \node[cas] (l1) at (2.4,0.45) {lettera};
    \node[cas] (u1) at (2.4,-0.45) {\code{\_}};
    \draw[bin] (1.0,0) to[out=0, in=180] (l1.west);
    \draw[bin] (1.0,0) to[out=0, in=180] (u1.west);
    \draw[bin] (l1.east) to[out=0, in=180] (3.8,0);
    \draw[bin] (u1.east) to[out=0, in=180] (3.8,0);
    % binario dritto: zero ripetizioni
    \draw[fr] (3.8,0) -- (11.2,0);
    \node[anchor=west, font=\Lato\scriptsize\bfseries] at (11.2,0) {fine};
    % ciclo delle ripetizioni
    \node[cas] (l2) at (7.5,-0.9) {lettera};
    \node[cas] (u2) at (7.5,-1.6) {\code{\_}};
    \node[cas] (d2) at (7.5,-2.3) {cifra};
    \draw[bin] (4.6,0) to[out=0, in=180] (5.6,-0.9);
    \foreach \n in {l2,u2,d2} {
      \draw[bin] (5.6,-0.9) to[out=-90, in=180] (\n.west);
      \draw[bin] (\n.east) to[out=0, in=-90] (9.4,-0.9);
    }
    \draw[fr] (9.4,-0.9) to[out=0, in=180] (10.4,0);
    \draw[fr] (9.4,-0.9) -- (9.4,-2.85) -- (5.6,-2.85) -- (5.6,-1.0);
    \node[mnota, anchor=north] at (7.5,-2.9) {si può ripetere quante volte si vuole};
    \node[mnota, anchor=south] at (7.5,0.05) {oppure nessuna volta};
    \node[mnota, anchor=north] at (2.4,-0.8) {esattamente una volta};
  \end{tikzpicture}
  \caption{Il diagramma sintattico di un identificatore.}
  \label{fig:sintattico}
\end{figure}

\begin{center}
\renewcommand{\arraystretch}{1.3}
\begin{tabular}{@{}>{\FiraCodeNL}l l@{\hspace{1.4cm}}>{\FiraCodeNL}l l@{}}
  \toprule
  \multicolumn{2}{@{}l}{\textbf{identificatori leciti}} & \multicolumn{2}{l@{}}{\textbf{identificatori non leciti}} \\
  \midrule
  K                & una lettera sola      & not\#me     & \code{\#} non è ammesso \\
  \_id             & comincia con \code{\_} & 101\_south & comincia con una cifra \\
  iamanidentifier2 & cifre dopo l'inizio   & -plus       & \code{-} è un operatore \\
  so\_am\_i        &                       & int         & è una parola chiave \\
  \bottomrule
\end{tabular}
\end{center}

\subsection{Maiuscole, lunghezza, scelta dei nomi}

Il C è \textit{case sensitive}: le lettere maiuscole e minuscole sono
considerate \textbf{diverse}, quindi \code{sum}, \code{Sum} e \code{SUM} sono
tre identificatori distinti.

Sulla lunghezza: nei sistemi più vecchi contavano solo i primi 8 caratteri,
per cui \codel{i\_am\_an\_identifier} e \codel{i\_am\_an\_elephant} sarebbero
stati lo stesso nome (i primi 8 caratteri sono \code{i\_am\_an\_} in
entrambi). L'ANSI C garantisce che siano considerati almeno i primi 31
caratteri, e i compilatori moderni ne considerano molti di più: in pratica,
la lunghezza non è un problema.

Gli identificatori vanno scelti in modo che contribuiscano alla
\textbf{leggibilità} e alla \textbf{documentazione} del programma. Per
calcolare l'IVA di un prodotto,

\begin{lstlisting}[style=c]
float tax_iva, price, tax_iva_rate;

tax_iva = price * tax_iva_rate;
\end{lstlisting}

si capisce da sé, mentre \code{t = p * r;} richiede un commento per ogni
variabile. Infine, per convenzione gli identificatori di sistema (della
libreria, del compilatore) cominciano con uno o due underscore: per evitare
conflitti, conviene non creare identificatori che cominciano con \code{\_}.

\section{Gli operatori}

\subsection{Precedenza e associatività}

Gli operatori hanno regole di \term{precedenza} e di \term{associatività}
che determinano come vengono calcolate le espressioni. Le espressioni fra
parentesi sono calcolate per prime; le parentesi si possono usare sia per
\textbf{chiarire} l'ordine delle operazioni, sia per \textbf{modificarlo}.

\begin{itemize}
  \item In \code{1 + 2 * 3} l'operatore \code{*} ha una precedenza più alta
        di \code{+}: il risultato è 7, come in \code{1 + (2 * 3)}.
  \item In \code{(1 + 2) * 3} le parentesi hanno la priorità: il risultato è
        9.
  \item In \code{1 + 2 - 3 + 4 - 5} gli operatori binari \code{+} e
        \code{-} hanno la \textbf{stessa} precedenza: l'ordine lo stabilisce
        la regola di associatività, che per loro è \textbf{da sinistra a
        destra}. L'espressione equivale a \code{(((1 + 2) - 3) + 4) - 5},
        cioè $-1$.
\end{itemize}

La figura~\ref{fig:alberi} mostra gli alberi delle tre espressioni:
precedenza e associatività servono proprio a stabilire, senza ambiguità,
quale albero corrisponde a una sequenza di token.

\begin{figure}[H]
  \centering
  \begin{tikzpicture}[font=\FiraCodeNL\small,
      op/.style={circle, draw=accent!65!black, fill=accent!8!white, line width=0.8pt,
        minimum size=0.58cm, inner sep=0pt, font=\FiraCodeNL\small\bfseries},
      nu/.style={font=\FiraCodeNL\small},
      val/.style={font=\Lato\tiny\bfseries, text=green-gradient-6!50!black, anchor=west},
      tit/.style={font=\FiraCodeNL\scriptsize\bfseries, text=primary!55!black}]
    % 1 + 2 * 3
    \node[tit] at (0.6,0.8) {1 + 2 * 3};
    \node[op] (a) at (0,0) {+};
    \node[nu] (a1) at (-0.7,-0.9) {1};
    \node[op] (a2) at (0.7,-0.9) {*};
    \node[nu] (a3) at (0.2,-1.8) {2};
    \node[nu] (a4) at (1.2,-1.8) {3};
    \foreach \x/\y in {a/a1, a/a2, a2/a3, a2/a4} \draw[fsramo] (\x) -- (\y);
    \node[val] at ([xshift=3pt]a.east) {7};
    \node[val] at ([xshift=3pt]a2.east) {6};
    % (1 + 2) * 3
    \begin{scope}[xshift=4.2cm]
      \node[tit] at (0.2,0.8) {(1 + 2) * 3};
      \node[op] (b) at (0,0) {*};
      \node[op] (b1) at (-0.7,-0.9) {+};
      \node[nu] (b2) at (0.7,-0.9) {3};
      \node[nu] (b3) at (-1.2,-1.8) {1};
      \node[nu] (b4) at (-0.2,-1.8) {2};
      \foreach \x/\y in {b/b1, b/b2, b1/b3, b1/b4} \draw[fsramo] (\x) -- (\y);
      \node[val] at ([xshift=3pt]b.east) {9};
      \node[val] at ([xshift=3pt]b1.east) {3};
    \end{scope}
    % 1 + 2 - 3 + 4 - 5
    \begin{scope}[xshift=9.6cm]
      \node[tit] at (0.4,0.8) {1 + 2 - 3 + 4 - 5};
      \node[op] (c) at (0,0) {-};
      \node[op] (c1) at (-0.7,-0.75) {+};
      \node[op] (c2) at (-1.4,-1.5) {-};
      \node[op] (c3) at (-2.1,-2.25) {+};
      \node[nu] (n1) at (-2.6,-3.0) {1};
      \node[nu] (n2) at (-1.6,-3.0) {2};
      \node[nu] (n3) at (-0.9,-2.25) {3};
      \node[nu] (n4) at (-0.2,-1.5) {4};
      \node[nu] (n5) at (0.5,-0.75) {5};
      \foreach \x/\y in {c/c1, c/n5, c1/c2, c1/n4, c2/c3, c2/n3, c3/n1, c3/n2} \draw[fsramo] (\x) -- (\y);
      \node[val] at ([xshift=3pt]c.east) {$-1$};
      \node[val] at ([xshift=3pt]c1.east) {4};
      \node[val] at ([xshift=3pt]c2.east) {0};
      \node[val] at ([xshift=3pt]c3.east) {3};
    \end{scope}
  \end{tikzpicture}
  \caption{Gli alberi delle tre espressioni, con il valore di ogni
           operazione: si calcola dal basso verso l'alto.}
  \label{fig:alberi}
\end{figure}

\subsection{La tabella degli operatori}

La tabella seguente riassume precedenza e associatività degli operatori
visti finora, dalla precedenza più alta (in alto) alla più bassa. Gli
operatori sulla stessa riga hanno la stessa precedenza.

\begin{center}
\renewcommand{\arraystretch}{1.35}
\begin{tabular}{@{}>{\FiraCodeNL}l l l@{}}
  \toprule
  \normalfont\textbf{operatori} & \textbf{significato} & \textbf{associatività} \\
  \midrule
  () \ ++ \ -{}- & parentesi; incremento e decremento \textbf{postfissi} & da sinistra a destra \\
  + \ - \ ++ \ -{}- & più e meno \textbf{unari}; incremento e decremento \textbf{prefissi} & da destra a sinistra \\
  * \ / \ \% & moltiplicazione, divisione, resto & da sinistra a destra \\
  + \ - & somma e sottrazione & da sinistra a destra \\
  = \ += \ -= \ *= \ /= & assegnazioni & da destra a sinistra \\
  \bottomrule
\end{tabular}
\end{center}

Qualche commento:

\begin{itemize}
  \item gli operatori \textbf{unari} hanno un solo operando: in \code{-x} il
        meno è unario (cambia segno), in \code{a - x} è binario (sottrae).
        Gli unari hanno precedenza più alta dei binari: \code{-a * b} è
        \code{(-a) * b};
  \item \code{\%} è il \term{resto} della divisione fra interi (detto anche
        \textit{modulo}): \code{7 / 2} vale 3 e \code{7 \% 2} vale 1, perché
        $7 = 3 \cdot 2 + 1$. Si usa spesso per sapere se un numero è pari
        (\code{n \% 2 == 0}) o divisibile per un altro;
  \item le assegnazioni hanno la precedenza più bassa (si calcola prima
        tutto il lato destro) e associano da \textbf{destra} a sinistra: è
        il motivo per cui \code{min = max = sum = avg = x;} del
        capitolo~\ref{chap:primi-programmi} funziona, perché equivale a
        \code{min = (max = (sum = (avg = x)));}.
\end{itemize}

\begin{affianco}[0.62][c]
\begin{lstlisting}[style=c]
printf("%d %d %d\n", 1 + 2 * 3, (1 + 2) * 3, 1 + 2 - 3 + 4 - 5);
printf("%d %d %d\n", 7 / 2, 7 % 2, -7 / 2);
\end{lstlisting}
\risultato
\begin{console}
7 9 -1
3 1 -3
\end{console}
\end{affianco}

Si noti l'ultimo valore: la divisione fra interi \textbf{tronca verso lo
zero}, quindi $-7/2 = -3{,}5$ diventa $-3$ (non $-4$).

\info{Nel dubbio, parentesi}{
  Nessuno ricorda a memoria tutta la tabella delle precedenze del C (che
  conta una quindicina di livelli). Quando l'ordine non è ovvio, le
  parentesi lo rendono esplicito per il compilatore e, soprattutto, per chi
  legge: \code{(a * b) + c} non costa nulla e non lascia dubbi.}

\section{Incremento e decremento}

\subsection{Prefisso e postfisso}

Gli operatori di incremento \code{++} e decremento \code{-{}-} sono
\textbf{unari} e si possono usare in due modi: come operatore
\textbf{prefisso} (\code{++i}, prima della variabile) o \textbf{postfisso}
(\code{i++}, dopo). Si applicano a \textbf{variabili}, non a costanti o a
espressioni: \code{++i} e \code{cnt-{}-} sono usi leciti, \code{777++} e
\code{++(a * b - 1)} no, perché non c'è una variabile da modificare. Il
compilatore lo dice così (\textit{lvalue}, “valore a sinistra”, è
qualcosa che può stare a sinistra di un'assegnazione, cioè una variabile):

\begin{console}
$ gcc -Wall -c costanti.c
costanti.c:4:8: error: lvalue required as increment operand
costanti.c:5:5: error: lvalue required as increment operand
\end{console}

Le due forme non hanno lo stesso comportamento, e per capire la differenza
bisogna ricordare che in C \textbf{ogni espressione ha un valore}:

\begin{itemize}
  \item \code{++i}: la variabile \code{i} viene incrementata di 1, e
        l'espressione vale il \textbf{nuovo} valore di \code{i} (“prima
        incrementa, poi usa”);
  \item \code{i++}: l'espressione vale il \textbf{vecchio} valore di
        \code{i}, e poi la variabile viene incrementata (“prima usa, poi
        incrementa”).
\end{itemize}

Allo stesso modo \code{-{}-i} decrementa e vale il nuovo valore,
\code{i-{}-} vale il vecchio valore e poi decrementa. Come istruzioni a sé
(\code{i++;}) le due forme sono equivalenti, perché il valore
dell'espressione non viene usato da nessuno.

\subsection{Un esempio da simulare}

\begin{affianco}[0.48][c]
\begin{lstlisting}[style=c, name=incr.c]
#include <stdio.h>

int main(void)
{
    int a, b, c = 0;

    a = ++c;
    b = c++;
    printf("%d %d %d\n", a, b, ++c);
    return 0;
}
\end{lstlisting}
\risultato
\begin{console}
$ ./incr
1 1 3
\end{console}
\wtnota{Nelle slide la dichiarazione è scritta \code{int a, b, c==0;}: è
un refuso, \code{==} è il confronto; serve \code{=}.}
\end{affianco}

Simuliamo l'esecuzione, istruzione per istruzione:

\begin{center}
\renewcommand{\arraystretch}{1.35}
\begin{tabular}{@{}>{\FiraCodeNL\small}l>{\raggedright\arraybackslash}p{6.6cm}ccc@{}}
  \toprule
  \normalfont\textbf{istruzione} & \textbf{che cosa succede} & \textbf{\code{a}} & \textbf{\code{b}} & \textbf{\code{c}} \\
  \midrule
  int a, b, c = 0; & dichiarazioni, solo \code{c} è inizializzata & ? & ? & 0 \\
  a = ++c; & \code{c} diventa 1; \code{++c} vale 1, che va in \code{a} & 1 & ? & 1 \\
  b = c++; & \code{c++} vale 1 (il vecchio valore), che va in \code{b}; poi \code{c} diventa 2 & 1 & 1 & 2 \\
  printf(\dots, ++c); & \code{c} diventa 3; \code{++c} vale 3 e viene stampato & 1 & 1 & 3 \\
  \bottomrule
\end{tabular}
\end{center}

\subsection{Quando \code{++} diventa pericoloso}

\code{++} e \code{-{}-} \textbf{modificano} la variabile su cui operano.
Usati da soli non danno problemi; in espressioni più complesse bisogna
prestare molta attenzione. Le slide mostrano questo programma, con l'intento
di calcolare $a \cdot (a+1)$ in tre modi:

\begin{lstlisting}[style=c, name=ub.c]
#include <stdio.h>

int main(void)
{
    int a = 1;

    printf("a*(a+1) = %d\n", a * (a + 1));
    printf("a*(a+1) = %d\n", a * a++);
    a = 1;
    printf("a*(a+1) = %d\n", a * ++a);
    return 0;
}
\end{lstlisting}

Ecco che cosa succede compilandolo con gcc (su Linux) e con clang (il
compilatore di macOS):

\begin{affianco}[0.5]
\begin{console}[gcc]
$ gcc -Wall -o ub ub.c
ub.c:8:35: warning: operation on (*@‘@*)a(*@’@*) may be undefined [-Wsequence-point]
ub.c:10:34: warning: operation on (*@‘@*)a(*@’@*) may be undefined [-Wsequence-point]
$ ./ub
a*(a+1) = 2
a*(a+1) = 2
a*(a+1) = 4
\end{console}
\risultato
\begin{console}[clang]
$ clang -Wall -o ub ub.c
ub.c:8:35: warning: unsequenced modification and access to 'a' [-Wunsequenced]
ub.c:10:34: warning: unsequenced modification and access to 'a' [-Wunsequenced]
$ ./ub
a*(a+1) = 2
a*(a+1) = 1
a*(a+1) = 2
\end{console}
\end{affianco}

\textbf{Lo stesso programma dà risultati diversi con compilatori diversi}, e
nessuno dei due ha torto. In \code{a * a++} il C non stabilisce se il primo
\code{a} vada letto prima o dopo l'incremento: il programma legge e modifica
la stessa variabile nella stessa espressione, senza un ordine definito, e lo
standard dice che in questo caso il comportamento è \textbf{indefinito}
(\textit{undefined behavior}): il compilatore può fare qualunque cosa. Le
slide mostrano i risultati di gcc (2, 2, 4) come se fossero “quelli
giusti”; in realtà non esistono risultati giusti, e infatti entrambi i
compilatori avvisano.

\error{Una variabile modificata con \code{++} non va usata altrove nella stessa espressione}{
  Espressioni come \code{a * a++}, \code{i = i++} o
  \code{printf("\%d \%d", n, n++)} hanno comportamento indefinito. La regola
  pratica delle slide è la migliore: \textbf{meglio evitarle} e scrivere
  esplicitamente ciò che si intende, per esempio \code{a * (a + 1)}. E i
  warning di \code{-Wall} vanno sempre letti.}

\section{La libreria standard e i nomi riservati}

Che cosa succede se si \textbf{ridefinisce} una funzione della libreria
standard? Le slide lo provano con \code{scanf}, scrivendo una funzione con lo
stesso nome che stampa la stringa ricevuta. Per non avere conflitti con il
prototipo vero, il file \textbf{non} include \code{stdio.h}:

\begin{lstlisting}[style=c, name=main.c]
void scanf(const char *s)
{
    printf(s);
}

int main()
{
    scanf("The printing scanf\n");
    return 0;
}
\end{lstlisting}

\begin{console}
$ gcc -o ridef main.c
main.c:1:6: warning: conflicting types for built-in function (*@‘@*)scanf(*@’@*); expected (*@‘@*)int(const char *, ...)(*@’@*) [-Wbuiltin-declaration-mismatch]
main.c:3:5: warning: implicit declaration of function (*@‘@*)printf(*@’@*) [-Wimplicit-function-declaration]
main.c:3:5: warning: incompatible implicit declaration of built-in function (*@‘@*)printf(*@’@*)
main.c:3:5: warning: format not a string literal and no format arguments [-Wformat-security]
$ ./ridef
The printing scanf
\end{console}

Il risultato è istruttivo:

\begin{itemize}
  \item anche senza \code{stdio.h}, il compilatore si accorge che si sta
        usando il nome di una funzione “predefinita” (\textit{built-in})
        con un tipo diverso da quello che conosce;
  \item dà però solo dei \textbf{warning}, non degli errori: l'eseguibile
        viene creato;
  \item e viene usata la \textbf{nuova} definizione: la “\code{scanf}”
        stampa invece di leggere.
\end{itemize}

Le slide mostrano anche la variante con la funzione in un file separato
(\code{scanf.c}) e il solo prototipo in \code{main.c}: il risultato è lo
stesso, perché il linker, trovando \code{scanf} già definita nei nostri file
oggetto, non va a cercarla nella libreria. (Si notino anche i warning su
\code{printf}: senza \code{stdio.h} il compilatore non ne conosce il
prototipo. Le versioni di gcc dalla 14 in poi considerano questa
dichiarazione implicita un errore, e il programma non si compila più.)

La morale: i nomi delle funzioni della libreria standard \textbf{non sono
parole chiave}, quindi il compilatore non ne impedisce l'uso, ma vanno
trattati come se lo fossero. Un programma che ridefinisce \code{scanf} fa
qualcosa di diverso da ciò che chiunque si aspetta leggendolo.

\insintesi{
  Il compilatore legge il programma come una sequenza di \textbf{token}:
  identificatori, parole chiave, operatori, punteggiatura, costanti; gli
  spazi separano i token quando serve e altrimenti sono ignorati, i
  commenti valgono come spazi. Le 32 parole chiave dell'ANSI C sono
  riservate e minuscole; un identificatore è fatto di lettere, cifre e
  \code{\_} e non comincia con una cifra, e maiuscole e minuscole contano.
  \textbf{Precedenza} e \textbf{associatività} stabiliscono l'albero di
  un'espressione (le parentesi lo rendono esplicito). \code{++i} vale il
  nuovo valore, \code{i++} il vecchio; leggere e modificare la stessa
  variabile nella stessa espressione dà un comportamento
  \textbf{indefinito}. I nomi della libreria standard vanno trattati come
  riservati.}

\section{Riferimenti}

\begin{itemize}
  \item \textbf{A. Kelley, I. Pohl}, \textit{A Book on C}: capitolo~2
        (elementi lessicali, operatori e il sistema C).
  \item \textbf{C Operator Precedence} --- cppreference.com. \\
        \urlx{https://en.cppreference.com/w/c/language/operator_precedence}
  \item \textbf{Order of evaluation} (comportamento indefinito di
        \code{i++ + i++} e simili) --- cppreference.com. \\
        \urlx{https://en.cppreference.com/w/c/language/eval_order}
\end{itemize}


  \chapter{Programmi su più file e make}
\label{chap:make}

Finora ogni programma stava in un solo file. I programmi veri sono fatti di
molti file, a volte migliaia, e compilarli tutti a ogni modifica sarebbe
lentissimo. Questo capitolo mostra come si divide un programma in più file,
come si compilano separatamente e si collegano (riprendendo il
capitolo~\ref{chap:dal-problema}), e come il programma \code{make} automatizza
il lavoro, ricompilando solo ciò che serve.

\section{Perché più file}

Dividere un programma in più file sorgente ha diversi vantaggi:

\begin{itemize}
  \item \textbf{organizzazione}: ogni file raccoglie le funzioni che
        riguardano un aspetto del problema (l'input, i calcoli, l'output),
        come i capitoli di un libro;
  \item \textbf{riuso}: un file con funzioni utili (per esempio di
        geometria) si può usare in più programmi;
  \item \textbf{lavoro di gruppo}: persone diverse lavorano su file
        diversi;
  \item \textbf{compilazione più rapida}: se si modifica un solo file, si
        ricompila solo quello e si rifà il collegamento.
\end{itemize}

\subsection{Un piccolo progetto}

Prendiamo un programma che calcola l'area di un cerchio e di un quadrato, e
dividiamolo in tre file: le funzioni di calcolo in \code{aree.c}, il
programma principale in \code{main.c}, e un \textbf{file di intestazione}
\code{aree.h} con ciò che serve a entrambi.

\begin{lstlisting}[style=h, name=aree.h]
/* aree.h: costanti e prototipi delle funzioni sulle aree */

#define PI 3.14159

double area_cerchio(double raggio);
double area_quadrato(double lato);
\end{lstlisting}

\begin{affianco}[0.46]
\begin{lstlisting}[style=c, name=aree.c]
/* aree.c: definizioni delle funzioni sulle aree */

#include "aree.h"

double area_cerchio(double raggio)
{
    return PI * raggio * raggio;
}

double area_quadrato(double lato)
{
    return lato * lato;
}
\end{lstlisting}
\risultato
\begin{lstlisting}[style=c, name=main.c]
#include <stdio.h>
#include "aree.h"

int main(void)
{
    double x;

    printf("Inserisci una misura: ");
    scanf("%lf", &x);
    printf("cerchio di raggio %g: area %g\n",
           x, area_cerchio(x));
    printf("quadrato di lato %g: area %g\n",
           x, area_quadrato(x));
    return 0;
}
\end{lstlisting}
\end{affianco}

Il ruolo di ciascun file:

\begin{itemize}
  \item \code{aree.h} contiene le \textbf{dichiarazioni}: la costante e i
        prototipi. Chi lo include sa \textit{che} le funzioni esistono e come
        si chiamano;
  \item \code{aree.c} contiene le \textbf{definizioni}: il codice delle
        funzioni, scritto una volta sola. Include anche lui \code{aree.h},
        così il compilatore controlla che le definizioni corrispondano ai
        prototipi;
  \item \code{main.c} include \code{aree.h} e può quindi chiamare le
        funzioni, anche se il loro codice sta in un altro file.
\end{itemize}

\subsection{Compilare separatamente, collegare insieme}

Ogni file \code{.c} si compila per conto suo con l'opzione \code{-c}, che
produce il file oggetto; poi il linker unisce i file oggetto
nell'eseguibile:

\begin{console}
$ gcc -Wall -c main.c
$ gcc -Wall -c aree.c
$ gcc -o aree main.o aree.o
$ ./aree
Inserisci una misura: <<2>>
cerchio di raggio 2: area 12.5664
quadrato di lato 2: area 4
\end{console}

Che cosa contengono i due file oggetto, e come fa il linker a metterli
insieme? Lo mostra \code{nm} (capitolo~\ref{chap:dal-problema}): in
\code{main.o} le due funzioni sono \textit{undefined} (\code{U}), cioè usate
ma senza codice; in \code{aree.o} sono presenti (\code{T}). Il linker
abbina ogni nome mancante con la sua definizione (schema a destra).

\begin{affianco}[0.5][c]
\begin{console}
$ nm main.o
                 U __isoc99_scanf
                 U __stack_chk_fail
                 U area_cerchio
                 U area_quadrato
0000000000000000 T main
                 U printf
$ nm aree.o
0000000000000000 T area_cerchio
0000000000000025 T area_quadrato
\end{console}
\risultato
\begin{immagine}
\begin{tikzpicture}[font=\FiraCodeNL\scriptsize,
    sim/.style={anchor=west, inner sep=1.5pt},
    box/.style={draw=primary!65!black, fill=primary!4!white, line width=0.7pt,
      rounded corners=3pt}]
  \draw[box] (0,0) rectangle (2.9,-2.3);
  \node[font=\Lato\scriptsize\bfseries, anchor=south west] at (0,0.02) {main.o};
  \node[sim] at (0.1,-0.35) {T main};
  \node[sim, text=merr] (u1) at (0.1,-0.85) {U area\_cerchio};
  \node[sim, text=merr] (u2) at (0.1,-1.35) {U area\_quadrato};
  \node[sim, text=neutral!45!black] at (0.1,-1.85) {U printf, scanf\dots};
  \draw[box] (4.6,-0.4) rectangle (7.6,-1.8);
  \node[font=\Lato\scriptsize\bfseries, anchor=south west] at (4.6,-0.38) {aree.o};
  \node[sim, text=mok] (t1) at (4.7,-0.85) {T area\_cerchio};
  \node[sim, text=mok] (t2) at (4.7,-1.35) {T area\_quadrato};
  \draw[wtarrowacc] (u1.east) -- (t1.west);
  \draw[wtarrowacc] (u2.east) -- (t2.west);
  \node[mnota, anchor=north, align=center] at (3.8,-2.4) {il linker abbina ogni nome usato (\code{U})\\alla sua definizione (\code{T}); \code{printf}\\e \code{scanf} le trova nella libreria standard};
\end{tikzpicture}
\end{immagine}
\end{affianco}

Che cosa succede se si dimentica \code{aree.o} al momento del linking? La
compilazione di \code{main.c} era già riuscita (grazie a \code{aree.h} il
compilatore sapeva tutto ciò che serviva), ma il linker non trova il codice
delle due funzioni:

\begin{console}
$ gcc -o aree main.o
/usr/bin/ld: main.o: in function `main':
main.c:(.text+0x54): undefined reference to `area_cerchio'
/usr/bin/ld: main.c:(.text+0x89): undefined reference to `area_quadrato'
collect2: error: ld returned 1 exit status
\end{console}

È lo stesso errore di \code{pow} senza \code{-lm} del
capitolo~\ref{chap:primi-programmi}: \textit{undefined reference} significa
sempre “il linker non trova il codice di questa funzione”.

\section{Il progetto delle lezioni}

\subsection{I file}

A lezione si usa un progetto su tre file, che gestisce un elenco di esami
sostenuti da alcuni studenti (letti da un file di testo): \code{main.c},
\code{funzioni.c} e \code{studente.h} (figura~\ref{fig:progetto}). Il codice
usa strumenti che vedremo più avanti (strutture, liste, file), e verrà
spiegato per intero nel capitolo~\ref{chap:liste}; per capire \code{make}
basta sapere chi dipende da chi.

\begin{figure}[H]
  \centering
  \begin{tikzpicture}[font=\Lato\scriptsize,
      f/.style={wtnodeplain, minimum width=3.2cm, minimum height=1.15cm, align=center}]
    \node[f, draw=accent!60!black, fill=accent!7!white] (h) at (0,0) {{\FiraCodeNL\bfseries studente.h}\\la struttura \code{studente}\\e i prototipi delle funzioni};
    \node[f] (m) at (-3.6,-1.9) {{\FiraCodeNL\bfseries main.c}\\legge i dati, filtra,\\scrive il risultato};
    \node[f] (fu) at (3.6,-1.9) {{\FiraCodeNL\bfseries funzioni.c}\\le funzioni sulla lista\\di studenti};
    \draw[wtarrowacc] (h) -- node[wtlab, above left, fill=none] {incluso da} (m);
    \draw[wtarrowacc] (h) -- node[wtlab, above right, fill=none] {incluso da} (fu);
    \draw[wtarrow, dashed] (m) -- node[wtlab, below=2pt, fill=none] {chiama le funzioni di} (fu);
  \end{tikzpicture}
  \caption{Il progetto delle lezioni: due sorgenti e un header condiviso.}
  \label{fig:progetto}
\end{figure}

L'eseguibile si crea come abbiamo visto: la compilazione
(\code{gcc -c main.c funzioni.c}, che si può fare anche con un solo comando
per più file) produce \code{main.o} e \code{funzioni.o}; il linking
(\code{gcc -o esercizio main.o funzioni.o}) produce l'eseguibile
\code{esercizio}.

\subsection{Il problema}

Con molti file la compilazione può essere lunga, e non sempre è necessario
ricompilare tutti i sorgenti: se si modifica solo \code{funzioni.c},
\code{main.o} è ancora valido. Vorremmo rendere le operazioni più semplici
e automatiche.

Un primo passo è uno \term{script}: un file di testo che contiene i comandi,
da eseguire tutti in fila.

\begin{affianco}[0.5][c]
\begin{lstlisting}[style=shell, name=compila.sh]
#!/bin/bash
gcc -c main.c
gcc -c funzioni.c
gcc -o esercizio main.o funzioni.o
\end{lstlisting}
\risultato
\begin{console}
$ chmod +x compila.sh
$ ./compila.sh
\end{console}
\end{affianco}

La prima riga (\code{\#!/bin/bash}, detta \textit{shebang}) dice al sistema
quale programma deve eseguire lo script: la shell bash. \code{chmod +x}
rende il file eseguibile, una volta per tutte; da lì in poi
\code{./compila.sh} esegue i tre comandi. Il procedimento è automatico, ma i
sorgenti vengono comunque compilati \textbf{tutti}, ogni volta. Vorremmo un
meccanismo che compila solo i file che sono stati modificati.

\section{make}

\subsection{L'idea: dipendenze e date}

Il comando \code{make} esegue la compilazione in modo automatico seguendo,
fondamentalmente, due regole:

\begin{itemize}
  \item la \textbf{dipendenza} fra i file: \code{esercizio} si costruisce da
        \code{main.o} e \code{funzioni.o}; \code{main.o} si costruisce da
        \code{main.c} (e da \code{studente.h}, che \code{main.c} include), e
        così via (figura~\ref{fig:dipendenze});
  \item l'\textbf{ordine temporale} in cui sono stati creati o modificati i
        file sorgente, oggetto ed eseguibile: ogni file ha una data e ora di
        ultima modifica, che \code{ls -l} mostra.
\end{itemize}

\begin{figure}[H]
  \centering
  \begin{tikzpicture}[font=\Lato\scriptsize,
      t/.style={wtnode, font=\FiraCodeNL\scriptsize, minimum width=2.3cm},
      s/.style={wtnodeplain, font=\FiraCodeNL\scriptsize, minimum width=2.3cm},
      ora/.style={font=\Lato\tiny, text=neutral!40!black}]
    \node[t, draw=red-gradient-6!70!black, fill=red-gradient-1!20!white] (e) at (0,0) {esercizio};
    \node[t] (mo) at (-3.0,-1.5) {main.o};
    \node[t] (fo) at (3.0,-1.5) {funzioni.o};
    \node[s] (mc) at (-4.6,-3.0) {main.c};
    \node[s, draw=accent!60!black, fill=accent!7!white] (h) at (0,-3.0) {studente.h};
    \node[s] (fc) at (4.6,-3.0) {funzioni.c};
    \foreach \a/\b in {e/mo, e/fo, mo/mc, mo/h, fo/h, fo/fc}
      \draw[wtarrow, <-] (\a) -- (\b);
    \node[ora, anchor=west] at ([xshift=3pt]e.east) {modificato alle 10:07};
    \node[ora, anchor=east] at ([xshift=-3pt]mo.west) {10:05};
    \node[ora, anchor=west] at ([xshift=3pt]fo.east) {10:06};
    \node[ora, anchor=north] at ([yshift=-2pt]mc.south) {10:00};
    \node[ora, anchor=north] at ([yshift=-2pt]h.south) {9:55};
    \node[ora, anchor=north, text=merr] at ([yshift=-2pt]fc.south) {\textbf{10:12}: modificato dopo \code{funzioni.o}!};
  \end{tikzpicture}
  \caption{Il grafo delle dipendenze del progetto: le frecce vanno da ciò
           che serve a ciò che si costruisce. \code{funzioni.c} è più
           recente di \code{funzioni.o}, quindi \code{funzioni.o} va
           ricostruito, e di conseguenza anche \code{esercizio}.}
  \label{fig:dipendenze}
\end{figure}

Il ragionamento di \code{make} è quello che faremmo noi: per costruire un
file, prima si assicura che siano aggiornati i file da cui dipende (e così
via, all'indietro); poi, se il file \textbf{non esiste} oppure è \textbf{più
vecchio} di almeno uno dei file da cui dipende, lo ricostruisce eseguendo i
comandi indicati. Altrimenti non fa niente. Nella situazione della
figura~\ref{fig:dipendenze}, \code{make} ricompila solo \code{funzioni.c} e
rifà il linking; \code{main.c} non viene toccato.

\subsection{Il makefile}

\code{make} esegue le operazioni descritte in un file, chiamato di solito
\code{makefile} o \code{Makefile}, che contiene la descrizione delle
dipendenze, i comandi da eseguire e, eventualmente, delle definizioni di
\textit{macro}. Per usarlo basta scrivere \code{make} nella cartella che lo
contiene; se il file ha un altro nome, per esempio \code{makeEsercizio}, lo
si indica con l'opzione \code{-f}:

\begin{console}
$ make
$ make -f makeEsercizio
\end{console}

Ecco il primo makefile delle lezioni. Nei listati di questo capitolo il
carattere di tabulazione (TAB) è disegnato come una freccia
\tikz[baseline=-2.6pt]\draw[vsclinenum, -{Stealth[length=3pt,width=3pt]}, line width=0.5pt] (0,0) -- (0.5,0);,
come fa VS Code quando si chiede di mostrare gli spazi bianchi: il motivo
sarà chiaro fra poco.

\begin{lstlisting}[style=make, name=makefile1]
# makefile di esempio
# versione 1
#

eseguibile: main.o funzioni.o
(*@\TAB@*)gcc -o esercizio main.o funzioni.o

main.o: main.c
(*@\TAB@*)gcc -c main.c

funzioni.o: funzioni.c
(*@\TAB@*)gcc -c funzioni.c
\end{lstlisting}

\subsection{I tipi di riga}

In un makefile si distinguono quattro tipi di righe:

\begin{itemize}
  \item le righe di \textbf{dipendenza}: definiscono una relazione fra un
        \term{target} (l'obiettivo, ciò che si costruisce) e uno o più file
        da cui dipende, separati dal carattere \code{:}. Per esempio
        \code{main.o: main.c};
  \item le righe di \textbf{comando}: definiscono l'operazione da eseguire
        per ottenere il target, e seguono la sua riga di dipendenza.
        \textbf{Devono cominciare con un TAB};
  \item le righe di \textbf{commento}: cominciano con \code{\#};
  \item le righe di definizione di \textbf{macro}, che vedremo fra poco.
\end{itemize}

\error{Il TAB non è uno spazio}{
  Le righe di comando devono cominciare con un carattere di
  \textbf{tabulazione}, non con degli spazi, anche se a occhio sono
  identici. Il file \code{makefile4} delle lezioni contiene proprio questo
  errore (quattro spazi al posto del TAB), e \code{make} si rifiuta di
  proseguire con un messaggio poco chiaro:
  \textit{makefile4:9: *** missing separator. Stop.} VS Code inserisce
  davvero un TAB solo se riconosce il file come makefile, cioè se il file si
  chiama \code{makefile} o \code{Makefile}: in basso a destra, nella barra
  di stato, compare allora \textit{Tab Size} invece di \textit{Spaces}.}

\subsection{Che cosa fa \code{make}, passo per passo}

Proviamo il primo makefile. La prima volta mancano tutti i file oggetto e
l'eseguibile, quindi \code{make} esegue tutto, cominciando dai file da cui
dipende il primo target:

\begin{console}
$ make -f makefile1
gcc -c main.c
gcc -c funzioni.c
gcc -o esercizio main.o funzioni.o
\end{console}

\code{make} stampa ogni comando prima di eseguirlo. Si noti l'ordine: per
costruire \code{eseguibile} servono \code{main.o} e \code{funzioni.o}, quindi
prima vengono costruiti loro. Se ora modifichiamo soltanto
\code{funzioni.c} (il comando \code{touch} cambia la data di un file, come
se lo avessimo modificato e salvato) e rilanciamo \code{make}, viene
ricompilato solo quello:

\begin{console}
$ touch funzioni.c
$ make -f makefile1
gcc -c funzioni.c
gcc -o esercizio main.o funzioni.o
\end{console}

\subsection{Le dipendenze dall'header}

Che cosa succede se modifichiamo \code{studente.h}? Con il primo makefile,
niente di buono:

\begin{console}
$ touch studente.h
$ make -f makefile1
gcc -o esercizio main.o funzioni.o
\end{console}

Nessun file viene ricompilato, perché nel makefile i file oggetto dipendono
solo dai rispettivi \code{.c}. Ma entrambi i \code{.c} \textbf{includono}
\code{studente.h}: se nell'header cambia, per esempio, la struttura
\code{studente}, i due file oggetto sono stati compilati con la versione
vecchia e il programma risultante è incoerente, con errori imprevedibili.
La soluzione è dichiarare la dipendenza:

\begin{lstlisting}[style=make, name=makefile2]
# makefile di esempio
# versione 2

eseguibile: main.o funzioni.o
(*@\TAB@*)gcc -o esercizio main.o funzioni.o

main.o: main.c studente.h
(*@\TAB@*)gcc -c main.c

funzioni.o: funzioni.c studente.h
(*@\TAB@*)gcc -c funzioni.c
\end{lstlisting}

\begin{console}
$ make -f makefile2
gcc -c main.c
gcc -c funzioni.c
gcc -o esercizio main.o funzioni.o
\end{console}

Ora, dopo una modifica all'header, vengono ricompilati entrambi i sorgenti.
La regola generale: \textbf{un file oggetto dipende dal suo \code{.c} e da
tutti gli header (nostri) che questo include}.

\subsection{Il nome del target}

Un'altra stranezza: con i primi due makefile il linking viene eseguito
\textbf{ogni volta}, anche se nulla è cambiato. Possiamo evitarlo? Il motivo
è il nome del target: \code{eseguibile}. \code{make} controlla se esiste un
file con quel nome e se è aggiornato; ma il comando produce un file chiamato
\code{esercizio}, quindi un file \code{eseguibile} non esiste mai, e
\code{make} rifà il linking per crearlo. Basta che il target abbia lo stesso
nome del file prodotto:

\begin{lstlisting}[style=make, name=makefile3]
# makefile di esempio
# versione 3

esercizio: main.o funzioni.o
(*@\TAB@*)gcc -o esercizio main.o funzioni.o

main.o: main.c studente.h
(*@\TAB@*)gcc -c main.c

funzioni.o: funzioni.c studente.h
(*@\TAB@*)gcc -c funzioni.c
\end{lstlisting}

\begin{console}
$ make -f makefile3
make: 'esercizio' is up to date.
\end{console}

Tutto è aggiornato, e \code{make} lo dice senza eseguire nulla.

\subsection{Le macro}

Una riga di definizione di \term{macro} definisce un nome simbolico per un
insieme di oggetti, con la sintassi \code{NOME = valore}; il nome si usa poi
scrivendo \code{\$(NOME)}. È l'equivalente del \code{\#define} del C: si
scrive l'elenco una volta sola, e se cambia si modifica in un solo punto.

\begin{lstlisting}[style=make, name=makefile4]
# makefile di esempio
# versione 4

OBJ = funzioni.o main.o

esercizio: $(OBJ)
(*@\TAB@*)gcc -o esercizio $(OBJ)

main.o: main.c studente.h
(*@\TAB@*)gcc -c main.c

funzioni.o: funzioni.c studente.h
(*@\TAB@*)gcc -c funzioni.c
\end{lstlisting}

(Questa è la versione corretta, con il TAB; nel file distribuito con le
lezioni la riga del linking comincia con quattro spazi e produce l'errore
\textit{missing separator}.)

\section{Un makefile completo}

Nella pratica i makefile usano quasi sempre qualche macro in più e un
target di servizio per fare pulizia. Ecco un makefile per il progetto
\code{aree} di inizio capitolo:

\begin{affianco}[0.5]
\begin{lstlisting}[style=make, name=Makefile]
CC = gcc
CFLAGS = -Wall
OBJ = main.o aree.o

aree: $(OBJ)
(*@\TAB@*)$(CC) $(CFLAGS) -o aree $(OBJ)

main.o: main.c aree.h
(*@\TAB@*)$(CC) $(CFLAGS) -c main.c

aree.o: aree.c aree.h
(*@\TAB@*)$(CC) $(CFLAGS) -c aree.c

clean:
(*@\TAB@*)rm -f aree $(OBJ)
\end{lstlisting}
\risultato
\begin{console}
$ make
gcc -Wall -c main.c
gcc -Wall -c aree.c
gcc -Wall -o aree main.o aree.o
$ make
make: 'aree' is up to date.
$ touch main.c
$ make
gcc -Wall -c main.c
gcc -Wall -o aree main.o aree.o
$ make clean
rm -f aree main.o aree.o
\end{console}
\end{affianco}

\begin{itemize}
  \item \code{CC} e \code{CFLAGS} sono i nomi tradizionali delle macro per
        il compilatore e le sue opzioni: per aggiungere un'opzione a tutte
        le compilazioni (per esempio \code{-g}, che serve al debugger) basta
        cambiare una riga;
  \item senza argomenti, \code{make} costruisce il \textbf{primo} target del
        file (qui \code{aree}); con \code{make main.o} si costruisce solo
        quello;
  \item \code{clean} non corrisponde a nessun file e non ha dipendenze: è un
        target “di servizio”, che con \code{make clean} cancella
        l'eseguibile e i file oggetto, per ripartire da zero.
\end{itemize}

\info{make non è solo per il C}{
  \code{make} non sa nulla di C: sa solo che certi file si costruiscono da
  altri file eseguendo dei comandi, se sono più vecchi. Per questo si usa
  anche per compilare documenti, generare siti, elaborare dati: ovunque ci
  sia una catena di file da ricostruire quando qualcosa cambia.}

\insintesi{
  Un programma si divide in più file \code{.c}, ciascuno compilato
  separatamente in un file oggetto (\code{gcc -c}), e in header \code{.h}
  con le dichiarazioni condivise; il linker unisce gli oggetti, abbinando
  ogni funzione usata alla sua definizione (\textit{undefined reference}
  vuol dire che una definizione manca). \code{make} legge un
  \textbf{makefile} fatto di righe di dipendenza (\code{target: file}),
  righe di comando (che cominciano con un \textbf{TAB}), commenti e macro, e
  ricostruisce un target solo se non esiste o è più vecchio di qualcosa da
  cui dipende. Gli header vanno elencati fra le dipendenze, e il nome del
  target deve essere quello del file prodotto.}

\section{Riferimenti}

\begin{itemize}
  \item \textbf{GNU make --- An Introduction to Makefiles}, il manuale
        ufficiale. \\
        \urlx{https://www.gnu.org/software/make/manual/html_node/Introduction.html}
  \item \textbf{Deitel \& Deitel}, \textit{Il linguaggio C}: appendice sulla
        compilazione di programmi su più file.
\end{itemize}


  \chapter{Tipi di dato}
\label{chap:tipi}

Nel capitolo~\ref{chap:calcolatore} abbiamo visto che la memoria contiene
soltanto bit, e che lo stesso byte può essere letto come un numero, come una
lettera o come un pezzo di qualcos'altro: è il \textbf{tipo} a stabilire
quanti byte leggere e come interpretarli. Questo capitolo guarda da vicino i
tipi fondamentali del C: i caratteri, gli interi e i reali. Per ciascuno
vedremo quanta memoria occupa, quali valori può contenere e che cosa succede
quando si superano i suoi limiti; poi come si leggono e si scrivono i
caratteri uno alla volta, come si usano le funzioni matematiche e con quali
regole il C converte un tipo in un altro.

\section{Dichiarazioni, espressioni, assegnazioni}

Tre concetti del capitolo~\ref{chap:primi-programmi} torneranno di
continuo, quindi conviene richiamarli (figura~\ref{fig:dich-espr}):

\begin{itemize}
  \item una \textbf{dichiarazione}, come \code{int i;}, riserva in memoria lo
        spazio per una variabile del tipo indicato e le dà un nome; finché
        non le si assegna nulla, la variabile contiene un valore a caso;
  \item un'\textbf{assegnazione}, come \code{i = 7;}, scrive un valore nella
        variabile;
  \item un'\textbf{espressione}, come \code{i * 2 + 1}, viene calcolata e
        produce un risultato che ha un \textbf{valore} (qui 15) e un
        \textbf{tipo} (qui \code{int}). Anche un'assegnazione è
        un'espressione: \code{i = 7} vale 7.
\end{itemize}

\begin{figure}[H]
  \centering
  \begin{tikzpicture}[font=\Lato\scriptsize,
      tit/.style={font=\Lato\scriptsize\bfseries, text=primary!60!black}]
    \node[tit] at (0.8,1.3) {dichiarazione};
    \node[mcod] at (0.8,0.8) {int i;};
    \variabile{p1}{(1.0,0)}{i}{\mvuota}
    \node[mnota, align=center, anchor=north] at (0.8,-0.4) {si riserva lo spazio;\\il valore è casuale};
    \draw[wtarrow] (2.35,0) -- (3.35,0);
    \node[tit] at (4.9,1.3) {assegnazione};
    \node[mcod] at (4.9,0.8) {i = 7;};
    \variabile{p2}{(5.1,0)}{i}{7}
    \node[mnota, align=center, anchor=north] at (4.9,-0.4) {il valore 7 viene\\scritto nella casella};
    \draw[wtarrow] (6.45,0) -- (7.45,0);
    \node[tit] at (9.4,1.3) {espressione};
    \node[mcod] at (9.4,0.8) {i * 2 + 1};
    \node[wtnodeplain, font=\Lato\scriptsize] at (9.4,0) {valore \textbf{15}\qquad tipo \textbf{int}};
    \node[mnota, align=center, anchor=north] at (9.4,-0.4) {si calcola: il risultato ha\\un valore e un tipo};
  \end{tikzpicture}
  \caption{Dichiarazione, assegnazione ed espressione.}
  \label{fig:dich-espr}
\end{figure}

\section{I tipi fondamentali}

\subsection{Forma completa e forma abbreviata}

Il C ha dodici tipi \term{fondamentali}. Ciascuno ha una forma completa, che
dice tutto esplicitamente, e una forma abbreviata, quella che si usa nella
pratica:

\begin{center}
\renewcommand{\arraystretch}{1.12}
\begin{tabular}{@{}>{\FiraCodeNL}l>{\FiraCodeNL}l@{\hspace{1.6cm}}>{\FiraCodeNL}l>{\FiraCodeNL}l@{}}
  \toprule
  \normalfont\textbf{forma completa} & \normalfont\textbf{si scrive} &
  \normalfont\textbf{forma completa} & \normalfont\textbf{si scrive} \\
  \midrule
  char               & char          & unsigned short int & unsigned short \\
  signed char        & signed char   & unsigned int       & unsigned \\
  unsigned char      & unsigned char & unsigned long int  & unsigned long \\
  signed short int   & short         & float              & float \\
  signed int         & int           & double             & double \\
  signed long int    & long          & long double        & long double \\
  \bottomrule
\end{tabular}
\end{center}

\begin{itemize}
  \item Sono tutte \textbf{parole chiave} (capitolo~\ref{chap:sintassi}):
        non si possono usare come nomi di variabili.
  \item Per \code{short}, \code{int} e \code{long} il segno è sottinteso:
        \code{signed int} è la stessa cosa di \code{int}, quindi la parola
        \code{signed} di solito non si scrive. Anche \code{signed} da solo
        equivale a \code{int}, ma si usa raramente.
  \item \code{short int}, \code{long int} e \code{unsigned int} si
        abbreviano in \code{short}, \code{long} e \code{unsigned}.
  \item Per i caratteri, invece, \code{char}, \code{signed char} e
        \code{unsigned char} sono tre tipi distinti: vedremo fra poco perché.
\end{itemize}

\subsection{Tipi interi e tipi reali}

In base a ciò che rappresentano, i tipi fondamentali formano due famiglie
(figura~\ref{fig:famiglie}). I \term{tipi interi} (\textit{integral types})
contengono numeri interi e, come vedremo, caratteri; i \term{tipi reali}
(\textit{floating types}) contengono numeri con la virgola. Insieme formano i
\term{tipi aritmetici}, quelli su cui si fanno i calcoli.

\begin{figure}[H]
  \centering
  \begin{tikzpicture}[font=\FiraCodeNL\scriptsize,
      tipo/.style={draw=primary!55!black, fill=white, rounded corners=2pt,
        line width=0.6pt, minimum width=2.3cm, minimum height=0.5cm, inner sep=2pt},
      lab/.style={font=\Lato\scriptsize, anchor=west, inner sep=1pt}]
    \draw[rounded corners=6pt, line width=0.8pt, draw=neutral!45!white, fill=neutral!5!white]
      (-1.7,1.55) rectangle (10.25,-2.15);
    \draw[rounded corners=5pt, line width=0.8pt, draw=primary!45!white, fill=primary!7!white]
      (-1.45,0.95) rectangle (6.55,-1.9);
    \draw[rounded corners=5pt, line width=0.8pt, draw=accent!45!white, fill=accent!7!white]
      (7.15,0.95) rectangle (10.0,-1.9);
    \node[lab] at (-1.55,1.28) {\textbf{tipi aritmetici}\ \ {\color{neutral!45!black}ci si fanno i calcoli}};
    \node[lab] at (-1.3,0.67) {\textbf{tipi interi}\ \ {\color{neutral!45!black}numeri interi e caratteri}};
    \node[lab] at (7.3,0.67) {\textbf{tipi reali}};
    \foreach \testo [count=\k from 0] in {char, signed char, unsigned char,
                                          short, int, long,
                                          unsigned short, unsigned, unsigned long} {
      \pgfmathtruncatemacro{\ifam}{\k/3}
      \pgfmathtruncatemacro{\jfam}{mod(\k,3)}
      \node[tipo] at (\jfam*2.55, -\ifam*0.7) {\testo};
    }
    \foreach \testo [count=\k from 0] in {float, double, long double}
      \node[tipo, draw=accent!60!black] at (8.575, -\k*0.7) {\testo};
  \end{tikzpicture}
  \caption{I tipi fondamentali raggruppati per famiglia.}
  \label{fig:famiglie}
\end{figure}

\subsection{I tipi derivati}

A partire dai tipi fondamentali se ne costruiscono altri, i \term{tipi
derivati} (figura~\ref{fig:derivati}): gli \textbf{array} e i
\textbf{puntatori} (capitolo~\ref{chap:puntatori}), le \textbf{strutture}
(capitolo~\ref{chap:strutture}), le enumerazioni e le unioni. In questo
capitolo ci occupiamo solo dei fondamentali.

\begin{figure}[H]
  \centering
  \begin{tikzpicture}[font=\Lato\scriptsize,
      cat/.style={wtnode, font=\Lato\scriptsize\bfseries, inner sep=4pt},
      foglia/.style={font=\Lato\scriptsize, inner sep=2pt, anchor=west}]
    \node[cat] (t) at (0,0) {tipi di dato};
    \node[cat] (f) at (-3.0,-1.15) {fondamentali};
    \node[cat, draw=accent!60!black, fill=accent!8!white] (d) at (3.0,-1.15) {derivati};
    \draw[wtarrow] (t.south) -- ++(0,-0.3) -| (f.north);
    \draw[wtarrow] (t.south) -- ++(0,-0.3) -| (d.north);
    \foreach \testo [count=\k] in {\code{char}, {\code{int} (e \code{short}, \code{long}, \code{unsigned})},
                                   \code{float}, {\code{double} (e \code{long double})}} {
      \node[foglia] (f\k) at ([xshift=16pt, yshift=-\k*0.45cm]f.south west) {\testo};
      \draw[fsramo] ([xshift=8pt]f.south west) |- (f\k.west);
    }
    \foreach \testo [count=\k] in {array (capitolo~\ref{chap:puntatori}),
                                   puntatori (capitolo~\ref{chap:puntatori}),
                                   {strutture, \code{struct} (capitolo~\ref{chap:strutture})},
                                   {enumerazioni, \code{enum}},
                                   {unioni, \code{union}}} {
      \node[foglia] (d\k) at ([xshift=16pt, yshift=-\k*0.45cm]d.south west) {\testo};
      \draw[fsramo] ([xshift=8pt]d.south west) |- (d\k.west);
    }
  \end{tikzpicture}
  \caption{Tipi fondamentali e tipi derivati.}
  \label{fig:derivati}
\end{figure}

\section{Caratteri e il tipo \code{char}}

Cominciamo dal tipo più piccolo, \code{char}, che serve a memorizzare i
caratteri: lettere, cifre, simboli.

\subsection{Un carattere è un numero}

Un computer non conosce le lettere: conosce solo numeri. Per memorizzare un
carattere lo si sostituisce con un numero, secondo una tabella stabilita una
volta per tutte, la \term{codifica}. Un \code{char} occupa \textbf{1 byte},
quindi può assumere $2^8 = 256$ valori diversi: abbastanza per le lettere
maiuscole e minuscole, le cifre, la punteggiatura e alcuni caratteri
speciali come lo spazio, la tabulazione e l'a capo. Non tutti i 256 valori,
però, corrispondono a caratteri stampabili.

In C ogni tipo intero può rappresentare un carattere; in pratica si usano
\code{char} e \code{int}. Le funzioni \code{getchar} e \code{putchar}, che
leggono e scrivono un carattere, lavorano con degli \code{int}; e una
costante carattere come \code{'a'} è di tipo \code{int}, non \code{char}
(in C++ invece è un \code{char}). Lo verificheremo con l'operatore
\code{sizeof}.

\subsection{La codifica ASCII}

Quasi tutti i computer usano la codifica \term{ASCII} (\textit{American
Standard Code for Information Interchange}), che assegna un numero da 0 a 127
a ogni carattere (figura~\ref{fig:ascii}). Esistono altre codifiche, come la
EBCDIC di IBM, usata sui grandi calcolatori (\textit{mainframe}), ma sono rare.

% Una casella della tabella ASCII: \asc{colonna}{riga}{categoria}{carattere}
\newcommand{\asc}[4]{%
  \node[acella, #3] at (#1*2.2, -#2*0.4) {};
  \pgfmathtruncatemacro{\ascn}{#1*16+#2}%
  \node[anum] at (#1*2.2-0.22, -#2*0.4) {\ascn};
  \node[acar] at (#1*2.2+0.42, -#2*0.4) {#4};}
\newcommand{\ascesc}[1]{{\color{accent!75!black}\textbackslash #1}}
\newcommand{\ascnome}[1]{{\fontsize{5}{6}\selectfont\color{neutral!45!black}#1}}

\begin{figure}[H]
  \centering
  \begin{tikzpicture}[
      acella/.style={minimum width=2.1cm, minimum height=0.37cm, rounded corners=1.5pt,
        line width=0.4pt, inner sep=0pt},
      actrl/.style={fill=neutral!9!white, draw=neutral!30!white},
      adig/.style={fill=green-gradient-1!50!white, draw=green-gradient-6!35!white},
      amai/.style={fill=primary!12!white, draw=primary!30!white},
      amin/.style={fill=accent!12!white, draw=accent!30!white},
      apun/.style={fill=white, draw=neutral!30!white},
      anum/.style={font=\FiraCodeNL\tiny, text=neutral!45!black, anchor=east, inner sep=0pt},
      acar/.style={font=\FiraCodeNL\scriptsize\bfseries, inner sep=0pt}]
    \asc{0}{0}{actrl}{\ascesc{0}} \asc{0}{1}{actrl}{\ascnome{SOH}} \asc{0}{2}{actrl}{\ascnome{STX}} \asc{0}{3}{actrl}{\ascnome{ETX}}
    \asc{0}{4}{actrl}{\ascnome{EOT}} \asc{0}{5}{actrl}{\ascnome{ENQ}} \asc{0}{6}{actrl}{\ascnome{ACK}} \asc{0}{7}{actrl}{\ascesc{a}}
    \asc{0}{8}{actrl}{\ascesc{b}} \asc{0}{9}{actrl}{\ascesc{t}} \asc{0}{10}{actrl}{\ascesc{n}} \asc{0}{11}{actrl}{\ascesc{v}}
    \asc{0}{12}{actrl}{\ascesc{f}} \asc{0}{13}{actrl}{\ascesc{r}} \asc{0}{14}{actrl}{\ascnome{SO}} \asc{0}{15}{actrl}{\ascnome{SI}}
    \asc{1}{0}{actrl}{\ascnome{DLE}} \asc{1}{1}{actrl}{\ascnome{DC1}} \asc{1}{2}{actrl}{\ascnome{DC2}} \asc{1}{3}{actrl}{\ascnome{DC3}}
    \asc{1}{4}{actrl}{\ascnome{DC4}} \asc{1}{5}{actrl}{\ascnome{NAK}} \asc{1}{6}{actrl}{\ascnome{SYN}} \asc{1}{7}{actrl}{\ascnome{ETB}}
    \asc{1}{8}{actrl}{\ascnome{CAN}} \asc{1}{9}{actrl}{\ascnome{EM}} \asc{1}{10}{actrl}{\ascnome{SUB}} \asc{1}{11}{actrl}{\ascnome{ESC}}
    \asc{1}{12}{actrl}{\ascnome{FS}} \asc{1}{13}{actrl}{\ascnome{GS}} \asc{1}{14}{actrl}{\ascnome{RS}} \asc{1}{15}{actrl}{\ascnome{US}}
    \asc{2}{0}{apun}{\char"2423} \asc{2}{1}{apun}{!} \asc{2}{2}{apun}{"} \asc{2}{3}{apun}{\#}
    \asc{2}{4}{apun}{\$} \asc{2}{5}{apun}{\%} \asc{2}{6}{apun}{\&} \asc{2}{7}{apun}{'}
    \asc{2}{8}{apun}{(} \asc{2}{9}{apun}{)} \asc{2}{10}{apun}{*} \asc{2}{11}{apun}{+}
    \asc{2}{12}{apun}{,} \asc{2}{13}{apun}{-} \asc{2}{14}{apun}{.} \asc{2}{15}{apun}{/}
    \asc{3}{0}{adig}{0} \asc{3}{1}{adig}{1} \asc{3}{2}{adig}{2} \asc{3}{3}{adig}{3}
    \asc{3}{4}{adig}{4} \asc{3}{5}{adig}{5} \asc{3}{6}{adig}{6} \asc{3}{7}{adig}{7}
    \asc{3}{8}{adig}{8} \asc{3}{9}{adig}{9} \asc{3}{10}{apun}{:} \asc{3}{11}{apun}{;}
    \asc{3}{12}{apun}{<} \asc{3}{13}{apun}{=} \asc{3}{14}{apun}{>} \asc{3}{15}{apun}{?}
    \asc{4}{0}{apun}{@} \asc{4}{1}{amai}{A} \asc{4}{2}{amai}{B} \asc{4}{3}{amai}{C}
    \asc{4}{4}{amai}{D} \asc{4}{5}{amai}{E} \asc{4}{6}{amai}{F} \asc{4}{7}{amai}{G}
    \asc{4}{8}{amai}{H} \asc{4}{9}{amai}{I} \asc{4}{10}{amai}{J} \asc{4}{11}{amai}{K}
    \asc{4}{12}{amai}{L} \asc{4}{13}{amai}{M} \asc{4}{14}{amai}{N} \asc{4}{15}{amai}{O}
    \asc{5}{0}{amai}{P} \asc{5}{1}{amai}{Q} \asc{5}{2}{amai}{R} \asc{5}{3}{amai}{S}
    \asc{5}{4}{amai}{T} \asc{5}{5}{amai}{U} \asc{5}{6}{amai}{V} \asc{5}{7}{amai}{W}
    \asc{5}{8}{amai}{X} \asc{5}{9}{amai}{Y} \asc{5}{10}{amai}{Z} \asc{5}{11}{apun}{[}
    \asc{5}{12}{apun}{\textbackslash} \asc{5}{13}{apun}{]} \asc{5}{14}{apun}{\textasciicircum} \asc{5}{15}{apun}{\_}
    \asc{6}{0}{apun}{`} \asc{6}{1}{amin}{a} \asc{6}{2}{amin}{b} \asc{6}{3}{amin}{c}
    \asc{6}{4}{amin}{d} \asc{6}{5}{amin}{e} \asc{6}{6}{amin}{f} \asc{6}{7}{amin}{g}
    \asc{6}{8}{amin}{h} \asc{6}{9}{amin}{i} \asc{6}{10}{amin}{j} \asc{6}{11}{amin}{k}
    \asc{6}{12}{amin}{l} \asc{6}{13}{amin}{m} \asc{6}{14}{amin}{n} \asc{6}{15}{amin}{o}
    \asc{7}{0}{amin}{p} \asc{7}{1}{amin}{q} \asc{7}{2}{amin}{r} \asc{7}{3}{amin}{s}
    \asc{7}{4}{amin}{t} \asc{7}{5}{amin}{u} \asc{7}{6}{amin}{v} \asc{7}{7}{amin}{w}
    \asc{7}{8}{amin}{x} \asc{7}{9}{amin}{y} \asc{7}{10}{amin}{z} \asc{7}{11}{apun}{\{}
    \asc{7}{12}{apun}{|} \asc{7}{13}{apun}{\}} \asc{7}{14}{apun}{\textasciitilde} \asc{7}{15}{actrl}{\ascnome{DEL}}
    % legenda
    \foreach \sty/\testo/\xx in {actrl/{caratteri di controllo (non stampabili)}/-0.6,
                               apun/{spazio e simboli}/4.6, adig/cifre/8.0,
                               amai/maiuscole/10.3, amin/minuscole/12.9} {
      \node[acella, \sty, minimum width=0.5cm] at (\xx, -6.85) {};
      \node[font=\Lato\scriptsize, anchor=west] at (\xx+0.35, -6.85) {\testo};
    }
  \end{tikzpicture}
  \caption{La tabella ASCII: accanto a ogni codice (da 0 a 127) il carattere
           corrispondente. I caratteri di controllo hanno un nome
           convenzionale, o la sequenza di escape con cui si scrivono in C;
           \code{\char"2423} è lo spazio.}
  \label{fig:ascii}
\end{figure}

La tabella non va imparata a memoria: in C si scrive \code{'A'} e il
compilatore lo sostituisce con 65. Conviene però capirne la struttura:

\begin{itemize}
  \item i codici da 0 a 31 e il 127 sono \textbf{caratteri di controllo}:
        non si stampano, ma comandano il dispositivo (il 10, per esempio, è
        l'a capo); il 32 è lo spazio;
  \item la scelta dei numeri è una \textbf{convenzione}: nella lettera
        \textit{A} non c'è nulla che la leghi a 65;
  \item però le cifre, le maiuscole e le minuscole occupano ciascuna dei
        codici \textbf{consecutivi} e in ordine: \code{'0'} < \code{'1'} <
        \dots{} < \code{'9'}, \code{'a'} < \code{'b'} < \dots{} <
        \code{'z'}. Confrontare i codici equivale quindi a confrontare
        nell'ordine alfabetico, e questo servirà per ordinare i testi;
  \item le maiuscole vengono prima delle minuscole, e ogni minuscola dista
        dalla sua maiuscola esattamente 32 posizioni: \code{'a'} vale 97,
        \code{'A'} vale 65;
  \item il carattere \code{'7'} non è il numero 7: vale 55. Il numero si
        ottiene con la differenza \code{'7' - '0'}, cioè $55 - 48 = 7$.
\end{itemize}

\info{E le lettere accentate?}{
  La tabella ASCII non contiene \textit{è}, \textit{à}, \textit{ù}. I sistemi
  di oggi (Linux, macOS, VS Code) usano la codifica \textbf{UTF-8}, che
  estende l'ASCII: i primi 128 codici sono gli stessi, mentre gli altri
  caratteri occupano da 2 a 4 byte. Una lettera accentata, quindi,
  \textbf{non sta in un \code{char}}: \code{char c = 'è';} produce due avvisi
  (\textit{multi-character character constant} e \textit{overflow in
  conversion}) e un valore senza senso. Dentro una stringa, invece, non c'è
  problema: \code{printf("Perché?\textbackslash n");} stampa correttamente,
  perché la stringa contiene semplicemente i due byte della \textit{é}.}

\subsection{Caratteri come interi, interi come caratteri}

Poiché un \code{char} contiene un numero, lo si può stampare in due modi:
con \code{\%c} \code{printf} disegna il carattere che ha quel codice, con
\code{\%d} scrive il numero. E siccome è un numero, lo si può usare nelle
espressioni aritmetiche: \code{c + 1} è il codice del carattere successivo.

\begin{affianco}[0.56][c]
\begin{lstlisting}[style=c]
char c = 'a';

printf("%c\n", c);
printf("%d\n", c);
printf("%c%c%c\n", c, c + 1, c + 2);
\end{lstlisting}
\risultato
\begin{console}[Terminale]
a
97
abc
\end{console}
\end{affianco}

Con la stessa idea si scorrono i caratteri con un ciclo: la variabile di
controllo parte da un codice e arriva a un altro. Nel programma seguente il
primo ciclo usa un \code{int} che va da \code{'a'} a \code{'z'}, il secondo
un \code{char} che va da 65 a 90 (i codici di \code{'A'} e \code{'Z'}), il
terzo stampa con \code{\%d} i codici delle dieci cifre.

\begin{affianco}[0.5]
\begin{lstlisting}[style=c, name=caratteri.c]
#include <stdio.h>

int main(void)
{
   int i;
   char c;

   /* stampa abc...z */
   for (i = 'a'; i <= 'z'; i++)
      printf("%c", i);
   printf("\n");

   /* stampa ABC...Z */
   for (c = 65; c <= 90; c++)
      printf("%c", c);
   printf("\n");

   /* stampa 48 49 ... 57 */
   for (c = '0'; c <= '9'; c++)
      printf("%d ", c);
   printf("\n");
   return 0;
}
\end{lstlisting}
\risultato
\begin{console}
$ ./caratteri
abcdefghijklmnopqrstuvwxyz
ABCDEFGHIJKLMNOPQRSTUVWXYZ
48 49 50 51 52 53 54 55 56 57
\end{console}
\end{affianco}

\warning{Una svista nelle slide}{
  Nelle slide il secondo ciclo è scritto \code{for (c=65; i<=90; c++)}: la
  condizione controlla \code{i} invece di \code{c}. Alla fine del primo
  ciclo \code{i} vale 123 (il primo valore per cui \code{i <= 'z'} è
  falsa), quindi \code{i <= 90} è falsa fin dall'inizio e il ciclo non stampa
  nulla. Il compilatore non può accorgersene: il codice è corretto dal punto
  di vista della sintassi, e fa semplicemente un'altra cosa.}

\section{Le sequenze di escape}

\subsection{Caratteri che non si possono scrivere}

Alcuni caratteri non si possono scrivere direttamente fra apici o fra
virgolette: o perché non sono stampabili (l'a capo, la tabulazione), o
perché hanno già un significato (le virgolette chiudono una stringa,
l'apice chiude una costante carattere). Per questi il C usa le
\term{sequenze di escape}: un backslash seguito da un carattere. Sono
\textbf{due simboli} nel codice, ma \textbf{un solo carattere} in memoria:
\code{'\textbackslash n'} è il carattere di codice 10.

\begin{center}
\renewcommand{\arraystretch}{1.12}
\begin{tabular}{@{}l>{\FiraCodeNL}c c>{\raggedright\arraybackslash}p{9.2cm}@{}}
  \toprule
  \textbf{carattere} & \normalfont\textbf{in C} & \textbf{codice} & \textbf{che cosa fa} \\
  \midrule
  a capo (\textit{newline})                   & \textbackslash n & 10 & porta il cursore all'inizio della riga successiva \\
  tabulazione (\textit{horizontal tab})       & \textbackslash t & 9  & sposta il cursore alla prossima tabulazione (ogni 8 colonne) \\
  ritorno carrello (\textit{carriage return}) & \textbackslash r & 13 & riporta il cursore all'inizio della riga \\
  tabulazione verticale (\textit{vertical tab}) & \textbackslash v & 11 & scende di una riga, restando nella stessa colonna \\
  salto pagina (\textit{formfeed})            & \textbackslash f & 12 & sulla stampante cambia pagina; nel terminale fa come \code{\textbackslash v} \\
  indietro (\textit{backspace})               & \textbackslash b & 8  & sposta il cursore indietro di una posizione \\
  campanello (\textit{alert})                 & \textbackslash a & 7  & emette un segnale acustico (o visivo) \\
  carattere nullo (\textit{null})             & \textbackslash 0 & 0  & il carattere di codice 0, che chiude le stringhe \\
  backslash                                   & \textbackslash\textbackslash & 92 & il carattere \code{\textbackslash} \\
  virgolette (\textit{double quote})          & \textbackslash" & 34 & il carattere \code{"} \\
  apice (\textit{single quote})               & \textbackslash' & 39 & il carattere \code{'} \\
  punto interrogativo                         & \textbackslash? & 63 & il carattere \code{?} (serve di rado) \\
  \bottomrule
\end{tabular}
\end{center}

\subsection{Virgolette, apici e campanello}

Per stampare delle virgolette dentro una stringa bisogna scriverle con il
backslash, altrimenti chiuderebbero la stringa. Il programma delle lezioni
lo fa in un modo un po' tortuoso, stampando con \code{\%c} il carattere
\code{'\textbackslash"'}; il modo diretto è la seconda riga:

\begin{affianco}[0.56][c]
\begin{lstlisting}[style=c]
printf("Di nome %cNemo%c\n", '\"', '\"');
printf("Di nome \"Nemo\"\n");
\end{lstlisting}
\risultato
\begin{console}
Di nome "Nemo"
Di nome "Nemo"
\end{console}
\end{affianco}

Fra apici le virgolette non danno fastidio, quindi anche \code{'"'} va bene;
viceversa, l'apice va protetto fra apici (\code{'\textbackslash'{}'}) ma non
in una stringa (\code{"l'una"}).

\code{printf("\%c", '\textbackslash a');} fa suonare il
\textbf{campanello} del terminale: un'eredità delle telescriventi, che ne
avevano uno vero. Nel terminale di VS Code, a seconda delle impostazioni, si
sente un suono, la finestra lampeggia o non succede nulla.

\subsection{Il cursore e i caratteri di controllo}

Il programma seguente stampa la stessa stringa con in mezzo, a turno,
diversi caratteri di controllo. Ecco che cosa compare sul terminale:

\begin{affianco}[0.5][c]
\begin{lstlisting}[style=c, name=escape.c]
#include <stdio.h>

int main(void)
{
  printf("abc \t abc\n");
  printf("abc \v abc\n");
  printf("abc \f abc\n");
  printf("abc \r abc\n");
  printf("abc \r \v abc\n");
  return 0;
}
\end{lstlisting}
\risultato
\begin{console}
$ ./escape
abc      abc
abc
     abc
abc
     abc
 abc
 bc
  abc
\end{console}
\end{affianco}

Per capirlo basta seguire il cursore (figura~\ref{fig:cursore}). Ogni
carattere stampabile viene scritto dove si trova il cursore, che poi avanza
di una colonna; i caratteri di controllo, invece, spostano soltanto il
cursore.

\begin{figure}[H]
  \centering
  \begin{tikzpicture}[font=\FiraCodeNL\scriptsize,
      tc/.style={mcella, minimum width=0.38cm, minimum height=0.46cm, inner xsep=0pt},
      tcn/.style={tc, fill=accent!14!white},
      tit/.style={font=\FiraCodeNL\scriptsize, anchor=west, inner sep=1pt},
      esc/.style={font=\FiraCodeNL\scriptsize\bfseries, text=accent!75!black, inner sep=1.5pt},
      lab/.style={font=\Lato\scriptsize, text=neutral!40!black, anchor=east}]
    \noindici
    % \t
    \node[tit] at (0,1.35) {printf("abc \textbackslash t abc\textbackslash n");};
    \celle[tc]{ta}{(0,0)}{a,b,c,\char"2423,{},{},{},{},\char"2423,a,b,c}
    \draw[wtarrowacc] ([yshift=1pt]ta-4.north west) to[bend left=55]
      node[esc, above] {\textbackslash t} ([yshift=1pt]ta-8.north west);
    \node[mindice, anchor=north] at ([yshift=-1pt]ta-0.south) {0};
    \node[mindice, anchor=north] at ([yshift=-1pt]ta-4.south) {4};
    \node[mindice, anchor=north] at ([yshift=-1pt]ta-8.south) {8};
    \node[mnota, anchor=north west, text width=4.4cm] at (-0.1,-0.6)
      {salta alla prossima colonna multipla di 8: qui dalla 4 alla 8};
    % \v
    \node[tit] at (6.0,1.35) {printf("abc \textbackslash v abc\textbackslash n");};
    \celle[tc]{tb}{(6.0,0)}{a,b,c,\char"2423,{},{},{},{}}
    \celle[tc]{tv}{(6.0,-0.85)}{{},{},{},{},\char"2423,a,b,c}
    \draw[wtarrowacc] (tb-4.center) -- node[esc, right=3pt, pos=0.45] {\textbackslash v} (tv-4.north);
    \node[mnota, anchor=north west, text width=4.4cm] at (5.9,-1.25)
      {scende di una riga e resta nella stessa colonna};
    % \r
    \node[tit] at (11.6,1.35) {printf("abc \textbackslash r abc\textbackslash n");};
    \celle[tc]{tr}{(12.4,0)}{a,b,c,\char"2423}
    \celle[tcn]{ts}{(12.4,-0.85)}{\char"2423,a,b,c}
    \node[lab] at (12.3,0) {prima};
    \node[lab] at (12.3,-0.85) {dopo};
    \draw[wtarrowacc] ([yshift=1pt]tr-3.north east) to[bend right=55]
      node[esc, above] {\textbackslash r} ([yshift=1pt]tr-0.north west);
    \node[mnota, anchor=north west, text width=4.4cm] at (11.5,-1.25)
      {torna all'inizio della riga: ciò che segue sovrascrive i caratteri già scritti};
  \end{tikzpicture}
  \caption{Il percorso del cursore con \code{\textbackslash t},
           \code{\textbackslash v} e \code{\textbackslash r}
           (\code{\char"2423} è uno spazio).}
  \label{fig:cursore}
\end{figure}

\begin{itemize}
  \item \code{\textbackslash t} sposta il cursore alla prossima
        \textbf{tabulazione}; le tabulazioni sono ogni 8 colonne (0, 8,
        16\dots), e per questo servono ad allineare i dati in colonna.
  \item \code{\textbackslash v} e \code{\textbackslash f} fanno scendere il
        cursore di una riga, ma senza riportarlo a sinistra.
  \item \code{\textbackslash r} riporta il cursore all'inizio della riga
        \textbf{senza} scendere: ciò che si stampa dopo sovrascrive. Così
        \code{"abc "} diventa \code{" abc"}.
  \item L'ultima riga combina le due cose: \code{\textbackslash r} torna
        all'inizio, lo spazio sovrascrive la \code{a} (si legge \code{" bc"}),
        \code{\textbackslash v} scende di una riga restando nella colonna 1,
        e lì viene scritto \code{" abc"}.
\end{itemize}

E \code{\textbackslash n}? Nel terminale fa due cose insieme: scende di una
riga \textbf{e} torna all'inizio. Sulle macchine da scrivere erano due
movimenti distinti, riportare il carrello all'inizio (\textit{carriage
return}) e far scorrere il foglio di una riga (\textit{line feed}): per
questo i file di testo di Windows terminano ogni riga con due caratteri,
\code{\textbackslash r\textbackslash n}, mentre Linux e macOS usano solo
\code{\textbackslash n}.

\subsection{Un'animazione con \code{\textbackslash r}}

Il ritorno carrello permette di ridisegnare più volte lo stesso punto dello
schermo. Il programma seguente stampa a ripetizione i caratteri
\code{|}, \code{/}, \code{-} e \code{\textbackslash}, ciascuno seguito da
\code{\textbackslash r}, con una pausa di un secondo fra l'uno e l'altro:
ogni simbolo prende il posto del precedente e si vede una stanghetta che
ruota.

\begin{affianco}[0.5][c]
\begin{lstlisting}[style=c, name=spinner.c]
#include <stdio.h>
#include <unistd.h>

#define PAUSA 1

void pausa(int l)
{
    sleep(l);
}

int main()
{
    setbuf(stdout, NULL);
    while (1) {
        printf("|\r");  pausa(PAUSA);
        printf("/\r");  pausa(PAUSA);
        printf("-\r");  pausa(PAUSA);
        printf("\\\r"); pausa(PAUSA);
    }
}
\end{lstlisting}
\risultato
\begin{itemize}
  \item \code{"\textbackslash\textbackslash\textbackslash r"} contiene due
        caratteri: il backslash (scritto \code{\textbackslash\textbackslash})
        e il ritorno carrello.
  \item \code{while (1)} è un ciclo infinito, perché la condizione è sempre
        vera: il programma si ferma con \tasti{Ctrl,C}.
  \item \code{sleep(l)}, dichiarata in \code{unistd.h}, sospende il
        programma per \code{l} secondi. È disponibile su Linux, su macOS e
        in WSL, non su Windows senza WSL.
\end{itemize}
\end{affianco}

Nella cartella degli esempi c'è anche una versione senza \code{sleep}, in
cui \code{pausa} perde tempo con un ciclo vuoto:
\code{for (i = 0; i < l; i++);} (il \code{;} da solo è l'istruzione vuota).
Funziona, ma la durata della pausa dipende dalla velocità del computer; e
compilando con l'ottimizzazione (\code{gcc -O2}) il compilatore si accorge
che il ciclo non fa nulla e lo elimina: in due secondi la versione ottimizzata
ha stampato oltre 2,8 milioni di caratteri, quella normale 70.

\subsection{Il buffer di \code{printf}}

Resta da capire \code{setbuf(stdout, NULL)}. Per efficienza, \code{printf}
non manda ogni carattere allo schermo appena lo riceve: lo accumula in una
zona di memoria, il \term{buffer}, e la svuota tutta insieme quando incontra
un \code{\textbackslash n}, quando è piena o quando il programma termina
(figura~\ref{fig:buffer}). Il nostro programma non stampa mai un
\code{\textbackslash n} e non termina mai: senza \code{setbuf} il buffer non
verrebbe mai svuotato e sullo schermo \textbf{non comparirebbe nulla}.
L'abbiamo verificato: senza quella riga, in tre secondi il programma non
mostra niente. \code{setbuf(stdout, NULL)} elimina il buffer dell'output
standard (\code{stdout}), e ogni carattere arriva subito allo schermo.

\begin{figure}[H]
  \centering
  \begin{tikzpicture}[font=\Lato\scriptsize]
    \noindici
    \node[font=\Lato\scriptsize\bfseries, anchor=west] at (-0.9,0.85) {senza \code{setbuf}};
    \node[wtnode, font=\Lato\scriptsize] (p1) at (0,0) {programma};
    \celle[mcella, minimum width=0.55cm]{bf}{(1.7,0)}{|,\textbackslash r,/,\textbackslash r,-,\textbackslash r,{}}
    \node[mnota, anchor=south] at (bf-box.north) {buffer di \code{stdout}: i caratteri si accumulano};
    \node[wtnodeplain] (s1) at (9.2,0) {schermo};
    \draw[wtarrow] (p1) -- (bf-box.west);
    \draw[wtarrow, dashed, draw=merr] (bf-box.east) -- node[above, mnota, text=merr] {solo con \code{\textbackslash n}} (s1);
    \node[font=\Lato\scriptsize\bfseries, anchor=west] at (-0.9,-0.95) {con \code{setbuf(stdout, NULL)}};
    \node[wtnode, font=\Lato\scriptsize] (p2) at (0,-1.8) {programma};
    \node[wtnodeplain] (s2) at (9.2,-1.8) {schermo};
    \draw[wtarrow, draw=mok] (p2) -- node[above, mnota] {ogni carattere arriva subito} (s2);
  \end{tikzpicture}
  \caption{Con il buffer l'output arriva allo schermo a righe intere; senza,
           un carattere alla volta.}
  \label{fig:buffer}
\end{figure}

\section{Attenzione a \code{++} dentro \code{printf}}

Le slide mettono in guardia: con gli operatori \code{++} e \code{--} a volte
i risultati \textbf{non sono prevedibili}. Il programma seguente stampa più
volte \code{c} e \code{a} incrementandoli dentro la stessa chiamata di
\code{printf}. Compilato con \code{gcc} dà i risultati delle slide; compilato
con \code{clang} dà risultati del tutto diversi:

\begin{affianco}[0.5][t]
\begin{lstlisting}[style=c, name=printf++.c]
#include <stdio.h>

int main()
{
  char c;  int a;

  c = 'a'; a=10;
  printf("%c%c%c\n", c, c, c++);
  printf ("%d %d %d \n", a , a , a++);

  c = 'a'; a=10;
  printf("%c%c%c\n", c, c, ++c);
  printf("%d %d %d \n", a , a , ++a);

  c = 'a'; a=10;
  printf("%c%c%c\n", c, ++c, ++c);
  printf("%d %d %d \n", a , ++a , ++a);

  c = 'a'; a=10;
  printf("%c%c%c\n", c, c++, c++);
  printf("%d %d %d \n", a , a++ , a++);

  c = 'a'; a=10;
  printf("%c%c%c \n", ++c , c++ , c);
  printf("%d %d %d \n", ++a , a++ , a);
}
\end{lstlisting}
\risultato
\begin{minipage}[t]{0.48\linewidth}
\begin{console}[gcc]
bba
11 11 10
bbb
11 11 11
ccb
12 12 12
cba
12 11 10
caa
12 10 12
\end{console}
\end{minipage}\hfill
\begin{minipage}[t]{0.48\linewidth}
\begin{console}[clang]
aaa
10 10 10
aab
10 10 11
abc
10 11 12
aab
10 10 11
bbc
11 11 12
\end{console}
\end{minipage}
\end{affianco}

Nessuno dei due compilatori sbaglia. Il C non stabilisce in che ordine
vengano calcolati gli argomenti di una funzione (qui \code{gcc} li ha
calcolati da destra a sinistra, \code{clang} da sinistra a destra); e
modificare una variabile con \code{++} e usarla di nuovo nella stessa
chiamata dà un comportamento \textbf{indefinito}, come \code{a * a++} nel
capitolo~\ref{chap:sintassi}. Con \code{-Wall}, \code{gcc} lo segnala per
ogni riga:

\begin{console}
$ gcc -Wall -o printf++ printf++.c
printf++.c: In function (*@‘@*)main(*@’@*):
printf++.c:8:29: warning: operation on (*@‘@*)c(*@’@*) may be undefined [-Wsequence-point]
    8 |   printf("%c%c%c\n", c, c, c++);
      |                            ~^~
\end{console}

La regola pratica: una variabile incrementata con \code{++} o decrementata
con \code{--} \textbf{non va usata altrove} nella stessa espressione o nella
stessa chiamata. Si incrementa prima, in un'istruzione a parte.

\section{Come è fatto un \code{char}}

Vediamo ora che cosa c'è davvero nel byte di un \code{char}, e perché
esistono tre tipi diversi per i caratteri.

\subsection{I bit di un carattere}

La notazione binaria (capitolo~\ref{chap:calcolatore}) è
\textbf{posizionale} come quella decimale. In base 10 il numero
$d_n d_{n-1} \dots d_1 d_0$ vale
$d_n \cdot 10^n + \dots + d_1 \cdot 10^1 + d_0 \cdot 10^0$ (per esempio
$20753 = 2 \cdot 10^4 + 0 \cdot 10^3 + 7 \cdot 10^2 + 5 \cdot 10^1 + 3 \cdot
10^0$: il 3 è la cifra meno significativa, il 2 la più significativa). In
base 2 ogni cifra $b_i$ è un bit, 0 o 1, e il numero vale
$b_n \cdot 2^n + \dots + b_1 \cdot 2^1 + b_0 \cdot 2^0$. Dopo la
dichiarazione \code{char c = 'a';} il byte di \code{c} contiene
\code{01100001} (figura~\ref{fig:bit-a}).

\begin{figure}[H]
  \centering
  \begin{tikzpicture}[font=\FiraCodeNL\small]
    \foreach \p [count=\i from 0] in {128,64,32,16,8,4,2,1} {
      \node[mcella, minimum width=1.1cm, minimum height=0.65cm] (b\i) at (\i*1.1,0) {};
      \node[font=\FiraCodeNL\scriptsize, text=neutral!40!black] at (\i*1.1,0.58) {\p};
      \pgfmathtruncatemacro{\e}{7-\i}
      \node[font=\Lato\tiny, text=neutral!50!black] at (\i*1.1,0.95) {$2^{\e}$};
    }
    \foreach \val [count=\i from 0] in {0,1,1,0,0,0,0,1} {
      \ifnum\val=1
        \node[font=\FiraCodeNL\normalsize\bfseries, text=accent!80!black] at (b\i) {1};
      \else
        \node[font=\FiraCodeNL\normalsize, text=neutral!55!black] at (b\i) {0};
      \fi
    }
    \node[anchor=west, font=\Lato\small] at (8.4,0)
      {$= 64 + 32 + 1 = \mathbf{97}$, il codice di \code{'a'}};
  \end{tikzpicture}
  \caption{Il byte del carattere \code{'a'}: ogni bit a 1 aggiunge il peso
           della sua posizione.}
  \label{fig:bit-a}
\end{figure}

\subsection{Tre tipi per i caratteri}

Il C ha tre tipi da un byte: \code{char}, \code{signed char} e
\code{unsigned char}. Ognuno ha 256 valori distinti, ma non gli stessi:

\begin{itemize}
  \item \code{signed char} va da $-128$ a $127$;
  \item \code{unsigned char} va da $0$ a $255$;
  \item \code{char} è uno dei due, a seconda dell'\textbf{implementazione}
        (cioè del compilatore e del processore).
\end{itemize}

Per i caratteri ASCII (da 0 a 127) la differenza non conta. Conta quando il
byte si usa come numero: per capirla bisogna vedere come si rappresentano i
numeri negativi.

\subsection{I numeri negativi}

L'idea più semplice è il \textbf{bit di segno}. Limitiamoci a 4 bit: il
bit più significativo indica il segno e gli altri tre il valore, quindi
$5 = $~\code{0101} e $-5 = $~\code{1101}. Lo zero però ha due
rappresentazioni, \code{0000} e \code{1000}, e così un valore va sprecato.

Per questo i computer usano il \term{complemento a due}: i numeri positivi si
rappresentano normalmente, e per un negativo si \textbf{invertono} i bit del
corrispondente positivo e si \textbf{somma 1}. Così \code{1000} non è più
uno zero, ma $-8$ (figura~\ref{fig:complemento}).

\begin{figure}[H]
  \centering
  \begin{tikzpicture}[font=\FiraCodeNL\scriptsize,
      vn/.style={mcella, minimum width=0.95cm, font=\FiraCodeNL\scriptsize\bfseries, fill=accent!10!white},
      vp/.style={mcella, minimum width=0.95cm, font=\FiraCodeNL\scriptsize\bfseries},
      pn/.style={mcella, minimum width=0.95cm, fill=accent!10!white},
      pp/.style={mcella, minimum width=0.95cm},
      lab/.style={font=\Lato\scriptsize, anchor=east, text=neutral!40!black}]
    \noindici
    \celle[vn]{vn}{(0,0)}{$-8$,$-7$,$-6$,$-5$,$-4$,$-3$,$-2$,$-1$}
    \celle[vp]{vp}{(vn-7.east)}{0,1,2,3,4,5,6,7}
    \celle[pn]{pn}{(0,-0.5)}{1000,1001,1010,1011,1100,1101,1110,1111}
    \celle[pp]{pp}{(pn-7.east)}{0000,0001,0010,0011,0100,0101,0110,0111}
    \node[lab] at (-0.1,0) {valore};
    \node[lab] at (-0.1,-0.5) {bit};
    \draw[wtbrace] ([yshift=-3pt]pn-0.south west) -- ([yshift=-3pt]pn-7.south east);
    \draw[wtbrace] ([yshift=-3pt]pp-0.south west) -- ([yshift=-3pt]pp-7.south east);
    \node[mnota, anchor=north] at ([yshift=-0.32cm]$(pn-0.south)!0.5!(pn-7.south)$) {il primo bit è 1: numeri negativi};
    \node[mnota, anchor=north] at ([yshift=-0.32cm]$(pp-0.south)!0.5!(pp-7.south)$) {il primo bit è 0: zero e numeri positivi};
  \end{tikzpicture}
  \caption{I 16 valori di 4 bit in complemento a due: da $-8$ a $+7$, con
           un solo zero.}
  \label{fig:complemento}
\end{figure}

\begin{affianco}[0.36][c]
\begin{center}
\renewcommand{\arraystretch}{1.15}
\begin{tabular}{@{}l>{\FiraCodeNL}r@{}}
  \toprule
  $5$ & 0101 \\
  bit invertiti & 1010 \\
  sommo 1 & + 0001 \\
  \midrule
  $-5$ & 1011 \\
  \bottomrule
\end{tabular}
\end{center}
\risultato
Per esempio, $-5$ si ottiene da \code{0101}: invertendo i bit si ha
\code{1010}, e sommando 1 si arriva a \code{1011}. Perché funziona? Un numero
più il suo inverso dà sempre \code{1111} (in ogni colonna c'è un 1 e uno 0,
senza riporti), cioè $-1$. Quindi l'inverso di $x$ vale $-1 - x$, e
aggiungendo 1 si ottiene proprio $-x$.
\end{affianco}

\subsection{Lo stesso byte, due letture}

Su 8 bit funziona allo stesso modo. I byte da \code{00000000} a
\code{01111111} hanno lo stesso valore, da 0 a 127, per \code{signed char} e
per \code{unsigned char}. Quelli che cominciano con un 1 no: per
\code{unsigned char} vanno da 128 a 255, per \code{signed char} da $-128$ a
$-1$ (figura~\ref{fig:due-letture}). Il byte \code{11111111}, per esempio,
vale 255 oppure $-1$.

\begin{figure}[H]
  \centering
  \begin{tikzpicture}[font=\FiraCodeNL\scriptsize,
      lab/.style={font=\Lato\scriptsize\bfseries, anchor=east}]
    \fill[green-gradient-1!40!white] (0,0) rectangle (6,0.55);
    \fill[accent!12!white] (6,0) rectangle (12,0.55);
    \draw[primary!65!black, line width=0.7pt] (0,0) rectangle (12,0.55) (6,0) -- (6,0.55);
    \node at (3,0.275) {00000000 \ldots\ 01111111};
    \node at (9,0.275) {10000000 \ldots\ 11111111};
    \node[lab] at (-0.35,0.275) {byte};
    \node[lab] at (-0.35,1.05) {unsigned char};
    \node[lab] at (-0.35,-0.5) {signed char};
    \foreach \xx/\val in {0.2/0, 5.7/127, 6.3/128, 11.8/255}
      \node at (\xx,1.05) {\val};
    \foreach \xx/\val in {0.2/0, 5.7/127, 6.3/$-128$, 11.8/$-1$}
      \node at (\xx,-0.5) {\val};
    \draw[wtarrow] (0.6,1.05) -- (5.25,1.05);
    \draw[wtarrow] (6.75,1.05) -- (11.4,1.05);
    \draw[wtarrow] (0.6,-0.5) -- (5.25,-0.5);
    \draw[wtarrow, draw=accent!70!black] (6.85,-0.5) -- (11.35,-0.5);
    \node[mnota, align=center, anchor=north] at (3,-0.75) {le due letture coincidono};
    \node[mnota, align=center, anchor=north] at (9,-0.75) {le due letture sono diverse};
  \end{tikzpicture}
  \caption{I 256 byte letti come \code{unsigned char} (sopra) e come
           \code{signed char} (sotto).}
  \label{fig:due-letture}
\end{figure}

\subsection{Il \code{char} del mio computer ha il segno?}

Come si scopre se \code{char} è \code{signed} o \code{unsigned}? Si assegna
$-1$ ai tre tipi e si stampano i valori con \code{\%d}:

\begin{affianco}[0.5][c]
\begin{lstlisting}[style=c, name=char.c]
#include <stdio.h>

int main()
{
    char c = -1;
    unsigned char u = -1;
    signed char s = -1;

    printf("c = %d  u = %d   s = %d\n", c, u, s);

    return 0;
}
\end{lstlisting}
\risultato
\begin{console}
$ gcc -Wall -o char char.c
$ ./char
c = -1  u = 255   s = -1
$ gcc -Wall -funsigned-char -o char char.c
$ ./char
c = 255  u = 255   s = -1
\end{console}
\end{affianco}

In tutte e tre le variabili il byte è \code{11111111}: $-1$ assegnato a un
\code{unsigned char}, che non ha negativi, diventa $-1 + 256 = 255$. La
differenza sta nella lettura: per stampare, il valore viene convertito in
\code{int} e conserva il suo significato, quindi $-1$ per il tipo con segno e
255 per quello senza. Sui PC (processori Intel e AMD) \code{char} ha il
segno, e infatti \code{c} vale $-1$. Su altri sistemi, per esempio Linux sui
processori ARM come quelli del Raspberry Pi, \code{char} è senza segno; lo si
può simulare con l'opzione \code{-funsigned-char} di \code{gcc}, ed ecco
\code{c = 255}.

La conclusione pratica: per i caratteri si usa \code{char}; quando un byte
serve come \textbf{numero} piccolo e il segno conta, si scrive esplicitamente
\code{signed char} o \code{unsigned char}.

\section{Il tipo \code{int}}

\code{int} è il tipo principale del C, quello che si usa per qualunque
numero intero quando non ci sono motivi per sceglierne un altro.

\subsection{Quanti valori}

Un \code{int} occupa 2 byte (16 bit) sui sistemi vecchi e 4 byte (32 bit)
su quasi tutti i sistemi moderni: il numero di valori distinti, quindi,
dipende dal sistema. In complemento a due, con $n$ bit si rappresentano gli
interi da $-2^{n-1}$ a $2^{n-1}-1$:

\begin{itemize}
  \item con 16 bit: da $-2^{15} = -32\,768$ a $2^{15} - 1 = 32\,767$;
  \item con 32 bit: da $-2^{31} = -2\,147\,483\,648$ a
        $2^{31} - 1 = 2\,147\,483\,647$, circa più o meno 2 miliardi.
\end{itemize}

Il massimo è $2^{n-1}-1$ e non $2^{n-1}$ perché fra i valori con il primo
bit a 0 c'è anche lo zero. Le slide chiamano questi estremi $N_{min\_int}$ e
$N_{max\_int}$: ogni variabile \code{int i} soddisfa
$N_{min\_int} \le i \le N_{max\_int}$. In C sono disponibili come costanti
simboliche, \code{INT\_MIN} e \code{INT\_MAX}, definite nell'header
\code{limits.h}.

\subsection{L'overflow}

Che cosa succede se un calcolo supera \code{INT\_MAX}? Il programma seguente
somma due miliardi a due miliardi:

\begin{affianco}[0.5][c]
\begin{lstlisting}[style=c, name=overflow.c]
#include <stdio.h>

#define BIG 2000000000   /* 2 miliardi */

int main(void)
{
   int a, b = BIG, c = BIG;

   a = b + c;
   printf("a = %d\n", a);
   return 0;
}
\end{lstlisting}
\risultato
\begin{console}
$ gcc -Wall -o overflow overflow.c
$ ./overflow
a = -294967296
\end{console}
\end{affianco}

Il codice è sintatticamente corretto e il compilatore non segnala nulla; ma
in esecuzione $b + c$ vale 4 miliardi, più di \code{INT\_MAX}: è un
\term{integer overflow} (traboccamento). Il programma prosegue come se
niente fosse, con un risultato sbagliato. La figura~\ref{fig:overflow} mostra
da dove viene proprio quel numero: superato il massimo si ricomincia dal
minimo, come il contachilometri di un'auto che dopo 99999 torna a 00000 (la
ruota del capitolo~\ref{chap:calcolatore}). Contando, si arriva a
$4\,000\,000\,000 - 2^{32} = -294\,967\,296$.

\begin{figure}[H]
  \centering
  \begin{tikzpicture}[font=\Lato\scriptsize]
    \draw[line width=1.2pt, neutral!45!black] (0,0) -- (12,0);
    \foreach \xx/\lab in {0/{\code{INT\_MIN}\\$-2\,147\,483\,648$}, 6/{0},
                          12/{\code{INT\_MAX}\\$2\,147\,483\,647$}} {
      \draw[line width=0.8pt, neutral!45!black] (\xx,-0.12) -- (\xx,0.12);
      \node[anchor=north, align=center] at (\xx,-0.2) {\lab};
    }
    \fill[primary!70!black] (11.59,0) circle (2.3pt);
    \node[anchor=north, align=center, text=primary!60!black] at (10.2,-0.2) {\code{b}\\$2\,000\,000\,000$};
    \draw[primary!50!black, line width=0.5pt] (11.59,-0.06) -- (10.6,-0.25);
    \draw[wtarrowacc] (11.59,0.25) -- node[above, text=accent!70!black] {$+$ \code{c}} (12.7,0.25);
    \draw[accent!70!black, dashed, line width=0.9pt] (12.7,0.25) .. controls (14.2,2.3) and (-2.2,2.3) .. (-0.7,0.25);
    \node[fill=white, inner sep=2pt, text=accent!70!black] at (6,1.75) {superato \code{INT\_MAX}, si ricomincia da \code{INT\_MIN}};
    \draw[wtarrowacc] (-0.7,0.25) -- (5.18,0.25);
    \fill[accent!80!black] (5.18,0) circle (2.3pt);
    \node[anchor=north, align=center, text=accent!70!black] at (4.6,-0.2) {\code{a}\\$-294\,967\,296$};
  \end{tikzpicture}
  \caption{\code{b + c} su 32 bit: la somma esce a destra e rientra da
           sinistra.}
  \label{fig:overflow}
\end{figure}

Il risultato si spiega così: la somma viene fatta su 32 bit, e i bit che non
ci stanno si perdono. Per i tipi \code{unsigned} questo “giro” è garantito
dallo standard (i calcoli sono fatti modulo $2^n$); per i tipi con segno lo
standard parla invece di comportamento \textbf{indefinito}, e il valore
ottenuto qui è solo ciò che succede in pratica sui nostri computer. In ogni
caso tocca al programmatore assicurarsi che i valori restino
nell'intervallo, per esempio usando un tipo più grande.

È quello che fa la versione di \code{overflow.c} nella cartella degli
esempi: usa \code{unsigned long} (8 byte) e numeri ancora più grandi, 40
miliardi, e la somma, 80 miliardi, viene giusta.

\begin{affianco}[0.5][c]
\begin{lstlisting}[style=c, name=overflow.c]
#include <stdio.h>

#define BIG	40000000000

int main()
{
	unsigned long int a, b, c;

	b = c = BIG;

	a = b + c;

	printf("a = %ld\n", a);

	return 0;
}
\end{lstlisting}
\risultato
\begin{console}
$ ./overflow
a = 80000000000
\end{console}
(Per un \code{unsigned long} la conversione esatta di \code{printf} sarebbe
\code{\%lu}; qui \code{\%ld} funziona perché il valore sta anche in un
\code{long}.)
\end{affianco}

\section{\code{short}, \code{long} e \code{unsigned}}

\subsection{Più piccoli, più grandi, senza segno}

Oltre a \code{char} e \code{int}, gli altri tipi interi servono per usi più
specifici:

\begin{itemize}
  \item \code{short} occupa di solito meno memoria di un \code{int} (2
        byte): si usa quando lo spazio conta, per esempio in tabelle con
        milioni di numeri piccoli;
  \item \code{long} si usa per valori grandi (8 byte su Linux e macOS a 64
        bit);
  \item \code{unsigned} occupa quanto un \code{int}, ma non ha negativi: con
        4 byte va da 0 a $2^{32} - 1 = 4\,294\,967\,295$.
\end{itemize}

I limiti esatti del proprio sistema si leggono in \code{limits.h}. Questo
programma li stampa (l'output è quello di Linux su un PC a 64 bit):

\begin{lstlisting}[style=c, name=limiti.c]
#include <stdio.h>
#include <limits.h>

int main(void)
{
   printf("signed char:   %d .. %d\n", SCHAR_MIN, SCHAR_MAX);
   printf("unsigned char: %d .. %d\n", 0, UCHAR_MAX);
   printf("char:          %d .. %d\n", CHAR_MIN, CHAR_MAX);
   printf("short:         %d .. %d\n", SHRT_MIN, SHRT_MAX);
   printf("int:           %d .. %d\n", INT_MIN, INT_MAX);
   printf("unsigned:      %u .. %u\n", 0u, UINT_MAX);
   printf("long:          %ld .. %ld\n", LONG_MIN, LONG_MAX);
   printf("unsigned long: %lu .. %lu\n", 0ul, ULONG_MAX);
   return 0;
}
\end{lstlisting}

\begin{console}
$ ./limiti
signed char:   -128 .. 127
unsigned char: 0 .. 255
char:          -128 .. 127
short:         -32768 .. 32767
int:           -2147483648 .. 2147483647
unsigned:      0 .. 4294967295
long:          -9223372036854775808 .. 9223372036854775807
unsigned long: 0 .. 18446744073709551615
\end{console}

\info{Che cosa garantisce lo standard}{
  Lo standard del C non fissa le dimensioni dei tipi, ma solo dei
  \textbf{minimi}: \code{char} ha almeno 8 bit, \code{short} e \code{int}
  almeno 16, \code{long} almeno 32; inoltre un \code{short} non è più grande
  di un \code{int}, che non è più grande di un \code{long}. Le dimensioni
  reali le sceglie l'implementazione: su Windows, per esempio, \code{long}
  ha 4 byte come \code{int}, mentre su Linux e macOS a 64 bit ne ha 8. Chi
  ha bisogno di dimensioni esatte trova nell'header \code{stdint.h} (C99)
  tipi come \code{int32\_t} e \code{int64\_t}.}

\subsection{Il tipo delle costanti intere}

Anche le costanti hanno un tipo. A una costante intera scritta in decimale
viene dato il primo tipo in grado di contenerla fra \code{int} e \code{long}
(e, se non basta, \code{long long}, un tipo introdotto dal C99). Il tipo si
può anche imporre con un \textbf{suffisso}: \code{U} per \code{unsigned},
\code{L} per \code{long}, \code{UL} per \code{unsigned long} (vanno bene
anche le minuscole, ma la \code{l} si confonde facilmente con un 1). Per i
reali, come vedremo, i suffissi sono \code{F} e \code{L}. Ecco i tipi
assegnati dal compilatore ad alcune costanti:

\begin{center}
\renewcommand{\arraystretch}{1.12}
\begin{tabular}{@{}>{\FiraCodeNL}l l l@{}}
  \toprule
  \normalfont\textbf{costante} & \textbf{tipo} & \textbf{perché} \\
  \midrule
  37         & \code{int}           & sta in un \code{int} \\
  37U        & \code{unsigned}      & suffisso \code{U} \\
  37L        & \code{long}          & suffisso \code{L} \\
  37UL       & \code{unsigned long} & suffissi \code{U} e \code{L} \\
  2147483647 & \code{int}           & è \code{INT\_MAX}: sta ancora in un \code{int} \\
  2147483648 & \code{long}          & non sta più in un \code{int} \\
  3.7        & \code{double}        & le costanti reali sono \code{double} \\
  3.7F       & \code{float}         & suffisso \code{F} \\
  3.7L       & \code{long double}   & suffisso \code{L} \\
  'a'        & \code{int}           & le costanti carattere sono \code{int} \\
  \bottomrule
\end{tabular}
\end{center}

(Le slide indicano come candidati \code{int}, \code{long} e \code{unsigned}:
era all'incirca la regola del C più vecchio, il C89. Nel C99 una costante
decimale senza suffisso non diventa mai \code{unsigned}; possono diventarlo
quelle scritte in esadecimale o in ottale.)

\section{I tipi reali}

\subsection{\code{float}, \code{double} e \code{long double}}

L'ANSI C fornisce tre tipi reali: \code{float}, \code{double} e
\code{long double}. Le costanti reali, come \code{3.7}, sono di tipo
\code{double}; il tipo si può imporre con un suffisso: \code{3.7F} è un
\code{float}, \code{3.7L} un \code{long double}.

Di solito il compilatore dà più memoria ai \code{double} che ai
\code{float}, ma non è obbligatorio; su molte macchine un \code{float}
occupa 4 byte e un \code{double} 8. Allo stesso modo un \code{long double}
può avere più memoria di un \code{double} (sui PC ne occupa 16), ma non è
obbligatorio.

I valori che un tipo reale può contenere si descrivono con due grandezze:

\begin{itemize}
  \item la \textbf{precisione}: quante cifre significative, in notazione
        decimale, si possono rappresentare;
  \item l'\textbf{intervallo} (\textit{range}): il valore più grande e il più
        piccolo che si possono rappresentare.
\end{itemize}

\subsection{La virgola mobile}

Su molte macchine un \code{float} ha una precisione di 6 cifre significative
e un intervallo che va da circa $10^{-38}$ a $10^{38}$. Si può pensarlo come
un numero della forma
\[
  0.d_1 d_2 d_3 d_4 d_5 d_6 \times 10^{n}, \qquad -38 \le n \le 38,
\]
dove ogni $d_i$ è una cifra decimale (figura~\ref{fig:virgola-mobile}). Le
cifre stabiliscono la precisione, l'esponente $n$ dice dove va la virgola e
i suoi limiti stabiliscono l'intervallo. La virgola “galleggia” a seconda
dell'esponente: con le stesse sei cifre si rappresenta $123.451$ (esponente
3), $0.000123451$ (esponente $-3$) o $123\,451\,000$ (esponente 9). Da qui il
nome \textit{floating point}, virgola mobile.

\begin{figure}[H]
  \centering
  \begin{tikzpicture}[font=\FiraCodeNL\small]
    \noindici
    \node[anchor=east, font=\FiraCodeNL\small] at (-0.05,0) {123.451\Lato\ \ =\ \ \FiraCodeNL 0.};
    \celle[mcella, minimum width=0.6cm, minimum height=0.6cm, font=\FiraCodeNL\small\bfseries]{m}{(0,0)}{1,2,3,4,5,1}
    \node[anchor=west, font=\Lato\small] (x) at ([xshift=2pt]m-5.east) {$\times\ 10$};
    \node[mcella, anchor=south west, minimum width=0.45cm, minimum height=0.42cm,
          font=\FiraCodeNL\scriptsize\bfseries, fill=accent!12!white, draw=accent!65!black] (e)
      at ([xshift=-1pt, yshift=-3pt]x.north east) {3};
    \draw[wtbrace] ([yshift=-3pt]m-0.south west) -- ([yshift=-3pt]m-5.south east);
    \node[mnota, anchor=north, align=center] at ([yshift=-0.35cm]$(m-0.south)!0.5!(m-5.south)$)
      {le cifre significative:\\quante sono decide la \textbf{precisione}};
    \node[mnota, anchor=west, align=left] at ([xshift=0.35cm]e.east)
      {l'esponente dice dove va la virgola;\\i suoi limiti decidono l'\textbf{intervallo}};
  \end{tikzpicture}
  \caption{Un \code{float} come sei cifre più un esponente.}
  \label{fig:virgola-mobile}
\end{figure}

In realtà la macchina usa la base 2, non la base 10: le cifre sono binarie
(24 per un \code{float}), e “6 cifre decimali” è il modo di dire quanta
precisione garantiscono. I limiti esatti sono nell'header \code{float.h}:

\begin{lstlisting}[style=c, name=reali.c]
#include <stdio.h>
#include <float.h>

int main(void)
{
   printf("float:  min %e  max %e  cifre %d\n",
          FLT_MIN, FLT_MAX, FLT_DIG);
   printf("double: min %e  max %e  cifre %d\n",
          DBL_MIN, DBL_MAX, DBL_DIG);
   return 0;
}
\end{lstlisting}

\begin{console}
$ ./reali
float:  min 1.175494e-38  max 3.402823e+38  cifre 6
double: min 2.225074e-308  max 1.797693e+308  cifre 15
\end{console}

Un \code{double}, su molte macchine, ha quindi circa 15 cifre significative
e un intervallo da circa $10^{-308}$ a $10^{308}$: si può pensare come
$0.d_1 d_2 \dots d_{15} \times 10^{n}$ con $-308 \le n \le 308$.

\subsection{Ai confini del \code{float}}

Il programma delle lezioni assegna a un \code{float} valori sempre più
piccoli, poi due valori molto grandi e infine un numero con molte cifre:

\begin{affianco}[0.56][t]
\codicepiccolo
\begin{lstlisting}[style=c, name=float.c]
#include <stdio.h>

int main()
{
   float x;

   x = 1e-37; printf("x = 1e-37    x = %e\n", x);
   x = 1e-38; printf("x = 1e-38    x = %e\n", x);
   x = 1e-39; printf("x = 1e-39    x = %e\n", x);
   x = 1e-40; printf("x = 1e-40    x = %e\n", x);
   x = 1e-41; printf("x = 1e-41    x = %e\n", x);
   x = 1e-42; printf("x = 1e-42    x = %e\n", x);
   x = 1e-43; printf("x = 1e-43    x = %e\n", x);
   x = 1e-44; printf("x = 1e-44    x = %e\n", x);
   x = 1e-45; printf("x = 1e-45    x = %e\n", x);
   x = 1e-46; printf("x = 1e-46    x = %e\n", x);

   x = 1e38; printf("x = 1e38    x = %e\n", x);
   x = 1e39; printf("x = 1e39    x = %e\n", x);

   x = 123.45123451234512345;
   printf("x = 123.45123451234512345     x = %3.17f\n", x);
}
\end{lstlisting}
\risultato
\codicepiccolo
\begin{console}
$ ./float
x = 1e-37    x = 1.000000e-37
x = 1e-38    x = 9.999999e-39
x = 1e-39    x = 1.000000e-39
x = 1e-40    x = 9.999946e-41
x = 1e-41    x = 9.999666e-42
x = 1e-42    x = 1.000527e-42
x = 1e-43    x = 9.949219e-44
x = 1e-44    x = 9.809089e-45
x = 1e-45    x = 1.401298e-45
x = 1e-46    x = 0.000000e+00
x = 1e38    x = 1.000000e+38
x = 1e39    x = inf
x = 123.45123451234512345
    x = 123.45123291015625000
\end{console}
\end{affianco}

(L'ultima riga dell'output è stata spezzata in due per farla stare nella
pagina.) Che cosa è successo? La figura~\ref{fig:retta-float} riassume i
valori positivi di un \code{float}:

\begin{itemize}
  \item fino a \code{FLT\_MIN}, circa $1.18 \times 10^{-38}$, i valori sono
        rappresentati con tutta la precisione: \code{1e-38} diventa
        \code{9.999999e-39}, cioè il \code{float} più vicino;
  \item sotto \code{FLT\_MIN} i numeri si possono ancora rappresentare, ma
        con \textbf{sempre meno cifre esatte} (sono detti
        \textit{subnormali}): \code{9.999946e-41}, \code{9.949219e-44},
        \code{9.809089e-45};
  \item il più piccolo \code{float} positivo è circa $1.4 \times 10^{-45}$:
        \code{1e-45} diventa \code{1.401298e-45}, e \code{1e-46}, ancora più
        piccolo, diventa \textbf{0};
  \item sopra \code{FLT\_MAX}, circa $3.4 \times 10^{38}$, non c'è più
        posto: \code{1e39} diventa \code{inf}, \textbf{infinito}.
\end{itemize}

\begin{figure}[H]
  \centering
  \begin{tikzpicture}[font=\Lato\scriptsize,
      tick/.style={font=\FiraCodeNL\tiny, anchor=north, align=center, inner sep=1pt}]
    \fill[red-gradient-1!40!white] (0,0) rectangle (1.4,0.55);
    \fill[orange-gradient-1!50!white] (1.4,0) rectangle (5.0,0.55);
    \fill[green-gradient-1!45!white] (5.0,0) rectangle (11.0,0.55);
    \fill[red-gradient-1!40!white] (11.0,0) rectangle (12.6,0.55);
    \draw[neutral!50!black, line width=0.6pt] (0,0) rectangle (12.6,0.55)
      (1.4,0) -- (1.4,0.55) (5.0,0) -- (5.0,0.55) (11.0,0) -- (11.0,0.55);
    \node at (0.7,0.275) {\textbf{0}};
    \node at (3.2,0.275) {\textbf{subnormali}};
    \node at (8.0,0.275) {\textbf{normali}};
    \node at (11.8,0.275) {\textbf{inf}};
    \node[tick] at (1.4,-0.05) {1.4e-45};
    \node[tick] at (5.0,-0.05) {1.18e-38\\(FLT\_MIN)};
    \node[tick] at (11.0,-0.05) {3.4e+38\\(FLT\_MAX)};
    \node[anchor=south, align=center] at (0.7,0.6) {troppo piccoli:\\diventano 0};
    \node[anchor=south, align=center] at (3.2,0.6) {sempre meno\\cifre esatte};
    \node[anchor=south, align=center] at (8.0,0.6) {circa 7 cifre esatte};
    \node[anchor=south, align=center] at (11.8,0.6) {troppo grandi:\\infinito};
  \end{tikzpicture}
  \caption{I valori positivi di un \code{float} (la scala non è
           proporzionale). Per i negativi vale lo stesso, con il segno meno.}
  \label{fig:retta-float}
\end{figure}

L'ultima riga mostra la \textbf{precisione}: del numero
\code{123.45123451234512345} il \code{float} conserva le prime cifre, poi
inventa. La figura~\ref{fig:cifre} confronta le cifre con quelle di un
\code{double}, che di cifre giuste ne conserva più del doppio.

\begin{figure}[H]
  \centering
  \begin{tikzpicture}[font=\FiraCodeNL\small,
      lab/.style={font=\Lato\scriptsize, anchor=east, text=neutral!35!black},
      giu/.style={anchor=base west, inner sep=0pt, text=mok, font=\FiraCodeNL\small\bfseries},
      sba/.style={anchor=base west, inner sep=0pt, text=merr, font=\FiraCodeNL\small}]
    \node[lab] at (-0.2,0.05) {nel programma};
    \node[anchor=base west, inner sep=0pt] at (0,0) {123.45123451234512345};
    \node[lab] at (-0.2,-0.6) {float};
    \node[giu] (f1) at (0,-0.65) {123.45123};
    \node[sba] at (f1.base east) {291015625000};
    \node[lab] at (-0.2,-1.25) {double};
    \node[giu] (d1) at (0,-1.3) {123.45123451234512};
    \node[sba] at (d1.base east) {402};
    \node[mnota, anchor=west] at (4.6,-0.6) {8 cifre giuste (ne sono garantite 6)};
    \node[mnota, anchor=west] at (4.6,-1.25) {17 cifre giuste (ne sono garantite 15)};
  \end{tikzpicture}
  \caption{Lo stesso numero in un \code{float} e in un \code{double}: in
           verde le cifre esatte, in rosso quelle sbagliate.}
  \label{fig:cifre}
\end{figure}

\subsection{Ai confini del \code{double}}

La stessa prova con un \code{double} (nell'ultima riga le cifre stampate sono
17, con \code{\%.17e}). Questa volta alcune costanti sono così estreme che
\code{gcc} protesta già in compilazione:

\begin{affianco}[0.56][t]
\codicepiccolo
\begin{lstlisting}[style=c, name=double.c]
#include <stdio.h>

int main()
{
   double x;

   x = 1e-307; printf("x = 1e-307    x = %e\n", x);
   x = 1e-308; printf("x = 1e-308    x = %e\n", x);
   x = 1e-309; printf("x = 1e-309    x = %e\n", x);
   x = 1e-400; printf("x = 1e-400    x = %e\n", x);
   x = 1e-401; printf("x = 1e-401    x = %e\n", x);
   x = 1e-402; printf("x = 1e-402    x = %e\n", x);
   x = 1e-403; printf("x = 1e-403    x = %e\n", x);
   x = 1e-404; printf("x = 1e-404    x = %e\n", x);
   x = 1e-405; printf("x = 1e-405    x = %e\n", x);
   x = 1e-406; printf("x = 1e-406    x = %e\n", x);

   x = 1e308; printf("x = 1e308    x = %e\n", x);
   x = 1e309; printf("x = 1e309    x = %e\n", x);

   x = 123.45123451234512345;
   printf("x = 123.45123451234512345     x = %.17e\n", x);
}
\end{lstlisting}
\risultato
\codicepiccolo
\begin{console}
$ ./double
x = 1e-307    x = 1.000000e-307
x = 1e-308    x = 1.000000e-308
x = 1e-309    x = 1.000000e-309
x = 1e-400    x = 0.000000e+00
x = 1e-401    x = 0.000000e+00
x = 1e-402    x = 0.000000e+00
x = 1e-403    x = 0.000000e+00
x = 1e-404    x = 0.000000e+00
x = 1e-405    x = 0.000000e+00
x = 1e-406    x = 0.000000e+00
x = 1e308    x = 1.000000e+308
x = 1e309    x = inf
x = 123.45123451234512345
    x = 1.23451234512345124e+02
\end{console}
\end{affianco}

Le costanti da \code{1e-400} a \code{1e-406} sono troppo piccole anche per
un \code{double}, e \code{1e309} è troppo grande: il compilatore lo dice con
un avviso per ciascuna riga (qui ne mostriamo due su otto) e le sostituisce
con 0 e con infinito.

\begin{console}
$ gcc -Wall -o double double.c
double.c: In function (*@‘@*)main(*@’@*):
double.c:10:4: warning: floating constant truncated to zero [-Woverflow]
   10 |    x = 1e-400; printf("x = 1e-400    x = %e\n", x);
      |    ^
double.c:19:4: warning: floating constant exceeds range of (*@‘@*)double(*@’@*) [-Woverflow]
   19 |    x = 1e309; printf("x = 1e309    x = %e\n", x);
      |    ^
\end{console}

\subsection{I reali non sono esatti}

Con un numero finito di bit si rappresenta un numero finito di valori: la
maggior parte dei numeri reali \textbf{non è rappresentabile}, e viene
sostituita dal valore rappresentabile più vicino. Succede anche a numeri
semplicissimi come $0.1$: in base 2 è un numero periodico
($0.000110011001100\dots$), esattamente come $1/3 = 0.333\dots$ in base 10,
e da qualche parte va troncato. Di conseguenza le operazioni con i reali,
a differenza di quelle con gli interi, non sono necessariamente esatte:

\begin{affianco}[0.5][c]
\begin{lstlisting}[style=c, name=decimali.c]
#include <stdio.h>

int main(void)
{
   double a = 0.1, b = 0.2;

   printf("0.1       = %.20f\n", a);
   printf("0.1 + 0.2 = %.20f\n", a + b);
   printf("0.3       = %.20f\n", 0.3);
   if (a + b == 0.3)
      printf("uguali\n");
   else
      printf("diversi!\n");
   return 0;
}
\end{lstlisting}
\risultato
\begin{console}
$ ./decimali
0.1       = 0.10000000000000000555
0.1 + 0.2 = 0.30000000000000004441
0.3       = 0.29999999999999998890
diversi!
\end{console}
\end{affianco}

$0.1 + 0.2$ e $0.3$ differiscono nella diciassettesima cifra, ma per
\code{==} sono diversi. Per questo due reali \textbf{non si confrontano con
\code{==}}: si controlla che la loro differenza sia piccola, per esempio con
\code{fabs(a + b - 0.3) < 1e-9} (\code{fabs}, il valore assoluto, è descritta
più avanti). Per i calcoli di piccola taglia questi errori sono trascurabili;
nei problemi grandi (la risoluzione di grandi sistemi di equazioni
differenziali, la minimizzazione di funzioni complicate) si accumulano, e
controllarli è il compito dell'analisi numerica.

\section{Dare un nome ai tipi: \code{typedef}}

Il meccanismo del \code{typedef} associa un identificatore a un tipo; da quel
momento l'identificatore si usa per dichiarare variabili (o funzioni)
esattamente come i tipi standard:

\begin{affianco}[0.5][c]
\begin{lstlisting}[style=c]
typedef char          uppercase;
typedef int           inches;
typedef unsigned long size_t;
typedef unsigned char u_char;

uppercase u;
u_char    c;
inches    length, width;
\end{lstlisting}
\risultato
A che cosa serve?
\begin{itemize}
  \item ad \textbf{abbreviare} le dichiarazioni lunghe (\code{u\_char} al
        posto di \code{unsigned char});
  \item a scrivere nomi che dicono \textbf{l'uso} delle variabili:
        \code{length} e \code{width} sono lunghezze in pollici
        (\textit{inches});
  \item a gestire le \textbf{differenze fra sistemi}.
\end{itemize}
\end{affianco}

Il terzo punto si capisce con un esempio. Un programma scritto per un sistema
con \code{int} di 4 byte va portato su un sistema dove gli \code{int} hanno
2 byte, troppo pochi per i suoi dati. Se il programma dichiara le sue
variabili con un tipo proprio, basta cambiare \textbf{una riga}:

\begin{affianco}[0.5][c]
\begin{lstlisting}[style=c]
/* sul primo sistema */
typedef int tipoInt;
\end{lstlisting}
\risultato
\begin{lstlisting}[style=c]
/* sul secondo sistema */
typedef long int tipoInt;
\end{lstlisting}
\end{affianco}

La libreria standard fa proprio così: \code{size\_t}, il tipo dei risultati di
\code{sizeof}, è definito con un \code{typedef} e cambia da un sistema
all'altro. Il \code{typedef} sarà molto comodo con le strutture
(capitolo~\ref{chap:strutture}).

\section{L'operatore \code{sizeof}}

\subsection{Quanti byte}

Il C fornisce l'operatore unario \code{sizeof}, che calcola il numero di byte
necessari per rappresentare un oggetto. L'oggetto, fra parentesi, può essere
un tipo (come \code{int} o \code{float}), un'espressione (come
\code{a + b}), un array o una struttura.
Ha la stessa priorità degli altri operatori unari. Il programma delle lezioni
lo applica ai tipi fondamentali:

\begin{lstlisting}[style=c, name=sizeofTypes.c]
/* Compute the size of some fundamental types. */

#include <stdio.h>

int main(void)
{
   printf("The size of some fundamental types is computed.\n\n");
   printf("       char:%3ld byte \n", sizeof(char));
   printf("      short:%3ld bytes\n", sizeof(short));
   printf("        int:%3ld bytes\n", sizeof(int));
   printf("       long:%3ld bytes\n", sizeof(long));
   printf("   unsigned:%3ld bytes\n", sizeof(unsigned));
   printf("      float:%3ld bytes\n", sizeof(float));
   printf("     double:%3ld bytes\n", sizeof(double));
   printf("long double:%3ld bytes\n", sizeof(long double));
   return 0;
}
\end{lstlisting}

\begin{console}
$ ./sizeofTypes
The size of some fundamental types is computed.

       char:  1 byte
      short:  2 bytes
        int:  4 bytes
       long:  8 bytes
   unsigned:  4 bytes
      float:  4 bytes
     double:  8 bytes
long double: 16 bytes
\end{console}

\code{sizeof(char)} vale 1 per definizione; gli altri valori sono quelli di
Linux su un PC a 64 bit e possono cambiare altrove. Il risultato di
\code{sizeof} è di tipo \code{size\_t}, un intero senza segno; la conversione
di \code{printf} pensata apposta è \code{\%zu}, ma su Linux a 64 bit
\code{size\_t} è un \code{unsigned long} e \code{\%ld} funziona.

\subsection{Una sorpresa}

Con \code{sizeof} verifichiamo quello che avevamo anticipato sulle costanti
carattere:

\begin{affianco}[0.56][c]
\begin{lstlisting}[style=c, name=dimensioni.c]
#include <stdio.h>

int main(void)
{
   char c = 'a';

   printf("sizeof(c)     = %zu\n", sizeof(c));
   printf("sizeof('a')   = %zu\n", sizeof('a'));
   printf("sizeof(c + c) = %zu\n", sizeof(c + c));
   return 0;
}
\end{lstlisting}
\risultato
\begin{console}
$ ./dimensioni
sizeof(c)     = 1
sizeof('a')   = 4
sizeof(c + c) = 4
\end{console}
\end{affianco}

Una variabile \code{char} occupa 1 byte, ma la costante \code{'a'} ne
occupa 4: è un \code{int}. E anche \code{c + c} è un \code{int}, per una
regola che vedremo alla fine del capitolo, la \textit{promozione a intero}.
Si noti infine che \code{sizeof} guarda solo il \textbf{tipo}
dell'espressione, senza calcolarla: \code{sizeof(i++)} non incrementa
\code{i}.

\section{Leggere e scrivere un carattere alla volta}

\subsection{\code{getchar} e \code{putchar}}

Per leggere e scrivere un carattere alla volta la libreria standard offre
due funzioni, dichiarate in \code{stdio.h}:

\begin{itemize}
  \item \code{getchar()} legge il prossimo carattere dallo
        \textbf{standard input} (la tastiera, o un file se l'input è
        rediretto con \code{<}) e ne restituisce il codice;
  \item \code{putchar(c)} scrive sullo \textbf{standard output} il
        carattere di codice \code{c}.
\end{itemize}

Di solito sono realizzate come macro (direttive per il preprocessore), ma si
usano come funzioni. E lavorano con gli \code{int}, perché un \code{int} può
rappresentare qualsiasi carattere e, in più, un valore speciale che segnala
la fine dell'input.

Il programma seguente legge una sequenza di caratteri e scrive ciascuno
\textbf{due volte}. Lo proviamo redirigendo sullo standard input il file
\code{input.txt}:

\begin{affianco}[0.5][c]
\begin{lstlisting}[style=c, name=double_out.c]
#include <stdio.h>

int main(void)
{
   int c;

   while ((c = getchar()) != EOF) {
      putchar(c);
      putchar(c);
   }
   return 0;
}
\end{lstlisting}
\risultato
\begin{console}
$ cat input.txt
abcdefghijklmno
$ ./double_out < input.txt
aabbccddeeffgghhiijjkkllmmnnoo

$
\end{console}
\end{affianco}

Da dove viene la riga vuota? Il file termina con un a capo, e anche
\code{'\textbackslash n'} è un carattere: raddoppiato, produce una riga
vuota. Se invece si lancia il programma senza redirezione, i caratteri
arrivano dalla tastiera:

\begin{console}
$ ./double_out
<<ciao>>
cciiaaoo

<<Ctrl+D>>
$
\end{console}

Il programma riceve i caratteri solo dopo \tasti{Invio}, tutti insieme: il
terminale raccoglie la riga (così la si può correggere) e la consegna quando
la si conferma. Poi \code{getchar} li restituisce uno alla volta. Per finire,
come nel capitolo~\ref{chap:primi-programmi}, si preme \tasti{Ctrl,D} a inizio
riga.

\subsection{Il ciclo di lettura}

La condizione \code{(c = getchar()) != EOF} fa tre cose in una
(figura~\ref{fig:ciclo-getchar}): legge un carattere, lo assegna a \code{c} e
controlla che l'input non sia finito.

\begin{figure}[H]
  \centering
  \setlength{\fboxsep}{1.5pt}
  \renewcommand{\arraystretch}{1.6}
  \begin{tabular}{@{}c>{\FiraCodeNL\small}l>{\raggedright\arraybackslash}p{8.4cm}@{}}
    \wtpassonum{1} & while ((c = \colorbox{accent!16!white}{getchar()}) != EOF)
                   & \code{getchar()} legge un carattere, per esempio \code{a}, e restituisce il suo codice, 97 \\
    \wtpassonum{2} & while ((\colorbox{accent!16!white}{c = 97}) != EOF)
                   & il valore viene assegnato a \code{c}; l'assegnazione vale 97 \\
    \wtpassonum{3} & while (\colorbox{accent!16!white}{97 != EOF})
                   & 97 è diverso da \code{EOF}: la condizione è vera \\
    \wtpassonum{4} & while (1)
                   & si esegue il corpo, poi si ricomincia dal passo 1 \\
  \end{tabular}
  \caption{Il calcolo della condizione, passo per passo. Quando l'input
           finisce, \code{getchar()} restituisce \code{EOF}, il confronto
           \code{EOF != EOF} è falso e il ciclo termina.}
  \label{fig:ciclo-getchar}
\end{figure}

Tutte le parentesi sono necessarie. L'operatore \code{!=} ha la precedenza
su \code{=} (capitolo~\ref{chap:sintassi}), quindi senza le parentesi interne
\code{c = getchar() != EOF} verrebbe letto come
\code{c = (getchar() != EOF)}: a \code{c} andrebbe il risultato del
confronto, 1 o 0, e il programma stamperebbe, al posto di ogni carattere,
due caratteri di codice 1, invisibili. \code{gcc} con \code{-Wall} intuisce
l'errore e suggerisce le parentesi:

\begin{console}
$ gcc -Wall -o senza_parentesi senza_parentesi.c
senza_parentesi.c: In function (*@‘@*)main(*@’@*):
senza_parentesi.c:7:11: warning: suggest parentheses around assignment used as truth value [-Wparentheses]
    7 |    while (c = getchar() != EOF) {
      |           ^
\end{console}

\subsection{\code{EOF}}

\code{EOF} (\textit{end of file}, fine del file) è il valore che
\code{getchar} restituisce quando l'input è finito: alla fine di un file
rediretto, o quando si preme \tasti{Ctrl,D} a inizio riga. È una costante
simbolica definita in \code{stdio.h}, tipicamente come

\begin{lstlisting}[style=c]
#define EOF (-1)
\end{lstlisting}

ma su un altro sistema potrebbe valere diversamente: usare il nome
\code{EOF} invece di $-1$ garantisce la \textbf{portabilità}. Stampandolo:

\begin{affianco}[0.5][c]
\begin{lstlisting}[style=c, name=eof.c]
#include <stdio.h>

int main()
{
   printf("EOF = %d\n", EOF);
   printf("EOF = %c\n", EOF);
}
\end{lstlisting}
\risultato
\begin{console}
$ ./eof
EOF = -1
EOF = (*@\char"FFFD@*)
\end{console}
\end{affianco}

Con \code{\%d} si vede $-1$. Con \code{\%c} viene scritto un byte che vale
255 (\code{11111111}, la lettura senza segno di $-1$), che non corrisponde a
nessun carattere: il terminale mostra un simbolo di carattere non valido.

\subsection{Perché \code{int} e non \code{char}}

Nel programma, \code{c} è dichiarata \code{int} e non \code{char}. Il motivo
è che \code{getchar} può restituire \textbf{257 valori} diversi: i 256
possibili byte, da 0 a 255, più \code{EOF}. Un \code{char} ne contiene solo
256, e due valori finirebbero per confondersi (figura~\ref{fig:getchar-int}):
\code{char} può essere \code{signed} o \code{unsigned}, e \code{EOF} potrebbe
non valere $-1$; un \code{int} va bene in ogni caso.

\begin{figure}[H]
  \centering
  \begin{tikzpicture}[font=\FiraCodeNL\scriptsize,
      gc/.style={mcella, minimum width=0.95cm, minimum height=0.55cm}]
    \noindici
    \celle[gc]{g}{(0,0)}{EOF,0,1,2,\ldots,254,255}
    \node[mindice, anchor=north] at ([yshift=-1pt]g-0.south) {$-1$};
    \draw[mrif, draw=merr] (g-6.north) to[out=140, in=40]
      node[above, mnota, text=merr] {in un \code{signed char} il byte 255 diventa $-1$: sembra \code{EOF}} (g-0.north);
    \draw[mrif, draw=merr] ([yshift=-0.3cm]g-0.south) to[out=-40, in=-140]
      node[below, mnota, text=merr] {in un \code{unsigned char} $-1$ diventa 255: \code{c} non vale mai \code{EOF}} ([yshift=-0.1cm]g-6.south);
    \node[mnota, anchor=west, align=left] at ([xshift=0.5cm]g-6.east)
      {257 valori possibili:\\in un \code{int} restano tutti distinti \mok};
  \end{tikzpicture}
  \caption{I valori che \code{getchar} può restituire, e che cosa succede se
           li si mette in un \code{char}.}
  \label{fig:getchar-int}
\end{figure}

Con \code{unsigned char} il guasto è evidente:

\begin{affianco}[0.5][c]
\begin{lstlisting}[style=c, name=uchar.c]
#include <stdio.h>

int main(void)
{
   unsigned char c;

   while ((c = getchar()) != EOF) {
      putchar(c);
      putchar(c);
   }
   return 0;
}
\end{lstlisting}
\risultato
\begin{console}
$ ./uchar < input.txt
aabbccddeeffgghhiijjkkllmmnnoo

(*@\char"FFFD\char"FFFD\char"FFFD\char"FFFD\char"FFFD\char"FFFD\char"FFFD\char"FFFD\char"FFFD\char"FFFD\char"FFFD\char"FFFD\char"FFFD\char"FFFD\char"FFFD\char"FFFD\char"FFFD\char"FFFD\char"FFFD\char"FFFD@*)
^C
\end{console}
\end{affianco}

Dopo il contenuto del file il programma scrive all'infinito il byte 255
(mostrato come carattere non valido), finché non lo si ferma con
\tasti{Ctrl,C}. Il motivo: quando l'input finisce, \code{getchar}
restituisce $-1$, che in \code{c} diventa 255; il confronto
\code{255 != EOF} è sempre vero, e il ciclo non termina più. Il compilatore
lo segnala solo con l'opzione \code{-Wextra}, che attiva altri avvisi oltre a
quelli di \code{-Wall}:

\begin{console}
$ gcc -Wall -Wextra -o uchar uchar.c
uchar.c: In function (*@‘@*)main(*@’@*):
uchar.c:7:27: warning: comparison is always true due to limited range of data type [-Wtype-limits]
    7 |    while ((c = getchar()) != EOF) {
      |                           ^~
\end{console}

Con un \code{signed char} il guasto è più subdolo: il programma funziona,
salvo fermarsi in anticipo se l'input contiene un byte che vale 255.

\subsection{Maiuscole e minuscole}

I caratteri sono ordinati come gli interi che li rappresentano, e si
comportano come tali:

\begin{itemize}
  \item \code{'a' + 1} vale \code{'b'}, e \code{'b' + 1} vale \code{'c'};
  \item \code{'z' - 'a'} vale 25;
  \item \code{'A' < 'a'}, perché 65 è minore di 97;
  \item \code{'A' - 'a'} ha lo stesso valore di \code{'B' - 'b'}: la
        distanza fra maiuscola e minuscola è la stessa per tutte le lettere
        (figura~\ref{fig:maiuscole}).
\end{itemize}

\begin{figure}[H]
  \centering
  \begin{tikzpicture}[font=\FiraCodeNL\scriptsize,
      lc/.style={mcella, minimum width=0.8cm, minimum height=0.55cm, font=\FiraCodeNL\small\bfseries},
      cod/.style={font=\FiraCodeNL\tiny, text=neutral!45!black, inner sep=1pt}]
    \noindici
    \foreach \mi/\ma/\nmi/\nma [count=\k from 0] in {a/A/97/65, b/B/98/66, c/C/99/67,
        d/D/100/68, e/E/101/69, f/F/102/70, g/G/103/71, \ldots/\ldots/{}/{}, y/Y/121/89, z/Z/122/90} {
      \node[lc, fill=accent!12!white, draw=accent!60!black] (m\k) at (\k*1.15,0) {\mi};
      \node[cod, anchor=south] at (m\k.north) {\nmi};
      \node[lc, fill=primary!12!white] (M\k) at (\k*1.15,-1.5) {\ma};
      \node[cod, anchor=north] at (M\k.south) {\nma};
    }
    \foreach \k in {0,1,2,3,4,5,6,8,9}
      \draw[wtarrowacc] (m\k.south) -- (M\k.north);
    \node[mnota, anchor=west, align=left] at ([xshift=0.4cm]$(m9.east)!0.5!(M9.east)$)
      {ogni freccia somma\\\code{'A' - 'a'}, cioè $-32$};
  \end{tikzpicture}
  \caption{Minuscole e maiuscole (sotto ciascuna, il codice): ogni lettera
           dista 32 posizioni dalla sua maiuscola.}
  \label{fig:maiuscole}
\end{figure}

Su questa regolarità si basa \code{capitalize.c}, che trasforma in maiuscolo
le lettere minuscole e copia senza modifiche tutti gli altri caratteri:

\begin{affianco}[0.5][c]
\begin{lstlisting}[style=c, name=capitalize.c]
#include <stdio.h>

int main(void)
{
   int   c;

   while ((c = getchar()) != EOF)
      if (c >= 'a' && c <= 'z')
         putchar(c + 'A' - 'a');
      else
         putchar(c);
   return 0;
}
\end{lstlisting}
\risultato
\begin{console}
$ cat input.txt
abcDE6ygTr56
$ ./capitalize < input.txt
ABCDE6YGTR56
\end{console}
\end{affianco}

Per la \code{g}, per esempio: $103 + 65 - 97 = 71$, il codice di
\code{'G'}. Scrivere \code{'A' - 'a'} invece di $-32$ rende il programma più
leggibile e non obbliga a ricordare la tabella. Il corpo del \code{while} è
un'unica istruzione, l'\code{if}-\code{else}, quindi le graffe non servono.
(La libreria standard ha una funzione che fa la stessa cosa, \code{toupper},
dichiarata in \code{ctype.h}.)

\section{Funzioni matematiche}

\subsection{La libreria matematica}

Il linguaggio C non contiene funzioni matematiche predefinite: le funzioni
matematiche si trovano nella \textbf{libreria matematica}, che fa parte della
libreria standard. Per esempio \code{sqrt} (radice quadrata), \code{pow}
(potenza), \code{log}, \code{sin} e \code{cos}. In genere ricevono parametri
di tipo \code{double} e restituiscono un risultato \code{double}; i
prototipi sono in \code{math.h}, e su Linux la libreria va collegata con
l'opzione \code{-lm} (capitolo~\ref{chap:primi-programmi}). Il programma
seguente usa \code{sqrt}, che ha un parametro, e \code{pow}, che ne ha due:

\begin{affianco}[0.56][t]
\codicepiccolo
\begin{lstlisting}[style=c, name=radice.c]
#include <math.h>
#include <stdio.h>

int main(void)
{
   double x;

   printf("\n%s\n%s\n%s\n\n",
      "The square root of x and x raised",
      "to the x power will be computed.",
      "---");
   while (1) {                  /* do it forever */
      printf("Input x:  ");
      scanf("%lf", &x);
      if (x >= 0.0)
         printf("\n%15s%22.15e\n%15s%22.15e\n%15s%22.15e\n\n",
            "x = ", x,
            "sqrt(x) = ", sqrt(x),
            "pow(x, x) = ", pow(x, x));
      else
         printf("\nSorry, your number must be nonnegative.\n\n");
   }
   return 0;
}
\end{lstlisting}
\risultato
\begin{console}
$ gcc -Wall -o radice radice.c -lm
$ ./radice

The square root of x and x raised
to the x power will be computed.
---

Input x:  <<2>>

           x =  2.000000000000000e+00
     sqrt(x) =  1.414213562373095e+00
   pow(x, x) =  4.000000000000000e+00

Input x:  <<-1>>

Sorry, your number must be nonnegative.

Input x:  ^C
\end{console}
\end{affianco}

\begin{itemize}
  \item La prima \code{printf} stampa tre stringhe passate come argomenti, con
        la conversione \code{\%s}: è un modo per spezzare un testo lungo.
  \item \code{\%15s} scrive una stringa in 15 caratteri, allineata a destra;
        \code{\%22.15e} scrive un numero in notazione scientifica con 15
        cifre dopo il punto, in 22 caratteri (capitolo~\ref{chap:primi-programmi}).
        Per questo i risultati sono incolonnati.
  \item \code{while (1)} non termina mai: si esce con \tasti{Ctrl,C}. Qui
        \tasti{Ctrl,D} non basta: \code{scanf} non legge più nulla,
        \code{x} conserva l'ultimo valore e il ciclo si ripete all'infinito.
  \item Nelle slide l'esponente è scritto con tre cifre, \code{e+000}: è il
        formato usato su Windows; su Linux le cifre sono due.
\end{itemize}

\subsection{Il valore assoluto}

Per il valore assoluto la libreria standard ha due funzioni:
\code{abs}, per gli \code{int}, con il prototipo in \code{stdlib.h}, e
\code{fabs}, per i \code{double}, con il prototipo in \code{math.h}.
Scambiarle è un errore insidioso:

\begin{affianco}[0.5][c]
\begin{lstlisting}[style=c, name=assoluto.c]
#include <stdio.h>
#include <stdlib.h>
#include <math.h>

int main(void)
{
   printf("abs(-7)    = %d\n", abs(-7));
   printf("fabs(-2.7) = %f\n", fabs(-2.7));
   printf("abs(-2.7)  = %d\n", abs(-2.7));
   return 0;
}
\end{lstlisting}
\risultato
\begin{console}
$ ./assoluto
abs(-7)    = 7
fabs(-2.7) = 2.700000
abs(-2.7)  = 2
\end{console}
\end{affianco}

\code{abs(-2.7)} vale 2: il prototipo di \code{abs} dice che il parametro è
un \code{int}, quindi $-2.7$ viene convertito in intero ($-2$) prima della
chiamata. Con \code{-Wall} non c'è alcun avviso; con \code{-Wextra} sì:

\begin{console}
$ gcc -Wall -Wextra -o assoluto assoluto.c -lm
assoluto.c: In function (*@‘@*)main(*@’@*):
assoluto.c:9:32: warning: using integer absolute value function (*@‘@*)abs(*@’@*) when argument is of floating-point type (*@‘@*)double(*@’@*) [-Wabsolute-value]
    9 |    printf("abs(-2.7)  = %d\n", abs(-2.7));
      |                                ^~~
\end{console}

\section{Conversioni e cast}

Un'espressione aritmetica come \code{x + y} ha un valore e un tipo. Se
\code{x} e \code{y} sono entrambi \code{int}, anche \code{x + y} è un
\code{int}. Ma che succede se i tipi sono diversi, o più piccoli di un
\code{int}? Il C applica delle \textbf{conversioni implicite}, secondo regole
precise; e il programmatore può chiederne di \textbf{esplicite}, con il
\textit{cast}.

\subsection{La promozione a intero}

Un \code{char} o uno \code{short}, con o senza segno, si possono usare in
qualunque espressione in cui si può usare un \code{int} o un
\code{unsigned int}. Nelle espressioni vengono infatti \textbf{promossi}:
se tutti i valori del tipo originale sono rappresentabili da un \code{int}
(ed è il caso dei PC), il valore viene convertito in \code{int}; altrimenti
in \code{unsigned int}. Il processore, del resto, lavora in modo naturale su
valori della dimensione di un \code{int}.

Così, se \code{x} e \code{y} sono \code{short}, \code{x + y} è un
\code{int}; e in

\begin{lstlisting}[style=c]
char c = 'A';
printf("%c\n", c);
\end{lstlisting}

il parametro passato a \code{printf} è un'espressione, \code{c}, che per la
promozione a intero è di tipo \code{int}, con valore 65. Lo stesso motivo
rende \code{sizeof(c + c)} uguale a 4.

\subsection{Le conversioni aritmetiche}

Quando un operatore binario (\code{+}, \code{-}, \code{*}, \code{/}, un
confronto\dots) ha operandi di tipo diverso, prima del calcolo uno dei due
viene convertito nel tipo dell'altro. Con \code{int i;} e \code{float f;},
in \code{i + f} il valore di \code{i} viene convertito in \code{float}, e il
risultato è un \code{float}. La regola generale è che l'operando di tipo
\textbf{“più piccolo”} viene convertito nel tipo \textbf{“più grande”}
(figura~\ref{fig:scala}). Nell'ordine:

\begin{enumerate}
  \item se uno dei due operandi è \code{long double}, l'altro viene
        convertito in \code{long double};
  \item altrimenti, se uno è \code{double}, l'altro diventa \code{double};
  \item altrimenti, se uno è \code{float}, l'altro diventa \code{float};
  \item altrimenti si applica la promozione a intero a entrambi, e poi:
        \begin{itemize}
          \item se uno è \code{unsigned long}, l'altro diventa
                \code{unsigned long};
          \item altrimenti, se uno è \code{long} e l'altro
                \code{unsigned}: se un \code{long} può rappresentare tutti i
                valori di un \code{unsigned}, l'\code{unsigned} diventa
                \code{long}, altrimenti entrambi diventano
                \code{unsigned long};
          \item altrimenti, se uno è \code{long}, l'altro diventa
                \code{long};
          \item altrimenti, se uno è \code{unsigned}, l'altro diventa
                \code{unsigned};
          \item altrimenti entrambi sono \code{int}.
        \end{itemize}
\end{enumerate}

\begin{figure}[H]
  \centering
  \begin{tikzpicture}[font=\FiraCodeNL\scriptsize,
      gr/.style={draw=primary!60!black, fill=primary!6!white, rounded corners=2pt,
        line width=0.7pt, minimum width=2.7cm, minimum height=0.5cm},
      grr/.style={gr, draw=accent!60!black, fill=accent!8!white}]
    \foreach \testo [count=\k from 0] in {int, unsigned, long, unsigned long}
      \node[gr] (s\k) at (0,\k*0.8) {\testo};
    \foreach \testo [count=\k from 4] in {float, double, long double}
      \node[grr] (s\k) at (0,\k*0.8) {\testo};
    \foreach \k [evaluate=\k as \kk using int(\k+1)] in {0,...,5}
      \draw[wtarrow] (s\k.north) -- (s\kk.south);
    \node[wtnodeplain, font=\FiraCodeNL\scriptsize] (cs) at (0,-1.15) {char, short};
    \draw[wtarrow, draw=accent!70!black] (cs.north) -- node[right, mnota] {promozione a intero} (s0.south);
    \draw[wtbrace] ([xshift=0.3cm]s0.south east) -- ([xshift=0.3cm]s3.north east);
    \node[mnota, anchor=west] at ([xshift=0.55cm]$(s0.east)!0.5!(s3.east)$) {tipi interi};
    \draw[wtbrace] ([xshift=0.3cm]s4.south east) -- ([xshift=0.3cm]s6.north east);
    \node[mnota, anchor=west] at ([xshift=0.55cm]s5.east) {tipi reali};
    \node[mnota, anchor=east, align=right] at ([xshift=-0.5cm]s3.west)
      {fra due operandi, quello più in basso\\viene convertito nel tipo dell'altro};
    \node[mnota, anchor=east, align=right, text=neutral!45!black] at ([xshift=-0.5cm]$(s1.west)!0.5!(s2.west)$)
      {\code{unsigned} con \code{long}: \code{long} su Linux e macOS\\a 64 bit, \code{unsigned long} su Windows};
  \end{tikzpicture}
  \caption{La scala delle conversioni aritmetiche (come su Linux a 64 bit).}
  \label{fig:scala}
\end{figure}

Ecco l'esempio delle slide. Con le dichiarazioni

\begin{lstlisting}[style=c]
char c;  short s;  int i;  long l;  unsigned u;  unsigned long ul;
float f;  double d;  long double ld;
\end{lstlisting}

le espressioni seguenti hanno questi tipi (verificati con il compilatore,
su Linux a 64 bit):

\begin{center}
\renewcommand{\arraystretch}{1.12}
\begin{tabular}{@{}>{\FiraCodeNL}l l@{\hspace{1.4cm}}>{\FiraCodeNL}l l@{}}
  \toprule
  \normalfont\textbf{espressione} & \textbf{tipo} & \normalfont\textbf{espressione} & \textbf{tipo} \\
  \midrule
  c - s / i   & \code{int}    & u * 7 - i   & \code{unsigned} \\
  u * 2.0 + 1 & \code{double} & f * 7 - i   & \code{float} \\
  c + 3       & \code{int}    & 7 * s * ul  & \code{unsigned long} \\
  c + 5.0     & \code{double} & ld + c      & \code{long double} \\
  d + s       & \code{double} & u - ul      & \code{unsigned long} \\
  2 * i / l   & \code{long}   & u - l       & dipende dal sistema \\
  \bottomrule
\end{tabular}
\end{center}

Qualche caso commentato: in \code{c - s / i}, \code{c} e \code{s} vengono
promossi a \code{int} e tutto il calcolo è fra \code{int}; in
\code{u * 7 - i} la costante 7 e poi \code{i} vengono convertiti in
\code{unsigned}; in \code{u * 2.0 + 1} la costante \code{2.0} è un
\code{double}, e trascina tutto in \code{double}. L'ultimo dipende davvero dal
sistema: su Linux e macOS a 64 bit un \code{long} (8 byte) contiene tutti
i valori di un \code{unsigned} (4 byte), quindi \code{u - l} è un
\code{long}; su Windows, dove \code{long} ha 4 byte, è un
\code{unsigned long}.

\error{Interi con e senza segno}{
  Mescolare \code{int} e \code{unsigned} può dare risultati sorprendenti.
  In \code{-1 < 1u}, per la regola vista, $-1$ viene convertito in
  \code{unsigned} e diventa $4\,294\,967\,295$: il confronto è
  \textbf{falso}. Un programma con \code{int i = -1; unsigned u = 1;} e
  \code{if (i < u)} sceglie quindi il ramo \code{else}. Anche qui l'avviso
  arriva solo con \code{-Wextra}: \textit{comparison of integer expressions
  of different signedness}. Conviene compilare sempre con
  \code{-Wall -Wextra}.}

\subsection{Le conversioni nelle assegnazioni}

Anche l'assegnazione converte: il valore viene trasformato nel tipo della
variabile che lo riceve. Da reale a intero la parte decimale viene
\textbf{troncata}, non arrotondata (cioè si va verso lo zero); da intero a
reale non si perde nulla.

\begin{affianco}[0.5][c]
\begin{lstlisting}[style=c, name=tronca.c]
#include <stdio.h>

int main(void)
{
   int i = 3.9;
   int j = -3.9;
   double x = 7;

   printf("i = %d   j = %d   x = %f\n", i, j, x);
   return 0;
}
\end{lstlisting}
\risultato
\begin{console}
$ ./tronca
i = 3   j = -3   x = 7.000000
\end{console}
\end{affianco}

\subsection{Il cast}

Le conversioni viste finora sono \textbf{implicite}: le fa il compilatore da
solo. Esistono anche le conversioni \textbf{esplicite}, chiamate
\term{cast}: si scrive il tipo desiderato fra parentesi davanti a
un'espressione.

\begin{lstlisting}[style=c]
int i;
double f;

f = (double) i;
\end{lstlisting}

La variabile \code{i} \textbf{non viene modificata}: cambia il tipo del
valore prodotto dall'espressione \code{i}. L'uso tipico è evitare la
divisione intera, per esempio nel calcolo di una media:

\begin{affianco}[0.56][c]
\begin{lstlisting}[style=c, name=media.c]
#include <stdio.h>

int main(void)
{
   int somma = 7, n = 2;
   double media;

   media = somma / n;
   printf("somma / n            = %f\n", media);
   media = (double) somma / n;
   printf("(double) somma / n   = %f\n", media);
   media = (double) (somma / n);
   printf("(double) (somma / n) = %f\n", media);
   printf("somma vale ancora %d\n", somma);
   return 0;
}
\end{lstlisting}
\risultato
\begin{console}
$ ./media
somma / n            = 3.000000
(double) somma / n   = 3.500000
(double) (somma / n) = 3.000000
somma vale ancora 7
\end{console}
\end{affianco}

Nel primo caso la divisione è fra due \code{int} e vale 3; il risultato viene
convertito in \code{double} solo dopo, nell'assegnazione. Il cast è un
operatore unario e ha la precedenza su \code{/}: in
\code{(double) somma / n} si converte prima \code{somma} in $7.0$, poi la
divisione è fra un \code{double} e un \code{int}, quindi anche \code{n} viene
convertito e il risultato è $3.5$. Nel terzo caso le parentesi fanno
eseguire prima la divisione intera: il cast arriva quando il danno è fatto
(figura~\ref{fig:cast}).

\begin{figure}[H]
  \centering
  \begin{tikzpicture}[font=\FiraCodeNL\small,
      op/.style={circle, draw=accent!65!black, fill=accent!8!white, line width=0.8pt,
        minimum size=0.62cm, inner sep=0pt, font=\FiraCodeNL\small\bfseries},
      ca/.style={draw=accent!65!black, fill=accent!8!white, line width=0.8pt,
        rounded corners=0.3cm, minimum height=0.62cm, inner xsep=5pt, font=\FiraCodeNL\scriptsize\bfseries},
      fo/.style={wtnodeplain, font=\FiraCodeNL\scriptsize, minimum width=1.2cm},
      val/.style={font=\Lato\scriptsize\bfseries, text=green-gradient-6!50!black, anchor=west},
      ord/.style={circle, fill=primary!75!black, text=white, font=\Lato\tiny\bfseries,
        inner sep=0pt, minimum size=0.36cm},
      tit/.style={font=\FiraCodeNL\small}]
    % (double) somma / n
    \node[tit] at (0.4,1.0) {(double) somma / n};
    \node[op] (d1) at (0.4,0) {/};
    \node[ca] (c1) at (-0.9,-1.1) {(double)};
    \node[fo] (s1) at (-0.9,-2.2) {somma\\[-1pt]{\color{neutral!50!black}(7)}};
    \node[fo] (n1) at (1.7,-1.1) {n\\[-1pt]{\color{neutral!50!black}(2)}};
    \foreach \a/\b in {d1/c1, d1/n1, c1/s1} \draw[fsramo] (\a) -- (\b);
    \node[val] at ([xshift=4pt]c1.east) {= 7.0};
    \node[val] at ([xshift=4pt]d1.east) {= 7.0 / 2.0 = 3.5};
    \node[ord] at ([shift={(-0.25,0.3)}]c1.north west) {1};
    \node[ord] at ([shift={(-0.42,0.32)}]d1.center) {2};
    % (double) (somma / n)
    \node[tit] at (8.0,1.0) {(double) (somma / n)};
    \node[ca] (c2) at (8.0,0) {(double)};
    \node[op] (d2) at (8.0,-1.1) {/};
    \node[fo] (s2) at (6.8,-2.2) {somma\\[-1pt]{\color{neutral!50!black}(7)}};
    \node[fo] (n2) at (9.2,-2.2) {n\\[-1pt]{\color{neutral!50!black}(2)}};
    \foreach \a/\b in {c2/d2, d2/s2, d2/n2} \draw[fsramo] (\a) -- (\b);
    \node[val] at ([xshift=4pt]d2.east) {= 7 / 2 = 3 (divisione intera)};
    \node[val] at ([xshift=4pt]c2.east) {= 3.0};
    \node[ord] at ([shift={(-0.42,0.32)}]d2.center) {1};
    \node[ord] at ([shift={(-0.25,0.3)}]c2.north west) {2};
  \end{tikzpicture}
  \caption{Dove sta il cast cambia tutto: a sinistra converte \code{somma}
           prima della divisione, a destra il risultato già troncato.}
  \label{fig:cast}
\end{figure}

\insintesi{
  Il C ha dodici tipi fondamentali: i tipi \textbf{interi} (\code{char},
  \code{short}, \code{int}, \code{long}, con e senza segno) e i tipi
  \textbf{reali} (\code{float}, \code{double}, \code{long double}). Un
  \code{char} è un intero di 1 byte che contiene il \textbf{codice} di un
  carattere (ASCII: le cifre, le maiuscole e le minuscole hanno codici
  consecutivi); le \textbf{sequenze di escape} scrivono i caratteri
  speciali. Ogni tipo ha un intervallo di valori (\code{limits.h},
  \code{float.h}): superarlo con gli interi provoca un \textbf{overflow}
  silenzioso, i reali hanno una \textbf{precisione} limitata e non sono
  esatti, quindi non si confrontano con \code{==}. \code{sizeof} dà i byte
  di un tipo o di un'espressione, \code{typedef} dà un nuovo nome a un tipo.
  \code{getchar} e \code{putchar} leggono e scrivono un carattere; la
  variabile che riceve \code{getchar} deve essere \code{int}, per
  distinguere \code{EOF}. Le funzioni matematiche sono in \code{math.h}
  (\code{-lm}). Nelle espressioni \code{char} e \code{short} diventano
  \code{int}, e l'operando di tipo più piccolo viene convertito nel tipo
  dell'altro; il \textbf{cast} \code{(tipo)} converte esplicitamente.}

\section{Riferimenti}

\begin{itemize}
  \item \textbf{A. Kelley, I. Pohl}, \textit{A Book on C}: capitolo~3 (i
        tipi di dato fondamentali).
  \item \textbf{Arithmetic types} --- cppreference.com. \\
        \urlx{https://en.cppreference.com/w/c/language/arithmetic_types}
  \item \textbf{Implicit conversions} --- cppreference.com. \\
        \urlx{https://en.cppreference.com/w/c/language/conversion}
  \item \textbf{D. Goldberg}, \textit{What Every Computer Scientist Should
        Know About Floating-Point Arithmetic} (il classico sugli errori dei
        reali). \\
        \urlx{https://docs.oracle.com/cd/E19957-01/806-3568/ncg_goldberg.html}
\end{itemize}


  \chapter{Il flusso di controllo}
\label{chap:flusso}

Le istruzioni di un programma vengono eseguite normalmente in sequenza, una
dopo l'altra. Quasi sempre, però, bisogna alterare questo flusso: per
\textbf{scegliere} fra azioni diverse (con \code{if}, \code{if}-\code{else} e
\code{switch}) o per \textbf{ripetere} un'azione (con \code{while},
\code{do}-\code{while} e \code{for}). Nel capitolo~\ref{chap:strutturata}
abbiamo visto queste strutture nei diagrammi di flusso, e nel
capitolo~\ref{chap:primi-programmi} le prime istruzioni del C; qui le vediamo
tutte, a partire da ciò che le governa: le \textbf{condizioni}.

\section{Vero e falso in C}

\subsection{Le condizioni sono numeri}

Il flusso di controllo è regolato da espressioni logiche, che possono essere
vere (\textit{TRUE}) o false (\textit{FALSE}). Il C non ha un tipo apposito
per i valori logici (il C99 ha aggiunto \code{\_Bool}, ma la regola resta
quella che segue): vero e falso sono \textbf{numeri}.

\begin{itemize}
  \item \textbf{Falso} è un valore nullo: \code{0}, \code{0.0}, il carattere
        \code{'\textbackslash 0'}.
  \item \textbf{Vero} è qualunque valore diverso da zero: 1, ma anche 7,
        $-3$ o \code{'A'}.
  \item Gli operatori che producono un valore logico restituiscono sempre un
        \code{int}: \textbf{1} se il risultato è vero, \textbf{0} se è falso.
\end{itemize}

Le condizioni si costruiscono con tre famiglie di operatori:

\begin{center}
\renewcommand{\arraystretch}{1.12}
\begin{tabular}{@{}l l>{\FiraCodeNL}c@{\hspace{1.4cm}}l l>{\FiraCodeNL}c@{}}
  \toprule
  \textbf{famiglia} & \textbf{operatore} & \normalfont\textbf{in C} &
  \textbf{famiglia} & \textbf{operatore} & \normalfont\textbf{in C} \\
  \midrule
  relazionali  & minore             & <  & di uguaglianza & uguale a           & == \\
               & maggiore           & >  &                & diverso da         & != \\
               & minore o uguale    & <= & logici         & negazione (unario) & !  \\
               & maggiore o uguale  & >= &                & e (\textit{and})   & \&\& \\
               &                    &    &                & o (\textit{or})    & || \\
  \bottomrule
\end{tabular}
\end{center}

Sono tutti operatori binari (con due operandi), tranne la negazione
\code{!}, che è unaria.

\section{Gli operatori relazionali}

\subsection{Confrontare due espressioni}

Gli operatori relazionali prendono due espressioni come operandi:
\code{a < 3}, \code{a > b}, \code{-1.3 >= (2.0 * x + 3.3)}. Sono scritti con
uno o due caratteri, sempre in quest'ordine e senza spazi in mezzo: per
questo \code{a =< b} è sbagliato (l'ordine è \code{<=}) e \code{a < = b}
pure (lo spazio spezza il token, capitolo~\ref{chap:sintassi}). Attenzione
anche ad \code{a >> b}: è corretto, ma \code{>>} è un altro operatore (lo
spostamento dei bit), con tutt'altro significato.

Come li calcola la macchina? Su molte macchine \code{a < b} viene valutato
come \code{a - b < 0}, e lo stesso per gli altri: si calcola la differenza
e se ne guarda il segno.

\begin{center}
\renewcommand{\arraystretch}{1.12}
\begin{tabular}{@{}l>{\FiraCodeNL}c>{\FiraCodeNL}c>{\FiraCodeNL}c>{\FiraCodeNL}c@{}}
  \toprule
  \textbf{\code{a - b}} & \normalfont\textbf{a > b} & \normalfont\textbf{a >= b} &
  \normalfont\textbf{a < b} & \normalfont\textbf{a <= b} \\
  \midrule
  positivo & 1 & 1 & 0 & 0 \\
  zero     & 0 & 1 & 0 & 1 \\
  negativo & 0 & 0 & 1 & 1 \\
  \bottomrule
\end{tabular}
\end{center}

\subsection{Qualche esempio}

Con le dichiarazioni

\begin{lstlisting}[style=c]
char   c = 'w';
int    i = 1, j = 2, k = -7;
double x = 7e+33, y = 0.001;
\end{lstlisting}

ecco alcune espressioni, l'espressione equivalente con tutte le parentesi
(che mostra in che ordine si calcola, capitolo~\ref{chap:sintassi}) e il
valore, verificato con il compilatore:

\begin{center}
\renewcommand{\arraystretch}{1.12}
\begin{tabular}{@{}>{\FiraCodeNL}l>{\FiraCodeNL}l c@{}}
  \toprule
  \normalfont\textbf{espressione} & \normalfont\textbf{equivalente} & \textbf{valore} \\
  \midrule
  'a' + 1 < c          & ('a' + 1) < c                 & 1 \\
  -i - 5 * j >= k + 1  & ((-i) - (5 * j)) >= (k + 1)   & 0 \\
  3 < j < 5            & (3 < j) < 5                   & 1 \\
  x - 3.333 <= x + y   & (x - 3.333) <= (x + y)        & 1 \\
  x < x + y            & x < (x + y)                   & 0 \\
  \bottomrule
\end{tabular}
\end{center}

\code{'a' + 1} vale 98 e \code{c} vale 119 (il codice di \code{'w'}):
quindi 1. Nella seconda, $-1 - 10 = -11$ non è maggiore o uguale a $-6$:
quindi 0. Le ultime tre meritano una spiegazione a parte.

\subsection{\code{a < b < c} non è la matematica}

In matematica $a < b < c$ significa “$a < b$ \textbf{e} $b < c$”. In C
l'espressione è sintatticamente corretta, ma è interpretata come
\code{(a < b) < c}: prima si calcola \code{a < b}, che vale 0 o 1, e poi si
confronta \textbf{quel} risultato con \code{c}. Per questo \code{3 < j < 5}
vale 1: \code{3 < j} è falso (0), e \code{0 < 5} è vero. Con numeri scelti
apposta, il risultato è chiaramente sbagliato:

\begin{affianco}[0.5][c]
\begin{lstlisting}[style=c, name=catena.c]
#include <stdio.h>

int main(void)
{
   int a = 5, b = 2, c = 3;

   if (a < b < c)
      printf("5 < 2 < 3 è vero?!\n");
   if (a < b && b < c)
      printf("questo non viene stampato\n");
   return 0;
}
\end{lstlisting}
\risultato
\begin{console}
$ ./catena
5 < 2 < 3 è vero?!
\end{console}
\end{affianco}

\begin{console}
$ gcc -Wall -o catena catena.c
catena.c: In function (*@‘@*)main(*@’@*):
catena.c:7:10: warning: comparisons like (*@‘@*)X<=Y<=Z(*@’@*) do not have their mathematical meaning [-Wparentheses]
    7 |    if (a < b < c)
      |        ~~^~~
\end{console}

Il modo giusto di scrivere la condizione matematica è il secondo
\code{if}: \code{a < b \&\& b < c}, con l'operatore logico \textit{and} che
vedremo fra poco. Come si vede, \code{gcc} con \code{-Wall} si accorge
dell'errore.

\subsection{I reali e la precisione}

In matematica $x < x + y$ equivale a $y > 0$. Con i reali del computer non
sempre: se \code{x} è molto grande e \code{y} molto piccolo, l'incremento
può andare perso. È il caso di $x = 7 \cdot 10^{33}$ e $y = 0.001$:

\begin{affianco}[0.5][c]
\begin{lstlisting}[style=c, name=precisione.c]
#include <stdio.h>

int main()
{
   double x = 7e+33;
   double y = 0.001;

   if (x < x + y) printf("x < x + y\n");
   if (x > x + y) printf("x > x + y\n");
   if (x == x + y) printf("x = x + y\n");
}
\end{lstlisting}
\risultato
\begin{console}
$ ./precisione
x = x + y
\end{console}
\end{affianco}

Per scrivere esattamente $7 \cdot 10^{33} + 0.001$ servirebbero 37 cifre
significative; un \code{double} ne ha circa 16
(capitolo~\ref{chap:tipi}). L'incremento dato da \code{y} è troppo piccolo
rispetto a \code{x}: la precisione non basta a registrarlo, e la
rappresentazione di \code{x + y} è identica a quella di \code{x}.

\section{Gli operatori di uguaglianza}

\subsection{\code{==} e \code{!=}}

Gli operatori di uguaglianza \code{==} e \code{!=} sono binari e agiscono su
espressioni: \code{c == 'A'}, \code{k != -2}, \code{x + y == 3 * z - 7}. Non
sono invece esempi corretti:

\begin{itemize}
  \item \code{a = b}, che è un'\textbf{assegnazione};
  \item \code{a = = b - 1}, perché lo spazio spezza l'operatore;
  \item \code{(x + y) =! 44}, che il compilatore legge come
        \code{(x + y) = (!44)}: un'assegnazione a \code{x + y}, che non è una
        variabile. Il messaggio di errore lo dice a modo suo:
\end{itemize}

\begin{console}
$ gcc -Wall -o errori errori.c
errori.c: In function (*@‘@*)main(*@’@*):
errori.c:7:12: error: lvalue required as left operand of assignment
    7 |    (x + y) =! 44;
      |            ^
\end{console}

(Un \textit{lvalue} è qualcosa che può stare a sinistra di un'assegnazione,
come una variabile.) Anche l'uguaglianza si calcola con una differenza:
\code{expr1 == expr2} equivale a \code{expr1 - expr2 == 0}, e
\code{expr1 != expr2} a \code{expr1 - expr2 != 0}.

\begin{affianco}[0.4][c]
\begin{center}
\renewcommand{\arraystretch}{1.12}
\begin{tabular}{@{}l>{\FiraCodeNL}c>{\FiraCodeNL}c@{}}
  \toprule
  \textbf{\code{expr1 - expr2}} & \normalfont\textbf{==} & \normalfont\textbf{!=} \\
  \midrule
  zero     & 1 & 0 \\
  non zero & 0 & 1 \\
  \bottomrule
\end{tabular}
\end{center}
\risultato
\begin{center}
\renewcommand{\arraystretch}{1.12}
\begin{tabular}{@{}>{\FiraCodeNL}l>{\FiraCodeNL}l c@{}}
  \toprule
  \multicolumn{3}{@{}l}{con \code{int i = 1, j = 2, k = 3;}} \\
  \midrule
  i == j          & i - j == 0              & 0 \\
  i != j          & i - j != 0              & 1 \\
  i+j+k == -2*-k  & ((i+j)+k) == ((-2)*(-k)) & 1 \\
  \bottomrule
\end{tabular}
\end{center}
\end{affianco}

\subsection{L'errore più comune: \code{=} al posto di \code{==}}

Un errore frequentissimo è confondere \code{a == b} con \code{a = b},
scrivendo \code{if (a = b)} invece di \code{if (a == b)}. Il compilatore lo
accetta, perché un'assegnazione è un'espressione con un valore: il valore di
\code{b}. Così la condizione è vera o falsa \textbf{a seconda del solo valore
di \code{b}}, e per giunta \code{a} viene modificata:

\begin{affianco}[0.5][c]
\begin{lstlisting}[style=c, name=uguale.c]
#include <stdio.h>

int main(void)
{
   int a = 5, b = 7;

   if (a = b)
      printf("uguali?! ora a vale %d\n", a);
   return 0;
}
\end{lstlisting}
\risultato
\begin{console}
$ ./uguale
uguali?! ora a vale 7
\end{console}
\end{affianco}

Per fortuna \code{gcc} con \code{-Wall} sospetta l'errore, e suggerisce di
mettere l'assegnazione fra parentesi se è davvero voluta:

\begin{console}
$ gcc -Wall -o uguale uguale.c
uguale.c: In function (*@‘@*)main(*@’@*):
uguale.c:7:8: warning: suggest parentheses around assignment used as truth value [-Wparentheses]
    7 |    if (a = b)
      |        ^
\end{console}

\section{Gli operatori logici}

\subsection{La negazione}

L'operatore logico di negazione \code{!} è unario: se un'espressione vale 0,
la sua negata vale 1, e viceversa (se vale qualunque cosa diversa da 0, la
negata vale 0). Sono corretti \code{!a}, \code{!(x + 7.7)},
\code{!(a < b || c < d)}; non lo sono \code{a!} (l'ordine è sbagliato: il
\code{!} va prima) e \code{a  !=  b} se si voleva una negazione, perché
\code{!=} è un altro operatore.

In logica, negare due volte un'affermazione la lascia invariata:
\textit{not}(\textit{not} $A$) $= A$. In C non è sempre così: \code{!!5} vale
\textbf{1}, non 5. La prima negazione trasforma 5 in 0, la seconda trasforma
0 in 1: la doppia negazione “normalizza” qualunque valore vero a 1.

\begin{affianco}[0.36][c]
\begin{lstlisting}[style=c]
char   c = 'A';
int    i = 7, j = 7;
double x = 0.0, y = 2.3;
\end{lstlisting}
\risultato
\begin{center}
\renewcommand{\arraystretch}{1.12}
\begin{tabular}{@{}>{\FiraCodeNL}l>{\FiraCodeNL}l c@{}}
  \toprule
  \normalfont\textbf{espressione} & \normalfont\textbf{equivalente} & \textbf{valore} \\
  \midrule
  !c        & !c                 & 0 \\
  !(i - j)  & !(i - j)           & 1 \\
  !i - j    & (!i) - j           & $-7$ \\
  !!(x + y) & !(!(x + y))        & 1 \\
  !x * !!y  & (!x) * (!(!y))     & 1 \\
  \bottomrule
\end{tabular}
\end{center}
\end{affianco}

\code{!i - j} vale $-7$ perché la negazione, come tutti gli operatori unari,
ha la precedenza sulla sottrazione: \code{!i} vale 0, e $0 - 7 = -7$.

\subsection{\textit{And} e \textit{or}}

Gli operatori \code{\&\&} (\textit{and}, “e”) e \code{||} (\textit{or},
“o”) sono binari, agiscono su espressioni e restituiscono 1 (vero) o 0
(falso): \code{a \&\& b}, \code{a || b}, \code{!(a < b) \&\& c},
\code{3 \&\& (-2 * a + 7)}. Non sono corretti \code{a \&\&} (manca un
operando), \code{a |  | b} (gli spazi spezzano l'operatore) e \code{\&b}, che
è l'indirizzo di \code{b} (capitolo~\ref{chap:primi-programmi}). Il loro
comportamento è descritto dalla \textbf{tabella di verità}:

\begin{center}
\renewcommand{\arraystretch}{1.12}
\begin{tabular}{@{}l l>{\FiraCodeNL}c>{\FiraCodeNL}c@{}}
  \toprule
  \textbf{\code{expr1}} & \textbf{\code{expr2}} &
  \normalfont\textbf{expr1 \&\& expr2} & \normalfont\textbf{expr1 || expr2} \\
  \midrule
  zero     & zero     & 0 & 0 \\
  zero     & non zero & 0 & 1 \\
  non zero & zero     & 0 & 1 \\
  non zero & non zero & 1 & 1 \\
  \bottomrule
\end{tabular}
\end{center}

\code{\&\&} è vero solo se \textbf{entrambi} gli operandi sono veri;
\code{||} è vero se \textbf{almeno uno} lo è.

\subsection{Precedenza}

\code{\&\&} ha la precedenza su \code{||}, ed entrambi hanno priorità più
bassa degli operatori unari, aritmetici, relazionali e di uguaglianza;
l'associatività va da sinistra a destra. La figura~\ref{fig:priorita}
riassume l'ordine degli operatori che si usano nelle condizioni (la tabella
completa è nel capitolo~\ref{chap:sintassi}).

\begin{figure}[H]
  \centering
  \begin{tikzpicture}[font=\Lato\scriptsize,
      liv/.style={draw=primary!60!black, fill=primary!6!white, rounded corners=2pt,
        line width=0.7pt, minimum width=3.6cm, minimum height=0.5cm, font=\FiraCodeNL\scriptsize},
      des/.style={font=\Lato\scriptsize, anchor=west, text=neutral!35!black}]
    \foreach \oper/\testo [count=\k from 0] in {{!\ \ ++\ \ --\ \ -}/{unari},
        {*\ \ /\ \ \%}/{moltiplicativi}, {+\ \ -}/{additivi},
        {<\ \ <=\ \ >\ \ >=}/{relazionali}, {==\ \ !=}/{di uguaglianza},
        {\&\&}/{\textit{and} logico}, {||}/{\textit{or} logico},
        {=\ \ +=\ \ -=\ \ \ldots}/{di assegnazione}} {
      \node[liv] (l\k) at (0,-\k*0.62) {\oper};
      \node[des] at ([xshift=0.3cm]l\k.east) {\testo};
    }
    \draw[wtarrowacc] (-2.3,0.1) -- node[left, align=right, text=accent!70!black] {precedenza\\più alta} ++(0,-1.0);
    \draw[wtarrowacc] (-2.3,-3.3) -- node[left, align=right, text=accent!70!black] {precedenza\\più bassa} ++(0,-1.0);
  \end{tikzpicture}
  \caption{Gli operatori delle condizioni, dalla precedenza più alta alla più
           bassa.}
  \label{fig:priorita}
\end{figure}

Con le dichiarazioni \code{char c = 'B';}, \code{int i = 3, j = 3, k = 3;}
e \code{double x = 0.0, y = 2.3;}:

\begin{center}
\renewcommand{\arraystretch}{1.12}
\begin{tabular}{@{}>{\FiraCodeNL}l>{\FiraCodeNL}l c@{}}
  \toprule
  \normalfont\textbf{espressione} & \normalfont\textbf{equivalente} & \textbf{valore} \\
  \midrule
  i \&\& j \&\& k                 & (i \&\& j) \&\& k                       & 1 \\
  x || i \&\& j - 3               & x || (i \&\& (j - 3))                   & 0 \\
  i < j \&\& x < y                & (i < j) \&\& (x < y)                    & 0 \\
  i < j || x < y                  & (i < j) || (x < y)                      & 1 \\
  'A' <= c \&\& c <= 'Z'          & ('A' <= c) \&\& (c <= 'Z')              & 1 \\
  c - 1 == 'A' || c + 1 == 'Z'    & ((c - 1) == 'A') || ((c + 1) == 'Z')    & 1 \\
  \bottomrule
\end{tabular}
\end{center}

Per esempio \code{x || i \&\& j - 3}: \code{x} vale 0.0, falso; \code{j - 3}
vale 0, quindi \code{i \&\& 0} è falso; e “falso o falso” è falso, 0. La
quinta riga è il modo corretto di scrivere la condizione matematica
$\text{'A'} \le c \le \text{'Z'}$: “\code{c} è una lettera maiuscola”.

\subsection{La valutazione in cortocircuito}

La valutazione delle espressioni con \code{\&\&} e \code{||} si
\textbf{ferma appena il risultato è noto} (\textit{short circuit},
cortocircuito):

\begin{itemize}
  \item in \code{expr1 \&\& expr2}, se \code{expr1} è falsa il risultato è
        falso comunque, e \code{expr2} \textbf{non viene valutata}; solo se
        \code{expr1} è vera si valuta anche \code{expr2};
  \item in \code{expr1 || expr2}, se \code{expr1} è vera il risultato è vero
        comunque, e \code{expr2} non viene valutata; solo se \code{expr1} è
        falsa si valuta anche \code{expr2}.
\end{itemize}

La differenza si vede quando \code{expr2} ha un effetto, come un
incremento. I due programmi seguenti differiscono solo per l'ordine degli
operandi di \code{\&\&}, eppure stampano risultati diversi:

\begin{affianco}[0.5][t]
\begin{lstlisting}[style=c, name=sc1.c]
#include <stdio.h>

int main()
{
   int a, b;

   a = 0;
   b = 0;
   while (a < 3 && ++b) {
      a++;
   }
   printf("a = %d b = %d\n", a, b);
}
\end{lstlisting}
\begin{console}
$ ./sc1
a = 3 b = 3
\end{console}
\risultato
\begin{lstlisting}[style=c, name=sc2.c]
#include <stdio.h>

int main()
{
   int a, b;

   a = 0;
   b = 0;
   while (++b && a < 3) {
      a++;
   }
   printf("a = %d b = %d\n", a, b);
}
\end{lstlisting}
\begin{console}
$ ./sc2
a = 3 b = 4
\end{console}
\end{affianco}

Seguiamo i controlli della condizione (figura~\ref{fig:cortocircuito}). Nei
primi tre entrambi i programmi valutano tutte e due le parti, e \code{b}
cresce insieme ad \code{a}. Al quarto controllo \code{a} vale 3: in
\code{sc1.c} la prima parte, \code{a < 3}, è già falsa, e \code{++b} non
viene eseguito; in \code{sc2.c} la prima parte è \code{++b}, che viene
eseguita comunque, e \code{b} arriva a 4.

\begin{figure}[H]
  \centering
  \setlength{\tabcolsep}{5pt}
  \renewcommand{\arraystretch}{1.12}
  \begin{tabular}{@{}c c>{\FiraCodeNL}c>{\FiraCodeNL}c c@{}}
    \multicolumn{5}{c}{\code{while (a < 3 \&\& ++b)}} \\
    \toprule
    \textbf{controllo} & \textbf{\code{a}} & \normalfont\textbf{a < 3} & \normalfont\textbf{++b} & \textbf{risultato} \\
    \midrule
    1 & 0 & vero  & b = 1 & vero \\
    2 & 1 & vero  & b = 2 & vero \\
    3 & 2 & vero  & b = 3 & vero \\
    4 & 3 & falso & \textcolor{accent!75!black}{\Lato non valutato} & falso \\
    \bottomrule
  \end{tabular}
  \hspace{1cm}
  \begin{tabular}{@{}c>{\FiraCodeNL}c c>{\FiraCodeNL}c c@{}}
    \multicolumn{5}{c}{\code{while (++b \&\& a < 3)}} \\
    \toprule
    \textbf{controllo} & \normalfont\textbf{++b} & \textbf{\code{a}} & \normalfont\textbf{a < 3} & \textbf{risultato} \\
    \midrule
    1 & b = 1 & 0 & vero  & vero \\
    2 & b = 2 & 1 & vero  & vero \\
    3 & b = 3 & 2 & vero  & vero \\
    4 & \textcolor{accent!75!black}{b = 4} & 3 & falso & falso \\
    \bottomrule
  \end{tabular}
  \caption{I controlli della condizione nei due programmi: al quarto,
           \code{++b} viene eseguito solo nel secondo.}
  \label{fig:cortocircuito}
\end{figure}

Il cortocircuito è anche una protezione utile. In
\code{if (n != 0 \&\& somma / n > 10)} la divisione viene eseguita solo se
\code{n} è diverso da zero: se \code{n} vale 0, la prima parte è falsa e la
seconda, che dividerebbe per zero, non viene neppure calcolata.

\section{\code{if} e \code{if}-\code{else}}

\subsection{La forma generale}

La forma generale dell'\code{if} è

\begin{lstlisting}[style=c]
if (expr)
    statement
statementSucc
\end{lstlisting}

Se \code{expr} è diversa da zero (vera) viene eseguito \code{statement}; se
è zero (falsa), \code{statement} viene saltato. In entrambi i casi poi si
prosegue con l'istruzione successiva, \code{statementSucc}: è la selezione a
una via della figura~\ref{fig:una-via}. \code{expr} può essere qualsiasi
espressione, ma di solito è relazionale, di uguaglianza o logica.

\begin{affianco}[0.5][c]
\begin{lstlisting}[style=c, name=voto.c]
#include <stdio.h>

int main(void)
{
   int grade;

   printf("Grade: ");
   scanf("%d", &grade);
   if (grade >= 90)
      printf("Congratulations!\n");
   printf("Your grade is %d.\n", grade);
   return 0;
}
\end{lstlisting}
\risultato
\begin{console}
$ ./voto
Grade: <<95>>
Congratulations!
Your grade is 95.
$ ./voto
Grade: <<70>>
Your grade is 70.
\end{console}
\end{affianco}

Le congratulazioni dipendono dalla condizione; la riga con il voto viene
stampata sempre, perché è l'istruzione successiva all'\code{if}.

\subsection{Le graffe, anche per un'istruzione sola}

Il corpo di un \code{if} è \textbf{una} istruzione. Per eseguirne più di
una si usa un'istruzione composta, cioè un blocco fra graffe:

\begin{lstlisting}[style=c]
if (c == ' ') {
    ++blank_cnt;
    printf("found another blank\n");
}
\end{lstlisting}

Le slide consigliano di usare le graffe \textbf{sempre}, anche quando
l'istruzione è una sola: \code{if (j < k) \{ min = j; \}} invece di
\code{if (j < k) min = j;}. Il motivo si vede quando, più avanti, si
aggiunge una seconda istruzione e ci si dimentica delle graffe:
l'indentazione inganna chi legge, ma non il compilatore.

\begin{affianco}[0.5][c]
\begin{lstlisting}[style=c, name=graffe.c]
#include <stdio.h>

int main(void)
{
   int j = 5, k = 3, min = k;

   if (j < k)
      min = j;
      printf("min aggiornato\n");
   printf("min = %d\n", min);
   return 0;
}
\end{lstlisting}
\risultato
\begin{console}
$ ./graffe
min aggiornato
min = 3
\end{console}
\end{affianco}

La condizione è falsa (5 non è minore di 3), eppure \code{"min aggiornato"}
viene stampato: solo \code{min = j;} dipende dall'\code{if}, la
\code{printf} è l'istruzione successiva. Per fortuna \code{gcc} con
\code{-Wall} lo nota:

\begin{console}
$ gcc -Wall -o graffe graffe.c
graffe.c: In function (*@‘@*)main(*@’@*):
graffe.c:7:4: warning: this (*@‘@*)if(*@’@*) clause does not guard... [-Wmisleading-indentation]
    7 |    if (j < k)
      |    ^~
graffe.c:9:7: note: ...this statement, but the latter is misleadingly indented as if it were guarded by the (*@‘@*)if(*@’@*)
    9 |       printf("min aggiornato\n");
      |       ^~~~~~
\end{console}

\subsection{\code{if}-\code{else}}

La forma generale dell'\code{if}-\code{else} è

\begin{lstlisting}[style=c]
if (expr)
    statement1
else
    statement2
\end{lstlisting}

Se \code{expr} è diversa da zero viene eseguito \code{statement1}, se è zero
\code{statement2}: è la selezione a due vie della figura~\ref{fig:due-vie}.
Esattamente una delle due istruzioni viene eseguita.

\subsection{L'\code{else} pendente}

Un \code{if} può contenere a sua volta un \code{if} (e lo stesso vale per
l'\code{if}-\code{else}). Nel blocco

\begin{lstlisting}[style=c]
if (a == 1)
    if (b == 2)
        printf("***\n");
\end{lstlisting}

l'istruzione controllata da \code{if (a == 1)} è proprio
\code{if (b == 2) printf("***\textbackslash n");}: le tre stelle vengono
stampate solo se \code{a} vale 1 \textbf{e} \code{b} vale 2. Aggiungiamo un
\code{else}. C'è differenza fra i due blocchi seguenti?

\begin{affianco}[0.5][c]
\begin{lstlisting}[style=c, name=Esempio 1]
if (a == 1)
    if (b == 2)
        printf("***\n");
    else
        printf("###\n");
\end{lstlisting}
\risultato
\begin{lstlisting}[style=c, name=Esempio 2]
if (a == 1)
    if (b == 2)
        printf("***\n");
else
    printf("###\n");
\end{lstlisting}
\end{affianco}

L'indentazione suggerisce due significati diversi: nel primo l'\code{else}
sembra legato a \code{if (b == 2)}, nel secondo a \code{if (a == 1)}. Ma per
il compilatore l'indentazione \textbf{non ha alcun significato}, e la regola
è una sola: \textbf{l'\code{else} si lega all'\code{if} più vicino} che non
ne abbia già uno. I due blocchi sono identici, ed è il secondo a essere
ingannevole (figura~\ref{fig:else-pendente}).

\begin{figure}[H]
  \centering
  \begin{tikzpicture}[font=\FiraCodeNL\small,
      riga/.style={anchor=west, inner sep=1.5pt}]
    \node[riga] (r1) at (0,0) {if (a == 1)};
    \node[riga] (r2) at (0.5,-0.55) {if (b == 2)};
    \node[riga] (r3) at (1.0,-1.1) {printf("***\textbackslash n");};
    \node[riga] (r4) at (0,-1.65) {else};
    \node[riga] (r5) at (0.5,-2.2) {printf("\#\#\#\textbackslash n");};
    \draw[wtarrowacc, rounded corners=4pt] (r4.east) -- (5.2,-1.65) -- (5.2,-0.55)
      node[pos=0.5, right, font=\Lato\scriptsize, text=accent!70!black, align=left]
        {si lega a questo \code{if},\\il più vicino \mok} -- (r2.east);
    \draw[mrif, draw=merr, dashed, rounded corners=4pt] (r4.west) -- (-0.5,-1.65) -- (-0.5,0)
      node[pos=0.5, left, font=\Lato\scriptsize, text=merr, align=right]
        {non a questo, come fa\\credere l'indentazione \merr} -- (r1.west);
  \end{tikzpicture}
  \caption{L'\code{else} dell'Esempio 2 appartiene all'\code{if} interno,
           nonostante l'indentazione.}
  \label{fig:else-pendente}
\end{figure}

Lo verifichiamo eseguendo il blocco con valori diversi di \code{a} e
\code{b} (con \code{-Wall}, \code{gcc} suggerisce di aggiungere le graffe:
\textit{suggest explicit braces to avoid ambiguous ‘else’}):

\begin{affianco}[0.5][c]
\begin{lstlisting}[style=c, name=dangling.c]
#include <stdio.h>

int main(void)
{
   int a, b;

   printf("a e b: ");
   scanf("%d%d", &a, &b);
   if (a == 1)
      if (b == 2)
         printf("***\n");
   else
      printf("###\n");
   return 0;
}
\end{lstlisting}
\risultato
\begin{console}
$ ./dangling
a e b: <<1 2>>
***
$ ./dangling
a e b: <<1 3>>
###
$ ./dangling
a e b: <<2 2>>
$
\end{console}
\end{affianco}

Con \code{a} uguale a 2 non viene stampato nulla: se l'\code{else} fosse
dell'\code{if} esterno, comparirebbe \code{\#\#\#}. Per ottenere davvero
quel significato servono le graffe, che chiudono l'\code{if} interno:

\begin{lstlisting}[style=c]
if (a == 1) {
    if (b == 2)
        printf("***\n");
} else
    printf("###\n");
\end{lstlisting}

Per chiarezza, la formattazione del testo deve sempre rispecchiare il
significato.

\section{Il \code{while}}

\subsection{La forma generale}

La forma generale del costrutto \code{while} è

\begin{affianco}[0.42][c]
\begin{lstlisting}[style=c]
while (expr)
    statement1
statement2
\end{lstlisting}
\risultato
\begin{enumerate}
  \item viene valutata \code{expr};
  \item se è diversa da zero (vera), viene eseguito \code{statement1} e il
        controllo torna all'inizio del ciclo, al punto 1;
  \item se è zero (falsa), il controllo passa a \code{statement2}.
\end{enumerate}
\end{affianco}

In sostanza \code{statement1} viene eseguito ripetutamente finché
\code{expr} non diventa falsa: è il ciclo con la condizione in testa
(figura~\ref{fig:while}, da confrontare con la figura~\ref{fig:testa}). Se
\code{expr} è falsa fin dall'inizio, \code{statement1} non viene eseguito
neppure una volta.

\begin{figure}[H]
  \centering
  \begin{tikzpicture}[font=\Lato\scriptsize]
    \node[fcgiunzione] (j) at (0,0) {};
    \node[fcdec] (t) at (0,-1.0) {\code{expr}?};
    \node[fcproc] (b) at (0,-2.25) {\code{statement1}};
    \node[fcproc] (s) at (0,-3.85) {\code{statement2}};
    \node[fcgiunzione] (u) at (0,-3.1) {};
    \draw[fcfreccia] (0,0.4) -- (j);
    \draw[fcfreccia] (j) -- (t);
    \draw[fcfreccia] (t) -- node[fcsi, right] {vero} (b);
    \draw[fcfreccia, draw=accent!70!black] (b.south) -- (0,-2.7) -- (1.9,-2.7) -- (1.9,0) -- (j.east);
    \draw[fcfreccia] (t.west) -- node[fcno, above] {falso} (-1.9,-1.0) -- (-1.9,-3.1) -- (u.west);
    \draw[fcfreccia] (u) -- (s);
    \node[mnota, anchor=west, text=accent!70!black] at (2.0,-1.35) {si ripete};
  \end{tikzpicture}
  \caption{Il \code{while}: la condizione si controlla prima di ogni giro.}
  \label{fig:while}
\end{figure}

\subsection{Attenti ai cicli infiniti}

Il frammento delle slide

\begin{lstlisting}[style=c]
printf("Input a positive integer: ");
scanf("%d", &n);
while (--n) {
    ...   /* do something */
}
\end{lstlisting}

dovrebbe eseguire il corpo del ciclo finché \code{--n} diventa 0, cioè
$n - 1$ volte. Si assume però che l'utente inserisca un intero positivo.
Che cosa succede se inserisce 0, o un numero negativo? Proviamo, stampando
\code{n} a ogni giro:

\begin{affianco}[0.5][c]
\begin{lstlisting}[style=c, name=meno.c]
#include <stdio.h>

int main(void)
{
   int n;

   printf("Input a positive integer: ");
   scanf("%d", &n);
   while (--n) {
      printf("%d ", n);
   }
   printf("\n");
   return 0;
}
\end{lstlisting}
\risultato
\begin{console}
$ ./meno
Input a positive integer: <<4>>
3 2 1
$ ./meno
Input a positive integer: <<0>>
-1 -2 -3 -4 -5 -6 -7 -8 -9 -10 -11 -12 -13 -14 -15 -16 -17 -18 -19 -20 -21 -22 -23
^C
\end{console}
\end{affianco}

Partendo da 0, \code{--n} vale $-1$, poi $-2$, $-3$\dots: non diventa
\textbf{mai} 0 (o meglio, ci arriverebbe solo dopo più di 4 miliardi di giri
e un overflow, che è un comportamento indefinito: in una prova senza
stampe, il programma compilato normalmente ha fatto 4\,294\,967\,295 giri
prima di fermarsi; compilato con \code{-O2} non si è fermato affatto). Il
ciclo va controllato meglio, in due modi:

\begin{affianco}[0.5][c]
\begin{lstlisting}[style=c]
while (--n > 0) {
    ...   /* do something */
}
\end{lstlisting}
\risultato
\begin{lstlisting}[style=c]
if (n <= 0) {
    errore();
}
while (--n > 0) {
    ...   /* do something */
}
\end{lstlisting}
\end{affianco}

A sinistra la condizione relazionale \code{--n > 0} ferma il ciclo anche se
\code{n} parte da un valore non positivo; a destra, in più, si controlla
l'input e si segnala il problema (\code{errore} sta per una funzione che lo
gestisce). Un programma non deve mai fidarsi di ciò che scrive l'utente.

\subsection{Il corpo vuoto}

A volte il corpo del \code{while} è un'istruzione vuota: tutto il lavoro lo
fa la condizione. Per esempio

\begin{lstlisting}[style=c]
while ((c = getchar()) == ' ');
\end{lstlisting}

salta tutti i caratteri spazio consecutivi: a ogni giro la condizione legge
un carattere, e il ciclo continua finché è uno spazio. Il \code{;} attaccato
alla parentesi però si nota poco, e potrebbe sembrare un errore di
battitura. Per rendere evidente che è intenzionale si suggerisce di metterlo
da solo su una riga:

\begin{affianco}[0.5][c]
\begin{lstlisting}[style=c, name=salta.c]
#include <stdio.h>

int main(void)
{
   int c;

   while ((c = getchar()) == ' ')
      ;
   printf("primo carattere non bianco: %c\n", c);
   return 0;
}
\end{lstlisting}
\risultato
\begin{console}
$ ./salta
<<     ciao>>
primo carattere non bianco: c
\end{console}
\end{affianco}

\subsection{Contare i caratteri}

Il programma seguente legge un testo e conta gli spazi, le cifre, le lettere,
le righe e gli altri caratteri. Il corpo del \code{while} è un'unica
istruzione, una catena di \code{if}-\code{else} (per questo le graffe non
sono necessarie).

\begin{lstlisting}[style=c, name=cnt_char.c]
/* Count blanks, digits, letters, newlines, and others. */

#include <stdio.h>

int main(void)
{
   int   blank_cnt = 0, c, digit_cnt = 0,
         letter_cnt = 0, nl_cnt = 0, other_cnt = 0;

   while ((c = getchar()) != EOF)     /* braces not necessary */
      if (c == ' ')
         ++blank_cnt;
      else if (c >= '0' && c <= '9')
         ++digit_cnt;
      else if (c >= 'a' && c <= 'z' || c >= 'A' && c <= 'Z')
         ++letter_cnt;
      else if (c == '\n')
         ++nl_cnt;
      else
         ++other_cnt;
   printf("%10s%10s%10s%10s%10s%10s\n\n",
      "blanks", "digits", "letters", "lines", "others", "total");
   printf("%10d%10d%10d%10d%10d%10d\n\n",
      blank_cnt, digit_cnt, letter_cnt, nl_cnt, other_cnt,
      blank_cnt + digit_cnt + letter_cnt + nl_cnt + other_cnt);
   return 0;
}
\end{lstlisting}


Eseguendolo sul suo stesso sorgente:

\begin{console}
$ ./cnt_char < cnt_char.c
    blanks    digits   letters     lines    others     total

       208        32       361        27       180       808

\end{console}

\begin{itemize}
  \item La condizione delle lettere,
        \code{c >= 'a' \&\& c <= 'z' || c >= 'A' \&\& c <= 'Z'}, sfrutta la
        precedenza di \code{\&\&} su \code{||}: “minuscola oppure maiuscola”.
        È corretta così, ma \code{gcc -Wall} suggerisce comunque di
        aggiungere le parentesi, per chiarezza.
  \item \code{\%10s} e \code{\%10d} scrivono ogni stringa e ogni numero in 10
        caratteri, allineati a destra: per questo i numeri stanno in colonna
        sotto i titoli.
  \item I numeri dipendono dal file letto: nelle slide il sorgente aveva una
        riga in più, e il risultato era 28 righe e 809 caratteri.
\end{itemize}

\subsection{La catena di \code{if}-\code{else}}

In \code{cnt\_char.c} ogni \code{else} contiene a sua volta un \code{if}: le
domande si fanno una dopo l'altra, e alla prima risposta vera si esegue
l'azione corrispondente e si salta tutto il resto. Si forma così una
selezione a più vie (figura~\ref{fig:elseif}). Per questo la catena si
scrive con gli \code{else if} allineati, invece di indentare ogni volta di
più: è un'unica scelta fra cinque possibilità.

\begin{figure}[H]
  \centering
  \begin{tikzpicture}[font=\Lato\scriptsize]
    \foreach \cond/\azione [count=\k from 0] in {
        {\code{c == ' '}?}/{\code{++blank\_cnt}},
        {è una cifra?}/{\code{++digit\_cnt}},
        {è una lettera?}/{\code{++letter\_cnt}},
        {\code{c == '\textbackslash n'}?}/{\code{++nl\_cnt}}} {
      \node[fcdec] (d\k) at (0,-\k*1.25) {\cond};
      \node[fcproc] (a\k) at (4.4,-\k*1.25) {\azione};
      \draw[fcfreccia] (d\k.east) -- node[fcsi, above] {vero} (a\k.west);
      \draw[fcfreccia] (a\k.east) -- (6.6,-\k*1.25);
    }
    \foreach \k [evaluate=\k as \kk using int(\k+1)] in {0,1,2}
      \draw[fcfreccia] (d\k.south) -- node[fcno, right] {falso} (d\kk.north);
    \node[fcproc] (a4) at (4.4,-5.0) {\code{++other\_cnt}};
    \draw[fcfreccia] (d3.south) |- node[fcno, pos=0.25, right] {falso} (a4.west);
    \draw[fcfreccia] (a4.east) -- (6.6,-5.0);
    \draw[fclinea] (6.6,0) -- (6.6,-5.6);
    \draw[fcfreccia] (6.6,-5.6) -- ++(0,-0.4);
    \node[mnota, anchor=west] at (6.7,-5.85) {prossimo carattere};
    \draw[fcfreccia] (0,0.85) -- (d0.north);
  \end{tikzpicture}
  \caption{La catena di \code{if}-\code{else} di \code{cnt\_char.c}: le
           domande si fanno in ordine, e alla prima risposta vera si
           incrementa un contatore.}
  \label{fig:elseif}
\end{figure}

\section{Il \code{for}}

\subsection{\code{for} e \code{while}}

La forma del \code{for}, a sinistra, è equivalente al \code{while} a
destra:

\begin{affianco}[0.5][c]
\begin{lstlisting}[style=c]
for (expr1; expr2; expr3)
    statementFor
statementSucc
\end{lstlisting}
\risultato
\begin{lstlisting}[style=c]
expr1;
while (expr2) {
    statementFor
    expr3;
}
statementSucc
\end{lstlisting}
\end{affianco}

\code{expr1} è l'\textbf{inizializzazione}, eseguita una volta sola;
\code{expr2} è la \textbf{condizione}, controllata prima di ogni giro;
\code{expr3} è l'\textbf{aggiornamento}, eseguito alla fine di ogni giro
(figura~\ref{fig:for}). Il \code{for} raccoglie in una riga le tre parti di
un ciclo su contatore, che con il \code{while} sarebbero sparse.

\begin{figure}[H]
  \centering
  \begin{tikzpicture}[font=\Lato\scriptsize]
    \node[fcproc] (e1) at (0,0) {\code{expr1}};
    \node[fcgiunzione] (j) at (0,-0.75) {};
    \node[fcdec] (t) at (0,-1.75) {\code{expr2}?};
    \node[fcproc] (b) at (0,-3.0) {\code{statementFor}};
    \node[fcproc] (e3) at (0,-4.0) {\code{expr3}};
    \node[fcgiunzione] (u) at (0,-4.75) {};
    \node[fcproc] (s) at (0,-5.5) {\code{statementSucc}};
    \draw[fcfreccia] (0,0.6) -- (e1);
    \draw[fcfreccia] (e1) -- (j);
    \draw[fcfreccia] (j) -- (t);
    \draw[fcfreccia] (t) -- node[fcsi, right] {vero} (b);
    \draw[fcfreccia] (b) -- (e3);
    \draw[fcfreccia, draw=accent!70!black] (e3.east) -- (2.0,-4.0) -- (2.0,-0.75) -- (j.east);
    \draw[fcfreccia] (t.west) -- node[fcno, above] {falso} (-2.0,-1.75) -- (-2.0,-4.75) -- (u.west);
    \draw[fcfreccia] (u) -- (s);
    \node[mnota, anchor=west] at (2.6,0) {\code{for (i = 1; i <= 3; i++)}:};
    \node[mnota, anchor=west] at (2.6,-0.45) {\code{i = 1} si esegue una volta sola};
    \node[mnota, anchor=west] at (2.6,-1.75) {\code{i <= 3} si controlla prima di ogni giro};
    \node[mnota, anchor=west] at (2.6,-4.4) {\code{i++} si esegue alla fine di ogni giro};
  \end{tikzpicture}
  \caption{Il \code{for}: inizializzazione, condizione e aggiornamento.}
  \label{fig:for}
\end{figure}

\begin{affianco}[0.5][c]
\begin{lstlisting}[style=c, name=for_while.c]
#include <stdio.h>

int main(void)
{
   int i;

   for (i = 1; i <= 3; i++)
      printf("for: i = %d\n", i);

   i = 1;
   while (i <= 3) {
      printf("while: i = %d\n", i);
      i++;
   }
   return 0;
}
\end{lstlisting}
\risultato
\begin{console}
$ ./for_while
for: i = 1
for: i = 2
for: i = 3
while: i = 1
while: i = 2
while: i = 3
\end{console}
\end{affianco}

Ognuna delle tre espressioni si può omettere (i \code{;} restano):
\code{for (;;)} è un ciclo infinito, come \code{while (1)}.

\subsection{Cicli annidati}

Il corpo di un ciclo può contenere un altro ciclo. Il programma seguente
stampa un triangolo di asterischi: il ciclo esterno scorre le righe, e per
ogni riga il ciclo interno stampa tanti asterischi quanto vale \code{riga}.

\begin{affianco}[0.5][c]
\begin{lstlisting}[style=c, name=triangolo.c]
#include <stdio.h>

int main(void)
{
   int riga, col, n = 4;

   for (riga = 1; riga <= n; riga++) {
      for (col = 1; col <= riga; col++)
         printf("*");
      printf("\n");
   }
   return 0;
}
\end{lstlisting}
\risultato
\begin{center}
\renewcommand{\arraystretch}{1.15}
\begin{tabular}{@{}>{\FiraCodeNL}l>{\FiraCodeNL}l>{\FiraCodeNL\bfseries}l@{}}
  \toprule
  \normalfont\textbf{riga} & \normalfont\textbf{col vale} & \normalfont\textbf{stampa} \\
  \midrule
  1 & 1             & * \\
  2 & 1, 2          & ** \\
  3 & 1, 2, 3       & *** \\
  4 & 1, 2, 3, 4    & **** \\
  \bottomrule
\end{tabular}
\end{center}
\vspace{4pt}
\wtnota{Per ogni giro del ciclo esterno, il ciclo interno ricomincia da
\code{col = 1}.}
\end{affianco}

\subsection{Meglio una relazione che un'uguaglianza}

Nelle condizioni dei cicli è spesso preferibile usare un operatore
relazionale (\code{<}, \code{<=}\dots) invece di un'uguaglianza (\code{==},
\code{!=}): il codice è più robusto. Con i \code{float} e i \code{double},
in particolare, un test di uguaglianza può andare oltre la precisione della
macchina. Il programma delle lezioni vorrebbe sommare i valori 0.0, 0.1,
0.2, \dots, 9.8, fermandosi quando \code{x} arriva a 9.9:

\begin{affianco}[0.56][c]
\begin{lstlisting}[style=c, name=test.c]
/* A test that fails. */

#include <stdio.h>

int main(void)
{
   int      cnt = 0;
   double   sum = 0.0, x;

   for (x = 0.0; x != 9.9; x += 0.1) {   /* trouble! */
      sum += x;
      printf("cnt = %5d\n", ++cnt);
   }
   printf("sum = %f\n", sum);
   return 0;
}
\end{lstlisting}
\risultato
\begin{console}
$ ./test
cnt =     1
cnt =     2
cnt =     3
(*@\textit{\dots}@*)
cnt =    99
cnt =   100
cnt =   101
cnt =   102
(*@\textit{\dots}@*)
^C
\end{console}
\end{affianco}

Il ciclo non termina: in un secondo ha stampato quasi 20 milioni di righe.
Il motivo è nel capitolo~\ref{chap:tipi}: 0.1 non è rappresentabile
esattamente, e sommandolo 99 volte si accumula un piccolo errore. \code{x} non
è mai \textbf{esattamente} uguale alla costante 9.9: le passa accanto e va
oltre (figura~\ref{fig:nove-nove}).

\begin{figure}[H]
  \centering
  \begin{tikzpicture}[font=\FiraCodeNL\scriptsize]
    \draw[line width=1pt, neutral!45!black] (0,0) -- (13.8,0);
    \foreach \xx/\som/\val in {1.2/98/9.79999999999998117, 5.0/99/9.89999999999998082, 12.6/100/9.99999999999998046} {
      \fill[primary!70!black] (\xx,0) circle (2.4pt);
      \node[anchor=north, align=center, text=primary!60!black] at (\xx,-0.15) {\val\\{\Lato dopo \som\ somme}};
    }
    \fill[merr] (6.2,0) circle (2.4pt);
    \draw[merr, line width=0.5pt] (6.2,-0.06) -- (7.7,-0.3);
    \node[anchor=north west, align=left, text=merr] at (7.6,-0.25) {9.90000000000000036\\{\Lato la costante 9.9}};
    \draw[wtarrowacc] (5.0,0.12) to[bend left=30] node[above, font=\Lato\scriptsize, text=accent!70!black] {$+0.1$: \code{x} scavalca 9.9} (12.6,0.12);
  \end{tikzpicture}
  \caption{I valori di \code{x} vicino a 9.9 (la scala non è proporzionale):
           dopo 99 somme \code{x} è più piccolo della costante, dopo 100 è più
           grande. \code{x != 9.9} è sempre vero.}
  \label{fig:nove-nove}
\end{figure}

Con la condizione relazionale \code{x < 9.9} il ciclo termina di sicuro. Si
noti però che i giri diventano 100, non 99: dopo 99 somme \code{x} vale
9.89999999999998, che è ancora minore di 9.9. Il modo più sicuro per contare
è usare un \textbf{intero}, e ricavare da esso il reale:
\codel{for (cnt = 0; cnt < 99; cnt++) \{ x = cnt * 0.1; \dots \}}.

\section{Il \code{do}-\code{while}}

\subsection{La condizione in coda}

Il \code{do}-\code{while} è il ciclo con la condizione \textbf{in coda}
(figura~\ref{fig:coda}): prima si esegue il corpo, poi si controlla la
condizione. Il corpo viene quindi eseguito \textbf{almeno una volta}.

\begin{affianco}[0.42][c]
\begin{lstlisting}[style=c]
do
    statement
while (expr);
statementSucc
\end{lstlisting}
\risultato
\begin{enumerate}
  \item viene eseguito \code{statement};
  \item viene valutata \code{expr}: se è vera si torna al punto 1;
  \item se è falsa si prosegue con \code{statementSucc}.
\end{enumerate}
Il \code{;} dopo la parentesi della condizione è obbligatorio.
\end{affianco}

È la struttura naturale per l'\textbf{input controllato} visto nel
capitolo~\ref{chap:strutturata}: si chiede un valore, e lo si richiede
finché non è accettabile. Chiedere ha senso solo dopo aver chiesto almeno una
volta.

\begin{affianco}[0.5][c]
\begin{lstlisting}[style=c, name=positivo.c]
#include <stdio.h>

int main(void)
{
   int n;

   do {
      printf("Inserisci un intero positivo: ");
      scanf("%d", &n);
   } while (n <= 0);
   printf("Hai inserito %d\n", n);
   return 0;
}
\end{lstlisting}
\risultato
\begin{console}
$ ./positivo
Inserisci un intero positivo: <<-3>>
Inserisci un intero positivo: <<0>>
Inserisci un intero positivo: <<5>>
Hai inserito 5
\end{console}
\end{affianco}

Con un \code{while} la domanda andrebbe scritta due volte, una prima del
ciclo e una dentro; il \code{do}-\code{while} la scrive una volta sola.

\section{Lo \code{switch}}

\subsection{La selezione a \textit{n} vie}

Quando una sola domanda decide fra molte strade si può usare una catena di
\code{if}-\code{else}, come in \code{cnt\_char.c}; ma se la domanda è
“quale valore ha questa espressione intera?”, il C offre un costrutto
apposito, lo \code{switch} (la selezione a $n$ vie della
figura~\ref{fig:n-vie}). Ecco una piccola calcolatrice:

\begin{affianco}[0.5][c]
\begin{lstlisting}[style=c, name=calcolatrice.c]
#include <stdio.h>

int main(void)
{
   double a, b;
   char op;

   printf("Operazione (es. 3 * 4): ");
   scanf("%lf %c %lf", &a, &op, &b);
   switch (op) {
   case '+':
      printf("%g\n", a + b);
      break;
   case '-':
      printf("%g\n", a - b);
      break;
   case '*':
      printf("%g\n", a * b);
      break;
   case '/':
      if (b != 0)
         printf("%g\n", a / b);
      else
         printf("divisione per zero\n");
      break;
   default:
      printf("operatore sconosciuto: %c\n", op);
   }
   return 0;
}
\end{lstlisting}
\risultato
\begin{console}
$ ./calcolatrice
Operazione (es. 3 * 4): <<3 * 4>>
12
$ ./calcolatrice
Operazione (es. 3 * 4): <<7 / 2>>
3.5
$ ./calcolatrice
Operazione (es. 3 * 4): <<5 / 0>>
divisione per zero
$ ./calcolatrice
Operazione (es. 3 * 4): <<2 ^ 8>>
operatore sconosciuto: ^
\end{console}
\vspace{2pt}
\wtnota{La conversione \code{\%g} sceglie da sola fra la notazione normale e
quella scientifica, e non scrive gli zeri inutili.}
\end{affianco}

\subsection{Come funziona}

\begin{itemize}
  \item Si calcola l'espressione fra parentesi, \code{op}, che deve essere
        di tipo intero (un \code{char} va bene).
  \item Si salta all'etichetta \code{case} con quel valore; se non ce n'è
        nessuna, a \code{default} (che è facoltativa: senza, non si esegue
        nulla).
  \item Da lì si eseguono le istruzioni \textbf{in sequenza}, finché si
        incontra un \code{break}, che fa uscire dallo \code{switch}, o si
        arriva alla graffa finale.
  \item Dopo \code{case} va una \textbf{costante} intera: non una
        variabile, non un intervallo, non una stringa.
\end{itemize}

Il terzo punto è il più importante: il \code{break} non è decorativo. Senza,
dopo il \code{case} giusto l'esecuzione “cade” nei successivi
(\textit{fall through}). Togliendo i \code{break} dalla calcolatrice, con
\code{3 + 4} si ottiene:

\begin{affianco}[0.5][c]
\begin{console}
$ ./senza_break
Operazione (es. 3 * 4): <<3 + 4>>
7
-1
12
0.75
operatore sconosciuto: +
\end{console}
\risultato
Dopo il caso \code{'+'} sono state eseguite tutte le istruzioni seguenti,
fino a \code{default} (figura~\ref{fig:switch}). Con \code{-Wextra} il
compilatore avvisa:
\end{affianco}

\begin{console}
$ gcc -Wall -Wextra -o senza_break senza_break.c
senza_break.c: In function (*@‘@*)main(*@’@*):
senza_break.c:12:7: warning: this statement may fall through [-Wimplicit-fallthrough=]
   12 |       printf("%g\n", a + b);
      |       ^~~~~~~~~~~~~~~~~~~~~
senza_break.c:13:4: note: here
   13 |    case '-':
      |    ^~~~
\end{console}

\begin{figure}[H]
  \centering
  \begin{tikzpicture}[font=\Lato\scriptsize,
      az/.style={fcproc, minimum width=2.4cm},
      et/.style={font=\FiraCodeNL\scriptsize, anchor=south, text=accent!70!black, inner sep=1.5pt}]
    \foreach \lab/\testo [count=\k from 0] in {{'+'}/{\code{a + b}}, {'-'}/{\code{a - b}},
                                               {'*'}/{\code{a * b}}, {'/'}/{\code{a / b}},
                                               {default}/{sconosciuto}} {
      \node[az] (s\k) at (2.5,-\k*0.85) {\testo};
      \draw[fcfreccia] (0,-\k*0.85) -- node[et] {\lab} (s\k.west);
      \draw[fcfreccia] (s\k.east) -- node[font=\FiraCodeNL\tiny, above, text=neutral!40!black] {\ifnum\k<4 break\fi} (4.9,-\k*0.85);
      \node[az] (n\k) at (8.9,-\k*0.85) {\testo};
      \draw[fcfreccia] (6.4,-\k*0.85) -- node[et] {\lab} (n\k.west);
    }
    \draw[fclinea] (0,0.4) -- (0,-3.4);
    \node[font=\FiraCodeNL\scriptsize, anchor=south] at (0,0.4) {switch (op)};
    \draw[fclinea] (4.9,0) -- (4.9,-4.0);
    \draw[fcfreccia] (4.9,-4.0) -- ++(0,-0.35) node[below, mnota] {dopo lo \code{switch}};
    \draw[fclinea] (6.4,0.4) -- (6.4,-3.4);
    \node[font=\FiraCodeNL\scriptsize, anchor=south] at (6.4,0.4) {switch (op)};
    \foreach \k [evaluate=\k as \kk using int(\k+1)] in {0,1,2,3}
      \draw[fcfreccia, draw=merr] (n\k.south) -- (n\kk.north);
    \draw[fcfreccia] (n4.south) -- ++(0,-0.6) node[below, mnota] {dopo lo \code{switch}};
    \node[font=\Lato\scriptsize\bfseries] at (2.5,1.05) {con \code{break}};
    \node[font=\Lato\scriptsize\bfseries, text=merr] at (8.3,1.05) {senza \code{break}: si “cade” nei casi seguenti};
  \end{tikzpicture}
  \caption{Lo \code{switch} salta al caso giusto; il \code{break} lo fa uscire,
           altrimenti l'esecuzione prosegue con i casi successivi.}
  \label{fig:switch}
\end{figure}

\subsection{Più casi, un'azione}

La caduta da un caso al successivo è però utile quando più valori devono
fare la stessa cosa: basta scrivere più etichette di seguito. Questo
programma conta le vocali minuscole, gli spazi e gli altri caratteri:

\begin{affianco}[0.5][c]
\begin{lstlisting}[style=c, name=vocali.c]
#include <stdio.h>

int main(void)
{
   int c, vocali = 0, spazi = 0, altri = 0;

   while ((c = getchar()) != EOF) {
      switch (c) {
      case 'a': case 'e': case 'i': case 'o': case 'u':
         vocali++;
         break;
      case ' ':
         spazi++;
         break;
      default:
         altri++;
      }
   }
   printf("vocali: %d  spazi: %d  altri: %d\n",
          vocali, spazi, altri);
   return 0;
}
\end{lstlisting}
\risultato
\begin{console}
$ ./vocali
<<ciao a tutti>>
<<Ctrl+D>>
vocali: 6  spazi: 2  altri: 5
\end{console}
Gli “altri” sono \code{c}, le tre \code{t} e l'a capo finale.
\end{affianco}

\section{\code{break} e \code{continue}}

\subsection{Uscire o saltare un giro}

Il \code{break} non serve solo nello \code{switch}: dentro un ciclo lo fa
terminare subito, e l'esecuzione prosegue dopo il ciclo. Il \code{continue},
invece, salta il resto del giro corrente e passa al controllo della
condizione (in un \code{for}, prima esegue l'aggiornamento \code{expr3}).
Il programma seguente somma i numeri positivi letti, salta i negativi e si
ferma allo 0:

\begin{affianco}[0.5][c]
\begin{lstlisting}[style=c, name=somma_positivi.c]
#include <stdio.h>

int main(void)
{
   int x, somma = 0;

   while (1) {
      printf("numero (0 per finire): ");
      scanf("%d", &x);
      if (x == 0)
         break;          /* esce dal ciclo */
      if (x < 0) {
         printf("  negativo, lo salto\n");
         continue;       /* giro successivo */
      }
      somma += x;
   }
   printf("somma dei positivi: %d\n", somma);
   return 0;
}
\end{lstlisting}
\risultato
\begin{console}
$ ./somma_positivi
numero (0 per finire): <<5>>
numero (0 per finire): <<-2>>
  negativo, lo salto
numero (0 per finire): <<10>>
numero (0 per finire): <<0>>
somma dei positivi: 15
\end{console}
\end{affianco}

La figura~\ref{fig:break-continue} mostra dove porta ciascuna delle due
istruzioni. Se i cicli sono annidati, \code{break} e \code{continue}
agiscono solo sul ciclo più interno che li contiene.

\begin{figure}[H]
  \centering
  \begin{tikzpicture}[font=\FiraCodeNL\small,
      riga/.style={anchor=west, inner sep=1.5pt}]
    \node[riga] (r1) at (0,0) {while (1) \{};
    \node[riga] (r2) at (0.6,-0.5) {scanf("\%d", \&x);};
    \node[riga] (r3) at (0.6,-1.0) {if (x == 0)};
    \node[riga, text=merr] (r4) at (1.2,-1.5) {break;};
    \node[riga] (r5) at (0.6,-2.0) {if (x < 0)};
    \node[riga, text=accent!75!black] (r6) at (1.2,-2.5) {continue;};
    \node[riga] (r7) at (0.6,-3.0) {somma += x;};
    \node[riga] (r8) at (0,-3.5) {\}};
    \node[riga] (r9) at (0,-4.0) {printf("somma ...");};
    \draw[mrif, draw=merr, rounded corners=4pt] (r4.east) -- (4.6,-1.5) -- (4.6,-4.0)
      node[pos=0.5, right, font=\Lato\scriptsize, text=merr] {\code{break}: esce dal ciclo} -- (r9.east);
    \draw[mrif, draw=accent!75!black, rounded corners=4pt] (r6.west) -- (-0.5,-2.5) -- (-0.5,0)
      node[pos=0.5, left, font=\Lato\scriptsize, text=accent!75!black, align=right] {\code{continue}: torna al\\controllo del ciclo} -- (r1.west);
  \end{tikzpicture}
  \caption{Dove saltano \code{break} e \code{continue}.}
  \label{fig:break-continue}
\end{figure}

\code{break} e \code{continue} vanno usati con misura: un ciclo con molte
uscite è difficile da seguire. Usati bene, però, evitano condizioni
complicate e variabili di comodo.

\insintesi{
  In C una condizione è un numero: \textbf{zero} è falso, qualunque altro
  valore è vero; gli operatori relazionali (\code{<} \code{>} \code{<=}
  \code{>=}), di uguaglianza (\code{==} \code{!=}) e logici (\code{!}
  \code{\&\&} \code{||}) restituiscono 1 o 0. \code{a < b < c} non ha il
  significato matematico; \code{=} non è \code{==}; \code{\&\&} e \code{||}
  si fermano appena il risultato è noto (\textbf{cortocircuito}). Il corpo di
  \code{if}, \code{while} e \code{for} è una sola istruzione: per più
  istruzioni servono le graffe, e conviene metterle sempre; l'\code{else} si
  lega all'\code{if} più vicino. \code{while} controlla la condizione prima
  di ogni giro, \code{do}-\code{while} dopo (almeno un giro); \code{for}
  equivale a un \code{while} con inizializzazione e aggiornamento. Nei cicli
  sono più robuste le condizioni relazionali, mai \code{==} o \code{!=} fra
  reali. Lo \code{switch} sceglie fra costanti intere e senza \code{break}
  prosegue nei casi seguenti; \code{break} esce da un ciclo, \code{continue}
  salta al giro successivo.}

\section{Riferimenti}

\begin{itemize}
  \item \textbf{A. Kelley, I. Pohl}, \textit{A Book on C}: capitolo~4
        (flusso di controllo).
  \item \textbf{Statements} (\code{if}, \code{switch}, \code{while},
        \code{do}-\code{while}, \code{for}, \code{break},
        \code{continue}) --- cppreference.com. \\
        \urlx{https://en.cppreference.com/w/c/language/statements}
  \item \textbf{Logical operators} --- cppreference.com. \\
        \urlx{https://en.cppreference.com/w/c/language/operator_logical}
\end{itemize}


  \chapter{Le funzioni}
\label{chap:funzioni}

Scomporre un problema in sotto-problemi più semplici è quasi sempre il modo
migliore per risolverlo (lo sviluppo top-down del
capitolo~\ref{chap:strutturata}), e per scrivere programmi grandi è vitale.
In C questo approccio si realizza con le \textbf{funzioni}. Nel
capitolo~\ref{chap:primi-programmi} ne abbiamo scritte alcune e abbiamo
incontrato prototipi e passaggio per valore; qui vediamo tutto in dettaglio:
come si definisce una funzione, come restituisce un risultato, che cosa
succede se il compilatore non la conosce, dove “vivono” le sue variabili.

\section{Programmi fatti di funzioni}

Un programma C è costituito da uno o più file (capitolo~\ref{chap:make}).
Ogni file contiene zero funzioni (come gli header) o più funzioni, e uno
solo contiene la funzione \code{main}, da cui comincia l'esecuzione;
\code{main} può chiamare altre funzioni, che a loro volta ne chiamano altre,
e tutte operano sulle variabili del programma.

Perché scrivere i programmi usando le funzioni?

\begin{itemize}
  \item è più facile scrivere \textbf{corrette} delle funzioni brevi;
  \item scrittura e ricerca degli errori (\textit{debugging}) sono più
        semplici, perché ogni funzione si prova da sola;
  \item le modifiche sono più semplici: si riscrivono solo le funzioni da
        cambiare;
  \item il codice è più comprensibile: funzioni brevi, con un nome che dice
        che cosa fanno, si leggono e si capiscono più facilmente.
\end{itemize}

\section{La definizione di una funzione}

\subsection{Intestazione e corpo}

Il codice che descrive che cosa fa una funzione è la sua
\term{definizione}. La forma generale è

\begin{lstlisting}[style=c]
type function_name(parameters)
{
    declarations
    statements
}
\end{lstlisting}

La parte prima della graffa aperta è l'\term{intestazione}
(\textit{header}); la parte fra le graffe è il \term{corpo} della funzione.
I \textbf{parametri} sono una lista di dichiarazioni di variabili separate da
virgole: sono detti anche \textbf{parametri formali}, perché ricevono i valori
(i \textit{parametri attuali}, o argomenti) solo al momento della chiamata.
La figura~\ref{fig:anatomia-funzione} mostra le parti nella funzione
fattoriale delle slide.

\begin{figure}[H]
  \centering
  \begin{tikzpicture}[font=\FiraCodeNL\small,
      pz/.style={anchor=base west, inner sep=0pt},
      lab/.style={font=\Lato\scriptsize, text=accent!70!black, align=center, inner sep=1pt}]
    \node[pz] (t) at (0,0) {int};
    \node[pz] (n) at ([xshift=5.4pt]t.base east) {factorial};
    \node[pz] (p) at (n.base east) {(int n)};
    \node[pz] (l2) at (0,-0.55) {\{};
    \node[pz] (l3) at (0,-1.1) {\ \ \ int i, product = 1;};
    \node[pz] (l4) at (0,-1.65) {\ \ \ for (i = 2; i <= n; i++)};
    \node[pz] (l5) at (0,-2.2) {\ \ \ \ \ \ product *= i;};
    \node[pz] (l6) at (0,-2.75) {\ \ \ return product;};
    \node[pz] (l7) at (0,-3.3) {\}};
    \draw[wtbrace] ([yshift=3pt]t.north east) -- ([yshift=3pt]t.north west);
    \draw[wtbrace] ([yshift=3pt]n.north east) -- ([yshift=3pt]n.north west);
    \draw[wtbrace] ([yshift=3pt]p.north east) -- ([yshift=3pt]p.north west);
    \node[lab, anchor=south] at ([yshift=0.32cm]t.north) {tipo del\\risultato};
    \node[lab, anchor=south] at ([yshift=0.32cm]n.north) {nome};
    \node[lab, anchor=south] at ([yshift=0.32cm]p.north) {parametro\\formale};
    \draw[wtbrace] (6.0,-0.12) -- (6.0,0.3);
    \node[lab, anchor=west] at (6.25,0.09) {intestazione (\textit{header})};
    \draw[wtbrace] (6.0,-3.4) -- (6.0,-0.4);
    \node[lab, anchor=west] at (6.25,-1.9) {corpo: dichiarazioni e istruzioni};
    \draw[wtarrowacc] (8.9,-2.75) -- ([xshift=4pt]l6.east);
    \node[lab, anchor=west] at (9.0,-2.75) {restituisce il risultato\\a chi ha chiamato};
  \end{tikzpicture}
  \caption{Le parti della definizione di una funzione.}
  \label{fig:anatomia-funzione}
\end{figure}

\code{factorial} riceve un \code{int} e restituisce un \code{int}: il
prodotto dei numeri da 2 a \code{n}, cioè $n!$ (per \code{n} uguale a 0 o 1
il ciclo non fa nessun giro, e il risultato è 1). Il valore restituito, se
necessario, viene convertito nel tipo della funzione, qui \code{int}; e
grazie alla lista dei parametri anche l'argomento, se necessario, viene
convertito in \code{int}.

\begin{affianco}[0.5][c]
\begin{lstlisting}[style=c, name=fattoriale.c]
#include <stdio.h>

int factorial(int n)              /* header */
{                                 /* body starts here */
   int i, product = 1;

   for (i = 2; i <= n; i++)
      product *= i;
   return product;
}

int main(void)
{
   int k;

   for (k = 0; k <= 5; k++)
      printf("%d! = %d\n", k, factorial(k));
   return 0;
}
\end{lstlisting}
\risultato
\begin{console}
$ ./fattoriale
0! = 1
1! = 1
2! = 2
3! = 6
4! = 24
5! = 120
\end{console}
\end{affianco}

Il corpo di una funzione è un blocco, e può contenere dichiarazioni di
variabili: \code{i} e \code{product} sono visibili \textbf{solo} dentro
\code{factorial}.

\subsection{\code{void}: niente risultato, niente parametri}

Una funzione può non restituire nulla, e può non avere parametri. In
entrambi i casi si usa la parola chiave \code{void} (“vuoto”):

\begin{affianco}[0.6][c]
\codicepiccolo
\begin{lstlisting}[style=c, name=santa.c]
#include <stdio.h>

void wrt_address(void)
{
  printf(" \n%s \n%s \n%s \n%s \n\n",
    "            ************************",
    "            **    SANTA CLAUS     **",
    "            **    NORTH POLE      **",
    "            **    EARTH           **");
}

int main(void)
{
  wrt_address();
  return 0;
}
\end{lstlisting}
\risultato
\begin{console}
$ ./santa

            ************************
            **    SANTA CLAUS     **
            **    NORTH POLE      **
            **    EARTH           **

\end{console}
\end{affianco}

Il primo \code{void} dice al compilatore che la funzione non restituisce
alcun valore; il secondo, fra le parentesi, che non ha parametri. La funzione
si chiama con l'istruzione \code{wrt\_address();}, con le parentesi vuote.
E se si scrivesse \code{a = wrt\_address();}? Non c'è nessun valore da
assegnare, e il compilatore si rifiuta:

\begin{console}
$ gcc -Wall -o santa2 santa2.c
santa2.c: In function (*@‘@*)main(*@’@*):
santa2.c:12:5: error: void value not ignored as it ought to be
   12 |   a = wrt_address();
      |     ^
\end{console}

\warning{Una svista nelle slide}{
  In una delle slide la funzione è scritta con \code{Printf}, con la P
  maiuscola. Il C distingue maiuscole e minuscole, e \code{Printf} non è
  \code{printf}: il compilatore lo tratta come una funzione sconosciuta
  (e suggerisce \textit{did you mean ‘printf’?}), poi il linker non la trova
  da nessuna parte e la compilazione fallisce con \textit{undefined reference
  to ‘Printf’}.}

\subsection{Altri esempi}

Ecco tre definizioni delle slide: una funzione che non fa nulla, una che
raddoppia un reale e una che somma tre interi.

\begin{affianco}[0.5][c]
\begin{lstlisting}[style=c, name=esempi.c]
#include <stdio.h>

void nothing(void)
{
}

double twice(double x)
{
   return (2.0 * x);
}

int all_add(int a, int b, int c)
{
   return (a + b + c);
}

int main(void)
{
   nothing();
   printf("%f\n", twice(3.5));
   printf("%d\n", all_add(1, 2, 3));
   return 0;
}
\end{lstlisting}
\risultato
\begin{console}
$ ./esempi
7.000000
6
\end{console}
\vspace{2pt}
\wtnota{\code{nothing} ha il corpo vuoto: chiamarla non fa nulla, ma è una
funzione lecita (può servire come segnaposto, in attesa di scriverla).}
\end{affianco}

Se il tipo della funzione non è indicato, il C tradizionale lo considera
\code{int}: \code{all\_add(int a, int b, int c) \{ \dots \}} equivale a
\code{int all\_add(\dots)}. Dal C99 questa regola non vale più, e
\code{gcc} protesta (\textit{return type defaults to ‘int’}); i compilatori
più recenti la considerano un errore. Specificare sempre il tipo della
funzione è comunque una buona pratica.

\subsection{Le definizioni nel C tradizionale}

Nel C tradizionale (quello precedente all'ANSI C) la definizione aveva una
sintassi diversa: fra le parentesi c'erano solo i nomi dei parametri, e le
loro dichiarazioni venivano dopo, prima della graffa aperta. L'ordine delle
dichiarazioni non contava, e senza parametri si mettevano solo le parentesi:

\begin{affianco}[0.5][c]
\begin{lstlisting}[style=c]
void f(a, b, c, x, y)
int     a, b, c;
double  x, y;
{
    ...
}
\end{lstlisting}
\risultato
\begin{lstlisting}[style=c]
void f()
{
    ...
}
\end{lstlisting}
\end{affianco}

Questa forma è obsoleta, e fra poco vedremo perché è pericolosa. Si
incontra ancora nei programmi vecchi.

\section{L'istruzione \code{return}}

\subsection{Terminare e restituire}

Con l'istruzione \code{return} la funzione termina e il controllo torna
all'ambiente da cui è stata chiamata. \code{return} non deve
necessariamente contenere un'espressione: \code{return;},
\code{return ++a;} e \code{return (a * b);} sono tutte istruzioni lecite (le
parentesi attorno all'espressione sono facoltative).

\begin{itemize}
  \item Se c'è un'espressione, il suo valore viene passato all'ambiente
        chiamante, convertito se necessario nel tipo della funzione.
  \item In una funzione ci possono essere zero o più \code{return}.
  \item Se non ce n'è nessuno, il controllo torna al chiamante quando si
        raggiunge la graffa di chiusura.
\end{itemize}

\begin{affianco}[0.5][c]
\begin{lstlisting}[style=c, name=assoluto.c]
#include <stdio.h>

double absolute_value(double x)
{
  if (x >= 0)
    return x;
  else
    return -x;
}

int main(void)
{
  printf("%f %f\n", absolute_value(-2.5),
                    absolute_value(4.0));
  return 0;
}
\end{lstlisting}
\risultato
\begin{console}
$ ./assoluto
2.500000 4.000000
\end{console}
\vspace{2pt}
\wtnota{Due \code{return}: a seconda del segno di \code{x}, la funzione
termina con l'uno o con l'altro.}
\end{affianco}

\subsection{Il valore si può ignorare}

Non sempre si usa il valore restituito. Nel frammento delle slide

\begin{lstlisting}[style=c]
while (...) {
    getchar();        /* il carattere letto viene scartato */
    c = getchar();
    ...
}
\end{lstlisting}

il primo \code{getchar()} serve solo a saltare un carattere. Lo facciamo
anche con \code{printf}, senza saperlo: anche \code{printf} restituisce un
valore, il numero di caratteri scritti, che di solito si ignora.

\begin{affianco}[0.5][c]
\begin{lstlisting}[style=c, name=ritorno.c]
#include <stdio.h>

int main(void)
{
   int n;

   n = printf("ciao\n");
   printf("caratteri scritti: %d\n", n);
   return 0;
}
\end{lstlisting}
\risultato
\begin{console}
$ ./ritorno
ciao
caratteri scritti: 5
\end{console}
\wtnota{I caratteri sono 5: \code{c}, \code{i}, \code{a}, \code{o} e l'a
capo.}
\end{affianco}

\section{I prototipi}

\subsection{Dichiarare prima di usare}

Le funzioni devono essere \textbf{dichiarate prima di essere usate}. L'ANSI C
fornisce per questo il \term{prototipo}, che informa il compilatore sul
numero e sul tipo dei parametri e sul tipo del valore restituito. Il
prototipo

\begin{lstlisting}[style=c]
double sqrt(double);
\end{lstlisting}

dice che \code{sqrt} è una funzione, che prende un solo argomento di tipo
\code{double} e che restituisce un valore di tipo \code{double}. La forma
generale è

\begin{lstlisting}[style=c]
type function_name(parameter type list);
\end{lstlisting}

Gli identificatori dei parametri sono facoltativi: \code{void f(char c, int
i);} e \code{void f(char, int);} sono equivalenti. Il compilatore non li
usa, ma servono da documentazione, soprattutto se i nomi sono significativi:
conviene metterli sempre.

Nel C tradizionale la dichiarazione non conteneva la lista dei parametri:
\code{double sqrt();}. I prototipi sono preferibili perché, conoscendo i
tipi dei parametri, il compilatore \textbf{converte} opportunamente i valori
passati alla funzione.

\subsection{Che cosa succede senza prototipo}

Consideriamo la chiamata \code{sqrt(a)}, con \code{a} intero e uguale a 4.
Con il prototipo il compilatore sa che \code{sqrt} vuole un \code{double}, e
converte 4 in 4.0; senza, non sa nulla, e il risultato non è corretto:

\begin{affianco}[0.5][t]
\begin{lstlisting}[style=c, name=prova.c]
#include <stdio.h>

double sqrt();

int main()
{
  int a = 4;

  printf("sqrt(4) = %f \n", sqrt(a));
}
\end{lstlisting}
\begin{console}
$ gcc -o prova prova.c -lm
$ ./prova
sqrt(4) = 0.000000
\end{console}
\risultato
\begin{lstlisting}[style=c, name=prova.c]
#include <stdio.h>

double sqrt(double);

int main()
{
  int a = 4;

  printf("sqrt(4) = %f \n", sqrt(a));
}
\end{lstlisting}
\begin{console}
$ gcc -o prova prova.c -lm
$ ./prova
sqrt(4) = 2.000000
\end{console}
\end{affianco}

Perché zero? Sui PC a 64 bit gli argomenti interi e quelli reali viaggiano in
due registri diversi del processore. Senza prototipo il valore 4 viene
consegnato come intero, nel registro degli interi; \code{sqrt} va a prendere
il suo \code{double} nell'altro registro, e ci trova un valore che non
c'entra (qui, per caso, 0). Con il prototipo il compilatore converte 4 in
4.0 e lo consegna dove \code{sqrt} lo cerca (figura~\ref{fig:registri}).

\begin{figure}[H]
  \centering
  \begin{tikzpicture}[font=\Lato\scriptsize,
      reg/.style={draw=primary!60!black, fill=primary!5!white, line width=0.7pt,
        minimum width=2.4cm, minimum height=0.55cm, font=\FiraCodeNL\scriptsize},
      vuoto/.style={reg, fill=white, draw=neutral!40!white},
      chi/.style={wtnode, font=\Lato\scriptsize, minimum width=1.9cm}]
    \node[font=\Lato\scriptsize\bfseries, anchor=west] at (-1.6,1.05) {senza prototipo};
    \node[chi] (m1) at (-0.6,0) {\code{main}};
    \node[reg] (ri1) at (3.6,0.35) {4};
    \node[reg] (rr1) at (3.6,-0.45) {0.0};
    \node[mnota, anchor=west] at ([xshift=3pt]ri1.east) {registro degli interi};
    \node[mnota, anchor=west] at ([xshift=3pt]rr1.east) {registro dei reali};
    \node[chi] (s1) at (9.4,0) {\code{sqrt}};
    \draw[wtarrow] (m1.east) -- node[above, mnota, pos=0.45] {consegna 4} (ri1.west);
    \draw[wtarrow, draw=merr] (rr1.south) |- node[pos=0.75, below, mnota, text=merr]
      {\code{sqrt} cerca qui il suo \code{double} e trova 0.0} (9.4,-1.0) -- (s1.south);
    \begin{scope}[yshift=-2.9cm]
      \node[font=\Lato\scriptsize\bfseries, anchor=west] at (-1.6,1.05) {con prototipo};
      \node[chi] (m2) at (-0.6,0) {\code{main}};
      \node[vuoto] (ri2) at (3.6,0.35) {};
      \node[reg] (rr2) at (3.6,-0.45) {4.0};
      \node[mnota, anchor=west] at ([xshift=3pt]ri2.east) {registro degli interi};
      \node[mnota, anchor=west] at ([xshift=3pt]rr2.east) {registro dei reali};
      \node[chi] (s2) at (9.4,0) {\code{sqrt}};
      \draw[wtarrow] (m2.east) -- node[below left, mnota, pos=0.55] {converte e consegna 4.0} (rr2.west);
      \draw[wtarrow, draw=mok] (rr2.south) |- node[pos=0.75, below, mnota, text=mok]
        {\code{sqrt} trova 4.0} (9.4,-1.0) -- (s2.south);
    \end{scope}
  \end{tikzpicture}
  \caption{Senza prototipo l'argomento viene consegnato nel posto sbagliato.}
  \label{fig:registri}
\end{figure}

\code{sqrt}, in realtà, è una funzione che \code{gcc} conosce bene: anche
senza \code{-Wall} avvisa che l'argomento è un \code{int} dove ci si
aspetta un \code{double}. Con \code{clang}, che la conosce altrettanto bene,
il primo programma stampa addirittura 2.000000. Con una funzione qualsiasi,
però, nessun compilatore può rimediare: senza prototipo il risultato non è
garantito.

\subsection{Come il compilatore scopre le funzioni}

Per il compilatore una funzione può essere dichiarata in tre modi: da una
\textbf{chiamata}, da una \textbf{definizione} o da un \textbf{prototipo}.

\begin{itemize}
  \item Se incontra una chiamata \code{f(x)} prima di qualunque altra
        dichiarazione, il compilatore assume la dichiarazione
        \code{int f();}: un risultato \code{int} e nessuna ipotesi sui
        parametri, che il programmatore deve passare del tipo giusto (dal
        C99 è un errore, che \code{gcc} segnala con un avviso, come abbiamo
        visto nel capitolo~\ref{chap:primi-programmi}).
  \item Una definizione in stile tradizionale dichiara solo il nome e il tipo
        del risultato, non i parametri.
  \item Una definizione in stile ANSI, come un prototipo, dichiara tutto.
\end{itemize}

Le slide lo mostrano con tre versioni della stessa funzione, che calcola la
radice quadrata del suo parametro, chiamata con un \code{int}, un
\code{double} e un \code{long double} che valgono tutti 4. Ecco la prima,
con la definizione tradizionale:

\begin{affianco}[0.5][c]
\begin{lstlisting}[style=c, name=kr1.c]
#include <stdio.h>
#include <math.h>

float f(x)
long double x;
{
  return sqrt(x);
}

int main()
{
  int a = 4;
  double b = 4;
  long double c = 4;

  printf("parametro intero: %f \n", f(a));
  printf("parametro double: %f \n", f(b));
  printf("parametro long double: %f \n", f(c));
}
\end{lstlisting}
\risultato
\begin{console}
$ ./kr1
parametro intero: 0.000000
parametro double: 0.000000
parametro long double: 2.000000
\end{console}
\end{affianco}

Il codice contiene la definizione di \code{f}, ma solo la riga
\code{float f(x)} vale come dichiarazione: non si assume nulla sui
parametri, che vanno passati del tipo giusto, perché il compilatore non può
fare nessuna conversione. Così solo la terza chiamata, che passa davvero un
\code{long double}, dà il risultato corretto.

\subsection{Un parametro senza tipo}

Nella seconda versione il tipo di \code{x} manca del tutto. Il compilatore
assume la dichiarazione \code{int x;} (e avvisa: \textit{type of ‘x’
defaults to ‘int’}); ora funziona solo la chiamata con l'intero:

\begin{affianco}[0.5][c]
\begin{lstlisting}[style=c, name=kr2.c]
#include <stdio.h>
#include <math.h>

float f(x)
{
  return sqrt(x);
}
\end{lstlisting}
\risultato
\begin{console}
$ ./kr2
parametro intero: 2.000000
parametro double: -nan
parametro long double: -nan
\end{console}
\wtnota{\code{main} è identico a quello di \code{kr1.c}.}
\end{affianco}

Nelle slide le due chiamate sbagliate stampano 0.000000; sul nostro
computer stampano \code{-nan} (\textit{not a number}: la radice di un
numero negativo). Sono valori presi a caso dalla memoria, e possono cambiare
da un computer all'altro.

\subsection{La definizione ANSI}

Nella terza versione l'intestazione è \code{float f(long double x)}: questa
volta la riga che fa da dichiarazione è un \textbf{prototipo}. Il
compilatore conosce la lista dei parametri, e può convertire l'\code{int} e
il \code{double} in \code{long double}:

\begin{affianco}[0.5][c]
\begin{lstlisting}[style=c, name=ansi.c]
#include <stdio.h>
#include <math.h>

float f(long double x)
{
  return sqrt(x);
}
\end{lstlisting}
\risultato
\begin{console}
$ ./ansi
parametro intero: 2.000000
parametro double: 2.000000
parametro long double: 2.000000
\end{console}
\end{affianco}

\begin{center}
\renewcommand{\arraystretch}{1.2}
\begin{tabular}{@{}l c c c@{}}
  \toprule
  \textbf{definizione di \code{f}} & \textbf{con un \code{int}} &
  \textbf{con un \code{double}} & \textbf{con un \code{long double}} \\
  \midrule
  \code{float f(x) long double x;} & \textcolor{merr}{0.000000} & \textcolor{merr}{0.000000} & \textcolor{mok}{2.000000} \\
  \code{float f(x)}                & \textcolor{mok}{2.000000}  & \textcolor{merr}{-nan}     & \textcolor{merr}{-nan} \\
  \code{float f(long double x)}    & \textcolor{mok}{2.000000}  & \textcolor{mok}{2.000000}  & \textcolor{mok}{2.000000} \\
  \bottomrule
\end{tabular}
\end{center}

In definitiva: la definizione della funzione in stile ANSI, oppure il suo
prototipo, deve comparire \textbf{prima} che la funzione sia usata. È anche
uno dei motivi per cui si includono gli header: contengono i prototipi delle
funzioni di libreria.

\section{Il passaggio per valore}

\subsection{Una copia per ogni argomento}

Gli argomenti delle funzioni sono passati \textbf{per valore}: ogni argomento
viene calcolato, e il suo valore viene usato localmente al posto del
parametro corrispondente. Se a una funzione si passa una variabile, il suo
valore nell'ambiente chiamante \textbf{non cambia}, anche se la funzione
modifica il valore che riceve (lo abbiamo visto nella
figura~\ref{fig:per-valore}). L'esempio delle slide:

\begin{affianco}[0.58][c]
\codicepiccolo
\begin{lstlisting}[style=c, name=compute_sum.c]
#include <stdio.h>

int compute_sum(int n)       /* sum the integers from 1 to n */
{
   int   sum = 0;

   for ( ; n > 0; --n)       /* stored value of n is changed */
      sum += n;
   return sum;
}

int main(void)
{
   int   n = 3, sum, compute_sum(int);

   printf("%d\n", n);
   sum = compute_sum(n);
   printf("%d\n", n);
   printf("%d\n", sum);
   return 0;
}
\end{lstlisting}
\risultato
\begin{console}
$ ./compute_sum
3
3
6
\end{console}
\wtnota{\code{compute\_sum(int)}, fra le dichiarazioni di \code{main}, è un
prototipo: lecito anche lì.}
\end{affianco}

\subsection{La pila delle chiamate}

Dove si trovano le variabili di una funzione? Ogni chiamata riceve una
propria zona di memoria (un \textit{record di attivazione}), che contiene i
parametri e le variabili locali; le zone si impilano una sull'altra, come
piatti, e ognuna viene tolta quando la sua funzione termina: è la
\term{pila delle chiamate} (\textit{call stack}). La
figura~\ref{fig:pila} la segue durante l'esecuzione di
\code{compute\_sum.c}.

\begin{figure}[H]
  \centering
  \begin{tikzpicture}[font=\Lato\scriptsize]
    \foreach \col/\testo in {0/{1. prima della chiamata}, 1/{2. inizio di \code{compute\_sum}},
                             2/{3. al \code{return}}, 3/{4. dopo il ritorno}}
      \node[font=\Lato\scriptsize\bfseries] at (\col*4.7+1.4,1.25) {\testo};
    % 1
    \variabile{a1}{(1.6,0.0)}{n}{3}
    \variabile{b1}{(1.6,-0.6)}{sum}{\mvuota}
    \scopo{2}{m1}{main}{(a1-nome) (a1) (b1-nome) (b1)}
    % 2
    \begin{scope}[xshift=4.7cm]
      \variabile{a2}{(1.6,0.0)}{n}{3}
      \variabile{b2}{(1.6,-0.6)}{sum}{\mvuota}
      \scopo{2}{m2}{main}{(a2-nome) (a2) (b2-nome) (b2)}
      \variabile[fill=accent!10!white]{c2}{(1.6,-1.95)}{n}{3}
      \variabile[fill=accent!10!white]{d2}{(1.6,-2.55)}{sum}{0}
      \scopo{3}{f2}{compute\_sum}{(c2-nome) (c2) (d2-nome) (d2)}
      \draw[mcerca] (a2.east) to[out=-10, in=10] node[wtlab, right, fill=none] {copia} (c2.east);
    \end{scope}
    % 3
    \begin{scope}[xshift=9.4cm]
      \variabile{a3}{(1.6,0.0)}{n}{3}
      \variabile{b3}{(1.6,-0.6)}{sum}{\mvuota}
      \scopo{2}{m3}{main}{(a3-nome) (a3) (b3-nome) (b3)}
      \variabile[fill=accent!10!white]{c3}{(1.6,-1.95)}{n}{0}
      \variabile[fill=accent!10!white]{d3}{(1.6,-2.55)}{sum}{6}
      \scopo{3}{f3}{compute\_sum}{(c3-nome) (c3) (d3-nome) (d3)}
      \draw[mcerca] (d3.east) to[out=10, in=-10] node[wtlab, right, fill=none] {6} (b3.east);
    \end{scope}
    % 4
    \begin{scope}[xshift=14.1cm]
      \variabile{a4}{(1.6,0.0)}{n}{3}
      \variabile{b4}{(1.6,-0.6)}{sum}{6}
      \scopo{2}{m4}{main}{(a4-nome) (a4) (b4-nome) (b4)}
      \node[mnota, align=center] at (1.4,-2.25) {la zona di \code{compute\_sum}\\è stata tolta dalla pila};
    \end{scope}
  \end{tikzpicture}
  \caption{La pila delle chiamate: \code{compute\_sum} ha la sua \code{n}
           (una copia) e la sua \code{sum}; modifica la propria \code{n}, e
           quando termina la sua zona scompare.}
  \label{fig:pila}
\end{figure}

\code{compute\_sum} porta la \textbf{sua} \code{n} da 3 a 0 mentre somma; la
\code{n} di \code{main} è un'altra variabile, e resta 3. Se si vuole che una
funzione modifichi una variabile del chiamante, ci sono due strade:
restituire il nuovo valore come risultato (e assegnarlo nel chiamante), oppure
passare l'\textbf{indirizzo} della variabile, cioè un puntatore
(capitolo~\ref{chap:puntatori}).

\section{Visibilità: variabili locali e globali}

\subsection{Locali e globali}

Le variabili dichiarate nel corpo di una funzione sono dette
\term{locali}: esistono e sono visibili solo lì. Le variabili dichiarate
fuori da ogni funzione sono dette \term{globali} (o esterne): sono visibili
da tutte le funzioni del file che vengono dopo la dichiarazione. Che cosa
stampa questo programma?

\begin{affianco}[0.5][c]
\begin{lstlisting}[style=c, name=globale.c]
#include <stdio.h>

int a = 33;

int all_add(int a, int b, int c)
{
  return (a+b+c);
}

int main(void)
{
  int b = 15, c = 34;
  printf("%d \n", all_add(a,b,c));
}
\end{lstlisting}
\risultato
\begin{console}
$ ./globale
82
\end{console}
\vspace{2pt}
\wtnota{In \code{main} non c'è una \code{a} locale: \code{a} è quella globale,
che vale 33, e $33 + 15 + 34 = 82$.}
\end{affianco}

Ora aggiungiamo una \code{a} locale in \code{main}, con valore 7:

\begin{affianco}[0.5][c]
\begin{lstlisting}[style=c, name=locale.c]
int main(void)
{
  int a=7, b = 15, c = 34;
  printf("%d \n", all_add(a,b,c));
}
\end{lstlisting}
\risultato
\begin{console}
$ ./locale
56
\end{console}
\vspace{2pt}
\wtnota{La \code{a} locale \textbf{nasconde} quella globale: $7 + 15 + 34 = 56$.}
\end{affianco}

Dentro \code{main}, il nome \code{a} indica la variabile locale, che
nasconde la globale; dentro \code{all\_add}, \code{a} è il parametro, che a
sua volta nasconde la globale (figura~\ref{fig:visibilita}). In ogni punto
del programma vale la dichiarazione \textbf{più vicina}.

\begin{figure}[H]
  \centering
  \begin{tikzpicture}[font=\Lato\scriptsize]
    % globale.c
    \node[font=\Lato\scriptsize\bfseries] at (2.4,2.4) {globale.c};
    \variabile{g1}{(1.0,1.55)}{a}{33}
    \variabile{mb1}{(1.0,0.2)}{b}{15}
    \variabile{mc1}{(1.0,-0.4)}{c}{34}
    \scopo{2}{mm1}{main}{(mb1-nome) (mb1) (mc1-nome) (mc1)}
    \variabile[fill=accent!10!white]{fa1}{(3.9,0.5)}{a}{33}
    \variabile[fill=accent!10!white]{fb1}{(3.9,-0.1)}{b}{15}
    \variabile[fill=accent!10!white]{fc1}{(3.9,-0.7)}{c}{34}
    \scopo{3}{ff1}{all\_add}{(fa1-nome) (fa1) (fb1-nome) (fb1) (fc1-nome) (fc1)}
    \scopo{1}{file1}{il file: variabili globali}{(g1-nome) (g1) (mm1) (ff1)}
    \draw[mcerca] (g1.east) to[out=0, in=90] (fa1.north);
    \node[mnota, text=mok] at (2.5,-1.75) {risultato: $33 + 15 + 34 = 82$};
    % locale.c
    \begin{scope}[xshift=7.6cm]
      \node[font=\Lato\scriptsize\bfseries] at (2.4,2.4) {locale.c};
      \variabile[draw=neutral!35!white, text=neutral!45!black]{g2}{(1.0,1.55)}{a}{33}
      \variabile{ma2}{(1.0,0.5)}{a}{7}
      \variabile{mb2}{(1.0,-0.1)}{b}{15}
      \variabile{mc2}{(1.0,-0.7)}{c}{34}
      \scopo{2}{mm2}{main}{(ma2-nome) (ma2) (mb2-nome) (mb2) (mc2-nome) (mc2)}
      \variabile[fill=accent!10!white]{fa2}{(3.9,0.5)}{a}{7}
      \variabile[fill=accent!10!white]{fb2}{(3.9,-0.1)}{b}{15}
      \variabile[fill=accent!10!white]{fc2}{(3.9,-0.7)}{c}{34}
      \scopo{3}{ff2}{all\_add}{(fa2-nome) (fa2) (fb2-nome) (fb2) (fc2-nome) (fc2)}
      \scopo{1}{file2}{il file: variabili globali}{(g2-nome) (g2) (mm2) (ff2)}
      \draw[mcerca] (ma2.east) -- (fa2-nome.west);
      \node[mnota, anchor=west, text=neutral!45!black] at ([xshift=4pt]g2.east) {nascosta in \code{main}};
      \node[mnota, text=mok] at (2.5,-1.75) {risultato: $7 + 15 + 34 = 56$};
    \end{scope}
  \end{tikzpicture}
  \caption{Da dove viene la \code{a} passata ad \code{all\_add}: a sinistra
           è la globale, a destra quella locale di \code{main}.}
  \label{fig:visibilita}
\end{figure}

\subsection{I blocchi}

Più in generale, un \textbf{blocco} è una parte di codice delimitata da
graffe, e può contenere dichiarazioni di variabili, visibili solo
all'interno del blocco. Il corpo di una funzione è un blocco, ma un blocco
può stare anche dentro un altro:

\begin{affianco}[0.5][c]
\begin{lstlisting}[style=c, name=blocchi.c]
#include <stdio.h>

int main(void)
{
   int a = 1;

   {
      int a = 2;              /* nasconde la a esterna */
      int b = 3;

      printf("dentro il blocco: a = %d, b = %d\n", a, b);
   }
   printf("fuori dal blocco: a = %d\n", a);
   return 0;
}
\end{lstlisting}
\risultato
\begin{console}
$ ./blocchi
dentro il blocco: a = 2, b = 3
fuori dal blocco: a = 1
\end{console}
\wtnota{Fuori dal blocco \code{b} non esiste più: usarla sarebbe un errore.
L'opzione \code{-Wshadow} di \code{gcc} segnala le variabili che ne
nascondono altre.}
\end{affianco}

\section{Le classi di memorizzazione}

Ogni variabile e ogni funzione del C ha due attributi: il \textbf{tipo} e la
\term{classe di memorizzazione} (\textit{storage class}), che stabilisce
\textbf{dove} e \textbf{per quanto tempo} esiste e \textbf{da dove} è
visibile. Le classi sono quattro, corrispondenti alle parole chiave
\code{auto}, \code{extern}, \code{register} e \code{static}.

\subsection{\code{auto}}

Le variabili dichiarate dentro un blocco sono \code{auto} per default.
Quando si entra nel blocco, il sistema alloca la memoria per le sue variabili
locali; quando se ne esce, la libera (è la pila della
figura~\ref{fig:pila}). Poiché è la classe più comune, la parola chiave
\code{auto} non si scrive quasi mai.

\subsection{\code{extern}}

Un modo semplice, ma sconsigliato, di passare informazioni fra blocchi e
funzioni sono le variabili esterne. Una variabile è esterna se è definita
fuori da ogni blocco, e riceve automaticamente la classe \code{extern}. La
parola chiave \code{extern}, invece, dice al compilatore di cercare la
variabile \textbf{da un'altra parte}. Un esempio su due file:

\begin{affianco}[0.5][t]
\begin{lstlisting}[style=c, name=file1.c]
#include <stdio.h>

int a = 1, b = 2, c = 3;        /* variabili esterne */
int f(void);

int main(void)
{
  printf("%3d\n", f());
  printf("%3d %3d %3d\n", a, b, c);
  return 0;
}
\end{lstlisting}
\risultato
\begin{lstlisting}[style=c, name=file2.c]
extern int a;              /* definita altrove */

int f(void)
{
  int b, c;                  /* variabili locali */

  a = b = c = 4;
  return (a + b + c);
}
\end{lstlisting}
\begin{console}
$ gcc -Wall -c file1.c
$ gcc -Wall -c file2.c
$ gcc -o esterne file1.o file2.o
$ ./esterne
 12
  4   2   3
\end{console}
\end{affianco}

\begin{itemize}
  \item I due file sono compilati separatamente
        (capitolo~\ref{chap:make}). In \code{file1.c}, \code{a}, \code{b} e
        \code{c} sono variabili esterne.
  \item Compilando \code{file2.c}, la riga \code{extern int a;} dice al
        compilatore che \code{a} è globale e definita in un altro file.
  \item Il linker mette insieme le informazioni e fa sì che \code{f} usi la
        cella di memoria giusta (figura~\ref{fig:extern}).
  \item \code{f} assegna 4 alla \code{a} globale e alle sue \code{b} e
        \code{c} locali, e restituisce 12; le \code{b} e \code{c} globali
        restano 2 e 3.
\end{itemize}

\begin{figure}[H]
  \centering
  \begin{tikzpicture}[font=\FiraCodeNL\scriptsize,
      sim/.style={anchor=west, inner sep=1.5pt},
      box/.style={draw=primary!65!black, fill=primary!4!white, line width=0.7pt,
        rounded corners=3pt}]
    \draw[box] (0,0) rectangle (3.2,-2.6);
    \node[font=\Lato\scriptsize\bfseries, anchor=south west] at (0,0.02) {file1.o};
    \node[sim, text=mok] (da) at (0.1,-0.35) {D a};
    \node[sim, text=mok] at (0.1,-0.75) {D b};
    \node[sim, text=mok] at (0.1,-1.15) {D c};
    \node[sim] at (0.1,-1.55) {T main};
    \node[sim, text=merr] (uf) at (0.1,-1.95) {U f};
    \node[sim, text=neutral!45!black] at (0.1,-2.35) {U printf};
    \draw[box] (5.6,-0.6) rectangle (8.4,-1.8);
    \node[font=\Lato\scriptsize\bfseries, anchor=south west] at (5.6,-0.58) {file2.o};
    \node[sim, text=merr] (ua) at (5.7,-0.95) {U a};
    \node[sim, text=mok] (tf) at (5.7,-1.45) {T f};
    \draw[wtarrowacc] (ua.west) -- (da.east);
    \draw[wtarrowacc] (uf.east) -- (tf.west);
    \node[mnota, anchor=north, align=center] at (4.2,-2.75)
      {\code{nm}: \code{D} è una variabile definita nel file, \code{T} una funzione definita,\\
       \code{U} un nome usato ma definito altrove: il linker collega ogni \code{U} alla sua definizione};
  \end{tikzpicture}
  \caption{Che cosa vede il linker: \code{file2.o} usa la \code{a} definita in
           \code{file1.o}, e \code{file1.o} usa la \code{f} definita in
           \code{file2.o}.}
  \label{fig:extern}
\end{figure}

Le informazioni si possono passare anche come parametri delle funzioni, ed
è la strada da preferire: bisogna cercare di evitare le variabili globali.
Con i parametri il codice è più chiaro e si capisce meglio dove i valori
vengono modificati; una variabile globale, invece, può essere cambiata da
qualunque funzione di qualunque file.

Anche le funzioni hanno per default la classe \code{extern}: sono visibili
ovunque nel programma, anche se è scritto su più file. Si può aggiungere
\code{extern} alla definizione o al prototipo
(\code{extern double sin(double);}), ma in generale non serve.

\subsection{\code{register}}

La classe \code{register} chiede al compilatore di tenere la variabile, se
possibile, in un \textbf{registro} del processore, la memoria più veloce
(capitolo~\ref{chap:calcolatore}). I registri sono pochi e servono anche al
sistema, quindi la richiesta non sempre è soddisfatta: in quel caso la
variabile diventa \code{auto}. Le variabili globali e le funzioni non
possono essere \code{register}.

\code{register} è un tentativo di velocizzare il programma, e le variabili
vanno scelte con cura: poche, usate molto spesso. I candidati tipici sono
le variabili di controllo dei cicli. Dichiarata all'inizio di \code{main},
però, la variabile occupa il registro per tutta l'esecuzione della
funzione; è più efficiente racchiudere dichiarazione e ciclo in un blocco:

\begin{affianco}[0.5][c]
\begin{lstlisting}[style=c]
int main(void)
{
    register int i;
    ...
    for (i = 0; i < LIMIT; i++) {
        ...
    }
}
\end{lstlisting}
\risultato
\begin{lstlisting}[style=c]
int main(void)
{
    ...
    {
        register int i;
        for (i = 0; i < LIMIT; i++) {
            ...
        }
    }
}
\end{lstlisting}
\end{affianco}

Una variabile \code{register} non ha un indirizzo in memoria, quindi non le
si può applicare \code{\&}: \textit{address of register variable ‘i’
requested}, risponde \code{gcc}. Oggi i compilatori scelgono da soli, e
molto bene, quali variabili tenere nei registri: la parola chiave
\code{register} è di fatto superata.

\subsection{\code{static} per le variabili locali}

Le dichiarazioni \code{static} hanno due usi distinti. Il primo è
\textbf{conservare il valore} di una variabile locale anche dopo l'uscita
dal blocco: una variabile locale \code{static} viene inizializzata una sola
volta, e non sta nella pila ma in una zona di memoria che esiste per tutta
l'esecuzione del programma.

\begin{affianco}[0.5][c]
\begin{lstlisting}[style=c, name=statica.c]
#include <stdio.h>

void f(void)
{
   static int cnt = 0;
   int volte = 0;

   ++cnt;
   ++volte;
   printf("cnt = %d   volte = %d\n", cnt, volte);
}

int main(void)
{
   f();
   f();
   f();
   return 0;
}
\end{lstlisting}
\risultato
\begin{console}
$ ./statica
cnt = 1   volte = 1
cnt = 2   volte = 1
cnt = 3   volte = 1
\end{console}
\end{affianco}

La prima volta \code{cnt} viene inizializzata a 0 e incrementata a 1. La
seconda volta che \code{f} viene chiamata, \code{cnt} \textbf{non} viene
inizializzata di nuovo: conserva il valore 1 e lo porta a 2. \code{volte},
che è \code{auto}, nasce e muore a ogni chiamata, e vale sempre 1
(figura~\ref{fig:statica}). Le slide usano questa idea per fare una cosa sì
e una no: \code{if (cnt \% 2 == 0)} distingue le chiamate pari da quelle
dispari.

\begin{figure}[H]
  \centering
  \begin{tikzpicture}[font=\Lato\scriptsize]
    \foreach \k in {1,2,3} {
      \node[font=\Lato\scriptsize\bfseries] at (\k*3.6-1.0,1.3) {chiamata \k};
      \variabile[fill=accent!10!white]{v\k}{(\k*3.6-0.8,-1.3)}{volte}{1}
      \scopo{3}{fr\k}{f}{(v\k-nome) (v\k)}
    }
    \variabile{c1}{(2.8,0.35)}{cnt}{1}
    \variabile{c2}{(6.4,0.35)}{cnt}{2}
    \variabile{c3}{(10.0,0.35)}{cnt}{3}
    \draw[wtarrowacc] (c1.east) -- node[above, mnota, text=accent!70!black] {$+1$} (c2-nome.west);
    \draw[wtarrowacc] (c2.east) -- node[above, mnota, text=accent!70!black] {$+1$} (c3-nome.west);
    \node[mnota, anchor=west, align=left] at (11.2,0.35) {\code{static}: la stessa variabile,\\fuori dalla pila, conserva il valore};
    \node[mnota, anchor=west, align=left] at (11.2,-1.3) {\code{auto}: una variabile nuova\\a ogni chiamata, che riparte da 0};
  \end{tikzpicture}
  \caption{\code{cnt} attraversa le chiamate; \code{volte} viene creata e
           distrutta ogni volta.}
  \label{fig:statica}
\end{figure}

\subsection{\code{static} per le variabili globali}

Il secondo uso riguarda le variabili esterne. Una variabile globale
\code{static} ha una visibilità più ristretta: è visibile solo dalle
funzioni del suo file definite dopo di lei, e non da quelle di altri file.
Tentare di raggiungerla con \code{extern} non serve. L'esempio delle slide è
una famiglia di generatori di numeri pseudo-casuali:

\begin{lstlisting}[style=c, name=random.c]
/* A family of pseudo random number generators. */

#define    INITIAL_SEED         17
#define    MULTIPLIER           25173
#define    INCREMENT            13849
#define    MODULUS              65536
#define    FLOATING_MODULUS     65536.0

static unsigned seed = INITIAL_SEED;   /* esterna, ma privata di questo file */

unsigned random(void)
{
   seed = (MULTIPLIER * seed + INCREMENT) % MODULUS;
   return seed;
}

double probability(void)
{
   seed = (MULTIPLIER * seed + INCREMENT) % MODULUS;
   return (seed / FLOATING_MODULUS);
}
\end{lstlisting}

\begin{itemize}
  \item \code{random()} produce un intero fra 0 e \code{MODULUS} (escluso),
        \code{probability()} un reale fra 0 e 1.
  \item Entrambe usano il vecchio valore di \code{seed} (il \textit{seme})
        per calcolare il nuovo, e lo aggiornano: ogni numero dipende dal
        precedente.
  \item Se \code{seed} venisse modificata per errore da un'altra parte del
        programma, la sequenza cambierebbe. Siccome è \code{static}, è
        privata di questo file e nessun altro modulo può toccarla.
\end{itemize}

Lo verifichiamo con un programma di prova, e poi con uno che prova a
“rubare” il seme:

\begin{affianco}[0.5][c]
\begin{lstlisting}[style=c, name=prova.c]
#include <stdio.h>

unsigned random(void);
double probability(void);

int main(void)
{
   int i;

   for (i = 0; i < 5; i++)
      printf("%u ", random());
   printf("\n%f\n", probability());
   return 0;
}
\end{lstlisting}
\risultato
\begin{console}
$ gcc -c random.c
$ gcc -o prova prova.c random.o
$ ./prova
48574 61999 40372 31453 39802
0.537521
$ nm random.o
000000000000002a T probability
0000000000000000 T random
0000000000000000 d seed
\end{console}
\end{affianco}

Il primo numero è proprio $(25173 \cdot 17 + 13849) \bmod 65536 = 48574$.
In \code{nm} il seme compare con la \code{d} \textbf{minuscola}: la
minuscola indica un simbolo locale al file, che il linker non mostra agli
altri. Infatti un programma che dichiara \code{extern unsigned seed;} e
prova ad assegnargli 0 non viene collegato:

\begin{affianco}[0.42][c]
\begin{lstlisting}[style=c, name=ruba.c]
#include <stdio.h>

extern unsigned seed;
unsigned random(void);

int main(void)
{
   seed = 0;
   printf("%u\n", random());
   return 0;
}
\end{lstlisting}
\risultato
Il compilatore accetta \code{ruba.c}: \code{extern} promette che \code{seed}
è definita da qualche parte. È il linker a non trovarla, perché in
\code{random.o} il seme non è visibile da fuori (le righe con
\textit{warning} sono avvisi secondari, conseguenza dello stesso problema).
\end{affianco}

\begin{console}
$ gcc -o ruba ruba.c random.o
/usr/bin/ld: /tmp/ccEe4NI7.o: warning: relocation against `seed' in read-only section `.text'
/usr/bin/ld: /tmp/ccEe4NI7.o: in function `main':
ruba.c:(.text+0xa): undefined reference to `seed'
/usr/bin/ld: warning: creating DT_TEXTREL in a PIE
collect2: error: ld returned 1 exit status
\end{console}

\subsection{\code{static} per le funzioni}

Anche una funzione può essere \code{static}: allora è visibile solo nel file
in cui è definita, e solo le funzioni dello stesso file possono chiamarla.
Serve a definire moduli con funzioni \textbf{private}. Nell'esempio delle
slide due file definiscono ciascuno la propria \code{sf}, \code{static}, e
non c'è alcun conflitto:

\begin{affianco}[0.5][t]
\begin{lstlisting}[style=c, name=main.c]
#include <stdio.h>

void a(void);
void f() { puts("main f"); }
static void sf() { puts("main sf"); }

void m() {
    f();
    sf();
}

int main() {
    m();
    a();
    return 0;
}
\end{lstlisting}
\risultato
\begin{lstlisting}[style=c, name=file.c]
#include <stdio.h>

void f(void);
static void sf() { puts("a sf"); }

void a() {
    f();
    sf();
}
\end{lstlisting}
\begin{console}
$ gcc -c file.c -o file.o
$ gcc -c main.c -o main.o
$ gcc -o main main.o file.o
$ ./main
main f
main sf
main f
a sf
\end{console}
\end{affianco}

\code{puts} scrive una stringa e va a capo. \code{main} chiama \code{m}, che
chiama \code{f} e la \code{sf} del suo file; poi chiama \code{a}, che sta in
\code{file.c}: \code{a} chiama \code{f}, che esiste solo in \code{main.c} ed
è visibile ovunque, e la \code{sf} del \textbf{proprio} file
(figura~\ref{fig:static-funzioni}).

\begin{figure}[H]
  \centering
  \begin{tikzpicture}[font=\FiraCodeNL\scriptsize,
      fn/.style={draw=primary!60!black, fill=white, rounded corners=2pt, line width=0.6pt,
        minimum width=1.5cm, minimum height=0.5cm},
      fs/.style={fn, draw=accent!65!black, fill=accent!10!white}]
    \node[fn] (main) at (0,0) {main};
    \node[fn] (m) at (0,-1.1) {m};
    \node[fn] (f) at (-1.3,-2.3) {f};
    \node[fs] (sf1) at (1.3,-2.3) {sf};
    \scopo{1}{fm}{main.c}{(main) (m) (f) (sf1)}
    \node[fn] (a) at (5.4,-1.1) {a};
    \node[fs] (sf2) at (5.4,-2.3) {sf};
    \scopo{1}{ff}{file.c}{(a) (sf2)}
    \draw[wtarrow] (main) -- (m);
    \draw[wtarrow] (m) -- (f);
    \draw[wtarrow] (m) -- (sf1);
    \draw[wtarrow] (main.east) -| (a.north);
    \draw[wtarrow] (a) -- (sf2);
    \draw[wtarrow, draw=accent!70!black, rounded corners=3pt] (a.west) -- (3.4,-1.1) |- (0.2,-3.15)
      node[pos=0.75, below, mnota, text=accent!70!black] {\code{a} chiama la \code{f} di \code{main.c}} -| (f.south);
    \node[mnota, anchor=west, align=left] at (7.0,-2.3) {\code{static}: visibile\\solo nel suo file};
  \end{tikzpicture}
  \caption{Chi chiama chi: \code{f} è visibile da entrambi i file, ciascuna
           \code{sf} solo dal proprio.}
  \label{fig:static-funzioni}
\end{figure}

\subsection{Riassunto}

\begin{center}
\renewcommand{\arraystretch}{1.2}
\begin{tabular}{@{}l>{\raggedright\arraybackslash}p{4.4cm}>{\raggedright\arraybackslash}p{5.2cm}>{\raggedright\arraybackslash}p{4.6cm}@{}}
  \toprule
  \textbf{classe} & \textbf{dove si dichiara} & \textbf{da dove è visibile} & \textbf{quanto vive} \\
  \midrule
  \code{auto}                 & in un blocco (è il default)       & nel blocco                          & dall'ingresso all'uscita dal blocco \\
  \code{register}             & in un blocco                      & nel blocco                          & come \code{auto} \\
  \code{static}               & in un blocco                      & nel blocco                          & per tutto il programma \\
  esterna (\code{extern})     & fuori dalle funzioni              & nel file; negli altri file con \code{extern} & per tutto il programma \\
  \code{static}               & fuori dalle funzioni              & solo nel suo file                   & per tutto il programma \\
  \bottomrule
\end{tabular}
\end{center}

\insintesi{
  Una funzione ha un'\textbf{intestazione} (tipo del risultato, nome,
  parametri formali) e un \textbf{corpo}; \code{void} indica che non
  restituisce nulla o che non ha parametri. \code{return} termina la funzione
  e restituisce un valore, convertito nel tipo della funzione; il valore si
  può ignorare. Una funzione va dichiarata \textbf{prima} dell'uso, con la
  definizione in stile ANSI o con un \textbf{prototipo}: senza, il
  compilatore non converte gli argomenti e i risultati sono sbagliati. Gli
  argomenti sono passati \textbf{per valore}: ogni chiamata ha la sua zona
  nella pila, con copie dei valori. Le variabili locali sono visibili solo
  nel loro blocco e nascondono quelle esterne con lo stesso nome. Le classi
  di memorizzazione stabiliscono durata e visibilità: \code{auto} (il
  default), \code{extern} (variabili definite altrove), \code{register}
  (superata), \code{static} (variabili locali che conservano il valore,
  variabili e funzioni private di un file).}

\section{Riferimenti}

\begin{itemize}
  \item \textbf{A. Kelley, I. Pohl}, \textit{A Book on C}: capitolo~5
        (funzioni e classi di memorizzazione).
  \item \textbf{Function declarations} e \textbf{Storage-class specifiers}
        --- cppreference.com. \\
        \urlx{https://en.cppreference.com/w/c/language/function_declaration} \\
        \urlx{https://en.cppreference.com/w/c/language/storage_duration}
\end{itemize}


  \chapter{Array e puntatori}
\label{chap:puntatori}

I programmi lavorano spesso su molti dati dello stesso tipo: i voti di una
classe, le misure di un esperimento, le parole di un testo. Gli
\textbf{array} sono il tipo di dato che li raccoglie in variabili
\textbf{indicizzate}. In C gli array sono strettamente legati ai
\textbf{puntatori}: il nome di un array è trattato come un puntatore
costante, i puntatori si possono indicizzare come gli array, e il C ha
un'\textbf{aritmetica dei puntatori}. I puntatori servono anche a ottenere
l'effetto del passaggio “per riferimento” dei parametri, e a chiedere memoria
mentre il programma è in esecuzione. Le stringhe, che in C sono array di
caratteri, e gli array a più dimensioni sono l'argomento del
capitolo~\ref{chap:stringhe}.

\section{Gli array}

\subsection{Tante variabili, un solo nome}

Un programma che gestisce dei voti potrebbe dichiarare
\code{int grade0, grade1, grade2;}. Ma se le variabili necessarie sono
cento, o mille, dichiararle tutte una per una diventa impossibile. Una
variabile \textbf{indicizzata} è molto più comoda:

\begin{lstlisting}[style=c]
int grade[3];                        /* un array di tre interi */

grade[0] = grade[1] = grade[2] = 0;
\end{lstlisting}

\code{grade} è un \term{array} di tre interi, che si chiamano
\code{grade[0]}, \code{grade[1]} e \code{grade[2]}: un solo nome, e un
\textbf{indice} fra parentesi quadre che seleziona l'elemento. Il primo
indice è \textbf{sempre 0}. In memoria gli elementi occupano celle
\textbf{consecutive} (figura~\ref{fig:tre-variabili}).

\begin{figure}[H]
  \centering
  \begin{tikzpicture}[font=\Lato\scriptsize]
    \node[font=\Lato\scriptsize\bfseries] at (1.4,1.0) {tre variabili separate};
    \variabile{g0}{(1.4,0.3)}{grade0}{0}
    \variabile{g1}{(1.4,-0.3)}{grade1}{0}
    \variabile{g2}{(1.4,-0.9)}{grade2}{0}
    \node[font=\Lato\scriptsize\bfseries] at (6.6,1.0) {un array: \code{int grade[3];}};
    \node[mnome, anchor=east] at (5.2,-0.3) {grade};
    \celle[minimum width=1.0cm]{g}{(5.3,-0.3)}{0,0,0}
    \node[mnota, anchor=west, align=left] at (8.6,-0.3) {celle consecutive;\\sotto, l'indice di ciascuna};
  \end{tikzpicture}
  \caption{Tre variabili indipendenti e un array di tre elementi.}
  \label{fig:tre-variabili}
\end{figure}

\subsection{La dichiarazione}

La dichiarazione di un array a una dimensione indica il \textbf{tipo} degli
elementi, l'\textbf{identificatore} e, fra parentesi quadre, un'espressione
\textbf{costante} intera: la dimensione. La dichiarazione
\code{int a[size];} riserva lo spazio per gli elementi da \code{a[0]} a
\code{a[size-1]}. È buona norma definire la dimensione come costante
simbolica, così che per cambiarla basti modificare una riga:

\begin{affianco}[0.5][c]
\begin{lstlisting}[style=c, name=somma.c]
#include <stdio.h>

#define N 4

int main(void)
{
  int a[N] = {3, 7, 1, 9}, i;
  int sum = 0;

  for (i=0; i<N; i++)
    sum += a[i];
  printf("Sum is %d \n", sum);
  return 0;
}
\end{lstlisting}
\risultato
\begin{console}
$ ./somma
Sum is 20
\end{console}
\vspace{2pt}
\wtnota{Il ciclo \code{for} con \code{i} da 0 a \code{N-1} è il modo tipico
di visitare tutti gli elementi.}
\end{affianco}

Come per ogni variabile, ogni elemento occupa i byte del suo tipo: gli
elementi di un array di \code{int} stanno a 4 byte di distanza l'uno
dall'altro (figura~\ref{fig:array-memoria}).

\begin{figure}[H]
  \centering
  \begin{tikzpicture}[font=\Lato\scriptsize]
    \memoria[minimum width=1.6cm]{m}{(2,0)}{1000/3/a[0], 1004/7/a[1], 1008/1/a[2], 1012/9/a[3]}
    \draw[wtbrace] ([xshift=1.9cm]m-1.north east) -- ([xshift=1.9cm]m-4.south east);
    \node[mnota, anchor=west, align=left] at ([xshift=2.25cm]$(m-1.east)!0.5!(m-4.east)$)
      {16 byte consecutivi: 4 elementi\\da 4 byte ciascuno};
  \end{tikzpicture}
  \caption{\code{int a[4] = \{3, 7, 1, 9\};} in memoria (indirizzi di
           esempio).}
  \label{fig:array-memoria}
\end{figure}

\subsection{L'inizializzazione}

Gli array possono avere classe di memorizzazione automatica, esterna o
statica, ma non \code{register} (capitolo~\ref{chap:funzioni}). Nel C
tradizionale solo gli array esterni e statici potevano essere inizializzati
esplicitamente; in ANSI C lo possono essere anche quelli automatici, con
una lista di valori fra graffe:

\begin{itemize}
  \item \code{float v[5] = \{0.0, 1.0, 2.0, 3.0, 4.0\};} assegna
        \code{v[0] = 0.0}, \code{v[1] = 1.0}, e così via;
  \item \code{int v[5] = \{0\};} azzera l'intero vettore;
  \item \code{int v[5] = \{1\};} assegna \code{v[0] = 1} e pone a 0 i
        rimanenti: se i valori sono meno degli elementi, gli altri valgono 0;
  \item se la dimensione manca e c'è l'inizializzazione, la dimensione è il
        numero dei valori: \code{int a[] = \{1, 2, 3, 4, 5\};} equivale a
        \code{int a[5] = \{1, 2, 3, 4, 5\};}.
\end{itemize}

Se invece un array \textbf{non} viene inizializzato, dipende dalla sua
classe: gli array esterni e statici vengono azzerati dal sistema, quelli
automatici no, e contengono ciò che c'era in memoria.

\begin{affianco}[0.5][c]
\begin{lstlisting}[style=c, name=init.c]
#include <stdio.h>

int v[5];

int main(void)
{
   static int a[5];
   int b[5];
   int i;

   for (i=0; i<5; i++)
     printf("%d ",a[i]);
   printf("\n");

   for (i=0; i<5; i++)
     printf("%d ",v[i]);
   printf("\n");

   for (i=0; i<5; i++)
     printf("%d ",b[i]);
   printf("\n");
}
\end{lstlisting}
\risultato
\begin{console}
$ ./init
0 0 0 0 0
0 0 0 0 0
0 0 -2037142800 32652 179344864
$ ./init
0 0 0 0 0
0 0 0 0 0
0 0 744688368 32742 1091613072
\end{console}
\wtnota{\code{a} (statico) e \code{v} (esterno) sono azzerati; \code{b}
(automatico) contiene valori a caso, diversi a ogni esecuzione. \code{gcc}
con \code{-Wall} avvisa: \textit{‘b’ may be used uninitialized}.}
\end{affianco}

\subsection{Gli indici e l'errore più comune}

Se \code{a} è un array, \code{a[expr]}, con \code{expr} espressione intera,
è un suo elemento. Con la dichiarazione \code{int a[N];} gli indici validi
vanno da 0 a \code{N-1}. Accedere ad \code{a[i]} con \code{i} fuori da
questo intervallo è un errore di programmazione molto comune, e il C
\textbf{non lo controlla}: si legge o si scrive memoria che non è stata
riservata per l'array. L'effetto dipende dal sistema: il programma può
fermarsi con un errore, oppure proseguire con risultati sbagliati.

\begin{affianco}[0.5][c]
\begin{lstlisting}[style=c, name=fuori.c]
#include <stdio.h>

int main(void)
{
   int a[5] = {10, 20, 30, 40, 50};
   int i;

   for (i = 0; i <= 5; i++)      /* un giro di troppo */
      printf("a[%d] = %d\n", i, a[i]);
   return 0;
}
\end{lstlisting}
\risultato
\begin{console}
$ ./fuori
a[0] = 10
a[1] = 20
a[2] = 30
a[3] = 40
a[4] = 50
a[5] = 32764
\end{console}
\wtnota{\code{a[5]} non esiste: il programma legge senza protestare ciò che
c'è dopo l'array.}
\end{affianco}

Errori come questo sono difficili da trovare, perché spesso il programma
sembra funzionare. Un aiuto prezioso è l'opzione
\code{-fsanitize=address} di \code{gcc}, che aggiunge al programma dei
controlli sugli accessi alla memoria e lo ferma al primo accesso sbagliato,
dicendo dove (qui il rapporto è abbreviato, e al posto del percorso della
cartella c'è un segnaposto):

\begin{console}
$ gcc -g -fsanitize=address -o fuori fuori.c
$ ./fuori
a[0] = 10
a[1] = 20
a[2] = 30
a[3] = 40
a[4] = 50
=================================================================
==2607==ERROR: AddressSanitizer: stack-buffer-overflow on address 0x7f364ce00034 at pc 0x5575a6d3d467 bp 0x7fff693a7980 sp 0x7fff693a7970
READ of size 4 at 0x7f364ce00034 thread T0
    #0 0x5575a6d3d466 in main (*@\textit{[percorso del file]}@*)/fuori.c:9
(*@\textit{\dots}@*)
    [32, 52) 'a' (line 5) <== Memory access at offset 52 overflows this variable
(*@\textit{\dots}@*)
\end{console}

Il rapporto dice che alla riga 9 si è letto (\textit{READ}) un \code{int}
(4 byte) appena oltre la fine dell'array \code{a}, dichiarato alla riga 5.

\section{I puntatori}

\subsection{Indirizzi}

Una variabile è memorizzata in un certo numero di byte, in una particolare
posizione della memoria identificata da un \textbf{indirizzo}
(capitolo~\ref{chap:calcolatore}). L'indirizzo di una variabile \code{v} si
ottiene con l'operatore \code{\&}: \code{\&v}. Gli indirizzi sono valori, e
come tali si possono memorizzare e manipolare: le variabili che contengono
indirizzi si chiamano \term{puntatori}, e servono ad accedere alla memoria e
a manipolare gli indirizzi.

\subsection{Dichiarare e usare un puntatore}

La dichiarazione

\begin{lstlisting}[style=c]
int *p;
\end{lstlisting}

dichiara la variabile \code{p} come \textbf{puntatore a \code{int}}: può
contenere l'indirizzo di una variabile \code{int}. I suoi valori possibili
sono lo zero e gli interi positivi che sul sistema possono essere
interpretati come indirizzi. Alcuni esempi di assegnazione:

\begin{lstlisting}[style=c]
p = 0;
p = NULL;              /* equivalente a p = 0 */
p = &i;
p = (int *) 1776;      /* un indirizzo assoluto in memoria */
\end{lstlisting}

\code{NULL}, definita negli header standard, vale 0 e indica un puntatore
che \textbf{non punta a niente}. L'ultima assegnazione impone un indirizzo
preciso: non si fa mai nei programmi normali, perché quella cella
probabilmente non appartiene al programma.

L'operatore \code{*}, applicato a un puntatore, dà il \textbf{valore della
variabile puntata}: \code{*p} è il valore della variabile di cui \code{p}
contiene l'indirizzo (si dice che \code{p} viene \textit{dereferenziato}).
\code{*} e \code{\&} sono uno l'inverso dell'altro: \code{i} e \code{*\&i}
sono equivalenti. Seguiamo le istruzioni delle slide
(figura~\ref{fig:puntatore}):

\begin{lstlisting}[style=c]
int a, b, *p;

a = 1;
b = 2;
p = &a;
b = *p;               /* equivale a b = a */
\end{lstlisting}

\begin{figure}[H]
  \centering
  \begin{tikzpicture}[font=\Lato\scriptsize]
    \foreach \col/\testo in {0/{dopo \code{a = 1; b = 2;}}, 1/{dopo \code{p = \&a;}}, 2/{dopo \code{b = *p;}}}
      \node[font=\Lato\scriptsize\bfseries] at (\col*5.4+1.5,1.0) {\testo};
    \variabile{a1}{(1.4,0.3)}{a}{1}
    \variabile{b1}{(1.4,-0.3)}{b}{2}
    \variabile{p1}{(1.4,-0.9)}{p}{\mvuota}
    \begin{scope}[xshift=5.4cm]
      \variabile{a2}{(1.4,0.3)}{a}{1}
      \variabile{b2}{(1.4,-0.3)}{b}{2}
      \variabile{p2}{(1.4,-0.9)}{p}{}
      \riferimento{p2}{to[out=0, in=0, looseness=1.6] (a2.east)}
      \node[mnota, anchor=west] at (2.9,-0.9) {\code{p} contiene l'indirizzo di \code{a}};
    \end{scope}
    \begin{scope}[xshift=10.8cm]
      \variabile{a3}{(1.4,0.3)}{a}{1}
      \variabile[fill=accent!10!white]{b3}{(1.4,-0.3)}{b}{1}
      \variabile{p3}{(1.4,-0.9)}{p}{}
      \riferimento{p3}{to[out=0, in=0, looseness=1.6] (a3.east)}
      \node[mnota, anchor=west, align=left] at (2.9,-0.3) {\code{*p} è il valore\\di \code{a}: 1};
    \end{scope}
  \end{tikzpicture}
  \caption{Un puntatore si disegna come una freccia verso la variabile di
           cui contiene l'indirizzo.}
  \label{fig:puntatore}
\end{figure}

\subsection{Che cosa contiene un puntatore}

Un indirizzo si stampa con la conversione \code{\%p}, in esadecimale:

\begin{affianco}[0.5][c]
\begin{lstlisting}[style=c, name=indirizzo.c]
#include <stdio.h>

int main(void)
{
   int   i = 7, *p;

   p = &i;
   printf("%s%d\n%s%p\n",
      "   Value of i: ", *p,
      "Location of i: ",  p);
   return 0;
}
\end{lstlisting}
\risultato
\begin{console}
$ ./indirizzo
   Value of i: 7
Location of i: 0x7ffda49df4ac
$ ./indirizzo
   Value of i: 7
Location of i: 0x7ffe989ec3bc
\end{console}
\end{affianco}

L'indirizzo cambia a ogni esecuzione: per sicurezza, il sistema operativo
colloca la memoria del programma in posizioni diverse ogni volta. Nelle
slide le due stringhe contengono anche un \code{\%d}
(\code{"   Value of i: \%d"}): così, però, \code{\%d} non è una conversione
(sta in un argomento, non nella stringa di controllo) e viene stampato tale
e quale, \code{   Value of i: \%d7}. La versione qui sopra è quella corretta.

\subsection{Espressioni con i puntatori}

Con le dichiarazioni

\begin{lstlisting}[style=c]
int i = 3, j = 5, *p = &i, *q = &j, *r;
double x;
\end{lstlisting}

ecco alcune espressioni, con l'espressione equivalente e il valore
(verificato con il compilatore, calcolando ciascuna partendo da questi
valori):

\begin{center}
\renewcommand{\arraystretch}{1.15}
\begin{tabular}{@{}>{\FiraCodeNL}l>{\FiraCodeNL}l l@{}}
  \toprule
  \normalfont\textbf{espressione} & \normalfont\textbf{equivalente} & \textbf{valore} \\
  \midrule
  p == \&i            & p == (\&i)                     & 1 \\
  **\&p               & *(*(\&p))                      & 3 \\
  *(r = \&j) *= *p    & (*(r = (\&j))) *= *p           & 15 \\
  7 * *p / *q + 7     & (((7 * (*p))) / (*q)) + 7      & 11 \\
  r = \&x             & r = \&x                        & illegale \\
  \bottomrule
\end{tabular}
\end{center}

\begin{itemize}
  \item \code{**\&p}: \code{\&p} è l'indirizzo di \code{p}; \code{*} lo
        riporta a \code{p}, e un altro \code{*} dà il valore puntato da
        \code{p}, cioè \code{i}: 3.
  \item \code{*(r = \&j) *= *p}: \code{r} riceve l'indirizzo di \code{j}, e
        \code{*r *= *p} moltiplica \code{j} per \code{i}: \code{j} diventa 15.
  \item \code{7 * *p / *q + 7}: $(7 \cdot 3) / 5 + 7 = 21 / 5 + 7 = 4 + 7 =
        11$ (divisione intera). Nelle slide è scritta \code{7**p / *q+7}: gli
        spazi non contano, ma rendono leggibile.
  \item \code{r = \&x} è illegale: \code{r} punta a \code{int}, \code{\&x} è
        l'indirizzo di un \code{double}. \code{gcc} lo segnala (e i
        compilatori più recenti lo considerano un errore):
\end{itemize}

\begin{console}
$ gcc -Wall -o rx rx.c
rx.c: In function (*@‘@*)main(*@’@*):
rx.c:9:5: warning: assignment to (*@‘@*)int *(*@’@*) from incompatible pointer type (*@‘@*)double *(*@’@*) [-Wincompatible-pointer-types]
    9 |   r = &x;
      |     ^
\end{console}

\subsection{Che cosa non si può puntare}

Non tutto ha un indirizzo. Non si può puntare a una \textbf{costante}
(\code{\&3}), a un'\textbf{espressione} (\code{\&(k + 9)}) né a una
variabile \code{register} (\code{\&v}). Le prime due non sono celle di
memoria con un nome (sono \textit{lvalue} mancati); la terza, come abbiamo
visto, sta in un registro:

\begin{console}
$ gcc -Wall -o nonpuntabili nonpuntabili.c
nonpuntabili.c: In function (*@‘@*)main(*@’@*):
nonpuntabili.c:6:8: error: lvalue required as unary (*@‘@*)&(*@’@*) operand
nonpuntabili.c:7:8: error: lvalue required as unary (*@‘@*)&(*@’@*) operand
nonpuntabili.c:8:4: error: address of register variable (*@‘@*)v(*@’@*) requested
\end{console}

Lo stesso vale per un cast: nella cartella degli esempi c'è un tentativo,
\code{r = \&((int) x);}, di far puntare un \code{int *} a un
\code{double} “convertito”. Il risultato di un cast è un valore, non una
variabile, e il compilatore risponde ancora \textit{lvalue required as
unary ‘\&’ operand}.

\section{Il passaggio per indirizzo}

\subsection{Modificare le variabili del chiamante}

Una variabile passata come parametro a una funzione non cambia valore dopo
la chiamata: è il passaggio \textbf{per valore}
(capitolo~\ref{chap:funzioni}). Altri linguaggi hanno anche il passaggio
\textbf{per riferimento}, in cui la funzione può modificare la variabile
passata. In C si ottiene un effetto simile con i puntatori: si passa
l'\textbf{indirizzo} della variabile, e la funzione lavora sulla variabile
attraverso \code{*}.

\begin{affianco}[0.5][t]
\begin{lstlisting}[style=c, name=valore.c]
int f(int a, int b)
{
  a = a+1;
  b = b+1;
  return (a*b);
}

int main()
{
  int a = 1, b = 2, c;

  c = f(a,b);
  printf("a = %d\nb = %d\nc = %d \n", a, b, c);
}
\end{lstlisting}
\begin{console}
$ ./valore
a = 1
b = 2
c = 6
\end{console}
\risultato
\begin{lstlisting}[style=c, name=indirizzi.c]
int f(int *a, int *b)
{
  *a = *a + 1;
  *b = *b + 1;
  return (*a * *b);
}

int main()
{
  int a = 1, b = 2, c;

  c = f(&a, &b);
  printf("a = %d\nb = %d\nc = %d \n", a, b, c);
}
\end{lstlisting}
\begin{console}
$ ./indirizzi
a = 2
b = 3
c = 6
\end{console}
\end{affianco}

A sinistra \code{f} incrementa le proprie copie; a destra riceve gli
indirizzi delle variabili di \code{main} e incrementa \textbf{quelle}
(figura~\ref{fig:per-indirizzo}). Si noti che \code{a} e \code{b}, dentro
\code{f}, ora sono puntatori: \code{*a} è l'intero puntato.

\begin{figure}[H]
  \centering
  \begin{tikzpicture}[font=\Lato\scriptsize]
    \variabile{ma}{(1.6,0.0)}{a}{1 \textcolor{accent!75!black}{$\rightarrow$ 2}}
    \variabile{mb}{(1.6,-0.6)}{b}{2 \textcolor{accent!75!black}{$\rightarrow$ 3}}
    \variabile{mc}{(1.6,-1.2)}{c}{6}
    \scopo{2}{m}{main}{(ma-nome) (ma) (mb-nome) (mb) (mc-nome) (mc)}
    \node[mvar, fill=accent!10!white] (fa) at (6.4,-0.2) {};
    \node[mnome, anchor=west] (fa-nome) at ([xshift=2pt]fa.east) {a};
    \node[mvar, fill=accent!10!white] (fb) at (6.4,-0.8) {};
    \node[mnome, anchor=west] (fb-nome) at ([xshift=2pt]fb.east) {b};
    \scopo{3}{f}{f}{(fa-nome) (fa) (fb-nome) (fb)}
    \riferimento{fa}{-- (ma.east)}
    \riferimento{fb}{-- (mb.east)}
    \node[mnota, anchor=west, align=left] at (7.9,-0.5) {i parametri di \code{f} contengono gli indirizzi:\\\code{*a = *a + 1} incrementa la \code{a} di \code{main}};
  \end{tikzpicture}
  \caption{Il passaggio per indirizzo: \code{f} arriva alle variabili di
           \code{main} attraverso i puntatori.}
  \label{fig:per-indirizzo}
\end{figure}

\warning{Una svista nelle slide}{
  Nella versione delle slide il corpo di \code{f} è rimasto quello di
  prima: \code{a = a+1; b = b+1; return (a*b);}. Ma ora \code{a} e
  \code{b} sono puntatori: \code{a = a+1} sposterebbe il puntatore, non
  incrementerebbe il valore, e moltiplicare due puntatori non ha senso. Il
  compilatore infatti si ferma: \textit{invalid operands to binary *
  (have ‘int *’ and ‘int *’)}. Bisogna usare \code{*} ovunque si intenda il
  valore puntato, come nella versione qui sopra.}

\subsection{Lo scambio di due variabili}

Lo scambio (\textit{swap}) dei valori di due variabili si fa di solito con
una variabile di appoggio: \code{tmp = a; a = b; b = tmp;}. Una funzione che
scambi due variabili del chiamante deve riceverne gli indirizzi:

\begin{affianco}[0.5][c]
\begin{lstlisting}[style=c, name=swap.c]
#include <stdio.h>

void   swap(int *, int *);

int main(void)
{
   int    i = 3, j = 5;

   swap(&i, &j);
   printf("%d  %d\n", i, j);       /* 5  3 is printed */
   return 0;
}

void swap(int *p, int *q)
{
   int   tmp;

   tmp = *p;
   *p = *q;
   *q = tmp;
}
\end{lstlisting}
\risultato
\begin{console}
$ ./swap
5  3
\end{console}
\end{affianco}

\begin{figure}[H]
  \centering
  \begin{tikzpicture}[font=\Lato\scriptsize]
    \foreach \col/\testo/\vi/\vj/\vt in {0/{\code{tmp = *p;}}/3/5/3, 1/{\code{*p = *q;}}/5/5/3, 2/{\code{*q = tmp;}}/5/3/3} {
      \begin{scope}[xshift=\col*5.6cm]
        \node[font=\Lato\scriptsize\bfseries] at (2.2,1.0) {dopo \testo};
        \variabile{i\col}{(1.0,0.3)}{i}{\vi}
        \variabile{j\col}{(1.0,-0.3)}{j}{\vj}
        \node[mvar, fill=accent!10!white] (p\col) at (3.6,0.3) {};
        \node[mnome, anchor=west] at ([xshift=2pt]p\col.east) {p};
        \node[mvar, fill=accent!10!white] (q\col) at (3.6,-0.3) {};
        \node[mnome, anchor=west] at ([xshift=2pt]q\col.east) {q};
        \node[mvar, fill=accent!10!white] (t\col) at (3.6,-0.9) {\vt};
        \node[mnome, anchor=west] at ([xshift=2pt]t\col.east) {tmp};
        \riferimento{p\col}{-- (i\col.east)}
        \riferimento{q\col}{-- (j\col.east)}
      \end{scope}
    }
  \end{tikzpicture}
  \caption{\code{swap} passo per passo: \code{p} e \code{q} puntano a
           \code{i} e \code{j} di \code{main}.}
  \label{fig:swap}
\end{figure}

È anche il motivo per cui \code{scanf} riceve \code{\&x}: per scrivere in
\code{x}, deve sapere dove si trova (capitolo~\ref{chap:primi-programmi}).

\section{Array e puntatori}

\subsection{Il nome di un array è un indirizzo}

Il nome di un array è un indirizzo, cioè il valore di un puntatore: quello
del suo primo elemento. Da qui le relazioni fondamentali:

\begin{itemize}
  \item se \code{a} è un array e \code{i} un intero, \code{a[i]} (l'elemento
        di indice \code{i}) è equivalente a \code{*(a + i)} (il contenuto
        dell'indirizzo \code{a + i});
  \item se \code{p} è un puntatore, anche \code{p[i]} è equivalente a
        \code{*(p + i)}: i puntatori si possono indicizzare come gli array;
  \item con \code{int a[N], *p;}, l'istruzione \code{p = a;} equivale a
        \code{p = \&a[0];}, e \code{p = a + 1;} equivale a
        \code{p = \&a[1];}.
\end{itemize}

Il punto chiave è l'\textbf{aritmetica dei puntatori}: sommare 1 a un
puntatore non lo sposta di un byte, ma di \textbf{un elemento}, cioè di
tanti byte quanti ne occupa il tipo puntato (4 per un \code{int}).

\begin{affianco}[0.55][c]
\begin{lstlisting}[style=c, name=aritmetica.c]
#include <stdio.h>

int main(void)
{
   int a[3] = {10, 20, 30};
   int *p = a;

   printf("a     = %p\n", (void *) a);
   printf("&a[0] = %p\n", (void *) &a[0]);
   printf("a + 1 = %p\n", (void *) (a + 1));
   printf("&a[1] = %p\n", (void *) &a[1]);
   printf("a + 2 = %p\n", (void *) (a + 2));
   printf("*(a + 2) = %d   p[2] = %d   *(p + 2) = %d\n",
          *(a + 2), p[2], *(p + 2));
   return 0;
}
\end{lstlisting}
\risultato
\begin{console}
$ ./aritmetica
a     = 0x7ffc5944463c
&a[0] = 0x7ffc5944463c
a + 1 = 0x7ffc59444640
&a[1] = 0x7ffc59444640
a + 2 = 0x7ffc59444644
*(a + 2) = 30   p[2] = 30   *(p + 2) = 30
\end{console}
\wtnota{Da \code{a} ad \code{a + 1} l'indirizzo cresce di 4
(\code{...63c}, \code{...640}): un \code{int}.}
\end{affianco}

\begin{figure}[H]
  \centering
  \begin{tikzpicture}[font=\FiraCodeNL\scriptsize]
    \noindici
    \celle[minimum width=2.0cm]{c}{(0,0)}{10,20,30}
    \foreach \k/\ind in {0/1000, 1/1004, 2/1008}
      \node[mindir, anchor=north] at ([yshift=-1pt]c-\k.south) {\ind};
    \foreach \k/\testo in {0/{a\\\&a[0]\\p}, 1/{a + 1\\\&a[1]}, 2/{a + 2\\\&a[2]}} {
      \node[font=\FiraCodeNL\scriptsize, text=mref, anchor=south, align=center] (l\k) at ([yshift=0.5cm]c-\k.north) {\testo};
      \draw[mrif] (l\k.south) -- (c-\k.north);
    }
    \node[mnota, anchor=west, align=left] at ([xshift=0.5cm]c-2.east) {\code{a[i]} è \code{*(a + i)}:\\il contenuto della cella \code{a + i}};
  \end{tikzpicture}
  \caption{\code{a + i} è l'indirizzo dell'elemento \code{i} (indirizzi di
           esempio): ogni passo vale 4 byte.}
  \label{fig:aritmetica}
\end{figure}

\subsection{Quattro modi di sommare un array}

Grazie a queste equivalenze, la somma degli elementi di un array si può
scrivere in molti modi, tutti equivalenti:

\begin{affianco}[0.55][c]
\begin{lstlisting}[style=c, name=quattro.c]
#include <stdio.h>

#define N 5

int main(void)
{
   int a[N] = {1, 2, 3, 4, 5}, i, *p, sum;

   sum = 0;
   for (i = 0; i < N; i++)
      sum += a[i];
   printf("a[i]:      %d\n", sum);

   sum = 0;
   for (p = a, i = 0; i < N; i++)
      sum += p[i];
   printf("p[i]:      %d\n", sum);

   sum = 0;
   for (i = 0; i < N; i++)
      sum += *(a + i);
   printf("*(a + i):  %d\n", sum);

   sum = 0;
   for (p = a; p < &a[N]; p++)
      sum += *p;
   printf("*p, p++:   %d\n", sum);
   return 0;
}
\end{lstlisting}
\risultato
\begin{console}
$ ./quattro
a[i]:      15
p[i]:      15
*(a + i):  15
*p, p++:   15
\end{console}
\vspace{2pt}
\wtnota{Nell'ultimo ciclo è il puntatore a scorrere l'array: parte da
\code{a} e avanza di un elemento a ogni giro, finché arriva a
\code{\&a[N]}, l'indirizzo subito dopo l'ultimo elemento.}
\end{affianco}

\subsection{Un array non si può spostare}

C'è però una differenza fra un array e un puntatore: un puntatore può
assumere indirizzi diversi, mentre il nome di un array è un indirizzo
\textbf{costante}, che non si può cambiare. Le istruzioni \code{a = p;},
\code{++a;} e \code{a += 2;} sono illegali:

\begin{console}
$ gcc -Wall -o costante costante.c
costante.c: In function (*@‘@*)main(*@’@*):
costante.c:5:6: error: assignment to expression with array type
costante.c:6:4: error: lvalue required as increment operand
costante.c:7:6: error: assignment to expression with array type
\end{console}

\section{Gli array come parametri}

\subsection{Si passa l'indirizzo}

Nella definizione di una funzione, un parametro dichiarato come array è in
realtà un \textbf{puntatore}. Quando si passa un array come argomento, viene
passato per valore il suo \textbf{indirizzo base}: gli elementi
\textbf{non vengono copiati}. Come esempio, una funzione che somma gli
elementi di un array di \code{double}: le due intestazioni seguenti sono
equivalenti.

\begin{affianco}[0.5][c]
\begin{lstlisting}[style=c]
double sum_array(double a[], int n)
\end{lstlisting}
\risultato
\begin{lstlisting}[style=c]
double sum_array(double *a, int n)
\end{lstlisting}
\end{affianco}

Siccome la funzione riceve solo un indirizzo, non sa quanto è lungo
l'array: la dimensione va passata come parametro (qui \code{n}). E siccome
riceve un indirizzo qualsiasi, può lavorare anche su una \textbf{parte}
dell'array. Con \code{double v[100];} e \code{v[i]} uguale a \code{i}:

\begin{affianco}[0.6][c]
\codicepiccolo
\begin{lstlisting}[style=c, name=sum_array.c]
#include <stdio.h>

double sum_array(double a[], int n)   /* n e' la dimensione di a[] */
{
   int    i;
   double sum = 0.0;

   for (i = 0; i < n; ++i)
      sum += a[i];
   return sum;
}

int main(void)
{
   double v[100];
   int    i, k = 10;

   for (i = 0; i < 100; i++)
      v[i] = i;                       /* v[i] vale i */
   printf("%.0f\n", sum_array(v, 100));        /* v[0] + ... + v[99]   */
   printf("%.0f\n", sum_array(v, 88));         /* v[0] + ... + v[87]   */
   printf("%.0f\n", sum_array(&v[7], k - 7));  /* v[7] + ... + v[k-1]  */
   printf("%.0f\n", sum_array(v + 7, 2 * k));  /* v[7] + ... + v[2k+6] */
   return 0;
}
\end{lstlisting}
\risultato
\begin{console}
$ ./sum_array
4950
3828
24
330
\end{console}
\wtnota{Con \code{k} uguale a 10: $7 + 8 + 9 = 24$, e la somma da 7 a 26
vale 330.}
\end{affianco}

Nella chiamata \code{sum\_array(v + 7, 2 * k)} il parametro \code{a} punta
a \code{v[7]}: per la funzione \code{a[0]} è \code{v[7]}, \code{a[1]} è
\code{v[8]}, e così via (figura~\ref{fig:parte-array}).

\begin{figure}[H]
  \centering
  \begin{tikzpicture}[font=\FiraCodeNL\scriptsize]
    \noindici
    \celle[minimum width=0.9cm]{v}{(0,0)}{0,\ldots,6,7,8,9,\ldots,26,27,\ldots}
    \foreach \k/\ind in {0/v[0], 2/v[6], 3/v[7], 4/v[8], 5/v[9], 7/v[26], 8/v[27]}
      \node[mindir, anchor=north] at ([yshift=-1pt]v-\k.south) {\ind};
    \foreach \k/\ind in {3/a[0], 4/a[1], 5/a[2], 7/a[19]}
      \node[font=\FiraCodeNL\tiny, text=accent!75!black, anchor=south] at ([yshift=1pt]v-\k.north) {\ind};
    \draw[wtbrace, draw=accent!70!black] ([yshift=-0.35cm]v-3.south west) -- ([yshift=-0.35cm]v-7.south east);
    \node[mnota, anchor=north, text=accent!70!black] at ([yshift=-0.65cm]$(v-3.south)!0.5!(v-7.south)$) {i 20 elementi sommati da \code{sum\_array(v + 7, 20)}};
    \node[font=\FiraCodeNL\scriptsize, text=mref, anchor=south] (pa) at ([yshift=0.8cm]v-3.north) {a};
    \draw[mrif] (pa.south) -- ([yshift=0.3cm]v-3.north);
  \end{tikzpicture}
  \caption{La funzione vede l'array a partire dall'indirizzo che riceve.}
  \label{fig:parte-array}
\end{figure}

\subsection{\code{sizeof} su un parametro}

Una conseguenza che sorprende: dentro la funzione, \code{sizeof} applicato
al parametro dà la dimensione di un \textbf{puntatore}, non dell'array.

\begin{affianco}[0.55][c]
\begin{lstlisting}[style=c, name=dimensione.c]
#include <stdio.h>

void stampa_dimensione(int a[])
{
   printf("nella funzione: sizeof(a) = %zu\n", sizeof(a));
}

int main(void)
{
   int v[10];

   printf("in main:        sizeof(v) = %zu\n", sizeof(v));
   stampa_dimensione(v);
   return 0;
}
\end{lstlisting}
\risultato
\begin{console}
$ ./dimensione
in main:        sizeof(v) = 40
nella funzione: sizeof(a) = 8
\end{console}
\wtnota{40 byte sono i dieci \code{int}; 8 byte sono un indirizzo.
\code{gcc} avvisa: \textit{‘sizeof’ on array function parameter ‘a’ will
return size of ‘int *’}.}
\end{affianco}

\section{Ordinare un array}

Un'applicazione classica degli array (e dei puntatori, per gli scambi) è
l'\textbf{ordinamento}: disporre gli elementi in ordine crescente. Le
slide presentano due algoritmi semplici. Entrambi usano la \code{swap} vista
sopra, a cui passano gli indirizzi di due elementi, \code{a + j - 1} e
\code{a + j}, cioè \code{\&a[j-1]} e \code{\&a[j]}.

\subsection{Il bubble sort}

\begin{affianco}[0.55][c]
\begin{lstlisting}[style=c, name=bubblesort.c]
void bubble(int *a, int n)     /* n is the size of a */
{
   int    i, j;

   for (i = 0; i < n - 1; ++i)
      for (j = n - 1; j > i; --j)
         if (a[j-1] > a[j])
            swap(a+j-1, a+j);
}
\end{lstlisting}
\risultato
\begin{console}
$ ./bubblesort
7 3 66 3 -5 22 -77 2
-77 -5 2 3 3 7 22 66
\end{console}
\wtnota{\code{main} stampa l'array prima e dopo la chiamata
\code{bubble(v, 8)}.}
\end{affianco}

A ogni passo (un giro del ciclo esterno), il ciclo interno scorre l'array
\textbf{dal fondo verso l'inizio} confrontando ogni coppia di elementi
vicini, e li scambia se sono in ordine sbagliato. Così il più piccolo degli
elementi rimasti “risale” come una bolla fino alla posizione \code{i}:
dopo il passo $k$, i primi $k$ elementi sono definitivamente al loro posto
(figura~\ref{fig:bubble}).

\begin{figure}[H]
  \centering
  \setlength{\tabcolsep}{6pt}
  \renewcommand{\arraystretch}{1.12}
  \newcommand{\ok}[1]{\cellcolor{green-gradient-1!45!white}#1}
  \begin{tabular}{@{}l>{\FiraCodeNL}r>{\FiraCodeNL}r>{\FiraCodeNL}r>{\FiraCodeNL}r>{\FiraCodeNL}r>{\FiraCodeNL}r>{\FiraCodeNL}r>{\FiraCodeNL}r@{}}
    \toprule
    Dati    & 7 & 3 & 66 & 3 & -5 & 22 & -77 & 2 \\
    Passo 1 & \ok{-77} & 7 & 3 & 66 & 3 & -5 & 22 & 2 \\
    Passo 2 & \ok{-77} & \ok{-5} & 7 & 3 & 66 & 3 & 2 & 22 \\
    Passo 3 & \ok{-77} & \ok{-5} & \ok{2} & 7 & 3 & 66 & 3 & 22 \\
    Passo 4 & \ok{-77} & \ok{-5} & \ok{2} & \ok{3} & 7 & 3 & 66 & 22 \\
    Passo 5 & \ok{-77} & \ok{-5} & \ok{2} & \ok{3} & \ok{3} & 7 & 22 & 66 \\
    Passo 6 & \ok{-77} & \ok{-5} & \ok{2} & \ok{3} & \ok{3} & \ok{7} & 22 & 66 \\
    Passo 7 & \ok{-77} & \ok{-5} & \ok{2} & \ok{3} & \ok{3} & \ok{7} & \ok{22} & \ok{66} \\
    \bottomrule
  \end{tabular}
  \caption{Il bubble sort passo per passo (verificato stampando l'array dopo
           ogni giro del ciclo esterno): in verde la parte già ordinata.}
  \label{fig:bubble}
\end{figure}

Gli ultimi passi non cambiano più nulla: l'array è già ordinato dal passo 5,
ma l'algoritmo non se ne accorge e fa comunque tutti gli $n - 1$ passi.

\subsection{Il selection sort}

\begin{affianco}[0.55][c]
\begin{lstlisting}[style=c, name=selectionsort.c]
void selection_sort(int *a, int n)     /* n is the size of a */
{
   int    i, j;

   for (i = 0; i < n - 1; ++i)
      for (j = i + 1; j < n ; j++)
         if (a[j] < a[i])
            swap(a+j, a+i);
}
\end{lstlisting}
\risultato
\begin{console}
$ ./selectionsort
7 3 66 3 -5 22 -77 2
-77 -5 2 3 3 7 22 66
\end{console}
\end{affianco}

Qui al passo \code{i} si confronta l'elemento in posizione \code{i} con
tutti quelli che seguono, scambiandolo ogni volta che se ne trova uno più
piccolo: alla fine del passo, in posizione \code{i} c'è il più piccolo degli
elementi rimasti, cioè quello che va lì (figura~\ref{fig:selection}).

\begin{figure}[H]
  \centering
  \setlength{\tabcolsep}{6pt}
  \renewcommand{\arraystretch}{1.12}
  \newcommand{\ok}[1]{\cellcolor{green-gradient-1!45!white}#1}
  \begin{tabular}{@{}l>{\FiraCodeNL}r>{\FiraCodeNL}r>{\FiraCodeNL}r>{\FiraCodeNL}r>{\FiraCodeNL}r>{\FiraCodeNL}r>{\FiraCodeNL}r>{\FiraCodeNL}r@{}}
    \toprule
    Dati    & 7 & 3 & 66 & 3 & -5 & 22 & -77 & 2 \\
    Passo 1 & \ok{-77} & 7 & 66 & 3 & 3 & 22 & -5 & 2 \\
    Passo 2 & \ok{-77} & \ok{-5} & 66 & 7 & 3 & 22 & 3 & 2 \\
    Passo 3 & \ok{-77} & \ok{-5} & \ok{2} & 66 & 7 & 22 & 3 & 3 \\
    Passo 4 & \ok{-77} & \ok{-5} & \ok{2} & \ok{3} & 66 & 22 & 7 & 3 \\
    Passo 5 & \ok{-77} & \ok{-5} & \ok{2} & \ok{3} & \ok{3} & 66 & 22 & 7 \\
    Passo 6 & \ok{-77} & \ok{-5} & \ok{2} & \ok{3} & \ok{3} & \ok{7} & 66 & 22 \\
    Passo 7 & \ok{-77} & \ok{-5} & \ok{2} & \ok{3} & \ok{3} & \ok{7} & \ok{22} & \ok{66} \\
    \bottomrule
  \end{tabular}
  \caption{Il selection sort passo per passo: a ogni passo il minimo dei
           rimanenti va nella sua posizione definitiva.}
  \label{fig:selection}
\end{figure}

\section{L'allocazione dinamica}

\subsection{\code{malloc}, \code{calloc} e \code{free}}

Finora la dimensione di un array era una costante, decisa scrivendo il
programma. Spesso però la si conosce solo durante l'esecuzione (quanti
numeri inserirà l'utente?). Per questo il C fornisce le funzioni
\code{malloc} (\textit{memory allocation}) e \code{calloc}
(\textit{contiguous allocation}), con i prototipi in \code{stdlib.h}: creano
spazio \textbf{dinamicamente}, mentre il programma è in esecuzione, per
array, strutture e unioni.

\begin{itemize}
  \item \code{calloc(n, el\_size)} riserva spazio contiguo per un array di
        \code{n} elementi da \code{el\_size} byte ciascuno (la dimensione
        si ottiene con \code{sizeof}), e lo inizializza con tutti i bit a
        zero. I due parametri sono di tipo \code{size\_t}.
  \item \code{malloc(size)} riserva \code{size} byte e non li inizializza:
        può essere un po' più veloce.
  \item Se riescono, entrambe restituiscono un puntatore di tipo
        \code{void *} (un indirizzo “generico”) alla base della memoria;
        altrimenti restituiscono \code{NULL}.
\end{itemize}

\begin{lstlisting}[style=c]
a = (int *) calloc(n, sizeof(int));
a = (int *) malloc(n * sizeof(int));
\end{lstlisting}

Il cast \code{(int *)} si può anche omettere, perché in C un
\code{void *} si converte da solo in qualsiasi altro puntatore. Lo spazio
allocato dinamicamente \textbf{non viene liberato} automaticamente
all'uscita dalle funzioni: va restituito esplicitamente con \code{free(ptr)}.
Se \code{ptr} è \code{NULL}, \code{free} non fa nulla; altrimenti
\code{ptr} deve essere esattamente l'indirizzo base restituito da
\code{malloc} o \code{calloc}, o il comportamento è indefinito (e l'effetto
dipende dal sistema).

\subsection{Un esempio}

Il programma seguente chiede ripetutamente un numero \code{n}, alloca un
array di \code{n} interi, lo riempie di numeri casuali fra $-9$ e 9, lo
stampa con la somma e lo libera. \code{rand()} (in \code{stdlib.h})
restituisce un intero pseudo-casuale, e \code{srand(time(NULL))} sceglie il
seme in base all'ora, così che i numeri cambino a ogni esecuzione.

\begin{lstlisting}[style=c, name=dinamico.c]
#include <stdio.h>
#include <stdlib.h>
#include <time.h>

void   fill_array(int *a, int n);
int    sum_array(int *a, int n);
void   wrt_array(int *a, int n);

int main(void)
{
   int   *a, n;

   srand(time(NULL));     /* seed the random number generator */
   printf("Input n: ");
   while (scanf("%d", &n) == 1 && n >= 1) {
      putchar('\n');
      a = calloc(n, sizeof(int));     /* allocate space for a[] */
      fill_array(a, n);
      wrt_array(a, n);
      printf("sum = %d\n\n", sum_array(a, n));
      free(a);
      printf("Input n: ");
   }
   printf("\nBye!\n\n");
   return 0;
}

void fill_array(int *a, int n)
{
   int   i;

   for (i = 0; i < n; ++i)
      a[i] = rand() % 19 - 9;
}

int sum_array(int *a, int n)
{
   int   i, sum = 0;

   for (i = 0; i < n; ++i)
      sum += a[i];
   return sum;
}

void wrt_array(int *a, int n)
{
   int   i;

   printf("a = [");
   for (i = 0; i < n; ++i)
      printf("%d%s", a[i], ((i < n - 1) ? ", " : "]\n"));
}
\end{lstlisting}

\subsection{Il programma in esecuzione}

Ecco un'esecuzione (i numeri, naturalmente, cambiano ogni volta):

\begin{affianco}[0.42][c]
\begin{console}
$ ./dinamico
Input n: <<5>>

a = [8, -1, -5, -9, -3]
sum = -10

Input n: <<3>>

a = [-4, -1, 7]
sum = 2

Input n: <<0>>

Bye!
\end{console}
\risultato
\begin{itemize}
  \item \code{rand() \% 19} vale da 0 a 18; togliendo 9 si va da $-9$ a 9.
  \item In \code{wrt\_array}, \code{(i < n - 1) ? ", " : "]\textbackslash
        n"} è l'\textbf{operatore condizionale}: vale \code{", "} se la
        condizione è vera, \code{"]\textbackslash n"} altrimenti. Così dopo
        l'ultimo elemento si chiude la parentesi invece di mettere la
        virgola.
  \item Il ciclo continua finché \code{scanf} legge un numero
        \textbf{e} il numero è almeno 1.
\end{itemize}
\end{affianco}

\warning{Una svista nelle slide}{
  Nelle slide la condizione è \code{scanf("\%d", \&n) == 1 || n >= 1}, con
  l'\textit{or}. Così, se si inserisce 0, la prima parte è vera e il ciclo
  prosegue con un array vuoto; e se l'input finisce (\tasti{Ctrl,D}) dopo un
  numero valido, la prima parte è falsa ma \code{n} è rimasto quello di
  prima, la seconda è vera, e il ciclo non termina più (in una prova ha
  stampato oltre 10 milioni di righe in un secondo). La condizione giusta,
  quella del libro da cui viene l'esempio, è con \code{\&\&}.}

\subsection{Lo stack e lo heap}

Dove finisce la memoria allocata dinamicamente? La memoria di un programma
in esecuzione è divisa in zone (figura~\ref{fig:zone-memoria}):

\begin{itemize}
  \item lo \term{stack} (la pila) è lo spazio di lavoro delle funzioni:
        a ogni chiamata viene riservato un blocco in cima allo stack per le
        variabili locali e alcuni dati, e viene tolto al ritorno
        (figura~\ref{fig:pila});
  \item lo \term{heap} è la memoria riservata all'allocazione dinamica:
        \code{malloc} e \code{calloc} prendono spazio da lì, e \code{free}
        lo restituisce;
  \item le variabili globali e \code{static} hanno una loro zona, che esiste
        per tutta l'esecuzione, e il codice del programma un'altra ancora.
\end{itemize}

Lo stack è tipicamente molto più piccolo dello heap (su Linux, di solito, 8
MB): un array enorme dichiarato come variabile locale può esaurirlo
(\textit{stack overflow}), mentre lo stesso array allocato con
\code{malloc} non dà problemi.

\begin{figure}[H]
  \centering
  \begin{tikzpicture}[font=\Lato\scriptsize,
      zona/.style={draw=primary!60!black, line width=0.7pt, minimum width=4.2cm, align=center}]
    \node[zona, fill=accent!10!white, minimum height=1.3cm] (st) at (0,2.7) {\textbf{stack}\\variabili locali, parametri};
    \node[zona, fill=white, minimum height=1.2cm, draw=neutral!40!white, dashed] (lib) at (0,1.45) {\color{neutral!45!black}spazio libero};
    \node[zona, fill=green-gradient-1!40!white, minimum height=1.3cm] (hp) at (0,0.2) {\textbf{heap}\\\code{malloc}, \code{calloc}};
    \node[zona, fill=primary!8!white, minimum height=0.75cm] (gl) at (0,-0.83) {variabili globali e \code{static}};
    \node[zona, fill=neutral!10!white, minimum height=0.75cm] (cd) at (0,-1.6) {codice del programma};
    \draw[wtarrowacc] ([xshift=-1.5cm]st.south) -- ++(0,-0.45) node[left, mnota, pos=0.3] {cresce verso il basso};
    \draw[wtarrowacc] ([xshift=-1.5cm]hp.north) -- ++(0,0.45) node[left, mnota, pos=0.3] {cresce verso l'alto};
    \node[mnota, anchor=west] at (2.4,2.7) {un blocco per ogni chiamata, tolto al ritorno};
    \node[mnota, anchor=west] at (2.4,0.2) {resta finché non si chiama \code{free}};
    \node[mnota, anchor=west] at (2.4,-1.2) {esistono per tutta l'esecuzione};
  \end{tikzpicture}
  \caption{Le zone della memoria di un programma (lo schema è semplificato).}
  \label{fig:zone-memoria}
\end{figure}

\insintesi{
  Un \textbf{array} raccoglie elementi dello stesso tipo in celle
  consecutive; gli indici vanno da 0 a $N-1$, e il C non controlla gli
  accessi fuori dai limiti (\code{-fsanitize=address} aiuta a trovarli). Gli
  array automatici non inizializzati contengono valori casuali. Un
  \textbf{puntatore} contiene un indirizzo: \code{\&v} è l'indirizzo di
  \code{v}, \code{*p} il valore puntato, \code{NULL} il puntatore che non
  punta a niente. Passando gli indirizzi una funzione può modificare le
  variabili del chiamante (\code{swap}). Il nome di un array è l'indirizzo
  costante del primo elemento: \code{a[i]} equivale a \code{*(a + i)}, e
  \code{p + 1} avanza di un elemento. Un array passato a una funzione è
  passato come indirizzo (la dimensione va passata a parte). Bubble sort e
  selection sort ordinano con confronti e scambi. \code{malloc} e
  \code{calloc} allocano memoria nello heap durante l'esecuzione, e
  \code{free} la restituisce.}

\section{Riferimenti}

\begin{itemize}
  \item \textbf{A. Kelley, I. Pohl}, \textit{A Book on C}: capitolo~6
        (array, puntatori e stringhe).
  \item \textbf{AddressSanitizer} --- documentazione di GCC
        (\code{-fsanitize=address}). \\
        \urlx{https://gcc.gnu.org/onlinedocs/gcc/Instrumentation-Options.html}
  \item \textbf{Dynamic memory management} --- cppreference.com. \\
        \urlx{https://en.cppreference.com/w/c/memory}
\end{itemize}


  \chapter{Stringhe e array a più dimensioni}
\label{chap:stringhe}

Questo capitolo completa la lezione sugli array e i puntatori con due
applicazioni importantissime. Le \textbf{stringhe}, cioè i testi, in C sono
array di caratteri: vedremo come sono fatte, le funzioni della libreria
standard che le trattano e come si ordinano le parole di un testo con un
array di puntatori. Gli \textbf{array a più dimensioni}, come le matrici,
sono array di array: vedremo come stanno in memoria, come si passano alle
funzioni e come si allocano dinamicamente.

\section{Le stringhe}

\subsection{Un array di caratteri con un terminatore}

Le stringhe sono array a una dimensione di caratteri. In C una stringa è
\textbf{terminata dal carattere nullo} \code{'\textbackslash 0'}, un byte con
tutti i bit a zero, il cui valore decimale è 0: è così che le funzioni
riconoscono dove finisce il testo. Per questo la dimensione di una stringa
deve includere lo spazio per il terminatore: la stringa \code{"a"} ha
\textbf{due} elementi, \code{'a'} e \code{'\textbackslash 0'}. La
dichiarazione

\begin{lstlisting}[style=c]
char s[] = "abc";
\end{lstlisting}

è trattata dal compilatore come
\code{char s[] = \{'a', 'b', 'c', '\textbackslash 0'\};}: un array di 4
caratteri (figura~\ref{fig:stringa}). Il programmatore deve stare attento a
non superare la dimensione dell'array, perché il C non lo controlla.

\begin{figure}[H]
  \centering
  \begin{tikzpicture}[font=\FiraCodeNL\small]
    \celle[minimum width=1.0cm, minimum height=0.6cm]{s}{(0,0)}{a,b,c,\textbackslash 0}
    \foreach \k/\cod in {0/97, 1/98, 2/99, 3/0}
      \node[font=\FiraCodeNL\tiny, text=neutral!45!black, anchor=south] at ([yshift=1pt]s-\k.north) {\cod};
    \node[mnome, anchor=east] at ([xshift=-3pt]s-0.west) {s};
    \node[mnota, anchor=west, align=left] at ([xshift=0.4cm]s-3.east)
      {4 byte: i 3 caratteri e il terminatore\\(sopra, i codici ASCII; sotto, gli indici)};
  \end{tikzpicture}
  \caption{La stringa \code{"abc"} in memoria.}
  \label{fig:stringa}
\end{figure}

\subsection{Le stringhe costanti sono puntatori}

Una stringa costante, come \code{"abc"}, è trattata dal compilatore come un
\textbf{puntatore} al suo primo carattere. Così si può scrivere

\begin{lstlisting}[style=c]
char *p = "abc";

printf("%s %s\n", p, p + 1);
\end{lstlisting}

che stampa \code{abc bc}: la conversione \code{\%s} riceve un indirizzo e
stampa i caratteri da lì fino al \code{'\textbackslash 0'}, e \code{p + 1}
punta alla \code{b} (figura~\ref{fig:stringa-puntatore}). Per lo stesso motivo
\code{"abc"[1]} vale \code{'b'} e \code{*("abc" + 2)} vale \code{'c'}:
espressioni corrette, che però non vanno mai usate nel codice vero.

\begin{figure}[H]
  \centering
  \begin{tikzpicture}[font=\FiraCodeNL\small]
    \noindici
    \celle[minimum width=1.0cm, minimum height=0.6cm]{t}{(0,0)}{a,b,c,\textbackslash 0}
    \foreach \k/\testo in {0/p, 1/p + 1} {
      \node[font=\FiraCodeNL\scriptsize, text=mref, anchor=south] (l\k) at ([yshift=0.6cm]t-\k.north) {\testo};
      \draw[mrif] (l\k.south) -- (t-\k.north);
    }
    \draw[wtbrace] ([yshift=-3pt]t-0.south west) -- ([yshift=-3pt]t-2.south east);
    \node[mnota, anchor=north] at ([yshift=-0.32cm]t-1.south) {\code{\%s} con \code{p}: \code{abc}};
    \draw[wtbrace, draw=accent!70!black] ([yshift=-0.95cm]t-1.south west) -- ([yshift=-0.95cm]t-2.south east);
    \node[mnota, anchor=north, text=accent!70!black] at ([yshift=-1.25cm]$(t-1.south)!0.5!(t-2.south)$) {\code{\%s} con \code{p + 1}: \code{bc}};
  \end{tikzpicture}
  \caption{\code{\%s} stampa dall'indirizzo ricevuto fino al terminatore.}
  \label{fig:stringa-puntatore}
\end{figure}

\subsection{Array o puntatore?}

\code{char s[] = "abc";} e \code{char *p = "abc";} sembrano uguali, ma non lo
sono. Il primo crea un array di 4 caratteri, una copia della stringa che il
programma può modificare. Il secondo crea solo un puntatore alla stringa
costante, che sta in una zona di memoria \textbf{di sola lettura}: tentare di
modificarla è un comportamento indefinito, e su Linux il programma viene
fermato.

\begin{affianco}[0.5][c]
\begin{lstlisting}[style=c, name=modifica.c]
#include <stdio.h>

int main(void)
{
   char s[] = "abc";
   char *p = "abc";

   s[0] = 'x';
   printf("%s\n", s);
   p[0] = 'x';
   printf("%s\n", p);
   return 0;
}
\end{lstlisting}
\risultato
\begin{console}
$ ./modifica
xbc
Segmentation fault
\end{console}
\wtnota{\textit{Segmentation fault} (a volte seguito da \textit{core dumped})
è il messaggio della shell quando il sistema ferma un programma che ha
toccato memoria non sua.}
\end{affianco}

\subsection{La stringa vuota}

Il carattere nullo non va scritto in output: è un segnale per i programmi,
non un carattere da mostrare, e scriverlo in un file potrebbe confondere i
programmi che lo elaborano. Scrivere stringhe \textbf{vuote} (\code{""}, che
contiene solo il terminatore), invece, può essere utile. Per esempio, per
mettere il plurale solo quando serve:

\begin{affianco}[0.6][c]
\begin{lstlisting}[style=c]
for (nfrogs = 1; nfrogs <= 3; nfrogs++)
   printf("%d frog%s\n", nfrogs, (nfrogs == 1) ? "" : "s");
\end{lstlisting}
\risultato
\begin{console}
1 frog
2 frogs
3 frogs
\end{console}
\end{affianco}

\section{Contare le parole}

La funzione seguente conta le parole di una stringa, intese come sequenze di
caratteri separate da spazi. Il parametro \code{s} è un puntatore che
\textbf{scorre} la stringa: \code{++s} lo sposta sul carattere successivo, e
\code{*s} è il carattere su cui si trova. \code{isspace} (in \code{ctype.h})
è vera per gli spazi, le tabulazioni e gli a capo.

\begin{affianco}[0.58][c]
\begin{lstlisting}[style=c, name=word_cnt.c]
/* Count the number of words in a string. */

#include <stdio.h>
#include <ctype.h>

int word_cnt(const char *s)
{
   int    cnt = 0;

   while (*s != '\0') {
      while (isspace(*s))         /* skip white space */
         ++s;
      if (*s != '\0') {           /* found a word */
         ++cnt;
         while (!isspace(*s) && *s != '\0')
            ++s;                  /* skip the word */
      }
   }
   return cnt;
}

int main(void)
{
   printf("%d\n", word_cnt("  Il linguaggio   C  "));
   printf("%d\n", word_cnt(""));
   printf("%d\n", word_cnt("una"));
   return 0;
}
\end{lstlisting}
\risultato
\begin{console}
$ ./word_cnt
3
0
1
\end{console}
\end{affianco}

A ogni giro del ciclo esterno si saltano gli spazi; se non si è arrivati
alla fine, si è trovata una parola: si incrementa il contatore e si salta la
parola fino al prossimo spazio (figura~\ref{fig:word-cnt}).

\begin{figure}[H]
  \centering
  \begin{tikzpicture}[font=\FiraCodeNL\scriptsize,
      sp/.style={mcella, minimum width=0.42cm, minimum height=0.5cm, inner xsep=0pt, fill=neutral!10!white},
      ch/.style={mcella, minimum width=0.42cm, minimum height=0.5cm, inner xsep=0pt}]
    \noindici
    \celle[ch]{w}{(0,0)}{\char"2423,\char"2423,I,l,\char"2423,l,i,n,g,u,a,g,g,i,o,\char"2423,\char"2423,\char"2423,C,\char"2423,\char"2423,\textbackslash 0}
    \foreach \ia/\ib/\nparola in {2/3/1, 5/14/2, 18/18/3} {
      \draw[wtbrace, draw=mok] ([yshift=-3pt]w-\ia.south west) -- ([yshift=-3pt]w-\ib.south east);
      \node[font=\Lato\scriptsize\bfseries, text=mok, anchor=north] at ([yshift=-0.3cm]$(w-\ia.south)!0.5!(w-\ib.south)$) {parola \nparola};
    }
    \node[mnota, anchor=west, align=left] at ([xshift=0.3cm]w-21.east) {\code{\char"2423} sono spazi:\\vengono saltati};
  \end{tikzpicture}
  \caption{\code{word\_cnt} alterna due fasi: salta gli spazi, poi salta una
           parola contandola.}
  \label{fig:word-cnt}
\end{figure}

\section{La parola chiave \code{const}}

\subsection{Variabili che non cambiano}

La parola chiave \code{const}, introdotta dall'ANSI C, indica che una
variabile può essere inizializzata ma non modificata
(capitolo~\ref{chap:primi-programmi}). Con \code{const int k = 3;}
l'istruzione \code{k = 1;} è un errore: \textit{assignment of read-only
variable ‘k’}.

Con i puntatori, \code{const} può riguardare il dato puntato o il puntatore
stesso:

\begin{itemize}
  \item \code{const int *p;} dichiara un puntatore a un \code{const int}:
        il \textbf{dato puntato} non si può modificare attraverso \code{p}
        (\code{*p = 1;} è un errore), ma il puntatore sì (\code{p = \&j;} va
        bene);
  \item \code{int *const p = \&i;} è il contrario: un puntatore
        \textbf{costante}, che punterà sempre a \code{i}, attraverso cui
        però \code{i} si può modificare.
\end{itemize}

\subsection{Una promessa verificata dal compilatore}

Il parametro di \code{word\_cnt} è \code{const char *s}: la funzione
\textbf{promette} di non modificare la stringa che riceve (sposta \code{s},
ma non tocca i caratteri). È un'informazione preziosa per chi la usa, e il
compilatore la fa rispettare. Se nella funzione si prova a scrivere nella
stringa:

\begin{lstlisting}[style=c, name=word_cnt2.c]
/* Count the number of words in a string. */

#include <ctype.h>

int word_cnt(const char *s)
{
   int    cnt = 0;

   while (*s != '\0') {
      while (isspace(*s))         /* skip white space */
         ++s;
      if (*s != '\0') {           /* found a word */
         ++cnt;
         while (!isspace(*s) && *s != '\0') {
            *s = 'a';             /* the compiler will complain */
            ++s;                  /* skip the word */
         }
      }
   }
   return cnt;
}
\end{lstlisting}

\begin{console}
$ gcc -Wall -c word_cnt2.c
word_cnt2.c: In function (*@‘@*)word_cnt(*@’@*):
word_cnt2.c:15:16: error: assignment of read-only location (*@‘@*)*s(*@’@*)
   15 |             *s = 'a';             /* the compiler will complain */
      |                ^
\end{console}

\section{Funzioni per caratteri e stringhe}

\subsection{I caratteri: \code{ctype.h}}

La libreria standard ha molte funzioni per classificare e trasformare i
caratteri, con i prototipi in \code{ctype.h}. Ricevono un \code{int} (come
\code{getchar}, capitolo~\ref{chap:tipi}) e le funzioni di verifica
restituiscono un valore diverso da zero (vero) o zero (falso).

\begin{center}
\renewcommand{\arraystretch}{1.12}
\begin{tabular}{@{}>{\FiraCodeNL}l l@{}}
  \toprule
  \normalfont\textbf{funzione} & \textbf{che cosa fa} \\
  \midrule
  int isblank(int c) & vero se \code{c} è uno spazio o una tabulazione (non c'era nei vecchi compilatori Microsoft) \\
  int isdigit(int c) & vero se \code{c} è una cifra \\
  int isalpha(int c) & vero se \code{c} è una lettera \\
  int isalnum(int c) & vero se \code{c} è una lettera o una cifra \\
  int islower(int c) & vero se \code{c} è una lettera minuscola \\
  int isupper(int c) & vero se \code{c} è una lettera maiuscola \\
  int tolower(int c) & se \code{c} è maiuscola restituisce la minuscola, altrimenti \code{c} \\
  int toupper(int c) & se \code{c} è minuscola restituisce la maiuscola, altrimenti \code{c} \\
  \bottomrule
\end{tabular}
\end{center}

\begin{lstlisting}[style=c, name=ctipi.c]
#include <stdio.h>
#include <ctype.h>

int main(void)
{
   const char *s = "aZ 7!";
   int i;

   printf("car  isalpha isdigit isalnum islower isupper isblank toupper\n");
   for (i = 0; s[i] != '\0'; i++)
      printf(" '%c' %5d %7d %7d %7d %7d %7d %7c\n", s[i],
             isalpha(s[i]) != 0, isdigit(s[i]) != 0, isalnum(s[i]) != 0,
             islower(s[i]) != 0, isupper(s[i]) != 0, isblank(s[i]) != 0,
             toupper(s[i]));
   return 0;
}
\end{lstlisting}

\begin{console}
$ ./ctipi
car  isalpha isdigit isalnum islower isupper isblank toupper
 'a'     1       0       1       1       0       0       A
 'Z'     1       0       1       0       1       0       Z
 ' '     0       0       0       0       0       1
 '7'     0       1       1       0       0       0       7
 '!'     0       0       0       0       0       0       !
\end{console}

Nel programma, \code{!= 0} trasforma qualunque “vero” in 1, per avere una
tabella leggibile.

\subsection{Le stringhe: \code{string.h}}

Le funzioni per le stringhe hanno i prototipi in \code{string.h}. Il
parametro \code{const char *} indica una stringa che la funzione non
modifica; il \code{char *} senza \code{const} è l'array in cui la funzione
scrive, che \textbf{deve avere spazio sufficiente}: nessuna di queste
funzioni lo controlla.

\begin{center}
\renewcommand{\arraystretch}{1.12}
\begin{tabular}{@{}>{\FiraCodeNL\small}l>{\raggedright\arraybackslash}p{8.4cm}@{}}
  \toprule
  \normalfont\textbf{prototipo} & \textbf{che cosa fa} \\
  \midrule
  size\_t strlen(const char *s) & la lunghezza di \code{s}, senza contare il terminatore \\
  char *strcpy(char *s1, const char *s2) & copia \code{s2} nell'array \code{s1}; restituisce \code{s1} \\
  char *strncpy(char *s1, const char *s2, size\_t n) & copia al massimo \code{n} caratteri di \code{s2} in \code{s1} \\
  char *strcat(char *s1, const char *s2) & aggiunge \code{s2} in coda a \code{s1}; restituisce \code{s1} \\
  char *strncat(char *s1, const char *s2, size\_t n) & aggiunge in coda al massimo \code{n} caratteri di \code{s2} \\
  int strcmp(const char *s1, const char *s2) & confronta: restituisce 0 se sono uguali, un negativo se \code{s1} viene prima, un positivo se viene dopo \\
  int strncmp(const char *s1, const char *s2, size\_t n) & come \code{strcmp}, ma solo sui primi \code{n} caratteri \\
  char *strchr(const char *s, int c) & un puntatore alla prima \code{c} in \code{s}, o \code{NULL} \\
  char *strrchr(const char *s, int c) & un puntatore all'ultima \code{c} in \code{s}, o \code{NULL} \\
  char *strstr(const char *s1, const char *s2) & un puntatore alla prima occorrenza di \code{s2} in \code{s1}, o \code{NULL} \\
  char *strpbrk(const char *s1, const char *s2) & un puntatore al primo carattere di \code{s1} che compare in \code{s2}, o \code{NULL} \\
  size\_t strspn(const char *s1, const char *s2) & la lunghezza del tratto iniziale di \code{s1} fatto \textbf{solo} di caratteri di \code{s2} \\
  size\_t strcspn(const char *s1, const char *s2) & la lunghezza del tratto iniziale di \code{s1} fatto di caratteri \textbf{non} in \code{s2} \\
  \bottomrule
\end{tabular}
\end{center}

\code{strcmp} confronta i codici dei caratteri uno alla volta, finché trova
una differenza: per questo l'ordine è quello alfabetico per le minuscole (o
per le maiuscole), ma tutte le maiuscole vengono prima delle minuscole
(capitolo~\ref{chap:tipi}). Ecco gli esempi delle slide per le ultime due
funzioni:

\begin{affianco}[0.55][c]
\begin{lstlisting}[style=c, name=spn.c]
printf("%zu\n", strcspn("abc", "abc"));       /* 0 */
printf("%zu\n", strcspn("abc", "abcde"));     /* 0 */
printf("%zu\n", strcspn("djpc", "lqrabc"));   /* 3 */
printf("%zu\n", strspn("abc", "abc"));        /* 3 */
printf("%zu\n", strspn("abc", "abcde"));      /* 3 */
printf("%zu\n", strspn("djpc", "lqrabc"));    /* 0 */
\end{lstlisting}
\risultato
\begin{itemize}
  \item \code{strcspn("djpc", "lqrabc")} vale 3: \code{d}, \code{j} e
        \code{p} non sono in \code{"lqrabc"}, la \code{c} sì.
  \item \code{strspn("abc", "abcde")} vale 3: tutti i caratteri di
        \code{"abc"} sono in \code{"abcde"}.
  \item I valori sono di tipo \code{size\_t}, che si stampa con \code{\%zu}
        (nelle slide c'è \code{\%d}, e \code{gcc} avvisa).
\end{itemize}
\end{affianco}

\subsection{Come sono fatte}

Non c'è niente di speciale in queste funzioni: si possono scrivere in C, in
poche righe. Ecco le versioni delle slide (con il prefisso \code{my\_}: i
nomi della libreria standard vanno trattati come riservati,
capitolo~\ref{chap:sintassi}):

\begin{affianco}[0.5][c]
\begin{lstlisting}[style=c, name=mystr.c]
size_t my_strlen(const char *s)
{
   size_t n;

   for (n = 0; *s != '\0'; ++s)
      ++n;
   return n;
}

char *my_strcpy(char *s1, const char *s2)
{
   char *p = s1;

   while ((*p++ = *s2++))
      ;
   return s1;
}

char *my_strcat(char *s1, const char *s2)
{
   char *p = s1;

   while (*p)
      ++p;
   while ((*p++ = *s2++))
      ;
   return s1;
}
\end{lstlisting}
\risultato
\begin{lstlisting}[style=c]
int main(void)
{
   char *s1;

   s1 = malloc(10 * sizeof(char));
   my_strcpy(s1, "abc");
   printf("%s %zu\n", s1, my_strlen(s1));
   my_strcat(s1, "efg");
   printf("%s %zu\n", s1, my_strlen(s1));
   free(s1);
   return 0;
}
\end{lstlisting}
\begin{console}
$ ./mystr
abc 3
abcefg 6
\end{console}
\end{affianco}

Il cuore di \code{my\_strcpy} è il ciclo \code{while ((*p++ = *s2++))}, che
fa tutto nella condizione (figura~\ref{fig:strcpy}): copia il carattere
puntato da \code{s2} nella cella puntata da \code{p}, sposta entrambi i
puntatori in avanti (\code{++} postfisso: si usa il valore e poi si
incrementa), e il valore dell'assegnazione, cioè il carattere copiato,
fa da condizione. Quando si copia il \code{'\textbackslash 0'}, che vale 0,
il ciclo si ferma: il terminatore è stato copiato anche lui. Le doppie
parentesi dicono al compilatore (e al lettore) che l'assegnazione nella
condizione è voluta: le slide hanno le parentesi singole, e \code{gcc -Wall}
suggerirebbe di aggiungerle. \code{my\_strcat} prima porta \code{p} sul
terminatore di \code{s1}, poi copia da lì.

\begin{figure}[H]
  \centering
  \begin{tikzpicture}[font=\FiraCodeNL\small]
    \noindici
    \celle[minimum width=0.9cm, minimum height=0.55cm]{src}{(0,0)}{a,b,c,\textbackslash 0}
    \celle[minimum width=0.9cm, minimum height=0.55cm, fill=accent!8!white]{dst}{(0,-1.6)}{a,b,c,\textbackslash 0,\mvuota,\mvuota}
    \node[mnome, anchor=east] at ([xshift=-3pt]src-0.west) {s2};
    \node[mnome, anchor=east] at ([xshift=-3pt]dst-0.west) {s1};
    \foreach \k in {0,1,2,3}
      \draw[wtarrowacc] (src-\k.south) -- node[right, font=\Lato\tiny, text=accent!70!black] {\the\numexpr\k+1\relax} (dst-\k.north);
    \node[mnota, anchor=west, align=left] at ([xshift=0.5cm]src-3.east)
      {un giro per carattere: \code{*p = *s2}, poi\\\code{p} e \code{s2} avanzano; al giro 4 si copia\\il \code{'\textbackslash 0'}, che vale 0, e il ciclo finisce};
  \end{tikzpicture}
  \caption{\code{my\_strcpy(s1, s2)}: quattro giri, l'ultimo copia il
           terminatore.}
  \label{fig:strcpy}
\end{figure}

\subsection{Lo spazio per la copia}

Nelle slide, il programma di prova di \code{strcpy} riserva per \code{s1}
\code{malloc(3*sizeof(char))}, cioè 3 byte, e poi ci copia \code{"abc"}, che
ne occupa 4 con il terminatore. È proprio l'errore da cui mettono in guardia
le funzioni di \code{string.h}: si scrive un byte oltre la memoria riservata.
Il programma sembra funzionare, ma \code{gcc} se ne accorge già in
compilazione, e \code{-fsanitize=address} lo conferma in esecuzione:

\begin{console}
$ gcc -Wall -o strcpy3 strcpy3.c
strcpy3.c: In function (*@‘@*)main(*@’@*):
strcpy3.c:11:8: warning: (*@‘@*)strcpy(*@’@*) writing 4 bytes into a region of size 3 overflows the destination [-Wstringop-overflow=]
   11 |   s1 = strcpy(s1, s2);
      |        ^~~~~~~~~~~~~~
$ ./strcpy3
abc abc
$ gcc -g -fsanitize=address -o strcpy3 strcpy3.c
$ ./strcpy3
=================================================================
==3088==ERROR: AddressSanitizer: heap-buffer-overflow on address 0x502000000013 at pc 0x7f5d792a7923 bp 0x7ffec73255a0 sp 0x7ffec7324d48
WRITE of size 4 at 0x502000000013 thread T0
(*@\textit{\dots}@*)
\end{console}

La regola: per copiare una stringa \code{s2} servono
\code{strlen(s2) + 1} byte, cioè \code{malloc(strlen(s2) + 1)}.

\section{Gli array di puntatori}

\subsection{Un array di stringhe}

La dichiarazione \code{int *a[20];} definisce \code{a} come un array di 20
elementi, ognuno dei quali è un \textbf{puntatore a intero}: può essere
pensato come un array di vettori, ciascuno con la sua lunghezza. Il caso più
comune è l'array di stringhe, \code{char *w[N];}: ogni \code{w[i]} punta a
una stringa, e le stringhe possono essere lunghe quanto serve
(figura~\ref{fig:array-puntatori}).

\begin{figure}[H]
  \centering
  \begin{tikzpicture}[font=\FiraCodeNL\scriptsize,
      cc/.style={mcella, minimum width=0.45cm, minimum height=0.5cm, inner xsep=0pt}]
    \noindici
    \memoria[minimum width=1.0cm]{w}{(0,0)}{w[0]/{}/, w[1]/{}/, w[2]/{}/, w[3]/{}/}
    \celle[cc]{s0}{(3.2,0)}{l,e,\textbackslash 0}
    \celle[cc]{s1}{(3.2,-0.52)}{p,a,r,o,l,e,\textbackslash 0}
    \celle[cc]{s2}{(3.2,-1.04)}{d,i,\textbackslash 0}
    \celle[cc]{s3}{(3.2,-1.56)}{u,n,\textbackslash 0}
    \foreach \k [evaluate=\k as \r using int(\k+1)] in {0,1,2,3}
      \riferimento{w-\r}{-- (s\k-0.west)}
    \node[mnota, anchor=west, align=left] at (7.2,-0.8) {ogni stringa occupa\\solo lo spazio che le serve};
  \end{tikzpicture}
  \caption{Un array di puntatori a stringhe di lunghezza diversa.}
  \label{fig:array-puntatori}
\end{figure}

\subsection{Ordinare le parole di un testo}

L'esempio delle slide è un programma, diviso in più file
(capitolo~\ref{chap:make}), che legge le parole dallo standard input e le
stampa in ordine lessicografico (alfabetico, secondo i codici ASCII).

\begin{affianco}[0.5][t]
\codicepiccolo
\begin{lstlisting}[style=h, name=sort.h]
#define   MAXWORD   50        /* max word size */
#define   N         300       /* array size of w[] */

void   error_exit_calloc_failed(void);
void   error_exit_too_many_words(void);
void   error_exit_word_too_long(void);
void   sort_words(char *w[], int n);
void   swap(char **p, char **q);
void   wrt_words(char *w[], int n);
\end{lstlisting}
\begin{lstlisting}[style=c, name=sort_words.c]
#include <string.h>
#include "sort.h"

void sort_words(char *w[], int n) /* n elements are to be sorted */
{
   int   i, j;

   for (i = 0; i < n; ++i)
      for (j = i + 1; j < n; ++j)
         if (strcmp(w[i], w[j]) > 0)
            swap(&w[i], &w[j]);
}
\end{lstlisting}
\begin{lstlisting}[style=c, name=swap.c]
#include "sort.h"

void swap(char **p, char **q)
{
   char   *tmp;

   tmp = *p;
   *p = *q;
   *q = tmp;
}
\end{lstlisting}
\risultato
\begin{lstlisting}[style=c, name=main.c]
/* Sort words lexicographically. */
#include <stdio.h>
#include <stdlib.h>
#include <string.h>
#include "sort.h"

int main(void)
{
   char   word[MAXWORD];      /* work space */
   char   *w[N];              /* an array of pointers */
   int    n;                  /* number of words to be sorted */
   int    i;

   for (i = 0; scanf("%s", word) == 1; ++i) {
      if (i >= N)
         error_exit_too_many_words();
      if (strlen(word) >= MAXWORD)
         error_exit_word_too_long();
      w[i] = calloc(strlen(word) + 1, sizeof(char));
      if (w[i] == NULL)
         error_exit_calloc_failed();
      strcpy(w[i], word);
   }
   n = i;
   sort_words(w, n);          /* sort the words */
   wrt_words(w, n);           /* write sorted list of words */
   return 0;
}
\end{lstlisting}
\begin{lstlisting}[style=c, name=wrt.c]
#include <stdio.h>
#include "sort.h"

void wrt_words(char *w[], int n)
{
   int   i;

   for (i = 0; i < n; ++i)
      printf("%s\n", w[i]);
}
\end{lstlisting}
\end{affianco}

\subsection{Come funziona}

\begin{itemize}
  \item \code{scanf("\%s", word)} legge una parola (una sequenza di
        caratteri senza spazi) nell'array \code{word}, aggiungendo il
        terminatore. Il ciclo continua finché \code{scanf} riesce a leggere.
  \item Per ogni parola si controlla che non ce ne siano troppe e che non
        sia troppo lunga; poi \code{calloc} riserva esattamente
        \code{strlen(word) + 1} caratteri, e \code{strcpy} ci copia la
        parola. \code{w[i]} punta alla copia: \code{word} verrà riusato per
        la parola successiva.
  \item \code{sort\_words} è un ordinamento per selezione con \code{strcmp}:
        se \code{w[i]} viene dopo \code{w[j]}, si scambiano.
  \item \code{swap} riceve gli \textbf{indirizzi} di due elementi di
        \code{w}, cioè dei \code{char **}, e scambia i due
        \textbf{puntatori}: le stringhe non si spostano, cambia solo
        l'ordine delle frecce (figura~\ref{fig:scambio-puntatori}).
  \item Il file \code{error.c} contiene le tre funzioni di errore, che
        stampano un messaggio e terminano il programma con \code{exit(1)}.
\end{itemize}

\begin{figure}[H]
  \centering
  \begin{tikzpicture}[font=\FiraCodeNL\scriptsize,
      cc/.style={mcella, minimum width=0.45cm, minimum height=0.5cm, inner xsep=0pt}]
    \noindici
    \node[font=\Lato\scriptsize\bfseries] at (2.2,0.75) {prima dello scambio};
    \memoria[minimum width=1.0cm]{a}{(0,0)}{w[0]/{}/, w[1]/{}/}
    \celle[cc]{x0}{(2.6,0)}{p,e,r,a,\textbackslash 0}
    \celle[cc]{x1}{(2.6,-0.52)}{m,e,l,a,\textbackslash 0}
    \riferimento{a-1}{-- (x0-0.west)}
    \riferimento{a-2}{-- (x1-0.west)}
    \begin{scope}[xshift=7.6cm]
      \node[font=\Lato\scriptsize\bfseries] at (2.2,0.75) {dopo \code{swap(\&w[0], \&w[1])}};
      \memoria[minimum width=1.0cm]{b}{(0,0)}{w[0]/{}/, w[1]/{}/}
      \celle[cc]{y0}{(2.6,0)}{p,e,r,a,\textbackslash 0}
      \celle[cc]{y1}{(2.6,-0.52)}{m,e,l,a,\textbackslash 0}
      \riferimento[draw=accent!75!black]{b-1}{to[out=0, in=180] (y1-0.west)}
      \riferimento[draw=accent!75!black]{b-2}{to[out=0, in=180] (y0-0.west)}
    \end{scope}
  \end{tikzpicture}
  \caption{\code{swap} scambia i puntatori, non i caratteri: ora
           \code{w[0]} punta a \code{"mela"}.}
  \label{fig:scambio-puntatori}
\end{figure}

\subsection{Il programma in esecuzione}

Si compilano insieme i cinque file sorgente e si prova con una frase:

\begin{console}
$ gcc -Wall -o sort_words main.c error.c wrt.c swap.c sort_words.c
$ echo "le parole di un Testo da Ordinare" | ./sort_words
Ordinare
Testo
da
di
le
parole
un
\end{console}

\code{echo} scrive la frase sullo standard output, e il simbolo \code{|}
(\textit{pipe}) la collega allo standard input di \code{sort\_words}: è un
altro modo di fornire l'input, oltre alla tastiera e alla redirezione
\code{<}. Le parole con la maiuscola vengono prima di tutte le altre, perché
in ASCII le maiuscole precedono le minuscole.

Nella cartella degli esempi c'è anche un file \code{input.txt} con una frase
di 50 parole (una citazione sugli array a lunghezza variabile, che vedremo
alla fine del capitolo); \code{./sort\_words < input.txt} le stampa
ordinate, da \code{According} e \code{Arrays} fino a \code{world}.

\section{Gli array a più dimensioni}

\subsection{Array di array}

In C si possono definire array di array:

\begin{lstlisting}[style=c]
int a[100];           /* array a una dimensione */
int b[2][7];          /* array a due dimensioni */
int c[5][3][2];       /* array a tre dimensioni */
\end{lstlisting}

Tutti gli elementi di un array sono memorizzati in memoria
\textbf{contigua}, a partire dal primo: \code{a[0]}, \code{b[0][0]},
\code{c[0][0][0]}. Un array a due dimensioni si può pensare come una
\textbf{matrice}: con \code{int a[3][5];}, 3 è il numero di righe e 5 il
numero di colonne. In memoria, però, la matrice è “srotolata”: prima tutta
la riga 0, poi la riga 1, e così via (figura~\ref{fig:matrice}).

\begin{affianco}[0.55][c]
\begin{lstlisting}[style=c, name=matrice.c]
#include <stdio.h>

int main(void)
{
   int a[3][5], i, j;

   for (i = 0; i < 3; i++)
      for (j = 0; j < 5; j++)
         a[i][j] = 10 * i + j;

   for (i = 0; i < 3; i++) {
      for (j = 0; j < 5; j++)
         printf("%4d", a[i][j]);
      printf("\n");
   }
   printf("&a[0][0] = %p\n", (void *) &a[0][0]);
   printf("&a[0][4] = %p\n", (void *) &a[0][4]);
   printf("&a[1][0] = %p\n", (void *) &a[1][0]);
   i = 2; j = 3;
   printf("%d %d %d %d %d\n", a[i][j], *(a[i] + j), (*(a + i))[j],
          *((*(a + i)) + j), *(&a[0][0] + 5 * i + j));
   return 0;
}
\end{lstlisting}
\risultato
\begin{console}
$ ./matrice
   0   1   2   3   4
  10  11  12  13  14
  20  21  22  23  24
&a[0][0] = 0x7ffff4d0aaa0
&a[0][4] = 0x7ffff4d0aab0
&a[1][0] = 0x7ffff4d0aab4
23 23 23 23 23
\end{console}
\wtnota{\code{a[1][0]} sta 4 byte dopo \code{a[0][4]}: subito dopo
l'ultimo elemento della riga 0.}
\end{affianco}

\begin{figure}[H]
  \centering
  \begin{tikzpicture}[font=\FiraCodeNL\scriptsize,
      r0/.style={mcella, minimum width=0.62cm},
      r1/.style={mcella, minimum width=0.62cm, fill=accent!10!white},
      r2/.style={mcella, minimum width=0.62cm, fill=green-gradient-1!40!white}]
    \noindici
    \node[font=\Lato\scriptsize\bfseries] at (1.6,0.75) {come la pensiamo};
    \celle[r0]{m0}{(0,0)}{0,1,2,3,4}
    \celle[r1]{m1}{(0,-0.5)}{10,11,12,13,14}
    \celle[r2]{m2}{(0,-1.0)}{20,21,22,23,24}
    \foreach \k in {0,1,2}
      \node[mindir, anchor=east] at ([xshift=-2pt]m\k-0.west) {riga \k};
    \node[font=\Lato\scriptsize\bfseries] at (10.0,0.75) {come sta in memoria};
    \celle[r0]{l0}{(5.0,-0.5)}{0,1,2,3,4}
    \celle[r1]{l1}{(l0-4.east)}{10,11,12,13,14}
    \celle[r2]{l2}{(l1-4.east)}{20,21,22,23,24}
    \node[mnota, anchor=north, align=center] at ([yshift=-0.2cm]l1-2.south) {\code{a[i][j]} sta a $5 \cdot i + j$ posti da \code{a[0][0]}};
  \end{tikzpicture}
  \caption{La matrice \code{int a[3][5]}: righe una dopo l'altra in memoria.}
  \label{fig:matrice}
\end{figure}

\subsection{La storage mapping function}

La funzione che, per ogni array, associa agli indici la posizione
dell'elemento si chiama \term{storage mapping function}. Per un array a una
dimensione è \code{a[i]} $\rightarrow$ \code{*(\&a[0] + i)}. Per un array a
due dimensioni, con $N$ colonne,

\[
  \code{a[i][j]} \;\rightarrow\; \code{*(\&a[0][0] + } N \cdot i + j\code{)}
\]

perché per arrivare alla riga $i$ bisogna saltare $i$ righe intere da $N$
elementi. Il numero di righe (la prima dimensione) \textbf{non viene usato}.
Per lo stesso motivo \code{a[i][j]} si può scrivere anche
\code{*(a[i] + j)}, \code{(*(a + i))[j]} o \code{*((*(a + i)) + j)}: l'ultima
riga di \code{matrice.c} li stampa tutti, e valgono tutti 23.

Lo stesso vale con più dimensioni. Con \code{int A[7][9][2];} il compilatore
riserva spazio contiguo per $7 \cdot 9 \cdot 2 = 126$ interi, l'indirizzo
base è \code{\&A[0][0][0]}, e

\[
  \code{A[i][j][k]} \;\rightarrow\; \code{*(\&A[0][0][0] + } 9 \cdot 2 \cdot i + 2 \cdot j + k\code{)}
\]

Anche qui la prima dimensione, 7, non serve.

\subsection{L'inizializzazione}

Un array a più dimensioni si può inizializzare in molti modi. Queste
inizializzazioni sono equivalenti:

\begin{lstlisting}[style=c]
int A[2][3] = {1, 2, 3, 4, 5, 6};
int A[2][3] = {{1, 2, 3}, {4, 5, 6}};
int A[][3] = {{1, 2, 3}, {4, 5, 6}};
int A[][3] = {1, 2, 3, 4, 5, 6};
\end{lstlisting}

\begin{itemize}
  \item Senza graffe interne, gli elementi vengono riempiti per riga:
        \code{A[0][0]}, \code{A[0][1]}, \code{A[0][2]}, \code{A[1][0]},
        \dots{} (\code{gcc -Wall} suggerisce comunque le graffe:
        \textit{missing braces around initializer}).
  \item Se i valori sono meno degli elementi, i rimanenti valgono 0:
        \code{\{\{1\}, \{4, 5\}\}} dà le righe \code{1 0 0} e \code{4 5 0}.
  \item Se la prima dimensione è vuota, il compilatore la ricava dal numero
        di graffe interne (o dei valori); tutte le altre dimensioni vanno
        indicate.
  \item \code{int A[2][2][3] = \{0\};} azzera tutto l'array; qui vanno
        dichiarate tutte le dimensioni. È il modo più semplice di inizializzare
        un array automatico, che altrimenti conterrebbe valori a caso.
\end{itemize}

Anche queste due inizializzazioni di un array a tre dimensioni sono
equivalenti (verificato stampando gli elementi: \code{1 1 0 2 0 0 3 0 0 4 4
0}):

\begin{lstlisting}[style=c]
int A[2][2][3] = {
    {{1, 1, 0}, {2, 0, 0}},
    {{3, 0, 0}, {4, 4, 0}}
};
int A[][2][3] = {{{1, 1}, {2}}, {{3}, {4, 4}}};
\end{lstlisting}

\subsection{Le matrici come parametri}

Una matrice si può passare a una funzione, ma per generare il codice il
compilatore ha bisogno della storage mapping function: nella dichiarazione
del parametro vanno quindi indicate \textbf{tutte le dimensioni tranne la
prima}. Queste tre intestazioni sono equivalenti:

\begin{lstlisting}[style=c]
int sum(int a[3][5])
int sum(int a[][5])
int sum(int (*a)[5])        /* a punta ad array di 5 int */
\end{lstlisting}

ma, se il corpo usa \code{i < 3} e \code{j < 5}, vanno bene solo per le
matrici $3 \times 5$. Passando il numero di righe come parametro, la funzione
va bene per tutte le matrici $m \times 5$:

\begin{affianco}[0.5][c]
\begin{lstlisting}[style=c, name=sum2d.c]
int sum(int a[][5], int m)
{
   int   i, j, sum = 0;

   for (i = 0; i < m; ++i)
      for (j = 0; j < 5; ++j)
         sum += a[i][j];
   return sum;
}
\end{lstlisting}
\risultato
Il numero di colonne, però, resta fisso. Come scrivere una funzione per
matrici di qualunque dimensione? La prima idea, \code{int **a}, è sbagliata,
come mostra la prossima sezione.
\end{affianco}

\subsection{\code{int **} non è una matrice}

Si potrebbe pensare di generalizzare quanto succede per gli array a una
dimensione, e dichiarare il parametro come puntatore a puntatori:

\begin{affianco}[0.5][c]
\begin{lstlisting}[style=c, name=sumpp.c]
#include <stdio.h>

int sum(int **a, int m, int n)
{
   int   i, j, sum = 0;

   for (i = 0; i < m; ++i)
      for (j = 0; j < n; ++j)
         sum += a[i][j];
   return sum;
}

int main()
{
   int A[3][5] = {{1, 2, 3, 4, 5}}, s;

   s = sum(A, 3, 5);
   printf("%d\n", s);
   return 0;
}
\end{lstlisting}
\risultato
\begin{console}
$ ./sumpp
Segmentation fault
\end{console}
\end{affianco}

\begin{console}
$ gcc -Wall -o sumpp sumpp.c
sumpp.c: In function (*@‘@*)main(*@’@*):
sumpp.c:17:12: warning: passing argument 1 of (*@‘@*)sum(*@’@*) from incompatible pointer type [-Wincompatible-pointer-types]
   17 |    s = sum(A, 3, 5);
      |            ^
      |            |
      |            int (*)[5]
sumpp.c:3:15: note: expected (*@‘@*)int **(*@’@*) but argument is of type (*@‘@*)int (*)[5](*@’@*)
\end{console}

Il compilatore avvisa, e il programma si blocca. Il motivo è che \code{A},
dichiarato \code{int A[3][5]}, è un puntatore ad \textbf{array} di 5 interi
(\code{int (*)[5]}), mentre un \code{int **} è un puntatore a
\textbf{puntatori}. Dentro \code{sum}, \code{a[i]} viene letto come un
indirizzo: ma nella memoria di \code{A} non ci sono indirizzi, ci sono gli
interi 1, 2, 3\dots, e usarli come indirizzi porta in zone di memoria
proibite (figura~\ref{fig:non-matrice}). Con un \code{int **}, invece,
\code{a[i]} è un puntatore che si può anche assegnare (\code{a[i] = \dots;}):
le due cose non sono equivalenti.

\begin{figure}[H]
  \centering
  \begin{tikzpicture}[font=\FiraCodeNL\scriptsize,
      cc/.style={mcella, minimum width=0.55cm}]
    \noindici
    \node[font=\Lato\scriptsize\bfseries] at (2.0,0.8) {\code{int A[3][5]}: 15 interi contigui};
    \celle[cc]{a}{(0,0)}{1,2,3,4,5,0,0,\ldots}
    \node[font=\Lato\scriptsize\bfseries] at (9.8,0.8) {\code{int **a}: puntatori a righe};
    \memoria[minimum width=0.9cm]{p}{(7.6,0)}{a[0]/{}/, a[1]/{}/, a[2]/{}/}
    \foreach \r in {1,2,3} {
      \celle[cc]{r\r}{(9.6,-\r*0.52+0.52)}{\mvuota,\mvuota,\mvuota,\mvuota,\mvuota}
      \riferimento{p-\r}{-- (r\r-0.west)}
    }
    \node[mnota, anchor=north, align=center, text=merr] at (2.2,-0.45)
      {se \code{sum} legge queste celle come indirizzi,\\“salta” all'indirizzo 1, 2\dots: memoria proibita};
  \end{tikzpicture}
  \caption{Un array a due dimensioni non contiene puntatori; un
           \code{int **} sì.}
  \label{fig:non-matrice}
\end{figure}

\subsection{Passare l'indirizzo base}

Un modo per scrivere una funzione valida per qualunque dimensione è passarle
l'\textbf{indirizzo base} e le dimensioni, e applicare a mano la storage
mapping function:

\begin{affianco}[0.5][c]
\begin{lstlisting}[style=c]
int sum_base(int *a, int m, int n)
{
   int   i, j, sum = 0;

   for (i = 0; i < m; ++i)
      for (j = 0; j < n; ++j)
         sum += *(a + n*i + j);
   return sum;
}
\end{lstlisting}
\risultato
Si chiama con \code{sum\_base(\&A[0][0], 3, 5)}. Con una matrice
\code{A[3][5]} tutta di 1, sia \code{sum(A, 3)} sia
\code{sum\_base(\&A[0][0], 3, 5)} restituiscono 15 (verificato). Funziona,
ma si perde la comoda scrittura \code{a[i][j]}.
\end{affianco}

\section{Le matrici dinamiche}

\subsection{Un array di righe}

Gli array a più dimensioni, insomma, non sono molto flessibili: non si
possono passare alle funzioni senza perdere in generalità (dimensioni fisse)
o in semplicità (indirizzo base). Sarebbe più comodo poter scrivere funzioni
come \code{int f(int **a, int m, int n);}, a cui si passa la matrice e le
sue dimensioni. Si può, allocando la matrice dinamicamente come un
\textbf{array di puntatori alle righe}:

\begin{affianco}[0.5][c]
\begin{lstlisting}[style=c, name=matrici.c]
double **matrix_alloc(int m, int n)
{
   int i;
   double **p;

   p = malloc(m * sizeof(double *));
   for (i = 0; i < m; i++)
      p[i] = malloc(n * sizeof(double));
   return p;
}

void free_matrix(double **a, int m)
{
   int i;

   for (i = 0; i < m; i++)
      free(a[i]);
   free(a);
}
\end{lstlisting}
\risultato
\begin{itemize}
  \item Prima si alloca l'array di \code{m} puntatori (ognuno occupa
        \code{sizeof(double *)}), poi ogni riga, di \code{n} \code{double}.
  \item All'elemento si accede con la comoda scrittura \code{a[i][j]}:
        \code{a[i]} è il puntatore alla riga $i$, e \code{[j]} sceglie
        l'elemento nella riga.
  \item La memoria della matrice \textbf{non è contigua}: ogni riga sta
        dove \code{malloc} l'ha messa.
  \item Si libera con una \code{free} per ogni riga, e infine quella
        dell'array dei puntatori.
\end{itemize}
\end{affianco}

\subsection{Memoria contigua}

Una variante alloca tutti gli elementi in un'unica zona contigua, e fa
puntare le righe nei punti giusti: \code{a[0]} all'inizio, e ogni riga $n$
elementi dopo la precedente (figura~\ref{fig:matrici-dinamiche}).

\begin{affianco}[0.5][c]
\begin{lstlisting}[style=c]
double **matrix_alloc2(int m, int n)
{
   int i;
   double **a;

   a = malloc(m * sizeof(double *));
   a[0] = malloc(m * n * sizeof(double));
   for (i = 1; i < m; i++)
      a[i] = a[i-1] + n;
   return a;
}

void free_matrix2(double **a)
{
   free(a[0]);
   free(a);
}
\end{lstlisting}
\risultato
\begin{console}
$ ./matrici
a[2][3] = 23   b[2][3] = 23
b: riga 1 inizia 4 double dopo la riga 0
\end{console}
\wtnota{\code{a} è allocata riga per riga, \code{b} in modo contiguo; entrambe
sono $3 \times 4$, con \code{[i][j]} uguale a $10 i + j$. Compilato con
\code{-fsanitize=address}, il programma non dà errori.}
\end{affianco}

Nelle slide la seconda funzione di liberazione è dichiarata
\code{double free\_matrix(double **a)}: non restituisce nulla, quindi il
tipo giusto è \code{void} (con \code{double}, \code{gcc -Wall} avvisa
\textit{control reaches end of non-void function}).

\begin{figure}[H]
  \centering
  \begin{tikzpicture}[font=\FiraCodeNL\scriptsize,
      cc/.style={mcella, minimum width=0.55cm}]
    \noindici
    \node[font=\Lato\scriptsize\bfseries] at (2.3,0.75) {riga per riga};
    \memoria[minimum width=0.9cm]{p}{(0,0)}{a[0]/{}/, a[1]/{}/, a[2]/{}/}
    \foreach \r/\y in {1/0.25, 2/-0.55, 3/-1.35} {
      \celle[cc]{q\r}{(2.0,\y)}{{},{},{},{}}
      \riferimento{p-\r}{-- (q\r-0.west)}
    }
    \node[mnota, anchor=north] at (2.6,-1.75) {tre zone separate};
    \begin{scope}[xshift=7.6cm]
      \node[font=\Lato\scriptsize\bfseries] at (2.6,0.75) {memoria contigua};
      \memoria[minimum width=0.9cm]{s}{(0,0)}{b[0]/{}/, b[1]/{}/, b[2]/{}/}
      \celle[cc]{z}{(2.0,-1.3)}{{},{},{},{},{},{},{},{},{},{},{},{}}
      \riferimento{s-1}{to[out=0, in=90] (z-0.north)}
      \riferimento{s-2}{to[out=0, in=90] (z-4.north)}
      \riferimento{s-3}{to[out=0, in=90] (z-8.north)}
      \node[mnota, anchor=north] at ([yshift=-0.15cm]z-6.south) {un'unica zona: \code{b[i] = b[i-1] + n}};
    \end{scope}
  \end{tikzpicture}
  \caption{Due modi di allocare una matrice $3 \times 4$.}
  \label{fig:matrici-dinamiche}
\end{figure}

\section{Gli array a lunghezza variabile}

\subsection{Le dimensioni decise durante l'esecuzione}

Il C99 ha introdotto gli \term{array a lunghezza variabile} (\textit{variable
length array}, VLA): array automatici la cui dimensione è un'espressione
calcolata durante l'esecuzione, per esempio un valore letto da tastiera.
Ecco una versione ridotta dell'esempio delle slide (che a sua volta viene
dal libro di Deitel):

\begin{affianco}[0.55][c]
\begin{lstlisting}[style=c, name=vla.c]
#include <stdio.h>

void print2DArray(int row, int col, int array[row][col]);

int main(void)
{
   int row, col;

   printf("Righe e colonne: ");
   scanf("%d %d", &row, &col);

   int m[row][col];                /* un VLA */

   for (int i = 0; i < row; ++i)
      for (int j = 0; j < col; ++j)
         m[i][j] = i + j;
   printf("sizeof(m) = %zu byte\n", sizeof(m));
   print2DArray(row, col, m);
   return 0;
}

void print2DArray(int row, int col, int array[row][col])
{
   for (int i = 0; i < row; ++i) {
      for (int j = 0; j < col; ++j)
         printf("%5d", array[i][j]);
      puts("");
   }
}
\end{lstlisting}
\risultato
\begin{console}
$ ./vla
Righe e colonne: <<2 5>>
sizeof(m) = 40 byte
    0    1    2    3    4
    1    2    3    4    5
$ ./vla
Righe e colonne: <<4 3>>
sizeof(m) = 48 byte
    0    1    2
    1    2    3
    2    3    4
    3    4    5
\end{console}
\end{affianco}

\begin{itemize}
  \item La dichiarazione \code{int m[row][col];} viene dopo la lettura delle
        dimensioni: in C99 le dichiarazioni si possono mettere in mezzo alle
        istruzioni, e così le variabili dei \code{for}
        (\code{for (int i = 0; \dots)}).
  \item \code{sizeof(m)} viene calcolato durante l'esecuzione: $2 \cdot 5
        \cdot 4 = 40$ byte la prima volta, $4 \cdot 3 \cdot 4 = 48$ la
        seconda.
  \item Nel parametro \code{int array[row][col]} le dimensioni sono altri
        parametri, dichiarati prima: la funzione va bene per matrici di
        qualunque dimensione.
\end{itemize}

\subsection{Pro e contro}

I VLA sono comodi, ma hanno un difetto: sono allocati nello
\textbf{stack} (capitolo~\ref{chap:puntatori}), che è piccolo rispetto allo
heap, e la loro dimensione non è nota in compilazione. Con dimensioni grandi
c'è il rischio concreto di esaurire lo stack: con \code{2000 2000} la matrice
occupa 16 MB, il doppio degli 8 MB che Linux assegna di solito allo stack, e
il programma viene fermato con \textit{Segmentation fault}. I VLA sono
quindi utilizzabili per piccole dimensioni; per le grandi si usa
\code{malloc}. È anche il parere di Linus Torvalds, il creatore di Linux,
riportato nel file \code{input.txt} dell'ordinamento delle parole: i VLA
\textit{bring nothing of positive value}, meglio un array fisso
sufficientemente grande o l'allocazione nello heap.

\insintesi{
  Una \textbf{stringa} è un array di caratteri terminato da
  \code{'\textbackslash 0'}, che va contato nella dimensione. Una stringa
  costante è un puntatore a memoria di sola lettura; \code{\%s} stampa fino
  al terminatore. \code{const} impedisce di modificare una variabile o il
  dato puntato, e il compilatore lo verifica. \code{ctype.h} classifica e
  trasforma i caratteri, \code{string.h} copia, concatena, confronta e cerca
  nelle stringhe, senza controllare lo spazio: per una copia servono
  \code{strlen + 1} byte. Un \textbf{array di puntatori} raccoglie stringhe
  di lunghezza diversa, e si ordina scambiando i puntatori. Gli array a
  \textbf{più dimensioni} stanno in memoria riga dopo riga (storage mapping
  function \code{\&a[0][0] + N*i + j}); come parametri vanno dichiarate
  tutte le dimensioni tranne la prima, e un \code{int **} non è una matrice.
  Le matrici dinamiche sono array di puntatori alle righe; i VLA hanno
  dimensioni decise in esecuzione, ma stanno nello stack.}

\section{Riferimenti}

\begin{itemize}
  \item \textbf{A. Kelley, I. Pohl}, \textit{A Book on C}: capitolo~6
        (stringhe, array di puntatori, array a più dimensioni).
  \item \textbf{H. M. Deitel, P. J. Deitel}, \textit{C: How to Program}:
        array a lunghezza variabile.
  \item \textbf{Null-terminated byte strings} --- cppreference.com. \\
        \urlx{https://en.cppreference.com/w/c/string/byte}
\end{itemize}


  \chapter{Le strutture}
\label{chap:strutture}

Gli array raccolgono molti dati \textbf{dello stesso tipo}. Spesso, però,
un'informazione è fatta di parti di tipo diverso: una carta da gioco ha un
valore (un numero) e un seme (un carattere); uno studente ha un nome (una
stringa), una matricola (un intero) e un voto (un carattere). Le
\textbf{strutture} (\code{struct}) permettono di aggregare variabili non
omogenee in un'unica struttura dati, da usare come un tutto. Sono anche il
mattone delle liste collegate del capitolo~\ref{chap:liste}.

\section{Definire una struttura}

\subsection{Il tipo \code{struct carta}}

Una carta da gioco italiana è identificata da un valore numerico, da 1 a 10,
e da un seme: oro, bastoni, coppe o spade. Possiamo dichiarare il tipo

\begin{lstlisting}[style=c]
struct carta {
    int  valore;          /* da 1 a 10 */
    char seme;            /* 'o', 'b', 'c', 's' */
};
\end{lstlisting}

\begin{itemize}
  \item \code{struct} è una parola chiave;
  \item \code{carta} è il nome assegnato alla struttura, detto
        \textbf{tag} (etichetta);
  \item \code{valore} e \code{seme} sono i \textbf{campi} (o membri) della
        struttura, ciascuno con il suo tipo.
\end{itemize}

Questa dichiarazione definisce il tipo di dato derivato \code{struct carta}
(capitolo~\ref{chap:tipi}): è un modello, e \textbf{nessuna memoria} viene
allocata in questo momento. La memoria viene allocata dalla dichiarazione
delle variabili:

\begin{lstlisting}[style=c]
struct carta c1, c2;     /* due variabili di tipo struct carta */
\end{lstlisting}

(Nelle slide, a questo punto, il tipo è chiamato \code{struct card}: è
rimasto il nome dell'esempio originale in inglese, ma il tipo dichiarato è
\code{struct carta}.)

\subsection{L'operatore punto}

Ai campi di una variabile struttura si accede con l'operatore \code{.}
(punto): \code{c1.valore} è il campo \code{valore} di \code{c1}, e si usa
come qualsiasi variabile \code{int}. Le strutture si possono anche
\textbf{assegnare} in blocco: \code{c2 = c1;} copia tutti i campi.

\begin{affianco}[0.5][c]
\begin{lstlisting}[style=c, name=carta.c]
#include <stdio.h>

typedef struct carta {
   int  valore;
   char seme;
} carta;

int main(void)
{
   carta c1, c2;

   c1.valore = 7;
   c1.seme = 'c';
   c2 = c1;                  /* copia tutti i campi */
   c2.valore = 1;
   printf("c1: %d di %c\n", c1.valore, c1.seme);
   printf("c2: %d di %c\n", c2.valore, c2.seme);
   printf("sizeof(carta) = %zu\n", sizeof(carta));
   return 0;
}
\end{lstlisting}
\risultato
\begin{console}
$ ./carta
c1: 7 di c
c2: 1 di c
sizeof(carta) = 8
\end{console}
\wtnota{Dopo la copia \code{c1} e \code{c2} sono indipendenti: cambiare
\code{c2.valore} non tocca \code{c1}.}
\end{affianco}

Il programma usa il \code{typedef} che vedremo fra poco. Ma perché la
struttura occupa 8 byte, se un \code{int} ne occupa 4 e un \code{char} 1? Il
processore accede più facilmente agli \code{int} che stanno a indirizzi
multipli di 4 (su alcuni processori è addirittura obbligatorio), e il
compilatore rispetta questa regola anche dentro le strutture. Pensiamo a un
array di carte: se ogni carta occupasse 5 byte, la seconda inizierebbe al
byte 5 e il suo \code{valore} non starebbe a un multiplo di 4. Per questo il
compilatore aggiunge alla fine della struttura 3 byte di
\textbf{riempimento} (\textit{padding}): ogni carta occupa 8 byte, e ogni
\code{valore} cade al posto giusto (figura~\ref{fig:carta}). Ecco perché la
dimensione di una struttura va sempre chiesta a \code{sizeof}, senza sommare
a mano quelle dei campi.

\begin{figure}[H]
  \centering
  \begin{tikzpicture}[font=\Lato\scriptsize, x=0.78cm, y=1cm]
    \foreach \base/\val/\sem in {0/7/c, 8/1/s} {
      \draw[draw=primary!65!black, fill=primary!3!white, line width=0.7pt]
        (\base,0) rectangle ++(4,0.55);
      \foreach \j in {1,2,3}
        \draw[draw=primary!35!white, densely dashed, line width=0.4pt]
          (\base+\j,0) -- ++(0,0.55);
      \node[font=\FiraCodeNL\footnotesize] at (\base+2,0.275) {\val};
      \draw[draw=accent!65!black, fill=accent!12!white, line width=0.7pt]
        (\base+4,0) rectangle ++(1,0.55);
      \node[font=\FiraCodeNL\scriptsize] at (\base+4.5,0.275) {'\sem'};
      \draw[draw=neutral!45!white, fill=neutral!12!white, line width=0.7pt]
        (\base+5,0) rectangle ++(3,0.55);
    }
    \foreach \k in {0,...,15}
      \node[mindir, anchor=north] at (\k+0.5,-0.02) {\k};
    % che cosa contiene la prima carta
    \draw[wtbrace] (4,0.68) -- (0,0.68);
    \node[mnota, anchor=south, align=center] at (2,0.92) {\code{valore}\\4 byte};
    \draw[draw=neutral!40!black, line width=0.6pt] (4.5,0.6) -- (4.5,0.9);
    \node[mnota, anchor=south, align=center] at (4.5,0.92) {\code{seme}\\1 byte};
    \draw[wtbrace] (8,0.68) -- (5,0.68);
    \node[mnota, anchor=south, align=center] at (6.5,0.92) {riempimento\\3 byte};
    \draw[wtbrace] (12,0.68) -- (8,0.68);
    \node[mnota, anchor=south, align=center, text=accent!70!black] at (10,0.92)
      {\code{mazzo[1].valore}\\inizia al byte 8};
    % le due carte
    \draw[wtbrace] (0,-0.42) -- (8,-0.42);
    \node[mnota, anchor=north] at (4,-0.66) {\code{mazzo[0]}: byte da 0 a 7};
    \draw[wtbrace] (8,-0.42) -- (16,-0.42);
    \node[mnota, anchor=north] at (12,-0.66) {\code{mazzo[1]}: byte da 8 a 15};
  \end{tikzpicture}
  \caption{Un array \code{carta mazzo[2]}: ogni carta occupa 8 byte (4 per
           \code{valore}, 1 per \code{seme} e 3 di riempimento), così il
           \code{valore} di ogni carta inizia a un multiplo di 4.}
  \label{fig:carta}
\end{figure}

\subsection{\code{typedef} e strutture}

Scrivere ogni volta \code{struct carta} è scomodo. Si può usare il
\code{typedef} (capitolo~\ref{chap:tipi}) per dare al tipo un nome più
breve, separatamente o direttamente nella definizione:

\begin{affianco}[0.5][c]
\begin{lstlisting}[style=c]
struct carta {
    int  valore;
    char seme;
};

typedef struct carta carta;
\end{lstlisting}
\risultato
\begin{lstlisting}[style=c]
typedef struct carta {
    int  valore;
    char seme;
} carta;
\end{lstlisting}
\end{affianco}

In entrambi i casi \code{carta} diventa sinonimo di \code{struct carta}, e le
variabili si dichiarano con \code{carta c1, c2;}. Il tag e il nome del
\code{typedef} possono coincidere, perché il C li tiene in elenchi separati.

\subsection{I nomi dei campi}

All'interno di una struttura i campi devono avere nomi diversi; campi di
strutture diverse, invece, possono avere lo stesso nome, perché si
distinguono dalla variabile a cui appartengono:

\begin{affianco}[0.5][c]
\begin{lstlisting}[style=c]
struct fruit {
   char   *name;
   int    calories;
};

struct vegetable {
   char    *name;
   int     calories;
};
\end{lstlisting}
\risultato
\begin{lstlisting}[style=c]
struct fruit      a;
struct vegetable  b;

a.calories = 38;
b.calories = 57;
\end{lstlisting}
\end{affianco}

\subsection{Strutture con e senza tag}

Il tag si può omettere, dichiarando subito le variabili:

\begin{lstlisting}[style=c]
struct {
   int    day, month, year;
   char   day_name[4];
   char   month_name[4];
} yesterday, today, tomorrow;
\end{lstlisting}

Senza tag, però, non si possono dichiarare altre variabili di quel tipo oltre
a \code{yesterday}, \code{today} e \code{tomorrow}: il tipo non ha un nome.
Con un tag, \code{struct date \{ \dots \};}, si possono dichiarare
liberamente variabili \code{struct date}. E con il \code{typedef} il tag non
serve: il nome lo dà il \code{typedef}.

\begin{affianco}[0.55][c]
\begin{lstlisting}[style=c, name=data.c]
#include <stdio.h>

typedef struct {
   int    day, month, year;
   char   day_name[4];
   char   month_name[4];
} date;

int main(void)
{
   date d1 = {31, 10, 2026, "Sat", "Oct"};

   printf("%s %d %s %d\n", d1.day_name, d1.day,
          d1.month_name, d1.year);
   return 0;
}
\end{lstlisting}
\risultato
\begin{console}
$ ./data
Sat 31 Oct 2026
\end{console}
\wtnota{Una struttura si inizializza come un array: i valori fra graffe, nello
stesso ordine dei campi.}
\end{affianco}

\subsection{Che cosa non si può fare}

Le strutture si possono assegnare, passare alle funzioni e restituire, ma non
si possono \textbf{confrontare} con \code{==}, né stampare con un'unica
conversione di \code{printf}. Proviamo a confrontare due carte:

\begin{affianco}[0.5][c]
\begin{lstlisting}[style=c, name=confronta.c]
typedef struct carta {
   int  valore;
   char seme;
} carta;

int main(void)
{
   carta c1 = {7, 'c'}, c2 = {7, 'c'};

   if (c1 == c2)
      return 1;
   return 0;
}
\end{lstlisting}
\risultato
Il compilatore non sa che cosa significhi \code{==} fra due strutture, e
rifiuta il programma (qui sotto). Il confronto va scritto campo per campo:

\vspace{4pt}
\begin{lstlisting}[style=c]
if (c1.valore == c2.valore
        && c1.seme == c2.seme)
   return 1;
\end{lstlisting}
\end{affianco}

\begin{console}
$ gcc -Wall -o confronta confronta.c
confronta.c: In function (*@‘@*)main(*@’@*):
confronta.c:10:11: error: invalid operands to binary == (have (*@‘@*)carta(*@’@*) and (*@‘@*)carta(*@’@*))
   10 |    if (c1 == c2)
      |           ^~
\end{console}

\section{Array di strutture}

\subsection{Un tipo come gli altri}

Una volta definito, un tipo struttura si usa come qualsiasi altro tipo: si
possono dichiarare array di strutture, puntatori a strutture, e allocarle
dinamicamente.

\begin{lstlisting}[style=c]
typedef struct {
   float   re;
   float   im;
} complex;

complex a, b, c[100], *d;
...
d = (complex *) malloc(n * sizeof(complex));
\end{lstlisting}

\code{c[100]} è un array di 100 numeri complessi: \code{c[i].re} è la parte
reale dell'elemento \code{i}. \code{d} punta a un array di \code{n} complessi
allocato nello heap, e anche lì si scrive \code{d[i].re}.

\subsection{Il programma \code{class\_info}}

L'esempio delle slide gestisce le informazioni su una classe. In un header,
incluso da ogni modulo del programma, si definiscono la dimensione massima e
il tipo \code{student}; la funzione \code{fail} conta gli studenti con il
voto \code{'F'} (insufficiente, nel sistema americano):

\begin{affianco}[0.5][t]
\begin{lstlisting}[style=h, name=class_info.h]
#define   CLASS_SIZE    100

typedef struct {
   char     *last_name;
   int      student_id;
   char     grade;
} student;

int fail(student *class, int N);
\end{lstlisting}
\begin{lstlisting}[style=c, name=fail.c]
/* Count the failing grades. */
#include "class_info.h"

int fail(student *class, int N)
{
   int   i, cnt = 0;

   for (i = 0; i < N; ++i)
      cnt += class[i].grade == 'F';
   return cnt;
}
\end{lstlisting}
\risultato
\begin{lstlisting}[style=c, name=main.c]
#include <stdio.h>
#include <stdlib.h>
#include "class_info.h"

int main(void)
{
   int       cnt, i;
   student   tmp, *class;

   class = malloc(CLASS_SIZE * sizeof(student));
   tmp.grade = 'A';
   tmp.last_name = "Casanova";
   tmp.student_id = 910017;
   class[0] = tmp;
   tmp.grade = 'F';
   tmp.last_name = "Superman";
   tmp.student_id = 910018;
   class[1] = tmp;
   tmp.grade = 'C';
   tmp.last_name = "Wonderwoman";
   tmp.student_id = 910019;
   class[2] = tmp;
   /* Enter nothing for the remainder of the class. */
   for (i = 3; i < CLASS_SIZE; ++i) {
      class[i].grade = ' ';
      class[i].last_name = "";
      class[i].student_id = 0;
   }
   cnt = fail(class, CLASS_SIZE);
   printf("Number of failing grades: %d\n", cnt);
   free(class);
   return 0;
}
\end{lstlisting}
\end{affianco}

\begin{affianco}[0.5][c]
\begin{console}
$ gcc -Wall -o class_info main.c fail.c
$ ./class_info
Number of failing grades: 1
\end{console}
\risultato
Nelle slide ci sono due versioni di \code{main}: una con l'array
\code{student class[CLASS\_SIZE]} dichiarato normalmente, e questa, con
l'array allocato dinamicamente. \code{fail} va bene per entrambe, perché
riceve un puntatore (capitolo~\ref{chap:puntatori}).
\end{affianco}

\subsection{Come funziona}

\begin{itemize}
  \item \code{tmp} è una variabile di appoggio: se ne riempiono i campi e
        la si copia in \code{class[0]}, poi la si riempie di nuovo e la si
        copia in \code{class[1]}, e così via (figura~\ref{fig:classe}).
  \item La copia \code{class[0] = tmp} copia \textbf{i campi}: per
        \code{last\_name}, che è un puntatore, copia l'indirizzo, non la
        stringa. Qui va bene, perché ogni studente punta a una stringa
        costante diversa.
  \item Gli altri 97 studenti hanno voto \code{' '}, cognome \code{""} (la
        stringa vuota) e matricola 0.
  \item In \code{fail}, \code{class[i].grade == 'F'} vale 1 o 0, e
        \code{cnt += \dots} lo somma: si contano i voti \code{'F'} senza
        bisogno di un \code{if}.
\end{itemize}

\begin{figure}[H]
  \centering
  \begin{tikzpicture}[font=\Lato\scriptsize]
    \mchiave{2.0cm}
    \oggetto[anchor=north west]{s0}{(0,0)}{3.3cm}{class[0]}{last\_name/{}, student\_id/910017, grade/{\code{'A'}}}
    \oggetto[anchor=north west]{s1}{([yshift=0.7pt]s0-box.south west)}{3.3cm}{class[1]}{last\_name/{}, student\_id/910018, grade/{\code{'F'}}}
    \oggetto[anchor=north west]{s2}{([yshift=0.7pt]s1-box.south west)}{3.3cm}{class[2]}{last\_name/{}, student\_id/910019, grade/{\code{'C'}}}
    \node[mnota, anchor=north, align=center] at ([yshift=-0.05cm]s2-box.south) {$\vdots$\\fino a \code{class[99]}};
    % le stringhe a cui puntano i campi last_name
    \node[mvar, anchor=west, minimum width=2.5cm] (t0) at ([xshift=1.3cm]s0-1.east) {"Casanova"};
    \node[mvar, anchor=west, minimum width=2.5cm] (t1) at ([xshift=1.3cm]s1-1.east) {"Superman"};
    \node[mvar, anchor=west, minimum width=2.5cm] (t2) at ([xshift=1.3cm]s2-1.east) {"Wonderwoman"};
    \foreach \ist in {0,1,2} {
      \draw[mrif] ([xshift=-0.45cm]s\ist-1.east) -- (t\ist.west);
      \node[mpunto] at ([xshift=-0.45cm]s\ist-1.east) {};
    }
    \node[mnota, anchor=south] at ([yshift=2pt]t0.north) {stringhe costanti};
    % il puntatore class
    \node[mvar, minimum width=0.7cm] (pc) at ([xshift=-1.3cm]s0.west) {};
    \node[mnome, anchor=east] at ([xshift=-2pt]pc.west) {class};
    \riferimento{pc}{-- (s0.west)}
  \end{tikzpicture}
  \caption{L'array \code{class} nello heap: 100 strutture \code{student}
           consecutive, di 16 byte ciascuna, e le prime tre contengono i dati
           degli studenti. Il campo \code{last\_name} contiene solo
           l'indirizzo della stringa.}
  \label{fig:classe}
\end{figure}

\section{L'operatore freccia}

\subsection{\code{p->campo}}

Quando si ha un \textbf{puntatore} a una struttura, per accedere a un campo
si dovrebbe scrivere \code{(*p).grade}: prima si dereferenzia il puntatore,
poi si sceglie il campo. Le parentesi sono necessarie, perché \code{.} ha la
precedenza su \code{*}: \code{*p.grade} vorrebbe dire \code{*(p.grade)}. Per
questo il C ha un operatore apposito, la \textbf{freccia} \code{->}:

\[
  \code{p->grade} \quad\text{equivale a}\quad \code{(*p).grade}
\]

Le slide riscrivono \code{fail} in altri due modi, entrambi con la freccia:

\begin{affianco}[0.5][c]
\begin{lstlisting}[style=c, name=fail2.c]
int fail(student *class, int N)
{
   int   i, cnt = 0;

   for (i = 0; i < N; ++i)
      cnt += (class + i)->grade == 'F';
   return cnt;
}
\end{lstlisting}
\risultato
\begin{lstlisting}[style=c, name=fail3.c]
int fail(student *class, int N)
{
   int       i, cnt = 0;
   student  *tmp = class;

   for (i = 0; i < N; ++i, tmp++)
      cnt += tmp->grade == 'F';
   return cnt;
}
\end{lstlisting}
\end{affianco}

Nella prima \code{class + i} è l'indirizzo dell'elemento \code{i} (per
l'aritmetica dei puntatori avanza di una \textbf{struttura} intera alla
volta), e \code{->grade} ne prende il voto. Nella seconda un puntatore
\code{tmp} scorre l'array: \code{tmp++} lo porta sullo studente successivo.
Tutte e tre le versioni danno lo stesso risultato (verificato).

\subsection{Espressioni con i puntatori a strutture}

Le slide riportano una tabella del libro di Kelley e Pohl, nel quale il tipo
è dichiarato con un tag, \code{struct student \{ \dots \};}, e ha gli stessi
campi di \code{class\_info.h}. Con le dichiarazioni e le assegnazioni

\begin{lstlisting}[style=c]
struct student  tmp, *p = &tmp;

tmp.grade = 'A';
tmp.last_name = "Casanova";
tmp.student_id = 910017;
\end{lstlisting}

valgono queste equivalenze (figura~\ref{fig:freccia}):

\begin{center}
\renewcommand{\arraystretch}{1.15}
\begin{tabular}{@{}>{\FiraCodeNL}l>{\FiraCodeNL}l l@{}}
  \toprule
  \normalfont\textbf{espressione} & \normalfont\textbf{equivalente} & \textbf{valore} \\
  \midrule
  tmp.grade               & p->grade               & \code{A} \\
  tmp.last\_name          & p->last\_name          & \code{Casanova} \\
  (*p).student\_id        & p->student\_id         & 910017 \\
  *p->last\_name + 1      & (*(p->last\_name)) + 1 & \code{D} \\
  *(p->last\_name + 2)    & (p->last\_name)[2]     & \code{s} \\
  \bottomrule
\end{tabular}
\end{center}

Le ultime due sono le più istruttive. In \code{*p->last\_name + 1} la freccia
lega più di tutti: si prende \code{p->last\_name}, la stringa; \code{*} ne
dà il primo carattere, \code{'C'}; e \code{+ 1} dà il carattere successivo,
\code{'D'}. In \code{*(p->last\_name + 2)}, invece, si avanza di due caratteri
nella stringa e si prende quello che c'è: la \code{s} di Ca\textbf{s}anova.

\begin{figure}[H]
  \centering
  \begin{tikzpicture}[font=\Lato\scriptsize]
    \mchiave{2.3cm}
    \oggetto{t}{(3.0,0)}{3.6cm}{tmp}{last\_name/{}, student\_id/910017, grade/{\code{'A'}}}
    \node[mnome, anchor=east] at (-0.15,-0.6) {p};
    \node[mvar, minimum width=0.7cm] (pp) at (0.3,-0.6) {};
    \riferimento{pp}{-- (t-box.west)}
    \noindici
    \celle[mcella, minimum width=0.45cm, inner xsep=0pt]{nm}{(7.0,-0.1)}{C,a,s,a,n,o,v,a,\textbackslash 0}
    \draw[mrif] ([xshift=-0.45cm]t-1.east) to[out=0, in=180] (nm-0.west);
    \node[mpunto] at ([xshift=-0.45cm]t-1.east) {};
    \draw[wtarrow] ([yshift=-0.6cm]nm-0.south) -- ([yshift=-1pt]nm-0.south);
    \node[mnota, anchor=north west] at ([xshift=-0.3cm, yshift=-0.6cm]nm-0.south)
      {\code{*p->last\_name} è la \code{C}; \code{+ 1} dà la \code{D}};
    \draw[wtarrowacc] ([yshift=0.55cm]nm-2.north) -- ([yshift=1pt]nm-2.north);
    \node[mnota, anchor=south, text=accent!70!black] at ([yshift=0.55cm]nm-2.north)
      {\code{p->last\_name + 2}: la \code{s}};
  \end{tikzpicture}
  \caption{\code{p} punta alla struttura \code{tmp}; il campo
           \code{last\_name} punta a sua volta alla stringa.}
  \label{fig:freccia}
\end{figure}

\subsection{Una verifica}

Il programma seguente dichiara il tipo con il tag, come nel libro, e stampa i
valori delle espressioni della tabella:

\begin{affianco}[0.6][c]
\begin{lstlisting}[style=c, name=puntatori.c]
#include <stdio.h>

struct student {
   char *last_name;
   int   student_id;
   char  grade;
};

int main(void)
{
   struct student tmp, *p = &tmp;

   tmp.grade = 'A';
   tmp.last_name = "Casanova";
   tmp.student_id = 910017;
   printf("%c\n", p -> grade);
   printf("%s\n", p -> last_name);
   printf("%d\n", (*p).student_id);
   printf("%c\n", *p -> last_name + 1);
   printf("%c\n", *(p -> last_name + 2));
   return 0;
}
\end{lstlisting}
\risultato
\begin{console}
$ ./puntatori
A
Casanova
910017
D
s
\end{console}
\wtnota{Come nella tabella delle slide, attorno a \code{->} ci sono degli
spazi: sono permessi, perché \code{->} è un unico token, ma di solito si
scrive \code{p->grade}.}
\end{affianco}

\code{*p -> last\_name + 1} è un numero, il codice del carattere
\code{'C'} più 1, cioè quello di \code{'D'}: per vederlo come carattere si
stampa con \code{\%c}.

\section{Strutture di strutture}

\subsection{Campi che sono strutture}

Un campo può essere a sua volta una struttura, o un puntatore a una
struttura. L'esempio delle slide descrive un dipendente:

\begin{lstlisting}[style=c]
typedef struct {
   char           name[20];
   int            number;
} department;

typedef struct {
   char           street_name[30];
   int            no;
   char           postcode[10];
} address;

typedef struct {
   char           name[25];
   int            employee_id;
   department     dept;          /* una struttura dentro la struttura */
   address        *home;         /* un puntatore a un'altra struttura */
   double         salary;
} employee_data;
\end{lstlisting}

(Nelle slide \code{name} di \code{department} e \code{street\_name} di
\code{address} sono dichiarati \code{char}, cioè un solo carattere: per
contenere un nome servono degli array, come qui.) Il reparto è contenuto
nella struttura del dipendente; l'indirizzo di casa, invece, sta altrove, e
il dipendente ne contiene solo il puntatore (figura~\ref{fig:annidate}).
Per arrivare a un campo interno si concatenano gli operatori, da sinistra a
destra: \code{e.dept.number} equivale a \code{(e.dept).number}, e con un
puntatore \code{p->dept.number} equivale a \code{(p->dept).number};
\code{e.home->no} è il numero civico.

\begin{figure}[H]
  \centering
  \begin{tikzpicture}[font=\Lato\scriptsize,
      riga/.style={mriga, minimum width=5.9cm,
        text width=\dimexpr5.9cm-6pt\relax, anchor=north}]
    % la struttura e, con la struttura dept al suo interno
    \node[mtesta, minimum width=5.9cm, anchor=north west] (e) at (0,0) {employee\_data e};
    \node[riga] (e-1) at ([yshift=0.7pt]e.south)
      {\makebox[2.1cm][l]{\color{primary!55!black}name}\code{"Mario Rossi"}};
    \node[riga] (e-2) at ([yshift=0.7pt]e-1.south)
      {\makebox[2.1cm][l]{\color{primary!55!black}employee\_id}1234};
    \node[riga, minimum height=1.85cm] (e-3) at ([yshift=0.7pt]e-2.south) {};
    \node[anchor=north west, font=\FiraCodeNL\scriptsize, text=primary!55!black,
          inner sep=3pt] at ([yshift=-2pt]e-3.north west) {dept};
    \mchiave{1.2cm}
    \oggetto[anchor=north west, fill=accent!12!white]{d}{([xshift=2.2cm, yshift=-0.2cm]e-3.north west)}{3.2cm}{department}{name/{\code{"Vendite"}}, number/1}
    \node[riga] (e-4) at ([yshift=0.7pt]e-3.south)
      {\makebox[2.1cm][l]{\color{primary!55!black}home}};
    \node[riga] (e-5) at ([yshift=0.7pt]e-4.south)
      {\makebox[2.1cm][l]{\color{primary!55!black}salary}1850.0};
    % la struttura address, altrove in memoria
    \mchiave{2.0cm}
    \oggetto[anchor=west]{a}{([xshift=1.6cm]e-4.east)}{4.1cm}{address}{street\_name/{\code{"Via Roma"}}, no/12, postcode/{\code{"80100"}}}
    \draw[mrif] ([xshift=-0.5cm]e-4.east) -- (a.west);
    \node[mpunto] at ([xshift=-0.5cm]e-4.east) {};
    % come si arriva a due campi interni
    \node[mnota, anchor=west, text=accent!70!black] (l1) at ([xshift=0.6cm]e-3.east |- d-2) {\code{e.dept.number}};
    \draw[wtarrowacc] (l1.west) -- (d-2.east);
    \node[mnota, anchor=west, text=accent!70!black] (l2) at ([xshift=0.6cm]a-2.east) {\code{e.home->no}};
    \draw[wtarrowacc] (l2.west) -- (a-2.east);
  \end{tikzpicture}
  \caption{\code{dept} è una struttura dentro \code{e}; \code{home} punta a una
           struttura che sta altrove. Le etichette mostrano le espressioni che
           arrivano a due campi interni.}
  \label{fig:annidate}
\end{figure}

\section{Strutture e funzioni}

\subsection{Due modi di passare una struttura}

Supponiamo di voler scrivere una funzione che aggiorna le informazioni di un
dipendente, per esempio il numero del reparto. Si può procedere in due
maniere:

\begin{affianco}[0.5][c]
\begin{lstlisting}[style=c]
employee_data update(employee_data e)
{
   ...
   e.dept.number = n;
   ...
   return e;
}

/* uso: */
e = update(e);
\end{lstlisting}
\risultato
\begin{lstlisting}[style=c]
void update(employee_data *p)
{
   ...
   p->dept.number = n;
   ...
}

/* uso: */
update(&e);
\end{lstlisting}
\end{affianco}

\begin{itemize}
  \item Nel primo metodo \code{e} è passata \textbf{per valore}: la funzione
        riceve una copia di tutta la struttura, la modifica e la restituisce,
        e il chiamante deve assegnare il risultato. Se la struttura è grande,
        creare la copia (due volte: all'andata e al ritorno) può essere
        oneroso.
  \item Nel secondo metodo si passa l'\textbf{indirizzo} di \code{e}: non si
        crea nessuna copia, e la funzione modifica direttamente la struttura
        del chiamante. È il metodo in generale più efficiente.
\end{itemize}

\subsection{Una prova}

Il programma seguente completa le due funzioni (con il nuovo numero di
reparto come parametro \code{n}) e mostra la differenza. In cima al file ci
sono i tre \code{typedef} della sezione precedente, che qui non sono
ripetuti:

\begin{affianco}[0.58][c]
\codicepiccolo
\begin{lstlisting}[style=c, name=impiegati.c]
#include <stdio.h>

/* ... i typedef di department, address ed employee_data ... */

employee_data update(employee_data e, int n)    /* per valore */
{
   e.dept.number = n;
   return e;
}

void update_p(employee_data *p, int n)          /* per indirizzo */
{
   p->dept.number = n;
}

int main(void)
{
   employee_data e = {"Mario Rossi", 1234, {"Vendite", 1}, NULL, 1850.0};

   update(e, 2);                    /* il risultato non viene usato */
   printf("dopo update(e, 2):       reparto %d\n", e.dept.number);
   e = update(e, 3);
   printf("dopo e = update(e, 3):   reparto %d\n", e.dept.number);
   update_p(&e, 4);
   printf("dopo update_p(&e, 4):    reparto %d\n", e.dept.number);
   printf("sizeof(employee_data) = %zu byte\n", sizeof(employee_data));
   return 0;
}
\end{lstlisting}
\risultato
\begin{console}
$ ./impiegati
dopo update(e, 2):       reparto 1
dopo e = update(e, 3):   reparto 3
dopo update_p(&e, 4):    reparto 4
sizeof(employee_data) = 72 byte
\end{console}
\end{affianco}

\begin{itemize}
  \item \code{update(e, 2)} da sola non cambia nulla: modifica la copia, che
        poi viene buttata via perché il risultato non è assegnato.
  \item \code{e = update(e, 3)} funziona, ma copia 72 byte all'andata e 72
        al ritorno.
  \item \code{update\_p(\&e, 4)} copia solo un indirizzo, 8 byte, e modifica
        \code{e} direttamente (figura~\ref{fig:struct-funzioni}).
  \item La struttura si inizializza con le graffe, anche quella interna
        (\code{\{"Vendite", 1\}}); \code{home} vale \code{NULL}, perché in
        questa prova l'indirizzo non c'è.
\end{itemize}

\subsection{Copia o indirizzo}

La figura~\ref{fig:struct-funzioni} mette a confronto i due modi. Per valore,
\code{update} lavora su una copia di tutti i 72 byte della struttura, e il
risultato va ricopiato in \code{e}; per indirizzo, \code{update\_p} riceve
soltanto gli 8 byte di un puntatore e, attraverso di esso, modifica proprio
la \code{e} di \code{main}.

\begin{figure}[H]
  \centering
  \begin{tikzpicture}[font=\Lato\scriptsize]
    \node[font=\Lato\scriptsize\bfseries] at (2.2,1.1) {\code{e = update(e, 3)}: per valore};
    \node[wtnodeplain, minimum width=2.0cm, minimum height=1.0cm] (m1) at (0.9,0) {\code{e}\\72 byte};
    \node[wtnodeplain, minimum width=2.0cm, minimum height=1.0cm, fill=accent!10!white] (f1) at (3.9,0) {copia di \code{e}\\72 byte};
    \draw[mcerca] ([yshift=0.15cm]m1.east) -- node[above, mnota] {copia} ([yshift=0.15cm]f1.west);
    \draw[mcerca] ([yshift=-0.15cm]f1.west) -- node[below, mnota] {copia} ([yshift=-0.15cm]m1.east);
    \node[mnota] at (0.9,-0.8) {\code{main}};
    \node[mnota] at (3.9,-0.8) {\code{update}};
    \begin{scope}[xshift=7.6cm]
      \node[font=\Lato\scriptsize\bfseries] at (2.2,1.1) {\code{update\_p(\&e, 4)}: per indirizzo};
      \node[wtnodeplain, minimum width=2.0cm, minimum height=1.0cm] (m2) at (0.9,0) {\code{e}\\72 byte};
      \node[mvar, fill=accent!10!white] (p2) at (3.9,0) {};
      \node[mnome, anchor=west] at ([xshift=2pt]p2.east) {p};
      \riferimento{p2}{-- (m2.east)}
      \node[mnota] at (0.9,-0.8) {\code{main}};
      \node[mnota] at (3.9,-0.8) {\code{update\_p}: 8 byte};
    \end{scope}
  \end{tikzpicture}
  \caption{Passare una struttura per valore copia tutti i suoi byte;
           passarla per indirizzo copia solo un puntatore.}
  \label{fig:struct-funzioni}
\end{figure}

Quando una funzione deve solo \textbf{leggere} una struttura grande, si passa
comunque l'indirizzo, ma dichiarando il parametro \code{const}:

\begin{lstlisting}[style=c]
void stampa(const employee_data *p);    /* legge *p, ma non può modificarlo */
\end{lstlisting}

Così si evita la copia, e il compilatore garantisce che la funzione non
modifichi la struttura (capitolo~\ref{chap:stringhe}).

\insintesi{
  Una \textbf{struttura} aggrega campi di tipo diverso: la dichiarazione
  \code{struct tag \{ \dots \};} definisce un tipo (senza allocare memoria),
  le variabili si dichiarano con \code{struct tag v;} o, con un
  \code{typedef}, con il solo nome del tipo. Ai campi si accede con il
  \textbf{punto} (\code{v.campo}) e, attraverso un puntatore, con la
  \textbf{freccia} (\code{p->campo}, cioè \code{(*p).campo}). Le strutture si
  inizializzano con le graffe, si assegnano in blocco (copiando i campi, e
  per i puntatori solo l'indirizzo) ma non si confrontano con \code{==}; la
  loro dimensione, che può includere byte di riempimento, va chiesta a
  \code{sizeof}. Si possono fare array di strutture, anche dinamici, e
  strutture che contengono strutture o puntatori a strutture. Alle funzioni
  conviene passare l'indirizzo di una struttura, invece di una copia.}

\section{Riferimenti}

\begin{itemize}
  \item \textbf{A. Kelley, I. Pohl}, \textit{A Book on C}: capitolo~9
        (strutture e unioni).
  \item \textbf{Struct declaration} --- cppreference.com. \\
        \urlx{https://en.cppreference.com/w/c/language/struct}
\end{itemize}


  \chapter{Le liste collegate}
\label{chap:liste}

Gli array hanno una dimensione fissa, e i loro elementi stanno uno accanto
all'altro: per far posto a un elemento nel mezzo bisogna spostare tutti
quelli che seguono. Le \textbf{liste collegate} superano questi limiti con
un'idea semplice: ogni elemento sa dove si trova il successivo. In questo
capitolo mettiamo insieme ciò che abbiamo visto finora (strutture, puntatori,
\code{malloc} e \code{free}) per costruire una lista e scrivere le operazioni
elencate nelle slide; poi introduciamo la \textbf{ricorsione} e i
\textbf{file}, e con questi strumenti leggiamo per intero il progetto delle
lezioni incontrato nel capitolo~\ref{chap:make}.

\section{Perché le liste}

\subsection{I limiti degli array}

Un array ha due limiti. Il primo: la sua dimensione si sceglie quando lo si
crea; se poi i dati crescono, bisogna allocare un array più grande e
ricopiarli. Il secondo: gli elementi sono \textbf{consecutivi}, quindi per
inserire un elemento in mezzo bisogna spostare di un posto tutti quelli che
seguono, e per toglierne uno bisogna spostarli indietro
(figura~\ref{fig:array-lista}, in alto). In un array di un milione di
elementi, inserirne uno all'inizio richiede un milione di spostamenti.

\subsection{L'idea: ogni elemento sa dov'è il successivo}

Una \textbf{lista collegata} (\textit{linked list}) è una struttura dati
\textbf{dinamica} in cui ogni elemento, detto \textbf{nodo}, è collegato al
successivo. Ogni nodo contiene

\begin{itemize}
  \item i \textbf{dati};
  \item un \textbf{puntatore al nodo successivo}.
\end{itemize}

I nodi possono stare ovunque in memoria: l'ordine della lista non dipende da
dove si trovano, ma dai puntatori. Per inserire un elemento basta creare un
nodo e cambiare due puntatori, senza spostare niente
(figura~\ref{fig:array-lista}, in basso).

\begin{figure}[H]
  \centering
  \begin{tikzpicture}[font=\Lato\scriptsize]
    % l'array
    \node[mnota, anchor=west, font=\Lato\scriptsize\bfseries] at (-1.0,1.25)
      {array: per inserire 30 in posizione 2 si spostano 50 e 40};
    \celle{a}{(0,0)}{10,20,40,50,{},{}}
    \draw[mcerca] (a-3.north) to[bend left=70] node[mpasso, above=1pt] {1} (a-4.north);
    \draw[mcerca] (a-2.north) to[bend left=70] node[mpasso, above=1pt] {2} (a-3.north);
    \node[mvar, minimum width=0.6cm, fill=accent!10!white] (x) at (-0.6,0.75) {30};
    \draw[mcerca] (x.east) to[bend left=20] node[mpasso, above=1pt] {3} (a-2.north west);
    \node[mnota, anchor=west] at ([xshift=0.25cm]a-5.east) {prima};
    \celle{b}{(0,-1.15)}{10,20,30,40,50,{}}
    \node[mnota, anchor=west] at ([xshift=0.25cm]b-5.east) {dopo};
    % la lista
    \node[mnota, anchor=west, font=\Lato\scriptsize\bfseries] at (-1.0,-2.3)
      {lista: si crea un nodo e si cambiano due puntatori};
    \variabile[minimum width=0.7cm]{h}{(0,-3.1)}{head}{}
    \nodo{n1}{(0.8,-3.1)}{10}
    \nodo{n2}{(2.7,-3.1)}{20}
    \nodo{n4}{(6.0,-3.1)}{40}
    \nodo{n5}{(7.9,-3.1)}{50}
    \nodonull{n5}
    \nodo[mlnuovo]{n3}{(4.3,-4.25)}{30}
    \riferimento{h}{-- (n1.west)}
    \collega{n1}{n2}
    \collega{n4}{n5}
    \draw[mrifvecchio] (n2-next.center) -- (n4.west);
    \draw[mrifnuovo] (n3-next.center) -- node[mpasso, pos=0.45, right=3pt] {1} (n4.west);
    \draw[mrifnuovo] (n2-next.center) -- node[mpasso, pos=0.55, left=3pt] {2} (n3.west);
    \node[mpunto] at (n2-next.center) {};
    \node[mpunto] at (n3-next.center) {};
  \end{tikzpicture}
  \caption{Inserire 30 fra 20 e 40. In un array (in alto) gli elementi
           successivi si spostano di un posto, a partire dall'ultimo. In una
           lista (in basso) il nuovo nodo viene collegato cambiando due
           puntatori, nell'ordine indicato; la freccia tratteggiata è il
           collegamento che non c'è più.}
  \label{fig:array-lista}
\end{figure}

\subsection{Vantaggi e svantaggi}

\begin{center}
\renewcommand{\arraystretch}{1.15}
\begin{tabular}{@{}lll@{}}
  \toprule
                                 & \textbf{array}              & \textbf{lista} \\
  \midrule
  dimensione                     & fissata quando lo si crea   & cresce e cala un nodo alla volta \\
  inserire o togliere in testa   & si spostano tutti gli elementi & si cambiano due puntatori \\
  arrivare all'elemento $i$      & immediato: \code{a[i]}      & bisogna attraversare $i$ nodi \\
  memoria                        & solo i dati                 & i dati e un puntatore per nodo \\
  \bottomrule
\end{tabular}
\end{center}

Il vantaggio delle liste è la \textbf{flessibilità}: la dimensione cambia a
ogni inserimento, e inserimenti e cancellazioni sono efficienti perché non
spostano dati. Lo svantaggio è l'\textbf{accesso sequenziale}: non si può
saltare direttamente all'elemento $i$, perché il suo indirizzo si scopre
solo arrivandoci, seguendo i puntatori a partire dal primo nodo. In un
array, invece, l'indirizzo di \code{a[i]} si calcola con una moltiplicazione
e una somma (capitolo~\ref{chap:puntatori}).

\section{Tipi di liste}

Esistono diverse varianti (figura~\ref{fig:tipi-liste}):

\begin{itemize}
  \item nelle liste \textbf{semplicemente collegate} (\textit{singly linked
        list}) ogni nodo punta al successivo. Sono quelle di questo capitolo;
  \item nelle liste \textbf{doppiamente collegate} (\textit{doubly linked
        list}) ogni nodo punta al successivo e al precedente, così la lista
        si può percorrere nei due versi (al prezzo di un puntatore in più per
        nodo, da tenere aggiornato);
  \item nelle liste \textbf{circolari} l'ultimo nodo punta al primo,
        formando un ciclo: non c'è un \code{NULL} che segna la fine, e lo
        scorrimento si ferma quando si torna al nodo da cui si è partiti.
\end{itemize}

\begin{figure}[H]
  \centering
  \begin{tikzpicture}[font=\Lato\scriptsize]
    % semplicemente collegata
    \node[mnota, anchor=west, font=\Lato\scriptsize\bfseries] at (-1.0,0.7)
      {semplicemente collegata: ogni nodo punta al successivo};
    \variabile[minimum width=0.7cm]{h1}{(0,0)}{head}{}
    \nodo{s1}{(0.8,0)}{10}
    \nodo{s2}{(2.7,0)}{20}
    \nodo{s3}{(4.6,0)}{30}
    \nodonull{s3}
    \riferimento{h1}{-- (s1.west)}
    \collega{s1}{s2}
    \collega{s2}{s3}
    % doppiamente collegata
    \node[mnota, anchor=west, font=\Lato\scriptsize\bfseries] at (-1.0,-1.0)
      {doppiamente collegata: ogni nodo punta al successivo e al precedente};
    \variabile[minimum width=0.7cm]{h2}{(0,-1.75)}{head}{}
    \foreach \k/\ix/\val in {1/0.8/10, 2/3.45/20, 3/6.1/30} {
      \node[mlnodo, minimum width=0.45cm, anchor=west] (dp\k) at (\ix,-1.75) {};
      \node[mlnodo, minimum width=0.9cm, anchor=west] (dd\k) at ([xshift=-0.7pt]dp\k.east) {\val};
      \node[mlnodo, minimum width=0.45cm, anchor=west] (dn\k) at ([xshift=-0.7pt]dd\k.east) {};
    }
    \node[mindir, anchor=south] at ([yshift=1pt]dp1.north) {prev};
    \node[mindir, anchor=south] at ([yshift=1pt]dd1.north) {data};
    \node[mindir, anchor=south] at ([yshift=1pt]dn1.north) {next};
    \riferimento{h2}{-- (dp1.west)}
    \foreach \k/\j in {1/2, 2/3} {
      \draw[mrif] ([yshift=0.09cm]dn\k.center) -- ([yshift=0.09cm]dp\j.west);
      \node[mpunto] at ([yshift=0.09cm]dn\k.center) {};
      \draw[mrif] ([yshift=-0.09cm]dp\j.center) -- ([yshift=-0.09cm]dn\k.east);
      \node[mpunto] at ([yshift=-0.09cm]dp\j.center) {};
    }
    \foreach \c in {dp1, dn3}
      \draw[line width=0.7pt, draw=primary!65!black]
        ([shift={(1.5pt,1.5pt)}]\c.south west) -- ([shift={(-1.5pt,-1.5pt)}]\c.north east);
    % circolare
    \node[mnota, anchor=west, font=\Lato\scriptsize\bfseries] at (-1.0,-2.75)
      {circolare: l'ultimo nodo punta al primo};
    \variabile[minimum width=0.7cm]{h3}{(0,-3.45)}{head}{}
    \nodo{c1}{(0.8,-3.45)}{10}
    \nodo{c2}{(2.7,-3.45)}{20}
    \nodo{c3}{(4.6,-3.45)}{30}
    \riferimento{h3}{-- (c1.west)}
    \collega{c1}{c2}
    \collega{c2}{c3}
    \draw[mrif] (c3-next.center) -- ++(0.55,0) |- ([yshift=-0.4cm]c1-dato.south) -- (c1-dato.south);
    \node[mpunto] at (c3-next.center) {};
  \end{tikzpicture}
  \caption{Le tre varianti delle liste. La barra in un campo indica il
           puntatore \code{NULL}.}
  \label{fig:tipi-liste}
\end{figure}

In C le varianti differiscono solo per i campi del nodo e per il modo in cui
si riconosce la fine della lista. Il nodo di una lista doppiamente collegata,
per esempio, ha due puntatori:

\begin{lstlisting}[style=c]
struct DNode {
   int data;
   struct DNode *prev;     /* il nodo precedente */
   struct DNode *next;     /* il nodo successivo */
};
\end{lstlisting}

\section{Una lista in C}

\subsection{Il nodo: una struttura che punta a se stessa}

Il nodo di una lista di interi, come nelle slide, è una struttura con due
campi: il dato e il puntatore al nodo successivo.

\begin{lstlisting}[style=c]
struct Node {
   int data;               /* il dato */
   struct Node *next;      /* il puntatore al nodo successivo */
};
\end{lstlisting}

Il campo \code{next} è un puntatore a \code{struct Node}, cioè al tipo che si
sta definendo: una struttura come questa si dice
\textbf{autoreferenziale}. È lecito perché un puntatore occupa sempre lo
stesso spazio (8 byte), qualunque cosa punti: il compilatore sa quanto posto
riservare a \code{next} anche se \code{struct Node} non è ancora completa.
(Le slide scrivono \code{struct Node* next}: lo spazio prima o dopo
l'asterisco non conta.) Non si può invece mettere dentro un nodo un
\textbf{nodo intero}: conterrebbe un nodo, che conterrebbe un nodo, e così
via all'infinito. Il compilatore lo rifiuta, perché mentre la definizione è
in corso il tipo è ancora incompleto:

\begin{console}
$ gcc -Wall -o nodo_intero nodo_intero.c
nodo_intero.c:3:16: error: field (*@‘@*)next(*@’@*) has incomplete type
    3 |    struct Node next;      /* un nodo intero, non un puntatore */
      |                ^~~~
\end{console}

\subsection{La lista vuota}

Una lista si rappresenta con un puntatore al suo primo nodo, che per
tradizione si chiama \code{head} (testa). La lista vuota non ha nodi, e il
puntatore vale \code{NULL}; le slide la creano così:

\begin{lstlisting}[style=c]
struct Node *head = NULL;      /* la lista vuota */
\end{lstlisting}

L'inizializzazione è essenziale: un puntatore non inizializzato contiene un
indirizzo casuale, e il programma lo scambierebbe per il primo nodo di una
lista.

\subsection{Collegare tre nodi a mano}

Per vedere com'è fatta una lista, costruiamone una di tre nodi a mano:
li allochiamo con \code{malloc}, ne riempiamo i dati e li colleghiamo uno
per uno. Poi scorriamo la lista stampando i dati, e stampiamo anche
l'indirizzo di ogni nodo e il valore del suo campo \code{next}:

\begin{affianco}[0.62][c]
\begin{lstlisting}[style=c, name=nodi.c]
#include <stdio.h>
#include <stdlib.h>

struct Node {
   int data;
   struct Node *next;
};

int main(void)
{
   struct Node *head = NULL;
   struct Node *a, *b, *c, *p;

   a = malloc(sizeof(struct Node));
   b = malloc(sizeof(struct Node));
   c = malloc(sizeof(struct Node));
   if (a == NULL || b == NULL || c == NULL)
      return 1;
   a->data = 10;
   b->data = 20;
   c->data = 30;
   a->next = b;            /* dopo a viene b */
   b->next = c;            /* dopo b viene c */
   c->next = NULL;         /* c è l'ultimo */
   head = a;               /* la lista inizia da a */

   for (p = head; p != NULL; p = p->next)
      printf("%d -> ", p->data);
   printf("NULL\n");
   printf("head = %p\n", (void *) head);
   for (p = head; p != NULL; p = p->next)
      printf("nodo %p: data = %d, next = %p\n",
             (void *) p, p->data, (void *) p->next);
   free(a);
   free(b);
   free(c);
   return 0;
}
\end{lstlisting}
\risultato
Ogni nodo è una \code{struct Node} allocata nello heap; \code{a}, \code{b} e
\code{c} sono puntatori ai tre nodi, e servono solo a costruire la lista e a
liberarla alla fine. Le tre assegnazioni ai campi \code{next} creano i
collegamenti; \code{head = a} dice dove comincia la lista. I due cicli
\code{for} scorrono la lista dall'inizio alla fine, come vedremo fra poco.
\end{affianco}

\begin{console}
$ gcc -Wall -o nodi nodi.c
$ ./nodi
10 -> 20 -> 30 -> NULL
head = 0x5608b5fd32a0
nodo 0x5608b5fd32a0: data = 10, next = 0x5608b5fd32c0
nodo 0x5608b5fd32c0: data = 20, next = 0x5608b5fd32e0
nodo 0x5608b5fd32e0: data = 30, next = (nil)
\end{console}

\subsection{Che cosa mostra l'output}

\begin{figure}[H]
  \centering
  \begin{tikzpicture}[font=\Lato\scriptsize,
      dato/.style={mlnodo, minimum width=0.75cm},
      ind/.style={mlnodo, minimum width=2.25cm, font=\FiraCodeNL\tiny}]
    \node[mvar, minimum width=2.25cm, font=\FiraCodeNL\tiny] (h) at (0,0) {0x5608b5fd32a0};
    \node[mnome, anchor=east] at ([xshift=-2pt]h.west) {head};
    \foreach \k/\ix/\val/\nx/\ad/\nome in {1/1.85/10/0x5608b5fd32c0/0x5608b5fd32a0/a,
        2/5.55/20/0x5608b5fd32e0/0x5608b5fd32c0/b, 3/9.25/30/{(nil)}/0x5608b5fd32e0/c} {
      \node[dato, anchor=west] (d\k) at (\ix,0) {\val};
      \node[ind, anchor=west] (x\k) at ([xshift=-0.7pt]d\k.east) {\nx};
      \node[mnota, anchor=south, align=center, inner sep=1pt]
        at ([yshift=2pt]$(d\k.north west)!0.5!(x\k.north east)$)
        {nodo \code{*\nome}, all'indirizzo\\{\FiraCodeNL\tiny\color{neutral!40!black}\ad}};
      \node[mindir, anchor=north] at ([yshift=-1pt]d\k.south) {data};
      \node[mindir, anchor=north] at ([yshift=-1pt]x\k.south) {next};
    }
    \draw[mrif] (h.east) -- (d1.west);
    \draw[mrif] (x1.east) -- (d2.west);
    \draw[mrif] (x2.east) -- (d3.west);
  \end{tikzpicture}
  \caption{La lista di \code{nodi.c} con gli indirizzi reali dell'esecuzione:
           ogni campo \code{next} contiene l'indirizzo del nodo successivo, e
           le frecce sono solo un modo di disegnarlo.}
  \label{fig:nodi}
\end{figure}

\begin{itemize}
  \item \code{head} contiene l'indirizzo del primo nodo, e il campo
        \code{next} di ogni nodo contiene l'\textbf{indirizzo} del nodo
        successivo (figura~\ref{fig:nodi}). Le frecce con cui disegneremo
        le liste rappresentano proprio questi indirizzi.
  \item Il \code{next} dell'ultimo nodo vale \code{NULL}, che \code{printf}
        con \code{\%p} stampa come \code{(nil)}.
  \item Nella prova i nodi distano 32 byte l'uno dall'altro (\code{0x20} in
        esadecimale): 16 sono del nodo, gli altri sono usati da
        \code{malloc} per gestire i blocchi. Sono vicini perché sono stati
        allocati uno dopo l'altro, ma nulla lo garantisce: in un programma
        che inserisce e cancella, i nodi finiscono sparsi nello heap, e
        l'unico modo per sapere dov'è il successivo è leggere \code{next}.
\end{itemize}

\subsection{Scorrere una lista}

Il ciclo che stampa i dati è lo schema di quasi tutte le operazioni sulle
liste, e conviene capirlo bene:

\begin{lstlisting}[style=c]
for (p = head; p != NULL; p = p->next)
   printf("%d -> ", p->data);
\end{lstlisting}

\begin{itemize}
  \item \code{p = head}: un puntatore di appoggio parte dal primo nodo. Si
        usa \code{p} e non \code{head} perché \code{head} non deve cambiare:
        è l'unico modo di ritrovare la lista.
  \item \code{p != NULL}: finché \code{p} punta a un nodo, il corpo del ciclo
        lo elabora (qui ne stampa il dato, \code{p->data}).
  \item \code{p = p->next}: \code{p} passa al nodo successivo, prendendone
        l'indirizzo dal nodo corrente. Dopo l'ultimo nodo \code{p} diventa
        \code{NULL}, e il ciclo termina (figura~\ref{fig:scorri}).
\end{itemize}

\begin{figure}[H]
  \centering
  \begin{tikzpicture}[font=\Lato\scriptsize]
    \variabile[minimum width=0.7cm]{h}{(0,0)}{head}{}
    \nodo{n1}{(0.8,0)}{10}
    \nodo{n2}{(3.0,0)}{20}
    \nodo{n3}{(5.2,0)}{30}
    \nodonull{n3}
    \riferimento{h}{-- (n1.west)}
    \collega{n1}{n2}
    \collega{n2}{n3}
    \foreach \k/\testo in {1/{primo giro:\\stampa 10}, 2/{secondo giro:\\stampa 20},
        3/{terzo giro:\\stampa 30}} {
      \node[mvar, minimum width=0.6cm] (p\k) at ([yshift=-1.0cm]n\k-dato.south) {};
      \node[mnome, anchor=east] at ([xshift=-2pt]p\k.west) {p};
      \riferimento{p\k}{-- (n\k-dato.south)}
      \node[mnota, anchor=north, align=center] at ([yshift=-2pt]p\k.south) {\testo};
    }
    \node[mvar, minimum width=0.9cm, font=\FiraCodeNL\tiny] (p4) at ([xshift=2.2cm]p3) {NULL};
    \node[mnome, anchor=east] at ([xshift=-2pt]p4.west) {p};
    \node[mnota, anchor=north, align=center] at ([yshift=-2pt]p4.south) {\code{p == NULL}:\\il ciclo finisce};
  \end{tikzpicture}
  \caption{Il puntatore \code{p} durante lo scorrimento: a ogni giro prende
           l'indirizzo contenuto nel campo \code{next} del nodo su cui si
           trova.}
  \label{fig:scorri}
\end{figure}

Lo stesso ciclo, con un contatore al posto della \code{printf}, calcola la
\textbf{lunghezza} della lista. Nel modulo che scriveremo fra poco diventano
due funzioni:

\begin{affianco}[0.5][c]
\begin{lstlisting}[style=c, name=lista.c]
void stampaLista(struct Node *head)
{
   struct Node *p;

   for (p = head; p != NULL; p = p->next)
      printf("%d -> ", p->data);
   printf("NULL\n");
}
\end{lstlisting}
\risultato
\begin{lstlisting}[style=c, name=lista.c]
int lunghezza(struct Node *head)
{
   struct Node *p;
   int n = 0;

   for (p = head; p != NULL; p = p->next)
      n++;
   return n;
}
\end{lstlisting}
\end{affianco}

Qui \code{head} è un parametro, cioè una copia del puntatore del chiamante:
la funzione potrebbe anche scorrere la lista usando direttamente
\code{head}, senza cambiare la \code{head} di chi la chiama. Usare \code{p}
rende però il codice più chiaro, ed è l'abitudine giusta.

\section{Le operazioni sulle liste}

\subsection{Un modulo per le liste}

Le slide elencano le operazioni fondamentali: creazione della lista,
inserimento di un nodo, eliminazione di un nodo, ricerca di un elemento,
lunghezza della lista, unione di due liste. Le scriviamo come funzioni e le
raccogliamo in un \textbf{modulo}, come nel capitolo~\ref{chap:make}: i
prototipi in \code{lista.h}, le definizioni in \code{lista.c}. Un programma
che usa le liste include \code{lista.h} e si compila insieme a
\code{lista.c}.

\begin{lstlisting}[style=h, name=lista.h]
struct Node {
   int data;               /* il dato */
   struct Node *next;      /* il puntatore al nodo successivo */
};

struct Node *nuovoNodo(int data);
void stampaLista(struct Node *head);
int lunghezza(struct Node *head);
struct Node *inserisciTesta(struct Node *head, int data);
struct Node *inserisciCoda(struct Node *head, int data);
struct Node *inserisciPosizione(struct Node *head, int data, int pos);
struct Node *inserisciOrdinato(struct Node *head, int data);
struct Node *cerca(struct Node *head, int data);
struct Node *eliminaTesta(struct Node *head);
struct Node *eliminaCoda(struct Node *head);
struct Node *eliminaValore(struct Node *head, int data);
struct Node *eliminaPosizione(struct Node *head, int pos);
struct Node *unisci(struct Node *a, struct Node *b);
void liberaLista(struct Node *head);
\end{lstlisting}

\code{lista.c} comincia includendo \code{stdio.h}, \code{stdlib.h} (per
\code{malloc} e \code{free}) e \code{lista.h}; le sue funzioni sono mostrate
una alla volta nelle prossime sezioni. Sono state tutte provate anche nei
casi limite (lista vuota, posizioni che non esistono) con
\code{-fsanitize=address}, e con valgrind, che vedremo alla fine del
capitolo, per controllare che ogni nodo venga liberato.

\subsection{Restituire la nuova testa}

Molte operazioni possono cambiare il primo nodo: inserire in testa, eliminare
il primo nodo, inserire in una lista vuota. Ma \code{head} è una variabile
del chiamante, e la funzione ne riceve una \textbf{copia}
(capitolo~\ref{chap:funzioni}): assegnare un nuovo valore al parametro non
cambia la \code{head} del chiamante. Per questo le funzioni del modulo, come
quelle del progetto delle lezioni, \textbf{restituiscono la testa} della
lista dopo l'operazione, e il chiamante la riassegna:

\begin{lstlisting}[style=c]
head = inserisciTesta(head, 10);
\end{lstlisting}

È lo stesso schema di \code{e = update(e, 3)} del
capitolo~\ref{chap:strutture}. Un errore tipico è dimenticare
l'assegnazione, scrivendo soltanto \code{inserisciTesta(head, 10);}: il nodo
viene creato, ma \code{head} continua a puntare al vecchio primo nodo, e il
nuovo nodo è perso, con la sua memoria, che non potrà più essere liberata.

\subsection{Creare un nodo}

Ogni inserimento comincia creando un nodo nello heap:

\begin{affianco}[0.55][c]
\begin{lstlisting}[style=c, name=lista.c]
struct Node *nuovoNodo(int data)
{
   struct Node *nodo = malloc(sizeof(struct Node));

   if (nodo == NULL)
      return NULL;
   nodo->data = data;
   nodo->next = NULL;
   return nodo;
}
\end{lstlisting}
\risultato
\code{malloc} restituisce \code{NULL} se la memoria è finita: in quel caso
anche \code{nuovoNodo} restituisce \code{NULL}, e chi la chiama deve
controllarlo. Il campo \code{next} parte da \code{NULL}: il nodo appena
creato non è ancora collegato a niente.
\end{affianco}

\section{Inserimento}

Le slide prevedono quattro modi di inserire un nodo: in testa, in coda, in
una posizione specificata e in ordine (per dati che si possono ordinare,
come gli interi).

\subsection{In testa}

\begin{lstlisting}[style=c, name=lista.c]
struct Node *inserisciTesta(struct Node *head, int data)
{
   struct Node *nodo = nuovoNodo(data);

   if (nodo == NULL)
      return head;          /* niente memoria: la lista resta com'era */
   nodo->next = head;       /* 1. il nuovo nodo punta al vecchio primo */
   return nodo;             /* 2. e diventa il primo */
}
\end{lstlisting}

\begin{figure}[H]
  \centering
  \begin{tikzpicture}[font=\Lato\scriptsize]
    \variabile[minimum width=0.7cm]{h}{(0,0)}{head}{}
    \nodo[mlnuovo]{n1}{(1.0,1.05)}{10}
    \nodo{n2}{(3.1,0)}{20}
    \nodo{n3}{(5.0,0)}{30}
    \nodonull{n3}
    \collega{n2}{n3}
    \draw[mrifvecchio] (h.center) -- (n2.west);
    \draw[mrifnuovo] (n1-next.center) -- node[mpasso, pos=0.45, above right=0pt] {1} (n2.west);
    \draw[mrifnuovo] (h.center) -- node[mpasso, pos=0.5, above left=0pt] {2} (n1.west);
    \node[mpunto] at (n1-next.center) {};
    \node[mpunto] at (h.center) {};
    \node[mnota, anchor=south] at ([yshift=2pt]n1.north) {il nuovo nodo};
  \end{tikzpicture}
  \caption{Inserimento in testa: prima il nuovo nodo punta al vecchio primo
           nodo, poi \code{head} punta al nuovo nodo.}
  \label{fig:ins-testa}
\end{figure}

I passi sono due (figura~\ref{fig:ins-testa}): \wtpassonum{1} il nuovo nodo
punta al vecchio primo nodo, il cui indirizzo è in \code{head};
\wtpassonum{2} il nuovo nodo diventa il primo: la funzione ne restituisce
l'indirizzo, e il chiamante lo assegna a \code{head}. La funzione va bene
anche per la lista vuota: \code{nodo->next} riceve \code{NULL}, e il nuovo
nodo è l'unico della lista. Non c'è nessun ciclo: l'inserimento in testa
richiede lo stesso tempo, qualunque sia la lunghezza della lista.

\subsection{In coda}

Per inserire in coda serve l'ultimo nodo, cioè quello il cui \code{next}
vale \code{NULL}:

\begin{lstlisting}[style=c, name=lista.c]
struct Node *inserisciCoda(struct Node *head, int data)
{
   struct Node *nodo = nuovoNodo(data);
   struct Node *p;

   if (nodo == NULL)
      return head;
   if (head == NULL)        /* lista vuota: il nodo è l'unico */
      return nodo;
   for (p = head; p->next != NULL; p = p->next)
      ;                     /* si ferma sull'ultimo nodo */
   p->next = nodo;
   return head;
}
\end{lstlisting}

\begin{figure}[H]
  \centering
  \begin{tikzpicture}[font=\Lato\scriptsize]
    \variabile[minimum width=0.7cm]{h}{(0,0)}{head}{}
    \nodo{n1}{(0.8,0)}{10}
    \nodo{n2}{(2.7,0)}{20}
    \nodo{n3}{(4.6,0)}{30}
    \nodo[mlnuovo]{n4}{(6.9,0)}{40}
    \nodonull{n4}
    \riferimento{h}{-- (n1.west)}
    \collega{n1}{n2}
    \collega{n2}{n3}
    \draw[mrifnuovo] (n3-next.center) -- node[mpasso, above=1pt] {2} (n4.west);
    \node[mpunto] at (n3-next.center) {};
    \node[mvar, minimum width=0.6cm] (p) at ([yshift=-1.0cm]n3-dato.south) {};
    \node[mnome, anchor=east] (pn) at ([xshift=-2pt]p.west) {p};
    \riferimento{p}{-- (n3-dato.south)}
    \node[mpasso, anchor=east] at ([xshift=-3pt]pn.west) {1};
    \node[mnota, anchor=west] at ([xshift=4pt]p.east)
      {il ciclo si ferma sull'ultimo nodo: \code{p->next == NULL}};
    \node[mnota, anchor=south] at ([yshift=2pt]n4.north) {il nuovo nodo};
  \end{tikzpicture}
  \caption{Inserimento in coda: \code{p} si ferma sull'ultimo nodo, poi il
           suo \code{next} punta al nuovo nodo.}
  \label{fig:ins-coda}
\end{figure}

\begin{itemize}
  \item Il ciclo ha il corpo vuoto (il \code{;} da solo,
        capitolo~\ref{chap:flusso}): serve solo a far avanzare \code{p}. A
        differenza del ciclo di scorrimento, la condizione è
        \code{p->next != NULL}, quindi il ciclo si ferma \textbf{sull'ultimo
        nodo}, non dopo (figura~\ref{fig:ins-coda}). Poi \code{p->next = nodo}
        aggancia il nuovo nodo, il cui \code{next} è già \code{NULL}.
  \item La lista vuota va trattata a parte: non ha un ultimo nodo, e con
        \code{p} uguale a \code{NULL} la condizione \code{p->next != NULL}
        leggerebbe un campo di \code{NULL}, e il programma verrebbe fermato.
        In quel caso il nuovo nodo è tutta la lista.
  \item Per arrivare all'ultimo nodo bisogna attraversarli tutti:
        l'inserimento in coda richiede un tempo proporzionale alla lunghezza
        della lista.
\end{itemize}

\subsection{In una posizione}

Le posizioni si contano da 0, come gli indici degli array: inserire in
posizione 0 è inserire in testa. Per inserire in posizione \code{pos} bisogna
arrivare al nodo \textbf{precedente}, quello in posizione \code{pos - 1},
perché è il suo campo \code{next} che deve cambiare:

\begin{lstlisting}[style=c, name=lista.c]
struct Node *inserisciPosizione(struct Node *head, int data, int pos)
{
   struct Node *nodo, *prec;
   int i;

   if (pos == 0)
      return inserisciTesta(head, data);
   prec = head;             /* cerca il nodo in posizione pos - 1 */
   for (i = 0; prec != NULL && i < pos - 1; i++)
      prec = prec->next;
   if (prec == NULL)        /* la lista è troppo corta */
      return head;
   nodo = nuovoNodo(data);
   if (nodo == NULL)
      return head;
   nodo->next = prec->next; /* 1. il nuovo nodo punta al successivo */
   prec->next = nodo;       /* 2. il precedente punta al nuovo nodo */
   return head;
}
\end{lstlisting}

\begin{figure}[H]
  \centering
  \begin{tikzpicture}[font=\Lato\scriptsize]
    \variabile[minimum width=0.7cm]{h}{(0,0)}{head}{}
    \nodo{n1}{(0.8,0)}{10}
    \nodo{n2}{(2.7,0)}{20}
    \nodo{n4}{(6.0,0)}{40}
    \nodonull{n4}
    \nodo[mlnuovo]{n3}{(4.3,-1.15)}{30}
    \riferimento{h}{-- (n1.west)}
    \collega{n1}{n2}
    \draw[mrifvecchio] (n2-next.center) -- (n4.west);
    \draw[mrifnuovo] (n3-next.center) -- node[mpasso, pos=0.45, right=3pt] {1} (n4.west);
    \draw[mrifnuovo] (n2-next.center) -- node[mpasso, pos=0.5, right=2pt] {2} (n3.west);
    \node[mpunto] at (n2-next.center) {};
    \node[mpunto] at (n3-next.center) {};
    \node[mvar, minimum width=0.6cm] (pr) at ([yshift=1.0cm]n2-dato.north) {};
    \node[mnome, anchor=east] at ([xshift=-2pt]pr.west) {prec};
    \riferimento{pr}{-- (n2-dato.north)}
    \node[mindir, anchor=north] at ([yshift=-1pt]n1-dato.south) {posizione 0};
    \node[mindir, anchor=north] at ([yshift=-1pt]n2-dato.south) {posizione 1};
    \node[mnota, anchor=north] at ([yshift=-2pt]n3.south) {il nuovo nodo, in posizione 2};
  \end{tikzpicture}
  \caption{Inserimento di 30 in posizione 2: \code{prec} si ferma sul nodo
           in posizione 1; poi il nuovo nodo punta al successivo di
           \code{prec}, e \code{prec} punta al nuovo nodo.}
  \label{fig:ins-pos}
\end{figure}

Il ciclo \code{for} fa avanzare \code{prec} di \code{pos - 1} nodi a partire
dal primo; se la lista finisce prima (\code{prec} diventa \code{NULL}), la
posizione non esiste, e la lista resta com'è. Poi vengono i due passi della
figura~\ref{fig:ins-pos}: \wtpassonum{1} il nuovo nodo punta al successivo
di \code{prec}; \wtpassonum{2} \code{prec} punta al nuovo nodo. La testa non
cambia, e la funzione restituisce \code{head} così com'era.

\warning{L'ordine dei due passi conta}{
  Se si scambiano i passi, scrivendo prima \code{prec->next = nodo} e poi
  \code{nodo->next = prec->next}, il secondo assegnamento copia in
  \code{nodo->next} l'indirizzo di \code{nodo} stesso: il nuovo nodo punta
  a se stesso, e l'indirizzo del resto della lista, che stava solo in
  \code{prec->next}, è perso per sempre. Uno scorrimento della lista non
  terminerebbe più. La regola: prima si collega il nuovo nodo al resto
  della lista, poi lo si aggancia al precedente.}

\subsection{In ordine}

Se la lista è ordinata (qui in ordine crescente), un nuovo dato va inserito
al suo posto, in modo che la lista resti ordinata. È un inserimento in
posizione, in cui la posizione si trova confrontando i dati:

\begin{lstlisting}[style=c, name=lista.c]
struct Node *inserisciOrdinato(struct Node *head, int data)
{
   struct Node *nodo, *prec;

   if (head == NULL || data <= head->data)
      return inserisciTesta(head, data);
   prec = head;             /* avanza finché il successivo è minore */
   while (prec->next != NULL && prec->next->data < data)
      prec = prec->next;
   nodo = nuovoNodo(data);
   if (nodo == NULL)
      return head;
   nodo->next = prec->next;
   prec->next = nodo;
   return head;
}
\end{lstlisting}

\begin{itemize}
  \item Se la lista è vuota, o il nuovo dato non supera il primo, il nodo
        va in testa. (Anche qui conta l'ordine delle condizioni: se
        \code{head} è \code{NULL}, la valutazione in cortocircuito,
        capitolo~\ref{chap:flusso}, evita di leggere \code{head->data}.)
  \item Altrimenti \code{prec} avanza finché il nodo \textbf{successivo} ha
        un dato minore di quello da inserire; quando si ferma, il nuovo nodo
        va subito dopo \code{prec}, con gli stessi due passi
        dell'inserimento in posizione.
  \item Se \code{prec} arriva all'ultimo nodo, \code{prec->next} vale
        \code{NULL}: la prima condizione del \code{while} è falsa, la
        seconda non viene valutata (non si legge il campo \code{data} di
        \code{NULL}), e il nuovo nodo va in coda.
\end{itemize}

Seguiamo a mano l'inserimento di 25 nella lista 10, 20, 30, 40. Il primo
dato, 10, è minore di 25, quindi non si inserisce in testa, e \code{prec}
parte dal primo nodo:

\begin{center}
\renewcommand{\arraystretch}{1.15}
\begin{tabular}{@{}cccl@{}}
  \toprule
  \textbf{\code{prec}} & \textbf{\code{prec->next->data}} & \textbf{minore di 25?} & \textbf{che cosa succede} \\
  \midrule
  10 & 20 & sì & \code{prec} avanza \\
  20 & 30 & no & il ciclo si ferma: 25 va fra 20 e 30 \\
  \bottomrule
\end{tabular}
\end{center}

\subsection{Una prova}

Il programma seguente usa le quattro funzioni, partendo da una lista vuota:

\begin{affianco}[0.6][c]
\begin{lstlisting}[style=c, name=inserimenti.c]
#include <stdio.h>
#include "lista.h"

int main(void)
{
   struct Node *head = NULL;

   head = inserisciTesta(head, 20);
   head = inserisciTesta(head, 10);
   printf("in testa 20 e 10:    ");
   stampaLista(head);
   head = inserisciCoda(head, 40);
   printf("in coda 40:          ");
   stampaLista(head);
   head = inserisciPosizione(head, 30, 2);
   printf("30 in posizione 2:   ");
   stampaLista(head);
   head = inserisciOrdinato(head, 25);
   head = inserisciOrdinato(head, 5);
   head = inserisciOrdinato(head, 50);
   printf("in ordine 25, 5, 50: ");
   stampaLista(head);
   printf("lunghezza: %d\n", lunghezza(head));
   liberaLista(head);
   return 0;
}
\end{lstlisting}
\risultato
\begin{itemize}
  \item Inserendo in testa prima 20 e poi 10, la lista è 10, 20: ogni nodo
        inserito in testa va davanti a tutti gli altri, quindi l'ordine
        della lista è l'opposto di quello degli inserimenti.
  \item 30 va in posizione 2, fra 20 e 40.
  \item Con \code{inserisciOrdinato} il 25 va fra 20 e 30, il 5 in testa e
        il 50 in coda: la lista resta ordinata.
  \item \code{liberaLista}, che vedremo alla fine, restituisce la memoria
        di tutti i nodi.
\end{itemize}
\end{affianco}

\begin{console}
$ gcc -Wall -o inserimenti inserimenti.c lista.c
$ ./inserimenti
in testa 20 e 10:    10 -> 20 -> NULL
in coda 40:          10 -> 20 -> 40 -> NULL
30 in posizione 2:   10 -> 20 -> 30 -> 40 -> NULL
in ordine 25, 5, 50: 5 -> 10 -> 20 -> 25 -> 30 -> 40 -> 50 -> NULL
lunghezza: 7
\end{console}

\section{Ricerca}

Cercare un elemento vuol dire trovare il nodo che contiene un certo dato. La
funzione scorre la lista e restituisce l'\textbf{indirizzo} del primo nodo
con quel dato, oppure \code{NULL} se non c'è:

\begin{affianco}[0.55][c]
\begin{lstlisting}[style=c, name=lista.c]
struct Node *cerca(struct Node *head, int data)
{
   struct Node *p;

   for (p = head; p != NULL; p = p->next)
      if (p->data == data)
         return p;          /* trovato: l'indirizzo del nodo */
   return NULL;             /* non trovato */
}
\end{lstlisting}
\risultato
Restituire l'indirizzo, e non un semplice sì o no, permette al chiamante di
usare il nodo trovato: leggerne o cambiarne i dati, guardare il successivo.
\end{affianco}

La ricerca è necessariamente \textbf{sequenziale}: anche se la lista è
ordinata, non c'è modo di saltare a metà lista, e nel caso peggiore (dato
assente, o nell'ultimo nodo) bisogna attraversarla tutta. Una prova della
ricerca è nella sezione seguente, insieme a quelle delle eliminazioni.

\section{Eliminazione}

Eliminare un nodo richiede due cose: \textbf{scollegarlo}, cioè fare in modo
che nessun nodo della lista punti più a lui, e \textbf{liberarne} la memoria
con \code{free}. Le slide elencano quattro casi: in testa, in coda, per
valore (il primo nodo che ha una certa proprietà, qui un certo dato) e data
la posizione.

\subsection{In testa}

\begin{lstlisting}[style=c, name=lista.c]
struct Node *eliminaTesta(struct Node *head)
{
   struct Node *secondo;

   if (head == NULL)
      return NULL;
   secondo = head->next;    /* 1. ricorda il secondo nodo */
   free(head);              /* 2. libera il primo */
   return secondo;          /* 3. il secondo diventa il primo */
}
\end{lstlisting}

\begin{figure}[H]
  \centering
  \begin{tikzpicture}[font=\Lato\scriptsize]
    \variabile[minimum width=0.7cm]{h}{(0,0)}{head}{}
    \nodo[mlvia]{n1}{(0.9,0)}{10}
    \nodo{n2}{(2.9,0)}{20}
    \nodo{n3}{(4.8,0)}{30}
    \nodonull{n3}
    \collega[mrifvecchio]{n1}{n2}
    \collega{n2}{n3}
    \draw[mrifvecchio] (h.center) -- (n1.west);
    \draw[mrifnuovo] (h.center) to[bend right=40] node[mpasso, pos=0.5, below=1pt] {3} (n2.south west);
    \node[mpunto] at (h.center) {};
    \node[mvar, minimum width=0.6cm] (s) at ([yshift=1.0cm]n2-dato.north) {};
    \node[mnome, anchor=east] at ([xshift=-2pt]s.west) {secondo};
    \riferimento{s}{-- (n2-dato.north)}
    \node[mpasso, anchor=west] at ([xshift=3pt]s.east) {1};
    \node[mpasso] (q2) at ([yshift=0.45cm]n1-dato.north) {2};
    \node[mnota, anchor=east] at ([xshift=-1pt]q2.west) {\code{free(head)}};
  \end{tikzpicture}
  \caption{Eliminazione in testa: si ricorda il secondo nodo, si libera il
           primo, e il secondo diventa la nuova testa.}
  \label{fig:el-testa}
\end{figure}

Il passo \wtpassonum{1} è indispensabile: dopo \code{free(head)} il nodo non
esiste più, e il suo campo \code{next} non si può più leggere. L'indirizzo
del secondo nodo va quindi salvato \textbf{prima}. Poi il chiamante assegna
il risultato a \code{head} (passo \wtpassonum{3}, figura~\ref{fig:el-testa}).

\subsection{In coda}

\begin{lstlisting}[style=c, name=lista.c]
struct Node *eliminaCoda(struct Node *head)
{
   struct Node *p;

   if (head == NULL)
      return NULL;
   if (head->next == NULL) {   /* un solo nodo */
      free(head);
      return NULL;
   }
   for (p = head; p->next->next != NULL; p = p->next)
      ;                     /* si ferma sul penultimo nodo */
   free(p->next);
   p->next = NULL;
   return head;
}
\end{lstlisting}

\begin{figure}[H]
  \centering
  \begin{tikzpicture}[font=\Lato\scriptsize]
    \variabile[minimum width=0.7cm]{h}{(0,0)}{head}{}
    \nodo{n1}{(0.8,0)}{10}
    \nodo{n2}{(2.7,0)}{20}
    \nodo[mlvia]{n3}{(5.0,0)}{30}
    \nodonull{n3}
    \riferimento{h}{-- (n1.west)}
    \collega{n1}{n2}
    \draw[mrifvecchio] (n2-next.east) -- (n3.west);
    \draw[line width=0.8pt, draw=mok] ([shift={(1.5pt,1.5pt)}]n2-next.south west)
      -- ([shift={(-1.5pt,-1.5pt)}]n2-next.north east);
    \node[mvar, minimum width=0.6cm] (p) at ([yshift=-1.0cm]n2-dato.south) {};
    \node[mnome, anchor=east] (pn) at ([xshift=-2pt]p.west) {p};
    \riferimento{p}{-- (n2-dato.south)}
    \node[mpasso, anchor=east] at ([xshift=-3pt]pn.west) {1};
    \node[mnota, anchor=west] at ([xshift=4pt]p.east)
      {si ferma sul penultimo nodo: \code{p->next->next == NULL}};
    \node[mpasso] (q2) at ([yshift=0.45cm]n3-dato.north) {2};
    \node[mnota, anchor=west] at ([xshift=1pt]q2.east) {\code{free(p->next)}};
    \node[mpasso] (q3) at ([yshift=0.45cm]n2-next.north) {3};
    \node[mnota, anchor=east] at ([xshift=-1pt]q3.west) {\code{p->next = NULL}};
  \end{tikzpicture}
  \caption{Eliminazione in coda: \code{p} si ferma sul penultimo nodo, che
           dopo il \code{free} dell'ultimo diventa l'ultimo.}
  \label{fig:el-coda}
\end{figure}

Per eliminare l'ultimo nodo bisogna cambiare il \code{next} del
\textbf{penultimo}, che deve diventare \code{NULL}. Il ciclo si ferma quando
il successivo del successivo di \code{p} non esiste:
\code{p->next->next == NULL} vuol dire che \code{p->next} è l'ultimo nodo,
e quindi \code{p} è il penultimo (figura~\ref{fig:el-coda}). Una lista di un
solo nodo non ha un penultimo, e va trattata a parte: dopo l'eliminazione è
vuota, e la funzione restituisce \code{NULL}.

\subsection{Per valore}

\begin{lstlisting}[style=c, name=lista.c]
struct Node *eliminaValore(struct Node *head, int data)
{
   struct Node *prec, *p;

   if (head == NULL)
      return NULL;
   if (head->data == data)
      return eliminaTesta(head);
   prec = head;             /* avanza finché il successivo non è quello */
   while (prec->next != NULL && prec->next->data != data)
      prec = prec->next;
   if (prec->next != NULL) {   /* trovato: è prec->next */
      p = prec->next;
      prec->next = p->next;    /* 1. il precedente scavalca il nodo */
      free(p);                 /* 2. il nodo viene liberato */
   }
   return head;
}
\end{lstlisting}

\begin{figure}[H]
  \centering
  \begin{tikzpicture}[font=\Lato\scriptsize]
    \variabile[minimum width=0.7cm]{h}{(0,0)}{head}{}
    \nodo{n1}{(0.8,0)}{20}
    \nodo{n2}{(2.7,0)}{30}
    \nodo[mlvia]{n3}{(4.6,0)}{40}
    \nodo{n4}{(6.5,0)}{50}
    \nodonull{n4}
    \riferimento{h}{-- (n1.west)}
    \collega{n1}{n2}
    \collega[mrifvecchio]{n2}{n3}
    \collega[mrifvecchio]{n3}{n4}
    \draw[mrifnuovo] (n2-next.center) to[bend right=45]
      node[mpasso, pos=0.5, below=1pt] {1} (n4.south west);
    \node[mvar, minimum width=0.6cm] (pr) at ([yshift=1.0cm]n2-dato.north) {};
    \node[mnome, anchor=east] at ([xshift=-2pt]pr.west) {prec};
    \riferimento{pr}{-- (n2-dato.north)}
    \node[mvar, minimum width=0.6cm] (pp) at ([yshift=1.0cm]n3-dato.north) {};
    \node[mnome, anchor=east] at ([xshift=-2pt]pp.west) {p};
    \riferimento{pp}{-- (n3-dato.north)}
    \node[mpasso, anchor=west] (q2) at ([xshift=4pt]pp.east) {2};
    \node[mnota, anchor=west] at ([xshift=1pt]q2.east) {\code{free(p)}};
  \end{tikzpicture}
  \caption{Eliminazione del 40: \code{prec} scavalca il nodo da eliminare,
           puntando al successivo, e poi il nodo viene liberato.}
  \label{fig:el-valore}
\end{figure}

Come per l'inserimento, serve il nodo \textbf{precedente} a quello da
eliminare, perché è il suo \code{next} che cambia. Il ciclo avanza finché il
successivo di \code{prec} non contiene il dato cercato; se lo trova,
\wtpassonum{1} \code{prec} scavalca il nodo, puntando a quello dopo, e
\wtpassonum{2} il nodo viene liberato (figura~\ref{fig:el-valore}). Anche
qui l'ordine conta: dopo \code{free(p)} non si potrebbe più leggere
\code{p->next}. Se il dato non c'è, la lista resta com'è; se è nel primo
nodo, è un'eliminazione in testa.

\subsection{Data la posizione}

\begin{lstlisting}[style=c, name=lista.c]
struct Node *eliminaPosizione(struct Node *head, int pos)
{
   struct Node *prec, *p;
   int i;

   if (head == NULL)
      return NULL;
   if (pos == 0)
      return eliminaTesta(head);
   prec = head;             /* cerca il nodo in posizione pos - 1 */
   for (i = 0; prec->next != NULL && i < pos - 1; i++)
      prec = prec->next;
   if (prec->next == NULL)  /* non c'è un nodo in posizione pos */
      return head;
   p = prec->next;
   prec->next = p->next;
   free(p);
   return head;
}
\end{lstlisting}

È la stessa operazione, ma \code{prec} si ferma sul nodo in posizione
\code{pos - 1} contando i nodi, invece di confrontare i dati. Se il nodo in
posizione \code{pos} non esiste, la lista resta com'è.

\subsection{Una prova}

\begin{affianco}[0.6][c]
\begin{lstlisting}[style=c, name=eliminazioni.c]
#include <stdio.h>
#include "lista.h"

int main(void)
{
   struct Node *head = NULL, *p;
   int i;

   for (i = 1; i <= 7; i++)
      head = inserisciCoda(head, i * 10);
   printf("all'inizio:          ");
   stampaLista(head);
   p = cerca(head, 40);
   if (p != NULL)
      printf("cerca(40): trovato, il successivo vale %d\n",
             p->next->data);
   if (cerca(head, 45) == NULL)
      printf("cerca(45): non trovato\n");
   head = eliminaTesta(head);
   printf("elimina in testa:    ");
   stampaLista(head);
   head = eliminaCoda(head);
   printf("elimina in coda:     ");
   stampaLista(head);
   head = eliminaValore(head, 40);
   printf("elimina il 40:       ");
   stampaLista(head);
   head = eliminaPosizione(head, 1);
   printf("elimina posizione 1: ");
   stampaLista(head);
   liberaLista(head);
   return 0;
}
\end{lstlisting}
\risultato
\begin{itemize}
  \item \code{cerca(head, 40)} restituisce l'indirizzo del nodo con il 40,
        e attraverso di esso il programma legge il dato del nodo
        successivo, 50.
  \item Le eliminazioni tolgono, nell'ordine, il primo nodo (10), l'ultimo
        (70), il nodo con il 40 e quello in posizione 1, che a quel punto è
        il 30.
\end{itemize}
\end{affianco}

\begin{console}
$ gcc -Wall -o eliminazioni eliminazioni.c lista.c
$ ./eliminazioni
all'inizio:          10 -> 20 -> 30 -> 40 -> 50 -> 60 -> 70 -> NULL
cerca(40): trovato, il successivo vale 50
cerca(45): non trovato
elimina in testa:    20 -> 30 -> 40 -> 50 -> 60 -> 70 -> NULL
elimina in coda:     20 -> 30 -> 40 -> 50 -> 60 -> NULL
elimina il 40:       20 -> 30 -> 50 -> 60 -> NULL
elimina posizione 1: 20 -> 50 -> 60 -> NULL
\end{console}

\section{Unione di due liste}

\subsection{Agganciare due liste}

L'unione (concatenazione) di due liste \code{a} e \code{b} si ottiene
agganciando il primo nodo di \code{b} all'ultimo nodo di \code{a}: non si
copia nessun nodo, si cambia un solo puntatore (figura~\ref{fig:unione}). Se
\code{a} è vuota, l'unione è semplicemente \code{b}.

\begin{affianco}[0.55][c]
\begin{lstlisting}[style=c, name=lista.c]
struct Node *unisci(struct Node *a, struct Node *b)
{
   struct Node *p;

   if (a == NULL)
      return b;
   for (p = a; p->next != NULL; p = p->next)
      ;                     /* l'ultimo nodo di a */
   p->next = b;             /* punta al primo nodo di b */
   return a;
}
\end{lstlisting}
\risultato
\begin{lstlisting}[style=c, name=unione.c]
#include <stdio.h>
#include "lista.h"

int main(void)
{
   struct Node *a = NULL, *b = NULL;

   a = inserisciCoda(a, 1);
   a = inserisciCoda(a, 2);
   a = inserisciCoda(a, 3);
   b = inserisciCoda(b, 7);
   b = inserisciCoda(b, 8);
   printf("a: ");
   stampaLista(a);
   printf("b: ");
   stampaLista(b);
   a = unisci(a, b);
   printf("a dopo unisci(a, b): ");
   stampaLista(a);
   printf("lunghezza: %d\n", lunghezza(a));
   liberaLista(a);    /* anche i nodi di b */
   return 0;
}
\end{lstlisting}
\end{affianco}

\begin{figure}[H]
  \centering
  \begin{tikzpicture}[font=\Lato\scriptsize]
    \variabile[minimum width=0.7cm]{ha}{(0,0)}{a}{}
    \nodo{a1}{(0.8,0)}{1}
    \nodo{a2}{(2.7,0)}{2}
    \nodo{a3}{(4.6,0)}{3}
    \riferimento{ha}{-- (a1.west)}
    \collega{a1}{a2}
    \collega{a2}{a3}
    \variabile[minimum width=0.7cm]{hb}{(5.75,-1.3)}{b}{}
    \nodo{b1}{(6.55,-1.3)}{7}
    \nodo{b2}{(8.45,-1.3)}{8}
    \nodonull{b2}
    \riferimento{hb}{-- (b1.west)}
    \collega{b1}{b2}
    \draw[mrifnuovo] (a3-next.center) -| node[mnota, pos=0.75, right=2pt] {\code{p->next = b}} (b1-dato.north);
    \node[mpunto] at (a3-next.center) {};
  \end{tikzpicture}
  \caption{\code{unisci(a, b)}: il \code{next} dell'ultimo nodo di \code{a},
           che valeva \code{NULL}, punta al primo nodo di \code{b}.}
  \label{fig:unione}
\end{figure}

\begin{console}
$ gcc -Wall -o unione unione.c lista.c
$ ./unione
a: 1 -> 2 -> 3 -> NULL
b: 7 -> 8 -> NULL
a dopo unisci(a, b): 1 -> 2 -> 3 -> 7 -> 8 -> NULL
lunghezza: 5
\end{console}

Dopo l'unione i nodi di \code{b} fanno parte di \code{a}, e \code{b} punta
ancora al nodo 7, che ora sta in mezzo alla lista \code{a}. Per questo il
programma libera solo \code{a}: liberare anche \code{b} vorrebbe dire
liberare due volte gli stessi nodi (\textit{double free}), un errore grave
quanto usare un nodo già liberato.

\section{Liberare la lista}

\subsection{Un nodo alla volta}

Quando una lista non serve più, ogni suo nodo va restituito con \code{free}:
è un'eliminazione in testa (figura~\ref{fig:el-testa}) ripetuta finché la
lista non è vuota.

\begin{affianco}[0.55][c]
\begin{lstlisting}[style=c, name=lista.c]
void liberaLista(struct Node *head)
{
   struct Node *succ;

   while (head != NULL) {
      succ = head->next;    /* salva il successivo */
      free(head);
      head = succ;
   }
}
\end{lstlisting}
\risultato
Il punto delicato è lo stesso dell'eliminazione: dopo \code{free(head)} il
nodo non esiste più, e il suo campo \code{next} non si può leggere.
L'indirizzo del successivo si salva quindi \textbf{prima}, in \code{succ}.
\end{affianco}

\subsection{L'errore tipico}

Sembra più semplice scrivere il ciclo di scorrimento con un \code{free} nel
corpo, come in \code{liberaSbagliata}:

\begin{lstlisting}[style=c, name=libera_male.c]
#include <stdlib.h>
#include "lista.h"

void liberaSbagliata(struct Node *head)
{
   struct Node *p;

   for (p = head; p != NULL; p = p->next)
      free(p);             /* poi p = p->next legge un nodo già liberato */
}

int main(void)
{
   struct Node *head = NULL;

   head = inserisciCoda(head, 10);
   head = inserisciCoda(head, 20);
   liberaSbagliata(head);
   return 0;
}
\end{lstlisting}

Ma \code{p = p->next} viene eseguito \textbf{dopo} \code{free(p)}: legge il
campo \code{next} di un nodo che non esiste più. gcc 13 se ne accorge già in
compilazione, con \code{-Wall}:

\begin{console}
$ gcc -Wall -o libera_male libera_male.c lista.c
libera_male.c: In function (*@‘@*)liberaSbagliata(*@’@*):
libera_male.c:8:32: warning: pointer (*@‘@*)p(*@’@*) used after (*@‘@*)free(*@’@*) [-Wuse-after-free]
    8 |    for (p = head; p != NULL; p = p->next)
      |                              ~~^~~~~~~~~
libera_male.c:9:7: note: call to (*@‘@*)free(*@’@*) here
    9 |       free(p);             /* poi p = p->next legge un nodo già liberato */
      |       ^~~~~~~
$ ./libera_male
Segmentation fault
\end{console}

Che cosa succede in esecuzione dipende da che cosa c'è nella memoria
liberata: nella prova, \code{free} aveva già usato quei byte per le proprie
informazioni interne, il ciclo ha letto come \code{next} un indirizzo senza
senso, e il programma è stato fermato. Altre volte un errore del genere
passa inosservato, e per questo è insidioso. Con \code{-fsanitize=address}
(capitolo~\ref{chap:puntatori}) il rapporto è preciso (qui abbreviato, con
un segnaposto al posto del percorso della cartella):

\begin{console}
$ gcc -Wall -g -fsanitize=address -o libera_male libera_male.c lista.c
(*@\textit{\dots{} lo stesso avviso di prima \dots}@*)
$ ./libera_male
=================================================================
==4236==ERROR: AddressSanitizer: heap-use-after-free on address 0x502000000018 at pc 0x5645579ac2f4 bp 0x7ffe8d8b8390 sp 0x7ffe8d8b8380
READ of size 8 at 0x502000000018 thread T0
    #0 0x5645579ac2f3 in liberaSbagliata (*@\textit{[percorso del file]}@*)/libera_male.c:8
    #1 0x5645579ac354 in main (*@\textit{[percorso del file]}@*)/libera_male.c:18
(*@\textit{\dots}@*)
0x502000000018 is located 8 bytes inside of 16-byte region [0x502000000010,0x502000000020)
freed by thread T0 here:
    #0 0x7f8074efc4d8 in free ../../../../src/libsanitizer/asan/asan_malloc_linux.cpp:52
    #1 0x5645579ac2ce in liberaSbagliata (*@\textit{[percorso del file]}@*)/libera_male.c:9
(*@\textit{\dots}@*)
\end{console}

Il rapporto si legge così: alla riga 8 si è letto (\textit{READ}) un valore
di 8 byte, un puntatore, a 8 byte dall'inizio di un blocco di 16 byte, cioè
proprio il campo \code{next} di un nodo; e quel blocco era stato liberato
alla riga 9.

\section{La ricorsione}

\subsection{Una funzione che chiama se stessa}

Una funzione è \textbf{ricorsiva} quando, per risolvere un problema, chiama
se stessa su un problema più piccolo dello stesso tipo. Con le liste viene
naturale, perché una lista ha una struttura ricorsiva: una lista

\begin{itemize}
  \item o è \textbf{vuota} (\code{NULL});
  \item oppure è un \textbf{nodo seguito da una lista}: quella che comincia
        da \code{head->next}, che ha un nodo in meno.
\end{itemize}

Ne segue una definizione ricorsiva della lunghezza: la lunghezza della lista
vuota è 0; la lunghezza di un nodo seguito da una lista è 1 più la lunghezza
di quella lista. In C si scrive quasi parola per parola:

\begin{lstlisting}[style=c, name=ricorsione.c]
int lunghezzaRic(struct Node *head)
{
   if (head == NULL)               /* caso base: la lista vuota */
      return 0;
   return 1 + lunghezzaRic(head->next);
}
\end{lstlisting}

Ogni funzione ricorsiva ha due ingredienti:

\begin{itemize}
  \item il \textbf{caso base}, che si risolve senza altre chiamate (qui la
        lista vuota, la cui lunghezza è 0);
  \item il \textbf{passo ricorsivo}, che riduce il problema e chiama la
        funzione sul problema ridotto (qui il resto della lista, che ha un
        nodo in meno).
\end{itemize}

Poiché a ogni chiamata la lista si accorcia di un nodo, prima o poi si arriva
alla lista vuota, e le chiamate si fermano. Senza il caso base, arrivata in
fondo alla lista la funzione leggerebbe \code{head->next} con \code{head}
uguale a \code{NULL}, e il programma verrebbe fermato.

\subsection{Una traccia delle chiamate}

Che una funzione chiami se stessa non ha niente di speciale: ogni chiamata è
indipendente dalle altre, e ha la sua copia dei parametri e delle variabili
locali sulla pila delle chiamate (capitolo~\ref{chap:funzioni}). La versione
seguente di \code{lunghezzaRic} racconta che cosa succede, con un rientro
che cresce a ogni chiamata aperta: la variabile globale \code{livello} conta
le chiamate aperte, e \code{rientra} stampa quattro spazi per ognuna.

\begin{lstlisting}[style=c, name=traccia.c]
#include <stdio.h>
#include "lista.h"

int livello = 0;                   /* quante chiamate sono aperte */

void rientra(void)
{
   int i;

   for (i = 0; i < livello; i++)
      printf("    ");
}

int lunghezzaRic(struct Node *head)
{
   int n;

   rientra();
   if (head == NULL) {
      printf("lunghezzaRic(NULL): caso base, restituisce 0\n");
      return 0;
   }
   printf("lunghezzaRic(lista da %d): chiama lunghezzaRic sul resto\n",
          head->data);
   livello++;
   n = 1 + lunghezzaRic(head->next);
   livello--;
   rientra();
   printf("lunghezzaRic(lista da %d): restituisce %d\n", head->data, n);
   return n;
}

int main(void)
{
   struct Node *head = NULL;

   head = inserisciCoda(head, 10);
   head = inserisciCoda(head, 20);
   head = inserisciCoda(head, 30);
   printf("risultato: %d\n", lunghezzaRic(head));
   liberaLista(head);
   return 0;
}
\end{lstlisting}

\begin{console}
$ gcc -Wall -o traccia traccia.c lista.c
$ ./traccia
lunghezzaRic(lista da 10): chiama lunghezzaRic sul resto
    lunghezzaRic(lista da 20): chiama lunghezzaRic sul resto
        lunghezzaRic(lista da 30): chiama lunghezzaRic sul resto
            lunghezzaRic(NULL): caso base, restituisce 0
        lunghezzaRic(lista da 30): restituisce 1
    lunghezzaRic(lista da 20): restituisce 2
lunghezzaRic(lista da 10): restituisce 3
risultato: 3
\end{console}

\subsection{Che cosa succede}

\begin{itemize}
  \item La prima chiamata riceve la lista intera, che comincia da 10. Non è
        vuota, quindi deve calcolare il valore di
        \code{1 + lunghezzaRic(head->next)}, e per farlo chiama la funzione
        sulla lista che comincia da 20. Il
        calcolo della prima chiamata resta in sospeso: la chiamata è
        \textbf{aperta}, e la sua copia di \code{head} resta sulla pila.
  \item Lo stesso succede per la lista da 20 e per quella da 30: alla fine
        quattro chiamate sono aperte insieme, ognuna con il suo \code{head}
        (figura~\ref{fig:pila-ric}).
  \item La quarta chiamata riceve \code{NULL}: è il caso base, e restituisce
        0 senza chiamare nessuno.
  \item Ora le chiamate si chiudono a ritroso: quella sulla lista da 30
        completa \code{1 + 0} e restituisce 1; quella da 20 restituisce 2; la
        prima restituisce 3 a \code{main}.
\end{itemize}

\begin{figure}[H]
  \centering
  \begin{tikzpicture}[font=\Lato\scriptsize]
    \nodo{n1}{(0,0)}{10}
    \nodo{n2}{(2.4,0)}{20}
    \nodo{n3}{(4.8,0)}{30}
    \nodonull{n3}
    \collega{n1}{n2}
    \collega{n2}{n3}
    \foreach \k/\ix/\iy in {1/0.475/-1.45, 2/2.875/-2.65, 3/5.275/-3.85} {
      \node[mvar, minimum width=0.6cm] (h\k) at (\ix,\iy) {};
      \node[mnome, anchor=west] (h\k-n) at ([xshift=2pt]h\k.east) {head};
      \node[mscopolab, anchor=north west] (h\k-l) at ([xshift=-2pt, yshift=-3pt]h\k.south west) {chiamata \k};
      \begin{pgfonlayer}{liv2}
        \node[mscopo2, fit=(h\k)(h\k-n)(h\k-l)] (f\k) {};
      \end{pgfonlayer}
      \riferimento{h\k}{-- (n\k-dato.south)}
    }
    \node[mvar, minimum width=0.9cm, font=\FiraCodeNL\tiny] (h4) at (7.675,-5.05) {NULL};
    \node[mnome, anchor=west] (h4-n) at ([xshift=2pt]h4.east) {head};
    \node[mscopolab, anchor=north west] (h4-l) at ([xshift=-2pt, yshift=-3pt]h4.south west) {chiamata 4: caso base};
    \begin{pgfonlayer}{liv2}
      \node[mscopo2, fit=(h4)(h4-n)(h4-l)] (f4) {};
    \end{pgfonlayer}
    % i valori restituiti
    \draw[mcerca] (f4.west) -| node[mnota, pos=0.25, below] {restituisce 0} (f3.south);
    \draw[mcerca] (f3.west) -| node[mnota, pos=0.25, below] {restituisce 1} (f2.south);
    \draw[mcerca] (f2.west) -| node[mnota, pos=0.25, below] {restituisce 2} (f1.south);
    \draw[mcerca] (f1.west) -- ++(-0.8,0) node[mnota, anchor=east] {3 a \code{main}};
  \end{tikzpicture}
  \caption{Le quattro chiamate di \code{lunghezzaRic} aperte insieme, ognuna
           con la sua copia di \code{head}. Le frecce tratteggiate sono i
           valori restituiti, dall'ultima chiamata alla prima.}
  \label{fig:pila-ric}
\end{figure}

\subsection{Altri esempi}

Il programma \code{ricorsione.c} contiene, oltre a \code{lunghezzaRic}, tre
funzioni ricorsive sulle liste, e le usa su una lista di quattro nodi:

\begin{lstlisting}[style=c, name=ricorsione.c]
#include <stdio.h>
#include <stdlib.h>
#include "lista.h"

int lunghezzaRic(struct Node *head)
{
   if (head == NULL)               /* caso base: la lista vuota */
      return 0;
   return 1 + lunghezzaRic(head->next);
}

void stampaRovescia(struct Node *head)
{
   if (head == NULL)
      return;
   stampaRovescia(head->next);     /* prima il resto della lista */
   printf("%d ", head->data);      /* poi questo nodo */
}

struct Node *inserisciCodaRic(struct Node *head, struct Node *nodo)
{
   if (head == NULL)               /* in fondo: qui va il nodo */
      return nodo;
   head->next = inserisciCodaRic(head->next, nodo);
   return head;
}

void liberaListaRic(struct Node *head)
{
   if (head != NULL) {
      liberaListaRic(head->next);  /* prima il resto della lista */
      free(head);                  /* poi questo nodo */
   }
}

int main(void)
{
   struct Node *head = NULL;
   int i;

   for (i = 1; i <= 4; i++)
      head = inserisciCodaRic(head, nuovoNodo(i * 10));
   stampaLista(head);
   printf("lunghezza: %d\n", lunghezzaRic(head));
   printf("al contrario: ");
   stampaRovescia(head);
   printf("\n");
   liberaListaRic(head);
   return 0;
}
\end{lstlisting}

\begin{console}
$ gcc -Wall -o ricorsione ricorsione.c lista.c
$ ./ricorsione
10 -> 20 -> 30 -> 40 -> NULL
lunghezza: 4
al contrario: 40 30 20 10
\end{console}

\subsection{Come lavorano}

\begin{itemize}
  \item \code{stampaRovescia} stampa la lista al contrario senza mai tornare
        indietro: ogni chiamata \textbf{prima} fa stampare tutto il resto
        della lista, \textbf{poi} stampa il proprio nodo. Il primo a essere
        stampato è quindi il nodo della chiamata più profonda, l'ultimo. Con
        un ciclo sarebbe molto più complicato, perché da un nodo non si può
        tornare al precedente: qui sono le chiamate aperte sulla pila a
        ricordare la strada.
  \item \code{inserisciCodaRic} inserisce in coda un nodo già creato: se la
        lista è vuota, il nodo è tutta la lista; altrimenti lo inserisce in
        coda al resto, e aggancia il risultato a \code{head->next}. È la
        funzione \code{inserisciLista} del progetto delle lezioni.
  \item \code{liberaListaRic} libera prima il resto della lista e poi il
        proprio nodo: così il campo \code{next} viene letto quando il nodo
        esiste ancora.
\end{itemize}

\subsection{Ricorsione o ciclo?}

Tutto ciò che si fa con la ricorsione si può fare con un ciclo, e viceversa.
La ricorsione rende spesso il codice più corto e più vicino alla definizione
del problema; ma ogni chiamata aperta occupa spazio sulla pila, e con una
lista molto lunga le chiamate aperte tutte insieme possono esaurirla. Il
programma \code{profonda.c} crea una lista di \code{n} nodi e ne calcola la
lunghezza in tutti e due i modi:

\begin{affianco}[0.55][c]
\begin{lstlisting}[style=c, name=profonda.c]
#include <stdio.h>
#include "lista.h"

int lunghezzaRic(struct Node *head)
{
   if (head == NULL)
      return 0;
   return 1 + lunghezzaRic(head->next);
}

int main(void)
{
   struct Node *head = NULL;
   int i, n;

   if (scanf("%d", &n) != 1)
      return 1;
   for (i = 0; i < n; i++)
      head = inserisciTesta(head, i);
   printf("con un ciclo: %d\n", lunghezza(head));
   printf("con la ricorsione: %d\n", lunghezzaRic(head));
   liberaLista(head);
   return 0;
}
\end{lstlisting}
\risultato
\begin{console}
$ echo 100000 | ./profonda
con un ciclo: 100000
con la ricorsione: 100000
$ echo 1000000 | ./profonda
con un ciclo: 1000000
Segmentation fault
\end{console}
\wtnota{\code{echo} scrive il numero, e la barra \code{|} lo manda
all'input del programma, come se lo si digitasse.}
\end{affianco}

Con un milione di nodi servono un milione di chiamate aperte insieme, e gli
8 MB che Linux assegna di solito alla pila (capitolo~\ref{chap:stringhe})
non bastano: il programma viene fermato. Il ciclo, invece, usa sempre le
stesse due variabili. Per questo, sulle liste che possono diventare molto
lunghe, si preferiscono i cicli.

\section{I file}

\subsection{Aprire e chiudere un file}

Finora i programmi hanno letto dalla tastiera e scritto sullo schermo, e
hanno usato dei file solo grazie alla redirezione \code{<} e \code{>} del
terminale (capitolo~\ref{chap:primi-programmi}). Il progetto delle lezioni,
invece, apre i file direttamente. Gli strumenti sono in \code{stdio.h}:

\begin{itemize}
  \item \code{FILE *fp;} dichiara un puntatore a \code{FILE}, una struttura
        definita in \code{stdio.h} che descrive un file aperto: quale file
        è, a che punto della lettura o della scrittura si è arrivati, il
        buffer, gli eventuali errori. Non serve sapere com'è fatta: la si usa
        solo attraverso le funzioni della libreria.
  \item \code{fp = fopen("dati.txt", "r");} apre il file \code{dati.txt} in
        \textbf{lettura} (\code{"r"}, \textit{read}). Con \code{"w"}
        (\textit{write}) lo apre in \textbf{scrittura}, creandolo se non
        esiste e \textbf{svuotandolo} se esiste; con \code{"a"}
        (\textit{append}) scrive in fondo a ciò che contiene già. Se il file
        non si può aprire (per esempio non esiste, o manca il permesso),
        \code{fopen} restituisce \code{NULL}, e questo va sempre controllato.
  \item \code{fclose(fp);} chiude il file quando non serve più: scrive su
        disco ciò che è rimasto nel buffer e libera la struttura \code{FILE}.
\end{itemize}

Un nome come \code{"dati.txt"}, senza percorso, indica un file nella cartella
in cui ci si trova quando si lancia il programma: in VS Code, quella del
terminale (capitolo~\ref{chap:ambiente}).

\subsection{Leggere e scrivere}

\code{fscanf} e \code{fprintf} funzionano esattamente come \code{scanf} e
\code{printf}, con un primo argomento in più: il file da cui leggere o su cui
scrivere. Anche il valore restituito da \code{fscanf} è quello di
\code{scanf}: il numero di conversioni riuscite, oppure \code{EOF} se il file
è finito prima della prima conversione. Del resto \code{printf(...)} equivale
a \code{fprintf(stdout, ...)}, e \code{scanf(...)} a
\code{fscanf(stdin, ...)}: \code{stdin} e \code{stdout} sono due
\code{FILE *} già aperti all'avvio del programma, collegati alla tastiera e
allo schermo (o ai file indicati con \code{<} e \code{>}). Il programma
seguente legge dei numeri da \code{numeri.txt} e ne scrive i quadrati in
\code{quadrati.txt}:

\begin{affianco}[0.55][c]
\begin{lstlisting}[style=c, name=quadrati.c]
#include <stdio.h>

int main(void)
{
   FILE *in, *out;
   int x, letti = 0;

   in = fopen("numeri.txt", "r");
   if (in == NULL) {
      printf("non riesco ad aprire numeri.txt\n");
      return 1;
   }
   out = fopen("quadrati.txt", "w");
   if (out == NULL) {
      printf("non riesco a creare quadrati.txt\n");
      fclose(in);
      return 1;
   }
   while (fscanf(in, "%d", &x) == 1) {
      fprintf(out, "%d al quadrato fa %d\n", x, x * x);
      letti++;
   }
   fclose(in);
   fclose(out);
   printf("letti %d numeri\n", letti);
   return 0;
}
\end{lstlisting}
\risultato
\begin{console}
$ cat numeri.txt
3 7
12
$ gcc -Wall -o quadrati quadrati.c
$ ./quadrati
letti 3 numeri
$ cat quadrati.txt
3 al quadrato fa 9
7 al quadrato fa 49
12 al quadrato fa 144
$ rm numeri.txt
$ ./quadrati
non riesco ad aprire numeri.txt
\end{console}
\end{affianco}

\begin{itemize}
  \item Il ciclo \code{while (fscanf(in, "\%d", \&x) == 1)} legge un numero
        alla volta, finché la lettura riesce. \code{fscanf} salta spazi e a
        capo, quindi non importa come i numeri siano disposti nelle righe.
        Alla fine del file restituisce \code{EOF}, e il ciclo termina.
  \item Il comando \code{cat} stampa il contenuto di un file, \code{rm} lo
        cancella. Senza \code{numeri.txt}, \code{fopen} restituisce
        \code{NULL}, e il programma lo dice invece di usare un puntatore
        nullo.
  \item Se la creazione di \code{quadrati.txt} fallisce, il programma chiude
        \code{numeri.txt} prima di uscire: ogni file aperto va chiuso.
\end{itemize}

\section{Il progetto delle lezioni}

\subsection{Che cosa fa}

Torniamo al progetto del capitolo~\ref{chap:make}, che ora abbiamo tutti gli
strumenti per capire. Il file \code{dati.txt} contiene gli esami sostenuti da
alcuni studenti, uno per riga: nome, cognome, matricola, esame e voto (in
fondo al file ci sono anche tre righe vuote, che la lettura salta).

\begin{affianco}[0.55][c]
\begin{lstlisting}[style=plain, name=dati.txt]
Luca      Verdi     364     AnalisiI            16
Enzo      Blu       367     analisiI            14
Mario     Rossi     272     programmazioneI     27
Luca      Verdi     364     basididati1         16
Enzo      Blu       367     basididati1         14
Anna      Bianco    471     algebra             23
Mario     Rossi     272     algebra             30
Luca      Verdi     364     geometria           16
Enzo      Blu       367     algebra             14
\end{lstlisting}
\risultato
Il programma:
\begin{passi}
  \item legge il file e costruisce una lista, con un nodo per ogni esame
        nell'ordine del file, e la stampa;
  \item toglie dalla lista gli esami con voto minore di 18, e la stampa;
  \item scrive nel file \code{dati2.txt} gli esami rimasti, raggruppati per
        studente, e stampa di nuovo la lista;
  \item libera la memoria.
\end{passi}
\end{affianco}

\subsection{\code{studente.h}}

\begin{affianco}[0.55][c]
\codicepiccolo
\begin{lstlisting}[style=h, name=studente.h]
#include <stdio.h>
#include <stdlib.h>
#include <string.h>

#define STRING_LENGTH  512

struct studente{
  char nome[STRING_LENGTH];
  char cognome[STRING_LENGTH];
  int matricola;
  char esame[STRING_LENGTH];
  int voto;
  struct studente *next;
};

void gestioneErrore();
struct studente *leggiFile(FILE *fp, struct studente *lista);
void scriviLista(struct studente *lista);
struct studente *rimuoviSoglia(struct studente *lista, int soglia);
void scriviFile(FILE *fp, struct studente *lista);
void freeLista(struct studente *lista);
\end{lstlisting}
\risultato
\begin{itemize}
  \item Il nodo della lista è una \code{struct studente}: i dati di un esame
        (cinque campi) e il puntatore \code{next}, come in
        \code{struct Node}.
  \item Nome, cognome ed esame sono array di \code{STRING\_LENGTH}
        caratteri dentro il nodo, che occupa così 1552 byte
        (\code{sizeof}), anche se i nomi sono corti.
  \item L'header include \code{stdio.h} perché i prototipi usano il tipo
        \code{FILE}. Sono i prototipi delle funzioni chiamate da
        \code{main.c}; le altre funzioni di \code{funzioni.c} sono
        \code{static} (capitolo~\ref{chap:funzioni}), visibili solo dentro
        quel file.
  \item \code{gestioneErrore()} ha le parentesi vuote, che fino al C17
        vogliono dire \textit{parametri non specificati}: la forma precisa
        è \code{gestioneErrore(void)}.
\end{itemize}
\end{affianco}

\subsection{Creare, inserire e togliere un nodo}

Le prime funzioni di \code{funzioni.c} lavorano su un nodo alla volta:

\begin{lstlisting}[style=c, name=funzioni.c]
#include <stdio.h>
#include <stdlib.h>
#include <string.h>
#include "studente.h"

static struct studente *nuovoNodo(char *nome,char *cognome, int matricola, char *esame, int voto)
{
  struct studente *nodo;

  nodo = (struct studente *) malloc(sizeof(struct studente));
  if (nodo == NULL) return NULL;

  strcpy(nodo->nome, nome);
  strcpy(nodo->cognome, cognome);
  strcpy(nodo->esame, esame);
  nodo->matricola = matricola;
  nodo->voto = voto;
  nodo->next = NULL;

  return nodo;
}

static void freeNodo(struct studente *nodo)
{
  free((void *) nodo);
}

static struct studente *inserisciLista(struct studente *nodo, struct studente *lista)
{
  if (lista==NULL) return nodo;
  lista->next = inserisciLista(nodo, lista->next);
  return lista;
}

static struct studente *rimuoviNodo(struct studente *lista, struct studente *nodo)
{
  struct studente *tmp;


  if (lista == nodo) {
    tmp = lista->next;
    freeNodo(lista);
    return tmp;
  }

  lista->next = rimuoviNodo(lista->next,nodo);
}
\end{lstlisting}

\begin{itemize}
  \item \code{nuovoNodo} è come quella del modulo \code{lista.c}, ma i dati
        sono cinque, e le stringhe vanno \textbf{copiate} nel nodo con
        \code{strcpy} (capitolo~\ref{chap:stringhe}): gli array del
        chiamante verranno riusati per la riga successiva del file. Il cast
        \code{(struct studente *)} davanti a \code{malloc} non è necessario
        in C, come quello a \code{void *} in \code{freeNodo}.
  \item \code{inserisciLista} inserisce \code{nodo} in coda alla lista,
        ricorsivamente: è \code{inserisciCodaRic}, con gli argomenti in
        ordine inverso.
  \item \code{rimuoviNodo} toglie \code{nodo} dalla lista, ricorsivamente:
        se il primo nodo è quello da togliere, lo libera e restituisce il
        resto; altrimenti lo toglie dal resto, e aggancia il risultato a
        \code{lista->next}.
\end{itemize}

Compilando con \code{-Wall}, gcc segnala un problema proprio in
\code{rimuoviNodo}:

\begin{console}
$ gcc -Wall -c main.c funzioni.c
funzioni.c: In function (*@‘@*)rimuoviNodo(*@’@*):
funzioni.c:47:1: warning: control reaches end of non-void function [-Wreturn-type]
   47 | }
      | ^
\end{console}

\warning{Una svista nel codice delle lezioni}{
  Nel ramo ricorsivo di \code{rimuoviNodo} manca il \code{return}: la
  funzione aggancia il risultato a \code{lista->next}, poi arriva alla
  graffa finale senza restituire niente, e il chiamante riceve un valore
  indefinito (in pratica, ciò che si trova per caso nel registro del
  processore usato per i risultati: capitolo~\ref{chap:funzioni}). Nel
  programma l'errore non si manifesta, perché \code{rimuoviNodo} viene
  chiamata solo da \code{rimuoviSoglia}, sempre con \code{nodo} uguale al
  primo nodo della lista: il ramo ricorsivo non viene mai eseguito. Ma
  basterebbe usare la funzione in un altro modo per perdere la lista. La
  correzione è una riga, \code{return lista;} prima della graffa finale:
  con questa, gcc non segnala più niente, e l'output del programma è lo
  stesso.}

\subsection{Le funzioni di servizio}

\begin{lstlisting}[style=c, name=funzioni.c, firstnumber=50]
static int lengthLista(struct studente *lista)
{
  int n = 0;

  while(lista!=NULL) {
    n++;
    lista = lista->next;
  }

  return n;
}


void gestioneErrore()
{
  printf("Errore\n");

  exit(0);
}


void freeLista(struct studente *lista)
{
  if (lista != NULL) {
    freeLista(lista->next);
    freeNodo(lista);
  }
}


void scriviLista(struct studente *lista)
{
  while (lista!=NULL) {
    printf("%s %s %d %s %d\n",lista->nome,lista->cognome,lista->matricola,lista->esame,lista->voto);
    lista = lista->next;
  }
  printf("\n");
}
\end{lstlisting}

\begin{itemize}
  \item \code{lengthLista} conta i nodi con un ciclo, come
        \code{lunghezza}.
  \item \code{gestioneErrore} stampa un messaggio e termina il programma con
        \code{exit} (capitolo~\ref{chap:primi-programmi}), da qualunque
        punto venga chiamata. Termina però con \code{exit(0)}, che alla
        shell comunica un successo: per un errore sarebbe più corretto
        \code{exit(1)}.
  \item \code{freeLista} libera la lista ricorsivamente, prima il resto e poi
        il nodo, come \code{liberaListaRic}.
  \item \code{scriviLista} stampa un nodo per riga, e una riga vuota alla
        fine.
\end{itemize}

\subsection{Leggere il file: \code{leggiFile}}

\begingroup
\codicepiccolo
% numeri di riga a tre cifre: una colonna dei numeri un po' piu' larga
\lstset{xleftmargin=19pt, framexleftmargin=15pt}
\begin{lstlisting}[style=c, name=funzioni.c, firstnumber=89]
struct studente *leggiFile(FILE *fp, struct studente *lista)
{
  char nome[STRING_LENGTH], cognome[STRING_LENGTH], esame[STRING_LENGTH];
  int matricola, voto;
  struct studente *nodo;

  while(fscanf(fp,"%s%s%d%s%d",nome,cognome,&matricola,esame,&voto)==5){
    nodo = nuovoNodo(nome,cognome,matricola,esame,voto);
    if (nodo==NULL) gestioneErrore();
    lista = inserisciLista(nodo, lista);
  }
  return lista;
}
\end{lstlisting}
\endgroup

\code{leggiFile} riceve un file già aperto e la lista a cui aggiungere i
nodi, all'inizio vuota. Il ciclo legge una riga alla volta: \code{fscanf}
legge nome, cognome, matricola, esame e voto, e restituisce 5 se tutte le
conversioni sono riuscite. Per ogni riga si crea un nodo e lo si inserisce in
coda, così la lista ha lo stesso ordine del file; se manca la memoria, si
chiama \code{gestioneErrore}. Quando il file finisce, \code{fscanf}
restituisce \code{EOF}, il ciclo si ferma, e la funzione restituisce la
lista. Il \code{\%s} legge una parola, cioè si ferma al primo spazio: per
questo nel file gli esami sono scritti senza spazi, come
\code{programmazioneI}.

\subsection{Filtrare: \code{rimuoviSoglia}}

\begingroup
\codicepiccolo
% numeri di riga a tre cifre: una colonna dei numeri un po' piu' larga
\lstset{xleftmargin=19pt, framexleftmargin=15pt}
\begin{lstlisting}[style=c, name=funzioni.c, firstnumber=104]
struct studente *rimuoviSoglia(struct studente *lista, int soglia)
{
  if (lista==NULL) return lista;
  if (lista->voto<soglia) {
    lista = rimuoviNodo(lista, lista);
    lista = rimuoviSoglia(lista, soglia);
  } else lista->next = rimuoviSoglia(lista->next, soglia);

  return lista;
}
\end{lstlisting}
\endgroup

\code{rimuoviSoglia} toglie dalla lista tutti i nodi con voto minore di
\code{soglia}, ricorsivamente:

\begin{itemize}
  \item se la lista è vuota, non c'è niente da togliere;
  \item se il primo nodo ha un voto insufficiente, lo toglie
        (\code{rimuoviNodo(lista, lista)} lo libera e restituisce il resto),
        e poi filtra il resto, che è diventato la lista;
  \item altrimenti il primo nodo resta, e la funzione filtra il resto,
        agganciando il risultato a \code{lista->next}.
\end{itemize}

\begin{figure}[H]
  \centering
  \begin{tikzpicture}[font=\Lato\scriptsize, stretto/.style={minimum width=0.62cm}]
    \node[mnota, anchor=west, font=\Lato\scriptsize\bfseries] at (-1.0,0.8)
      {la lista letta da \code{dati.txt}};
    \variabile[minimum width=0.6cm]{h}{(0,0)}{lista}{}
    \nodo[stretto, mlvia]{s1}{(0.75,0)}{16}
    \nodo[stretto, mlvia]{s2}{(2.3,0)}{14}
    \nodo[stretto]{s3}{(3.85,0)}{27}
    \nodo[stretto, mlvia]{s4}{(5.4,0)}{16}
    \nodo[stretto, mlvia]{s5}{(6.95,0)}{14}
    \nodo[stretto]{s6}{(8.5,0)}{23}
    \nodo[stretto]{s7}{(10.05,0)}{30}
    \nodo[stretto, mlvia]{s8}{(11.6,0)}{16}
    \nodo[stretto, mlvia]{s9}{(13.15,0)}{14}
    \nodonull{s9}
    \riferimento{h}{-- (s1.west)}
    \foreach \k/\j in {1/2, 2/3, 3/4, 4/5, 5/6, 6/7, 7/8, 8/9} {\collega{s\k}{s\j}}
    \foreach \k/\nome in {1/Luca, 2/Enzo, 3/Mario, 4/Luca, 5/Enzo, 6/Anna, 7/Mario, 8/Luca, 9/Enzo}
      \node[mindir, anchor=south] at ([yshift=1pt]s\k.north) {\nome};
    \node[mnota, anchor=west, font=\Lato\scriptsize\bfseries] at (-1.0,-0.95)
      {dopo \code{lista = rimuoviSoglia(lista, 18)}};
    \variabile[minimum width=0.6cm]{h2}{(0,-1.75)}{lista}{}
    \nodo[stretto]{r1}{(0.75,-1.75)}{27}
    \nodo[stretto]{r2}{(2.3,-1.75)}{23}
    \nodo[stretto]{r3}{(3.85,-1.75)}{30}
    \nodonull{r3}
    \riferimento{h2}{-- (r1.west)}
    \collega{r1}{r2}
    \collega{r2}{r3}
    \foreach \k/\nome in {1/Mario, 2/Anna, 3/Mario}
      \node[mindir, anchor=south] at ([yshift=1pt]r\k.north) {\nome};
  \end{tikzpicture}
  \caption{I voti dei nove esami letti da \code{dati.txt}, con il nome dello
           studente sopra ogni nodo: \code{rimuoviSoglia} libera i sei nodi
           con voto minore di 18, evidenziati, e collega fra loro quelli che
           restano.}
  \label{fig:soglia}
\end{figure}

\subsection{Raggruppare per studente: \code{scriviFile}}

\begingroup
\codicepiccolo
% numeri di riga a tre cifre: una colonna dei numeri un po' piu' larga
\lstset{xleftmargin=19pt, framexleftmargin=15pt}
\begin{lstlisting}[style=c, name=funzioni.c, firstnumber=116]
void scriviFile(FILE *fp, struct studente *lista)
{
  int *buffer, n, ind1, ind2;
  struct studente *ptr1, *ptr2;

  n = lengthLista(lista);
  buffer = (int *) calloc(n,sizeof(int));

  for (ptr1=lista, ind1=0; ptr1!=NULL; ptr1=ptr1->next, ind1++) {
    if (!buffer[ind1]) {
      buffer[ind1] = 1;
      fprintf(fp,"%s %s %d %s %d\n",ptr1->nome, ptr1->cognome, ptr1->matricola, ptr1->esame, ptr1->voto);
    } else continue;
    for (ptr2=ptr1->next, ind2=ind1+1; ptr2!=NULL; ptr2=ptr2->next, ind2++) {
      if (ptr2->matricola==ptr1->matricola) {
        buffer[ind2] = 1;
        fprintf(fp,"%s %s %d %s %d\n",ptr2->nome, ptr2->cognome, ptr2->matricola, ptr2->esame, ptr2->voto);
      }
    }
  }

  free((void *) buffer);
}
\end{lstlisting}
\endgroup

\code{scriviFile} scrive i nodi nel file \textbf{raggruppati per matricola}:
tutti gli esami di uno studente uno dopo l'altro, nell'ordine in cui
compaiono nella lista.

\begin{itemize}
  \item \code{buffer} è un array di \code{n} interi, uno per nodo, allocato
        con \code{calloc} (capitolo~\ref{chap:puntatori}), quindi tutti a
        zero: \code{buffer[i]} diventa 1 quando il nodo in posizione
        \code{i} è stato scritto. Una lista non ha indici, quindi i due
        \code{for} fanno avanzare insieme due variabili, separate dalla
        virgola: il puntatore al nodo e la sua posizione (\code{ptr1} e
        \code{ind1}, \code{ptr2} e \code{ind2}).
  \item Il \code{for} esterno prende i nodi uno alla volta. Se il nodo è già
        stato scritto, passa al successivo (\code{continue}); altrimenti lo
        scrive, e il \code{for} interno scorre i nodi \textbf{seguenti},
        scrivendo (e segnando) quelli con la stessa matricola.
\end{itemize}

Seguiamo a mano la scrittura della lista filtrata, che ha tre nodi:

\begin{center}
\renewcommand{\arraystretch}{1.2}
\begin{tabular}{@{}l c p{8.4cm}@{}}
  \toprule
  \textbf{\code{ptr1} (\code{ind1})} & \textbf{\code{buffer}} & \textbf{che cosa succede} \\
  \midrule
  Mario Rossi, programmazioneI (0) & \code{0 0 0} & lo scrive; il \code{for}
    interno trova in posizione 2 l'altro esame della matricola 272 e scrive
    anche quello: \code{buffer} diventa \code{1 0 1} \\
  Anna Bianco, algebra (1) & \code{1 0 1} & lo scrive; nessun nodo
    successivo ha la matricola 471: \code{buffer} diventa \code{1 1 1} \\
  Mario Rossi, algebra (2) & \code{1 1 1} & è già stato scritto:
    \code{continue} \\
  \bottomrule
\end{tabular}
\end{center}

\subsection{\code{main.c} e l'esecuzione}

\begin{affianco}[0.5][c]
\begin{lstlisting}[style=c, name=main.c]
#include <stdio.h>
#include <stdlib.h>
#include <string.h>
#include "studente.h"

int main()
{
  struct studente *lista=NULL;
  FILE *fp;

  fp = fopen("dati.txt","r");
  if (fp==NULL) gestioneErrore();
  lista = leggiFile(fp, lista);
  fclose(fp);
  scriviLista(lista);

  lista = rimuoviSoglia(lista, 18);
  scriviLista(lista);

  fp = fopen("dati2.txt","w");
  if (fp==NULL) gestioneErrore();
  scriviFile(fp, lista);
  fclose(fp);

  scriviLista(lista);

  freeLista(lista);
}
\end{lstlisting}
\risultato
\begin{console}
$ gcc -Wall -c main.c funzioni.c
(*@\textit{\dots{} l'avviso su rimuoviNodo \dots}@*)
$ gcc -o esercizio main.o funzioni.o
$ ./esercizio
Luca Verdi 364 AnalisiI 16
Enzo Blu 367 analisiI 14
Mario Rossi 272 programmazioneI 27
Luca Verdi 364 basididati1 16
Enzo Blu 367 basididati1 14
Anna Bianco 471 algebra 23
Mario Rossi 272 algebra 30
Luca Verdi 364 geometria 16
Enzo Blu 367 algebra 14

Mario Rossi 272 programmazioneI 27
Anna Bianco 471 algebra 23
Mario Rossi 272 algebra 30

Mario Rossi 272 programmazioneI 27
Anna Bianco 471 algebra 23
Mario Rossi 272 algebra 30

$ cat dati2.txt
Mario Rossi 272 programmazioneI 27
Mario Rossi 272 algebra 30
Anna Bianco 471 algebra 23
\end{console}
\end{affianco}

\begin{itemize}
  \item \code{main} apre \code{dati.txt} in lettura (se non ci riesce chiama
        \code{gestioneErrore}), costruisce la lista con \code{leggiFile},
        chiude il file e stampa i nove esami.
  \item \code{rimuoviSoglia(lista, 18)} lascia i tre esami superati, e la
        lista viene stampata di nuovo.
  \item \code{dati2.txt} viene aperto in scrittura e riempito da
        \code{scriviFile}: i due esami di Mario Rossi sono uno dopo l'altro,
        come voleva il raggruppamento. La terza stampa è uguale alla
        seconda, perché \code{scriviFile} non modifica la lista.
  \item Alla fine \code{freeLista} libera i nodi. \code{main} non ha un
        \code{return}: dal C99, arrivare alla graffa finale di \code{main}
        equivale a \code{return 0}.
\end{itemize}

\subsection{Nessuna perdita di memoria}

\textbf{valgrind} è un programma che ne esegue un altro tenendo traccia di
ogni allocazione e di ogni \code{free}, e alla fine dice se qualche blocco è
rimasto allocato: una \textbf{perdita di memoria} (\textit{memory leak}). Sul
progetto (qui l'output del programma, che si mescola a quello di valgrind, è
omesso):

\begin{console}
$ valgrind ./esercizio
==451== Memcheck, a memory error detector
==451== Copyright (C) 2002-2022, and GNU GPL'd, by Julian Seward et al.
==451== Using Valgrind-3.22.0 and LibVEX; rerun with -h for copyright info
==451== Command: ./esercizio
==451==
(*@\textit{\dots{} l'output del programma \dots}@*)
==451==
==451== HEAP SUMMARY:
==451==     in use at exit: 0 bytes in 0 blocks
==451==   total heap usage: 15 allocs, 15 frees, 27,212 bytes allocated
==451==
==451== All heap blocks were freed -- no leaks are possible
==451==
==451== For lists of detected and suppressed errors, rerun with: -s
==451== ERROR SUMMARY: 0 errors from 0 contexts (suppressed: 0 from 0)
\end{console}

Le 15 allocazioni comprendono i 9 nodi, il \code{buffer} di
\code{scriviFile} e quelle che la libreria fa per gestire i file e l'output:
sono state liberate tutte. Proviamo a commentare la chiamata
\code{freeLista(lista);} in \code{main.c}, e a compilare con \code{-g}
(che aggiunge i numeri di riga, utili a strumenti come questo); con
\code{--leak-check=full}, valgrind dice anche dove sono stati allocati i
blocchi persi:

\begin{console}
$ gcc -g -o esercizio main.c funzioni.c
$ valgrind --leak-check=full ./esercizio
(*@\textit{\dots}@*)
==442== HEAP SUMMARY:
==442==     in use at exit: 4,656 bytes in 3 blocks
==442==   total heap usage: 15 allocs, 12 frees, 27,212 bytes allocated
==442==
==442== 4,656 (1,552 direct, 3,104 indirect) bytes in 1 blocks are definitely lost in loss record 2 of 2
==442==    at 0x4846828: malloc (in /usr/libexec/valgrind/vgpreload_memcheck-amd64-linux.so)
==442==    by 0x1093E3: nuovoNodo (funzioni.c:10)
==442==    by 0x1096A9: leggiFile (funzioni.c:96)
==442==    by 0x10931D: main (main.c:13)
==442==
==442== LEAK SUMMARY:
==442==    definitely lost: 1,552 bytes in 1 blocks
==442==    indirectly lost: 3,104 bytes in 2 blocks
(*@\textit{\dots}@*)
\end{console}

I tre nodi rimasti dopo il filtro non sono stati liberati: 1552 byte
ciascuno, come diceva \code{sizeof}. Il primo è \textit{definitely lost}
(perso di sicuro): nessun puntatore lo raggiunge più, perché la variabile
\code{lista} di \code{main} non esiste più. Gli altri due sono
\textit{indirectly lost} (persi indirettamente): sono raggiungibili solo
attraverso il primo, che è perso. È la struttura della lista vista da
valgrind.

\info{Installare valgrind}{
  Su Ubuntu e su WSL valgrind si installa con
  \code{sudo apt install valgrind}, e si usa scrivendo \code{valgrind} prima
  del comando che lancia il programma. Su macOS valgrind non funziona: per
  gli accessi sbagliati alla memoria c'è \code{-fsanitize=address}, e per le
  perdite il comando \code{leaks} (\code{leaks --atExit -- ./esercizio}).}

\insintesi{
  Una \textbf{lista collegata} è una catena di nodi allocati nello heap: ogni
  nodo contiene i dati e l'indirizzo del successivo, l'ultimo ha \code{next}
  uguale a \code{NULL}, e la lista è individuata dal puntatore al primo nodo
  (\code{NULL} se è vuota). Rispetto a un array, la dimensione è flessibile e
  inserire o togliere un nodo costa il cambio di un paio di puntatori, ma
  l'accesso è sequenziale. Quasi ogni operazione segue lo stesso schema: un
  puntatore di appoggio scorre la lista e si ferma sul nodo giusto, spesso il
  \textbf{precedente} di quello su cui si opera; l'ordine delle assegnazioni
  conta, e un nodo si libera solo dopo averne letto il \code{next}. Le
  funzioni che possono cambiare il primo nodo restituiscono la nuova testa.
  Una lista è anche una struttura \textbf{ricorsiva} (vuota, oppure un nodo
  seguito da una lista), e molte operazioni si scrivono con funzioni che
  chiamano se stesse, con un caso base e un passo ricorsivo; ogni chiamata
  aperta, però, occupa spazio sulla pila. I file si aprono con \code{fopen},
  si usano con \code{fscanf} e \code{fprintf} e si chiudono con
  \code{fclose}. valgrind e \code{-fsanitize=address} aiutano a scoprire le
  perdite di memoria e gli accessi alla memoria già liberata.}

\section{Riferimenti}

\begin{itemize}
  \item \textbf{A. Kelley, I. Pohl}, \textit{A Book on C}: capitolo~10
        (strutture e liste), capitolo~11 (input/output) e, per la
        ricorsione, capitolo~5 (funzioni).
  \item \textbf{fopen} --- cppreference.com. \\
        \urlx{https://en.cppreference.com/w/c/io/fopen}
  \item \textbf{Valgrind Quick Start Guide} --- valgrind.org. \\
        \urlx{https://valgrind.org/docs/manual/quick-start.html}
\end{itemize}


\end{document}
